% Options for packages loaded elsewhere
% Options for packages loaded elsewhere
\PassOptionsToPackage{unicode}{hyperref}
\PassOptionsToPackage{hyphens}{url}
\PassOptionsToPackage{dvipsnames,svgnames,x11names}{xcolor}
%
\documentclass[
  spanish,
  a4paper,
]{scrbook}
\usepackage{xcolor}
\usepackage[top=25mm,bottom=25mm,inner=28mm,outer=22mm]{geometry}
\usepackage{amsmath,amssymb}
\setcounter{secnumdepth}{5}
\usepackage{iftex}
\ifPDFTeX
  \usepackage[T1]{fontenc}
  \usepackage[utf8]{inputenc}
  \usepackage{textcomp} % provide euro and other symbols
\else % if luatex or xetex
  \usepackage{unicode-math} % this also loads fontspec
  \defaultfontfeatures{Scale=MatchLowercase}
  \defaultfontfeatures[\rmfamily]{Ligatures=TeX,Scale=1}
\fi
\usepackage{lmodern}
\ifPDFTeX\else
  % xetex/luatex font selection
\fi
% Use upquote if available, for straight quotes in verbatim environments
\IfFileExists{upquote.sty}{\usepackage{upquote}}{}
\IfFileExists{microtype.sty}{% use microtype if available
  \usepackage[]{microtype}
  \UseMicrotypeSet[protrusion]{basicmath} % disable protrusion for tt fonts
}{}
\makeatletter
\@ifundefined{KOMAClassName}{% if non-KOMA class
  \IfFileExists{parskip.sty}{%
    \usepackage{parskip}
  }{% else
    \setlength{\parindent}{0pt}
    \setlength{\parskip}{6pt plus 2pt minus 1pt}}
}{% if KOMA class
  \KOMAoptions{parskip=half}}
\makeatother
% Make \paragraph and \subparagraph free-standing
\makeatletter
\ifx\paragraph\undefined\else
  \let\oldparagraph\paragraph
  \renewcommand{\paragraph}{
    \@ifstar
      \xxxParagraphStar
      \xxxParagraphNoStar
  }
  \newcommand{\xxxParagraphStar}[1]{\oldparagraph*{#1}\mbox{}}
  \newcommand{\xxxParagraphNoStar}[1]{\oldparagraph{#1}\mbox{}}
\fi
\ifx\subparagraph\undefined\else
  \let\oldsubparagraph\subparagraph
  \renewcommand{\subparagraph}{
    \@ifstar
      \xxxSubParagraphStar
      \xxxSubParagraphNoStar
  }
  \newcommand{\xxxSubParagraphStar}[1]{\oldsubparagraph*{#1}\mbox{}}
  \newcommand{\xxxSubParagraphNoStar}[1]{\oldsubparagraph{#1}\mbox{}}
\fi
\makeatother


\usepackage{longtable,booktabs,array}
\newcounter{none} % for unnumbered tables
\usepackage{calc} % for calculating minipage widths
% Correct order of tables after \paragraph or \subparagraph
\usepackage{etoolbox}
\makeatletter
\patchcmd\longtable{\par}{\if@noskipsec\mbox{}\fi\par}{}{}
\makeatother
% Allow footnotes in longtable head/foot
\IfFileExists{footnotehyper.sty}{\usepackage{footnotehyper}}{\usepackage{footnote}}
\makesavenoteenv{longtable}
\usepackage{graphicx}
\makeatletter
\newsavebox\pandoc@box
\newcommand*\pandocbounded[1]{% scales image to fit in text height/width
  \sbox\pandoc@box{#1}%
  \Gscale@div\@tempa{\textheight}{\dimexpr\ht\pandoc@box+\dp\pandoc@box\relax}%
  \Gscale@div\@tempb{\linewidth}{\wd\pandoc@box}%
  \ifdim\@tempb\p@<\@tempa\p@\let\@tempa\@tempb\fi% select the smaller of both
  \ifdim\@tempa\p@<\p@\scalebox{\@tempa}{\usebox\pandoc@box}%
  \else\usebox{\pandoc@box}%
  \fi%
}
% Set default figure placement to htbp
\def\fps@figure{htbp}
\makeatother



\ifLuaTeX
\usepackage[bidi=basic,shorthands=off]{babel}
\else
\usepackage[bidi=default,shorthands=off]{babel}
\fi
\ifLuaTeX
  \usepackage{selnolig} % disable illegal ligatures
\fi


\setlength{\emergencystretch}{3em} % prevent overfull lines

\providecommand{\tightlist}{%
  \setlength{\itemsep}{0pt}\setlength{\parskip}{0pt}}



 


\usepackage{amsmath,amssymb}
% babel-spanish deja "Figura 8.1.:" con dos signos de puntuacion
\usepackage{caption}
\captionsetup{labelsep=space}
\makeatletter
\@ifpackageloaded{tcolorbox}{}{\usepackage[skins,breakable]{tcolorbox}}
\@ifpackageloaded{fontawesome5}{}{\usepackage{fontawesome5}}
\definecolor{quarto-callout-color}{HTML}{909090}
\definecolor{quarto-callout-note-color}{HTML}{0758E5}
\definecolor{quarto-callout-important-color}{HTML}{CC1914}
\definecolor{quarto-callout-warning-color}{HTML}{EB9113}
\definecolor{quarto-callout-tip-color}{HTML}{00A047}
\definecolor{quarto-callout-caution-color}{HTML}{FC5300}
\definecolor{quarto-callout-color-frame}{HTML}{acacac}
\definecolor{quarto-callout-note-color-frame}{HTML}{4582ec}
\definecolor{quarto-callout-important-color-frame}{HTML}{d9534f}
\definecolor{quarto-callout-warning-color-frame}{HTML}{f0ad4e}
\definecolor{quarto-callout-tip-color-frame}{HTML}{02b875}
\definecolor{quarto-callout-caution-color-frame}{HTML}{fd7e14}
\makeatother
\makeatletter
\@ifpackageloaded{bookmark}{}{\usepackage{bookmark}}
\makeatother
\makeatletter
\@ifpackageloaded{caption}{}{\usepackage{caption}}
\AtBeginDocument{%
\ifdefined\contentsname
  \renewcommand*\contentsname{Tabla de contenidos}
\else
  \newcommand\contentsname{Tabla de contenidos}
\fi
\ifdefined\listfigurename
  \renewcommand*\listfigurename{Listado de Figuras}
\else
  \newcommand\listfigurename{Listado de Figuras}
\fi
\ifdefined\listtablename
  \renewcommand*\listtablename{Listado de Tablas}
\else
  \newcommand\listtablename{Listado de Tablas}
\fi
\ifdefined\figurename
  \renewcommand*\figurename{Figura}
\else
  \newcommand\figurename{Figura}
\fi
\ifdefined\tablename
  \renewcommand*\tablename{Tabla}
\else
  \newcommand\tablename{Tabla}
\fi
}
\@ifpackageloaded{float}{}{\usepackage{float}}
\floatstyle{ruled}
\@ifundefined{c@chapter}{\newfloat{codelisting}{h}{lop}}{\newfloat{codelisting}{h}{lop}[chapter]}
\floatname{codelisting}{Listado}
\newcommand*\listoflistings{\listof{codelisting}{Listado de Listados}}
\usepackage{amsthm}
\theoremstyle{definition}
\newtheorem{definition}{Definición}[chapter]
\theoremstyle{plain}
\newtheorem{proposition}{Proposición}[chapter]
\theoremstyle{definition}
\newtheorem{example}{Ejemplo}[chapter]
\theoremstyle{plain}
\newtheorem{theorem}{Teorema}[chapter]
\theoremstyle{definition}
\newtheorem{exercise}{Ejercicio}[chapter]
\theoremstyle{remark}
\AtBeginDocument{\renewcommand*{\proofname}{Demostración}}
\newtheorem*{remark}{Observación}
\newtheorem*{solution}{Solución}
\newtheorem{refremark}{Observación}[chapter]
\newtheorem{refsolution}{Solución}[chapter]
\makeatother
\makeatletter
\makeatother
\makeatletter
\@ifpackageloaded{caption}{}{\usepackage{caption}}
\@ifpackageloaded{subcaption}{}{\usepackage{subcaption}}
\makeatother
\usepackage{bookmark}
\IfFileExists{xurl.sty}{\usepackage{xurl}}{} % add URL line breaks if available
\urlstyle{same}
% fallback for those not using the hyperref driver hyperxmp:
\makeatletter
\@ifundefined{xmpquote}{\newcommand{\xmpquote}[1]{#1}}{}
\makeatother
\hypersetup{
  pdftitle={Métodos Matemáticos para la Física},
  pdfauthor={Julieta Améndola},
  pdflang={es},
  colorlinks=true,
  linkcolor={violet},
  filecolor={Maroon},
  citecolor={Blue},
  urlcolor={violet},
  pdfcreator={LaTeX via pandoc}}


\title{Métodos Matemáticos para la Física}
\author{Julieta Améndola}
\date{septiembre 2026}
\begin{document}
\frontmatter
\maketitle

\renewcommand*\contentsname{Tabla de contenidos}
{
\setcounter{tocdepth}{2}
\tableofcontents
}

\mainmatter
\bookmarksetup{startatroot}

\chapter*{Antes de empezar}\label{antes-de-empezar}
\addcontentsline{toc}{chapter}{Antes de empezar}

\markboth{Antes de empezar}{Antes de empezar}

Métodos Matemáticos para la Física es una materia rara: no tiene un
objeto propio de estudio. No es ``la teoría de X''. Es una caja de
herramientas, y como toda caja de herramientas se justifica por lo que
se puede hacer con ella.

Eso trae un problema de aprendizaje muy concreto. Uno puede aprender a
resolver una ecuación de Sturm--Liouville, a calcular una función de
Green y a invertir una transformada de Laplace, y salir de la materia
sin haber entendido nunca que son tres formas de contestar la misma
pregunta. Este libro está escrito para que esa pregunta esté siempre a
la vista.

\section*{Los temas}\label{los-temas}
\addcontentsline{toc}{section}{Los temas}

\markright{Los temas}

Son catorce capítulos repartidos en dos partes.

La primera parte va de las ecuaciones diferenciales ordinarias a los
métodos generales que después se usan en todo lo demás: ecuaciones
lineales, puntos singulares y el método de Frobenius; el problema de
Sturm--Liouville, con sus autofunciones y sus desarrollos; las series de
Fourier; la función de Green; las distribuciones y la delta de Dirac; y
las transformadas de Laplace y de Fourier como herramientas para
resolver ecuaciones.

La segunda parte son las ecuaciones en derivadas parciales: qué son,
cómo se clasifican y con qué se resuelven, y después los tres tipos uno
por uno ---hiperbólicas, parabólicas y elípticas---, con los problemas
físicos que les corresponden. Cierran el libro las funciones especiales
que aparecen al separar variables en dominios con simetría, los
desarrollos asintóticos y las ecuaciones integrales.

\section*{Cómo está pensado}\label{cuxf3mo-estuxe1-pensado}
\addcontentsline{toc}{section}{Cómo está pensado}

\markright{Cómo está pensado}

Cada capítulo va del problema físico a la herramienta matemática, y no
al revés. Primero aparece una situación ---una cuerda que vibra, una
barra que se calienta, un potencial que hay que reconstruir--- y recién
después la maquinaria que la resuelve. La razón no es pedagógica sino
histórica: casi todos estos métodos se inventaron así.

Los teoremas están enunciados con precisión y las demostraciones están
cuando aportan algo. Cuando una demostración es larga y no enseña nada
que no se vea mejor en un ejemplo, la reemplazo por el ejemplo y digo
dónde encontrarla.

Al final de cada capítulo hay ejercicios. Algunos son de cálculo, para
fijar la mecánica. Otros piden interpretar un resultado.

Cada capítulo está hecho de \textbf{teoría}, \textbf{figuras generadas
con código} ---el código está a la vista, plegado, en cada figura---,
\textbf{videos} cuando alguien ya explicó algo mejor de lo que yo
podría, \textbf{ejercicios interactivos} que corrigen y explican, y
ejercicios de lápiz y papel con la solución escondida.

Todo el libro se puede descargar en PDF desde el ícono de la barra
lateral. En papel, las partes interactivas se convierten en ejercicios
comunes.

\begin{tcolorbox}[enhanced jigsaw, arc=.35mm, bottomrule=.15mm, breakable, colback=white, colframe=quarto-callout-note-color-frame, left=2mm, leftrule=.75mm, opacityback=0, rightrule=.15mm, toprule=.15mm]
\begin{minipage}[t]{5.5mm}
\textcolor{quarto-callout-note-color}{\faInfo}
\end{minipage}%
\begin{minipage}[t]{\textwidth - 5.5mm}

\vspace{-3mm}\textbf{Material en construcción}\vspace{3mm}

Este libro se está escribiendo. Hay capítulos terminados, capítulos con
partes escritas y capítulos que por ahora son solo el temario; cada uno
avisa al empezar en qué estado está.

\end{minipage}%
\end{tcolorbox}

\part{Ecuaciones diferenciales ordinarias y métodos generales}

\chapter{Ecuaciones diferenciales ordinarias}\label{sec-edo}

\begin{tcolorbox}[enhanced jigsaw, arc=.35mm, bottomrule=.15mm, breakable, colback=white, colframe=quarto-callout-note-color-frame, left=2mm, leftrule=.75mm, opacityback=0, rightrule=.15mm, toprule=.15mm]
\begin{minipage}[t]{5.5mm}
\textcolor{quarto-callout-note-color}{\faInfo}
\end{minipage}%
\begin{minipage}[t]{\textwidth - 5.5mm}

\vspace{-3mm}\textbf{Capítulo en construcción}\vspace{3mm}

Abajo está el temario. El contenido se agrega en orden de necesidad.

\end{minipage}%
\end{tcolorbox}

\section*{Temario}\label{temario}
\addcontentsline{toc}{section}{Temario}

\markright{Temario}

Ecuaciones no lineales de primer orden: Bernoulli y Riccati. Ecuaciones
diferenciales lineales. Wronskiano. Soluciones de ecuaciones homogéneas
e inhomogéneas. Lema de Abel. Método de variación de parámetros. Método
de reducción de orden. Puntos ordinarios. Puntos singulares regulares e
irregulares. Método de Frobenius. Ecuaciones diferenciales especiales:
hipergeométrica, de Bessel, de Legendre. Desarrollos en el infinito.

\chapter{El problema de Sturm--Liouville}\label{sec-sturm}

\begin{tcolorbox}[enhanced jigsaw, arc=.35mm, bottomrule=.15mm, breakable, colback=white, colframe=quarto-callout-note-color-frame, left=2mm, leftrule=.75mm, opacityback=0, rightrule=.15mm, toprule=.15mm]
\begin{minipage}[t]{5.5mm}
\textcolor{quarto-callout-note-color}{\faInfo}
\end{minipage}%
\begin{minipage}[t]{\textwidth - 5.5mm}

\vspace{-3mm}\textbf{Capítulo en construcción}\vspace{3mm}

Abajo está el temario. El contenido se agrega en orden de necesidad.

\end{minipage}%
\end{tcolorbox}

\section*{Temario}\label{temario-1}
\addcontentsline{toc}{section}{Temario}

\markright{Temario}

Producto interno entre funciones. Sistemas de funciones ortogonales.
Desigualdad de Bessel. Propiedades de soluciones de ecuaciones
diferenciales ordinarias. Teoremas de separación, oscilación y
comparación de Sturm. Autofunciones y autovalores. Operadores
autoadjuntos. Fórmulas de Lagrange y Green. Problema de Sturm--Liouville
regular.

\chapter{Series de Fourier}\label{sec-fourier}

\begin{tcolorbox}[enhanced jigsaw, arc=.35mm, bottomrule=.15mm, breakable, colback=white, colframe=quarto-callout-note-color-frame, left=2mm, leftrule=.75mm, opacityback=0, rightrule=.15mm, toprule=.15mm]
\begin{minipage}[t]{5.5mm}
\textcolor{quarto-callout-note-color}{\faInfo}
\end{minipage}%
\begin{minipage}[t]{\textwidth - 5.5mm}

\vspace{-3mm}\textbf{Capítulo en construcción}\vspace{3mm}

Abajo está el temario. Después hay un \textbf{adelanto}: una sección
escrita, para ver cómo va a estar armado cada capítulo cuando esté
terminado.

\end{minipage}%
\end{tcolorbox}

\section*{Temario}\label{temario-2}
\addcontentsline{toc}{section}{Temario}

\markright{Temario}

Definición de serie de Fourier. Condiciones para que una serie de
Fourier converja a la función. Series en seno y en coseno. Forma
compleja. Convergencias puntual y en media. Desigualdad de Bessel e
identidad de Parseval. Aplicaciones.

\section*{Adelanto: qué pasa en un
salto}\label{adelanto-quuxe9-pasa-en-un-salto}
\addcontentsline{toc}{section}{Adelanto: qué pasa en un salto}

\markright{Adelanto: qué pasa en un salto}

Una serie de Fourier es una suma de senos y cosenos, que son funciones
perfectamente suaves. La pregunta natural es qué puede hacer una suma
así frente a una función que \textbf{salta}.

La onda cuadrada es el caso de manual. Vale \(-1\) en \((-\pi, 0)\) y
\(+1\) en \((0, \pi)\), y su serie de Fourier resulta

\begin{equation}\protect\phantomsection\label{eq-cuadrada}{
f(x) \;\sim\; \frac{4}{\pi} \sum_{k \text{ impar}} \frac{\sin kx}{k}
\;=\; \frac{4}{\pi}\left( \sin x + \frac{\sin 3x}{3} + \frac{\sin 5x}{5} + \cdots \right).
}\end{equation}

El teorema de Dirichlet dice que esta serie converge a \(f(x)\) en todo
punto donde \(f\) es continua, y en el salto converge al
\textbf{promedio de los dos límites laterales}. Acá, a cero.

Eso es la convergencia puntual. Lo que no dice es qué forma tienen las
sumas parciales por el camino, y ahí aparece algo que sorprende la
primera vez.

\begin{figure}

\centering{

\pandocbounded{\includegraphics[keepaspectratio]{03-fourier_files/figure-pdf/fig-gibbs-output-1.pdf}}

}

\caption{\label{fig-gibbs}Izquierda: las sumas parciales de la
\ref{eq-cuadrada} con 1, 5 y 25 términos. Derecha: un acercamiento al
salto en \(x=0\). La oscilación se aprieta contra la discontinuidad a
medida que crece \(N\), pero el pico \textbf{no baja}: se estaciona en
\(1{,}17898\).}

\end{figure}%

Esto es el \textbf{fenómeno de Gibbs}. La suma parcial se pasa del valor
límite en un pico que no se achica nunca: se corre hacia el salto y se
angosta, pero su altura tiende a

\begin{equation}\protect\phantomsection\label{eq-gibbs}{
\frac{2}{\pi}\int_0^{\pi} \frac{\sin t}{t}\,dt \;=\; 1{,}178979\ldots
}\end{equation}

Es decir, casi un \(18\%\) por encima del valor \(1\) hacia el que
converge la serie, o un \(9\%\) del salto total, que vale \(2\). Y no es
un defecto de la onda cuadrada: aparece en el mismo porcentaje en
\textbf{cualquier} función con un salto.

No hay contradicción con Dirichlet. La convergencia puntual dice que,
fijado un \(x\), la sucesión de sumas parciales tiende a \(f(x)\). El
pico también tiende a \(x=0\), así que para cada \(x\) fijo termina
quedando a la izquierda del punto. Lo que falla es la convergencia
\textbf{uniforme}: no hay ningún \(N\) a partir del cual todas las sumas
parciales estén dentro de una banda estrecha alrededor de \(f\).

Video recomendado: Grant Sanderson (3Blue1Brown), \emph{But what is a
Fourier series? From heat flow to circle drawings},
\url{https://www.youtube.com/watch?v=r6sGWTCMz2k}.

La \ref{eq-cuadrada} evaluada en el punto de salto \(x = 0\) da
\(\sin 0 + \tfrac{1}{3}\sin 0 + \cdots\), que es cero. Pero la onda
cuadrada no vale cero en \(x=0\): vale \(-1\) por izquierda y \(+1\) por
derecha. ¿Qué está pasando?

\begin{itemize}
\tightlist
\item
  {La serie está mal calculada.}
\item
  {La serie converge al promedio de los límites laterales, que es cero.}
\item
  {La serie diverge en el punto de salto.}
\item
  {La serie converge a \(+1\), porque es el valor por derecha.}
\end{itemize}

Es exactamente lo que garantiza el teorema de Dirichlet: en un punto de
salto la serie converge a \(\tfrac{1}{2}\big(f(x^-) + f(x^+)\big)\).
Acá, \(\tfrac{1}{2}(-1 + 1) = 0\). La serie no ``conoce'' el valor que
le asignemos a \(f\) en el punto: los coeficientes son integrales, y
cambiar una función en un punto aislado no cambia ninguna integral.

\begin{tcolorbox}[enhanced jigsaw, arc=.35mm, bottomrule=.15mm, bottomtitle=1mm, breakable, colback=white, colbacktitle=quarto-callout-tip-color!10!white, colframe=quarto-callout-tip-color-frame, coltitle=black, left=2mm, leftrule=.75mm, opacityback=0, opacitybacktitle=0.6, rightrule=.15mm, title=\textcolor{quarto-callout-tip-color}{\faLightbulb}\hspace{0.5em}{Para hacer a mano: el valor de \(\pi/4\)}, titlerule=0mm, toprule=.15mm, toptitle=1mm]

Evaluando la \ref{eq-cuadrada} en \(x = \pi/2\), deducir la fórmula de
Leibniz \[
\frac{\pi}{4} = 1 - \frac13 + \frac15 - \frac17 + \cdots
\] y estimar cuántos términos hacen falta para obtener tres decimales
correctos. La respuesta al segundo punto explica por qué nadie calcula
\(\pi\) así.

\end{tcolorbox}

\chapter{Función de Green}\label{sec-green}

\begin{tcolorbox}[enhanced jigsaw, arc=.35mm, bottomrule=.15mm, breakable, colback=white, colframe=quarto-callout-note-color-frame, left=2mm, leftrule=.75mm, opacityback=0, rightrule=.15mm, toprule=.15mm]
\begin{minipage}[t]{5.5mm}
\textcolor{quarto-callout-note-color}{\faInfo}
\end{minipage}%
\begin{minipage}[t]{\textwidth - 5.5mm}

\vspace{-3mm}\textbf{Capítulo en construcción}\vspace{3mm}

Abajo está el temario. El contenido se agrega en orden de necesidad.

\end{minipage}%
\end{tcolorbox}

\section*{Temario}\label{temario-3}
\addcontentsline{toc}{section}{Temario}

\markright{Temario}

Ecuaciones diferenciales lineales con funciones impulsivas. Función de
Green. Función de Green para el operador de Sturm--Liouville. Desarrollo
en serie para la función de Green. Aplicaciones a vigas cargadas y
conductores eléctricos cargados.

\chapter{Distribuciones}\label{sec-distribuciones}

\begin{tcolorbox}[enhanced jigsaw, arc=.35mm, bottomrule=.15mm, breakable, colback=white, colframe=quarto-callout-note-color-frame, left=2mm, leftrule=.75mm, opacityback=0, rightrule=.15mm, toprule=.15mm]
\begin{minipage}[t]{5.5mm}
\textcolor{quarto-callout-note-color}{\faInfo}
\end{minipage}%
\begin{minipage}[t]{\textwidth - 5.5mm}

\vspace{-3mm}\textbf{Capítulo en construcción}\vspace{3mm}

Abajo está el temario. El contenido se agrega en orden de necesidad.

\end{minipage}%
\end{tcolorbox}

\section*{Temario}\label{temario-4}
\addcontentsline{toc}{section}{Temario}

\markright{Temario}

Delta de Dirac. Sucesiones y distribuciones. Función salto.
Representaciones y cálculos con distribuciones. Aplicaciones.

\chapter{Transformada de Laplace}\label{sec-laplace}

\begin{tcolorbox}[enhanced jigsaw, arc=.35mm, bottomrule=.15mm, breakable, colback=white, colframe=quarto-callout-note-color-frame, left=2mm, leftrule=.75mm, opacityback=0, rightrule=.15mm, toprule=.15mm]
\begin{minipage}[t]{5.5mm}
\textcolor{quarto-callout-note-color}{\faInfo}
\end{minipage}%
\begin{minipage}[t]{\textwidth - 5.5mm}

\vspace{-3mm}\textbf{Capítulo en construcción}\vspace{3mm}

Abajo está el temario. El contenido se agrega en orden de necesidad.

\end{minipage}%
\end{tcolorbox}

\section*{Temario}\label{temario-5}
\addcontentsline{toc}{section}{Temario}

\markright{Temario}

Transformadas integrales. Transformada de Laplace. Fórmulas de
inversión. Transformada de Laplace de distribuciones. Convolución y
aplicaciones. Resolución de ecuaciones diferenciales con transformadas.

\chapter{Transformada de Fourier}\label{sec-tfourier}

\begin{tcolorbox}[enhanced jigsaw, arc=.35mm, bottomrule=.15mm, breakable, colback=white, colframe=quarto-callout-note-color-frame, left=2mm, leftrule=.75mm, opacityback=0, rightrule=.15mm, toprule=.15mm]
\begin{minipage}[t]{5.5mm}
\textcolor{quarto-callout-note-color}{\faInfo}
\end{minipage}%
\begin{minipage}[t]{\textwidth - 5.5mm}

\vspace{-3mm}\textbf{Capítulo en construcción}\vspace{3mm}

Abajo está el temario. El contenido se agrega en orden de necesidad.

\end{minipage}%
\end{tcolorbox}

\section*{Temario}\label{temario-6}
\addcontentsline{toc}{section}{Temario}

\markright{Temario}

Fórmula integral de Fourier. Fórmulas de inversión. Transformada de
Fourier directa e inversa. Propiedades, teorema de Parseval.
Transformada de Fourier de distribuciones. Transformada seno y coseno.
Convolución, correlación y autocorrelación. Resolución de ecuaciones
diferenciales con transformadas.

\part{Ecuaciones en derivadas parciales y métodos avanzados}

\chapter{Ecuaciones en derivadas parciales}\label{sec-edp}

\begin{tcolorbox}[enhanced jigsaw, arc=.35mm, bottomrule=.15mm, breakable, colback=white, colframe=quarto-callout-note-color-frame, left=2mm, leftrule=.75mm, opacityback=0, rightrule=.15mm, toprule=.15mm]
\begin{minipage}[t]{5.5mm}
\textcolor{quarto-callout-note-color}{\faInfo}
\end{minipage}%
\begin{minipage}[t]{\textwidth - 5.5mm}

\vspace{-3mm}\textbf{Qué hace falta para este capítulo}\vspace{3mm}

Para las seis primeras secciones alcanza con ecuaciones diferenciales
ordinarias de primero y segundo orden, la regla de la cadena para
funciones de varias variables, curvas de nivel y ---para la
sección~\ref{sec-muchas-variables}--- diagonalización de matrices
simétricas. Las tres últimas secciones son la continuación directa de
los capítulos anteriores en varias variables: transformadas integrales
retoma la Capítulo~\ref{sec-laplace} y la Capítulo~\ref{sec-tfourier},
función de Green retoma la Capítulo~\ref{sec-green}, y separación de
variables retoma la Capítulo~\ref{sec-sturm}. Cada una avisa al empezar
qué necesita.

\end{minipage}%
\end{tcolorbox}

\section{De la partícula al campo}\label{sec-motivacion}

Casi todo lo que vimos hasta acá tenía la misma forma: una magnitud que
depende de \textbf{una} variable y una ecuación que la relaciona con sus
derivadas. La segunda ley de Newton para una partícula es de ese tipo.
Hay una trayectoria \(x(t)\), y la ecuación dice cómo se curva.

Pero la mayor parte de la física no habla de partículas sino de
\textbf{campos}: la temperatura en cada punto de una barra, la altura de
una cuerda en cada punto y en cada instante, el potencial eléctrico en
cada punto de una región, la función de onda de un electrón. Un campo es
una función de varias variables, y las leyes que lo gobiernan son
\emph{locales}: no dicen qué hace el campo en todo el espacio de golpe,
sino cómo lo que pasa en un punto está atado a lo que pasa en sus
vecinos inmediatos un instante después.

«Cómo cambia acá respecto de sus vecinos» es exactamente una derivada
parcial. Por eso el lenguaje natural de la física de campos son las
ecuaciones en derivadas parciales, y por eso este capítulo abre la
segunda mitad del libro: de acá en adelante todo va a ser, de una forma
u otra, una EDP.

Video recomendado: 3Blue1Brown, \emph{But what is a partial differential
equation?}, \url{https://www.youtube.com/watch?v=ly4S0oi3Yz8}.

\subsection{Una ley de conservación, tres ecuaciones
distintas}\label{sec-conservacion}

Vale la pena ver una sola vez cómo \emph{aparece} una de estas
ecuaciones, porque el mismo argumento produce tres de las que vamos a
estudiar durante todo el capítulo.

Sea \(u(x,t)\) la densidad lineal de algo que se conserva ---masa,
carga, energía, cantidad de autos en una ruta--- y sea \(q(x,t)\) el
\textbf{flujo}: cuánto de esa cantidad atraviesa por segundo la posición
\(x\), hacia la derecha. Que la cantidad se conserve significa que en
cualquier tramo \([a,b]\) lo que hay adentro solo cambia por lo que
entra por un extremo y sale por el otro:

\[
\frac{d}{dt}\int_a^b u(x,t)\,dx \;=\; q(a,t) - q(b,t) \;=\; -\int_a^b \frac{\partial q}{\partial x}\,dx .
\]

Si \(u\) y \(q\) son suaves podemos meter la derivada temporal adentro
de la integral, y como la igualdad vale para \textbf{todo} tramo
\([a,b]\), los integrandos tienen que coincidir punto a punto:

\begin{equation}\protect\phantomsection\label{eq-conservacion}{
\frac{\partial u}{\partial t} + \frac{\partial q}{\partial x} = 0 .
}\end{equation}

\begin{definition}[Ley de
conservación]\protect\hypertarget{def-ley-conservacion}{}\label{def-ley-conservacion}

Una \textbf{ley de conservación} en una dimensión espacial es una
ecuación de la forma \ref{eq-conservacion}, donde \(u\) es la densidad
de una magnitud conservada y \(q\) su flujo. Cuando el flujo es una
función de la propia densidad, \(q = F(u)\), la ecuación se escribe
\(u_t + F(u)_x = 0\) y se dice que \(F\) es la \textbf{función de
flujo}.

\end{definition}

\ref{eq-conservacion} todavía no es una ecuación: tiene dos incógnitas,
\(u\) y \(q\). Lo que la convierte en una ecuación es decir \emph{de qué
manera} el flujo depende del estado. Esa segunda relación no sale de la
conservación sino de la física del material, y se la llama
\textbf{relación constitutiva}. Miremos tres:

\textbf{Uno.} Lo que hay es simplemente arrastrado por una corriente de
velocidad constante \(c\). Entonces por cada punto pasa por segundo lo
que hay ahí, multiplicado por lo rápido que va: \(q = c\,u\).
Reemplazando en \ref{eq-conservacion},

\[
u_t + c\,u_x = 0 ,
\]

la \textbf{ecuación de transporte}. Es la más simple de todas y es donde
vamos a aprender el método más importante del capítulo.

\textbf{Dos.} No hay corriente, pero la sustancia se mueve
espontáneamente de donde hay mucho hacia donde hay poco, tanto más
rápido cuanto más pronunciada sea la diferencia. Eso es la ley de
Fourier para el calor y la ley de Fick para una concentración:
\(q = -k\,u_x\), con \(k>0\). El signo menos es toda la física ---el
flujo apunta hacia abajo de la pendiente---. Ahora \ref{eq-conservacion}
da

\[
u_t = k\,u_{xx} ,
\]

la \textbf{ecuación de difusión}, o del calor.

\textbf{Tres.} Lo que hay es arrastrado, pero la velocidad de arrastre
no es constante: es la propia densidad, de modo que las regiones donde
hay más avanzan más rápido. Entonces \(q = u^2/2\), porque
\(\partial_x(u^2/2) = u\,u_x\), y queda

\[
u_t + u\,u_x = 0 ,
\]

la \textbf{ecuación de Burgers sin viscosidad}. Es la más inocente de
las tres a la vista y la única que va a romperse sola: si las crestas
viajan más rápido que los valles, tarde o temprano una cresta alcanza al
valle que tiene adelante y la solución deja de ser una función.
Volveremos sobre esto en la sección~\ref{sec-bien-planteado} y en la
sección~\ref{sec-casi-lineales}.

\begin{tcolorbox}[enhanced jigsaw, arc=.35mm, bottomrule=.15mm, breakable, colback=white, colframe=quarto-callout-tip-color-frame, left=2mm, leftrule=.75mm, opacityback=0, rightrule=.15mm, toprule=.15mm]
\begin{minipage}[t]{5.5mm}
\textcolor{quarto-callout-tip-color}{\faLightbulb}
\end{minipage}%
\begin{minipage}[t]{\textwidth - 5.5mm}

\vspace{-3mm}\textbf{Lo que hay que llevarse de acá}\vspace{3mm}

Una misma ley de conservación, tres relaciones constitutivas distintas,
tres ecuaciones con comportamientos irreconciliables: una transporta sin
deformar, otra suaviza y desparrama, la tercera se rompe. La forma de la
ecuación no lo grita, pero está todo escrito ahí. Aprender a leerlo es
la primera mitad de este capítulo.

\end{minipage}%
\end{tcolorbox}

Un gas de partículas se mueve por un tubo con velocidad constante \(c\),
pero además las partículas se desintegran a una tasa proporcional a
cuántas hay. La densidad ya no se conserva: por cada unidad de longitud
y de tiempo desaparecen \(\lambda u\) partículas. ¿Cuál es la ecuación?

\begin{itemize}
\tightlist
\item
  {\(u_t + c\,u_x = 0\)}
\item
  {\(u_t + c\,u_x = -\lambda u\)}
\item
  {\(u_t + c\,u_x = -\lambda u_x\)}
\item
  {\(u_t = k\,u_{xx} - \lambda u\)}
\end{itemize}

El balance de \ref{eq-conservacion} decía «lo que cambia adentro = lo
que entra menos lo que sale». Si además hay un sumidero que se come
\(\lambda u\) por unidad de longitud y de tiempo, el balance pasa a ser
\[\frac{d}{dt}\int_a^b u\,dx = q(a,t) - q(b,t) - \int_a^b \lambda u\,dx,\]
y el mismo argumento de antes deja \(u_t + q_x = -\lambda u\). Con
\(q = cu\) queda la segunda opción. La tercera confunde un sumidero con
un flujo, y la cuarta agrega una difusión que nadie pidió: no hay
ninguna ley de Fick en el enunciado.

\subsection{Las ecuaciones que vamos a encontrar}\label{sec-catalogo}

Antes de clasificar nada conviene tener a la vista el catálogo. Todas
estas van a aparecer en este capítulo o en los tres siguientes.

\begin{longtable}[]{@{}
  >{\raggedright\arraybackslash}p{(\linewidth - 4\tabcolsep) * \real{0.2000}}
  >{\raggedright\arraybackslash}p{(\linewidth - 4\tabcolsep) * \real{0.3800}}
  >{\raggedright\arraybackslash}p{(\linewidth - 4\tabcolsep) * \real{0.4200}}@{}}
\caption{Las ecuaciones del capítulo. Todas menos las dos últimas son de
segundo orden; todas menos Burgers son
lineales.}\label{tbl-catalogo}\tabularnewline
\toprule\noalign{}
\begin{minipage}[b]{\linewidth}\raggedright
Nombre
\end{minipage} & \begin{minipage}[b]{\linewidth}\raggedright
Ecuación
\end{minipage} & \begin{minipage}[b]{\linewidth}\raggedright
Qué describe
\end{minipage} \\
\midrule\noalign{}
\endfirsthead
\toprule\noalign{}
\begin{minipage}[b]{\linewidth}\raggedright
Nombre
\end{minipage} & \begin{minipage}[b]{\linewidth}\raggedright
Ecuación
\end{minipage} & \begin{minipage}[b]{\linewidth}\raggedright
Qué describe
\end{minipage} \\
\midrule\noalign{}
\endhead
\bottomrule\noalign{}
\endlastfoot
Transporte & \(u_t + c\,u_x = 0\) & algo que viaja sin deformarse \\
Onda & \(u_{tt} = c^2\,u_{xx}\) & una cuerda, el sonido, una barra que
vibra \\
Difusión o calor & \(u_t = k\,u_{xx}\) & temperatura, concentración \\
Laplace & \(\nabla^2 u = 0\) & potencial en equilibrio, sin fuentes \\
Poisson & \(\nabla^2 u = -\rho/\varepsilon_0\) & potencial en
equilibrio, con fuentes \\
Helmholtz & \(\nabla^2 u + \lambda u = 0\) & modos estacionarios,
régimen oscilatorio \\
Schrödinger & \(i\hbar\,\psi_t = -\dfrac{\hbar^2}{2m}\psi_{xx} + V\psi\)
& amplitud de probabilidad cuántica \\
Burgers & \(u_t + u\,u_x = 0\) & ondas que se empinan y rompen \\
Viga de Euler--Bernoulli & \(\rho\,w_{tt} + EI\,w_{xxxx} = f\) & flexión
de una viga elástica \\
\end{longtable}

\subsection{Un signo separa tres mundos}\label{sec-tres-mundos}

Miremos tres de ellas escritas una al lado de la otra, con todo del
mismo lado:

\[
u_{tt} - c^2 u_{xx} = 0, \qquad u_t - k\,u_{xx} = 0, \qquad u_{xx} + u_{yy} = 0 .
\]

Se parecen muchísimo. Las tres tienen una derivada segunda respecto de
la variable espacial; las diferencias son qué se le pone al lado ---una
derivada segunda, una primera, otra segunda--- y un signo. Y sin embargo
describen tres mundos que no se tocan:

\begin{itemize}
\tightlist
\item
  En la \textbf{ecuación de onda}, la información viaja a velocidad
  finita \(c\). Un pico con un quiebre sigue teniendo el quiebre para
  siempre. Si filmamos la solución y pasamos la película al revés, sigue
  siendo una solución.
\item
  En la \textbf{ecuación de difusión}, cualquier dato inicial se vuelve
  infinitamente derivable apenas \(t>0\): los quiebres se borran
  instantáneamente. Además la influencia es inmediata a cualquier
  distancia. Y la película al revés \textbf{no} es una solución:
  intentar reconstruir el pasado es un problema mal planteado, del que
  hablaremos en la sección~\ref{sec-bien-planteado}.
\item
  En la \textbf{ecuación de Laplace} no hay tiempo. No hay nada que
  evolucione. El valor en cada punto es el promedio de los valores a su
  alrededor, y por lo tanto depende de \emph{todo} el borde a la vez.
\end{itemize}

Nada en el aspecto de las tres ecuaciones anuncia esa diferencia. Lo que
la anuncia es un solo número, que se calcula en dos líneas, y aprender a
calcularlo es lo primero que vamos a hacer con las ecuaciones de segundo
orden.

\begin{figure}

\centering{

\pandocbounded{\includegraphics[keepaspectratio]{08-edp_files/figure-pdf/fig-tres-mundos-output-1.pdf}}

}

\caption{\label{fig-tres-mundos}El mismo perfil ---un pico con dos
quiebres--- bajo las tres ecuaciones modelo. La de onda lo parte en dos
mitades que se van en sentidos opuestos, con los quiebres intactos para
siempre. La de difusión los borra apenas arranca el tiempo (la lupa
muestra el vértice con \(t\) mil veces más chico) y después desparrama
el pico. La de Laplace ni siquiera tiene tiempo: el perfil, dibujado en
la tira de abajo, es un dato prescrito sobre el borde inferior del
cuadrado, y la solución ---el color y las curvas de nivel--- es el único
equilibrio compatible con ese borde.}

\end{figure}%

\section{Lo que se rompe al pasar de una variable a
varias}\label{sec-porque}

Una ecuación diferencial ordinaria relaciona una función de una variable
con sus derivadas. Su solución general depende de un número finito de
constantes: dos para una ecuación de segundo orden, \(n\) para una de
orden \(n\). Fijar esas constantes es fijar dos o \(n\) números, y para
eso alcanzan dos o \(n\) condiciones.

Con las derivadas parciales eso se rompe de entrada. Consideremos la
ecuación más simple que se pueda escribir en dos variables:

\begin{equation}\protect\phantomsection\label{eq-trivial}{
\frac{\partial u}{\partial y}(x,y) = 0 .
}\end{equation}

Nos dice que \(u\) no cambia cuando nos movemos en la dirección \(y\).
Su solución general es

\[
u(x,y) = f(x),
\]

con \(f\) \textbf{cualquier} función derivable de una variable. La
solución general no depende de constantes: depende de una función
arbitraria. Y una función arbitraria no se fija con dos números.

No es un accidente de esa ecuación en particular. Tomemos

\begin{equation}\protect\phantomsection\label{eq-trivial2}{
\frac{\partial^2 u}{\partial x\,\partial y} = 0 .
}\end{equation}

Leída como «\(\partial_y(u_x) = 0\)», dice que \(u_x\) no depende de
\(y\), o sea \(u_x = p(x)\); integrando en \(x\) queda

\[
u(x,y) = f(x) + g(y),
\]

con \(f\) y \(g\) arbitrarias. Una ecuación de segundo orden, dos
funciones arbitrarias. La regla empírica ---que no es un teorema, pero
orienta--- es que una EDP de orden \(n\) arrastra \(n\) funciones
arbitrarias, así como una EDO de orden \(n\) arrastra \(n\) constantes.

Vale la pena guardar \ref{eq-trivial2}, porque no es un ejemplo de
juguete: en la sección~\ref{sec-formas-canonicas} vamos a ver que
\textbf{toda} ecuación hiperbólica se convierte en ella con un cambio de
variables, y de ahí va a salir la solución \(u = f(x+ct) + g(x-ct)\) de
la ecuación de onda ---dos funciones arbitrarias, exactamente las que
anticipa la regla---.

Este hecho, que parece un detalle técnico, es el que organiza todo el
resto del capítulo. Si la indeterminación de una EDP es funcional, los
datos que la resuelven también tienen que ser funcionales: valores de
\(u\) sobre una curva, no en dos puntos. De ahí que en EDP se hable de
condiciones \emph{de contorno} e \emph{iniciales}, y que la pregunta de
cuáles alcanzan sea seria y no trivial.

\begin{tcolorbox}[enhanced jigsaw, arc=.35mm, bottomrule=.15mm, breakable, colback=white, colframe=quarto-callout-tip-color-frame, left=2mm, leftrule=.75mm, opacityback=0, rightrule=.15mm, toprule=.15mm]
\begin{minipage}[t]{5.5mm}
\textcolor{quarto-callout-tip-color}{\faLightbulb}
\end{minipage}%
\begin{minipage}[t]{\textwidth - 5.5mm}

\vspace{-3mm}\textbf{La versión corta}\vspace{3mm}

En EDO se busca una curva y alcanzan las condiciones puntuales. En EDP
se busca una superficie y hacen falta datos sobre curvas enteras. Casi
todas las dificultades del capítulo salen de ahí.

\end{minipage}%
\end{tcolorbox}

¿Cuál es la solución general de \(\;u_{xx} = 0\;\) para \(u = u(x,y)\)?

\begin{itemize}
\tightlist
\item
  {\(u(x,y) = ax + b\), con \(a\) y \(b\) constantes}
\item
  {\(u(x,y) = f(y)\,x + g(y)\), con \(f\) y \(g\) arbitrarias}
\item
  {\(u(x,y) = f(x) + g(y)\), con \(f\) y \(g\) arbitrarias}
\item
  {\(u(x,y) = f(y)\), con \(f\) arbitraria}
\end{itemize}

Integrar en \(x\) dos veces sí se puede, pero las «constantes» de
integración son constantes \emph{respecto de \(x\)}, no respecto de
todo. Cada una puede depender libremente de \(y\): \(u_x = f(y)\) y
después \(u = f(y)\,x + g(y)\). Es una ecuación de segundo orden y
aparecen dos funciones arbitrarias, como anticipa la regla empírica. La
primera opción es la respuesta que daríamos si \(u\) dependiera solo de
\(x\), y es el error típico: en EDP las constantes de integración son
funciones de las otras variables.

\section{Qué es exactamente una EDP}\label{sec-vocabulario}

Antes de seguir, fijemos el vocabulario y la notación. Todo lo que sigue
en el capítulo está escrito en estos términos.

\begin{definition}[Ecuación en derivadas
parciales]\protect\hypertarget{def-edp}{}\label{def-edp}

Sea \(u\) una función incógnita de \(n\) variables \(x_1,\dots,x_n\).
Una \textbf{ecuación en derivadas parciales} es una relación de la forma

\begin{equation}\protect\phantomsection\label{eq-general-edp}{
\Phi\!\left(x_1,\dots,x_n,\;u,\;\frac{\partial u}{\partial x_1},\dots,\frac{\partial u}{\partial x_n},\;
\frac{\partial^2 u}{\partial x_1^2},\frac{\partial^2 u}{\partial x_1\partial x_2},\dots,\;
\frac{\partial^N u}{\partial x_n^N}\right) = 0 .
}\end{equation}

Su \textbf{orden} es el mayor \(N\) tal que en la relación aparece
efectivamente una derivada de orden \(N\).

\end{definition}

Para no escribir cadenas de símbolos cada vez, adoptamos dos
abreviaturas de uso universal. La primera es poner las variables de
derivación como subíndices:

\[
\frac{\partial u}{\partial x_i} = u_{x_i}, \qquad
\frac{\partial^2 u}{\partial x_i\,\partial x_j} = u_{x_i x_j} .
\]

La segunda es la \textbf{notación multiíndice}, que sirve para hablar de
«todas las derivadas de orden \(\alpha\)» sin listarlas. Dado
\(\alpha = (\alpha_1,\dots,\alpha_n)\) con \(\alpha_i \ge 0\) enteros,
se escribe \(|\alpha| = \alpha_1 + \dots + \alpha_n\) y

\begin{equation}\protect\phantomsection\label{eq-multiindice}{
D^{\alpha}u = \frac{\partial^{|\alpha|} u}{\partial x_1^{\alpha_1}\cdots\partial x_n^{\alpha_n}} .
}\end{equation}

Así, \(D^{(2,0)}u = u_{xx}\), \(D^{(1,1)}u = u_{xy}\) y
\(D^{(0,2)}u = u_{yy}\) son las tres derivadas de orden \(2\) en dos
variables, y la suma \(\sum_{|\alpha|=2} a_\alpha D^\alpha u\) las
recorre todas de una.

\begin{definition}[Solución
clásica]\protect\hypertarget{def-solucion-clasica}{}\label{def-solucion-clasica}

Una \textbf{solución clásica} ---o \textbf{fuerte}--- de una EDP de
orden \(N\) en un dominio \(D\) es una función \(u\) que es \(N\) veces
continuamente derivable en \(D\) y que satisface la ecuación en todos
los puntos de \(D\).

\end{definition}

El calificativo «clásica» está avisando que hay otras. Ya vimos con la
ecuación de Burgers que hay problemas físicos perfectamente razonables
cuya solución deja de ser derivable, y a veces deja de ser continua.
Esas soluciones \textbf{generalizadas} o \textbf{débiles} son el tema de
la sección~\ref{sec-bien-planteado}, y son las que le dan sentido físico
a las ondas de choque.

¿De qué orden es la ecuación de la viga de la Tabla~\ref{tbl-catalogo},
\[\rho\,w_{tt} + \gamma\,w_t + EI\,w_{xxxx} = f(x,t)\;?\] Escribí solo
el número.

El orden es el de la derivada más alta que aparece, sin importar
respecto de qué variable ni cuántas veces aparezcan las otras. Acá hay
derivadas segundas en \(t\) y una derivada cuarta en \(x\), así que la
ecuación es de orden \(4\). Es la única de la tabla que no es de segundo
orden, y por eso la vamos a tratar aparte, en la
sección~\ref{sec-separacion}.

\section{Las tres preguntas del capítulo}\label{sec-plan}

Todo lo que sigue son respuestas a tres preguntas, siempre en este
orden.

\textbf{1. ¿De qué tipo es esta ecuación?} Hay dos clasificaciones y son
independientes entre sí. Una mira cuán lineal es la ecuación ---lineal,
semilineal, cuasi-lineal, completamente no lineal--- y decide qué
métodos están permitidos (sección~\ref{sec-linealidad}). La otra mira
solo los términos de orden más alto y decide cómo se comportan las
soluciones: hiperbólica, parabólica o elíptica
(sección~\ref{sec-clasificacion}, sección~\ref{sec-muchas-variables}).
Una misma ecuación tiene un tipo en cada clasificación.

\textbf{2. ¿Qué datos hacen falta, y alcanzan?} Una ecuación sola no
determina nada. Hay que agregarle condiciones, y no cualquier juego de
condiciones sirve: algunos dejan infinitas soluciones, otros ninguna, y
otros ---los peores, porque no se notan--- dejan exactamente una
solución que no se parece en nada a la del problema de al lado. Eso es
el problema bien planteado (sección~\ref{sec-bien-planteado}).

\textbf{3. ¿Cómo se construye la solución?} Cuatro métodos, y cada uno
es una sección:

\begin{longtable}[]{@{}
  >{\raggedright\arraybackslash}p{(\linewidth - 4\tabcolsep) * \real{0.2600}}
  >{\raggedright\arraybackslash}p{(\linewidth - 4\tabcolsep) * \real{0.4400}}
  >{\raggedright\arraybackslash}p{(\linewidth - 4\tabcolsep) * \real{0.3000}}@{}}
\caption{Los cuatro métodos del
capítulo.}\label{tbl-metodos}\tabularnewline
\toprule\noalign{}
\begin{minipage}[b]{\linewidth}\raggedright
Método
\end{minipage} & \begin{minipage}[b]{\linewidth}\raggedright
Idea
\end{minipage} & \begin{minipage}[b]{\linewidth}\raggedright
Dónde
\end{minipage} \\
\midrule\noalign{}
\endfirsthead
\toprule\noalign{}
\begin{minipage}[b]{\linewidth}\raggedright
Método
\end{minipage} & \begin{minipage}[b]{\linewidth}\raggedright
Idea
\end{minipage} & \begin{minipage}[b]{\linewidth}\raggedright
Dónde
\end{minipage} \\
\midrule\noalign{}
\endhead
\bottomrule\noalign{}
\endlastfoot
Características & convertir la EDP en EDOs sobre curvas &
sección~\ref{sec-caracteristicas}, sección~\ref{sec-casi-lineales} \\
Transformadas integrales & sacarse una variable de encima &
sección~\ref{sec-transformadas} \\
Función de Green & resolver una vez para una fuente puntual &
sección~\ref{sec-green-varias} \\
Separación de variables & buscar soluciones producto y sumarlas &
sección~\ref{sec-separacion} \\
\end{longtable}

Ninguno sirve siempre. Características vive de que existan curvas reales
sobre las cuales la ecuación se vuelve ordinaria, y por eso funciona de
maravilla en las hiperbólicas y no existe en las elípticas. Separación
de variables necesita que el dominio tenga la simetría adecuada. Green
necesita linealidad. Elegir bien es la mitad del trabajo, y la
clasificación de la pregunta 1 es justamente lo que permite elegir sin
probar a ciegas.

\section{Cuán lineal es una ecuación}\label{sec-linealidad}

Vamos con la primera de las dos clasificaciones. Esta mira \textbf{cómo
entra la incógnita en la ecuación} y decide qué métodos están
permitidos. La otra, que viene en la sección~\ref{sec-clasificacion},
mira solamente los términos de orden más alto y decide cómo se comportan
las soluciones. Son independientes: toda ecuación tiene un tipo en cada
una, y decir «hiperbólica» no dice nada sobre si es lineal, ni al revés.

Las dos, eso sí, giran alrededor del mismo objeto.

\begin{definition}[Parte
principal]\protect\hypertarget{def-parte-principal}{}\label{def-parte-principal}

Sea una EDP de orden \(N\) escrita como

\[
\sum_{|\alpha| = N} a_\alpha\, D^\alpha u \;+\; (\text{términos con derivadas de orden} < N) = 0 .
\]

Su \textbf{parte principal} es el operador formado solo por los términos
de orden exactamente \(N\):

\begin{equation}\protect\phantomsection\label{eq-parte-principal}{
L_P[u] = \sum_{|\alpha| = N} a_\alpha\, D^\alpha u .
}\end{equation}

\end{definition}

\begin{tcolorbox}[enhanced jigsaw, arc=.35mm, bottomrule=.15mm, breakable, colback=white, colframe=quarto-callout-important-color-frame, left=2mm, leftrule=.75mm, opacityback=0, rightrule=.15mm, toprule=.15mm]
\begin{minipage}[t]{5.5mm}
\textcolor{quarto-callout-important-color}{\faExclamation}
\end{minipage}%
\begin{minipage}[t]{\textwidth - 5.5mm}

\vspace{-3mm}\textbf{Por qué la parte principal se lleva toda la atención}\vspace{3mm}

Porque es la única que importa para casi todo. El tipo de la ecuación
---hiperbólica, parabólica o elíptica---, las curvas características, la
velocidad a la que viaja la información y si las discontinuidades
sobreviven o se borran: todo eso lo decide \(L_P\) sola. Los términos de
orden menor cambian la solución, por supuesto, pero no cambian el tipo.
Una ecuación de onda con rozamiento sigue siendo hiperbólica.

Con eso en mente, la escala de linealidad que sigue se lee de una sola
manera: \textbf{es una escala de cuánto contamina la incógnita a su
propia parte principal.}

\end{minipage}%
\end{tcolorbox}

\subsection{Lineales}\label{sec-lineales}

\begin{definition}[Ecuación
lineal]\protect\hypertarget{def-lineal}{}\label{def-lineal}

Una EDP de orden \(N\) es \textbf{lineal} si se puede escribir como

\begin{equation}\protect\phantomsection\label{eq-lineal}{
\sum_{|\alpha| \le N} a_\alpha(x)\, D^\alpha u = f(x),
}\end{equation}

es decir, si es lineal en \(u\) y en todas sus derivadas, con
coeficientes que dependen únicamente de las variables independientes.
Escribiendo \(L[u] = \sum_{|\alpha|\le N} a_\alpha(x)D^\alpha u\), la
ecuación es \(L[u] = f\), y \(L\) es un \textbf{operador lineal}:
\(L[\lambda u + \mu v] = \lambda L[u] + \mu L[v]\).

Si \(f \equiv 0\) la ecuación se dice \textbf{homogénea}; si no,
\textbf{inhomogénea} o \textbf{con término independiente}.

\end{definition}

En dos variables y orden dos, \ref{eq-lineal} se ve así:

\begin{equation}\protect\phantomsection\label{eq-lineal-2var}{
A\,u_{xx} + B\,u_{xy} + C\,u_{yy} + D\,u_x + E\,u_y + F\,u = G,
}\end{equation}

con todos los coeficientes funciones de \((x,y)\) solamente. De la
Tabla~\ref{tbl-catalogo} son lineales todas menos Burgers: calor, onda,
Laplace, Poisson, Helmholtz, Schrödinger ---que es lineal aunque sea
compleja--- y la viga.

\begin{tcolorbox}[enhanced jigsaw, arc=.35mm, bottomrule=.15mm, breakable, colback=white, colframe=quarto-callout-tip-color-frame, left=2mm, leftrule=.75mm, opacityback=0, rightrule=.15mm, toprule=.15mm]
\begin{minipage}[t]{5.5mm}
\textcolor{quarto-callout-tip-color}{\faLightbulb}
\end{minipage}%
\begin{minipage}[t]{\textwidth - 5.5mm}

\vspace{-3mm}\textbf{Por qué nos importa: superposición}\vspace{3mm}

Si \(L\) es lineal y \(L[u_1] = L[u_2] = 0\), entonces
\(L[\lambda u_1 + \mu u_2] = 0\). Suena modesto y es la propiedad más
rentable de todo el libro. \textbf{Todos} los métodos de la primera
parte son superposición disfrazada:

\begin{itemize}
\tightlist
\item
  las series de Fourier suman infinitas soluciones elementales
  (Capítulo~\ref{sec-fourier});
\item
  separación de variables fabrica esas soluciones elementales y después
  las suma (sección~\ref{sec-separacion});
\item
  la función de Green resuelve una vez para una fuente puntual y después
  integra, que es sumar infinitas soluciones con pesos
  (Capítulo~\ref{sec-green}, sección~\ref{sec-green-varias});
\item
  las transformadas integrales descomponen el dato en modos, resuelven
  modo por modo y vuelven a sumar (Capítulo~\ref{sec-tfourier},
  sección~\ref{sec-transformadas}).
\end{itemize}

Y para las inhomogéneas, \(u = u_h + u_p\): solución general de la
homogénea más una particular, exactamente como en EDO.

Fuera de las lineales \textbf{nada de esto funciona}. Es la razón de que
casi todo lo que se resuelve explícitamente en un curso sea lineal, y de
que el salto a lo no lineal sea tan brusco.

\end{minipage}%
\end{tcolorbox}

Dos pulsos viajan por una misma cuerda en sentidos opuestos. Se cruzan,
y después siguen de largo cada uno como si nada hubiera pasado. ¿De qué
propiedad de la ecuación de onda depende eso?

\begin{itemize}
\tightlist
\item
  {De que sea de segundo orden}
\item
  {De que sea lineal}
\item
  {De que sea hiperbólica}
\item
  {De que la cuerda sea perfectamente elástica}
\end{itemize}

De la linealidad. Si \(u_1\) y \(u_2\) son soluciones, \(u_1+u_2\)
también lo es, así que mientras los pulsos se superponen la altura de la
cuerda es simplemente la suma de las dos alturas, y cuando se separan
cada uno sigue con la forma que traía. \textbf{No se chocan: se suman.}

Que sea hiperbólica explica otra cosa ---por qué viajan a velocidad
finita y sin deformarse--- pero no por qué se atraviesan sin enterarse
el uno del otro. Y el orden no tiene nada que ver.

Es un experimento que se puede hacer con una soga en el piso, y vale la
pena hacerlo. Si la amplitud es grande la cosa cambia: la tensión deja
de ser constante, la ecuación se vuelve no lineal, y los pulsos empiezan
a interactuar de verdad. Eso es lo que pasa con las olas grandes en el
mar, que sí se afectan entre sí.

\subsection{Semilineales}\label{sec-semilineales}

\begin{definition}[Ecuación
semilineal]\protect\hypertarget{def-semilineal}{}\label{def-semilineal}

Una EDP de orden \(N\) es \textbf{semilineal} si su parte principal es
lineal con coeficientes que dependen solo de \(x\), y toda la no
linealidad está confinada a derivadas de orden estrictamente menor que
\(N\):

\begin{equation}\protect\phantomsection\label{eq-semilineal}{
\sum_{|\alpha| = N} a_\alpha(x)\, D^\alpha u \;+\; a_0\big(x,\,u,\,Du,\,\dots,\,D^{N-1}u\big) = 0 .
}\end{equation}

\end{definition}

En dos variables y orden dos:

\begin{equation}\protect\phantomsection\label{eq-semilineal-2var}{
A(x,y)\,u_{xx} + 2B(x,y)\,u_{xy} + C(x,y)\,u_{yy} = F(x,y,u,u_x,u_y),
}\end{equation}

con \(F\) tan fea como se quiera. Esta es exactamente la ecuación que
vamos a clasificar en la sección~\ref{sec-clasificacion}, y ahora se
entiende por qué se la elige a ella y no a la lineal: no hace falta
menos que esto, y no se puede pedir más.

\begin{example}[Tres semilineales de
verdad]\protect\hypertarget{exm-semilineales}{}\label{exm-semilineales}

~

\begin{itemize}
\tightlist
\item
  \(u_{xx} + u_{yy} = e^{u}\), la \textbf{ecuación de Liouville}, que
  aparece en geometría de superficies y en física de plasmas.
\item
  \(u_t = k\,u_{xx} + r\,u(1-u)\), la \textbf{ecuación de Fisher--KPP}:
  difusión más crecimiento logístico. Modela cómo se expande una
  población o un gen ventajoso.
\item
  \(u_{tt} - u_{xx} = -\sin u\), la \textbf{ecuación de sine--Gordon},
  que describe una fila de péndulos acoplados y tiene soluciones tipo
  solitón.
\end{itemize}

En las tres, la parte principal es la de siempre ---Laplace, calor,
onda--- y la no linealidad está toda en el término sin derivadas.

\end{example}

\begin{tcolorbox}[enhanced jigsaw, arc=.35mm, bottomrule=.15mm, breakable, colback=white, colframe=quarto-callout-tip-color-frame, left=2mm, leftrule=.75mm, opacityback=0, rightrule=.15mm, toprule=.15mm]
\begin{minipage}[t]{5.5mm}
\textcolor{quarto-callout-tip-color}{\faLightbulb}
\end{minipage}%
\begin{minipage}[t]{\textwidth - 5.5mm}

\vspace{-3mm}\textbf{Por qué nos importa: el tipo es una propiedad del punto}\vspace{3mm}

Se perdió la superposición. Pero se conservó algo casi tan valioso:
\textbf{como los coeficientes de la parte principal no dependen de
\(u\), el tipo de la ecuación se puede calcular sin resolver nada.}
Decir «esta ecuación es elíptica en la región \(x>0\)» es una afirmación
sobre la ecuación y sobre el punto, no sobre ninguna solución en
particular.

Eso es lo que va a permitir, en la sección~\ref{sec-clasificacion},
mirar \(A\), \(B\) y \(C\), calcular un discriminante y saber de
antemano si el problema se parece a una onda, a una difusión o a un
equilibrio. Con las cuasi-lineales ese privilegio se termina.

\end{minipage}%
\end{tcolorbox}

Y esta: \[x^2\,u_{xx} + \operatorname{sen}(y)\,u_{yy} + u^3 = 0.\]

\begin{itemize}
\tightlist
\item
  {Semilineal}
\item
  {Lineal}
\item
  {Cuasi-lineal pero no semilineal}
\item
  {Completamente no lineal}
\end{itemize}

Es no lineal ---\(u^3\) lo garantiza--- pero eso no es lo que se
pregunta. La parte principal es
\(x^2u_{xx} + \operatorname{sen}(y)\,u_{yy}\): lineal, y con
coeficientes que dependen solo de \(x\) e \(y\). Toda la no linealidad
está en un término sin derivadas, o sea de orden \(0 < 2\). Semilineal.

Ojo con la tercera opción: la ecuación \textbf{sí} es cuasi-lineal,
porque toda semilineal lo es ---las inclusiones de la
Figura~\ref{fig-linealidad} son inclusiones---. Lo que es falso es el
«pero no semilineal». Cuando se pide clasificar, se pide la clase
\textbf{más chica} a la que pertenece.

\subsection{Cuasi-lineales}\label{sec-cuasilineales-def}

\begin{definition}[Ecuación
cuasi-lineal]\protect\hypertarget{def-cuasilineal}{}\label{def-cuasilineal}

Una EDP de orden \(N\) es \textbf{cuasi-lineal} si es lineal en las
derivadas de orden \(N\) ---aparecen a la primera potencia y no
multiplicadas entre sí--- pero los coeficientes de esas derivadas pueden
depender de \(u\) y de sus derivadas de orden menor:

\begin{equation}\protect\phantomsection\label{eq-cuasilineal}{
\sum_{|\alpha| = N} a_\alpha\big(x,u,Du,\dots,D^{N-1}u\big)\, D^\alpha u
\;+\; a_0\big(x,u,Du,\dots,D^{N-1}u\big) = 0 .
}\end{equation}

\end{definition}

\begin{example}[Tres cuasi-lineales de
verdad]\protect\hypertarget{exm-cuasilineales}{}\label{exm-cuasilineales}

~

\begin{itemize}
\tightlist
\item
  \(u_t + u\,u_x = 0\), la \textbf{ecuación de Burgers sin viscosidad},
  que ya salió de la

  \begin{enumerate}
  \def\labelenumi{\arabic{enumi}.}
  \tightlist
  \item
    El coeficiente de \(u_x\) es la propia \(u\).
  \end{enumerate}
\item
  \((1+u_y^2)\,u_{xx} - 2u_xu_y\,u_{xy} + (1+u_x^2)\,u_{yy} = 0\), la
  \textbf{ecuación de las superficies mínimas}: la que satisface la
  forma de una película de jabón sostenida por un alambre.
\item
  \(u_t + \big(u\,v_{\max}(1-u/u_{\max})\big)_x = 0\), el \textbf{modelo
  de tráfico de Lighthill--Whitham}: \(u\) es la densidad de autos y la
  velocidad de cada auto baja a medida que la ruta se llena.
  Desarrollando la derivada queda \(u_t + q'(u)\,u_x = 0\), una ley de
  conservación como las de la 1, con la velocidad de propagación
  dependiendo de la densidad. Los embotellamientos que aparecen sin
  causa aparente son sus choques.
\end{itemize}

\end{example}

\begin{tcolorbox}[enhanced jigsaw, arc=.35mm, bottomrule=.15mm, breakable, colback=white, colframe=quarto-callout-warning-color-frame, left=2mm, leftrule=.75mm, opacityback=0, rightrule=.15mm, toprule=.15mm]
\begin{minipage}[t]{5.5mm}
\textcolor{quarto-callout-warning-color}{\faExclamationTriangle}
\end{minipage}%
\begin{minipage}[t]{\textwidth - 5.5mm}

\vspace{-3mm}\textbf{Mirá el orden antes que la apariencia}\vspace{3mm}

La ecuación de \textbf{Burgers viscosa}, \(u_t + u\,u_x = \nu\,u_{xx}\),
parece más no lineal que la de Burgers a secas, porque tiene un término
más. Y sin embargo es \textbf{semilineal}: al agregar \(\nu u_{xx}\) el
orden pasó a ser \(2\), así que la parte principal es ahora solo
\(\nu u_{xx}\) ---con coeficiente constante--- y el término \(u\,u_x\),
que es el que rompe la linealidad, quedó relegado a orden \(1\).

Moraleja: agregar un término de orden más alto puede \emph{bajar} una
ecuación en la escala de linealidad. Siempre hay que identificar primero
el orden, después la parte principal, y recién ahí clasificar.

\end{minipage}%
\end{tcolorbox}

\begin{tcolorbox}[enhanced jigsaw, arc=.35mm, bottomrule=.15mm, breakable, colback=white, colframe=quarto-callout-tip-color-frame, left=2mm, leftrule=.75mm, opacityback=0, rightrule=.15mm, toprule=.15mm]
\begin{minipage}[t]{5.5mm}
\textcolor{quarto-callout-tip-color}{\faLightbulb}
\end{minipage}%
\begin{minipage}[t]{\textwidth - 5.5mm}

\vspace{-3mm}\textbf{Por qué nos importa: el tipo pasa a depender de la solución}\vspace{3mm}

Dos consecuencias, una mala y una buena.

\textbf{La mala.} Los coeficientes de la parte principal ahora dependen
de \(u\), así que el discriminante también. La misma ecuación puede ser
elíptica donde la solución hace una cosa e hiperbólica donde hace otra,
y no se sabe cuál es cuál hasta haber resuelto ---o al menos estimado---
la solución. El Ejemplo~\ref{exm-transonico} muestra el caso clásico.

\textbf{La buena.} El método de las características \textbf{sobrevive}.
Sigue habiendo, en cada punto, una dirección privilegiada a lo largo de
la cual la ecuación se vuelve ordinaria, y sigue alcanzando con
seguirla. Es prácticamente lo único que sobrevive del arsenal, y por eso
la sección~\ref{sec-casi-lineales} le dedica una sección entera a las
cuasi-lineales de primer orden: es donde todavía se puede hacer algo
explícito sin linealidad.

Y es acá donde aparecen los choques. Si las características se cruzan
---cosa que pueden hacer, porque su dirección depende de la solución que
transportan--- la solución deja de ser una función, y hay que ampliar
qué se entiende por solución. De eso trata la
sección~\ref{sec-bien-planteado}.

\end{minipage}%
\end{tcolorbox}

\begin{example}[El tipo que cambia según la solución: flujo
transónico]\protect\hypertarget{exm-transonico}{}\label{exm-transonico}

El potencial de velocidades \(\varphi\) de un flujo estacionario,
irrotacional y compresible en dos dimensiones satisface

\[
\left(1 - \frac{\varphi_x^2}{c^2}\right)\varphi_{xx}
\;-\; \frac{2\varphi_x\varphi_y}{c^2}\,\varphi_{xy}
\;+\; \left(1 - \frac{\varphi_y^2}{c^2}\right)\varphi_{yy} = 0 ,
\]

donde \(c\) es la velocidad local del sonido. Es cuasi-lineal: las
derivadas segundas aparecen a la primera potencia, pero sus coeficientes
dependen de las derivadas primeras de la incógnita.

Calculemos el discriminante con la convención de
\ref{eq-semilineal-2var}, en la que el coeficiente del término cruzado
es \(2B\). Acá \(2B = -2\varphi_x\varphi_y/c^2\), o sea
\(B = -\varphi_x\varphi_y/c^2\), y

\[
\Delta = B^2 - AC
= \frac{\varphi_x^2\varphi_y^2}{c^4}
- \left(1-\frac{\varphi_x^2}{c^2}\right)\left(1-\frac{\varphi_y^2}{c^2}\right)
= \frac{\varphi_x^2 + \varphi_y^2}{c^2} - 1 = M^2 - 1 ,
\]

donde \(M = |\nabla\varphi|/c\) es el \textbf{número de Mach}: cuántas
veces la velocidad del sonido va el fluido en ese punto.

El resultado es imposible de mejorar. Donde el flujo es subsónico
(\(M<1\)) la ecuación es elíptica; donde es supersónico (\(M>1\)) es
hiperbólica; y la línea sónica \(M=1\) es donde se vuelve parabólica.
Que el tipo dependa de la solución no es un tecnicismo: \textbf{es} la
física del problema, y es la razón de que un ala a velocidad transónica
sea difícil de calcular. Compare esto con la ecuación de Tricomi del
Ejemplo~\ref{exm-tricomi}, que es el modelo lineal más simple con
exactamente este comportamiento ---y que se inventó justamente para
estudiarlo---.

Video recomendado: Russell Croman, \emph{Shock Wave Formation in
Transonic Flight}, \url{https://www.youtube.com/watch?v=HekbC6Pl4_Y}.

Una aclaración: \(c\) no es realmente una constante dada, porque la
velocidad del sonido local depende de \(|\nabla\varphi|\) a través de la
ecuación de Bernoulli. Tratarla como dada no cambia nada de la
conclusión sobre el tipo, pero conviene saber que es una simplificación.

\end{example}

Clasificá según su linealidad: \[u_t = \big(1 + u_x^2\big)\,u_{xx}.\]

\begin{itemize}
\tightlist
\item
  {Lineal}
\item
  {Semilineal}
\item
  {Cuasi-lineal}
\item
  {Completamente no lineal}
\end{itemize}

El orden es \(2\), así que la parte principal la forman los términos con
derivadas segundas: acá, \(\big(1+u_x^2\big)u_{xx}\). Y \(u_{xx}\)
aparece \textbf{a la primera potencia}: es cuasi-lineal.

Lo que la saca de semilineal es que el coeficiente \(\big(1+u_x^2\big)\)
depende de \(u_x\), una derivada de orden menor. Que ese coeficiente
tenga un cuadrado adentro no la vuelve completamente no lineal: el
cuadrado está sobre una derivada de orden \(1\), no sobre una de orden
\(2\). La pregunta a hacerse siempre es «¿cómo entra la derivada
\textbf{de orden más alto}?».

¿De qué orden es la \textbf{ecuación de Korteweg--de Vries},
\[u_t + \alpha\,u\,u_x + \beta\,u_{xxx} = 0\;?\] Escribí solo el número.
(Es la ecuación de las ondas solitarias en un canal poco profundo: la
primera vez que se observó un solitón, en 1834, fue en el canal de Union
en Escocia.)

Orden \(3\), por \(u_{xxx}\). Y de paso: es \textbf{cuasi-lineal}, no
completamente no lineal. La parte principal es \(\beta u_{xxx}\), que
tiene coeficiente constante; el término \(\alpha u u_x\), que es el que
rompe la linealidad, es de orden \(1 < 3\). Con el mismo criterio del
segundo ejercicio de esta sección, eso la vuelve incluso
\textbf{semilineal}.

\subsection{Completamente no lineales}\label{sec-no-lineales}

\begin{definition}[Ecuación completamente no
lineal]\protect\hypertarget{def-no-lineal}{}\label{def-no-lineal}

Una EDP es \textbf{completamente no lineal} si depende de las derivadas
de orden \(N\) de manera no lineal: alguna aparece elevada a una
potencia, dentro de una función, o multiplicada por otra derivada de
orden \(N\).

\end{definition}

\begin{example}[Tres completamente no
lineales]\protect\hypertarget{exm-no-lineales}{}\label{exm-no-lineales}

~

\begin{itemize}
\tightlist
\item
  \(u_x^2 + u_y^2 = 1\), la \textbf{ecuación eikonal}. Es la ecuación de
  los frentes de onda en óptica geométrica: sus soluciones son las
  funciones cuyo gradiente tiene módulo uno, es decir, las distancias a
  una curva. Los rayos de luz son sus características.
\item
  \(u_t + H(\nabla u,\, x,\, t) = 0\), la \textbf{ecuación de
  Hamilton--Jacobi}, que reformula toda la mecánica clásica y cuya
  versión ondulatoria es el puente histórico hacia la mecánica cuántica.
\item
  \(u_{xx}u_{yy} - u_{xy}^2 = f(x,y)\), la \textbf{ecuación de
  Monge--Ampère}, que aparece en geometría ---prescribir la curvatura de
  una superficie--- y en transporte óptimo.
\end{itemize}

\end{example}

\begin{tcolorbox}[enhanced jigsaw, arc=.35mm, bottomrule=.15mm, breakable, colback=white, colframe=quarto-callout-tip-color-frame, left=2mm, leftrule=.75mm, opacityback=0, rightrule=.15mm, toprule=.15mm]
\begin{minipage}[t]{5.5mm}
\textcolor{quarto-callout-tip-color}{\faLightbulb}
\end{minipage}%
\begin{minipage}[t]{\textwidth - 5.5mm}

\vspace{-3mm}\textbf{Por qué nos importa: hace falta más que una curva}\vspace{3mm}

Todavía hay método de características, pero con una vuelta de tuerca. Ya
no alcanza con seguir una curva y arrastrar el valor de \(u\): hay que
arrastrar además las derivadas primeras, porque la ecuación las mezcla.
El objeto que se propaga se llama \textbf{banda característica} y es una
curva con un plano tangente pegado en cada punto.

No lo vamos a desarrollar acá ---González lo menciona al pasar y remite
a la bibliografía---, pero vale saber que la eikonal y Hamilton--Jacobi
se resuelven así, y que en ambos casos las características son objetos
con nombre propio: los rayos de luz en una, las trayectorias del sistema
mecánico en la otra.

\end{minipage}%
\end{tcolorbox}

Ahora esta, que es de primer orden: \[u_t + u_x^2 = 0.\]

\begin{itemize}
\tightlist
\item
  {Cuasi-lineal, porque es lineal en \(u_t\)}
\item
  {Semilineal, porque el término no lineal no tiene derivadas de orden
  máximo}
\item
  {Completamente no lineal}
\item
  {Lineal, después de un cambio de variables}
\end{itemize}

Acá está la trampa que engancha a casi todos. La ecuación es de orden
\(1\), así que la parte principal la forman \textbf{todas} las derivadas
de orden \(1\): tanto \(u_t\) como \(u_x\). Y \(u_x^2\) no es lineal en
ellas. Completamente no lineal.

La primera opción falla porque mira solo una de las dos derivadas. La
segunda falla porque supone que \(u_x\) es «de orden menor», y no lo es:
es del mismo orden que \(u_t\). Comparala con \(u_t + u\,u_x = 0\), que
sí es cuasi-lineal: ahí el factor que multiplica a \(u_x\) es \(u\), de
orden \(0\).

Esta ecuación, de hecho, es una Hamilton--Jacobi con \(H(p) = p^2\).

\subsection{El panorama completo}\label{sec-panorama}

Las cuatro clases están encajadas una dentro de otra, y las inclusiones
son estrictas:

\[
\text{lineales} \;\subsetneq\; \text{semilineales} \;\subsetneq\; \text{cuasi-lineales}
\;\subsetneq\; \text{todas las EDP}.
\]

\begin{figure}

\centering{

\pandocbounded{\includegraphics[keepaspectratio]{08-edp_files/figure-pdf/fig-linealidad-output-1.pdf}}

}

\caption{\label{fig-linealidad}Las cuatro clases, encajadas. Cada anillo
agrega una forma de que la incógnita contamine su propia parte
principal, y cada paso hacia afuera cuesta un método. Adentro de todo
vale la superposición y con ella todo el arsenal del libro; en el
segundo anillo el tipo de la ecuación todavía se calcula sin resolverla;
en el tercero ya depende de la solución, pero las características
aguantan; afuera hace falta propagar también las derivadas.}

\end{figure}%

Y lo que se gana y se pierde en cada paso:

\begin{longtable}[]{@{}
  >{\raggedright\arraybackslash}p{(\linewidth - 8\tabcolsep) * \real{0.2600}}
  >{\centering\arraybackslash}p{(\linewidth - 8\tabcolsep) * \real{0.1600}}
  >{\raggedright\arraybackslash}p{(\linewidth - 8\tabcolsep) * \real{0.2400}}
  >{\centering\arraybackslash}p{(\linewidth - 8\tabcolsep) * \real{0.1400}}
  >{\raggedright\arraybackslash}p{(\linewidth - 8\tabcolsep) * \real{0.2000}}@{}}
\caption{Qué sobrevive en cada clase. La columna que importa es la
primera: es la que separa el libro en dos
mitades.}\label{tbl-linealidad}\tabularnewline
\toprule\noalign{}
\begin{minipage}[b]{\linewidth}\raggedright
\end{minipage} & \begin{minipage}[b]{\linewidth}\centering
Superposición
\end{minipage} & \begin{minipage}[b]{\linewidth}\raggedright
El tipo se calcula\ldots{}
\end{minipage} & \begin{minipage}[b]{\linewidth}\centering
Características
\end{minipage} & \begin{minipage}[b]{\linewidth}\raggedright
Ejemplo
\end{minipage} \\
\midrule\noalign{}
\endfirsthead
\toprule\noalign{}
\begin{minipage}[b]{\linewidth}\raggedright
\end{minipage} & \begin{minipage}[b]{\linewidth}\centering
Superposición
\end{minipage} & \begin{minipage}[b]{\linewidth}\raggedright
El tipo se calcula\ldots{}
\end{minipage} & \begin{minipage}[b]{\linewidth}\centering
Características
\end{minipage} & \begin{minipage}[b]{\linewidth}\raggedright
Ejemplo
\end{minipage} \\
\midrule\noalign{}
\endhead
\bottomrule\noalign{}
\endlastfoot
Lineal & sí & sin resolver & sí & \(u_{tt}=c^2u_{xx}\) \\
Semilineal & no & sin resolver & sí & \(u_{xx}+u_{yy}=e^u\) \\
Cuasi-lineal & no & solo con la solución & sí & \(u_t+u\,u_x=0\) \\
Completamente no lineal & no & solo con la solución & con bandas &
\(u_x^2+u_y^2=1\) \\
\end{longtable}

\begin{tcolorbox}[enhanced jigsaw, arc=.35mm, bottomrule=.15mm, breakable, colback=white, colframe=quarto-callout-warning-color-frame, left=2mm, leftrule=.75mm, opacityback=0, rightrule=.15mm, toprule=.15mm]
\begin{minipage}[t]{5.5mm}
\textcolor{quarto-callout-warning-color}{\faExclamationTriangle}
\end{minipage}%
\begin{minipage}[t]{\textwidth - 5.5mm}

\vspace{-3mm}\textbf{Cuidado con los nombres: cada libro los usa distinto}\vspace{3mm}

Esto genera más confusión que cualquier otra cosa del capítulo, así que
conviene tenerlo anotado. Nosotros seguimos la convención de González,
que es la de la cátedra.

{\def\LTcaptype{none} % do not increment counter
\begin{longtable}[]{@{}
  >{\raggedright\arraybackslash}p{(\linewidth - 6\tabcolsep) * \real{0.2500}}
  >{\raggedright\arraybackslash}p{(\linewidth - 6\tabcolsep) * \real{0.2500}}
  >{\raggedright\arraybackslash}p{(\linewidth - 6\tabcolsep) * \real{0.2500}}
  >{\raggedright\arraybackslash}p{(\linewidth - 6\tabcolsep) * \real{0.2500}}@{}}
\toprule\noalign{}
\begin{minipage}[b]{\linewidth}\raggedright
Lo que acá llamamos\ldots{}
\end{minipage} & \begin{minipage}[b]{\linewidth}\raggedright
González
\end{minipage} & \begin{minipage}[b]{\linewidth}\raggedright
Zachmanoglou (inglés)
\end{minipage} & \begin{minipage}[b]{\linewidth}\raggedright
Dallos
\end{minipage} \\
\midrule\noalign{}
\endhead
\bottomrule\noalign{}
\endlastfoot
lineal & lineal & \emph{linear} & lineal \\
semilineal & semilineal & \emph{semi-linear} & «casi lineal» (solo
primer orden) \\
cuasi-lineal & \textbf{casi--lineal} & \emph{quasi-linear} &
cuasilineal \\
completamente no lineal & completamente no lineal & \emph{fully
nonlinear} & no lineal \\
\end{longtable}
}

Los dos choques a tener presentes: González escribe
\textbf{«casi--lineal»} donde nosotros y casi todo el mundo decimos
\emph{cuasi-lineal}; y Dallos usa \textbf{«casi lineal»} para otra cosa
---para \(a(x,y)u_x + b(x,y)u_y = c(x,y,u)\), que es una semilineal de
primer orden---. Cuando leas «casi lineal» en un texto, fijate primero
de qué libro salió.

Hay además libros ---Dallos entre ellos--- que escriben la ecuación de
segundo orden con \(B\,u_{xy}\) y discriminante \(B^2-4AC\), mientras
que González escribe \(2B\,u_{xy}\) y discriminante \(B^2-AC\). Las dos
dan el mismo signo y por lo tanto el mismo tipo, pero no el mismo
número. Volvemos sobre esto en la sección~\ref{sec-clasificacion}.

\end{minipage}%
\end{tcolorbox}

\subsection*{Dos ejercicios largos}\label{dos-ejercicios-largos}
\addcontentsline{toc}{subsection}{Dos ejercicios largos}

Los ejercicios cortos de esta sección están repartidos por el texto.
Estos dos son más largos y conviene hacerlos con lápiz y papel.

\begin{exercise}[]\protect\hypertarget{exr-clasificar-linealidad}{}\label{exr-clasificar-linealidad}

Para cada una: orden, clase de linealidad (la más chica que la contenga)
y, si es lineal, si es homogénea.

\begin{enumerate}
\def\labelenumi{\alph{enumi}.}
\tightlist
\item
  \(u_{xx} + 2u_{xy} + u_{yy} = 0\)
\item
  \(u\,u_{xx} + u_{yy} = 0\)
\item
  \(u_{xx} + u_{yy} = \operatorname{sen}(xy)\)
\item
  \(u_t + u_x = u^2\)
\item
  \(u_x u_{xx} + u_{yy} = u_y\)
\item
  \(\big(u_{xx}\big)^2 + u_{yy} = 0\)
\item
  \(y\,u_{xx} - x\,u_{yy} + u_x = 0\)
\item
  \(u_t = \big(D(u)\,u_x\big)_x\), con \(D\) una función dada y
  derivable
\item
  \(u_{tt} - u_{xx} + u^3 - u = 0\)
\item
  \(u_x^2 + u_y^2 = u\)
\end{enumerate}

\end{exercise}

\begin{tcolorbox}[enhanced jigsaw, arc=.35mm, bottomrule=.15mm, bottomtitle=1mm, breakable, colback=white, colbacktitle=quarto-callout-tip-color!10!white, colframe=quarto-callout-tip-color-frame, coltitle=black, left=2mm, leftrule=.75mm, opacityback=0, opacitybacktitle=0.6, rightrule=.15mm, title=\textcolor{quarto-callout-tip-color}{\faLightbulb}\hspace{0.5em}{Solución del Ejercicio~\ref{exr-clasificar-linealidad}}, titlerule=0mm, toprule=.15mm, toptitle=1mm]

\begin{enumerate}
\def\labelenumi{\alph{enumi}.}
\tightlist
\item
  Orden 2, \textbf{lineal}, homogénea. Coeficientes constantes.
\item
  Orden 2, \textbf{cuasi-lineal}. El coeficiente de \(u_{xx}\) es \(u\),
  que es de orden \(0 < 2\), y \(u_{xx}\) aparece a la primera potencia.
\item
  Orden 2, \textbf{lineal}, inhomogénea. Que el término independiente
  sea una función complicada de \(x\) e \(y\) no afecta en nada: no
  involucra a \(u\).
\item
  Orden 1, \textbf{semilineal}. La parte principal es \(u_t + u_x\), con
  coeficientes constantes; el \(u^2\) es de orden \(0\).
\item
  Orden 2, \textbf{cuasi-lineal}. El coeficiente de \(u_{xx}\) es
  \(u_x\), de orden \(1 < 2\).
\item
  Orden 2, \textbf{completamente no lineal}. \(u_{xx}\) está elevada al
  cuadrado, y es de orden máximo.
\item
  Orden 2, \textbf{lineal}, homogénea. Coeficientes variables, pero solo
  de \(x\) e \(y\). (Compará con el Ejemplo~\ref{exm-tricomi}: esta
  también cambia de tipo según la región, y sigue siendo lineal. El tipo
  y la linealidad son cosas distintas.)
\item
  Orden 2, \textbf{cuasi-lineal}. Conviene desarrollar antes de decidir:
  \(\big(D(u)u_x\big)_x = D(u)\,u_{xx} + D'(u)\,u_x^2\). La parte
  principal es \(D(u)u_{xx}\), con coeficiente que depende de \(u\). Es
  la ecuación del calor con conductividad que depende de la temperatura,
  y es el ejemplo más común de difusión no lineal.
\item
  Orden 2, \textbf{semilineal}. Parte principal \(u_{tt}-u_{xx}\), de
  coeficientes constantes; toda la no linealidad en \(u^3-u\). Es una
  ecuación de onda no lineal muy estudiada, pariente de sine--Gordon.
\item
  Orden 1, \textbf{completamente no lineal}. Es la eikonal con el lado
  derecho cambiado. Igual que en el segundo ejercicio interactivo:
  \(u_x\) y \(u_y\) son de orden máximo, y aparecen al cuadrado.
\end{enumerate}

\end{tcolorbox}

\begin{exercise}[]\protect\hypertarget{exr-superficie-minima}{}\label{exr-superficie-minima}

La ecuación de las superficies mínimas del
Ejemplo~\ref{exm-cuasilineales} es cuasi-lineal, así que en principio su
tipo podría depender de la solución. Calculá su discriminante
\(\Delta = B^2 - AC\) ---cuidado con el factor \(2\) en el término
cruzado--- y mostrá que es siempre negativo, sea cual sea la solución.
¿Qué significa físicamente que una película de jabón satisfaga una
ecuación elíptica y no una hiperbólica?

\end{exercise}

\begin{tcolorbox}[enhanced jigsaw, arc=.35mm, bottomrule=.15mm, bottomtitle=1mm, breakable, colback=white, colbacktitle=quarto-callout-tip-color!10!white, colframe=quarto-callout-tip-color-frame, coltitle=black, left=2mm, leftrule=.75mm, opacityback=0, opacitybacktitle=0.6, rightrule=.15mm, title=\textcolor{quarto-callout-tip-color}{\faLightbulb}\hspace{0.5em}{Solución}, titlerule=0mm, toprule=.15mm, toptitle=1mm]

Con la convención de \ref{eq-general}, el término cruzado es
\(2B\,u_{xy}\), así que
\[A = 1+u_y^2, \qquad 2B = -2u_xu_y \ \Rightarrow\ B = -u_xu_y, \qquad C = 1+u_x^2 .\]
Entonces \[
\Delta = B^2 - AC = u_x^2u_y^2 - \big(1+u_y^2\big)\big(1+u_x^2\big)
= u_x^2u_y^2 - 1 - u_x^2 - u_y^2 - u_x^2u_y^2 = -\big(1 + u_x^2 + u_y^2\big),
\] que es \textbf{estrictamente negativo pase lo que pase}, porque es
menos uno menos dos cuadrados. La ecuación es elíptica para cualquier
solución. Es un caso interesante de cuasi-lineal en el que el tipo,
aunque en principio podría depender de la solución, resulta no depender.

\textbf{Lo físico.} Que sea elíptica quiere decir que no hay nada que se
propague. Y es lo correcto: una película de jabón es un problema de
\textbf{equilibrio}, no de evolución. Si movés un punto del alambre,
toda la película se reacomoda a la vez; no sale una onda desde ahí que
tarde un rato en llegar al otro extremo. La forma en cada punto depende
del alambre entero.

Si en cambio fuera hiperbólica, existirían direcciones privilegiadas
sobre la película por las que viajaría información, y una deformación
del alambre afectaría solo a la región que esas direcciones alcanzan.
Eso es lo que pasa en una membrana de tambor \textbf{vibrando} ---que sí
es hiperbólica--- y no en una película de jabón quieta.

\end{tcolorbox}

\begin{exercise}[]\protect\hypertarget{exr-mach}{}\label{exr-mach}

En el Ejemplo~\ref{exm-transonico} verificá la cuenta del discriminante
paso a paso, y después respondé: si un ala se mueve a \(M = 0{,}85\)
pero en el extremo superior del perfil el aire se acelera hasta
\(M = 1{,}1\), ¿qué tipo de ecuación gobierna cada región? ¿Por qué esto
vuelve al problema difícil de resolver numéricamente, si cada tipo por
separado se sabe tratar?

\end{exercise}

\begin{tcolorbox}[enhanced jigsaw, arc=.35mm, bottomrule=.15mm, bottomtitle=1mm, breakable, colback=white, colbacktitle=quarto-callout-tip-color!10!white, colframe=quarto-callout-tip-color-frame, coltitle=black, left=2mm, leftrule=.75mm, opacityback=0, opacitybacktitle=0.6, rightrule=.15mm, title=\textcolor{quarto-callout-tip-color}{\faLightbulb}\hspace{0.5em}{Solución}, titlerule=0mm, toprule=.15mm, toptitle=1mm]

La cuenta del discriminante está hecha en el
Ejemplo~\ref{exm-transonico}; lo que hay que verificar es que los
términos \(\varphi_x^2\varphi_y^2/c^4\) se cancelan al desarrollar el
producto, y que lo que queda es \((\varphi_x^2+\varphi_y^2)/c^2 - 1\),
que es \(M^2-1\) por definición de número de Mach.

Con eso, la respuesta es directa:

\begin{itemize}
\tightlist
\item
  donde \(M = 0{,}85\): \(\Delta = 0{,}7225 - 1 = -0{,}2775 < 0\),
  \textbf{elíptica};
\item
  donde \(M = 1{,}1\): \(\Delta = 1{,}21 - 1 = 0{,}21 > 0\),
  \textbf{hiperbólica}.
\end{itemize}

\textbf{Por qué eso es un problema serio.} Los dos tipos piden cosas
incompatibles, y las piden sobre el mismo dominio y al mismo tiempo:

\begin{itemize}
\tightlist
\item
  En la región elíptica el valor en cada punto depende de \textbf{toda}
  la frontera a la vez (Tabla~\ref{tbl-datos}). Numéricamente eso
  significa resolver un sistema global: no se puede avanzar de a poco,
  hay que resolver todo junto.
\item
  En la región hiperbólica la información viaja por características a
  velocidad finita, y el método natural es avanzar paso a paso
  siguiéndolas.
\end{itemize}

Y hay algo peor: \textbf{la frontera entre las dos regiones no se conoce
de antemano}. La línea sónica \(M=1\) es donde \(\Delta\) cambia de
signo, pero \(M\) depende de la solución, así que no se sabe dónde poner
la frontera hasta haber resuelto el problema, y no se sabe cómo
resolverlo hasta saber dónde está la frontera. Encima aparecen ondas de
choque, que son las soluciones débiles de la
sección~\ref{sec-generalizadas}.

Esa es la razón de que el régimen transónico ---el de un avión de línea
en crucero, casualmente--- haya sido durante décadas el más difícil de
la aerodinámica computacional, y de que la ecuación de Tricomi del
Ejemplo~\ref{exm-tricomi} se estudie tanto: es el modelo lineal más
simple con exactamente esta patología.

\end{tcolorbox}

\section{Qué datos hacen falta: el problema bien
planteado}\label{sec-bien-planteado}

Vimos en la sección~\ref{sec-porque} que una EDP sola no determina nada:
su indeterminación es funcional, y hace falta agregarle datos sobre
curvas o superficies enteras. Toca la segunda pregunta del capítulo:
\textbf{qué datos, y cuántos}.

La respuesta corta es incómoda. No hay un juego de condiciones que sirva
para todas las ecuaciones. Los datos que hacen que la ecuación de onda
tenga exactamente una solución arruinan a la de Laplace; los que hacen
andar a la de Laplace no le sirven a la del calor. Qué datos van con qué
ecuación es una de las cosas que decide la clasificación de la
sección~\ref{sec-clasificacion}, y es la razón por la que los capítulos
9, 10 y 11 están separados.

Pero antes hay que decir qué significa que un juego de datos «sirva».

\subsection{Los tres pedidos de Hadamard}\label{sec-hadamard-def}

\begin{definition}[Problema bien
planteado]\protect\hypertarget{def-bien-planteado}{}\label{def-bien-planteado}

Un problema formado por una ecuación en derivadas parciales en un
dominio \(D\) junto con condiciones sobre la frontera de \(D\)
---espacial, temporal, o ambas--- está \textbf{bien planteado} en el
sentido de Hadamard si cumple las tres cosas:

\begin{enumerate}
\def\labelenumi{\arabic{enumi}.}
\tightlist
\item
  \textbf{Existencia.} Hay al menos una solución.
\item
  \textbf{Unicidad.} Hay a lo sumo una solución.
\item
  \textbf{Estabilidad.} La solución depende de manera continua de los
  datos: si los datos cambian poco, la solución cambia poco.
\end{enumerate}

Si falla cualquiera de las tres, el problema está \textbf{mal
planteado}.

\end{definition}

Las dos primeras se entienden solas. La tercera es la que hay que
discutir, porque es la que más se pasa por alto y la única que tiene que
ver con que el problema sirva para algo.

\begin{tcolorbox}[enhanced jigsaw, arc=.35mm, bottomrule=.15mm, breakable, colback=white, colframe=quarto-callout-important-color-frame, left=2mm, leftrule=.75mm, opacityback=0, rightrule=.15mm, toprule=.15mm]
\begin{minipage}[t]{5.5mm}
\textcolor{quarto-callout-important-color}{\faExclamation}
\end{minipage}%
\begin{minipage}[t]{\textwidth - 5.5mm}

\vspace{-3mm}\textbf{Por qué la estabilidad no es un capricho de matemáticos}\vspace{3mm}

Los datos de un problema físico nunca son exactos. Salen de un
instrumento con error, o de un modelo que ya simplificó algo, o de una
condición inicial que se conoce solo en algunos puntos. Lo que uno
resuelve nunca es el problema verdadero: es un problema vecino.

Si el problema es estable, resolver el problema vecino da una solución
vecina, y eso es todo lo que se puede pedir. Si es inestable, la
solución del problema vecino puede no parecerse en nada a la verdadera,
y entonces \textbf{el hecho de que exista y sea única no sirve de nada}:
es una solución que no vamos a poder encontrar ni reconocer.

González lo dice desde la simetría, con un argumento que vale la pena
reproducir. El principio de razón suficiente de Leibniz asegura que
causas simétricas producen efectos simétricos, y así se usa todo el
tiempo en física teórica. Pero el principio \textbf{no} dice que causas
levemente asimétricas produzcan efectos levemente asimétricos. Eso sería
cierto solo si el problema fuera estable. Cuando no lo es, una asimetría
minúscula ---la que siempre hay--- puede producir un efecto
completamente asimétrico, y la solución simétrica que calculamos con
tanto cuidado es la que nunca se observa en el laboratorio.

\end{minipage}%
\end{tcolorbox}

Un problema tiene exactamente una solución para cada juego de datos, y
para el dato \(g \equiv 0\) la solución es \(u \equiv 0\). Alguien mide
\(g\) con un error de \(10^{-9}\) y obtiene una solución de amplitud
\(10^{6}\). ¿Qué falló?

\begin{itemize}
\tightlist
\item
  {La existencia}
\item
  {La unicidad}
\item
  {La estabilidad}
\item
  {Nada: el problema está bien planteado y el error es del instrumento}
\end{itemize}

El enunciado dice que hay solución (existencia) y que es única
(unicidad). Lo que falla es la tercera condición: un cambio de
\(10^{-9}\) en el dato produjo un cambio de \(10^{6}\) en la solución,
así que la solución no depende continuamente de los datos.

La última opción es la trampa interesante. El error del instrumento es
inevitable: siempre hay uno. Lo que hace inservible al problema no es
que exista el error sino que el problema lo amplifique. Un problema bien
planteado con el mismo instrumento habría dado una solución de amplitud
del orden de \(10^{-9}\).

\subsection{El ejemplo de Hadamard}\label{sec-hadamard}

Hadamard construyó, para mostrar que esto pasa de verdad, un problema de
apariencia completamente inofensiva. Tomó las condiciones de
Cauchy--Riemann, que cualquiera reconoce de variable compleja:

\begin{equation}\protect\phantomsection\label{eq-cauchy-riemann}{
\frac{\partial u}{\partial y} = \frac{\partial v}{\partial x}, \qquad
\frac{\partial v}{\partial y} = -\frac{\partial u}{\partial x} .
}\end{equation}

Es un sistema lineal de primer orden con coeficientes constantes. Para
cada entero \(n\), el par

\begin{equation}\protect\phantomsection\label{eq-hadamard-sol}{
u_n(x,y) = a_n\,e^{ny}\cos(nx), \qquad v_n(x,y) = a_n\,e^{ny}\operatorname{sen}(nx)
}\end{equation}

lo satisface, con \(a_n\) cualquier constante: es cuestión de derivar y
comparar. Ahora elijamos \(a_n = e^{-\sqrt{n}}\) y demos como dato lo
que valen sobre la recta \(y=0\):

\begin{equation}\protect\phantomsection\label{eq-hadamard-dato}{
u_n(x,0) = e^{-\sqrt{n}}\cos(nx), \qquad v_n(x,0) = e^{-\sqrt{n}}\operatorname{sen}(nx) .
}\end{equation}

Miremos qué pasa cuando \(n\) crece.

\textbf{Los datos se van a cero, y no de cualquier manera.} El factor
\(e^{-\sqrt{n}}\) aplasta todo:
\(e^{-\sqrt{100}} = e^{-10} \approx 4.5\times 10^{-5}\), y
\(e^{-\sqrt{400}} = e^{-20} \approx
2\times 10^{-9}\). Y no solo los datos: también todas sus derivadas,
porque derivar \(k\) veces multiplica por \(n^k\) y
\(n^k e^{-\sqrt n} \to 0\) igual. Para \(n\) grande, el dato es
indistinguible de la función nula con cualquier instrumento.

\textbf{La solución no se va a cero.} Para cualquier \(y > 0\) fijo, por
chico que sea, la amplitud de la solución es

\[
e^{-\sqrt{n}}\,e^{ny} = e^{\,ny - \sqrt{n}} \;\xrightarrow[n\to\infty]{}\; \infty ,
\]

porque \(ny\) le gana a \(\sqrt n\) apenas \(n\) pasa de \(1/y^2\). Con
\(y = 0.1\), la amplitud vale \(0.08\) para \(n=25\), vale \(1\) para
\(n=100\) y vale \(5\times 10^{8}\) para \(n=400\).

Y sin embargo, el dato \textbf{nulo} ---\(u \equiv v \equiv 0\) sobre
\(y=0\)--- tiene como solución \(u \equiv v \equiv 0\). Así que tenemos
una sucesión de datos que tiende al dato nulo cuyas soluciones no
tienden a la solución nula: no hay dependencia continua. El problema
tiene solución, la solución es única, y está mal planteado.

\begin{figure}

\centering{

\pandocbounded{\includegraphics[keepaspectratio]{08-edp_files/figure-pdf/fig-hadamard-output-1.pdf}}

}

\caption{\label{fig-hadamard}El ejemplo de Hadamard, leído como un
problema de medición. \textbf{Izquierda:} sobre la superficie accesible,
el dato verdadero es cero y lo que mide el instrumento es cero más un
ruido de amplitud \(4{,}5\times10^{-5}\) y longitud de onda
\(\lambda = 2\pi/100\). Cualquier aparato real da algo así, y a la
escala del problema es indistinguible de cero. \textbf{Derecha:} el
factor por el que ese ruido queda multiplicado en la solución, en
función de la distancia \(y\) a la superficie, para cuatro longitudes de
onda. El factor es \(e^{2\pi y/\lambda}\): crece exponencialmente al
alejarse, y tanto más rápido cuanto más fina sea la ondulación. La línea
punteada marca el factor \(1\), que sería no amplificar nada.}

\end{figure}%

\begin{tcolorbox}[enhanced jigsaw, arc=.35mm, bottomrule=.15mm, breakable, colback=white, colframe=quarto-callout-important-color-frame, left=2mm, leftrule=.75mm, opacityback=0, rightrule=.15mm, toprule=.15mm]
\begin{minipage}[t]{5.5mm}
\textcolor{quarto-callout-important-color}{\faExclamation}
\end{minipage}%
\begin{minipage}[t]{\textwidth - 5.5mm}

\vspace{-3mm}\textbf{De qué problema físico estamos hablando}\vspace{3mm}

El ejemplo puede parecer un artificio de laboratorio matemático. No lo
es: es un problema que se resuelve todos los días y que sale mal
exactamente por esta razón.

\textbf{La situación.} Uno mide un campo sobre la única superficie a la
que tiene acceso ---el potencial gravitatorio o el campo magnético al
nivel del suelo, el potencial eléctrico sobre una placa, el campo sobre
la piel de un paciente--- y quiere saber cuánto vale más adentro, donde
no puede medir. En la región sin fuentes el potencial satisface la
ecuación de Laplace, así que el problema es: \textbf{dados \(u\) y su
derivada normal sobre una superficie, hallar \(u\) del otro lado.} Eso
es exactamente el problema de Cauchy que acabamos de plantear. En
geofísica se lo llama \emph{continuación hacia abajo}.

\textbf{Lo que dice la cuenta.} Escribamos la ondulación del dato por su
longitud de onda: si el dato tiene una ondulación de longitud
\(\lambda\), entonces \(n = 2\pi/\lambda\), y el factor \(e^{ny}\) de
\ref{eq-hadamard-sol} se lee

\[
\text{factor de amplificación} \;=\; e^{\,2\pi y/\lambda} .
\]

Es decir: \textbf{el error se amplifica exponencialmente con la
distancia, y tanto más rápido cuanto más fino sea el detalle que uno
quiera resolver}. Con \(\lambda = 10\) m y \(y = 10\) m el factor es
\(e^{2\pi} \approx 535\); con \(\lambda = 1\) m y la misma profundidad,
\(e^{62{,}8} \approx
2\times10^{27}\).

\textbf{La traducción.} No se pueden resolver detalles finos a
distancia. Un instrumento con ruido de una parte en un millón, a diez
metros de profundidad, no distingue absolutamente nada de menos de unos
metros de tamaño. Y esto \textbf{no es una limitación del instrumento ni
del método numérico}: está en la ecuación, y ningún algoritmo la va a
esquivar.

Es la misma razón por la que una tomografía tiene la resolución que
tiene, por la que se puede detectar una cavidad grande y no una fisura,
y por la que existe toda la maquinaria de regularización de la que
hablamos unos párrafos más abajo.

\end{minipage}%
\end{tcolorbox}

\begin{tcolorbox}[enhanced jigsaw, arc=.35mm, bottomrule=.15mm, breakable, colback=white, colframe=quarto-callout-important-color-frame, left=2mm, leftrule=.75mm, opacityback=0, rightrule=.15mm, toprule=.15mm]
\begin{minipage}[t]{5.5mm}
\textcolor{quarto-callout-important-color}{\faExclamation}
\end{minipage}%
\begin{minipage}[t]{\textwidth - 5.5mm}

\vspace{-3mm}\textbf{Lo que el ejemplo está diciendo en realidad}\vspace{3mm}

El sistema \ref{eq-cauchy-riemann} no es cualquier sistema: si \(u\) y
\(v\) lo satisfacen, derivando una ecuación respecto de \(x\), la otra
respecto de \(y\) y sumando se obtiene

\[
u_{xx} + u_{yy} = 0 ,
\]

y lo mismo para \(v\). \textbf{El ejemplo de Hadamard es la ecuación de
Laplace disfrazada.} Y lo que mostró Hadamard no es que la ecuación de
Laplace sea mala, sino que \textbf{darle los datos equivocados la vuelve
inservible}.

Los datos que le dimos son datos de Cauchy: el valor de la función sobre
una curva abierta ---la recta \(y=0\)---, con la intención de propagar
desde ahí. Eso es lo que se hace con la ecuación de onda y funciona
perfecto. Con una ecuación elíptica no funciona nunca. La ecuación de
Laplace pide otra cosa: valores sobre una frontera \textbf{cerrada}, que
rodee al dominio entero.

Esto es lo que anticipábamos en la sección~\ref{sec-tres-mundos} cuando
decíamos que en una ecuación elíptica el valor en cada punto depende de
todo el borde a la vez. Ahora tiene consecuencias: si el valor en un
punto depende de todo el borde, no hay manera de calcularlo propagando
desde un pedazo del borde.

\end{minipage}%
\end{tcolorbox}

\begin{exercise}[]\protect\hypertarget{exr-hadamard-umbral}{}\label{exr-hadamard-umbral}

~

\begin{enumerate}
\def\labelenumi{\alph{enumi}.}
\tightlist
\item
  Para cada distancia \(y>0\) fija, hallar a partir de qué \(n\) la
  amplitud \(e^{ny-\sqrt n}\) de la solución supera a \(1\), es decir a
  partir de qué finura de detalle el ruido deja de ser despreciable.
  Comprobar que el umbral es \(n > 1/y^{2}\).
\item
  Traducirlo a longitudes de onda: ¿cuál es la \(\lambda\) más chica que
  todavía se puede recuperar a una distancia \(y\)?
\item
  Un magnetómetro mide el campo al nivel del suelo con un ruido de una
  parte en \(10^{6}\). Usando el factor \(e^{2\pi y/\lambda}\) de la
  Figura~\ref{fig-hadamard}, ¿qué tamaño mínimo de estructura se puede
  esperar distinguir a \(5\) metros de profundidad?
\end{enumerate}

\end{exercise}

\begin{tcolorbox}[enhanced jigsaw, arc=.35mm, bottomrule=.15mm, bottomtitle=1mm, breakable, colback=white, colbacktitle=quarto-callout-tip-color!10!white, colframe=quarto-callout-tip-color-frame, coltitle=black, left=2mm, leftrule=.75mm, opacityback=0, opacitybacktitle=0.6, rightrule=.15mm, title=\textcolor{quarto-callout-tip-color}{\faLightbulb}\hspace{0.5em}{Solución}, titlerule=0mm, toprule=.15mm, toptitle=1mm]

\textbf{a.} La amplitud supera a \(1\) cuando el exponente es positivo:
\[
ny - \sqrt n > 0 \iff \sqrt n\big(\sqrt n\,y - 1\big) > 0 \iff \sqrt n > \frac1y
\iff n > \frac{1}{y^{2}} .
\] Notá que \textbf{sobre la recta del dato, y solo sobre ella}, no hay
umbral: con \(y=0\) la amplitud \(e^{-\sqrt n}\) decrece siempre.
Cualquier \(y>0\), por chico que sea, tiene su umbral; solo hay que ir a
un \(n\) suficientemente grande.

\textbf{b.} Con \(n = 2\pi/\lambda\), la condición \(n > 1/y^2\) es
\(\lambda < 2\pi y^{2}\). O sea que las ondulaciones \textbf{más finas}
que ese valor son las que explotan. Al revés: para poder recuperar
detalles de tamaño \(\lambda\) a distancia \(y\), hace falta
\(\lambda \gtrsim y^2\) en las unidades en que se hizo la cuenta ---el
umbral exacto depende del nivel de ruido, que es lo que hace el punto
siguiente---.

\textbf{c.} Con ruido relativo \(\varepsilon = 10^{-6}\), la
reconstrucción es confiable mientras el factor de amplificación no
supere a \(1/\varepsilon\): \[
e^{2\pi y/\lambda} \lesssim \frac{1}{\varepsilon}
\quad\Longrightarrow\quad
\lambda \gtrsim \frac{2\pi y}{\ln(1/\varepsilon)} .
\] Con \(y = 5\) m y \(\ln(10^{6}) \approx 13{,}8\):
\[\lambda \gtrsim \frac{2\pi\cdot 5}{13{,}8} \approx 2{,}3\ \text{m}.\]

Es decir: a cinco metros de profundidad, con un instrumento excelente,
no se puede esperar distinguir estructuras de menos de un par de metros.
Y fijate en la forma de la fórmula: \(\lambda_{\min}\) es proporcional a
\(y\) y solo \textbf{logarítmica} en el ruido. Mejorar el instrumento en
un factor \(1000\) mejora la resolución en apenas un factor \(1{,}5\).
Por eso, en los problemas mal planteados, comprar un aparato mejor casi
nunca es la solución.

\end{tcolorbox}

\subsection{Las condiciones de contorno con
nombre}\label{sec-condiciones}

Conviene fijar el vocabulario de las condiciones que se usan, porque van
a aparecer en todos los capítulos que siguen. Sea \(S\) la frontera de
un dominio \(D\) y \(\hat n\) la normal exterior a \(S\).

\begin{definition}[Dirichlet, Neumann, Robin y
Cauchy]\protect\hypertarget{def-condiciones}{}\label{def-condiciones}

~

\begin{itemize}
\tightlist
\item
  \textbf{Condición de Dirichlet:} se prescribe el valor de la función
  sobre la frontera, \(u\big|_S = h\). Fijar la temperatura de los
  bordes de una placa; sostener los extremos de una cuerda.
\item
  \textbf{Condición de Neumann:} se prescribe la derivada normal,
  \(\dfrac{\partial u}{\partial n}\Big|_S = h\). Fijar el flujo de calor
  que entra; un borde aislado es \(\partial u/\partial n = 0\).
\item
  \textbf{Condición de Robin} o \textbf{mixta:} una combinación de las
  dos, \(\dfrac{\partial u}{\partial n} + \alpha u\Big|_S = h\). Es la
  condición de enfriamiento de Newton: el calor que sale por el borde es
  proporcional a la diferencia de temperatura con el ambiente.
\item
  \textbf{Condición de Cauchy:} se prescriben \textbf{a la vez} el valor
  y la derivada normal sobre una curva o superficie,
  \(u\big|_\Gamma = h_1\) y
  \(\dfrac{\partial u}{\partial n}\Big|_\Gamma = h_2\). Es la que se usa
  como condición inicial en los problemas de evolución: dar la posición
  y la velocidad iniciales de una cuerda es darle datos de Cauchy a la
  ecuación de onda sobre la recta \(t=0\).
\end{itemize}

Una condición se llama \textbf{homogénea} si \(h \equiv 0\).

\end{definition}

\begin{exercise}[]\protect\hypertarget{exr-condiciones-fisicas}{}\label{exr-condiciones-fisicas}

Traducir a una condición de Dirichlet, Neumann o Robin cada una de estas
situaciones físicas, y decir sobre qué parte de la frontera actúa:

\begin{enumerate}
\def\labelenumi{\alph{enumi}.}
\tightlist
\item
  Una barra metálica con los dos extremos sumergidos en hielo a
  \(0^\circ\)C.
\item
  Una barra metálica con un extremo perfectamente aislado.
\item
  Una placa que pierde calor por su borde hacia un ambiente a
  temperatura \(T_0\), a una tasa proporcional a la diferencia.
\item
  Una membrana de tambor sujeta a un aro rígido.
\item
  Una cuerda de la que se conoce la forma y la velocidad en el instante
  inicial.
\end{enumerate}

\end{exercise}

\begin{tcolorbox}[enhanced jigsaw, arc=.35mm, bottomrule=.15mm, bottomtitle=1mm, breakable, colback=white, colbacktitle=quarto-callout-tip-color!10!white, colframe=quarto-callout-tip-color-frame, coltitle=black, left=2mm, leftrule=.75mm, opacityback=0, opacitybacktitle=0.6, rightrule=.15mm, title=\textcolor{quarto-callout-tip-color}{\faLightbulb}\hspace{0.5em}{Solución}, titlerule=0mm, toprule=.15mm, toptitle=1mm]

\textbf{a.} Extremos en hielo a \(0^\circ\)C: la temperatura en los
bordes está fijada, así que es \textbf{Dirichlet homogénea},
\(u(0,t) = u(L,t) = 0\), sobre los dos extremos de la barra ---o sea
sobre la frontera espacial, para todo \(t\)---.

\textbf{b.} Extremo perfectamente aislado: aislado quiere decir que
\textbf{no pasa calor}, y el flujo de calor es \(-k\,u_x\) por la ley de
Fourier de la 1. Que no pase calor es \(u_x = 0\) ahí: \textbf{Neumann
homogénea}, sobre ese extremo. Es un buen ejemplo de que la condición
natural no siempre es sobre el valor: acá lo que se conoce es la
derivada.

\textbf{c.} La placa pierde calor hacia un ambiente a \(T_0\) a una tasa
proporcional a la diferencia ---es la ley de enfriamiento de Newton---.
El calor que sale por unidad de área es \(-k\,\partial u/\partial n\) y
tiene que ser igual a \(h\,(u - T_0)\), así que
\[\frac{\partial u}{\partial n} + \alpha\,u = \alpha\,T_0, \qquad \alpha = h/k,\]
una \textbf{condición de Robin inhomogénea}, sobre todo el borde de la
placa. Notá que los dos casos anteriores son los extremos de este:
\(\alpha \to \infty\) da Dirichlet (contacto perfecto con el ambiente) y
\(\alpha = 0\) da Neumann (aislamiento perfecto).

\textbf{d.} Membrana sujeta a un aro rígido: el desplazamiento es cero
en el aro, \textbf{Dirichlet homogénea} sobre la circunferencia. Es la
misma condición que la del caso \(a\) aunque la magnitud sea otra
---altura en vez de temperatura---, y por eso el problema de la membrana
y el de la barra comparten tanta matemática.

\textbf{e.} Forma y velocidad iniciales de una cuerda: son \textbf{dos}
datos sobre la misma curva, la recta \(t=0\). Eso es una
\textbf{condición de Cauchy}: \(u(x,0) = \varphi(x)\) y
\(u_t(x,0) = \psi(x)\). Y es exactamente lo que pide una ecuación
hiperbólica según la Tabla~\ref{tbl-datos}: para saber qué va a hacer
una cuerda no alcanza con verla, hay que saber además a qué velocidad se
está moviendo cada punto.

\end{tcolorbox}

Y ahora se puede decir de una vez qué va con qué. Esta tabla es, en un
renglón cada uno, el contenido de los capítulos 9, 10 y 11.

\begin{longtable}[]{@{}
  >{\raggedright\arraybackslash}p{(\linewidth - 6\tabcolsep) * \real{0.1400}}
  >{\raggedright\arraybackslash}p{(\linewidth - 6\tabcolsep) * \real{0.2000}}
  >{\raggedright\arraybackslash}p{(\linewidth - 6\tabcolsep) * \real{0.4200}}
  >{\raggedright\arraybackslash}p{(\linewidth - 6\tabcolsep) * \real{0.2400}}@{}}
\caption{Qué datos pide cada tipo. La columna de la derecha es la que
hay que mirar: son los tres errores clásicos, y el de la fila de abajo
es el ejemplo de Hadamard.}\label{tbl-datos}\tabularnewline
\toprule\noalign{}
\begin{minipage}[b]{\linewidth}\raggedright
Tipo
\end{minipage} & \begin{minipage}[b]{\linewidth}\raggedright
Modelo
\end{minipage} & \begin{minipage}[b]{\linewidth}\raggedright
Datos que la hacen bien planteada
\end{minipage} & \begin{minipage}[b]{\linewidth}\raggedright
Datos que la arruinan
\end{minipage} \\
\midrule\noalign{}
\endfirsthead
\toprule\noalign{}
\begin{minipage}[b]{\linewidth}\raggedright
Tipo
\end{minipage} & \begin{minipage}[b]{\linewidth}\raggedright
Modelo
\end{minipage} & \begin{minipage}[b]{\linewidth}\raggedright
Datos que la hacen bien planteada
\end{minipage} & \begin{minipage}[b]{\linewidth}\raggedright
Datos que la arruinan
\end{minipage} \\
\midrule\noalign{}
\endhead
\bottomrule\noalign{}
\endlastfoot
Hiperbólica & \(u_{tt}=c^2u_{xx}\) & Cauchy en \(t=0\) (\(u\) y
\(u_t\)), más contorno en los extremos si el dominio es acotado & --- \\
Parabólica & \(u_t=k\,u_{xx}\) & \(u\) en \(t=0\) más contorno,
\textbf{hacia adelante} en el tiempo & hacia atrás en el tiempo \\
Elíptica & \(u_{xx}+u_{yy}=0\) & Dirichlet, Neumann o Robin sobre
\textbf{toda} la frontera cerrada & Cauchy sobre una curva abierta \\
\end{longtable}

\begin{tcolorbox}[enhanced jigsaw, arc=.35mm, bottomrule=.15mm, breakable, colback=white, colframe=quarto-callout-note-color-frame, left=2mm, leftrule=.75mm, opacityback=0, rightrule=.15mm, toprule=.15mm]
\begin{minipage}[t]{5.5mm}
\textcolor{quarto-callout-note-color}{\faInfo}
\end{minipage}%
\begin{minipage}[t]{\textwidth - 5.5mm}

\vspace{-3mm}\textbf{La fila del medio: problemas inversos}\vspace{3mm}

Que la ecuación del calor esté mal planteada hacia atrás en el tiempo no
es una curiosidad. Es el enunciado matemático de que la difusión borra
información: dos distribuciones de temperatura muy distintas, al cabo de
un rato, se vuelven casi idénticas ---lo vimos en la
Figura~\ref{fig-tres-mundos}---. Reconstruir el pasado es entonces
separar cosas que quedaron casi iguales, y cualquier error de medición
se amplifica sin control.

Toda una rama, la de los \textbf{problemas inversos}, se dedica a esto:
recuperar una causa a partir de un efecto medido, cuando el problema
directo es estable y el inverso no. Aparece en tomografía, en
prospección sísmica, en procesamiento de imágenes. La estrategia general
es renunciar a resolver el problema exacto y resolver en cambio uno
cercano que sí sea estable, agregando información previa sobre cómo
debería ser la solución; eso se llama \textbf{regularización}. Nada de
esto entra en el programa, pero conviene saber que la fila del medio de
la Tabla~\ref{tbl-datos} es la puerta de entrada a un área entera.

\end{minipage}%
\end{tcolorbox}

¿Cuál de estos problemas está \textbf{mal} planteado?

\begin{itemize}
\tightlist
\item
  {La ecuación de onda en una cuerda, dados la posición y la velocidad
  iniciales}
\item
  {La ecuación de Laplace en un disco, dado el valor de \(u\) en toda la
  circunferencia}
\item
  {La ecuación del calor en una barra, dada la temperatura de hoy, para
  averiguar la de hace una hora}
\item
  {La ecuación del calor en una barra, dada la temperatura inicial y la
  de los dos extremos}
\end{itemize}

Es la fila del medio de la Tabla~\ref{tbl-datos} leída de derecha a
izquierda. La ecuación del calor está bien planteada hacia adelante en
el tiempo y mal planteada hacia atrás, porque la difusión borra las
diferencias: al cabo de una hora, dos temperaturas iniciales muy
distintas producen temperaturas casi iguales, y entonces retroceder
exige separar cosas casi iguales.

Las otras tres son los tres casos estándar. La primera es Cauchy para
una hiperbólica; la segunda es Dirichlet sobre una frontera cerrada para
una elíptica; la cuarta es dato inicial más contorno para una
parabólica, hacia adelante.

\subsection{Cuándo una solución no es
derivable}\label{sec-generalizadas}

Falta una discusión más, y es sobre qué se acepta como solución.

La Definición~\ref{def-solucion-clasica} pide que una solución de una
ecuación de orden \(N\) sea \(N\) veces continuamente derivable. Es lo
natural: si la ecuación tiene \(u_{xx}\), parece razonable pedir que
\(u_{xx}\) exista. Pero esa exigencia deja afuera fenómenos físicos
perfectamente reales.

Volvamos a la ley de conservación de la 1,

\begin{equation}\protect\phantomsection\label{eq-conservacion-F}{
u_t + F(u)_x = 0 ,
}\end{equation}

y recordemos de dónde salió: de la afirmación integral

\begin{equation}\protect\phantomsection\label{eq-conservacion-integral}{
\frac{d}{dt}\int_a^b u\,dx = F\big(u(a,t)\big) - F\big(u(b,t)\big) ,
}\end{equation}

que vale para todo tramo \([a,b]\). Ahora bien:
\textbf{\ref{eq-conservacion-integral} tiene sentido aunque \(u\) tenga
un salto}, porque solo pide que \(u\) sea integrable. Fue en el paso de
\ref{eq-conservacion-integral} a \ref{eq-conservacion-F} donde metimos
derivadas que quizás no existan.

Y la que es físicamente cierta es la primera. La ecuación diferencial es
un producto derivado: salió de suponer que la solución era suave para
poder escribir el balance punto a punto. Si la solución desarrolla una
discontinuidad ---y en la sección~\ref{sec-casi-lineales} vamos a ver
que la ecuación de Burgers lo hace sola, sin ninguna ayuda--- no es la
física la que falla, es nuestra forma de escribirla.

La salida es correr las derivadas de lugar.

Video recomendado: SmarterEveryDay, \emph{Shockwave Shadows in Ultra
Slow Motion (Bullet Schlieren)},
\url{https://www.youtube.com/watch?v=BPwdlEgLn5Q}.

\begin{definition}[Solución débil de una ley de
conservación]\protect\hypertarget{def-solucion-debil}{}\label{def-solucion-debil}

Sea \(g\) el dato inicial. Una función acotada \(u(x,t)\) es
\textbf{solución débil} o \textbf{generalizada} de
\ref{eq-conservacion-F} con \(u(x,0)=g(x)\) si para toda función de
prueba \(\varphi\) infinitamente derivable y de soporte compacto en
\(\{t \ge 0\}\) se cumple

\begin{equation}\protect\phantomsection\label{eq-solucion-debil}{
\int_0^{\infty}\!\!\int_{-\infty}^{\infty}
\Big( u\,\varphi_t + F(u)\,\varphi_x \Big)\,dx\,dt
\;+\; \int_{-\infty}^{\infty} g(x)\,\varphi(x,0)\,dx \;=\; 0 .
}\end{equation}

\end{definition}

La definición se obtiene multiplicando \ref{eq-conservacion-F} por
\(\varphi\), integrando e integrando por partes para pasarle
\textbf{todas} las derivadas a \(\varphi\). Después de eso, en la
expresión \ref{eq-solucion-debil} no queda ninguna derivada de \(u\): la
función incógnita puede saltar todo lo que quiera. Y toda solución
clásica es débil ---basta deshacer la integración por partes---, así que
no perdimos nada: ampliamos.

\begin{tcolorbox}[enhanced jigsaw, arc=.35mm, bottomrule=.15mm, breakable, colback=white, colframe=quarto-callout-tip-color-frame, left=2mm, leftrule=.75mm, opacityback=0, rightrule=.15mm, toprule=.15mm]
\begin{minipage}[t]{5.5mm}
\textcolor{quarto-callout-tip-color}{\faLightbulb}
\end{minipage}%
\begin{minipage}[t]{\textwidth - 5.5mm}

\vspace{-3mm}\textbf{Esto ya lo vimos}\vspace{3mm}

Pasarle las derivadas a una función de prueba suave para poder derivar
cosas que no son derivables es exactamente la idea de la
Capítulo~\ref{sec-distribuciones}. La delta de Dirac no es una función y
sin embargo se deriva sin problema, porque lo que se define es cómo
actúa contra una función de prueba. Acá se hace lo mismo con la solución
de una EDP en vez de con una función aislada, y por eso la teoría de
distribuciones y las EDP no lineales crecieron juntas.

\end{minipage}%
\end{tcolorbox}

\begin{tcolorbox}[enhanced jigsaw, arc=.35mm, bottomrule=.15mm, bottomtitle=1mm, breakable, colback=white, colbacktitle=quarto-callout-tip-color!10!white, colframe=quarto-callout-tip-color-frame, coltitle=black, left=2mm, leftrule=.75mm, opacityback=0, opacitybacktitle=0.6, rightrule=.15mm, title=\textcolor{quarto-callout-tip-color}{\faLightbulb}\hspace{0.5em}{Para el que quiera seguir: a qué velocidad viaja un choque}, titlerule=0mm, toprule=.15mm, toptitle=1mm]

Esto va más allá del programa, pero es una consecuencia bonita y directa
de \ref{eq-conservacion-integral}, y contesta la pregunta obvia de «¿y
entonces qué pasa después de que la solución se rompe?».

Supongamos que la solución tiene una discontinuidad que se mueve según
una curva \(x = s(t)\), con valores \(u_-\) a la izquierda y \(u_+\) a
la derecha. Apliquemos \ref{eq-conservacion-integral} a un tramo
\([a,b]\) que contenga al salto. La integral de la izquierda se parte en
dos pedazos, y al derivar respecto de \(t\) aparecen los términos de
borde móvil:

\[
\frac{d}{dt}\left(\int_a^{s(t)} u\,dx + \int_{s(t)}^b u\,dx\right)
= \dot s\,(u_- - u_+) + (\text{términos que se van al achicar el tramo}) .
\]

Igualando con \(F(u_-)-F(u_+)\) y achicando \([a,b]\) alrededor del
salto queda la \textbf{condición de Rankine--Hugoniot}:

\[
\dot s = \frac{F(u_+) - F(u_-)}{u_+ - u_-} .
\]

La velocidad del choque es el cociente entre el salto del flujo y el
salto de la densidad. Para Burgers, con \(F(u)=u^2/2\), da
\(\dot s = (u_+ + u_-)/2\): el choque viaja al promedio de las dos
velocidades que lo forman.

Una advertencia: las soluciones débiles, así definidas, \textbf{no son
únicas}. Hace falta un criterio adicional para quedarse con la
físicamente correcta, y en dinámica de gases ese criterio es que la
entropía aumente al cruzar la discontinuidad. Por eso se lo llama
condición de entropía. Zachmanoglou y Thoe lo discuten en el capítulo
III; Evans, con todo detalle, en el capítulo 3.

\end{tcolorbox}

\subsection*{Un ejercicio largo}\label{un-ejercicio-largo}
\addcontentsline{toc}{subsection}{Un ejercicio largo}

\begin{exercise}[]\protect\hypertarget{exr-debil-escalon}{}\label{exr-debil-escalon}

Sea \(F(u) = u^2/2\), o sea la ecuación de Burgers, y consideremos el
dato inicial escalón
\[g(x) = \begin{cases} 1 & x < 0,\\ 0 & x \ge 0.\end{cases}\] Usando la
condición de Rankine--Hugoniot, hallar la velocidad a la que se mueve el
salto y escribir la solución débil resultante. Verificar que \textbf{no}
es una solución clásica en ningún entorno de la recta del choque.

\end{exercise}

\begin{tcolorbox}[enhanced jigsaw, arc=.35mm, bottomrule=.15mm, bottomtitle=1mm, breakable, colback=white, colbacktitle=quarto-callout-tip-color!10!white, colframe=quarto-callout-tip-color-frame, coltitle=black, left=2mm, leftrule=.75mm, opacityback=0, opacitybacktitle=0.6, rightrule=.15mm, title=\textcolor{quarto-callout-tip-color}{\faLightbulb}\hspace{0.5em}{Solución}, titlerule=0mm, toprule=.15mm, toptitle=1mm]

Con \(F(u) = u^2/2\), y llamando \(u_- = 1\) al valor de la izquierda y
\(u_+ = 0\) al de la derecha, la condición de Rankine--Hugoniot da \[
\dot s = \frac{F(u_+) - F(u_-)}{u_+ - u_-} = \frac{0 - \tfrac12}{0 - 1} = \frac12 .
\] El salto arranca en \(x=0\) y viaja hacia la derecha a velocidad
\(1/2\), o sea \(s(t) = t/2\), y la solución débil es \[
u(x,t) = \begin{cases} 1 & x < t/2,\\ 0 & x > t/2.\end{cases}
\] Como era de esperar por la fórmula general de Burgers, la velocidad
del choque es el promedio de los dos valores: \((1+0)/2 = 1/2\).

\textbf{Por qué no es clásica.} En cualquier entorno de la recta
\(x = t/2\) la función tiene un salto, así que ni siquiera es continua,
y mucho menos derivable. No satisface la ecuación \(u_t + u\,u_x = 0\)
en esos puntos, porque las derivadas no existen. Sí satisface la
formulación integral \ref{eq-solucion-debil}, que es la que tiene
sentido físico.

\textbf{La lectura con características.} Del lado izquierdo, \(u=1\),
las características tienen pendiente \(dx/dt = 1\); del derecho,
\(u=0\), son verticales. Las de la izquierda alcanzan a las de la
derecha \textbf{de entrada}: acá \(t^{*} = 0\), el choque existe desde
el primer instante. Es el caso opuesto al del Ejemplo~\ref{exm-burgers},
donde el dato era suave y había que esperar hasta \(t^{*}\).

\end{tcolorbox}

\begin{exercise}[]\protect\hypertarget{exr-hadamard-verificar}{}\label{exr-hadamard-verificar}

Verificar que el par \ref{eq-hadamard-sol} satisface el sistema
\ref{eq-cauchy-riemann} para cualquier constante \(a_n\), y que \(u_n\)
y \(v_n\) son funciones armónicas. Después mostrar que no solo el dato
\ref{eq-hadamard-dato} tiende a cero cuando \(n \to \infty\), sino
también todas sus derivadas:
\(\max_x \big|\partial_x^k u_n(x,0)\big| = n^k e^{-\sqrt n} \to 0\) para
cada \(k\) fijo. ¿Por qué es importante que sea así, y no solo que el
dato tienda a cero?

\end{exercise}

\begin{tcolorbox}[enhanced jigsaw, arc=.35mm, bottomrule=.15mm, bottomtitle=1mm, breakable, colback=white, colbacktitle=quarto-callout-tip-color!10!white, colframe=quarto-callout-tip-color-frame, coltitle=black, left=2mm, leftrule=.75mm, opacityback=0, opacitybacktitle=0.6, rightrule=.15mm, title=\textcolor{quarto-callout-tip-color}{\faLightbulb}\hspace{0.5em}{Solución}, titlerule=0mm, toprule=.15mm, toptitle=1mm]

\textbf{Cauchy--Riemann.} Con \(u_n = a_n e^{ny}\cos(nx)\) y
\(v_n = a_n e^{ny}\operatorname{sen}(nx)\): \[
u_y = a_n n\,e^{ny}\cos(nx) = v_x \ \checkmark, \qquad
v_y = a_n n\,e^{ny}\operatorname{sen}(nx) = -u_x \ \checkmark,
\] porque \(u_x = -a_n n\,e^{ny}\operatorname{sen}(nx)\). La constante
\(a_n\) sale de factor y no interviene.

\textbf{Armónicas.} \(u_{xx} = -a_n n^2 e^{ny}\cos(nx)\) y
\(u_{yy} = +a_n n^2 e^{ny}\cos(nx)\), así que \(u_{xx}+u_{yy}=0\). Lo
mismo para \(v\). (Es el hecho general de que la parte real y la parte
imaginaria de una función holomorfa son armónicas.)

\textbf{Las derivadas.} Derivar \(k\) veces respecto de \(x\) multiplica
por \(n^k\) y cambia el seno por el coseno, así que el máximo del módulo
es \(n^k e^{-\sqrt n}\). Llamando \(s = \sqrt n\), eso es
\(s^{2k}e^{-s}\), y la exponencial le gana a cualquier potencia: tiende
a cero.

\textbf{Por qué importa.} Sin eso el ejemplo no probaría nada. Una
función chiquita pero que oscila rapidísimo tiene derivadas enormes, y
siempre se podría contestar que el dato «en realidad no era chico», que
lo que hay que mirar es alguna norma más exigente que incluya las
derivadas. Al mostrar que \textbf{todas} las derivadas tienden a cero,
Hadamard cierra esa salida: el dato es chico en cualquier sentido
razonable de la palabra, y aun así la solución explota. La inestabilidad
no es un problema de haber elegido mal cómo medir.

\end{tcolorbox}

\section{Primer orden: el método de las
características}\label{sec-caracteristicas}

Este es el método más importante del capítulo, y probablemente de la
materia. Vale la pena decir de entrada por qué.

Es el único de los cuatro que \textbf{no necesita linealidad}.
Superposición, transformadas, función de Green y separación de variables
se caen todos apenas la ecuación deja de ser lineal, como vimos en la
Tabla~\ref{tbl-linealidad}. Las características no: siguen funcionando
en las cuasi-lineales, y con más trabajo hasta en las completamente no
lineales. Además es el que explica \emph{por qué} la clasificación de la
sección~\ref{sec-clasificacion} es la que es ---hiperbólicas,
parabólicas y elípticas se distinguen por cuántas características reales
tienen--- y el que va a resolver la ecuación de onda en el capítulo 9.

La idea de fondo se puede decir en una línea: \textbf{convertir la
ecuación en derivadas parciales en un sistema de ecuaciones ordinarias,
eligiendo bien por dónde caminar.} El resto de la sección es entender
qué quiere decir «por dónde».

\subsection{La solución es una superficie}\label{sec-superficie}

Primero hay que cambiar la manera de mirar una solución.

Cuando escribimos \(u = u(x,y)\) solemos pensar en una fórmula. Conviene
pensar en un dibujo: el \textbf{gráfico} de \(u\), que es la superficie
de \(\mathbb{R}^3\) formada por los puntos \((x, y, u(x,y))\). El plano
\((x,y)\) es el \textbf{piso} ---el dominio, donde vive el problema--- y
la superficie es un \textbf{techo} que flota sobre él, a una altura
distinta sobre cada punto.

\begin{definition}[Superficie
solución]\protect\hypertarget{def-superficie-solucion}{}\label{def-superficie-solucion}

Dada una EDP de primer orden en dos variables, una \textbf{superficie
solución} es el gráfico \[
S = \{(x,y,u(x,y)) : (x,y) \in D\}
\] de una solución \(u\) definida en un dominio \(D\) del plano.

Como \(S\) es la superficie de nivel cero de
\(\Phi(x,y,z) = z - u(x,y)\), un vector normal a \(S\) en cada punto es
\[
\vec n = \nabla\Phi = \big(-u_x,\, -u_y,\, 1\big),
\] o, cambiándole el signo, \(\big(u_x, u_y, -1\big)\).

\end{definition}

Vamos a trabajar con la ecuación cuasi-lineal general de primer orden en
dos variables,

\begin{equation}\protect\phantomsection\label{eq-cuasilineal-1}{
a(x,y,u)\,u_x + b(x,y,u)\,u_y = c(x,y,u),
}\end{equation}

que incluye como casos particulares todo lo que vamos a ver: si \(a\),
\(b\), \(c\) no dependen de \(u\) es semilineal, y si además \(c\) es
lineal en \(u\) es lineal.

Y ahora una observación que va a organizar todo. Armemos con los
coeficientes el campo vectorial

\begin{equation}\protect\phantomsection\label{eq-campo}{
\vec V(x,y,u) = \big(a(x,y,u),\; b(x,y,u),\; c(x,y,u)\big) ,
}\end{equation}

que en cada punto del espacio \((x,y,u)\) señala una dirección. Entonces

\[
\vec V \cdot \big(u_x, u_y, -1\big) = a\,u_x + b\,u_y - c ,
\]

y \textbf{la ecuación \ref{eq-cuasilineal-1} dice exactamente que ese
producto escalar es cero.}

\begin{tcolorbox}[enhanced jigsaw, arc=.35mm, bottomrule=.15mm, breakable, colback=white, colframe=quarto-callout-important-color-frame, left=2mm, leftrule=.75mm, opacityback=0, rightrule=.15mm, toprule=.15mm]
\begin{minipage}[t]{5.5mm}
\textcolor{quarto-callout-important-color}{\faExclamation}
\end{minipage}%
\begin{minipage}[t]{\textwidth - 5.5mm}

\vspace{-3mm}\textbf{La ecuación, en una frase geométrica}\vspace{3mm}

\[
a\,u_x + b\,u_y = c
\quad\Longleftrightarrow\quad
\vec V \perp \vec n
\quad\Longleftrightarrow\quad
\textbf{el campo } \vec V \textbf{ es tangente a la superficie solución.}
\]

Eso es todo lo que dice la ecuación. No pide nada más. Resolver la EDP
es encontrar superficies a las que el campo \(\vec V\) les sea tangente
en cada uno de sus puntos.

Y encontrar curvas tangentes a un campo vectorial es un problema de
ecuaciones \textbf{ordinarias} que ya sabemos resolver. La pregunta pasa
a ser cómo se arma una superficie con curvas.

\end{minipage}%
\end{tcolorbox}

\subsection{El observador y el hilo que lleva
encima}\label{sec-observador}

Acá está el corazón del método, y conviene contarlo despacio porque es
una imagen que sirve para todo el resto de la materia.

Imaginemos un \textbf{observador que camina por el piso}, es decir por
el plano \((x,y)\), siguiendo alguna trayectoria
\(\big(x(s),\, y(s)\big)\) que todavía no elegimos. Mientras camina va
mirando hacia arriba y anotando la altura del techo justo encima suyo:

\[
\varphi(s) = u\big(x(s),\, y(s)\big) .
\]

¿Cómo cambia lo que anota? Por la regla de la cadena,

\begin{equation}\protect\phantomsection\label{eq-cadena-observador}{
\varphi'(s) = \frac{dx}{ds}\,u_x + \frac{dy}{ds}\,u_y .
}\end{equation}

Comparemos \ref{eq-cadena-observador} con la ecuación
\ref{eq-cuasilineal-1}. Se parecen muchísimo: las dos son una
combinación de \(u_x\) y \(u_y\). La única diferencia es que en una los
coeficientes son \(a\) y \(b\), y en la otra son las componentes de la
velocidad del observador, que \textbf{las elegimos nosotros}.

Elijámoslas iguales. Si el observador camina de modo que

\[
\frac{dx}{ds} = a(x,y,u), \qquad \frac{dy}{ds} = b(x,y,u),
\]

entonces \ref{eq-cadena-observador} se vuelve
\(\varphi'(s) = a\,u_x + b\,u_y\), que por la ecuación vale \(c\). Y ya
no queda ninguna derivada parcial:

\begin{equation}\protect\phantomsection\label{eq-sistema-caracteristico}{
\boxed{\;
\frac{dx}{ds} = a(x,y,u), \qquad
\frac{dy}{ds} = b(x,y,u), \qquad
\frac{du}{ds} = c(x,y,u).
\;}
}\end{equation}

Tres ecuaciones diferenciales \textbf{ordinarias}. La EDP desapareció.

\begin{definition}[Sistema característico, curva característica y
característica
proyectada]\protect\hypertarget{def-sistema-caracteristico}{}\label{def-sistema-caracteristico}

Al sistema anterior se lo llama \textbf{sistema característico} de la
ecuación \ref{eq-cuasilineal-1}. También se lo escribe, eliminando el
parámetro, como

\[
\frac{dx}{a} = \frac{dy}{b} = \frac{du}{c} .
\]

Sus soluciones \(\big(x(s),y(s),u(s)\big)\) son curvas en el espacio
\((x,y,u)\) y se llaman \textbf{curvas características} de la ecuación:
son, por construcción, las curvas integrales del campo \(\vec V\) de
\ref{eq-campo}.

La sombra de una curva característica sobre el plano \((x,y)\) ---o sea
el par \(\big(x(s),y(s)\big)\) sin la altura--- se llama
\textbf{característica proyectada}, o característica base. Es la
trayectoria del observador.

\end{definition}

\begin{tcolorbox}[enhanced jigsaw, arc=.35mm, bottomrule=.15mm, breakable, colback=white, colframe=quarto-callout-warning-color-frame, left=2mm, leftrule=.75mm, opacityback=0, rightrule=.15mm, toprule=.15mm]
\begin{minipage}[t]{5.5mm}
\textcolor{quarto-callout-warning-color}{\faExclamationTriangle}
\end{minipage}%
\begin{minipage}[t]{\textwidth - 5.5mm}

\vspace{-3mm}\textbf{Dos objetos, un solo nombre}\vspace{3mm}

Casi todos los textos ---González entre ellos--- llaman
«características» a las dos cosas y dejan que el contexto desambigüe.
Conviene tener presente cuál es cuál, porque viven en espacios
distintos: la característica \textbf{proyectada} es una curva en el
plano del dominio, la característica a secas es una curva en el espacio
y está apoyada sobre la superficie solución.

En el caso lineal y semilineal, \(a\) y \(b\) no dependen de \(u\), así
que las dos primeras ecuaciones del sistema característico se resuelven
solas: las características proyectadas se pueden dibujar \textbf{antes}
de conocer la solución. Ese es un privilegio que se pierde en el caso
cuasi-lineal, donde no se sabe por dónde camina el observador hasta no
saber qué altura está midiendo.

\end{minipage}%
\end{tcolorbox}

Ahora el nombre que le da el título a la sección. Mientras el observador
camina, el punto del techo que tiene encima también se mueve: describe
la curva \(\big(x(s), y(s), \varphi(s)\big)\) en el espacio. Llamemos a
esa curva \textbf{el hilo} del observador. Tres cosas sobre ella:

\begin{enumerate}
\def\labelenumi{\arabic{enumi}.}
\tightlist
\item
  \textbf{El hilo está apoyado sobre la superficie solución.} Por
  definición, su tercera coordenada es la altura del techo sobre el
  piso.
\item
  \textbf{El hilo es una curva característica.} Sus tres coordenadas
  cumplen el sistema característico.
\item
  \textbf{El hilo se calcula sin conocer la superficie.} El sistema
  característico solo usa los coeficientes de la ecuación. Si sabemos
  dónde arranca el hilo ---un punto del piso y la altura del techo
  ahí--- podemos seguirlo hasta el final resolviendo tres EDOs.
\end{enumerate}

El punto 3 es el que convierte a esto en un método. Conocer un solo
valor de \(u\), en un solo punto, nos deja conocerla sobre toda una
curva.

Un observador camina por el plano \((x,y)\) y mide la altura de la
superficie solución de \[2u_x + 3u_y = 0.\] ¿En qué dirección tiene que
caminar para que lo que mide no cambie nunca?

\begin{itemize}
\tightlist
\item
  {En la dirección \((3, 2)\)}
\item
  {En la dirección \((2, 3)\)}
\item
  {En la dirección \((-3, 2)\), perpendicular al vector de coeficientes}
\item
  {En cualquier dirección: la ecuación es homogénea}
\end{itemize}

El sistema característico es \(\dot x = 2\), \(\dot y = 3\), así que la
velocidad del observador tiene que ser \((2,3)\): \textbf{las mismas
componentes que los coeficientes}, en el mismo orden.

La primera opción los intercambia, que es el error más frecuente. La
tercera confunde la dirección característica con su perpendicular:
caminar perpendicular al campo es justamente la dirección en la que
\(u\) cambia \textbf{más} rápido, no menos. Y la cuarta confunde
«homogénea» con «constante».

Comprobación: sobre esa dirección, \(\varphi' = 2u_x + 3u_y = 0\). Las
características proyectadas son las rectas \(3x - 2y = \text{const}\) y
la solución general es \(u = F(3x-2y)\).

\subsection{El campo que guía}\label{sec-campo-guia}

Vale la pena frenar acá, porque el objeto que acabamos de usar es el que
hace funcionar todo el método y el que va a reaparecer durante el resto
de la materia.

Hasta recién, \(a\), \(b\) y \(c\) eran «los coeficientes de la
ecuación»: tres funciones sueltas. A partir de ahora son otra cosa. Son
\textbf{un campo vectorial}, y resolver la ecuación es seguir sus
flechas.

\subsubsection{Primero, el campo en el piso}\label{sec-campo-piso}

\begin{definition}[Campo vectorial y curva
integral]\protect\hypertarget{def-campo-integral}{}\label{def-campo-integral}

Un \textbf{campo vectorial} en un dominio \(D \subseteq \mathbb{R}^2\)
es una función que a cada punto le asigna un vector, \[
\vec v : D \to \mathbb{R}^2, \qquad \vec v(x,y) = \big(v_1(x,y),\, v_2(x,y)\big) .
\]

Una \textbf{curva integral} del campo es una curva
\(s \mapsto \big(x(s), y(s)\big)\) que en cada punto avanza en la
dirección que marca el campo, y con esa misma rapidez:
\begin{equation}\protect\phantomsection\label{eq-curva-integral}{
\frac{dx}{ds} = v_1\big(x(s), y(s)\big), \qquad
\frac{dy}{ds} = v_2\big(x(s), y(s)\big) .
}\end{equation}

A la familia de todas las curvas integrales se la llama \textbf{flujo}
del campo, y a su dibujo, \textbf{retrato de fases}.

\end{definition}

Ahora mirá las dos primeras ecuaciones del sistema característico y
compará con \ref{eq-curva-integral}. Son la misma cosa. En el caso
lineal y semilineal, donde \(a\) y \(b\) dependen solo de \((x,y)\), el
par

\begin{equation}\protect\phantomsection\label{eq-campo-piso}{
\vec v(x,y) = \big(a(x,y),\, b(x,y)\big)
}\end{equation}

es un campo vectorial honesto sobre el piso, y:

\begin{tcolorbox}[enhanced jigsaw, arc=.35mm, bottomrule=.15mm, breakable, colback=white, colframe=quarto-callout-important-color-frame, left=2mm, leftrule=.75mm, opacityback=0, rightrule=.15mm, toprule=.15mm]
\begin{minipage}[t]{5.5mm}
\textcolor{quarto-callout-important-color}{\faExclamation}
\end{minipage}%
\begin{minipage}[t]{\textwidth - 5.5mm}

\vspace{-3mm}\textbf{La característica proyectada es una curva integral}\vspace{3mm}

La trayectoria del observador no es una elección nuestra: es una
\textbf{curva integral del campo \(\vec v = (a,b)\)}. La ecuación ya
eligió por dónde hay que caminar; nosotros solo la obedecemos.

Dibujar el campo es, entonces, dibujar de antemano todas las rutas
posibles. Antes de resolver nada, antes de saber cuánto vale \(u\) en
ningún lado, ya sabemos por dónde va a viajar la información.

\end{minipage}%
\end{tcolorbox}

Y hay una segunda lectura, todavía más corta. El miembro izquierdo de la
ecuación es

\begin{equation}\protect\phantomsection\label{eq-derivada-direccional}{
a\,u_x + b\,u_y \;=\; \vec v \cdot \nabla u ,
}\end{equation}

que es la \textbf{derivada direccional} de \(u\) en la dirección de
\(\vec v\) (sin normalizar: pesada por la longitud de \(\vec v\)). O sea
que la ecuación \(a\,u_x + b\,u_y = c\) no dice más que esto:

\begin{quote}
La derivada de \(u\) en la dirección del campo vale \(c\).
\end{quote}

Y el caso \(c = 0\) pasa a ser transparente: \textbf{\(u\) es constante
a lo largo de las curvas integrales del campo}. Es exactamente la noción
de primera integral que vamos a definir en la
sección~\ref{sec-coef-variables}, vista desde el otro lado.

\subsubsection{Después, levantarlo al espacio}\label{sec-campo-espacio}

El campo del piso dice por dónde caminar. Falta la altura, y de eso se
encarga la tercera componente. El campo completo

\[
\vec V(x,y,u) = \big(a,\; b,\; c\big)
\]

es el de \ref{eq-campo}, y su \textbf{sombra sobre el piso es
\(\vec v\)}. La tercera componente \(c\) es el \textbf{ritmo de ascenso
del hilo}: cuánto sube el techo por cada unidad de camino recorrido.

\begin{longtable}[]{@{}
  >{\raggedright\arraybackslash}p{(\linewidth - 6\tabcolsep) * \real{0.2500}}
  >{\raggedright\arraybackslash}p{(\linewidth - 6\tabcolsep) * \real{0.2500}}
  >{\raggedright\arraybackslash}p{(\linewidth - 6\tabcolsep) * \real{0.2500}}
  >{\raggedright\arraybackslash}p{(\linewidth - 6\tabcolsep) * \real{0.2500}}@{}}
\caption{Qué hace la tercera componente del campo. Las tres filas tienen
el mismo piso: las rutas son las mismas, cambia solo cómo sube el hilo
encima de ellas.}\label{tbl-tercera-componente}\tabularnewline
\toprule\noalign{}
\begin{minipage}[b]{\linewidth}\raggedright
Si \(c\)\ldots{}
\end{minipage} & \begin{minipage}[b]{\linewidth}\raggedright
el hilo\ldots{}
\end{minipage} & \begin{minipage}[b]{\linewidth}\raggedright
la ecuación
\end{minipage} & \begin{minipage}[b]{\linewidth}\raggedright
ejemplo
\end{minipage} \\
\midrule\noalign{}
\endfirsthead
\toprule\noalign{}
\begin{minipage}[b]{\linewidth}\raggedright
Si \(c\)\ldots{}
\end{minipage} & \begin{minipage}[b]{\linewidth}\raggedright
el hilo\ldots{}
\end{minipage} & \begin{minipage}[b]{\linewidth}\raggedright
la ecuación
\end{minipage} & \begin{minipage}[b]{\linewidth}\raggedright
ejemplo
\end{minipage} \\
\midrule\noalign{}
\endhead
\bottomrule\noalign{}
\endlastfoot
es \(0\) & va a altura constante & \(a\,u_x + b\,u_y = 0\) & transporte,
sección~\ref{sec-transporte} \\
depende solo de \((x,y)\) & sube a un ritmo prefijado en cada punto &
\(a\,u_x + b\,u_y = c(x,y)\) & Ejemplo~\ref{exm-inhomogenea} \\
depende de \(u\) & sube a un ritmo que depende de su propia altura &
\(a\,u_x + b\,u_y = c(x,y,u)\) & decaimiento, reacción \\
\end{longtable}

Y acá aparece la diferencia de fondo entre los casos que vamos a
estudiar:

\begin{tcolorbox}[enhanced jigsaw, arc=.35mm, bottomrule=.15mm, breakable, colback=white, colframe=quarto-callout-warning-color-frame, left=2mm, leftrule=.75mm, opacityback=0, rightrule=.15mm, toprule=.15mm]
\begin{minipage}[t]{5.5mm}
\textcolor{quarto-callout-warning-color}{\faExclamationTriangle}
\end{minipage}%
\begin{minipage}[t]{\textwidth - 5.5mm}

\vspace{-3mm}\textbf{En las cuasi-lineales no hay campo en el piso}\vspace{3mm}

Si \(a\) o \(b\) dependen de \(u\), la flecha del punto \((x,y)\)
\textbf{depende de a qué altura estés}. Dos observadores parados en el
mismo punto del piso, midiendo alturas distintas, caminan en direcciones
distintas.

Es decir: el campo existe, pero vive en el espacio \((x,y,u)\) y no se
puede bajar al piso. No hay un retrato de fases que dibujar antes de
resolver. Esa ---y no otra--- es la razón por la que el caso
cuasi-lineal de la sección~\ref{sec-casi-lineales} es más difícil, y la
razón por la que ahí aparecen los choques: dos hilos que en el espacio
nunca se tocan pueden tener \textbf{sombras que se cruzan}.

\end{minipage}%
\end{tcolorbox}

\subsubsection{Tres campos que ya conocés}\label{sec-campos-fisicos}

La interpretación no es una metáfora que inventamos para explicar el
método: los campos vectoriales con los que ya trabajaste en física son
exactamente esto, y sus líneas dibujadas son curvas integrales.

\begin{example}[El colorante en el balde que
gira]\protect\hypertarget{exm-corriente}{}\label{exm-corriente}

Un fluido gira como un cuerpo rígido con velocidad angular \(\omega\).
Su campo de velocidades es

\[
\vec v(x,y) = \big(-\omega y,\; \omega x\big) ,
\]

y las \textbf{líneas de corriente} ---que en un flujo estacionario
coinciden con las trayectorias de las partículas--- son sus curvas
integrales.

Echamos una gota de colorante que no difunde ni pesa: simplemente se
deja llevar. Su concentración \(u\) cumple que la derivada a lo largo de
la trayectoria de cada partícula es cero, o sea
\(u_t + \vec v \cdot \nabla u = 0\). Si además el patrón es
\textbf{estacionario} ---el dibujo del colorante no cambia con el
tiempo, \(u_t = 0\)---, queda

\[
-\omega y\,u_x + \omega x\,u_y = 0 .
\]

El sistema característico es \(\dot x = -\omega y\),
\(\dot y = \omega x\), \(\dot u = 0\): las características proyectadas
son las circunferencias \(x^2 + y^2 = \text{const}\), recorridas con
velocidad angular \(\omega\), y sobre cada una \(u\) no cambia. Por lo
tanto

\[
u(x,y) = F\big(x^2 + y^2\big)
\]

para una \(F\) arbitraria. La solución es cualquier patrón con simetría
de revolución.

\textbf{Lo que dice la física es lo mismo.} Si el patrón tuviera un
brazo, el brazo giraría con el fluido y el patrón no sería estacionario.
Para quedarse quieto mientras todo rota, el dibujo tiene que ser
invariante frente a rotaciones: anillos. La cuenta y la intuición dan la
misma respuesta, y la cuenta sale en tres renglones porque el campo ya
estaba dibujado.

\end{example}

\begin{example}[Las líneas de campo de una carga
puntual]\protect\hypertarget{exm-campo-electrico}{}\label{exm-campo-electrico}

El campo eléctrico de una carga \(q\) en el origen es radial:
\(\vec E = \dfrac{q}{4\pi\varepsilon_0}\dfrac{\vec r}{r^{3}}\). Como
para dibujar las líneas de campo solo importa la dirección, podemos usar
el campo \(\vec v(x,y) = (x, y)\), que apunta para el mismo lado en cada
punto.

Una magnitud \(u\) que sea constante a lo largo de las líneas de campo
cumple

\[
x\,u_x + y\,u_y = 0 .
\]

Sistema característico: \(\dot x = x\), \(\dot y = y\), de donde
\(x = x_0e^{s}\), \(y = y_0e^{s}\). Las curvas integrales son las
\textbf{semirrectas que salen del origen}, y como \(y/x\) no cambia
sobre cada una,

\[
u(x,y) = F\!\left(\frac{y}{x}\right) :
\]

cualquier función del ángulo.

\textbf{Dos hechos que ya sabías sobre las líneas de campo, traducidos.}

\begin{itemize}
\tightlist
\item
  \emph{Las líneas de campo nunca se cruzan.} Eso es, palabra por
  palabra, la unicidad de las curvas integrales: si se cruzaran, en el
  punto de cruce el campo tendría que apuntar en dos direcciones a la
  vez. Y es lo que nos va a garantizar en la sección~\ref{sec-tejido}
  que los hilos tejen sin superponerse.
\item
  \emph{Salvo en los puntos donde el campo se anula o explota.} En el
  origen, \(\vec v = (0,0)\): ahí se juntan todas las semirrectas, cada
  una llegando con su propio valor de \(u\). La solución solo puede ser
  continua en el origen si es constante. El punto donde falla la
  hipótesis es el punto donde falla la conclusión, y eso no es
  casualidad.
\end{itemize}

\end{example}

\begin{example}[El tráfico, que es el caso
difícil]\protect\hypertarget{exm-trafico-campo}{}\label{exm-trafico-campo}

En una ruta, la velocidad a la que se propaga una perturbación de la
densidad de autos depende de \textbf{cuántos autos hay}: con la ruta
vacía la perturbación viaja para adelante, y con la ruta saturada viaja
para atrás. En el lenguaje de esta sección: la flecha en el punto
\((x,t)\) depende de la altura \(u\), y no hay campo que dibujar en el
piso.

Es el ejemplo típico de la advertencia de más arriba, y es el que vamos
a desarrollar en el Ejemplo~\ref{exm-trafico}. Guardate la imagen: los
embotellamientos que aparecen sin motivo aparente son sombras de hilos
que se cruzan.

\end{example}

\subsubsection{Para jugar: el campo y sus curvas
integrales}\label{sec-lienzo-campo}

\emph{(En la versión web hay acá un gráfico interactivo: cuatro campos
vectoriales dibujados con flechas ---transporte, rotación, radial y
cizalla--- con sus curvas integrales, una curva inicial y un observador
que se puede mover a lo largo de la curva integral que sale del punto
elegido. Para cada campo se muestran la ecuación asociada y su solución
general.)}

Y ahora el mismo campo, levantado. Abajo está el piso con sus flechas y
las sombras; arriba, los hilos que suben a ritmo \(c\) y tejen la
superficie.

\emph{(En la versión web hay acá un segundo gráfico interactivo, en tres
dimensiones y girable: el campo de rotación dibujado con flechas sobre
el piso, sus circunferencias integrales como sombras, y los hilos
levantados encima con un ritmo de ascenso \(c\) regulable. Con \(c=0\)
los hilos son anillos horizontales; con \(c \neq 0\) la superficie es
una rampa en espiral, y el piso no cambia.)}

Para la ecuación \(-y\,u_x + x\,u_y = 3\), ¿cómo son las
\textbf{características proyectadas}?

\begin{itemize}
\tightlist
\item
  {Circunferencias centradas en el origen, igual que si el miembro
  derecho fuera \(0\)}
\item
  {Espirales que se abren, porque el miembro derecho no es cero}
\item
  {Hélices que suben a ritmo \(3\)}
\item
  {No se pueden dibujar sin conocer la solución}
\end{itemize}

Las características proyectadas las dictan las \textbf{dos primeras}
ecuaciones del sistema característico, \(\dot x = -y\), \(\dot y = x\),
y ahí el \(3\) no aparece. Son circunferencias, las mismas que en el
caso homogéneo.

Lo que cambia con el \(3\) es el \textbf{hilo}: ahora \(\dot u = 3\),
así que en vez de quedarse a altura constante sube a ritmo \(3\). La
curva característica en el espacio sí es una hélice ---la tercera opción
confunde el hilo con su sombra---, pero su sombra sigue siendo una
circunferencia.

La cuarta opción sería correcta si la ecuación fuera cuasi-lineal, es
decir si \(a\) o \(b\) dependieran de \(u\). Acá no: \(a=-y\) y \(b=x\)
dependen solo del punto del piso.

Un campo vectorial en el plano vale \(\vec v(x,y) = (y, -x)\). ¿Cuál de
estas funciones es solución de \(y\,u_x - x\,u_y = 0\)?

\begin{itemize}
\tightlist
\item
  {\(u(x,y) = e^{-(x^2+y^2)}\)}
\item
  {\(u(x,y) = x + y\)}
\item
  {\(u(x,y) = xy\)}
\item
  {\(u(x,y) = \operatorname{arctg}(y/x)\)}
\end{itemize}

Las curvas integrales de \((y,-x)\) son circunferencias centradas en el
origen ---recorridas en sentido horario, al revés que las de \((-y,x)\),
pero las mismas curvas---. Entonces las soluciones son exactamente las
funciones constantes sobre circunferencias, o sea las de la forma
\(u = F(x^2+y^2)\).

La primera lo es, con \(F(w) = e^{-w}\). Comprobación directa:
\(u_x = -2x\,u\), \(u_y = -2y\,u\), y \(y(-2xu) - x(-2yu) = 0\).

Las otras dos primeras no son constantes sobre circunferencias: \(x+y\)
y \(xy\) cambian si girás el punto. Y la última es \textbf{la peor
trampa}: \(\operatorname{arctg}(y/x)\) es el ángulo, que es constante
sobre las \textbf{semirrectas}, no sobre las circunferencias. Es
solución del campo radial \((x,y)\) del
Ejemplo~\ref{exm-campo-electrico}, no de este. Los dos campos son
perpendiculares entre sí en cada punto, y por eso sus soluciones se
intercambian.

\subsection{El tejido}\label{sec-tejido}

Un hilo solo no alcanza para armar un techo. Hacen falta muchos, y hace
falta que se apoyen uno al lado del otro sin dejar huecos ni
superponerse.

Eso es exactamente lo que pasa. Si \(a\), \(b\) y \(c\) son de clase
\(C^1\), el teorema de existencia y unicidad para sistemas de EDOs dice
que \textbf{por cada punto del espacio pasa una y solo una curva
característica}. Los hilos no se cortan entre sí y no se bifurcan.

\begin{theorem}[La superficie solución está tejida con
hilos]\protect\hypertarget{thm-tejido}{}\label{thm-tejido}

Una superficie es superficie solución de \ref{eq-cuasilineal-1} si y
solo si está formada por curvas características, es decir, si por cada
uno de sus puntos pasa una curva característica que queda enteramente
contenida en ella.

\end{theorem}

La demostración es la cuenta de la sección~\ref{sec-observador} leída en
los dos sentidos. Si \(S\) es solución y arrancamos una característica
en un punto de \(S\), la derivada de \(u\) a lo largo de esa curva
coincide con \(c\), que es lo mismo que le pide la tercera ecuación del
sistema; por unicidad, la curva no se despega de \(S\). Recíprocamente,
si \(S\) está hecha de características, en cada punto el vector
\(\vec V\) es tangente a \(S\), y eso ---lo vimos--- es la ecuación.

Y de ahí sale la receta completa.

\begin{tcolorbox}[enhanced jigsaw, arc=.35mm, bottomrule=.15mm, breakable, colback=white, colframe=quarto-callout-tip-color-frame, left=2mm, leftrule=.75mm, opacityback=0, rightrule=.15mm, toprule=.15mm]
\begin{minipage}[t]{5.5mm}
\textcolor{quarto-callout-tip-color}{\faLightbulb}
\end{minipage}%
\begin{minipage}[t]{\textwidth - 5.5mm}

\vspace{-3mm}\textbf{Cómo se teje una solución}\vspace{3mm}

\textbf{Datos:} una curva inicial \(\Gamma\) en el espacio, dada
paramétricamente por \[
x = x_0(\sigma), \qquad y = y_0(\sigma), \qquad u = u_0(\sigma) ,
\] que es la forma precisa de decir «conocemos \(u\) sobre una curva del
piso». El parámetro \(\sigma\) recorre la curva.

\textbf{Receta:}

\begin{enumerate}
\def\labelenumi{\arabic{enumi}.}
\tightlist
\item
  Por \textbf{cada} punto de \(\Gamma\), largar el hilo que le
  corresponde: resolver el sistema característico con dato inicial
  \(x(0)=x_0(\sigma)\), \(y(0)=y_0(\sigma)\), \(u(0)=u_0(\sigma)\).
\item
  Barrer \(\sigma\) a lo largo de toda \(\Gamma\). Los hilos, uno al
  lado del otro, tejen una superficie \[
  x = x(\sigma,s), \qquad y = y(\sigma,s), \qquad u = u(\sigma,s) .
  \]
\item
  Despejar \(\sigma\) y \(s\) en función de \(x\) e \(y\) y reemplazar
  en la tercera, para obtener \(u = u(x,y)\).
\end{enumerate}

Cada punto del techo queda identificado por dos coordenadas con
significado: \textbf{\(\sigma\) dice en qué hilo estamos} y
\textbf{\(s\) dice cuánto avanzamos por ese hilo}. Son, literalmente, la
trama y la urdimbre del tejido.

\end{minipage}%
\end{tcolorbox}

\begin{figure}

\centering{

\pandocbounded{\includegraphics[keepaspectratio]{08-edp_files/figure-pdf/fig-tejido-output-1.pdf}}

}

\caption{\label{fig-tejido}El tejido, para la ecuación de transporte
\(u_t + c\,u_x = 0\) con \(c=1\). A la izquierda, el piso: las
características proyectadas son las trayectorias posibles del
observador, y la curva inicial es el eje \(t=0\). A la derecha, el
techo: sobre cada trayectoria del piso se levanta un hilo, y los hilos,
uno al lado del otro, son la superficie solución. Los dos hilos
resaltados salen de los dos puntos resaltados del piso. Como acá
\(du/ds=0\), cada hilo es horizontal: el observador siempre mide lo
mismo.}

\end{figure}%

\begin{tcolorbox}[enhanced jigsaw, arc=.35mm, bottomrule=.15mm, breakable, colback=white, colframe=quarto-callout-important-color-frame, left=2mm, leftrule=.75mm, opacityback=0, rightrule=.15mm, toprule=.15mm]
\begin{minipage}[t]{5.5mm}
\textcolor{quarto-callout-important-color}{\faExclamation}
\end{minipage}%
\begin{minipage}[t]{\textwidth - 5.5mm}

\vspace{-3mm}\textbf{Para tejer, los hilos tienen que cruzar la curva inicial}\vspace{3mm}

La receta tiene un supuesto escondido, y es el que va a explicar todos
los casos patológicos de la sección~\ref{sec-casi-lineales}.

Para que los hilos barran una superficie hace falta que \textbf{se
alejen} de la curva inicial. Si en algún punto la característica
proyectada resulta ser tangente a la curva inicial ---si el observador,
al salir, ya estaba caminando en la dirección de \(\Gamma\)--- entonces
el hilo que sale de ahí se queda pegado a \(\Gamma\) en vez de
despegarse, y no barre nada. En un telar, la trama tiene que cruzar a la
urdimbre; si van paralelas no hay tela.

Cuando eso pasa, pueden ocurrir dos cosas, y las dos son malas:

\begin{itemize}
\tightlist
\item
  si el dato \(u_0(\sigma)\) que pusimos sobre \(\Gamma\) no es el que
  la ecuación misma habría producido a lo largo de \(\Gamma\), hay una
  contradicción y \textbf{no existe solución};
\item
  si casualmente coincide, la ecuación no dice nada nuevo, y \textbf{hay
  infinitas soluciones}: cualquier tejido que contenga a ese hilo sirve.
\end{itemize}

Esa es la razón profunda de las curvas «prohibidas» como dato. Es una de
las cuatro maneras de romper el tejido, y las vemos todas juntas en la
sección~\ref{sec-patologias}; el teorema y la cuenta concreta están en
la sección~\ref{sec-casi-lineales}.

\end{minipage}%
\end{tcolorbox}

\emph{(En la versión web, la figura anterior se puede girar: hay un
gráfico interactivo con la misma superficie, cuatro sliders para mover
la cámara y la velocidad, y las sombras de los hilos dibujadas sobre el
piso.)}

\subsection{Cuando el tejido se rompe}\label{sec-patologias}

La receta de la sección~\ref{sec-tejido} parece infalible: se largan
hilos y sale una superficie. No lo es, y conviene ver dónde falla,
porque \textbf{cada falla está atada a una hipótesis} de los teoremas
que vamos a enunciar. Las hipótesis no son decoración: cada una está
tapando un agujero concreto.

Son cuatro, y las cuatro se entienden mirando el campo.

\subsubsection{1. Los hilos se
bifurcan}\label{sec-patologia-bifurcacion}

La sección~\ref{sec-tejido} usó, casi al pasar, que \textbf{por cada
punto pasa una sola curva característica}. Eso no es gratis: es el
teorema de existencia y unicidad para EDOs, y pide que el campo sea al
menos localmente Lipschitz ---por ejemplo, de clase \(C^1\)---.

Si no lo es, se rompe. Tomemos el campo

\[
\vec v(x,y) = \big(1,\; 3y^{2/3}\big) ,
\]

que es \textbf{continuo en todo el plano} pero no es derivable sobre el
eje \(x\): su segunda componente tiene derivada \(2y^{-1/3}\), que
explota en \(y=0\). Sus curvas integrales salen de resolver
\(y' = 3y^{2/3}\), y por el origen pasan:

\begin{itemize}
\tightlist
\item
  la recta \(y \equiv 0\);
\item
  la cúbica \(y = x^3\);
\item
  y, para cada \(\lambda > 0\), la curva que \textbf{se queda quieta}
  sobre el eje hasta \(x=\lambda\) y recién ahí despega como
  \(y = (x-\lambda)^3\).
\end{itemize}

Infinitas curvas integrales por un solo punto. El observador que sale
del origen no tiene una ruta: tiene un menú. Y entonces la ecuación
\(u_x + 3y^{2/3}u_y = 0\) deja de determinar \(u\), porque no se sabe
qué hilo levantar.

\begin{tcolorbox}[enhanced jigsaw, arc=.35mm, bottomrule=.15mm, breakable, colback=white, colframe=quarto-callout-note-color-frame, left=2mm, leftrule=.75mm, opacityback=0, rightrule=.15mm, toprule=.15mm]
\begin{minipage}[t]{5.5mm}
\textcolor{quarto-callout-note-color}{\faInfo}
\end{minipage}%
\begin{minipage}[t]{\textwidth - 5.5mm}

\vspace{-3mm}\textbf{Esto le pasa a la física, no solo a los ejemplos armados}\vspace{3mm}

Hay un sistema mecánico clásico con exactamente esta patología, la
\textbf{cúpula de Norton}: una cúpula con un perfil elegido de modo que
la ecuación de movimiento de una masa apoyada en su vértice sea
\(\ddot r = \sqrt{r}\).

El miembro derecho es continuo pero no Lipschitz en \(r=0\). La masa
quieta en el vértice es una solución; pero también lo es, para
\textbf{cada} \(T>0\), \[
r(t) = \frac{(t-T)^4}{144} \quad\text{para } t > T, \qquad r(t) = 0 \quad\text{para } t \le T ,
\] ---verificalo: \(\ddot r = (t-T)^2/12\) y
\(\sqrt r = (t-T)^2/12\)---. O sea: la masa puede quedarse quieta un
rato arbitrario y después empezar a deslizarse sola, sin que nada la
empuje, sin violar la ecuación de Newton.

Es un ejemplo famoso en filosofía de la física, y lo que muestra es que
el determinismo de la mecánica clásica no es un teorema: es una
hipótesis de regularidad sobre las fuerzas (Norton, J. D.,
\emph{Causation as Folk Science}, 2003).

\end{minipage}%
\end{tcolorbox}

\subsubsection{2. El dato está puesto a lo largo de un
hilo}\label{sec-patologia-transversalidad}

Es el que ya anunciamos al final de la sección~\ref{sec-tejido}. Para
que los hilos barran una superficie tienen que \textbf{cruzar} la curva
inicial; si salen paralelos a ella, no barren nada.

El caso más chico posible: \(u_x + u_y = 0\). El campo es el constante
\((1,1)\) y las características proyectadas son las rectas
\(y = x + k\). Si damos el dato sobre la recta \(y=x\) ---que es una de
ellas---, pasa lo siguiente:

\begin{itemize}
\tightlist
\item
  sobre esa recta la ecuación obliga a \(u\) a ser \textbf{constante},
  porque \(\dot u = 0\);
\item
  si el dato que pusimos no es constante, por ejemplo
  \(u(\sigma,\sigma) = \sigma^2\), hay una contradicción: \textbf{no
  existe solución};
\item
  si es constante, por ejemplo \(u(\sigma,\sigma)=5\), la ecuación no
  aporta nada nuevo fuera de esa recta: \textbf{hay infinitas
  soluciones}, una por cada \(F\) con \(F(0)=5\) en \(u = F(y-x)\).
\end{itemize}

La hipótesis que tapa este agujero es la \textbf{condición de
transversalidad} del Teorema~\ref{thm-transversalidad}, y el detalle de
los dos casos está en el Teorema~\ref{thm-degenerado}.

\subsubsection{3. Dos sombras que se cruzan}\label{sec-patologia-choque}

Este es el más interesante, y el que más se malentiende. Hay que tener
cuidado con qué se cruza con qué:

\begin{tcolorbox}[enhanced jigsaw, arc=.35mm, bottomrule=.15mm, breakable, colback=white, colframe=quarto-callout-important-color-frame, left=2mm, leftrule=.75mm, opacityback=0, rightrule=.15mm, toprule=.15mm]
\begin{minipage}[t]{5.5mm}
\textcolor{quarto-callout-important-color}{\faExclamation}
\end{minipage}%
\begin{minipage}[t]{\textwidth - 5.5mm}

\vspace{-3mm}\textbf{Los hilos no se cruzan; sus sombras sí pueden}\vspace{3mm}

En el espacio \((x,y,u)\) la unicidad sigue valiendo: por cada punto
pasa una sola curva característica y los hilos nunca se tocan. Lo que
puede pasar es que \textbf{dos hilos que van a alturas distintas
proyecten sombras que se cruzan en el piso}.

Cuando eso ocurre, al punto del piso donde se cruzan las sombras llegan
dos valores distintos de \(u\). El tejido sigue siendo una superficie
perfectamente lisa en el espacio, pero deja de ser el \textbf{gráfico de
una función}: se dobla sobre sí misma, como una ola que rompe.

En el caso lineal y semilineal esto es imposible, porque ahí las sombras
son curvas integrales de un campo del piso y por lo tanto tampoco se
cruzan. En el cuasi-lineal es el comportamiento típico.

\end{minipage}%
\end{tcolorbox}

La hipótesis que se viola no es ninguna de las del teorema: es la
\textbf{conclusión} la que está limitada. El
Teorema~\ref{thm-transversalidad} da una solución de clase \(C^1\)
\textbf{en un entorno}, y nada más. El Teorema~\ref{thm-quiebre} dice
exactamente cuándo se termina ese entorno, y la
sección~\ref{sec-generalizadas} dice qué hacer después: aflojar la
definición de solución y admitir discontinuidades.

\subsubsection{4. Hay piso que ningún hilo
alcanza}\label{sec-patologia-hueco}

El último es el más simple y el más fácil de olvidar en un parcial. Los
hilos solo llevan información \textbf{desde} la curva inicial. En la
parte del piso a la que no llega ningún hilo, el dato no dice
absolutamente nada, y la solución no queda determinada ahí.

Para el transporte \(u_t + c\,u_x = 0\) con \(c>0\) y un dato dado solo
sobre el segmento \(0 \le x \le 1\) del eje \(t=0\), la solución queda
determinada únicamente en la franja \(0 \le x - ct \le 1\). Afuera puede
valer cualquier cosa.

Esto no es una patología: es \textbf{causalidad}. Uno solo puede saber
lo que le trajo algo que llegó hasta él. Es la misma idea que el cono de
luz, y es la razón por la que en el capítulo 9 va a aparecer el
\emph{dominio de dependencia} de la ecuación de onda.

\begin{tcolorbox}[enhanced jigsaw, arc=.35mm, bottomrule=.15mm, breakable, colback=white, colframe=quarto-callout-tip-color-frame, left=2mm, leftrule=.75mm, opacityback=0, rightrule=.15mm, toprule=.15mm]
\begin{minipage}[t]{5.5mm}
\textcolor{quarto-callout-tip-color}{\faLightbulb}
\end{minipage}%
\begin{minipage}[t]{\textwidth - 5.5mm}

\vspace{-3mm}\textbf{El caso opuesto: los hilos se abren y dejan un hueco}\vspace{3mm}

En las cuasi-lineales también pasa lo contrario del choque. Si el dato
inicial tiene un salto hacia arriba, los hilos de la derecha salen más
rápido que los de la izquierda, se separan, y queda una cuña del piso
sin ningún hilo encima.

Eso se llama \textbf{onda de rarefacción}, y es lo que pasa cuando
arranca el auto de adelante en un semáforo: la fila no se estira toda
junta, se va soltando de a poco. La solución se completa ahí abanicando
hilos desde el punto del salto.

\end{minipage}%
\end{tcolorbox}

\subsubsection{El resumen de las cuatro}\label{sec-tabla-patologias}

\begin{longtable}[]{@{}
  >{\raggedright\arraybackslash}p{(\linewidth - 4\tabcolsep) * \real{0.3333}}
  >{\raggedright\arraybackslash}p{(\linewidth - 4\tabcolsep) * \real{0.3333}}
  >{\raggedright\arraybackslash}p{(\linewidth - 4\tabcolsep) * \real{0.3333}}@{}}
\caption{Las hipótesis del método, y el agujero que tapa cada una. Las
tres primeras son hipótesis de los teoremas; las dos últimas son límites
de lo que los teoremas pueden
prometer.}\label{tbl-patologias}\tabularnewline
\toprule\noalign{}
\begin{minipage}[b]{\linewidth}\raggedright
Hipótesis
\end{minipage} & \begin{minipage}[b]{\linewidth}\raggedright
Qué garantiza
\end{minipage} & \begin{minipage}[b]{\linewidth}\raggedright
Qué pasa si falla
\end{minipage} \\
\midrule\noalign{}
\endfirsthead
\toprule\noalign{}
\begin{minipage}[b]{\linewidth}\raggedright
Hipótesis
\end{minipage} & \begin{minipage}[b]{\linewidth}\raggedright
Qué garantiza
\end{minipage} & \begin{minipage}[b]{\linewidth}\raggedright
Qué pasa si falla
\end{minipage} \\
\midrule\noalign{}
\endhead
\bottomrule\noalign{}
\endlastfoot
\(a\), \(b\), \(c\) de clase \(C^1\) & por cada punto pasa un solo hilo
& los hilos se bifurcan: \(u\) no queda determinada \\
\((a,b) \neq (0,0)\) & el observador efectivamente avanza & el punto es
singular del campo; las sombras se juntan ahí \\
transversalidad: \(a\,y_0' - b\,x_0' \neq 0\) & los hilos cruzan
\(\Gamma\) y barren una superficie & o no hay solución, o hay
infinitas \\
la conclusión es \textbf{local} & la superficie es el gráfico de una
función & las sombras se cruzan: choque, y hay que pasar a soluciones
débiles \\
\(\Gamma\) toca todas las características de la región & el dato alcanza
para todo el dominio & quedan zonas del piso sin tejer \\
\end{longtable}

\subsubsection{Para jugar: las cuatro
fallas}\label{sec-lienzo-patologias}

\emph{(En la versión web hay acá un gráfico interactivo con las cuatro
fallas, una por vez: las infinitas curvas integrales que salen del
origen para un campo no Lipschitz; una curva inicial que se puede
inclinar hasta volverse característica, con el determinante de
transversalidad calculado en vivo; las características de Burgers
juntándose hasta tocarse en \(t^{*}=1\); y la franja del piso que queda
determinada cuando el dato se da solo sobre un segmento.)}

Para \(u_x + u_y = 0\) damos el dato \(u = \sigma^{2}\) sobre la recta
\(y = x\), parametrizada como \(x_0(\sigma)=\sigma\),
\(y_0(\sigma)=\sigma\). ¿Qué pasa?

\begin{itemize}
\tightlist
\item
  {No existe solución}
\item
  {Existe una única solución, \(u = \tfrac14(x+y)^2\)}
\item
  {Existen infinitas soluciones}
\item
  {Existe solución pero no es de clase \(C^1\)}
\end{itemize}

El determinante de transversalidad es \[
a\,y_0'(\sigma) - b\,x_0'(\sigma) = 1\cdot 1 - 1\cdot 1 = 0 ,
\] así que la recta \(y=x\) \textbf{es} una característica y el teorema
no se aplica. Hay que mirar el dato a mano.

Sobre una característica la ecuación obliga a \(\dot u = 0\), o sea que
\(u\) tiene que ser constante a lo largo de \(y=x\). El dato que
pusimos, \(\sigma^2\), no lo es. Contradicción: \textbf{no existe
solución}, ni clásica ni de ninguna otra clase.

La tercera opción sería la correcta si el dato hubiera sido constante
---por ejemplo \(u=5\)---: ahí la ecuación no se contradice pero tampoco
determina nada fuera de la recta, y sirve cualquier \(u = F(y-x)\) con
\(F(0)=5\). Las dos caras de la misma falla, y cuál de las dos te toca
depende del dato, no de la curva.

La ecuación \(u_t + u\,u_x = 0\) arranca con el dato \(u(x,0) = 1 - 4x\)
en el intervalo donde ese dato es decreciente. ¿En qué instante
\(t^{*}\) se tocan por primera vez dos características?

La característica que sale de \(x_0\) es la recta
\(x = x_0 + f(x_0)\,t\) con \(f(x_0) = 1-4x_0\). Dos características
vecinas se tocan cuando la separación entre ellas se anula. Derivando
respecto de \(x_0\): \[
\frac{\partial x}{\partial x_0} = 1 + f'(x_0)\,t = 1 - 4t ,
\] que se anula en \(t = 1/4\). Como \(f' = -4\) es constante,
\textbf{todas} las características se juntan en el mismo instante y en
el mismo punto: el dato lineal es el caso más simétrico posible.

En general \(t^{*} = -1/\min f'\), que es el Teorema~\ref{thm-quiebre}.
Acá \(\min f' = -4\) y \(t^{*} = 1/4 = 0{,}25\).

\begin{exercise}[Diagnosticar la
falla]\protect\hypertarget{exr-patologias-diagnostico}{}\label{exr-patologias-diagnostico}

Para cada uno de estos problemas, decidí cuál de las cuatro fallas del
Tabla~\ref{tbl-patologias} aparece ---si es que aparece alguna--- y qué
le pasa a la solución.

\begin{enumerate}
\def\labelenumi{\alph{enumi})}
\tightlist
\item
  \(u_x + 2u_y = 0\) con \(u = \operatorname{sen}\sigma\) sobre la curva
  \(x_0(\sigma)=\sigma\), \(y_0(\sigma)=2\sigma\).
\item
  \(x\,u_x + y\,u_y = 0\) con \(u = \sigma\) sobre la circunferencia
  \(x_0 = \cos\sigma\), \(y_0 = \operatorname{sen}\sigma\), cerca del
  origen.
\item
  \(u_t + c\,u_x = 0\) con \(c>0\) constante y dato conocido solo para
  \(x \ge 0\) en \(t=0\). ¿Dónde queda determinada la solución?
\item
  \(u_t + u\,u_x = 0\) con \(u(x,0) = -x\).
\end{enumerate}

\begin{tcolorbox}[enhanced jigsaw, arc=.35mm, bottomrule=.15mm, bottomtitle=1mm, breakable, colback=white, colbacktitle=quarto-callout-tip-color!10!white, colframe=quarto-callout-tip-color-frame, coltitle=black, left=2mm, leftrule=.75mm, opacityback=0, opacitybacktitle=0.6, rightrule=.15mm, title=\textcolor{quarto-callout-tip-color}{\faLightbulb}\hspace{0.5em}{Solución}, titlerule=0mm, toprule=.15mm, toptitle=1mm]

\textbf{a) Falla la transversalidad.} El campo es \((1,2)\) y la curva
inicial tiene vector tangente \((x_0', y_0') = (1,2)\): es la misma
dirección. El determinante da \(1\cdot 2 - 2\cdot 1 = 0\). La curva
inicial es una característica.

Sobre ella la ecuación pide \(u\) constante, y el dato
\(\operatorname{sen}\sigma\) no lo es: \textbf{no hay solución}.

\textbf{b) Falla la hipótesis \((a,b)\neq(0,0)\), pero no sobre la curva
inicial.} Sobre la circunferencia el campo es
\((\cos\sigma, \operatorname{sen}\sigma)\), radial, y la tangente es
\((-\operatorname{sen}\sigma, \cos\sigma)\): el determinante vale
\(\cos^2\sigma + \operatorname{sen}^2\sigma = 1 \neq 0\). Transversal,
todo bien.

El problema está en el \textbf{origen}, donde el campo se anula. Las
características son las semirrectas que salen de ahí, cada una llevando
su propio valor \(\sigma\), y todas llegan al mismo punto. La solución
\(u = \sigma(x,y)\) ---el ángulo--- existe y es suave en todo el plano
menos el origen, y ahí \textbf{no se puede extender con continuidad}: es
el Ejemplo~\ref{exm-campo-electrico}.

Como yapa: tampoco se puede extender con continuidad a la semirrecta
donde el ángulo pega el salto de \(2\pi\) a \(0\). Eso ya no es una
falla del método, es que el dato que elegimos sobre la circunferencia no
cierra: \(u(2\pi) \neq u(0)\) aunque sea el mismo punto.

\textbf{c) Queda piso sin tejer.} Las características son
\(x - ct = \text{const}\) y salen de \(x_0 \ge 0\), así que barren la
región \(x - ct \ge 0\). Ahí la solución vale \(u = f(x-ct)\).

En \(x - ct < 0\) ---a la izquierda de la recta \(x=ct\)--- no llega
ningún hilo y la solución no queda determinada. Con \(c>0\) eso es la
región que «viene de afuera» por la izquierda: para saber qué pasa ahí
haría falta un dato de contorno en \(x=0\) para todo \(t>0\). Es
exactamente la situación de un caño o una ruta semi-infinita, donde hay
que decir qué entra por la punta.

\textbf{d) No falla nada, y esa es la gracia.} Acá \(f(x) = -x\) es
decreciente, así que uno esperaría un choque. Pero la característica que
sale de \(x_0\) es \[
x = x_0 + f(x_0)\,t = x_0(1-t) ,
\] y todas se juntan \textbf{a la vez} en el origen, en \(t = 1\). Antes
de eso, \(x_0 = x/(1-t)\) y \[
u(x,t) = -\frac{x}{1-t} ,
\] que es solución clásica en \(0 \le t < 1\) y explota en \(t=1\).

O sea: el Teorema~\ref{thm-quiebre} da
\(t^{*} = -1/\min f' = -1/(-1) = 1\), y se cumple. Lo llamativo es que
acá la solución no desarrolla un salto sino que \textbf{la pendiente se
va a infinito en todos lados al mismo tiempo}. Es el mismo fenómeno del
ejercicio numérico de arriba, y muestra que «se cruzan las
características» y «aparece un choque» no son sinónimos: lo primero es
lo que mata a la solución clásica, y la forma que toma después depende
del dato.

\end{tcolorbox}

\end{exercise}

\subsection{El caso más simple: transporte a velocidad
constante}\label{sec-transporte}

Con la maquinaria armada, el caso testigo sale en cuatro renglones. Sea
\(c\) una constante y

\begin{equation}\protect\phantomsection\label{eq-transporte}{
\frac{\partial u}{\partial t} + c\,\frac{\partial u}{\partial x} = 0 .
}\end{equation}

Físicamente: \(u(x,t)\) es una magnitud que se transporta sin deformarse
a velocidad \(c\) ---una mancha de tinta que baja por un canal, un pulso
que viaja por un cable, una perturbación de densidad en un gas---. Ya la
habíamos deducido en la 1 como la ley de conservación con flujo
\(q = cu\).

Acá las variables son \((x,t)\) en vez de \((x,y)\), y comparando con
\ref{eq-cuasilineal-1} el campo es \(\vec V = (c, 1, 0)\): coeficiente
\(c\) para \(u_x\), coeficiente \(1\) para \(u_t\), y lado derecho nulo.
El sistema característico queda

\[
\frac{dx}{ds} = c, \qquad \frac{dt}{ds} = 1, \qquad \frac{du}{ds} = 0 .
\]

La segunda dice que podemos usar directamente \(t\) como parámetro. Con
eso:

\[
x(t) = x_0 + ct, \qquad u = \text{constante a lo largo de cada hilo.}
\]

Las características proyectadas son rectas paralelas de pendiente
\(1/c\) en el plano \((x,t)\), y sobre cada una el observador mide
siempre lo mismo. Los hilos son horizontales, como se ve en la
Figura~\ref{fig-tejido}.

Vale la pena guardar la frase que resume el método y que sirve para
reconocerlo en cualquier libro: \textbf{una característica es una curva
a lo largo de la cual la ecuación en derivadas parciales se reduce a una
ecuación diferencial ordinaria.} Para \ref{eq-transporte} esa ecuación
ordinaria es \(du/dt = 0\), y por eso \(u\) es constante sobre cada
característica.

Si conocemos \(u\) en el instante inicial, \(u(x,0) = g(x)\), el valor
en cualquier punto se obtiene retrocediendo por la característica que
pasa por él hasta llegar al piso inicial:

\begin{equation}\protect\phantomsection\label{eq-sol-transporte}{
u(x,t) = g(x - ct).
}\end{equation}

\begin{example}[Un pulso que
viaja]\protect\hypertarget{exm-pulso}{}\label{exm-pulso}

Tomemos \(g(x) = e^{-x^2}\) y \(c = 1\). La solución es
\(u(x,t) = e^{-(x-t)^2}\): la misma campana, corrida hacia la derecha.
No se ensancha ni se achata, porque \ref{eq-transporte} no tiene ningún
término que difunda. Compárese con lo que le pasará al mismo pulso bajo
la ecuación del calor en la Capítulo~\ref{sec-parabolicas}.

\end{example}

\begin{figure}

\centering{

\pandocbounded{\includegraphics[keepaspectratio]{08-edp_files/figure-pdf/fig-caracteristicas-output-1.pdf}}

}

\caption{\label{fig-caracteristicas}Izquierda: las características de la
ecuación de transporte son rectas paralelas de pendiente \(1/c\), y el
dato inicial se propaga sin deformarse a lo largo de ellas. Derecha: el
pulso en tres instantes.}

\end{figure}%

Y ahora, para jugar con las tres cosas a la vez.

\emph{(En la versión web hay acá un gráfico interactivo: se elige una
característica con un slider, se avanza en el tiempo con otro, y se ve
que la altura que mide el observador no cambia nunca.)}

\begin{exercise}[]\protect\hypertarget{exr-transporte}{}\label{exr-transporte}

Resolver \(u_t + 3u_x = 0\) con \(u(x,0) = \dfrac{1}{1+x^2}\) y
describir qué pasa con el máximo de la solución a medida que avanza el
tiempo. ¿Qué diferencia habría si la ecuación fuera \(u_t + 3u_x = -u\)?
Interpretar el resultado en términos de los hilos: ¿siguen siendo
horizontales?

\end{exercise}

\begin{tcolorbox}[enhanced jigsaw, arc=.35mm, bottomrule=.15mm, bottomtitle=1mm, breakable, colback=white, colbacktitle=quarto-callout-tip-color!10!white, colframe=quarto-callout-tip-color-frame, coltitle=black, left=2mm, leftrule=.75mm, opacityback=0, opacitybacktitle=0.6, rightrule=.15mm, title=\textcolor{quarto-callout-tip-color}{\faLightbulb}\hspace{0.5em}{Solución}, titlerule=0mm, toprule=.15mm, toptitle=1mm]

\textbf{Sin el término extra.} Las características son
\(x - 3t = \text{const}\) y sobre ellas \(u\) no cambia, así que
\[u(x,t) = \frac{1}{1 + (x-3t)^2}.\] El máximo vale \(1\) siempre y se
mueve hacia la derecha a velocidad \(3\). \textbf{La forma no cambia en
absoluto}: la campana no se ensancha ni se achata. Es lo que hace la
ecuación de transporte y nada más.

\textbf{Con el \(-u\).} Las características son las mismas ---el lado
izquierdo no cambió--- pero ahora la tercera ecuación es \(du/ds = -u\)
en vez de \(0\), con \(s = t\). Integrando, \(u = u_0 e^{-t}\) sobre
cada característica, y por lo tanto
\[u(x,t) = \frac{e^{-t}}{1 + (x-3t)^2}.\] El máximo ahora decae como
\(e^{-t}\), pero \textbf{sigue viajando a velocidad \(3\) y la forma
sigue sin deformarse}: lo único que cambió es que todo el perfil se
apaga parejo.

\textbf{Los hilos.} Ya no son horizontales: bajan exponencialmente
mientras avanzan. Es el Ejemplo~\ref{exm-inhomogenea} con el signo
cambiado.

Físicamente, el término \(-u\) es un sumidero proporcional a lo que hay
---una sustancia radiactiva que decae mientras la arrastra la corriente,
o un contaminante que se degrada---. Y el resultado es exactamente lo
que uno esperaría: el decaimiento y el transporte no se estorban entre
sí, porque actúan en direcciones distintas del problema.

\end{tcolorbox}

\subsection{Coeficientes variables}\label{sec-coef-variables}

El método no cambia en nada si los coeficientes dependen del punto. Sea

\begin{equation}\protect\phantomsection\label{eq-lineal-homogenea-1}{
a(x,y)\,u_x + b(x,y)\,u_y = 0 .
}\end{equation}

El sistema característico es \(\dot x = a\), \(\dot y = b\),
\(\dot u = 0\): sobre cada característica, \(u\) sigue siendo constante.
Lo único que cambió es que las características proyectadas ya no son
rectas, sino las curvas integrales del campo plano \((a,b)\) ---o sea,
las soluciones de

\[
\frac{dy}{dx} = \frac{b(x,y)}{a(x,y)} ,
\]

que es una EDO de primer orden como las de la Capítulo~\ref{sec-edo}---.
Pueden curvarse, juntarse, y en algunos puntos hacer cosas raras.

\begin{definition}[Primera
integral]\protect\hypertarget{def-primera-integral}{}\label{def-primera-integral}

Una función \(\phi(x,y)\) de clase \(C^1\), no constante, es una
\textbf{primera integral} del campo \((a,b)\) si es constante a lo largo
de cada curva integral del campo.

Derivando \(\phi(x(s),y(s))\) respecto de \(s\) y usando el sistema
característico, la condición «\(\phi\) es constante sobre las curvas
integrales» se escribe

\[
a\,\phi_x + b\,\phi_y = 0 ,
\]

es decir: \textbf{\(\phi\) es primera integral si y solo si \(\phi\) es
solución de \ref{eq-lineal-homogenea-1}.}

\end{definition}

Ese es un resultado muy rentable, porque da la solución general de una
sola vez.

\begin{theorem}[Solución general de la
homogénea]\protect\hypertarget{thm-general-homogenea}{}\label{thm-general-homogenea}

Si \(\phi\) es una primera integral del campo \((a,b)\) en un dominio
\(D\), entonces

\[
u(x,y) = F\big(\phi(x,y)\big)
\]

es solución de \ref{eq-lineal-homogenea-1} en \(D\) para
\textbf{cualquier} función \(F\) de una variable de clase \(C^1\).
Recíprocamente, toda solución tiene esa forma.

\end{theorem}

La demostración de la ida es la regla de la cadena:
\(a\,(F'\phi_x) + b\,(F'\phi_y) = F'\,(a\phi_x + b\phi_y) = 0\). La
vuelta es el Teorema~\ref{thm-tejido}: una solución es constante sobre
cada característica, y \(\phi\) etiqueta a las características, así que
el valor de \(u\) solo puede depender de la etiqueta.

Y acá reaparece, con nombre y apellido, la función arbitraria que nos
había asustado en la sección~\ref{sec-porque}: \textbf{es la \(F\)}. La
indeterminación funcional de una EDP de primer orden es exactamente la
libertad de elegir qué valor asignarle a cada característica.

\begin{example}[Dos ejemplos con características
curvas]\protect\hypertarget{exm-coef-variables}{}\label{exm-coef-variables}

\textbf{Uno.} \(x\,u_x + y\,u_y = 0\) en el primer cuadrante. Las
características proyectadas cumplen \(dy/dx = y/x\), o sea \(y = kx\):
son los rayos que salen del origen. Una primera integral es
\(\phi = y/x\), y la solución general es \(u(x,y) = F(y/x)\). Tiene
sentido: \(u\) debe ser constante sobre cada rayo, así que solo puede
depender del ángulo.

El origen es un punto conflictivo, y ahora se ve por qué: todas las
características se cruzan ahí. Como sobre cada una \(u\) vale algo
distinto, no hay ningún valor que se le pueda asignar a \(u(0,0)\) de
manera continua, salvo que \(F\) sea constante.

\textbf{Dos.} \(u_t + x\,u_x = 0\) con \(u(x,0)=g(x)\). El sistema es
\(\dot t = 1\), \(\dot x = x\), \(\dot u = 0\), así que
\(x(t) = x_0 e^{t}\): las características se abren exponencialmente
desde el eje \(t\). Como \(u\) es constante sobre cada una y
\(x_0 = x e^{-t}\), \[
u(x,t) = g\big(x\,e^{-t}\big) .
\] El perfil no viaja: se estira. Cada punto se aleja del origen a
velocidad proporcional a su distancia.

\end{example}

Para el caso inhomogéneo, \(a\,u_x + b\,u_y = c(x,y,u)\), cambia solo la
tercera ecuación: \(\dot u = c\) en vez de \(\dot u = 0\). El observador
ya no mide siempre lo mismo, sino que ve crecer o decrecer el techo a un
ritmo conocido, y el hilo deja de ser horizontal.

\begin{example}[Un hilo que
sube]\protect\hypertarget{exm-inhomogenea}{}\label{exm-inhomogenea}

Resolvamos \(a\,u_x + b\,u_y = u\) con \(a\), \(b\) constantes y
\(u(0,y) = f(y)\).

El sistema característico es \(\dot x = a\), \(\dot y = b\),
\(\dot u = u\). Parametricemos la curva inicial ---que es el eje
\(y\)--- por \(\sigma\): arrancamos en \(x=0\), \(y=\sigma\),
\(u=f(\sigma)\). Integrando, \[
x = a s, \qquad y = b s + \sigma, \qquad u = e^{s} f(\sigma) .
\] Ahora despejamos: \(s = x/a\) y \(\sigma = y - bx/a\). Reemplazando,
\[
u(x,y) = e^{x/a}\,f\!\left(y - \frac{bx}{a}\right).
\] El dato sigue viajando a lo largo de las rectas
\(y - bx/a = \text{const}\), como antes, pero ahora además se amplifica
por \(e^{x/a}\) mientras viaja. Los hilos suben.

\end{example}

\begin{exercise}[]\protect\hypertarget{exr-caracteristicas-basicas}{}\label{exr-caracteristicas-basicas}

Para cada ecuación, hallar las características proyectadas, una primera
integral y la solución general.

\begin{enumerate}
\def\labelenumi{\alph{enumi}.}
\tightlist
\item
  \(u_x + 2u_y = 0\)
\item
  \(y\,u_x - x\,u_y = 0\) (¿qué son las características,
  geométricamente?)
\item
  \(u_x + e^{x}u_y = 0\)
\item
  \(x\,u_x - y\,u_y = 0\) en \(x>0\), \(y>0\)
\end{enumerate}

\end{exercise}

\begin{tcolorbox}[enhanced jigsaw, arc=.35mm, bottomrule=.15mm, bottomtitle=1mm, breakable, colback=white, colbacktitle=quarto-callout-tip-color!10!white, colframe=quarto-callout-tip-color-frame, coltitle=black, left=2mm, leftrule=.75mm, opacityback=0, opacitybacktitle=0.6, rightrule=.15mm, title=\textcolor{quarto-callout-tip-color}{\faLightbulb}\hspace{0.5em}{Solución}, titlerule=0mm, toprule=.15mm, toptitle=1mm]

En los cuatro casos el sistema característico es \(\dot x = a\),
\(\dot y = b\), \(\dot u = 0\), así que \(u\) es constante sobre cada
característica y la solución general es \(u = F(\phi)\) con \(\phi\) una
primera integral.

\textbf{a.} \(\dot x = 1\), \(\dot y = 2\), o sea \(dy/dx = 2\): rectas
\(y = 2x + k\). Primera integral \(\phi = y - 2x\), solución general
\(u = F(y-2x)\).

\textbf{b.} \(\dot x = y\), \(\dot y = -x\), o sea \(dy/dx = -x/y\), que
separa: \(y\,dy = -x\,dx\) y \(x^2+y^2 = k\). \textbf{Las
características son circunferencias centradas en el origen}, y el campo
\((y,-x)\) es el campo de velocidades de una rotación rígida. Primera
integral \(\phi = x^2+y^2\) y solución \(u = F(x^2+y^2)\): las
soluciones son exactamente las funciones con simetría de rotación, lo
cual es lo que uno esperaría de una ecuación que dice «\(u\) no cambia
si girás».

\textbf{c.} \(dy/dx = e^{x}\), o sea \(y = e^{x} + k\). Primera integral
\(\phi = y - e^{x}\) y \(u = F(y - e^{x})\).

\textbf{d.} \(dy/dx = -y/x\), que separa: \(dy/y = -dx/x\),
\(\ln y = -\ln x + k\), o sea \(xy = k\). Las características son
hipérbolas. Primera integral \(\phi = xy\) y \(u = F(xy)\).

Un chequeo rápido que conviene hacer siempre: derivar la respuesta y
reemplazar. Para \(d\), \(u_x = F'\cdot y\) y \(u_y = F'\cdot x\), así
que \(xu_x - yu_y = xyF' - yxF' = 0\). ✓

\end{tcolorbox}

\begin{exercise}[]\protect\hypertarget{exr-caracteristicas-variables}{}\label{exr-caracteristicas-variables}

Hallar las características de \(x\,u_x + t\,u_t = 0\) para \(x>0\),
\(t>0\), y explicar por qué el origen es un punto conflictivo. Comparar
con el primer ejemplo del Ejemplo~\ref{exm-coef-variables}.

\end{exercise}

\begin{tcolorbox}[enhanced jigsaw, arc=.35mm, bottomrule=.15mm, bottomtitle=1mm, breakable, colback=white, colbacktitle=quarto-callout-tip-color!10!white, colframe=quarto-callout-tip-color-frame, coltitle=black, left=2mm, leftrule=.75mm, opacityback=0, opacitybacktitle=0.6, rightrule=.15mm, title=\textcolor{quarto-callout-tip-color}{\faLightbulb}\hspace{0.5em}{Solución}, titlerule=0mm, toprule=.15mm, toptitle=1mm]

El sistema es \(\dot x = x\), \(\dot t = t\), \(\dot u = 0\). Dividiendo
las dos primeras, \(dt/dx = t/x\), que separa y da \(t = kx\):
\textbf{las características son los rayos que salen del origen}. Una
primera integral es \(\phi = t/x\), y la solución general
\(u = F(t/x)\).

El origen es conflictivo porque \textbf{todas} las características pasan
por él. Sobre cada rayo \(u\) vale una constante distinta ---el valor
que \(F\) le asigna a esa pendiente--- y todos esos valores querrían
estar en el mismo punto \((0,0)\). La única manera de que \(u\) sea
continua ahí es que \(F\) sea constante, es decir que \(u\) sea
constante en todo el cuadrante.

Es la misma patología del primer ejemplo del
Ejemplo~\ref{exm-coef-variables}, y tiene un nombre geométrico: el campo
\((x,t)\) \textbf{se anula} en el origen, así que ahí no hay ninguna
dirección característica definida y el método simplemente no dice nada.

\end{tcolorbox}

En la Figura~\ref{fig-tejido} los hilos son horizontales. ¿Qué habría
que cambiarle a la ecuación \(u_t + c\,u_x = 0\) para que los hilos
suban a medida que el observador avanza?

\begin{itemize}
\tightlist
\item
  {Poner \(c\) variable, \(c = c(x,t)\)}
\item
  {Poner un lado derecho no nulo que dependa de \(u\) o del punto}
\item
  {Cambiar la curva inicial}
\item
  {Agregar un término con \(u_{xx}\)}
\end{itemize}

La altura del hilo la gobierna \textbf{la tercera} ecuación del sistema
característico, \(du/ds = c(x,y,u)\), que es el lado derecho de la EDP.
Mientras el lado derecho sea cero, los hilos son horizontales, sin
importar lo retorcidas que sean las características proyectadas.

Poner \(c\) variable curva las trayectorias del observador ---las
sombras en el piso--- pero no cambia que \(du/ds=0\). Cambiar la curva
inicial cambia de dónde salen los hilos, no cómo suben. Y agregar
\(u_{xx}\) cambia el orden de la ecuación: ya no hay sistema
característico de este tipo.

Mirá el Ejemplo~\ref{exm-inhomogenea}, donde el lado derecho es \(u\) y
los hilos crecen exponencialmente.

\subsection{Muchas variables}\label{sec-primer-orden-n}

Nada de lo anterior usó que hubiera dos variables. Para

\begin{equation}\protect\phantomsection\label{eq-primer-orden-n}{
\sum_{i=1}^{n} a_i(x_1,\dots,x_n)\, u_{x_i} = 0 ,
}\end{equation}

las características proyectadas son las curvas integrales del campo
\(\vec a = (a_1,\dots,a_n)\) en \(\mathbb{R}^n\),

\begin{equation}\protect\phantomsection\label{eq-sistema-n}{
\frac{dx_i}{ds} = a_i(x_1,\dots,x_n), \qquad i = 1,\dots,n,
}\end{equation}

y sobre ellas \(du/ds = \sum a_i u_{x_i} = 0\): la solución sigue siendo
constante a lo largo de cada característica. Exactamente el mismo
enunciado.

Lo que sí cambia es la contabilidad de las funciones arbitrarias. En
\(\mathbb{R}^n\), una curva tiene codimensión \(n-1\), así que hacen
falta \(n-1\) primeras integrales funcionalmente independientes
\(\phi_1,\dots,\phi_{n-1}\) para etiquetar unívocamente a las
características, y la solución general es

\[
u = F\big(\phi_1,\dots,\phi_{n-1}\big),
\]

con \(F\) arbitraria de \(n-1\) variables. Para \(n=3\) el sistema
\ref{eq-sistema-n} se escribe

\[
\frac{dx}{a} = \frac{dy}{b} = \frac{dz}{c} ,
\]

y hacen falta dos primeras integrales.

\begin{example}[Tres
variables]\protect\hypertarget{exm-tres-variables}{}\label{exm-tres-variables}

Hallemos la solución general de \(\;x\,u_x + y\,u_y + z\,u_z = 0\;\) en
\(x>0\).

El sistema es \(dx/x = dy/y = dz/z\). De la primera igualdad,
\(\ln y - \ln x = \text{const}\), o sea \(\phi_1 = y/x\). De la primera
con la tercera, \(\phi_2 = z/x\). Son funcionalmente independientes en
\(x>0\), así que la solución general es

\[
u(x,y,z) = F\!\left(\frac{y}{x},\, \frac{z}{x}\right).
\]

Geométricamente: las características son los rayos que salen del origen,
y \(u\) solo puede depender de la dirección del rayo, que en
\(\mathbb{R}^3\) necesita dos números.

\end{example}

\begin{exercise}[]\protect\hypertarget{exr-tres-variables-ej}{}\label{exr-tres-variables-ej}

Hallar la solución general de \(u_x + u_y + u_z = 0\) en
\(\mathbb{R}^3\), exhibiendo dos primeras integrales funcionalmente
independientes. ¿Cuántas funciones arbitrarias tiene la respuesta, y de
cuántas variables?

\end{exercise}

\begin{tcolorbox}[enhanced jigsaw, arc=.35mm, bottomrule=.15mm, bottomtitle=1mm, breakable, colback=white, colbacktitle=quarto-callout-tip-color!10!white, colframe=quarto-callout-tip-color-frame, coltitle=black, left=2mm, leftrule=.75mm, opacityback=0, opacitybacktitle=0.6, rightrule=.15mm, title=\textcolor{quarto-callout-tip-color}{\faLightbulb}\hspace{0.5em}{Solución}, titlerule=0mm, toprule=.15mm, toptitle=1mm]

El sistema es \(\dot x = \dot y = \dot z = 1\): las características son
las rectas de dirección \((1,1,1)\). Para etiquetarlas hacen falta dos
números, y sirven las diferencias:
\[\phi_1 = x - y, \qquad \phi_2 = y - z,\] que son constantes a lo largo
de cada recta y funcionalmente independientes. La solución general es
\[u(x,y,z) = F(x-y,\; y-z),\] con \textbf{una} función arbitraria de
\textbf{dos} variables.

La contabilidad general: en \(\mathbb{R}^n\) las características son
curvas, hacen falta \(n-1\) primeras integrales para distinguirlas, y la
solución general lleva una función arbitraria de \(n-1\) variables. Acá
\(n=3\).

(Otras elecciones son igual de válidas: \(\phi_1 = x-z\) y
\(\phi_2 = y-z\) también sirven, y dan la misma familia de soluciones
escrita distinto.)

\end{tcolorbox}

\subsection*{Dos ejercicios largos}\label{dos-ejercicios-largos-1}
\addcontentsline{toc}{subsection}{Dos ejercicios largos}

Los cortos están repartidos por la sección. Estos dos conviene hacerlos
enteros, porque el primero anticipa el problema central de la sección
que viene y el segundo es la receta de la sección~\ref{sec-tejido}
aplicada de principio a fin.

\begin{exercise}[]\protect\hypertarget{exr-dato-en-otra-curva}{}\label{exr-dato-en-otra-curva}

Resolver \(u_x + u_y = 0\) con el dato \(u = \operatorname{sen}\sigma\)
sobre la curva inicial \(x = \sigma\), \(y = 2\sigma\). Después intentar
lo mismo con el dato sobre la curva \(x = \sigma\), \(y = \sigma\), y
explicar con la imagen del telar qué sale mal.

\end{exercise}

\begin{tcolorbox}[enhanced jigsaw, arc=.35mm, bottomrule=.15mm, bottomtitle=1mm, breakable, colback=white, colbacktitle=quarto-callout-tip-color!10!white, colframe=quarto-callout-tip-color-frame, coltitle=black, left=2mm, leftrule=.75mm, opacityback=0, opacitybacktitle=0.6, rightrule=.15mm, title=\textcolor{quarto-callout-tip-color}{\faLightbulb}\hspace{0.5em}{Solución}, titlerule=0mm, toprule=.15mm, toptitle=1mm]

Las características de \(u_x + u_y = 0\) salen en la dirección
\((1,1)\), así que son las rectas \(y - x = \text{const}\), y sobre cada
una \(u\) es constante.

\textbf{Primer dato, sobre \(x=\sigma\), \(y=2\sigma\).} La curva
inicial va en la dirección \((1,2)\), que no es paralela a \((1,1)\):
los hilos cruzan la curva y tejen sin problema. Sobre el punto de
parámetro \(\sigma\) vale \(y - x = 2\sigma - \sigma = \sigma\), así que
la característica que pasa por ahí es la que tiene etiqueta \(\sigma\),
y lleva el valor \(\operatorname{sen}\sigma\). Entonces
\[u(x,y) = \operatorname{sen}(y-x).\]

\textbf{Segundo dato, sobre \(x=\sigma\), \(y=\sigma\).} Esa recta
\textbf{es} una característica: va en la dirección \((1,1)\). Y acá se
rompe todo. La ecuación dice que \(u\) tiene que ser \textbf{constante}
a lo largo de esa recta; el dato dice que sobre esa misma recta \(u\)
vale \(\operatorname{sen}\sigma\), que no es constante. Las dos cosas no
pueden pasar a la vez: \textbf{no hay solución}.

\textbf{En el telar.} El hilo que sale de cada punto de la curva inicial
sale en la misma dirección que la curva, así que en lugar de despegarse
y barrer una superficie se queda pegado a ella. No hay trama que cruce
la urdimbre, y no se teje nada. Peor todavía: como la ecuación ya había
decidido cuánto tiene que valer \(u\) sobre ese hilo ---una
constante---, el dato que le pusimos encima la contradice.

Si en cambio hubiéramos dado el dato constante, \(u = 5\) sobre \(y=x\),
no habría contradicción, pero tampoco unicidad: \textbf{cualquier}
\(u = F(y-x)\) con \(F(0)=5\) resuelve el problema. Infinitas
soluciones. Los dos casos son los del Teorema~\ref{thm-degenerado}, que
vemos en la sección siguiente.

\end{tcolorbox}

\begin{exercise}[]\protect\hypertarget{exr-hilo-sube}{}\label{exr-hilo-sube}

Resolver \(u_x + u_y = 2u\) con \(u(x,0) = e^{-x^2}\), siguiendo la
receta de la sección~\ref{sec-tejido} paso a paso: parametrizar la curva
inicial, integrar el sistema característico, y despejar. Dibujar a mano
cómo se vería la Figura~\ref{fig-tejido} para esta ecuación.

\end{exercise}

\begin{tcolorbox}[enhanced jigsaw, arc=.35mm, bottomrule=.15mm, bottomtitle=1mm, breakable, colback=white, colbacktitle=quarto-callout-tip-color!10!white, colframe=quarto-callout-tip-color-frame, coltitle=black, left=2mm, leftrule=.75mm, opacityback=0, opacitybacktitle=0.6, rightrule=.15mm, title=\textcolor{quarto-callout-tip-color}{\faLightbulb}\hspace{0.5em}{Solución}, titlerule=0mm, toprule=.15mm, toptitle=1mm]

\textbf{Paso 1: parametrizar la curva inicial.} El dato está sobre la
recta \(y=0\), así que
\[x_0(\sigma) = \sigma, \qquad y_0(\sigma) = 0, \qquad u_0(\sigma) = e^{-\sigma^2}.\]

\textbf{Paso 2: el sistema característico.} Comparando con
\ref{eq-cuasilineal-1}, acá \(a=1\), \(b=1\), \(c=2u\):
\[\frac{dx}{ds} = 1, \qquad \frac{dy}{ds} = 1, \qquad \frac{du}{ds} = 2u,\]
con \(x(0)=\sigma\), \(y(0)=0\), \(u(0)=e^{-\sigma^2}\). Integrando cada
una: \[x = s + \sigma, \qquad y = s, \qquad u = e^{2s}\,e^{-\sigma^2}.\]

\textbf{Paso 3: despejar.} De la segunda, \(s = y\); de la primera,
\(\sigma = x - y\). Reemplazando en la tercera,
\[\boxed{\;u(x,y) = e^{2y}\,e^{-(x-y)^2}.\;}\]

\textbf{Verificación.} \(u_x = -2(x-y)u\) y \(u_y = 2u + 2(x-y)u\), así
que \(u_x + u_y = 2u\). ✓ Y en \(y=0\) da \(e^{-x^2}\). ✓

\textbf{El dibujo.} Las sombras en el piso son las mismas rectas de
\(45^\circ\) que en la Figura~\ref{fig-tejido}, porque el lado izquierdo
no cambió. Lo que cambia es el techo: cada hilo, en vez de mantenerse a
altura constante, \textbf{sube como \(e^{2y}\)} mientras avanza. La
campana viaja en la dirección \((1,1)\) sin deformarse y a la vez se
infla exponencialmente. Comparalo con el Ejercicio~\ref{exr-transporte},
donde el término era \(-u\) y los hilos bajaban.

\end{tcolorbox}

\section{Primer orden casi-lineales}\label{sec-casi-lineales}

Toda la maquinaria de la sección~\ref{sec-caracteristicas} fue
construida para la ecuación cuasi-lineal general \ref{eq-cuasilineal-1},
así que en principio no hay nada nuevo que aprender: se arma el sistema
característico, se largan los hilos, se teje. Lo que cambia es qué tan
bien se porta el tejido.

\begin{tcolorbox}[enhanced jigsaw, arc=.35mm, bottomrule=.15mm, breakable, colback=white, colframe=quarto-callout-warning-color-frame, left=2mm, leftrule=.75mm, opacityback=0, rightrule=.15mm, toprule=.15mm]
\begin{minipage}[t]{5.5mm}
\textcolor{quarto-callout-warning-color}{\faExclamationTriangle}
\end{minipage}%
\begin{minipage}[t]{\textwidth - 5.5mm}

\vspace{-3mm}\textbf{El nombre en el libro de la cátedra}\vspace{3mm}

González llama \textbf{«casi--lineales»} a estas ecuaciones. Es lo mismo
que nosotros y casi toda la bibliografía llamamos \emph{cuasi-lineales},
y lo que en inglés se dice \emph{quasi-linear}. Ver la tabla de nombres
de la sección~\ref{sec-panorama}.

\end{minipage}%
\end{tcolorbox}

Recordemos qué se rompió. En una ecuación lineal o semilineal, las dos
primeras ecuaciones del sistema característico

\[
\frac{dx}{ds} = a(x,y,u), \qquad \frac{dy}{ds} = b(x,y,u), \qquad \frac{du}{ds} = c(x,y,u)
\]

no involucran a \(u\), así que se resuelven aparte: \textbf{las
características proyectadas se pueden dibujar antes de conocer la
solución}, y una vez dibujadas, el dato se transporta sobre ellas. En el
caso cuasi-lineal las tres ecuaciones están acopladas. La trayectoria
del observador depende de la altura que está midiendo.

Suena a una complicación técnica. Es mucho más que eso, y esta sección
trata de qué es.

\subsection{El caso homogéneo y la ecuación de
Burgers}\label{sec-burgers}

Empecemos con la forma más simple, que es la que aparece en casi todas
las aplicaciones:

\begin{equation}\protect\phantomsection\label{eq-cuasi-homogenea}{
u_t + a(x,t,u)\,u_x = 0, \qquad u(x,0) = f(x).
}\end{equation}

El sistema característico, usando \(t\) como parámetro, es

\[
\frac{dx}{dt} = a(x,t,u), \qquad \frac{du}{dt} = 0 .
\]

La segunda ecuación dice que \textbf{\(u\) sigue siendo constante a lo
largo de cada característica}, igual que en el transporte. Pero ahora
esa constante se mete en la primera: la velocidad con la que avanza el
observador es \(a\) evaluada en el valor que él mismo está midiendo.

En el caso más importante, \(a = a(u)\) solamente, la consecuencia es
notable. Si \(u\) es constante sobre la característica, entonces
\(a(u)\) también lo es, y por lo tanto \(dx/dt = \text{const}\):
\textbf{la característica es una recta}, pero cada una con su propia
pendiente, decidida por el dato inicial en su punto de salida.

\begin{equation}\protect\phantomsection\label{eq-caract-burgers}{
x(t) = x_0 + a\big(f(x_0)\big)\,t , \qquad u = f(x_0) .
}\end{equation}

Y rectas con pendientes distintas \textbf{se cruzan}.

\begin{example}[La ecuación de Burgers sin
viscosidad]\protect\hypertarget{exm-burgers}{}\label{exm-burgers}

Es el caso \(a(u) = u\):

\begin{equation}\protect\phantomsection\label{eq-burgers}{
u_t + u\,u_x = 0, \qquad u(x,0) = f(x).
}\end{equation}

La lectura física es la de la 1: \(u\) es a la vez la magnitud
transportada y la velocidad con la que se transporta. Donde \(u\) es
grande, el perfil avanza rápido; donde es chico, avanza despacio. Las
crestas persiguen a los valles.

De \ref{eq-caract-burgers}, las características son las rectas
\(x = x_0 + f(x_0)\,t\) y sobre cada una \(u = f(x_0)\). Eliminando
\(x_0\) queda la solución en forma implícita:

\begin{equation}\protect\phantomsection\label{eq-burgers-implicita}{
u = f\big(x - u\,t\big) .
}\end{equation}

Es una fórmula rara ---la incógnita aparece de los dos lados--- pero es
una fórmula: para cada \((x,t)\) define \(u\), al menos mientras se la
pueda despejar. El teorema de la función implícita dice cuándo se puede,
y la respuesta es lo interesante.

\end{example}

\begin{theorem}[Tiempo de
quiebre]\protect\hypertarget{thm-quiebre}{}\label{thm-quiebre}

Sea \(f\) de clase \(C^1\) y acotada, y sea \(u\) la solución de
\ref{eq-burgers} dada implícitamente por \ref{eq-burgers-implicita}.
Derivando \ref{eq-burgers-implicita} respecto de \(x\),

\[
u_x = f'\cdot\big(1 - u_x t\big) \quad\Longrightarrow\quad
u_x = \frac{f'(x_0)}{1 + t\,f'(x_0)} .
\]

La solución clásica existe mientras el denominador no se anule. Si
\(f' \ge 0\) en todas partes, eso no pasa nunca y la solución es clásica
para todo \(t>0\). Si en cambio \(f'\) toma algún valor negativo, la
derivada explota en el instante

\begin{equation}\protect\phantomsection\label{eq-tiempo-quiebre}{
t^{*} = \frac{-1}{\min_{x} f'(x)} ,
}\end{equation}

que es el primer momento en que dos características se cruzan. A partir
de ahí no hay solución clásica.

\end{theorem}

Video recomendado: UNSW, \emph{How do waves break?},
\url{https://www.youtube.com/watch?v=aXuQC1qRuEM}.

La interpretación es exactamente la del dibujo: dos características que
salen de puntos distintos llevan valores distintos de \(u\), y si se
cruzan, en el punto de cruce la solución tendría que valer dos cosas a
la vez. La pendiente del perfil se va poniendo vertical hasta que, en
\(t^{*}\), la onda \textbf{rompe}.

Después de \(t^*\) no es que no pase nada: pasa que la solución deja de
ser una función y hay que ampliar el concepto, como preparamos en la
sección~\ref{sec-generalizadas}. La solución débil que continúa a esta
es una discontinuidad que viaja a la velocidad que da la condición de
Rankine--Hugoniot.

\begin{figure}

\centering{

\pandocbounded{\includegraphics[keepaspectratio]{08-edp_files/figure-pdf/fig-burgers-output-1.pdf}}

}

\caption{\label{fig-burgers}La ecuación de Burgers con dato inicial
\(f(x)=e^{-2x^2}\). Izquierda: las características son rectas, pero cada
una con la pendiente que le dicta el dato en su punto de salida; las que
salen del flanco derecho de la joroba son más lentas que las de la
cresta, y se cruzan. La primera colisión ocurre en \(t^{*}\), marcado
con la línea punteada. Derecha: el perfil se empina hasta ponerse
vertical en \(t^{*}\) y después se dobla sobre sí mismo, dejando de ser
el gráfico de una función.}

\end{figure}%

\emph{(En la versión web hay acá un gráfico interactivo con dos sliders:
la amplitud del dato inicial y el tiempo. Permite ver cómo el tiempo de
quiebre baja al subir la amplitud, y cómo el perfil se dobla sobre sí
mismo pasado ese instante.)}

\emph{(En la versión web hay acá un gráfico interactivo en tres
dimensiones: la superficie solución de Burgers tejida con sus hilos, las
características cruzándose dibujadas en el piso, y el perfil en un
instante regulable apoyado sobre la superficie. Pasado el tiempo de
quiebre se ve cómo la superficie se pliega sobre sí misma.)}

\begin{example}[Por qué hay embotellamientos sin
motivo]\protect\hypertarget{exm-trafico}{}\label{exm-trafico}

Tomemos el modelo de tráfico del Ejemplo~\ref{exm-cuasilineales}. Sea
\(u(x,t)\) la densidad de autos ---autos por kilómetro---, normalizada
de modo que \(u=1\) sea la ruta llena y \(u=0\) la ruta vacía. Un
conductor razonable maneja más despacio cuanto más llena esté la ruta,
digamos \(v(u) = v_{\max}(1-u)\). El flujo de autos es entonces
\(q(u) = u\,v(u) = v_{\max}\,u(1-u)\) y la ley de conservación de la 1
da

\[
u_t + v_{\max}(1-2u)\,u_x = 0 .
\]

Es cuasi-lineal con \(a(u) = v_{\max}(1-2u)\). Dos observaciones que son
toda la física del asunto:

\textbf{Uno.} \(a(u)\) es la velocidad de las \textbf{características},
no la de los autos. Y como \(a(u) < v(u)\) siempre ---y es negativa
cuando \(u > 1/2\)---, las perturbaciones de densidad viajan más
despacio que los autos, y en una ruta más que medio llena viajan
\textbf{hacia atrás}. Por eso uno ve venir un frenón como una onda que
avanza en contra del tránsito.

Video recomendado: New Scientist, \emph{Shockwave traffic jams recreated
for first time}, \url{https://www.youtube.com/watch?v=Suugn-p5C1M}.

\textbf{Dos.} Donde la densidad crece hacia adelante (\(u_x > 0\)), como
al acercarse a una zona congestionada, se tiene
\(a'(u)\,u_x = -2v_{\max}u_x < 0\): las características convergen y hay
quiebre. El choque es el frenón: una discontinuidad en la densidad que
se propaga hacia atrás a la velocidad de Rankine--Hugoniot. No hace
falta ningún accidente ni ningún conductor distraído para que aparezca.
Lo produce sola la ecuación.

\end{example}

Para la ecuación de Burgers con dato \(u(x,0) = f(x)\), ¿cuál de estos
datos iniciales \textbf{no} produce quiebre en ningún \(t>0\)?

\begin{itemize}
\tightlist
\item
  {\(f(x) = e^{-x^2}\)}
\item
  {\(f(x) = \arctan(x)\)}
\item
  {\(f(x) = \cos x\)}
\item
  {\(f(x) = -x\)}
\end{itemize}

Por el Teorema~\ref{thm-quiebre}, hay quiebre si y solo si \(f'\) toma
algún valor negativo, y el tiempo de quiebre es \(t^* = -1/\min f'\).

\begin{itemize}
\tightlist
\item
  \(\arctan'(x) = 1/(1+x^2) > 0\) en todas partes: \textbf{no rompe
  nunca}. Es una rarefacción pura, el perfil se estira para siempre.
\item
  \(e^{-x^2}\) tiene derivada negativa en todo el flanco derecho: rompe,
  y es el caso de la

  \begin{enumerate}
  \def\labelenumi{\arabic{enumi}.}
  \setcounter{enumi}{1}
  \tightlist
  \item
  \end{enumerate}
\item
  \(\cos x\) tiene \(\min(-\operatorname{sen} x) = -1\), así que rompe
  en \(t^*=1\).
\item
  \(f(x) = -x\) tiene \(f' = -1\) constante: rompe en \(t^*=1\), y de la
  peor manera, porque \textbf{todas} las características se cruzan en el
  mismo punto a la vez.
\end{itemize}

Tomemos la ecuación de Burgers con dato inicial \(u(x,0) = 2\cos(x)\).
¿En qué instante \(t^{*}\) rompe la solución? Escribí el número con dos
decimales.

Con \(f(x) = 2\cos x\) se tiene \(f'(x) = -2\operatorname{sen}(x)\),
cuyo mínimo es \(-2\). Entonces
\[t^* = \frac{-1}{\min f'} = \frac{-1}{-2} = \frac{1}{2}.\] Notá la
relación general: si el dato es \(A\,g(x)\) con \(g\) fija, entonces
\(\min f' = A\,\min g'\) y por lo tanto \(t^*\) es \textbf{inversamente
proporcional a la amplitud}. Eso es exactamente lo que se ve moviendo el
slider \(A\) en el gráfico interactivo de esta sección: duplicar la
altura de la joroba parte el tiempo de quiebre al medio.

\begin{exercise}[]\protect\hypertarget{exr-burgers-explicito}{}\label{exr-burgers-explicito}

Resolver \(u_t + u\,u_x = 0\) con \(u(x,0) = 1 - x\) para \(x\in[0,1]\),
extendido por \(u=1\) para \(x<0\) y \(u=0\) para \(x>1\). Dibujar las
características, hallar \(t^*\) y describir qué pasa exactamente en ese
instante. (Este es el caso extremo en que todas las características del
tramo se cruzan en un solo punto.)

\end{exercise}

\begin{tcolorbox}[enhanced jigsaw, arc=.35mm, bottomrule=.15mm, bottomtitle=1mm, breakable, colback=white, colbacktitle=quarto-callout-tip-color!10!white, colframe=quarto-callout-tip-color-frame, coltitle=black, left=2mm, leftrule=.75mm, opacityback=0, opacitybacktitle=0.6, rightrule=.15mm, title=\textcolor{quarto-callout-tip-color}{\faLightbulb}\hspace{0.5em}{Solución}, titlerule=0mm, toprule=.15mm, toptitle=1mm]

\textbf{Las características.} Salen de \(t=0\) con pendiente
\(dx/dt = f(x_0)\):

\begin{itemize}
\tightlist
\item
  de \(x_0 < 0\), donde \(f=1\): rectas \(x = x_0 + t\), inclinadas;
\item
  de \(x_0 \in [0,1]\), donde \(f = 1-x_0\): rectas
  \(x = x_0 + (1-x_0)\,t\);
\item
  de \(x_0 > 1\), donde \(f = 0\): rectas \textbf{verticales}
  \(x = x_0\).
\end{itemize}

\textbf{El tiempo de quiebre.} En la rampa, \(f' = -1\), y ese es el
mínimo, así que \[t^{*} = \frac{-1}{\min f'} = \frac{-1}{-1} = 1 .\]

\textbf{Qué pasa exactamente en \(t=1\).} Mirá las rectas de la rampa:
\(x = x_0 + (1-x_0)t = x_0(1-t) + t\). Cuando \(t=1\) eso da \(x = 1\)
\textbf{para todo \(x_0\)}: las infinitas características de la rampa no
se cruzan de a dos en momentos distintos, sino que \textbf{se juntan
todas en el mismo punto \((1,1)\) al mismo tiempo}. Es el quiebre más
violento posible.

\textbf{La solución antes del quiebre.} Despejando \(x_0 = (x-t)/(1-t)\)
y usando \(u = f(x_0) = 1-x_0\): \[
u(x,t) = \frac{1-x}{1-t} \qquad \text{para } t \le x \le 1, \; t<1,
\] con \(u=1\) a la izquierda y \(u=0\) a la derecha. La rampa se va
poniendo cada vez más empinada ---su pendiente es \(-1/(1-t)\)--- hasta
hacerse vertical en \(t=1\).

\textbf{Después.} Queda un escalón de \(1\) a \(0\), que es exactamente
el Ejercicio~\ref{exr-debil-escalon}: un choque que viaja a velocidad
\((1+0)/2 = 1/2\), saliendo de \(x=1\), o sea \(s(t) = 1 + (t-1)/2\).

\textbf{La lectura física.} Esto es una onda de compresión: material
rápido detrás alcanzando a material lento adelante. La rampa lineal es
el caso en que todo el frente colapsa de golpe, y por eso es el ejemplo
estándar para introducir choques.

\end{tcolorbox}

\begin{exercise}[]\protect\hypertarget{exr-trafico}{}\label{exr-trafico}

En el Ejemplo~\ref{exm-trafico}, tomar \(v_{\max}=1\) y el dato inicial
\[u(x,0) = \begin{cases} 0{,}2 & x<0,\\ 0{,}9 & x\ge 0,\end{cases}\] que
representa autos que circulan sueltos y se encuentran de golpe con una
zona congestionada. Hallar la velocidad del choque con la condición de
Rankine--Hugoniot y comprobar que es \textbf{negativa}: el frenón se
propaga hacia atrás. ¿A qué velocidad van los autos a cada lado?

\end{exercise}

\begin{tcolorbox}[enhanced jigsaw, arc=.35mm, bottomrule=.15mm, bottomtitle=1mm, breakable, colback=white, colbacktitle=quarto-callout-tip-color!10!white, colframe=quarto-callout-tip-color-frame, coltitle=black, left=2mm, leftrule=.75mm, opacityback=0, opacitybacktitle=0.6, rightrule=.15mm, title=\textcolor{quarto-callout-tip-color}{\faLightbulb}\hspace{0.5em}{Solución}, titlerule=0mm, toprule=.15mm, toptitle=1mm]

Con \(v_{\max}=1\), el flujo es \(F(u) = u(1-u)\). A la izquierda
\(u_-=0{,}2\) y a la derecha \(u_+=0{,}9\), así que
\[F(u_-) = 0{,}2 \cdot 0{,}8 = 0{,}16, \qquad F(u_+) = 0{,}9\cdot 0{,}1 = 0{,}09,\]
y la condición de Rankine--Hugoniot da \[
\dot s = \frac{F(u_+)-F(u_-)}{u_+-u_-} = \frac{0{,}09 - 0{,}16}{0{,}9-0{,}2}
= \frac{-0{,}07}{0{,}7} = -0{,}1 .
\] \textbf{El choque viaja hacia atrás} a un décimo de la velocidad
máxima.

Mientras tanto, los autos van hacia adelante: \(v(u) = 1-u\) da
\(v = 0{,}8\) del lado suelto y \(v = 0{,}1\) del lado congestionado.

\textbf{Lo que esto describe.} El frente del embotellamiento se corre
río arriba mientras cada auto avanza. Un conductor viaja a \(0{,}8\), se
encuentra con el frente que viene hacia él a \(0{,}1\), frena, y a
partir de ahí avanza a \(0{,}1\). Nadie se movió hacia atrás: lo que se
mueve hacia atrás es \textbf{la discontinuidad}, que no es un objeto
sino un lugar donde la densidad salta.

Es exactamente lo que se ve en el video de Nagoya, y es la razón de que
uno pueda quedar atrapado en un embotellamiento que no tiene ninguna
causa adelante: la causa pasó por ahí hace rato y siguió de largo, en el
otro sentido.

\end{tcolorbox}

\subsection{El caso general y la condición de
transversalidad}\label{sec-transversalidad}

Volvamos a la ecuación cuasi-lineal completa en dos variables,

\begin{equation}\protect\phantomsection\label{eq-cuasilineal-general}{
a_1(x,y,u)\,u_x + a_2(x,y,u)\,u_y = a(x,y,u),
}\end{equation}

con dato sobre una curva \(\Gamma\) dada por \(x = x_0(\sigma)\),
\(y = y_0(\sigma)\), \(u = u_0(\sigma)\).

Ya sabemos, de la sección~\ref{sec-tejido}, cuál es el problema: para
que los hilos tejan una superficie tienen que \textbf{cruzar} a
\(\Gamma\), no correr paralelos a ella. Ahora le ponemos una cuenta.

La dirección en la que sale el hilo, proyectada al piso, es
\((a_1, a_2)\). La dirección de la curva inicial, proyectada al piso, es
\((x_0', y_0')\). Que crucen quiere decir que esos dos vectores no sean
paralelos, es decir que el determinante que forman no se anule.

\begin{theorem}[Existencia y unicidad
local]\protect\hypertarget{thm-transversalidad}{}\label{thm-transversalidad}

Sean \(a_1\), \(a_2\) y \(a\) de clase \(C^1\) en un dominio de
\(\mathbb{R}^3\) que contiene al punto
\(\big(x_0(\sigma_0),\,y_0(\sigma_0),\,u_0(\sigma_0)\big)\), y sea
\(\Gamma\) continuamente derivable. Si en ese punto

\begin{equation}\protect\phantomsection\label{eq-transversalidad}{
\begin{vmatrix}
a_1 & x_0'(\sigma_0)\\[2pt]
a_2 & y_0'(\sigma_0)
\end{vmatrix}
= a_1\,y_0'(\sigma_0) - a_2\,x_0'(\sigma_0) \;\neq\; 0 ,
}\end{equation}

entonces existe una única solución de \ref{eq-cuasilineal-general} de
clase \(C^1\) en un entorno de
\(\big(x_0(\sigma_0), y_0(\sigma_0)\big)\) que toma sobre \(\Gamma\) los
valores \(u_0\).

La condición \ref{eq-transversalidad} se llama \textbf{condición de
transversalidad}, y dice exactamente que la proyección del campo
\(\vec V=(a_1,a_2,a)\) sobre el plano \((x,y)\) no es tangente a la
curva inicial en ese punto.

\end{theorem}

\begin{tcolorbox}[enhanced jigsaw, arc=.35mm, bottomrule=.15mm, breakable, colback=white, colframe=quarto-callout-important-color-frame, left=2mm, leftrule=.75mm, opacityback=0, rightrule=.15mm, toprule=.15mm]
\begin{minipage}[t]{5.5mm}
\textcolor{quarto-callout-important-color}{\faExclamation}
\end{minipage}%
\begin{minipage}[t]{\textwidth - 5.5mm}

\vspace{-3mm}\textbf{Es local, y eso importa}\vspace{3mm}

El teorema garantiza la solución \textbf{en un entorno} de cada punto
donde vale la condición, y nada más. No dice hasta dónde llega ese
entorno, ni que los pedazos se peguen en una solución global. La
ecuación de Burgers cumple la condición de transversalidad en todo punto
de la recta \(t=0\) ---ahí \(\vec V\) proyecta a \((f(x_0), 1)\) y la
curva inicial va en la dirección \((1,0)\), con determinante
\(-1 \ne 0\)--- y sin embargo la solución global se rompe en \(t^*\).
Las dos cosas conviven sin contradicción: el problema de Burgers está
bien planteado localmente, y deja de estarlo cuando dos entornos que
venían andando bien se pisan.

\end{minipage}%
\end{tcolorbox}

¿Y si el determinante se anula? Entonces el hilo sale paralelo a
\(\Gamma\), y pasa lo que anticipamos. Vale la pena escribirlo porque es
lo que explica todos los ejemplos raros.

\begin{theorem}[Qué pasa cuando falla la
transversalidad]\protect\hypertarget{thm-degenerado}{}\label{thm-degenerado}

Supongamos que en \(\sigma_0\) se anula el determinante
\ref{eq-transversalidad}, o sea que
\(\big(x_0'(\sigma_0), y_0'(\sigma_0)\big) = \rho\,\big(a_1, a_2\big)\)
para algún \(\rho\). Derivando el dato a lo largo de \(\Gamma\) de las
dos maneras posibles ---usando la regla de la cadena y usando la
ecuación--- se obtiene la \textbf{condición de compatibilidad}

\begin{equation}\protect\phantomsection\label{eq-compatibilidad}{
u_0'(\sigma_0) \;=\; \rho\; a\big(x_0,y_0,u_0\big) .
}\end{equation}

Entonces:

\begin{itemize}
\tightlist
\item
  Si \ref{eq-compatibilidad} \textbf{no} se cumple, \textbf{no existe}
  ninguna solución en ningún entorno del punto. La ecuación y el dato se
  contradicen.
\item
  Si se cumple en toda la curva ---es decir, si \(\Gamma\) es ella misma
  una curva característica--- entonces hay \textbf{infinitas}
  soluciones: cualquier superficie tejida que contenga a ese hilo sirve.
\end{itemize}

\end{theorem}

La cuenta de la condición de compatibilidad es corta y conviene verla,
porque es la que explica de dónde sale. Sobre \(\Gamma\), la derivada
del dato es, por la regla de la cadena,

\[
u_0'(\sigma) = x_0'\,u_x + y_0'\,u_y ,
\]

y si \((x_0',y_0') = \rho(a_1,a_2)\) eso vale
\(\rho\,(a_1u_x + a_2u_y)\), que por la ecuación es \(\rho\,a\). Así que
\(u_0'\) está \textbf{obligado} a valer \(\rho a\): la ecuación ya
decidió cómo tiene que variar \(u\) a lo largo de \(\Gamma\), y no nos
dejó elegir el dato.

\begin{example}[El ejemplo que no tiene
solución]\protect\hypertarget{exm-sin-solucion}{}\label{exm-sin-solucion}

Este es el Ejemplo 8.4.4 del libro de González, que ahí aparece con el
comentario de que «no se determinan soluciones». Con el
Teorema~\ref{thm-degenerado} se puede decir bastante más. Consideremos

\[
x\,u_x + y\,u_y = u + 1, \qquad u(x,y) = x^2 \ \text{ sobre la recta } y = x .
\]

\textbf{Paso 1: parametrizar.} \(x_0(\sigma) = \sigma\),
\(y_0(\sigma) = \sigma\), \(u_0(\sigma) = \sigma^2\).

\textbf{Paso 2: la transversalidad.} Acá \(a_1 = x\), \(a_2 = y\),
\(a = u+1\). Sobre \(\Gamma\) vale \(a_1 = a_2 = \sigma\), y
\(x_0' = y_0' = 1\), así que

\[
a_1 y_0' - a_2 x_0' = \sigma\cdot 1 - \sigma\cdot 1 = 0
\]

\textbf{para todo \(\sigma\)}. El determinante se anula en toda la
curva: \(\Gamma\) es una característica proyectada. Y se ve por qué: las
características de esta ecuación cumplen \(\dot x = x\), \(\dot y = y\),
o sea \(x = \sigma e^{s}\), \(y = \sigma e^{s}\), que son los rayos
\(y=x\) que salen del origen. Le dimos el dato justo sobre uno de ellos.

\textbf{Paso 3: la compatibilidad.} El factor de proporcionalidad es
\((x_0',y_0') = (1,1) = \tfrac{1}{\sigma}(\sigma,\sigma)\), o sea
\(\rho = 1/\sigma\). La condición \ref{eq-compatibilidad} pide

\[
u_0'(\sigma) = \rho\,a = \frac{1}{\sigma}\big(\sigma^2 + 1\big) = \sigma + \frac{1}{\sigma},
\]

mientras que el dato que elegimos tiene \(u_0'(\sigma) = 2\sigma\). Las
dos coinciden solo si \(\sigma = 1/\sigma\), o sea en
\(\sigma = \pm 1\).

\textbf{Conclusión.} Salvo en esos dos puntos aislados, la
compatibilidad falla y \textbf{el problema no tiene solución}. Y ahora
se ve que no es mala suerte: el problema no tiene solución porque le
pedimos a \(u\) que sobre el rayo \(y=x\) valiera \(\sigma^2\), cuando
la propia ecuación ya había decidido ---integrando \(\dot u = u+1\)
sobre ese rayo--- que ahí \(u\) tenía que valer \(C e^{s} - 1\). Le
dimos dos órdenes contradictorias.

\textbf{Y de yapa.} Si en vez de \(x^2\) hubiéramos puesto el dato
\(u_0(\sigma) = \lambda\sigma - 1\), que sí cumple
\(u_0' = \lambda = \rho\,a\) para todo \(\sigma\) ---verificarlo---,
entonces la curva inicial sería una curva característica entera, y
habría \textbf{infinitas} soluciones.

\end{example}

\begin{example}[Dos que sí
salen]\protect\hypertarget{exm-dos-buenos}{}\label{exm-dos-buenos}

\textbf{Uno.} \(u_x + u_y = 1\) con \(u(x,0) = f(x)\). El sistema es
\(\dot x = 1\), \(\dot y = 1\), \(\dot u = 1\). Con
\(\Gamma: x=\sigma\), \(y=0\), \(u=f(\sigma)\), la transversalidad da
\(a_1y_0' - a_2x_0' = 1\cdot 0 - 1\cdot 1 = -1 \ne 0\). Integrando,
\(x = s+\sigma\), \(y = s\), \(u = s + f(\sigma)\), y despejando
\(s = y\), \(\sigma = x - y\):

\[
u(x,y) = y + f(x-y).
\]

Los hilos suben con pendiente \(1\) mientras se desplazan: es la
Figura~\ref{fig-tejido} con el techo inclinándose.

\textbf{Dos, cuasi-lineal de verdad.} \(u\,u_x + u_y = 0\) con
\(u(x,0) = x\). Es Burgers con los nombres cambiados. El sistema es
\(\dot x = u\), \(\dot y = 1\), \(\dot u = 0\). Sobre
\(\Gamma: x=\sigma\), \(y=0\), \(u=\sigma\), el determinante vale
\(u\cdot 0 - 1\cdot 1 = -1 \ne 0\). Integrando: \(u = \sigma\),
\(y = s\), \(x = \sigma s + \sigma\). Entonces \(\sigma = x/(1+y)\) y

\[
u(x,y) = \frac{x}{1+y}.
\]

Acá \(f(x)=x\) tiene \(f' = 1 > 0\), así que por el
Teorema~\ref{thm-quiebre} no hay quiebre para \(y>0\): las
características se abren en abanico en vez de juntarse. Sí hay problema
para \(y \to -1\), donde todas se cruzan a la vez.

\end{example}

\begin{exercise}[]\protect\hypertarget{exr-transversalidad-cuentas}{}\label{exr-transversalidad-cuentas}

Para cada problema, verificar si vale la condición de transversalidad
\ref{eq-transversalidad} sobre la curva dada; si no vale, decidir con el
Teorema~\ref{thm-degenerado} si no hay solución o si hay infinitas.

\begin{enumerate}
\def\labelenumi{\alph{enumi}.}
\tightlist
\item
  \(u_x + 2u_y = u\), dato \(u = \sigma\) sobre \(x=\sigma\),
  \(y=2\sigma\)
\item
  \(u_x + 2u_y = u\), dato \(u = e^{\sigma}\) sobre \(x=\sigma\),
  \(y=2\sigma\)
\item
  \(y\,u_x - x\,u_y = 0\), dato \(u = \sigma\) sobre la circunferencia
  \(x=\cos\sigma\), \(y=\operatorname{sen}\sigma\)
\item
  \(u\,u_x + u_y = 1\), dato \(u = \sigma/2\) sobre \(x=\sigma\),
  \(y=0\)
\end{enumerate}

\end{exercise}

\begin{tcolorbox}[enhanced jigsaw, arc=.35mm, bottomrule=.15mm, bottomtitle=1mm, breakable, colback=white, colbacktitle=quarto-callout-tip-color!10!white, colframe=quarto-callout-tip-color-frame, coltitle=black, left=2mm, leftrule=.75mm, opacityback=0, opacitybacktitle=0.6, rightrule=.15mm, title=\textcolor{quarto-callout-tip-color}{\faLightbulb}\hspace{0.5em}{Solución}, titlerule=0mm, toprule=.15mm, toptitle=1mm]

La cuenta es siempre la misma: calcular \(a_1y_0' - a_2x_0'\); si da
cero, hallar \(\rho\) y comparar \(u_0'\) con \(\rho\,a\).

\textbf{a.} \(a_1=1\), \(a_2=2\), \(a=u\). La curva va en la dirección
\((1,2)\), así que \(a_1y_0' - a_2x_0' = 1\cdot 2 - 2\cdot 1 = 0\):
\textbf{falla} la transversalidad, la curva inicial es una
característica proyectada. De \((1,2) = \rho\,(1,2)\) sale \(\rho = 1\),
y la compatibilidad pide \(u_0' = \rho\,a = u_0 = \sigma\). Pero el dato
tiene \(u_0' = 1\). Solo coinciden en \(\sigma=1\): \textbf{no hay
solución} en ningún entorno de los demás puntos.

\textbf{b.} Misma ecuación, misma curva, mismo \(\rho=1\); lo que cambia
es el dato. Ahora \(u_0 = e^{\sigma}\) y \(u_0' = e^{\sigma} = u_0\): la
compatibilidad se cumple \textbf{en toda la curva}. Eso significa que la
curva inicial no es solo una característica proyectada sino una curva
característica entera, y entonces hay \textbf{infinitas soluciones}: la
ecuación no aporta ninguna información nueva sobre ese hilo, y cualquier
tejido que lo contenga sirve.

Comparar \(a\) con \(b\) es el punto del ejercicio: la misma ecuación y
la misma curva pueden dar «ninguna» o «infinitas» según qué dato se les
ponga encima.

\textbf{c.} \(a_1=y\), \(a_2=-x\), \(a=0\). Sobre la circunferencia,
\(x_0'=-\operatorname{sen}\sigma\) y \(y_0'=\cos\sigma\), así que
\[a_1y_0' - a_2x_0' = \operatorname{sen}\sigma\cos\sigma - \cos\sigma\operatorname{sen}\sigma = 0 .\]
Falla, y era previsible: por el
Ejercicio~\ref{exr-caracteristicas-basicas}, las características de esta
ecuación \textbf{son} las circunferencias centradas en el origen, y le
dimos el dato justo sobre una. Con \(\rho = -1\), la compatibilidad pide
\(u_0' = -1\cdot 0 = 0\), o sea que \(u\) sea constante sobre la
circunferencia; pero pedimos \(u=\sigma\), que no lo es. \textbf{No hay
solución.}

\textbf{d.} \(a_1=u\), \(a_2=1\), \(a=1\). Sobre \(x=\sigma\), \(y=0\):
\(x_0'=1\), \(y_0'=0\), y \(a_1y_0' - a_2x_0' = 0 - 1 = -1 \neq 0\).
\textbf{Vale la transversalidad} en todo punto, así que por el
Teorema~\ref{thm-transversalidad} hay una única solución local. (Es el
caso general del corolario que mencionamos: cuando el dato está sobre la
recta \(y=0\) y el coeficiente de \(u_y\) es \(1\), la transversalidad
se cumple siempre.)

\end{tcolorbox}

\begin{exercise}[]\protect\hypertarget{exr-hilo-cuasilineal}{}\label{exr-hilo-cuasilineal}

La ecuación \(u\,u_x + u_y = 0\) del Ejemplo~\ref{exm-dos-buenos}, con
dato \(u(x,0)=x\), tiene solución \(u = x/(1+y)\). Verificar que
efectivamente resuelve la ecuación, y explicar en términos de los hilos
qué pasa cuando \(y \to -1^{+}\): ¿por qué todas las características se
juntan ahí, y qué tiene eso que ver con el Teorema~\ref{thm-quiebre}?

\end{exercise}

\begin{tcolorbox}[enhanced jigsaw, arc=.35mm, bottomrule=.15mm, bottomtitle=1mm, breakable, colback=white, colbacktitle=quarto-callout-tip-color!10!white, colframe=quarto-callout-tip-color-frame, coltitle=black, left=2mm, leftrule=.75mm, opacityback=0, opacitybacktitle=0.6, rightrule=.15mm, title=\textcolor{quarto-callout-tip-color}{\faLightbulb}\hspace{0.5em}{Solución}, titlerule=0mm, toprule=.15mm, toptitle=1mm]

\textbf{Verificación.} Con \(u = x/(1+y)\):
\[u_x = \frac{1}{1+y}, \qquad u_y = -\frac{x}{(1+y)^2},\] y entonces
\[u\,u_x + u_y = \frac{x}{1+y}\cdot\frac{1}{1+y} - \frac{x}{(1+y)^2} = 0 . \ \checkmark\]
Y en \(y=0\) da \(u = x\). ✓

\textbf{Qué pasa en \(y \to -1^{+}\).} Las características son
\(x = \sigma(1+y)\), donde \(\sigma\) es el valor de \(u\) que
transportan. Cuando \(y \to -1\), \textbf{todas} dan \(x \to 0\): los
hilos, que en \(y=0\) estaban repartidos por todo el eje, convergen
todos al mismo punto. Y como cada uno lleva un valor distinto de \(u\),
la solución no puede estar definida ahí: efectivamente
\(u = x/(1+y) \to \infty\) para cualquier \(x \neq 0\).

\textbf{La conexión con el Teorema~\ref{thm-quiebre}.} El teorema está
escrito con el tiempo corriendo hacia adelante, y acá el papel del
tiempo lo hace \(y\). Como \(f(x)=x\) tiene \(f' = 1 > 0\), hacia
adelante no hay quiebre: los hilos se \textbf{abren} en abanico y la
solución vive para todo \(y>0\). Pero la fórmula del teorema,
\[u_x = \frac{f'}{1 + y\,f'} = \frac{1}{1+y},\] tiene su denominador
nulo en \(y=-1\). O sea que el quiebre existe, solo que \textbf{hacia
atrás en el tiempo}: la solución que vemos abrirse viene de haber estado
toda concentrada en un punto.

Es un buen recordatorio de que las características se cruzan en la
dirección en la que convergen, y que con \(f'>0\) esa dirección es el
pasado.

\end{tcolorbox}

\subsection{Una nota sobre la solución
general}\label{sec-ecuacion-general}

Para las ecuaciones cuasi-lineales existe también un análogo del
Teorema~\ref{thm-general-homogenea}, con una salvedad importante. Si
\(\phi_1(x,y,u)\) y \(\phi_2(x,y,u)\) son dos primeras integrales
funcionalmente independientes del sistema característico ---es decir,
dos funciones constantes a lo largo de cada curva característica---,
entonces toda relación de la forma

\begin{equation}\protect\phantomsection\label{eq-ecuacion-general}{
F\big(\phi_1(x,y,u),\,\phi_2(x,y,u)\big) = 0 ,
}\end{equation}

con \(F\) arbitraria de clase \(C^1\), define implícitamente soluciones
de \ref{eq-cuasilineal-general}.

La salvedad, que Dallos subraya y conviene repetir: \textbf{no toda
solución se obtiene de \ref{eq-ecuacion-general}}. Por eso a esa
expresión se la llama \emph{ecuación general} de la EDP y no
\emph{solución general}, a diferencia del caso lineal homogéneo, donde
\(u = F(\phi)\) sí las agota todas.

\begin{example}[Un
ejemplo]\protect\hypertarget{exm-ecuacion-general}{}\label{exm-ecuacion-general}

Hallemos la ecuación general de \(\;xy\,u_x - xy\,u_y = (x-y)\,u\).

El sistema característico es \[
\frac{dx}{xy} = \frac{dy}{-xy} = \frac{du}{(x-y)u} .
\] De la primera igualdad, \(dx = -dy\), y por lo tanto \(\phi_1 = x+y\)
es constante a lo largo de cada característica. Para la segunda no hay
un par de fracciones que se relacionen solas, así que conviene probar
candidatos. Parametrizando el sistema como \(\dot x = xy\),
\(\dot y = -xy\), \(\dot u = (x-y)u\) y probando con \(\phi_2 = xyu\):
\[
\frac{d(xyu)}{ds} = yu\,\dot x + xu\,\dot y + xy\,\dot u
= yu(xy) + xu(-xy) + xy(x-y)u = xyu\,(y - x) + xyu\,(x-y) = 0 .
\] Entonces \(\phi_2 = xyu\) también es primera integral, y la ecuación
general es \(F(x+y,\,xyu) = 0\). Eligiendo \(F(m,n) = m-n\) sale la
solución particular \(u = (x+y)/(xy)\); despejando en general,
\(u = \dfrac{1}{xy}\,g(x+y)\) con \(g\) arbitraria.

\end{example}

\subsection{Más allá: completamente no
lineales}\label{sec-no-lineales-1}

Cerramos con lo que González menciona al pasar y vale la pena tener
ubicado. Existen ecuaciones de primer orden importantísimas en física
que no son cuasi-lineales, y la más notable es la de
\textbf{Hamilton--Jacobi}:

\begin{equation}\protect\phantomsection\label{eq-hamilton-jacobi}{
u_t + H\big(u_x, u_y, x, y, t\big) = 0 ,
}\end{equation}

que reformula toda la mecánica clásica: \(u\) es la función principal de
Hamilton y \(H\) el hamiltoniano del sistema. Su generalización
ondulatoria fue el puente histórico entre la teoría corpuscular y la
ondulatoria de la materia, y de ahí salió la ecuación de Schrödinger.

También se resuelve por características, pero con la vuelta de tuerca
que anticipamos en la sección~\ref{sec-no-lineales}: como la ecuación
mezcla las derivadas de manera no lineal, no alcanza con propagar el
punto y el valor de \(u\). Hay que propagar además \(p = u_x\) y
\(q = u_y\), y el objeto que viaja ---punto, valor y plano tangente---
se llama \textbf{banda característica}. Para \ref{eq-hamilton-jacobi},
ese sistema ampliado resulta ser exactamente \textbf{las ecuaciones de
Hamilton} del sistema mecánico, y las características son sus
trayectorias en el espacio de fases. Es uno de los resultados más lindos
de la mecánica analítica.

Zachmanoglou y Thoe lo tratan en el capítulo III; Evans, con todo
detalle, en el capítulo 3.

\subsection*{Un ejercicio largo}\label{un-ejercicio-largo-1}
\addcontentsline{toc}{subsection}{Un ejercicio largo}

\begin{exercise}[]\protect\hypertarget{exr-general}{}\label{exr-general}

Consideremos \(\;(y+u)\,u_x + y\,u_y = x-y\).

\begin{enumerate}
\def\labelenumi{\alph{enumi}.}
\tightlist
\item
  Verificar que \(\phi_1 = \dfrac{x+u}{y}\) es una primera integral.
  \emph{(Sugerencia: en el sistema
  \(\frac{dx}{y+u} = \frac{dy}{y} = \frac{du}{x-y}\), sumar el numerador
  y el denominador de la primera y la tercera fracción; el truco
  funciona porque \(\frac{dx}{P}=\frac{du}{R}\) implica
  \(\frac{d(x+u)}{P+R}\), y acá \(P+R = x+u\).)}
\item
  Verificar que \(\phi_2 = y\,(y+u-x)\) también lo es, derivándola a lo
  largo del sistema característico.
\item
  Escribir la ecuación general de la EDP, y obtener de ella al menos una
  solución explícita.
\end{enumerate}

\end{exercise}

\begin{tcolorbox}[enhanced jigsaw, arc=.35mm, bottomrule=.15mm, bottomtitle=1mm, breakable, colback=white, colbacktitle=quarto-callout-tip-color!10!white, colframe=quarto-callout-tip-color-frame, coltitle=black, left=2mm, leftrule=.75mm, opacityback=0, opacitybacktitle=0.6, rightrule=.15mm, title=\textcolor{quarto-callout-tip-color}{\faLightbulb}\hspace{0.5em}{Solución}, titlerule=0mm, toprule=.15mm, toptitle=1mm]

Parametrizamos el sistema característico:
\[\dot x = y+u, \qquad \dot y = y, \qquad \dot u = x-y .\]

\textbf{a.} Sumando la primera y la tercera:
\(\dot x + \dot u = (y+u) + (x-y) = x+u\), o sea
\(\dfrac{d(x+u)}{ds} = x+u\). Y por otro lado \(\dot y = y\). Las dos
cantidades \(x+u\) e \(y\) satisfacen la misma ecuación, así que su
cociente es constante: \[
\frac{d}{ds}\left(\frac{x+u}{y}\right)
= \frac{(\dot x + \dot u)\,y - (x+u)\,\dot y}{y^2}
= \frac{(x+u)y - (x+u)y}{y^2} = 0 . \ \checkmark
\]

\textbf{b.} Derivando directamente: \[
\frac{d}{ds}\Big[y\,(y+u-x)\Big] = \dot y\,(y+u-x) + y\,(\dot y + \dot u - \dot x) .
\] El primer paréntesis del segundo término es
\(y + (x-y) - (y+u) = x - y - u = -(y+u-x)\). Entonces \[
= y\,(y+u-x) - y\,(y+u-x) = 0 . \ \checkmark
\]

\textbf{c.} Con dos primeras integrales funcionalmente independientes,
la ecuación general es \[
F\!\left(\frac{x+u}{y},\; y\,(y+u-x)\right) = 0 ,
\] con \(F\) arbitraria de clase \(C^1\). Eligiendo \(F(m,n) = m - k\)
con \(k\) constante sale \((x+u)/y = k\), es decir la familia de
soluciones \[u(x,y) = k\,y - x ,\] que se verifica en un renglón:
\(u_x = -1\), \(u_y = k\), y
\((y+u)(-1) + y\,k = -y - ky + x + yk = x - y\). ✓

Recordá la advertencia de la sección~\ref{sec-ecuacion-general}: esa
expresión se llama \emph{ecuación general} y no \emph{solución general}
porque no toda solución sale de ella.

\end{tcolorbox}

\section{Segundo orden: la clasificación}\label{sec-clasificacion}

Las ecuaciones de la física matemática son casi todas de segundo orden,
y casi todas semilineales. En dos variables, la forma general es la
\ref{eq-semilineal-2var} de la sección~\ref{sec-semilineales}, que
repetimos acá:

\begin{equation}\protect\phantomsection\label{eq-general}{
A\,u_{xx} + 2B\,u_{xy} + C\,u_{yy} = F(x, y, u, u_x, u_y),
}\end{equation}

con \(A\), \(B\), \(C\) funciones de \((x,y)\) que no se anulan
simultáneamente. Toda la geometría del problema está en el lado
izquierdo ---en la parte principal--- y, como vamos a ver, en el signo
de una sola cantidad.

\begin{tcolorbox}[enhanced jigsaw, arc=.35mm, bottomrule=.15mm, breakable, colback=white, colframe=quarto-callout-warning-color-frame, left=2mm, leftrule=.75mm, opacityback=0, rightrule=.15mm, toprule=.15mm]
\begin{minipage}[t]{5.5mm}
\textcolor{quarto-callout-warning-color}{\faExclamationTriangle}
\end{minipage}%
\begin{minipage}[t]{\textwidth - 5.5mm}

\vspace{-3mm}\textbf{El factor 2, que es la fuente de errores número uno}\vspace{3mm}

Escribimos el término cruzado como \(2B\,u_{xy}\), siguiendo a González.
Muchos otros textos ---Dallos, Strauss, y buena parte de la bibliografía
en inglés--- lo escriben \(B\,u_{xy}\) a secas. Las dos convenciones
conviven y hay que saber en cuál está uno:

{\def\LTcaptype{none} % do not increment counter
\begin{longtable}[]{@{}
  >{\raggedright\arraybackslash}p{(\linewidth - 6\tabcolsep) * \real{0.2500}}
  >{\raggedright\arraybackslash}p{(\linewidth - 6\tabcolsep) * \real{0.2500}}
  >{\raggedright\arraybackslash}p{(\linewidth - 6\tabcolsep) * \real{0.2500}}
  >{\raggedright\arraybackslash}p{(\linewidth - 6\tabcolsep) * \real{0.2500}}@{}}
\toprule\noalign{}
\begin{minipage}[b]{\linewidth}\raggedright
\end{minipage} & \begin{minipage}[b]{\linewidth}\raggedright
Término cruzado
\end{minipage} & \begin{minipage}[b]{\linewidth}\raggedright
Discriminante
\end{minipage} & \begin{minipage}[b]{\linewidth}\raggedright
Ejemplo: \(3u_{xx}+2u_{xy}+u_{yy}\)
\end{minipage} \\
\midrule\noalign{}
\endhead
\bottomrule\noalign{}
\endlastfoot
\textbf{Nosotros y González} & \(2B\,u_{xy}\) & \(\Delta = B^2 - AC\) &
\(B=1\), \(\Delta = 1-3 = -2\) \\
Dallos, Strauss y otros & \(B\,u_{xy}\) & \(\Delta = B^2 - 4AC\) &
\(B=2\), \(\Delta = 4-12 = -8\) \\
\end{longtable}
}

Los números son distintos ---el segundo es cuatro veces el primero---
pero \textbf{el signo es el mismo}, y por lo tanto el tipo también. Lo
que no se puede hacer es mezclar: leer el coeficiente con una convención
y aplicar la fórmula de la otra. El primer ejercicio de esta sección es
exactamente esa trampa.

La razón de que exista el \(2\) es buena y la vamos a ver enseguida:
hace que \(A\), \(B\), \(C\) sean directamente las entradas de una
matriz simétrica.

\end{minipage}%
\end{tcolorbox}

\begin{definition}[Discriminante]\protect\hypertarget{def-discriminante}{}\label{def-discriminante}

Se llama \textbf{discriminante} de \ref{eq-general} a
\(\;\Delta = B^2 - AC\).

\end{definition}

\begin{theorem}[Clasificación]\protect\hypertarget{thm-clasificacion}{}\label{thm-clasificacion}

En un entorno de un punto donde \(A\), \(B\) y \(C\) son continuas, la
ecuación \ref{eq-general} se clasifica según el signo de
\(\Delta = B^2 - AC\):

{\def\LTcaptype{none} % do not increment counter
\begin{longtable}[]{@{}llll@{}}
\toprule\noalign{}
\(\Delta\) & Tipo & Características reales & Ecuación modelo \\
\midrule\noalign{}
\endhead
\bottomrule\noalign{}
\endlastfoot
\(>0\) & hiperbólica & dos familias & \(u_{tt} - c^2 u_{xx} = 0\) \\
\(=0\) & parabólica & una familia & \(u_t - k\,u_{xx} = 0\) \\
\(<0\) & elíptica & ninguna & \(u_{xx} + u_{yy} = 0\) \\
\end{longtable}
}

\end{theorem}

Antes de demostrarlo conviene entender de dónde salen los nombres,
porque el nombre \textbf{es} la demostración.

Antes de los ejercicios de lápiz y papel, uno para hacer mentalmente. La
ecuación

\[3\,u_{xx} + 2\,u_{xy} + u_{yy} = 0\]

es, en todo el plano:

\begin{itemize}
\tightlist
\item
  {hiperbólica}
\item
  {parabólica}
\item
  {elíptica}
\item
  {de tipo variable según la región}
\end{itemize}

La trampa está en la \ref{eq-general}: el coeficiente del término
cruzado es \(2B\), no \(B\). Acá \(2B = 2\), así que \(B = 1\), y no
\(2\).

Entonces \(A = 3\), \(B = 1\), \(C = 1\), y el discriminante vale
\(\Delta = B^2 - AC = 1 - 3 = -2 < 0\). Elíptica, y como \(A\), \(B\) y
\(C\) son constantes, lo es en todo el plano. Si uno se olvida del
factor y toma \(B = 2\), obtiene \(\Delta = 4 - 3 = 1 > 0\) y concluye
``hiperbólica'': el error más común de todo el capítulo.

Con la otra convención el resultado es el mismo: \(B_{\text{otra}} = 2\)
y \(\Delta_{\text{otra}} = B^2 - 4AC = 4 - 12 = -8 < 0\). Mismo signo,
mismo tipo.

\subsection{De dónde salen los nombres: cónicas}\label{sec-conicas}

Busquemos, como haremos en detalle en la
sección~\ref{sec-formas-canonicas}, un cambio de variables
\(\xi = \xi(x,y)\) que simplifique \ref{eq-general}. Al hacer el cambio
y agrupar, el coeficiente que acompaña a \(u_{\xi\xi}\) resulta ser

\begin{equation}\protect\phantomsection\label{eq-forma-cuadratica}{
A\,\xi_x^2 + 2B\,\xi_x\xi_y + C\,\xi_y^2 ,
}\end{equation}

---la cuenta está en la sección~\ref{sec-formas-canonicas}, pero se
puede anticipar: \(u_{\xi\xi}\) solo puede venir de aplicar la regla de
la cadena dos veces, y cada vez aporta un factor \(\xi_x\) o
\(\xi_y\)---. Y \ref{eq-forma-cuadratica} no es cualquier expresión:
\textbf{es una forma cuadrática en \((\xi_x, \xi_y)\)}, con exactamente
los mismos coeficientes que la parte principal.

\begin{definition}[Forma cuadrática y matriz
asociadas]\protect\hypertarget{def-forma-cuadratica-asociada}{}\label{def-forma-cuadratica-asociada}

A la parte principal de \ref{eq-general} le asociamos la forma
cuadrática

\begin{equation}\protect\phantomsection\label{eq-Q}{
Q(\xi,\eta) = A\,\xi^2 + 2B\,\xi\eta + C\,\eta^2
= \begin{pmatrix} \xi & \eta \end{pmatrix}
\begin{pmatrix} A & B \\ B & C \end{pmatrix}
\begin{pmatrix} \xi \\ \eta \end{pmatrix},
}\end{equation}

y la \textbf{matriz de coeficientes}
\(M = \begin{pmatrix} A & B \\ B & C\end{pmatrix}\), que es simétrica.
Notar que

\[
\det M = AC - B^2 = -\Delta .
\]

\end{definition}

Acá está el motivo del factor \(2\): con la convención de González,
\(A\), \(B\), \(C\) son directamente las entradas de \(M\), sin dividir
nada.

Ahora la pregunta que decide todo. Recordemos de la
sección~\ref{sec-caracteristicas} que las características son las curvas
donde la ecuación degenera. Buscar un cambio de variables que
\textbf{anule} el coeficiente de \(u_{\xi\xi}\) es pedir

\begin{equation}\protect\phantomsection\label{eq-direcciones-nulas}{
Q(\xi_x, \xi_y) = 0 ,
}\end{equation}

es decir, buscar las \textbf{direcciones nulas} de la forma cuadrática.
Dividiendo por \(\xi_y^2\) y llamando \(m = \xi_x/\xi_y\), eso es una
ecuación de segundo grado,

\[
A\,m^2 + 2B\,m + C = 0 \qquad\Longrightarrow\qquad m = \frac{-B \pm \sqrt{B^2 - AC}}{A} ,
\]

y \textbf{el radicando es el discriminante}. Ahí está todo:

\begin{itemize}
\tightlist
\item
  si \(\Delta > 0\) hay \textbf{dos} direcciones nulas reales y
  distintas, y por lo tanto dos familias de características;
\item
  si \(\Delta = 0\) hay \textbf{una} sola, doble;
\item
  si \(\Delta < 0\) \textbf{no hay ninguna} real: las raíces son
  complejas conjugadas.
\end{itemize}

\begin{tcolorbox}[enhanced jigsaw, arc=.35mm, bottomrule=.15mm, breakable, colback=white, colframe=quarto-callout-important-color-frame, left=2mm, leftrule=.75mm, opacityback=0, rightrule=.15mm, toprule=.15mm]
\begin{minipage}[t]{5.5mm}
\textcolor{quarto-callout-important-color}{\faExclamation}
\end{minipage}%
\begin{minipage}[t]{\textwidth - 5.5mm}

\vspace{-3mm}\textbf{Y los nombres son los de las cónicas}\vspace{3mm}

Miremos ahora el conjunto de nivel \(Q(\xi,\eta) = 1\), que es una curva
del plano \((\xi,\eta)\). Es exactamente la parte cuadrática de una
cónica

\[
A\xi^2 + 2B\xi\eta + C\eta^2 + D\xi + E\eta + F = 0 ,
\]

y de la geometría analítica sabemos que su tipo lo decide el mismo
número: hipérbola si \(B^2 - AC > 0\), parábola si \(=0\), elipse si
\(<0\).

\textbf{La clasificación de las EDP de segundo orden no se parece a la
de las cónicas: es la misma.} Los nombres «hiperbólica», «parabólica» y
«elíptica» no son una analogía ni una metáfora, son el nombre de la
cónica que le corresponde a la parte principal de la ecuación. Y lo que
en la cónica son las \textbf{asíntotas} de la hipérbola, en la EDP son
las \textbf{direcciones características}: las direcciones donde \(Q\) se
anula.

Una honestidad sobre la analogía: mirando solo la parte cuadrática, las
tres curvas que salen son cónicas \emph{degeneradas} ---dos rectas
cruzadas, dos rectas paralelas, o un punto---, y hacen falta los
términos lineales \(D\xi + E\eta + F\) para que aparezca una parábola de
verdad. En la EDP, esos términos lineales son los de orden menor, los
que están escondidos en \(F\). La correspondencia entre las dos
clasificaciones es exacta al nivel de la parte cuadrática, que es justo
el nivel que decide el tipo.

\end{minipage}%
\end{tcolorbox}

\begin{figure}

\centering{

\pandocbounded{\includegraphics[keepaspectratio]{08-edp_files/figure-pdf/fig-conicas-output-1.pdf}}

}

\caption{\label{fig-conicas}Las curvas de nivel \(Q(\xi,\eta)=\pm 1\) de
la forma cuadrática asociada a cada ecuación modelo, y sus direcciones
nulas en línea de trazos: las direcciones donde \(Q\) vale cero. Son las
que deciden todo. En el caso hiperbólico son las asíntotas de la
hipérbola ---dos direcciones reales, dos familias de características---;
en el parabólico las dos se juntan en una sola, doble; en el elíptico la
curva es una elipse, que no tiene asíntotas, y no hay ninguna dirección
real donde \(Q\) se anule.}

\end{figure}%

\begin{exercise}[]\protect\hypertarget{exr-conica-dibujo}{}\label{exr-conica-dibujo}

Dibujar a mano la cónica \(Q(\xi,\eta)=1\) asociada a cada una de estas
ecuaciones, marcando sus direcciones nulas si las tiene, y decir el
tipo:

\begin{enumerate}
\def\labelenumi{\alph{enumi}.}
\tightlist
\item
  \(u_{xx} - u_{yy} = 0\)
\item
  \(u_{xy} = 0\)
\item
  \(u_{xx} + 4u_{xy} + 4u_{yy} = 0\)
\item
  \(2u_{xx} + 2u_{xy} + 3u_{yy} = 0\)
\end{enumerate}

(El caso \(b\) es instructivo: \(A = C = 0\) y \(2B = 1\). ¿Qué cónica
es \(\xi\eta = \text{const}\), y cuáles son sus asíntotas?)

\end{exercise}

\begin{tcolorbox}[enhanced jigsaw, arc=.35mm, bottomrule=.15mm, bottomtitle=1mm, breakable, colback=white, colbacktitle=quarto-callout-tip-color!10!white, colframe=quarto-callout-tip-color-frame, coltitle=black, left=2mm, leftrule=.75mm, opacityback=0, opacitybacktitle=0.6, rightrule=.15mm, title=\textcolor{quarto-callout-tip-color}{\faLightbulb}\hspace{0.5em}{Solución}, titlerule=0mm, toprule=.15mm, toptitle=1mm]

\textbf{a.} \(A=1\), \(B=0\), \(C=-1\), \(\Delta = 1 > 0\):
\textbf{hiperbólica}. La forma es \(Q = \xi^2 - \eta^2\), y \(Q=1\) es
una hipérbola con las asíntotas \(\xi = \pm\eta\), a \(45^\circ\). Es la
ecuación de onda con \(c=1\), y esas dos direcciones son las que dan las
características \(x \pm y = \text{const}\).

\textbf{b.} Acá hay que tener cuidado: \(u_{xy}\) significa
\(2B\,u_{xy}\) con \(2B=1\), o sea \(B = 1/2\), y \(A = C = 0\).
Entonces \(\Delta = 1/4 > 0\): \textbf{hiperbólica}. La forma es
\(Q = \xi\eta\), y \(Q = 1\) es una hipérbola equilátera \textbf{cuyas
asíntotas son los propios ejes} \(\xi=0\) y \(\eta=0\).

Y eso cierra un círculo lindo: las direcciones nulas son las de los
ejes, así que las características son las rectas \(x=\text{const}\) e
\(y=\text{const}\). Que es exactamente lo que uno esperaba, porque ya
sabemos de la sección~\ref{sec-porque} que la solución general de
\(u_{xy}=0\) es \(u = f(x) + g(y)\): una función arbitraria por cada
familia de características. Esta ecuación \textbf{ya está en forma
canónica hiperbólica}.

\textbf{c.} \(A=1\), \(2B=4\) o sea \(B=2\), \(C=4\), y
\(\Delta = 4-4 = 0\): \textbf{parabólica}. La forma es
\(Q = \xi^2+4\xi\eta+4\eta^2 = (\xi+2\eta)^2\), un cuadrado perfecto.
\(Q=1\) son las dos rectas paralelas \(\xi+2\eta = \pm 1\) ---la cónica
degenerada de la que hablábamos--- y hay una sola dirección nula,
\(\xi = -2\eta\), doble. Una sola familia de características.

\textbf{d.} \(A=2\), \(B=1\), \(C=3\), \(\Delta = 1-6 = -5 < 0\):
\textbf{elíptica}. \(Q\) es definida positiva y \(Q=1\) es una elipse
inclinada, sin direcciones nulas reales. Con el gráfico interactivo de
esta sección se puede reproducir: poné \(A=2\), \(B=1\), \(C=3\).

\end{tcolorbox}

\subsection{Las rotaciones}\label{sec-rotaciones}

La otra mitad de la historia de las cónicas también se traslada. Sabemos
de álgebra lineal que una cónica con término cruzado se lleva a sus ejes
principales \textbf{rotando} el plano, y que la rotación diagonaliza la
matriz simétrica \(M\): en los ejes nuevos, la forma queda

\[
Q = \lambda_1\,\tilde\xi^{\,2} + \lambda_2\,\tilde\eta^{\,2} ,
\]

con \(\lambda_1\), \(\lambda_2\) los autovalores de \(M\), que son
reales por ser \(M\) simétrica. Y como
\(\lambda_1\lambda_2 = \det M = -\Delta\), los signos de los autovalores
dicen lo mismo que el discriminante:

\begin{longtable}[]{@{}
  >{\raggedright\arraybackslash}p{(\linewidth - 8\tabcolsep) * \real{0.2200}}
  >{\centering\arraybackslash}p{(\linewidth - 8\tabcolsep) * \real{0.1800}}
  >{\centering\arraybackslash}p{(\linewidth - 8\tabcolsep) * \real{0.1200}}
  >{\raggedright\arraybackslash}p{(\linewidth - 8\tabcolsep) * \real{0.1800}}
  >{\raggedright\arraybackslash}p{(\linewidth - 8\tabcolsep) * \real{0.3000}}@{}}
\caption{La clasificación leída con autovalores. Es la versión que se
generaliza a muchas variables.}\label{tbl-autovalores}\tabularnewline
\toprule\noalign{}
\begin{minipage}[b]{\linewidth}\raggedright
Autovalores
\end{minipage} & \begin{minipage}[b]{\linewidth}\centering
\(\det M = \lambda_1\lambda_2\)
\end{minipage} & \begin{minipage}[b]{\linewidth}\centering
\(\Delta\)
\end{minipage} & \begin{minipage}[b]{\linewidth}\raggedright
Cónica
\end{minipage} & \begin{minipage}[b]{\linewidth}\raggedright
Tipo de EDP
\end{minipage} \\
\midrule\noalign{}
\endfirsthead
\toprule\noalign{}
\begin{minipage}[b]{\linewidth}\raggedright
Autovalores
\end{minipage} & \begin{minipage}[b]{\linewidth}\centering
\(\det M = \lambda_1\lambda_2\)
\end{minipage} & \begin{minipage}[b]{\linewidth}\centering
\(\Delta\)
\end{minipage} & \begin{minipage}[b]{\linewidth}\raggedright
Cónica
\end{minipage} & \begin{minipage}[b]{\linewidth}\raggedright
Tipo de EDP
\end{minipage} \\
\midrule\noalign{}
\endhead
\bottomrule\noalign{}
\endlastfoot
mismo signo & \(>0\) & \(<0\) & elipse & elíptica \\
signos opuestos & \(<0\) & \(>0\) & hipérbola & hiperbólica \\
alguno nulo & \(=0\) & \(=0\) & degenerada & parabólica \\
\end{longtable}

\begin{tcolorbox}[enhanced jigsaw, arc=.35mm, bottomrule=.15mm, breakable, colback=white, colframe=quarto-callout-tip-color-frame, left=2mm, leftrule=.75mm, opacityback=0, rightrule=.15mm, toprule=.15mm]
\begin{minipage}[t]{5.5mm}
\textcolor{quarto-callout-tip-color}{\faLightbulb}
\end{minipage}%
\begin{minipage}[t]{\textwidth - 5.5mm}

\vspace{-3mm}\textbf{Rotar la cónica es hallar la forma canónica}\vspace{3mm}

Vale la pena decir explícitamente el paralelo, porque es la mejor manera
de entender qué va a hacer la sección~\ref{sec-formas-canonicas}.

{\def\LTcaptype{none} % do not increment counter
\begin{longtable}[]{@{}
  >{\raggedright\arraybackslash}p{(\linewidth - 2\tabcolsep) * \real{0.5000}}
  >{\raggedright\arraybackslash}p{(\linewidth - 2\tabcolsep) * \real{0.5000}}@{}}
\toprule\noalign{}
\begin{minipage}[b]{\linewidth}\raggedright
En la cónica
\end{minipage} & \begin{minipage}[b]{\linewidth}\raggedright
En la EDP
\end{minipage} \\
\midrule\noalign{}
\endhead
\bottomrule\noalign{}
\endlastfoot
rotar los ejes & cambiar de variables \((x,y) \to (\xi,\eta)\) \\
matar el término \(\xi\eta\) & matar el término \(u_{\xi\eta}\) (caso
elíptico) \\
las asíntotas & las características \\
tomar las asíntotas como ejes & tomar las características como
coordenadas (caso hiperbólico) \\
ecuación reducida \(\pm\tilde\xi^2 \pm \tilde\eta^2 = 1\) & forma
canónica \(u_{\xi\xi} \pm u_{\eta\eta} = G\) \\
\end{longtable}
}

Con una diferencia que conviene tener presente: en la EDP el cambio de
variables \textbf{no tiene por qué ser una rotación}. Puede ser
cualquier cambio invertible, incluso uno que dependa del punto, y eso da
más libertad ---por eso las formas canónicas tienen coeficientes
\(\pm 1\) y \(0\) y no los autovalores---. Lo que ningún cambio real
puede hacer es cambiar el \textbf{signo} del determinante, y por eso el
tipo es lo único que no se puede alterar.

\end{minipage}%
\end{tcolorbox}

\emph{(En la versión web hay acá un gráfico interactivo con tres
sliders, \(A\), \(B\) y \(C\), que dibuja a la vez la cónica asociada
con sus direcciones nulas y las familias de características, y muestra
el discriminante y el tipo. Sirve para ver cómo la elipse se abre en
hipérbola al crecer el término cruzado.)}

\begin{exercise}[]\protect\hypertarget{exr-autovalores}{}\label{exr-autovalores}

Para cada una de las ecuaciones \(a\), \(c\) y \(e\) del
Ejercicio~\ref{exr-clasificar}, escribir la matriz de coeficientes
\(M\), calcular sus autovalores y comprobar que los signos coinciden con
lo que predice la Tabla~\ref{tbl-autovalores}. Para el caso \(e\), ¿en
qué puntos del plano se anula algún autovalor?

\end{exercise}

\begin{tcolorbox}[enhanced jigsaw, arc=.35mm, bottomrule=.15mm, bottomtitle=1mm, breakable, colback=white, colbacktitle=quarto-callout-tip-color!10!white, colframe=quarto-callout-tip-color-frame, coltitle=black, left=2mm, leftrule=.75mm, opacityback=0, opacitybacktitle=0.6, rightrule=.15mm, title=\textcolor{quarto-callout-tip-color}{\faLightbulb}\hspace{0.5em}{Solución}, titlerule=0mm, toprule=.15mm, toptitle=1mm]

\textbf{a.} \(u_{xx}-4u_{xy}+3u_{yy}=0\) tiene \(A=1\), \(B=-2\),
\(C=3\), o sea \(M = \begin{pmatrix} 1 & -2\\ -2 & 3\end{pmatrix}\).
Traza \(4\) y determinante \(3-4=-1\), así que
\(\lambda = 2 \pm \sqrt{5}\), aproximadamente \(4{,}24\) y \(-0{,}24\):
\textbf{signos opuestos}. La Tabla~\ref{tbl-autovalores} predice
hiperbólica, y coincide con \(\Delta = B^2-AC = 4-3 = 1 > 0\). ✓ (Y se
ve la relación general: \(\det M = -\Delta\).)

\textbf{c.} \(u_{xx}+2u_{xy}+u_{yy}=u_x\) tiene
\(M = \begin{pmatrix} 1&1\\1&1\end{pmatrix}\), cuyos autovalores son
\(0\) y \(2\). \textbf{Un autovalor nulo}: parabólica, como predice la
tabla, y \(\Delta = 1-1 = 0\). ✓ (Fijate que el lado derecho \(u_x\) no
interviene: es de orden menor.)

\textbf{e.} \(u_{xx}+2\operatorname{sen}(x)\,u_{xy}+u_{yy}=0\) tiene
\(M = \begin{pmatrix} 1 & \operatorname{sen}x\\ \operatorname{sen}x & 1\end{pmatrix}\),
con autovalores \(\lambda = 1 \pm \operatorname{sen}x\). Los dos son
\textbf{mayores o iguales que cero}, así que la ecuación es elíptica
donde los dos son estrictamente positivos.

Algún autovalor se anula cuando \(\operatorname{sen}x = \pm 1\), o sea
en las rectas verticales
\[x = \frac{\pi}{2} + k\pi, \qquad k \in \mathbb{Z},\] donde la ecuación
se vuelve parabólica. Nunca es hiperbólica, porque
\(\Delta = \operatorname{sen}^2x - 1 = -\cos^2 x \le 0\).

\end{tcolorbox}

¿Cuál de estas afirmaciones sobre una ecuación \textbf{elíptica} es
falsa?

\begin{itemize}
\tightlist
\item
  {Su matriz de coeficientes tiene los dos autovalores del mismo signo}
\item
  {Su forma cuadrática asociada no se anula en ninguna dirección real no
  nula}
\item
  {No admite ningún cambio de variables que la lleve a una forma
  canónica}
\item
  {No tiene características reales}
\end{itemize}

Sí admite forma canónica: la elíptica se lleva a
\(u_{\sigma\sigma} + u_{\tau\tau} = G\), como vamos a ver en la
sección~\ref{sec-formas-canonicas}. Lo que pasa es que el cambio de
variables que la diagonaliza es \textbf{complejo} ---las direcciones
nulas son complejas conjugadas---, y hay que tomar parte real e
imaginaria para volver a variables reales. Que no haya características
\emph{reales} no significa que no haya forma canónica.

Las otras tres son equivalentes entre sí y son la definición de elíptica
leída de tres maneras: \(\lambda_1\lambda_2 = \det M = -\Delta > 0\),
\(Q\) definida (positiva o negativa), y \(\Delta < 0\).

\subsection{Por qué el discriminante decide
todo}\label{sec-demostracion-clasificacion}

Con lo anterior, el Teorema~\ref{thm-clasificacion} se cae de maduro. Se
busca un cambio de variables \(\xi = \xi(x,y)\), \(\eta = \eta(x,y)\)
que anule los coeficientes de \(u_{\xi\xi}\) y de \(u_{\eta\eta}\); por
\ref{eq-forma-cuadratica}, eso pide que \((\xi_x,\xi_y)\) y
\((\eta_x,\eta_y)\) sean direcciones nulas de \(Q\), y esas direcciones
son reales y distintas si y solo si \(\Delta > 0\). Las curvas
\(\xi = \text{const}\) y \(\eta = \text{const}\) resultantes son las
características. La sección~\ref{sec-formas-canonicas} hace la cuenta
completa en los tres casos.

\begin{tcolorbox}[enhanced jigsaw, arc=.35mm, bottomrule=.15mm, breakable, colback=white, colframe=quarto-callout-important-color-frame, left=2mm, leftrule=.75mm, opacityback=0, rightrule=.15mm, toprule=.15mm]
\begin{minipage}[t]{5.5mm}
\textcolor{quarto-callout-important-color}{\faExclamation}
\end{minipage}%
\begin{minipage}[t]{\textwidth - 5.5mm}

\vspace{-3mm}\textbf{Por qué esto importa físicamente}\vspace{3mm}

La clasificación no es una taxonomía. Predice el comportamiento:

\begin{itemize}
\tightlist
\item
  \textbf{Hiperbólicas.} Información que viaja a velocidad finita, sin
  suavizarse. Una discontinuidad en el dato inicial sigue siendo una
  discontinuidad para siempre ---viaja por una característica, como
  vimos en la sección~\ref{sec-caracteristicas}---. Ondas.
\item
  \textbf{Parabólicas.} Información que se propaga a velocidad infinita
  y se suaviza al instante. Cualquier dato inicial, por feo que sea, se
  vuelve infinitamente derivable apenas \(t>0\). Difusión, calor.
\item
  \textbf{Elípticas.} No hay tiempo ni propagación, y no hay ninguna
  dirección privilegiada por donde propagar: no hay características
  reales. El valor en un punto depende de \emph{todo} el contorno a la
  vez. Equilibrio, potenciales.
\end{itemize}

Es la misma diferencia que hay entre golpear una cuerda, calentar una
barra y sostener una membrana tensa. Y explica la Tabla~\ref{tbl-datos}:
solo las hiperbólicas y las parabólicas admiten datos de tipo inicial,
porque solo ellas tienen por dónde propagarlos.

\end{minipage}%
\end{tcolorbox}

\begin{example}[Una ecuación que cambia de
tipo]\protect\hypertarget{exm-tricomi}{}\label{exm-tricomi}

La ecuación de Tricomi, \[
u_{xx} + x\,u_{yy} = 0,
\] tiene \(A = 1\), \(B = 0\), \(C = x\), así que \(\Delta = -x\). Es
elíptica para \(x>0\), parabólica sobre el eje \(x=0\) e hiperbólica
para \(x<0\). Aparece en flujo transónico: la región elíptica es el
flujo subsónico y la hiperbólica el supersónico, y la recta que las
separa es exactamente donde el flujo alcanza la velocidad del sonido.

Vale la pena volver ahora sobre el Ejemplo~\ref{exm-transonico}. Ahí, la
ecuación del potencial de un flujo compresible resultaba tener
\(\Delta = M^2-1\), y era \textbf{cuasi-lineal}: el tipo dependía de la
solución. Tricomi es el modelo \textbf{lineal} más simple con el mismo
comportamiento: se cambia la dependencia complicada en \(\nabla\varphi\)
por una dependencia lineal en \(x\), y se conserva lo único que importa,
que es que \(\Delta\) cambie de signo al cruzar una recta. Por eso
Tricomi la estudió: es el laboratorio donde se pueden hacer cuentas.

\end{example}

\begin{exercise}[]\protect\hypertarget{exr-tricomi}{}\label{exr-tricomi}

Para la ecuación de Tricomi del Ejemplo~\ref{exm-tricomi}, hallar
explícitamente las dos familias de características en la región \(x<0\).
(Sugerencia: la ecuación de las características es
\(A(y')^2 - 2By' + C = 0\); resolver y después integrar.)

\end{exercise}

\begin{tcolorbox}[enhanced jigsaw, arc=.35mm, bottomrule=.15mm, bottomtitle=1mm, breakable, colback=white, colbacktitle=quarto-callout-tip-color!10!white, colframe=quarto-callout-tip-color-frame, coltitle=black, left=2mm, leftrule=.75mm, opacityback=0, opacitybacktitle=0.6, rightrule=.15mm, title=\textcolor{quarto-callout-tip-color}{\faLightbulb}\hspace{0.5em}{Solución}, titlerule=0mm, toprule=.15mm, toptitle=1mm]

Para \(u_{xx} + x\,u_{yy}=0\) tenemos \(A=1\), \(B=0\), \(C=x\), y la
ecuación de las características \(A\,(y')^2 - 2B\,y' + C = 0\) queda \[
(y')^2 + x = 0 \qquad\Longrightarrow\qquad \frac{dy}{dx} = \pm\sqrt{-x},
\] que tiene sentido real justamente donde \(x<0\), la región
hiperbólica.

Integrando, con \(w=-x\): \[
y = \pm\int \sqrt{-x}\,dx = \mp\frac{2}{3}\,(-x)^{3/2} + k ,
\] así que las dos familias son \[
\xi = y + \tfrac{2}{3}(-x)^{3/2} = \text{const}, \qquad
\eta = y - \tfrac{2}{3}(-x)^{3/2} = \text{const}.
\]

Son parábolas semicúbicas. Lo interesante es qué hacen al acercarse al
eje \(x=0\): como \(dy/dx = \pm\sqrt{-x} \to 0\), las dos familias
llegan a la recta sónica \textbf{tangentes a ella}, y ahí se juntan de a
pares formando una cúspide. Es la imagen geométrica de que las dos
familias de características se fusionan en una sola justo donde la
ecuación pasa de hiperbólica a parabólica, y desaparecen del lado
elíptico.

\end{tcolorbox}

\subsection*{Tres ejercicios largos}\label{tres-ejercicios-largos}
\addcontentsline{toc}{subsection}{Tres ejercicios largos}

Los cortos están repartidos por la sección. Estos tres son de recorrer.

\begin{exercise}[]\protect\hypertarget{exr-clasificar}{}\label{exr-clasificar}

Clasificar cada ecuación e indicar en qué región del plano vale la
clasificación:

\begin{enumerate}
\def\labelenumi{\alph{enumi}.}
\tightlist
\item
  \(u_{xx} - 4u_{xy} + 3u_{yy} = 0\)
\item
  \(y\,u_{xx} + u_{yy} = 0\)
\item
  \(u_{xx} + 2u_{xy} + u_{yy} = u_x\)
\item
  \(e^{x}u_{xx} + e^{y}u_{yy} = 0\)
\item
  \(u_{xx} + 2\operatorname{sen}(x)\,u_{xy} + u_{yy} = 0\)
\item
  \(x^2u_{xx} - y^2u_{yy} = 0\)
\item
  \(u_{xx} + xy\,u_{yy} = 0\)
\item
  \((1+x^2)u_{xx} + (1+y^2)u_{yy} = 0\)
\end{enumerate}

\end{exercise}

\begin{tcolorbox}[enhanced jigsaw, arc=.35mm, bottomrule=.15mm, bottomtitle=1mm, breakable, colback=white, colbacktitle=quarto-callout-tip-color!10!white, colframe=quarto-callout-tip-color-frame, coltitle=black, left=2mm, leftrule=.75mm, opacityback=0, opacitybacktitle=0.6, rightrule=.15mm, title=\textcolor{quarto-callout-tip-color}{\faLightbulb}\hspace{0.5em}{Solución}, titlerule=0mm, toprule=.15mm, toptitle=1mm]

Recordá que el coeficiente del término cruzado es \(2B\), y que
\(\Delta = B^2-AC\).

\textbf{a.} \(A=1\), \(B=-2\), \(C=3\): \(\Delta = 4-3 = 1 > 0\).
\textbf{Hiperbólica en todo el plano} (los coeficientes son constantes).

\textbf{b.} \(A=y\), \(B=0\), \(C=1\): \(\Delta = -y\).
\textbf{Elíptica} en \(y>0\), \textbf{parabólica} sobre el eje \(y=0\),
\textbf{hiperbólica} en \(y<0\). Es la ecuación de Tricomi con los
papeles de \(x\) e \(y\) cambiados.

\textbf{c.} \(A=1\), \(B=1\), \(C=1\): \(\Delta = 1-1 = 0\).
\textbf{Parabólica en todo el plano}. El lado derecho \(u_x\) no cambia
nada.

\textbf{d.} \(A=e^{x}\), \(B=0\), \(C=e^{y}\): \(\Delta = -e^{x+y}\),
que es negativo siempre. \textbf{Elíptica en todo el plano.} Un buen
recordatorio de que coeficientes variables no implican tipo variable.

\textbf{e.} \(A=1\), \(B=\operatorname{sen}x\), \(C=1\):
\(\Delta = \operatorname{sen}^2x - 1 = -\cos^2x\). \textbf{Elíptica}
salvo en las rectas \(x = \pi/2 + k\pi\), donde es \textbf{parabólica}.
Nunca hiperbólica.

\textbf{f.} \(A=x^2\), \(B=0\), \(C=-y^2\): \(\Delta = x^2y^2 \ge 0\).
\textbf{Hiperbólica} en todo el plano salvo sobre los ejes \(x=0\) e
\(y=0\), donde \(\Delta=0\) y es \textbf{parabólica}. Ojo con el origen:
ahí \(A=B=C=0\) y la ecuación no tiene parte principal, así que la
clasificación directamente no se aplica.

\textbf{g.} \(A=1\), \(B=0\), \(C=xy\): \(\Delta = -xy\).
\textbf{Elíptica} donde \(xy>0\), o sea en el primer y el tercer
cuadrante; \textbf{hiperbólica} donde \(xy<0\), en el segundo y el
cuarto; \textbf{parabólica} sobre los dos ejes.

\textbf{h.} \(A=1+x^2\), \(B=0\), \(C=1+y^2\):
\(\Delta = -(1+x^2)(1+y^2) < 0\) siempre. \textbf{Elíptica en todo el
plano.}

Una observación que vale para todos: el tipo lo decidieron \textbf{solo}
\(A\), \(B\) y \(C\). Los términos con \(u_x\), \(u_y\) y \(u\) no
aparecieron ni una vez.

\end{tcolorbox}

\begin{exercise}[]\protect\hypertarget{exr-adivinar-tipo}{}\label{exr-adivinar-tipo}

Antes de escribir ninguna ecuación, decidí qué tipo esperás en cada caso
y por qué. Después escribí la ecuación y verificalo.

\begin{enumerate}
\def\labelenumi{\alph{enumi}.}
\tightlist
\item
  Una cuerda de guitarra que se pulsa y se suelta.
\item
  Una taza de café que se enfría sobre la mesa.
\item
  El potencial eléctrico en la región entre dos placas cargadas, en
  equilibrio.
\item
  La deformación de una membrana de jabón sostenida por un alambre
  torcido.
\item
  Un pulso de presión que recorre un caño largo lleno de aire.
\item
  Un contaminante que baja por un río, arrastrado por la corriente y a
  la vez difundiéndose.
\end{enumerate}

\end{exercise}

\begin{tcolorbox}[enhanced jigsaw, arc=.35mm, bottomrule=.15mm, bottomtitle=1mm, breakable, colback=white, colbacktitle=quarto-callout-tip-color!10!white, colframe=quarto-callout-tip-color-frame, coltitle=black, left=2mm, leftrule=.75mm, opacityback=0, opacitybacktitle=0.6, rightrule=.15mm, title=\textcolor{quarto-callout-tip-color}{\faLightbulb}\hspace{0.5em}{Solución}, titlerule=0mm, toprule=.15mm, toptitle=1mm]

La pregunta que resuelve casi todos los casos es: \textbf{¿hay algo que
se propague a velocidad finita, algo que se difunda, o nada que
evolucione?}

\textbf{a. Hiperbólica.} Una perturbación en un punto de la cuerda tarda
un tiempo en llegar al otro extremo, y llega sin desdibujarse: es una
onda. \(u_{tt} = c^2u_{xx}\), con \(\Delta = c^2 > 0\).

\textbf{b. Parabólica.} El calor no viaja como una onda: se desparrama y
suaviza. Además el proceso tiene una flecha del tiempo clarísima ---el
café se enfría, nunca se calienta solo---, y eso es la firma de una
parabólica. \(u_t = k\,u_{xx}\), con \(\Delta = 0\).

\textbf{c.~Elíptica.} No hay tiempo en el problema: es equilibrio. El
potencial en cada punto queda determinado por lo que pasa en
\textbf{toda} la frontera a la vez, y si se mueve una placa todo se
reacomoda de una vez. \(\nabla^2 u = 0\), con \(\Delta = -1 < 0\).

\textbf{d.~Elíptica.} Igual que \(c\), aunque no lo parezca: es un
problema de equilibrio, la película adopta la forma de área mínima
compatible con el alambre. Es la ecuación del
Ejercicio~\ref{exr-superficie-minima}, y ahí calculamos
\(\Delta = -(1+u_x^2+u_y^2) < 0\). Ojo con la trampa: que la superficie
\emph{se vea} como una onda congelada no la vuelve hiperbólica.

\textbf{e. Hiperbólica.} Un pulso de presión es sonido, y el sonido
viaja a velocidad finita. Misma ecuación que \(a\) con \(c\) la
velocidad del sonido.

\textbf{f.~Parabólica.} Este es el interesante, porque tiene las dos
cosas: arrastre y difusión. La ecuación es \[u_t + v\,u_x = D\,u_{xx},\]
la \textbf{ecuación de advección--difusión}. Pero el tipo lo decide
\textbf{solo la parte principal}, que acá es \(-D\,u_{xx}\): \(A = -D\),
\(B=0\), \(C=0\) y \(\Delta = 0\), \textbf{parabólica}. El término de
arrastre \(v\,u_x\) es de orden menor y no cambia el tipo, aunque domine
el aspecto de la solución ---la mancha viaja río abajo mientras se
ensancha---.

Es un caso donde la intuición física y la clasificación parecen
pelearse: uno \emph{ve} algo propagándose, que suena a hiperbólico. Pero
lo que decide es si hay difusión, y con \(D>0\), por chico que sea, la
difusión gana: los quiebres se borran al instante y el problema hacia
atrás está mal planteado. Solo en el límite \(D=0\) la ecuación se
vuelve de primer orden y recupera el carácter de transporte puro.

\end{tcolorbox}

\begin{exercise}[]\protect\hypertarget{exr-cambio-tipo}{}\label{exr-cambio-tipo}

Inventar una ecuación de segundo orden con coeficientes que dependan de
\(x\) e \(y\) que sea hiperbólica dentro del disco unitario y elíptica
afuera. Verificar la respuesta calculando el discriminante.

\end{exercise}

\begin{tcolorbox}[enhanced jigsaw, arc=.35mm, bottomrule=.15mm, bottomtitle=1mm, breakable, colback=white, colbacktitle=quarto-callout-tip-color!10!white, colframe=quarto-callout-tip-color-frame, coltitle=black, left=2mm, leftrule=.75mm, opacityback=0, opacitybacktitle=0.6, rightrule=.15mm, title=\textcolor{quarto-callout-tip-color}{\faLightbulb}\hspace{0.5em}{Solución}, titlerule=0mm, toprule=.15mm, toptitle=1mm]

Hay muchas respuestas. La más simple sale de pedir directamente que el
discriminante sea \(1-x^2-y^2\): eso es positivo adentro del disco y
negativo afuera. Con \(B=0\) y \(C=1\) alcanza con tomar
\(A = -(1-x^2-y^2)\), o sea \[
\big(x^2+y^2-1\big)\,u_{xx} + u_{yy} = 0 .
\] Verificación: \(A = x^2+y^2-1\), \(B=0\), \(C=1\), y
\[\Delta = B^2 - AC = -\big(x^2+y^2-1\big) = 1-x^2-y^2,\] que es \(>0\)
adentro del disco unitario (\textbf{hiperbólica}), \(=0\) sobre la
circunferencia (\textbf{parabólica}) y \(<0\) afuera
(\textbf{elíptica}). ✓

Cualquier variante sirve: por ejemplo
\(u_{xx} + (x^2+y^2-1)\,u_{yy} = 0\) da el mismo discriminante. Y notá
que estas ecuaciones son \textbf{lineales} y sin embargo cambian de
tipo: es la misma situación que Tricomi, y es una buena prueba de que la
clasificación por linealidad y la clasificación por tipo son
independientes.

\end{tcolorbox}

\section{Formas canónicas}\label{sec-formas-canonicas}

Clasificar una ecuación sirve para elegir el método. Pero la
clasificación viene con un regalo extra: \textbf{para cada tipo existe
un sistema de coordenadas en el que la ecuación queda lo más simple que
puede quedar}, y ese sistema tiene un significado físico claro.

Empecemos por el caso donde eso se ve sin ninguna cuenta.

\subsection{La idea, en la cuerda}\label{sec-idea-canonica}

Una cuerda tensa satisface \(u_{tt} = c^2u_{xx}\). Ya sabemos qué hace:
las perturbaciones viajan a velocidad \(c\) para los dos lados.
Entonces, ¿qué pasa si en vez de mirarla desde el laboratorio nos
subimos a la onda?

Un observador que viaja hacia la derecha a velocidad \(c\) tiene, en
cada instante, la coordenada \(x - ct\) fija. Uno que viaja hacia la
izquierda tiene fija \(x + ct\). Tomemos esas dos cantidades como
coordenadas nuevas:

\begin{equation}\protect\phantomsection\label{eq-coord-onda}{
\xi = x + ct, \qquad \eta = x - ct .
}\end{equation}

Y hagamos el cambio. Por la regla de la cadena, \(u_x = u_\xi + u_\eta\)
y \(u_t = c\,u_\xi - c\,u_\eta\), así que

\[
u_{xx} = u_{\xi\xi} + 2u_{\xi\eta} + u_{\eta\eta}, \qquad
u_{tt} = c^2\big(u_{\xi\xi} - 2u_{\xi\eta} + u_{\eta\eta}\big),
\]

y restando,

\begin{equation}\protect\phantomsection\label{eq-onda-canonica}{
u_{tt} - c^2u_{xx} = -4c^2\,u_{\xi\eta} .
}\end{equation}

Toda la ecuación de onda colapsó a

\[
u_{\xi\eta} = 0 .
\]

Que es la \ref{eq-trivial2}, la que ya resolvimos en la
sección~\ref{sec-porque}: su solución general es
\(u = f(\xi) + g(\eta)\), o sea

\begin{equation}\protect\phantomsection\label{eq-dalembert}{
u(x,t) = f(x+ct) + g(x-ct) .
}\end{equation}

Esa es la \textbf{solución de d'Alembert}: toda solución de la ecuación
de onda es la suma de un perfil que viaja hacia la derecha sin
deformarse y otro que viaja hacia la izquierda. Las dos funciones
arbitrarias que la regla empírica de la sección~\ref{sec-porque}
anticipaba para una ecuación de segundo orden son exactamente esos dos
perfiles.

\begin{tcolorbox}[enhanced jigsaw, arc=.35mm, bottomrule=.15mm, breakable, colback=white, colframe=quarto-callout-important-color-frame, left=2mm, leftrule=.75mm, opacityback=0, rightrule=.15mm, toprule=.15mm]
\begin{minipage}[t]{5.5mm}
\textcolor{quarto-callout-important-color}{\faExclamation}
\end{minipage}%
\begin{minipage}[t]{\textwidth - 5.5mm}

\vspace{-3mm}\textbf{Qué acaba de pasar}\vspace{3mm}

No usamos ninguna técnica nueva. \textbf{Cambiamos a las coordenadas que
viajan con la onda}, y en esas coordenadas la ecuación no tenía nada que
decir salvo «hay una parte que va para allá y otra que va para acá».

Y las coordenadas que elegimos no salieron de la galera:
\(x \pm ct = \text{const}\) son las \textbf{características} de la
ecuación de onda. Es el mismo objeto de la
sección~\ref{sec-caracteristicas}, ahora usado de otra manera: en primer
orden seguíamos \emph{una} característica para volver ordinaria la
ecuación; en segundo orden usamos \emph{las dos familias} como ejes
coordenados.

Todo el resto de esta sección es hacer eso mismo para una ecuación
cualquiera.

\end{minipage}%
\end{tcolorbox}

\begin{tcolorbox}[enhanced jigsaw, arc=.35mm, bottomrule=.15mm, breakable, colback=white, colframe=quarto-callout-warning-color-frame, left=2mm, leftrule=.75mm, opacityback=0, rightrule=.15mm, toprule=.15mm]
\begin{minipage}[t]{5.5mm}
\textcolor{quarto-callout-warning-color}{\faExclamationTriangle}
\end{minipage}%
\begin{minipage}[t]{\textwidth - 5.5mm}

\vspace{-3mm}\textbf{Hacé la cuenta del factor una vez}\vspace{3mm}

El \(-4c^2\) de \ref{eq-onda-canonica} es fácil de perder por el camino.
Conviene derivarlo despacio una vez ---son cuatro renglones--- y
quedarse con el propio resultado: si al aplicar la fórmula a un problema
con fuente, \(u_{tt}-c^2u_{xx} = s(x,t)\), te da
\(u_{\xi\eta} = -s/(4c^2)\), la cuenta está bien. Distintos textos
escriben la constante distinto según cómo definan \(\xi\) y \(\eta\), y
algún factor \(2\) se pierde con facilidad.

\end{minipage}%
\end{tcolorbox}

En \(u(x,t) = f(x+ct) + g(x-ct)\), ¿qué hace cada término?

\begin{itemize}
\tightlist
\item
  {\(f\) viaja hacia la izquierda y \(g\) hacia la derecha}
\item
  {\(f\) viaja hacia la derecha y \(g\) hacia la izquierda}
\item
  {\(f\) es la posición inicial y \(g\) la velocidad inicial}
\item
  {\(f\) es la solución y \(g\) una corrección de orden menor}
\end{itemize}

El perfil \(f\) se mantiene constante donde \(x+ct\) es constante, o sea
sobre \(x = \text{const} - ct\): a medida que \(t\) crece, ese punto se
va hacia \textbf{la izquierda}. Con \(g\) pasa lo contrario.

El signo es la trampa clásica y conviene fijarlo con un caso concreto:
si \(g\) es una joroba centrada en el origen, entonces \(g(x-ct)\) es
esa joroba centrada en \(x=ct\), que se corre hacia la derecha. Es la
\ref{eq-sol-transporte} de la ecuación de transporte, y la ecuación de
onda es dos transportes al mismo tiempo, uno para cada lado.

La tercera opción confunde las dos funciones arbitrarias con las dos
condiciones iniciales. Están relacionadas ---\(\varphi\) y \(\psi\)
determinan \(f\) y \(g\), es el Ejercicio~\ref{exr-dalembert}--- pero no
son lo mismo.

\subsection{El cambio de variables, en
general}\label{sec-cambio-general}

Tomemos ahora la ecuación semilineal general \ref{eq-general} y un
cambio cualquiera

\[
\xi = \xi(x,y), \qquad \eta = \eta(x,y),
\]

invertible, es decir con jacobiano \(\xi_x\eta_y - \xi_y\eta_x \neq 0\).
Aplicando la regla de la cadena dos veces:

\begin{equation}\protect\phantomsection\label{eq-cadena-segundo}{
\begin{aligned}
u_x &= \xi_x u_\xi + \eta_x u_\eta, \\
u_y &= \xi_y u_\xi + \eta_y u_\eta, \\
u_{xx} &= \xi_x^2 u_{\xi\xi} + 2\xi_x\eta_x u_{\xi\eta} + \eta_x^2 u_{\eta\eta}
        + \xi_{xx} u_\xi + \eta_{xx} u_\eta, \\
u_{xy} &= \xi_x\xi_y u_{\xi\xi} + (\xi_x\eta_y + \xi_y\eta_x) u_{\xi\eta} + \eta_x\eta_y u_{\eta\eta}
        + \xi_{xy} u_\xi + \eta_{xy} u_\eta, \\
u_{yy} &= \xi_y^2 u_{\xi\xi} + 2\xi_y\eta_y u_{\xi\eta} + \eta_y^2 u_{\eta\eta}
        + \xi_{yy} u_\xi + \eta_{yy} u_\eta .
\end{aligned}
}\end{equation}

Reemplazando en \ref{eq-general} y agrupando, la ecuación queda

\begin{equation}\protect\phantomsection\label{eq-transformada}{
\alpha\,u_{\xi\xi} + \beta\,u_{\xi\eta} + \gamma\,u_{\eta\eta} = H(\xi,\eta,u,u_\xi,u_\eta),
}\end{equation}

donde todo lo de orden menor ---incluidos los términos con \(\xi_{xx}\),
\(\eta_{xy}\) y compañía--- se fue adentro de \(H\), y los tres
coeficientes de arriba son

\begin{equation}\protect\phantomsection\label{eq-alfa-beta-gamma}{
\begin{aligned}
\alpha &= A\,\xi_x^2 + 2B\,\xi_x\xi_y + C\,\xi_y^2, \\
\beta  &= 2\big[A\,\xi_x\eta_x + B(\xi_x\eta_y + \xi_y\eta_x) + C\,\xi_y\eta_y\big], \\
\gamma &= A\,\eta_x^2 + 2B\,\eta_x\eta_y + C\,\eta_y^2 .
\end{aligned}
}\end{equation}

\begin{tcolorbox}[enhanced jigsaw, arc=.35mm, bottomrule=.15mm, breakable, colback=white, colframe=quarto-callout-tip-color-frame, left=2mm, leftrule=.75mm, opacityback=0, rightrule=.15mm, toprule=.15mm]
\begin{minipage}[t]{5.5mm}
\textcolor{quarto-callout-tip-color}{\faLightbulb}
\end{minipage}%
\begin{minipage}[t]{\textwidth - 5.5mm}

\vspace{-3mm}\textbf{Esto ya lo teníamos}\vspace{3mm}

Mirá \(\alpha\) y \(\gamma\): son \textbf{la forma cuadrática \(Q\) del
Definición~\ref{def-forma-cuadratica-asociada}}, evaluada en
\((\xi_x,\xi_y)\) y en \((\eta_x,\eta_y)\) respectivamente.

\[\alpha = Q(\xi_x,\xi_y), \qquad \gamma = Q(\eta_x,\eta_y).\]

Así que pedir «que se anule el coeficiente de \(u_{\xi\xi}\)» es pedir
que el gradiente de \(\xi\) apunte en una \textbf{dirección nula} de
\(Q\) ---una asíntota de la cónica---, que es exactamente lo que
estudiamos en la sección~\ref{sec-conicas}. Todo lo que sigue ya está
decidido por el signo de \(\Delta\): si hay dos direcciones nulas reales
podemos anular \(\alpha\) y \(\gamma\) a la vez, si hay una sola podemos
anular una, y si no hay ninguna no podemos anular ninguna de las dos.

Se puede probar además ---es una cuenta--- que
\[\beta^2 - 4\alpha\gamma = 4\,\big(B^2-AC\big)\,\big(\xi_x\eta_y - \xi_y\eta_x\big)^2 ,\]
o sea que el discriminante nuevo tiene \textbf{el mismo signo} que el
viejo, porque el jacobiano al cuadrado es positivo. Ningún cambio de
variables real puede cambiar el tipo, y por eso la clasificación es una
propiedad de la ecuación y no de las coordenadas.

\end{minipage}%
\end{tcolorbox}

Las tres formas canónicas a las que apuntamos son

\[
\underbrace{u_{\xi\eta} = G}_{\text{hiperbólica}} \qquad
\underbrace{u_{\xi\xi} = G}_{\text{parabólica}} \qquad
\underbrace{u_{\sigma\sigma} + u_{\tau\tau} = G}_{\text{elíptica}}
\]

con \(G\) conteniendo solo términos de orden menor. Vamos una por una.

Una ecuación es elíptica en un punto. ¿Puede algún cambio de variables
real e invertible volverla hiperbólica ahí?

\begin{itemize}
\tightlist
\item
  {Sí, si el cambio depende del punto}
\item
  {Sí, si el jacobiano es negativo}
\item
  {No, nunca}
\item
  {Solo si los coeficientes son constantes}
\end{itemize}

Nunca. La relación
\[\beta^2 - 4\alpha\gamma = 4\,(B^2-AC)\,(\xi_x\eta_y - \xi_y\eta_x)^2\]
dice que el discriminante nuevo es el viejo multiplicado por el
\textbf{cuadrado} del jacobiano. Un cuadrado es positivo, y el jacobiano
no puede ser cero porque el cambio es invertible: el signo se conserva
pase lo que pase. Que el jacobiano sea negativo tampoco importa,
justamente porque está elevado al cuadrado.

Lo que sí puede hacer un cambio de variables es simplificar la ecuación
todo lo posible dentro de su tipo. Eso es la forma canónica, y por eso
hay \textbf{tres} y no una: el tipo es lo único que no se puede
negociar.

\subsection{Hiperbólicas}\label{sec-canonica-hiperbolica}

Con \(\Delta = B^2-AC > 0\) hay dos direcciones nulas reales, así que
podemos pedir \(\alpha = \gamma = 0\). Eso significa que tanto \(\xi\)
como \(\eta\) deben satisfacer

\[
A\,\phi_x^2 + 2B\,\phi_x\phi_y + C\,\phi_y^2 = 0 ,
\]

y despejando \(\phi_x/\phi_y\) con la resolvente,

\begin{equation}\protect\phantomsection\label{eq-dos-nulas}{
\frac{\xi_x}{\xi_y} = \frac{-B+\sqrt{B^2-AC}}{A}, \qquad
\frac{\eta_x}{\eta_y} = \frac{-B-\sqrt{B^2-AC}}{A} .
}\end{equation}

Cada una de esas es una ecuación en derivadas parciales de
\textbf{primer orden} para \(\xi\), y ya sabemos resolverla: es la
sección~\ref{sec-caracteristicas}. Sus características son las curvas
\(y(x)\) a lo largo de las cuales \(\xi\) es constante, y cumplen

\begin{equation}\protect\phantomsection\label{eq-edo-caracteristicas}{
\boxed{\;A\left(\frac{dy}{dx}\right)^2 - 2B\,\frac{dy}{dx} + C = 0\;}
}\end{equation}

\begin{definition}[Curvas características de una ecuación de segundo
orden]\protect\hypertarget{def-curvas-caracteristicas-2}{}\label{def-curvas-caracteristicas-2}

Las soluciones de la ecuación diferencial ordinaria
\ref{eq-edo-caracteristicas} se llaman \textbf{curvas características}
de \ref{eq-general}. Cuando \(\Delta>0\) hay dos familias reales; si
\(F_1(x,y) = \text{const}\) y \(F_2(x,y) = \text{const}\) las describen,
el cambio de variables \(\xi = F_1\), \(\eta = F_2\) lleva la ecuación a
la forma canónica hiperbólica

\[
u_{\xi\eta} = G(\xi,\eta,u,u_\xi,u_\eta).
\]

\end{definition}

Hay una segunda forma canónica hiperbólica, la que se usa cuando uno
quiere que se parezca a la ecuación de onda. Con

\[
\sigma = \xi + \eta, \qquad \tau = \xi - \eta
\]

se obtiene

\begin{equation}\protect\phantomsection\label{eq-canonica-hiperbolica-2}{
u_{\sigma\sigma} - u_{\tau\tau} = K(\sigma,\tau,u,u_\sigma,u_\tau).
}\end{equation}

Y acá hay una vuelta linda: si aplicamos esto a la ecuación de onda, con
\(\xi = x+ct\) y \(\eta = x-ct\), sale \(\sigma = 2x\) y \(\tau = 2ct\).
\textbf{Volvimos a las coordenadas originales.} La ecuación de onda ya
estaba escrita en su segunda forma canónica desde el principio, y eso es
otra manera de decir que es el modelo de todas las hiperbólicas.

\begin{figure}

\centering{

\pandocbounded{\includegraphics[keepaspectratio]{08-edp_files/figure-pdf/fig-coordenadas-output-1.pdf}}

}

\caption{\label{fig-coordenadas}La red de coordenadas características. A
la izquierda, la ecuación de onda: las dos familias son las rectas
\(x \pm ct = \text{const}\), y son rectas porque los coeficientes son
constantes. A la derecha, la ecuación \(y^2u_{xx}-x^2u_{yy}=0\): las
familias son las hipérbolas \(y^2-x^2=\text{const}\) y las
circunferencias \(y^2+x^2=\text{const}\). En los dos casos, tomar esas
curvas como ejes coordenados es lo que lleva la ecuación a su forma
canónica.}

\end{figure}%

\begin{example}[La cuerda con una fuerza
aplicada]\protect\hypertarget{exm-onda-fuente}{}\label{exm-onda-fuente}

Sea una cuerda forzada, \(\;u_{tt} - c^2u_{xx} = s(x,t)\), donde \(s\)
es la fuerza externa por unidad de masa.

Es hiperbólica: con \(A = -c^2\) (el coeficiente de \(u_{xx}\)), \(B=0\)
y \(C=1\), queda \(\Delta = c^2 > 0\). La \ref{eq-edo-caracteristicas}
da
\[-c^2\left(\frac{dt}{dx}\right)^2 + 1 = 0 \ \Longrightarrow\ \frac{dt}{dx} = \pm\frac{1}{c}
\ \Longrightarrow\ x \mp ct = \text{const},\] o sea las dos familias de
siempre, y con \(\xi = x+ct\), \(\eta = x-ct\) la cuenta de
\ref{eq-onda-canonica} da \[
u_{\xi\eta} = -\frac{s}{4c^2}, \qquad\text{con } x = \frac{\xi+\eta}{2}, \quad t = \frac{\xi-\eta}{2c} .
\] Con \(s\equiv 0\) recuperamos d'Alembert. Con \(s\) cualquiera, la
ecuación se integra dos veces ---una en \(\xi\) y otra en \(\eta\)--- y
de ahí sale la fórmula de d'Alembert con fuente, que vamos a usar en el
capítulo 9.

\end{example}

\begin{exercise}[]\protect\hypertarget{exr-dalembert}{}\label{exr-dalembert}

Usando la forma canónica del Ejemplo~\ref{exm-onda-fuente} con
\(s\equiv 0\), deducir la \textbf{fórmula de d'Alembert}: la solución de
\[u_{tt} = c^2u_{xx}, \qquad u(x,0) = \varphi(x), \qquad u_t(x,0) = \psi(x)\]
en toda la recta. \emph{(Escribí \(u = f(x+ct)+g(x-ct)\), imponé las dos
condiciones y despejá \(f\) y \(g\).)}

\end{exercise}

\begin{tcolorbox}[enhanced jigsaw, arc=.35mm, bottomrule=.15mm, bottomtitle=1mm, breakable, colback=white, colbacktitle=quarto-callout-tip-color!10!white, colframe=quarto-callout-tip-color-frame, coltitle=black, left=2mm, leftrule=.75mm, opacityback=0, opacitybacktitle=0.6, rightrule=.15mm, title=\textcolor{quarto-callout-tip-color}{\faLightbulb}\hspace{0.5em}{Solución}, titlerule=0mm, toprule=.15mm, toptitle=1mm]

De la forma canónica, \(u = f(x+ct) + g(x-ct)\). Imponiendo las
condiciones en \(t=0\): \[
f(x) + g(x) = \varphi(x), \qquad
c\,f'(x) - c\,g'(x) = \psi(x) .
\] De la segunda, \(f'-g' = \psi/c\), e integrando desde un punto
cualquiera \(x_1\):
\[f(x) - g(x) = \frac{1}{c}\int_{x_1}^{x}\psi(s)\,ds + K .\] Sumando y
restando con la primera, \[
f(x) = \frac{\varphi(x)}{2} + \frac{1}{2c}\int_{x_1}^{x}\psi + \frac{K}{2}, \qquad
g(x) = \frac{\varphi(x)}{2} - \frac{1}{2c}\int_{x_1}^{x}\psi - \frac{K}{2} .
\] Y evaluando en \(x+ct\) y \(x-ct\), la constante \(K\) se cancela y
las dos integrales se juntan en una sola: \[
\boxed{\;u(x,t) = \frac{\varphi(x+ct) + \varphi(x-ct)}{2}
\;+\; \frac{1}{2c}\int_{x-ct}^{x+ct}\psi(s)\,ds .\;}
\]

\textbf{Cómo leerla.} El primer término dice que la forma inicial se
parte en dos mitades que se van para cada lado ---lo que muestra el
primer panel de la Figura~\ref{fig-tres-mundos}---. El segundo dice que
la velocidad inicial no viaja: se \textbf{acumula} sobre el intervalo
que quedó entre las dos mitades. Por eso una cuerda pulsada (posición
inicial, velocidad nula) y una cuerda golpeada (posición nula, velocidad
inicial) suenan distinto, y por eso una guitarra y un piano no se
parecen.

La vamos a rehacer con más cuidado en el capítulo 9.

\end{tcolorbox}

\begin{exercise}[]\protect\hypertarget{exr-dominio-dependencia}{}\label{exr-dominio-dependencia}

De la fórmula del Ejercicio~\ref{exr-dalembert}: el valor de \(u\) en un
punto \((x_0,t_0)\), ¿de qué parte del dato inicial depende? Dibujarlo
en el plano \((x,t)\) y decir qué tienen que ver esas dos rectas con las
características.

Después, al revés: si a \(t=0\) alguien golpea la cuerda solo en el
intervalo \([-a,a]\), ¿en qué región del plano \((x,t)\) se siente el
golpe? Si \(c = 340\ \text{m/s}\) y el golpe fue en un metro de cuerda,
¿cuánto tarda en sentirse a diez metros?

\end{exercise}

\begin{tcolorbox}[enhanced jigsaw, arc=.35mm, bottomrule=.15mm, bottomtitle=1mm, breakable, colback=white, colbacktitle=quarto-callout-tip-color!10!white, colframe=quarto-callout-tip-color-frame, coltitle=black, left=2mm, leftrule=.75mm, opacityback=0, opacitybacktitle=0.6, rightrule=.15mm, title=\textcolor{quarto-callout-tip-color}{\faLightbulb}\hspace{0.5em}{Solución}, titlerule=0mm, toprule=.15mm, toptitle=1mm]

\textbf{Hacia atrás: el dominio de dependencia.} En la fórmula aparecen
\(\varphi\) evaluada en \(x_0\pm ct_0\) y \(\psi\) integrada entre esos
dos puntos. Así que \(u(x_0,t_0)\) depende \textbf{exclusivamente} del
dato inicial en el intervalo \[[\,x_0 - ct_0,\; x_0 + ct_0\,],\] que se
llama \textbf{dominio de dependencia} del punto. En el plano \((x,t)\)
es la base del triángulo que se obtiene bajando desde \((x_0,t_0)\) las
dos rectas de pendiente \(\pm 1/c\) ---o sea las dos
\textbf{características que pasan por el punto}---. Nada de lo que pase
fuera de ese intervalo puede afectar a \(u(x_0,t_0)\).

\textbf{Hacia adelante: el rango de influencia.} Si el dato es no nulo
solo en \([-a,a]\), la solución puede ser no nula solo donde el dominio
de dependencia toque ese intervalo, es decir en \[|x| \le a + c\,t ,\]
un cono que se abre desde el segmento inicial con pendiente \(\pm 1/c\).
Afuera de ese cono la cuerda está quieta: \textbf{la perturbación
todavía no llegó}.

\textbf{Los números.} Con \(c=340\) m/s y el golpe en
\([-0{,}5,\,0{,}5]\), el frente llega a \(x=10\) m cuando
\(10 = 0{,}5 + 340\,t\), o sea
\[t = \frac{9{,}5}{340} \approx 0{,}028\ \text{s} \approx 28\ \text{ms}.\]

\textbf{Y el contraste que importa.} Esto es lo que hace hiperbólica a
la ecuación de onda: hay un frente nítido y afuera no pasa nada. En la
ecuación del calor, en cambio, el dominio de dependencia de cualquier
punto es \textbf{toda la recta}: apenas \(t>0\), la temperatura en un
punto depende de la temperatura inicial en todos lados, por lejos que
estén. Es la sección~\ref{sec-tres-mundos} otra vez, ahora con una
fórmula que lo demuestra.

\end{tcolorbox}

\begin{example}[Una con coeficientes
variables]\protect\hypertarget{exm-canonica-variable}{}\label{exm-canonica-variable}

Llevemos \(\;y^2u_{xx} - x^2u_{yy} = 0\;\) a su forma canónica, en la
región \(x>0\), \(y>0\).

\textbf{Tipo.} \(A=y^2\), \(B=0\), \(C=-x^2\), y
\(\Delta = x^2y^2 > 0\): hiperbólica.

\textbf{Características.} \ref{eq-edo-caracteristicas} queda
\[y^2\left(\frac{dy}{dx}\right)^2 - x^2 = 0 \ \Longrightarrow\ \frac{dy}{dx} = \pm\frac{x}{y}
\ \Longrightarrow\ y\,dy = \pm x\,dx ,\] e integrando,
\(y^2 \mp x^2 = \text{const}\). Las dos familias son hipérbolas y
circunferencias, como en la Figura~\ref{fig-coordenadas}. Tomamos
\[\xi = y^2 - x^2, \qquad \eta = y^2 + x^2 .\]

\textbf{El cambio.} \(\xi_x = -2x\), \(\xi_y = 2y\), \(\eta_x = 2x\),
\(\eta_y = 2y\), y las segundas derivadas son \(\xi_{xx}=-2\),
\(\xi_{yy}=2\), \(\eta_{xx}=\eta_{yy}=2\), todas las cruzadas nulas.
Metiendo en \ref{eq-alfa-beta-gamma}: \[
\alpha = y^2(4x^2) - x^2(4y^2) = 0 \ \checkmark, \qquad
\gamma = y^2(4x^2) - x^2(4y^2) = 0 \ \checkmark,
\] \[
\beta = 2\big[y^2(-2x)(2x) - x^2(2y)(2y)\big] = -16x^2y^2 .
\] Los términos de orden uno tienen coeficientes
\(A\xi_{xx} + C\xi_{yy} = -2y^2-2x^2 = -2\eta\) y
\(A\eta_{xx}+C\eta_{yy} = 2y^2-2x^2 = 2\xi\). Y como
\(y^2 = \tfrac{\eta+\xi}{2}\) y \(x^2 = \tfrac{\eta-\xi}{2}\), resulta
\(\beta = -4(\eta^2-\xi^2)\). La ecuación transformada es \[
-4(\eta^2-\xi^2)\,u_{\xi\eta} - 2\eta\,u_\xi + 2\xi\,u_\eta = 0 ,
\] y dividiendo, \[
u_{\xi\eta} = \frac{\xi\,u_\eta - \eta\,u_\xi}{2\,(\eta^2-\xi^2)} .
\] Que ya no se puede simplificar más: los términos de la derecha son de
orden uno y no hay ningún cambio de variables que los mate. Eso es lo
que significa «forma canónica»: no que la ecuación se resuelva, sino que
la parte principal quedó lo más limpia posible.

Según cómo se ordenen y se les cambie el signo a \(\xi\) y \(\eta\), la
expresión final se ve distinta en distintos libros. Es la misma
ecuación.

\end{example}

\begin{example}[Dos hiperbólicas que salen del
laboratorio]\protect\hypertarget{exm-fisica-hiperbolicas}{}\label{exm-fisica-hiperbolicas}

\textbf{La cuerda tensa.} Una cuerda de densidad lineal \(\rho\)
constante, tensada con una fuerza \(T_0\) y sin fuerzas externas, tiene
un desplazamiento transversal \(u(x,t)\) que satisface \[
u_{tt} = a^2 u_{xx}, \qquad a^2 = \frac{T_0}{\rho} .
\] La velocidad de propagación sale sola del cociente entre la fuerza
que endereza la cuerda y la inercia que se le opone: \textbf{más tensa,
más rápido; más pesada, más lento}. Es lo que uno hace al afinar una
guitarra.

\textbf{La barra que se estira.} Si en vez de mover la cuerda de costado
empujamos una barra a lo largo de su eje, el desplazamiento longitudinal
\(u(x,t)\) del punto de abscisa \(x\) cumple, con la ley de Hooke
\(T = k(x)u_x\) ---siendo \(k\) el módulo de Young---, \[
u_{tt} = a^2u_{xx} + f(x,t), \qquad f = \frac{F(x,t)}{\rho} ,
\] con \(F\) la densidad de fuerza externa. La misma ecuación, otra
física: acá lo que viaja no es una ondulación sino una compresión, y es
literalmente el sonido dentro del metal.

\end{example}

\subsection{Parabólicas}\label{sec-canonica-parabolica}

Con \(\Delta = 0\) hay una sola dirección nula, así que no podemos
anular \(\alpha\) y \(\gamma\) a la vez. Lo que se hace es apuntar a la
forma

\[
u_{\xi\xi} = G(\xi,\eta,u,u_\xi,u_\eta),
\]

es decir anular \(\beta\) y \(\gamma\). Para \(\gamma = 0\) hay que
tomar \(\eta\) en la única dirección nula que existe: suponiendo
\(A \neq 0\), la resolvente da la raíz doble

\[
\frac{\eta_x}{\eta_y} = -\frac{B}{A} ,
\]

o sea \(\eta\) constante sobre las curvas
\(\dfrac{dy}{dx} = \dfrac{B}{A}\), la \textbf{única} familia de
características.

Y acá pasa algo notable: \textbf{\(\beta\) se anula solo}. Reemplazando
\(\eta_x = -\tfrac{B}{A}\eta_y\) en la expresión de \(\beta\) y usando
\(B^2 = AC\),

\[
\beta = 2\left[-B\,\xi_x\eta_y + B\Big(\xi_x\eta_y - \tfrac{B}{A}\xi_y\eta_y\Big) + C\,\xi_y\eta_y\right]
= 2\,\eta_y\,\xi_y\,\frac{AC-B^2}{A} = 0 .
\]

Así que \(\xi\) \textbf{se puede elegir libremente}, con la única
condición de que sea funcionalmente independiente de \(\eta\) ---para
que el cambio sea invertible---. Lo habitual es tomar \(\xi = x\) o
\(\xi = y\), lo que resulte más cómodo.

Si \(A = 0\) pero \(C \neq 0\) el papel de \(x\) e \(y\) se intercambia
y la condición pasa a ser \(\eta_x + \tfrac{B}{C}\eta_y = 0\); si
\(A = C = 0\) entonces \(B^2 = AC = 0\) y no hay parte principal.

\begin{tcolorbox}[enhanced jigsaw, arc=.35mm, bottomrule=.15mm, breakable, colback=white, colframe=quarto-callout-tip-color-frame, left=2mm, leftrule=.75mm, opacityback=0, rightrule=.15mm, toprule=.15mm]
\begin{minipage}[t]{5.5mm}
\textcolor{quarto-callout-tip-color}{\faLightbulb}
\end{minipage}%
\begin{minipage}[t]{\textwidth - 5.5mm}

\vspace{-3mm}\textbf{Una sola familia, una sola dirección de propagación}\vspace{3mm}

Que haya una sola familia de características es la razón geométrica de
que las parabólicas se comporten como se comportan. En una hiperbólica
hay dos direcciones privilegiadas y la información se reparte entre las
dos, viajando a velocidad finita. En una parabólica hay una sola, y
resulta ser \(t = \text{const}\): las características de la ecuación del
calor son las \textbf{líneas de tiempo constante}.

Eso quiere decir que «propagarse a lo largo de una característica» es,
para el calor, extenderse por todo el espacio en un mismo instante. De
ahí la velocidad infinita de propagación de la
sección~\ref{sec-tres-mundos}, y de ahí que el problema hacia atrás esté
mal planteado.

\end{minipage}%
\end{tcolorbox}

\begin{example}[Una parabólica con coeficientes
variables]\protect\hypertarget{exm-canonica-parabolica}{}\label{exm-canonica-parabolica}

Llevemos \(\;x^2u_{xx} + 2xy\,u_{xy} + y^2u_{yy} = 0\;\) a su forma
canónica, con \(x>0\).

\textbf{Tipo.} \(A=x^2\), \(B=xy\), \(C=y^2\), y
\(\Delta = x^2y^2 - x^2y^2 = 0\): parabólica en todo el plano.

\textbf{La única familia.}
\(\dfrac{dy}{dx} = \dfrac{B}{A} = \dfrac{xy}{x^2} = \dfrac{y}{x}\), que
da \(y = kx\): \textbf{los rayos que salen del origen}. Una función
constante sobre cada rayo es \(\eta = y/x\). Elegimos \(\xi = x\), que
es independiente.

\textbf{La cuenta.} \(\eta_x = -y/x^2\), \(\eta_y = 1/x\),
\(\xi_x = 1\), \(\xi_y = 0\). Con \ref{eq-alfa-beta-gamma}: \[
\alpha = x^2, \qquad
\beta = 2\Big[x^2\big(-\tfrac{y}{x^2}\big) + xy\,\tfrac1x\Big] = 2(-y+y) = 0 \ \checkmark,
\] \[
\gamma = x^2\frac{y^2}{x^4} + 2xy\Big(-\frac{y}{x^2}\Big)\frac1x + y^2\frac1{x^2}
= \frac{y^2}{x^2} - \frac{2y^2}{x^2} + \frac{y^2}{x^2} = 0 \ \checkmark .
\] Los coeficientes de orden uno también dan cero ---conviene
verificarlo---, así que la ecuación transformada es simplemente \[
x^2\,u_{\xi\xi} = 0 \qquad\Longrightarrow\qquad u_{\xi\xi} = 0 .
\]

\textbf{Y ahora se resuelve sola.} Integrando dos veces en \(\xi\):
\(u = \xi\,f(\eta) + g(\eta)\), o sea \[
u(x,y) = x\,f\!\left(\frac{y}{x}\right) + g\!\left(\frac{y}{x}\right) .
\] La lectura geométrica es clara: sobre cada rayo que sale del origen,
\(u\) es una función \textbf{lineal} de la distancia. Las dos funciones
arbitrarias eligen la pendiente y la ordenada al origen de cada rayo.

\end{example}

¿Cuáles son las curvas características de la ecuación del calor
\(u_t = k\,u_{xx}\)?

\begin{itemize}
\tightlist
\item
  {Las rectas \(x = \text{const}\)}
\item
  {Las rectas \(t = \text{const}\)}
\item
  {No tiene características}
\item
  {Las rectas \(x \pm \sqrt{k}\,t = \text{const}\)}
\end{itemize}

Escribamos la ecuación como \(-k\,u_{xx} + u_t = 0\) y miremos
\textbf{solo la parte principal}, que es la de orden dos: acá es
únicamente \(-k\,u_{xx}\). Con las variables \((x,t)\) eso da \(A=-k\),
\(B=0\) y \(C=0\), así que \(\Delta = 0\) ---parabólica, como
sabíamos--- y la \ref{eq-edo-caracteristicas} queda
\[-k\left(\frac{dt}{dx}\right)^2 = 0 \qquad\Longrightarrow\qquad \frac{dt}{dx} = 0,\]
o sea \(t = \text{const}\): una sola familia, las líneas de tiempo
constante.

La tercera opción es la que corresponde a una elíptica, no a esta. Y la
cuarta es la respuesta a una pregunta distinta ---esas serían las
características si el término dominante fuera \(u_{tt} - k\,u_{xx}\), es
decir si la ecuación fuera de onda---.

\textbf{Lo que significa.} Que la única característica sea
\(t=\text{const}\) es la razón geométrica de que el calor se propague a
velocidad infinita: «viajar por una característica» es, acá, extenderse
por todo el espacio dentro del mismo instante.

\subsection{Elípticas}\label{sec-canonica-eliptica}

Con \(\Delta < 0\) no hay direcciones nulas reales. Las raíces de la
resolvente son complejas conjugadas:

\[
\frac{\xi_x}{\xi_y} = \frac{-B + i\sqrt{AC-B^2}}{A}, \qquad
\frac{\eta_x}{\eta_y} = \frac{-B - i\sqrt{AC-B^2}}{A} .
\]

Formalmente se puede seguir el procedimiento hiperbólico y llegar a
\(u_{\xi\eta} = G\), pero con \(\xi\) y \(\eta\) complejas conjugadas,
lo cual no sirve para trabajar. La salida es quedarse con la parte real
y la parte imaginaria:

\[
\sigma = \frac{\xi+\eta}{2} = \operatorname{Re}\xi, \qquad
\tau = \frac{\xi-\eta}{2i} = \operatorname{Im}\xi ,
\]

que son funciones reales, y en esas coordenadas la ecuación queda

\begin{equation}\protect\phantomsection\label{eq-canonica-eliptica}{
u_{\sigma\sigma} + u_{\tau\tau} = G(\sigma,\tau,u,u_\sigma,u_\tau).
}\end{equation}

O sea: \textbf{toda ecuación elíptica es, en las coordenadas adecuadas,
la ecuación de Laplace más términos de orden menor}. Es la razón de que
la teoría del potencial sirva para tantas cosas que no son
electrostática.

\begin{example}[Una
elíptica]\protect\hypertarget{exm-canonica-eliptica}{}\label{exm-canonica-eliptica}

Llevemos \(\;u_{xx} + x^2u_{yy} = 0\;\) a forma canónica en \(x>0\).

\(A=1\), \(B=0\), \(C=x^2\), y \(\Delta = -x^2 < 0\): elíptica. La
\ref{eq-edo-caracteristicas} da
\[\left(\frac{dy}{dx}\right)^2 + x^2 = 0 \ \Longrightarrow\ \frac{dy}{dx} = \pm ix
\ \Longrightarrow\ y \mp i\frac{x^2}{2} = \text{const},\]
características complejas. Tomando \(\xi = y + i\frac{x^2}{2}\), sus
partes real e imaginaria son
\[\sigma = y, \qquad \tau = \frac{x^2}{2} .\]

Con \(\sigma_x=0\), \(\sigma_y=1\), \(\tau_x=x\), \(\tau_y=0\), la
\ref{eq-alfa-beta-gamma} da
\[\alpha = C\,\sigma_y^2 = x^2, \qquad \gamma = A\,\tau_x^2 = x^2, \qquad \beta = 0,\]
---los dos coeficientes iguales y sin término cruzado, que es la firma
de una elíptica--- y para el orden uno, \(A\tau_{xx} + C\tau_{yy} = 1\).
Entonces \[
x^2\big(u_{\sigma\sigma} + u_{\tau\tau}\big) + u_\tau = 0 ,
\] y como \(x^2 = 2\tau\), \[
u_{\sigma\sigma} + u_{\tau\tau} + \frac{u_\tau}{2\tau} = 0 .
\] Laplace más un término de orden uno, tal como anunciamos.

\end{example}

\begin{exercise}[]\protect\hypertarget{exr-canonica-practica}{}\label{exr-canonica-practica}

Llevar cada ecuación a su forma canónica, indicando primero el tipo y
las curvas características.

\begin{enumerate}
\def\labelenumi{\alph{enumi}.}
\tightlist
\item
  \(u_{xx} - 4u_{xy} + 3u_{yy} = 0\)
\item
  \(u_{xx} + 2u_{xy} + u_{yy} = 0\)
\item
  \(4u_{xx} + u_{yy} = 0\)
\end{enumerate}

\end{exercise}

\begin{tcolorbox}[enhanced jigsaw, arc=.35mm, bottomrule=.15mm, bottomtitle=1mm, breakable, colback=white, colbacktitle=quarto-callout-tip-color!10!white, colframe=quarto-callout-tip-color-frame, coltitle=black, left=2mm, leftrule=.75mm, opacityback=0, opacitybacktitle=0.6, rightrule=.15mm, title=\textcolor{quarto-callout-tip-color}{\faLightbulb}\hspace{0.5em}{Solución}, titlerule=0mm, toprule=.15mm, toptitle=1mm]

\textbf{a.} \(A=1\), \(B=-2\), \(C=3\), \(\Delta = 4-3 = 1 > 0\):
\textbf{hiperbólica}. La \ref{eq-edo-caracteristicas} queda
\[\left(\frac{dy}{dx}\right)^2 + 4\,\frac{dy}{dx} + 3 = 0 \ \Longrightarrow\
\frac{dy}{dx} = -1 \ \text{ o } \ -3 ,\] o sea \(y+x=\text{const}\) e
\(y+3x=\text{const}\). Tomamos \(\xi = x+y\), \(\eta = 3x+y\).

Con \(\xi_x=\xi_y=1\), \(\eta_x=3\), \(\eta_y=1\) y todas las segundas
derivadas nulas, \(\alpha = 1-4+3 = 0\) ✓, \(\gamma = 9-12+3 = 0\) ✓, y
\[\beta = 2\big[1\cdot 3 + (-2)(1+3) + 3\cdot 1\big] = 2(3-8+3) = -4 .\]
Entonces \(-4u_{\xi\eta} = 0\), es decir
\[u_{\xi\eta} = 0 \qquad\Longrightarrow\qquad u = f(x+y) + g(3x+y).\]
Chequeo: para \(u = g(3x+y)\), \(u_{xx}=9g''\), \(u_{xy}=3g''\),
\(u_{yy}=g''\), y \(9-12+3=0\) ✓.

\textbf{b.} \(A=B=C=1\), \(\Delta = 1-1 = 0\): \textbf{parabólica}. Una
sola familia, \(dy/dx = B/A = 1\), o sea \(y-x=\text{const}\). Tomamos
\(\eta = y-x\) y elegimos \(\xi = x\), que es independiente.

\(\alpha = 1\), y se verifica \(\beta = 2[1\cdot(-1) + 1(1+0) + 0] = 0\)
✓ y \(\gamma = 1 - 2 + 1 = 0\) ✓. Los términos de orden uno son nulos
porque todo es constante. Queda
\[u_{\xi\xi} = 0 \qquad\Longrightarrow\qquad u = \xi\,f(\eta) + g(\eta)
= x\,f(y-x) + g(y-x).\]

\textbf{c.} \(A=4\), \(B=0\), \(C=1\), \(\Delta = -4 < 0\):
\textbf{elíptica}. La ecuación característica es \(4(y')^2 + 1 = 0\), o
sea \(dy/dx = \pm i/2\) y \(y \mp i\,x/2 = \text{const}\). Tomando
\(\xi = y + i\,x/2\), sus partes real e imaginaria son
\[\sigma = y, \qquad \tau = \frac{x}{2}.\] Con \(\sigma_x=0\),
\(\sigma_y=1\), \(\tau_x=1/2\), \(\tau_y=0\):
\[\alpha = C\,\sigma_y^2 = 1, \qquad \gamma = A\,\tau_x^2 = 4\cdot\tfrac14 = 1, \qquad
\beta = 0 ,\] y los términos de orden uno se anulan. Queda
\[u_{\sigma\sigma} + u_{\tau\tau} = 0 :\] \textbf{la ecuación de
Laplace}, en coordenadas donde el eje \(x\) está reescalado por un
factor \(2\). Físicamente, \(4u_{xx}+u_{yy}=0\) describe un medio en
equilibrio que conduce mejor en una dirección que en la otra; el cambio
de variables no hace más que estirar el eje corto hasta que el medio se
vuelve isótropo.

\end{tcolorbox}

\subsection{Resumen operativo}\label{sec-receta-canonica}

\begin{longtable}[]{@{}
  >{\raggedright\arraybackslash}p{(\linewidth - 6\tabcolsep) * \real{0.1600}}
  >{\raggedright\arraybackslash}p{(\linewidth - 6\tabcolsep) * \real{0.2000}}
  >{\raggedright\arraybackslash}p{(\linewidth - 6\tabcolsep) * \real{0.3800}}
  >{\raggedright\arraybackslash}p{(\linewidth - 6\tabcolsep) * \real{0.2600}}@{}}
\caption{La receta, en tres renglones. En los tres casos el primer paso
es el mismo: resolver la ecuación diferencial ordinaria
\ref{eq-edo-caracteristicas}.}\label{tbl-canonicas}\tabularnewline
\toprule\noalign{}
\begin{minipage}[b]{\linewidth}\raggedright
Tipo
\end{minipage} & \begin{minipage}[b]{\linewidth}\raggedright
Qué se anula
\end{minipage} & \begin{minipage}[b]{\linewidth}\raggedright
Cómo se elige el cambio
\end{minipage} & \begin{minipage}[b]{\linewidth}\raggedright
Forma canónica
\end{minipage} \\
\midrule\noalign{}
\endfirsthead
\toprule\noalign{}
\begin{minipage}[b]{\linewidth}\raggedright
Tipo
\end{minipage} & \begin{minipage}[b]{\linewidth}\raggedright
Qué se anula
\end{minipage} & \begin{minipage}[b]{\linewidth}\raggedright
Cómo se elige el cambio
\end{minipage} & \begin{minipage}[b]{\linewidth}\raggedright
Forma canónica
\end{minipage} \\
\midrule\noalign{}
\endhead
\bottomrule\noalign{}
\endlastfoot
Hiperbólica & \(\alpha = \gamma = 0\) & \(\xi\), \(\eta\) constantes
sobre \textbf{cada} familia de características & \(u_{\xi\eta} = G\), o
bien \(u_{\sigma\sigma}-u_{\tau\tau}=K\) \\
Parabólica & \(\beta = \gamma = 0\) & \(\eta\) sobre la \textbf{única}
familia; \(\xi\) libre e independiente & \(u_{\xi\xi} = G\) \\
Elíptica & \(\beta = 0\), \(\alpha=\gamma\) & \(\sigma\), \(\tau\) =
parte real e imaginaria de las características complejas &
\(u_{\sigma\sigma}+u_{\tau\tau}=G\) \\
\end{longtable}

\section{Muchas variables}\label{sec-muchas-variables}

Hasta acá trabajamos con dos variables. La física real casi nunca tiene
dos: tiene tres coordenadas espaciales y el tiempo, y solo cuando hay
una simetría que permite ignorar alguna se puede reducir el problema.
Así que la pregunta es obligatoria: \textbf{¿se puede clasificar una
ecuación de segundo orden en \(n\) variables?}

La respuesta corta es que sí, que aparecen tipos nuevos que en dos
variables no existían, y que hay una complicación seria a partir de tres
variables. Vamos por partes.

\subsection{La matriz de coeficientes}\label{sec-matriz-n}

Sea \(D\) un dominio abierto de \(\mathbb{R}^n\) y consideremos una
ecuación semilineal de segundo orden cuya parte principal es

\begin{equation}\protect\phantomsection\label{eq-parte-principal-n}{
L_P[u] = \sum_{i=1}^{n}\sum_{j=1}^{n} a_{ij}(x)\,u_{x_ix_j}, \qquad x \in D .
}\end{equation}

Podemos suponer que los coeficientes son \textbf{simétricos},
\(a_{ij}=a_{ji}\), sin perder nada: como \(u_{x_ix_j} = u_{x_jx_i}\)
para funciones suficientemente suaves, cualquier parte antisimétrica se
cancela sola. Y entonces, exactamente como en la
sección~\ref{sec-conicas}, la parte principal tiene asociada una
\textbf{matriz simétrica} \(M = (a_{ij})\) y una forma cuadrática

\begin{equation}\protect\phantomsection\label{eq-Q-n}{
Q(\zeta) = \sum_{i,j} a_{ij}\,\zeta_i\zeta_j = \zeta^{T}M\,\zeta , \qquad \zeta \in \mathbb{R}^n .
}\end{equation}

En dos variables, \(M = \begin{pmatrix} A & B \\ B & C\end{pmatrix}\) y
\(Q(\xi,\eta)\) era la cónica del
Definición~\ref{def-forma-cuadratica-asociada}. En \(n\) variables,
\(Q(\zeta)=1\) es una \textbf{cuádrica} ---un elipsoide, un
hiperboloide, un cilindro--- y el tipo de la ecuación va a ser, otra
vez, el tipo de esa superficie.

La ecuación \(\;u_{xx} + 4u_{xy} + u_{yy} + 2u_{xz} + u_{zz} = 0\;\)
tiene matriz de coeficientes simétrica \(M=(a_{ij})\). ¿Cuánto vale
\(a_{12}\)?

\begin{itemize}
\tightlist
\item
  {\(4\)}
\item
  {\(2\)}
\item
  {\(1\)}
\item
  {\(0\)}
\end{itemize}

El término \(4u_{xy}\) aparece \textbf{dos veces} en la suma doble
\(\sum_{i,j}a_{ij}u_{x_ix_j}\): una como \(a_{12}u_{xy}\) y otra como
\(a_{21}u_{yx}\). Como pedimos que la matriz sea simétrica, esos dos
coeficientes son iguales y cada uno se lleva la mitad:
\(a_{12}=a_{21}=2\).

Es exactamente la misma historia que el factor \(2\) de la
sección~\ref{sec-clasificacion}, pero acá aparece sola: la simetría de
la matriz obliga a repartir. Por eso conviene escribir la ecuación de
dos variables como \(A u_{xx} + 2B u_{xy} + C u_{yy}\), para que \(A\),
\(B\) y \(C\) sean directamente las entradas.

La matriz completa de este ejemplo es
\[M = \begin{pmatrix} 1 & 2 & 1 \\ 2 & 1 & 0 \\ 1 & 0 & 1\end{pmatrix},\]
con \(a_{13}=a_{31}=1\) (la mitad de \(2\)) y \(a_{23}=0\) porque no hay
término \(u_{yz}\).

\subsection{Diagonalizar: la signatura}\label{sec-signatura}

En el punto \(x^0\) hacemos un cambio lineal de coordenadas

\begin{equation}\protect\phantomsection\label{eq-cambio-lineal}{
y_i = \sum_{j=1}^{n} c_{ij}\,\big(x_j - x_j^0\big) ,
}\end{equation}

y por la regla de la cadena la parte principal se transforma en
\(M \mapsto C\,M\,C^{T}\), con \(C=(c_{ij})\). Como \(M\) es simétrica,
el álgebra lineal garantiza que se puede elegir \(C\) de modo que la
matriz nueva sea \textbf{diagonal}, y además que reescalando cada eje
los elementos de la diagonal queden todos iguales a \(+1\), \(-1\) o
\(0\).

\begin{definition}[Signatura]\protect\hypertarget{def-signatura}{}\label{def-signatura}

Diagonalizada la matriz de coeficientes en un punto, llamamos

\begin{itemize}
\tightlist
\item
  \(n_+\) = cantidad de unos positivos,
\item
  \(n_-\) = cantidad de unos negativos,
\item
  \(n_0\) = cantidad de ceros,
\end{itemize}

con \(n_+ + n_- + n_0 = n\). La terna \((n_+, n_-, n_0)\) se llama
\textbf{signatura} de la forma cuadrática en ese punto.

La \textbf{ley de inercia de Sylvester} asegura que la signatura no
depende de cómo se haya diagonalizado: dos diagonalizaciones distintas
dan siempre los mismos tres números. Por eso tiene sentido usarla para
clasificar.

\end{definition}

\begin{definition}[Clasificación en \(n\)
variables]\protect\hypertarget{def-clasificacion-n}{}\label{def-clasificacion-n}

Sea una ecuación semilineal de segundo orden en un entorno de
\(x^0 \in \mathbb{R}^n\). En \(x^0\) se dice

\begin{itemize}
\tightlist
\item
  \textbf{elíptica} si \(n_+ = n\) o \(n_- = n\) (todos los signos
  iguales, ningún cero);
\item
  \textbf{parabólica} si \(n_0 > 0\) (hay al menos un cero);
\item
  \textbf{hiperbólica} si \(n_0 = 0\) y \(n_+ = 1\) o \(n_- = 1\)
  (ningún cero y un solo signo distinto);
\item
  \textbf{ultrahiperbólica} si \(n_0 = 0\), \(n_+ > 1\) y \(n_- > 1\).
\end{itemize}

\end{definition}

Para \(n=2\) esto reproduce exactamente lo de la
sección~\ref{sec-clasificacion} ---elíptica es \((2,0)\) o \((0,2)\),
hiperbólica es \((1,1)\), parabólica tiene un cero--- y no hay lugar
para ultrahiperbólicas, porque harían falta al menos cuatro variables.
La primera aparece en \(n=4\).

Una ecuación en cuatro variables tiene signatura \((2,2,0)\). ¿De qué
tipo es?

\begin{itemize}
\tightlist
\item
  {Elíptica}
\item
  {Hiperbólica}
\item
  {Parabólica}
\item
  {Ultrahiperbólica}
\end{itemize}

No hay ceros, así que no es parabólica. No están todos los signos
iguales, así que no es elíptica. Y para ser hiperbólica haría falta que
\textbf{uno solo} de los dos números fuera \(1\); acá los dos valen
\(2\). Es el caso ultrahiperbólico de la
Definición~\ref{def-clasificacion-n}, que aparece recién a partir de
cuatro variables.

Y no vas a encontrarte muchas: una signatura \((2,2,0)\) significa dos
direcciones «tipo tiempo» independientes, es decir dos tiempos. Ninguna
teoría física estándar los tiene, y los problemas de valores iniciales
para estas ecuaciones son notoriamente mal planteados. Es un tipo que
existe por completitud de la clasificación más que por utilidad.

\subsection{Las cuádricas}\label{sec-cuadricas}

Igual que en dos variables el tipo era el de la cónica, acá es el de la
cuádrica \(Q(\zeta)=1\), y la lista es corta.

\begin{longtable}[]{@{}
  >{\raggedright\arraybackslash}p{(\linewidth - 8\tabcolsep) * \real{0.1600}}
  >{\raggedright\arraybackslash}p{(\linewidth - 8\tabcolsep) * \real{0.2600}}
  >{\raggedright\arraybackslash}p{(\linewidth - 8\tabcolsep) * \real{0.1800}}
  >{\raggedright\arraybackslash}p{(\linewidth - 8\tabcolsep) * \real{0.1600}}
  >{\raggedright\arraybackslash}p{(\linewidth - 8\tabcolsep) * \real{0.2400}}@{}}
\caption{Las cuádricas en tres variables y las ecuaciones que les
corresponden.}\label{tbl-cuadricas}\tabularnewline
\toprule\noalign{}
\begin{minipage}[b]{\linewidth}\raggedright
Signatura
\end{minipage} & \begin{minipage}[b]{\linewidth}\raggedright
Cuádrica \(Q=1\)
\end{minipage} & \begin{minipage}[b]{\linewidth}\raggedright
Direcciones nulas
\end{minipage} & \begin{minipage}[b]{\linewidth}\raggedright
Tipo
\end{minipage} & \begin{minipage}[b]{\linewidth}\raggedright
Ejemplo
\end{minipage} \\
\midrule\noalign{}
\endfirsthead
\toprule\noalign{}
\begin{minipage}[b]{\linewidth}\raggedright
Signatura
\end{minipage} & \begin{minipage}[b]{\linewidth}\raggedright
Cuádrica \(Q=1\)
\end{minipage} & \begin{minipage}[b]{\linewidth}\raggedright
Direcciones nulas
\end{minipage} & \begin{minipage}[b]{\linewidth}\raggedright
Tipo
\end{minipage} & \begin{minipage}[b]{\linewidth}\raggedright
Ejemplo
\end{minipage} \\
\midrule\noalign{}
\endhead
\bottomrule\noalign{}
\endlastfoot
\((3,0,0)\) & elipsoide & ninguna & elíptica & \(\nabla^2u=0\) en el
espacio \\
\((2,1,0)\) & hiperboloide de una hoja & un \textbf{cono} & hiperbólica
& \(u_{tt}=c^2(u_{xx}+u_{yy})\) \\
\((2,0,1)\) & cilindro elíptico & una recta & parabólica &
\(u_t = k(u_{xx}+u_{yy})\) \\
\end{longtable}

\begin{figure}

\centering{

\pandocbounded{\includegraphics[keepaspectratio]{08-edp_files/figure-pdf/fig-cuadricas-output-1.pdf}}

}

\caption{\label{fig-cuadricas}Las tres cuádricas de la
Tabla~\ref{tbl-cuadricas}. Izquierda: el elipsoide de una elíptica
---cerrado, sin ninguna dirección en la que \(Q\) se anule---. Centro:
el cono de las direcciones nulas de una hiperbólica; para la ecuación de
onda en dos dimensiones espaciales, ese cono es \textbf{el cono de luz},
y separa las direcciones por las que la señal puede propagarse de las
que le están prohibidas. Derecha: el cilindro de una parabólica, abierto
en la dirección del tiempo, que es la dirección en la que la forma
cuadrática no dice nada.}

\end{figure}%

\emph{(En la versión web hay acá un gráfico interactivo con cinco
sliders: los tres autovalores \(\lambda_1,\lambda_2,\lambda_3\) y dos
para girar la cámara. Muestra la cuádrica \(Q=1\) junto con su cono de
direcciones nulas, y cómo el elipsoide se abre en hiperboloide al
cambiar de signo uno de los autovalores, pasando por el cilindro
parabólico.)}

\begin{tcolorbox}[enhanced jigsaw, arc=.35mm, bottomrule=.15mm, breakable, colback=white, colframe=quarto-callout-important-color-frame, left=2mm, leftrule=.75mm, opacityback=0, rightrule=.15mm, toprule=.15mm]
\begin{minipage}[t]{5.5mm}
\textcolor{quarto-callout-important-color}{\faExclamation}
\end{minipage}%
\begin{minipage}[t]{\textwidth - 5.5mm}

\vspace{-3mm}\textbf{El cono de luz es una cuádrica}\vspace{3mm}

Vale la pena detenerse en el panel del medio, porque es de las cosas más
lindas del capítulo.

La ecuación de onda en tres dimensiones espaciales,
\[u_{tt} - c^2\big(u_{xx}+u_{yy}+u_{zz}\big) = 0,\] tiene, en las
variables \((t,x,y,z)\), matriz de coeficientes
\(\operatorname{diag}(1,-c^2,-c^2,-c^2)\): signatura \((1,3,0)\),
hiperbólica. Su forma cuadrática es
\[Q = \tau^2 - c^2\big(\zeta_1^2+\zeta_2^2+\zeta_3^2\big),\] que es
\textbf{exactamente la métrica de Minkowski} de la relatividad especial,
y sus direcciones nulas ---donde \(Q=0\)--- son las que cumplen
\(|\zeta| = \tau/c\): el \textbf{cono de luz}.

O sea: la clasificación de la ecuación de onda y la estructura causal
del espacio-tiempo son el mismo objeto matemático. Que la ecuación sea
hiperbólica y que exista una velocidad máxima de propagación son dos
maneras de decir lo mismo, y la superficie característica que separa el
adentro del afuera del cono es la que hace que una señal emitida en un
punto llegue a unos lugares y a otros todavía no.

Y de paso queda claro por qué las \textbf{ultrahiperbólicas} casi no
aparecen en física: harían falta dos o más direcciones «tipo tiempo», o
sea dos tiempos independientes. No hay ninguna teoría física estándar
que los tenga.

\end{minipage}%
\end{tcolorbox}

¿Por qué la ecuación del calor \(u_t = k\,\nabla^2 u\) es parabólica en
las cuatro variables \((t,x,y,z)\)?

\begin{itemize}
\tightlist
\item
  {Porque toda la fila y la columna del tiempo de su matriz de
  coeficientes son cero}
\item
  {Porque tiene tres signos iguales y uno distinto}
\item
  {Porque \(k>0\)}
\item
  {Porque el tiempo no se puede tratar como una coordenada más}
\end{itemize}

La matriz de coeficientes se arma \textbf{solo con las derivadas de
orden dos}. En la ecuación del calor, el tiempo aparece únicamente en
\(u_t\), que es de orden uno, así que no entra en la matriz: la fila y
la columna correspondientes a \(t\) quedan en cero. Con la ordenación
\((t,x,y,z)\), \[M = \operatorname{diag}\big(0,\,k,\,k,\,k\big),\] así
que \(n_0 = 1\) y por la Definición~\ref{def-clasificacion-n} la
ecuación es parabólica.

La segunda opción describe a la ecuación de \textbf{onda}, cuya matriz
es \(\operatorname{diag}(1,-c^2,-c^2,-c^2)\): signatura \((1,3,0)\),
hiperbólica. Y ahí está la diferencia entre las dos ecuaciones más
importantes de la física, en una línea: \textbf{el calor es de primer
orden en el tiempo y la onda es de segundo}, y por eso una tiene un cero
en la signatura y la otra no.

Una aclaración sobre signos: si multiplicamos la ecuación por \(-1\) la
signatura pasa de \((3,0,1)\) a \((0,3,1)\). No importa, porque la
Definición~\ref{def-clasificacion-n} está escrita para que el tipo no
dependa del signo global.

\subsection{La complicación que aparece a partir de tres
variables}\label{sec-obstruccion}

Todo lo anterior fue \textbf{en un punto}. La pregunta siguiente es si
se puede hacer el cambio de variables de manera que la ecuación quede en
forma canónica en \textbf{todo un entorno} del punto, y acá aparece el
problema.

Contemos. Un cambio de coordenadas en un entorno da \(n\) funciones
libres ---una por coordenada nueva---. Lo que hay que pedir es que se
anulen todos los elementos de fuera de la diagonal de la matriz
transformada, y una matriz simétrica de \(n\times n\) tiene

\[
\frac{n(n-1)}{2}
\]

elementos fuera de la diagonal, contando cada par una sola vez.
Comparemos:

\begin{longtable}[]{@{}
  >{\centering\arraybackslash}p{(\linewidth - 6\tabcolsep) * \real{0.1000}}
  >{\centering\arraybackslash}p{(\linewidth - 6\tabcolsep) * \real{0.2800}}
  >{\centering\arraybackslash}p{(\linewidth - 6\tabcolsep) * \real{0.3000}}
  >{\raggedright\arraybackslash}p{(\linewidth - 6\tabcolsep) * \real{0.3200}}@{}}
\caption{El conteo que arruina todo.}\label{tbl-conteo}\tabularnewline
\toprule\noalign{}
\begin{minipage}[b]{\linewidth}\centering
\(n\)
\end{minipage} & \begin{minipage}[b]{\linewidth}\centering
funciones disponibles
\end{minipage} & \begin{minipage}[b]{\linewidth}\centering
condiciones a imponer
\end{minipage} & \begin{minipage}[b]{\linewidth}\raggedright
¿alcanza?
\end{minipage} \\
\midrule\noalign{}
\endfirsthead
\toprule\noalign{}
\begin{minipage}[b]{\linewidth}\centering
\(n\)
\end{minipage} & \begin{minipage}[b]{\linewidth}\centering
funciones disponibles
\end{minipage} & \begin{minipage}[b]{\linewidth}\centering
condiciones a imponer
\end{minipage} & \begin{minipage}[b]{\linewidth}\raggedright
¿alcanza?
\end{minipage} \\
\midrule\noalign{}
\endhead
\bottomrule\noalign{}
\endlastfoot
\(2\) & \(2\) & \(1\) & sí, sobra \\
\(3\) & \(3\) & \(3\) & justo \\
\(4\) & \(4\) & \(6\) & \textbf{no} \\
\(n\) & \(n\) & \(n(n-1)/2\) & no, para \(n \ge 4\) \\
\end{longtable}

Para \(n \ge 4\) el sistema está \textbf{sobredeterminado} y en general
no tiene solución. Para \(n=3\) la cuenta da justo y la matriz se puede
diagonalizar, pero los elementos de la diagonal pueden salir distintos
en puntos distintos del entorno, así que tampoco se consigue una forma
canónica con coeficientes fijos.

\begin{tcolorbox}[enhanced jigsaw, arc=.35mm, bottomrule=.15mm, breakable, colback=white, colframe=quarto-callout-warning-color-frame, left=2mm, leftrule=.75mm, opacityback=0, rightrule=.15mm, toprule=.15mm]
\begin{minipage}[t]{5.5mm}
\textcolor{quarto-callout-warning-color}{\faExclamationTriangle}
\end{minipage}%
\begin{minipage}[t]{\textwidth - 5.5mm}

\vspace{-3mm}\textbf{Qué se pierde y qué no}\vspace{3mm}

En resumen: \textbf{para \(n\ge 3\) no se puede reducir la ecuación a
forma canónica en un entorno}, solo en un punto. La única excepción es
que los coeficientes sean constantes, en cuyo caso el cambio lineal del
\ref{eq-cambio-lineal} sirve en todos lados a la vez.

Esto casi nunca se dice en los cursos, porque casi todos los ejemplos
que se hacen tienen coeficientes constantes y el problema no se nota.
Pero conviene tenerlo presente para no esperar de más del método.

Lo que \textbf{no} se pierde es la clasificación. El tipo sigue estando
bien definido en cada punto, sigue siendo el que decide cómo se
comportan las soluciones, y sigue diciendo qué datos hacen falta. Que no
exista un cambio de variables global no cambia nada de eso.

Y para \(n=2\) todo funciona incluso con coeficientes variables: por eso
pudimos hacer la sección~\ref{sec-formas-canonicas} sin ninguna
salvedad.

\end{minipage}%
\end{tcolorbox}

\begin{example}[Una lineal en tres
variables]\protect\hypertarget{exm-tres-variables-clasif}{}\label{exm-tres-variables-clasif}

Clasifiquemos \(\;u_{xx} - 2u_{xz} + 2u_{yy} + 4u_{yz} + 6u_{zz} = 0\).

Primero la matriz. Ojo con los términos cruzados: como \(u_{xz}\) y
\(u_{zx}\) son lo mismo, el coeficiente \(-2\) se reparte entre
\(a_{13}\) y \(a_{31}\), y a cada uno le toca \(-1\). Lo mismo con el
\(4\) de \(u_{yz}\), que da \(a_{23}=a_{32}=2\). Entonces

\[
M = \begin{pmatrix} 1 & 0 & -1 \\ 0 & 2 & 2 \\ -1 & 2 & 6 \end{pmatrix}.
\]

Para la signatura no hace falta calcular autovalores: alcanza con
\textbf{completar cuadrados} en
\[Q = \zeta_1^2 + 2\zeta_2^2 + 6\zeta_3^2 - 2\zeta_1\zeta_3 + 4\zeta_2\zeta_3 .\]
Juntando lo que tiene \(\zeta_1\): \[
\zeta_1^2 - 2\zeta_1\zeta_3 = (\zeta_1-\zeta_3)^2 - \zeta_3^2 ,
\] y con lo que queda,
\(2\zeta_2^2 + 4\zeta_2\zeta_3 + 5\zeta_3^2 = 2(\zeta_2+\zeta_3)^2 + 3\zeta_3^2\).
Sumando todo, \[
Q = (\zeta_1-\zeta_3)^2 + 2(\zeta_2+\zeta_3)^2 + 3\zeta_3^2 .
\] Tres cuadrados, los tres con signo \(+\): la signatura es \((3,0,0)\)
y la ecuación es \textbf{elíptica}.

Y como los coeficientes son constantes, esta forma vale en todo punto:
es la excepción de la que hablábamos. El cambio de variables que la
lleva a Laplace es, leyendo los paréntesis, \[
y_1 = x-z, \qquad y_2 = \sqrt{2}\,(y+z), \qquad y_3 = \sqrt{3}\,z .
\]

Completar cuadrados es, en general, mucho más rápido que diagonalizar, y
da la signatura directamente: \textbf{contar los signos de los
cuadrados}.

\end{example}

\begin{exercise}[]\protect\hypertarget{exr-signaturas}{}\label{exr-signaturas}

Hallar la signatura y el tipo de cada ecuación, completando cuadrados.

\begin{enumerate}
\def\labelenumi{\alph{enumi}.}
\tightlist
\item
  \(u_{xx} + u_{yy} - u_{zz} = 0\)
\item
  \(u_{xx} + 2u_{xy} + u_{yy} + u_{zz} = 0\)
\item
  \(u_{tt} - u_{xx} - u_{yy} - u_{zz} + u_t = 0\)
\item
  \(u_{xx} + 4u_{xy} + 5u_{yy} + 2u_{yz} + 2u_{zz} = 0\)
\end{enumerate}

\end{exercise}

\begin{tcolorbox}[enhanced jigsaw, arc=.35mm, bottomrule=.15mm, bottomtitle=1mm, breakable, colback=white, colbacktitle=quarto-callout-tip-color!10!white, colframe=quarto-callout-tip-color-frame, coltitle=black, left=2mm, leftrule=.75mm, opacityback=0, opacitybacktitle=0.6, rightrule=.15mm, title=\textcolor{quarto-callout-tip-color}{\faLightbulb}\hspace{0.5em}{Solución}, titlerule=0mm, toprule=.15mm, toptitle=1mm]

\textbf{a.} Ya está diagonal: signatura \((2,1,0)\). Ningún cero y un
solo signo negativo: \textbf{hiperbólica}. Es la ecuación de onda en dos
dimensiones espaciales, con \(z\) haciendo de tiempo y \(c=1\). Su cono
de direcciones nulas es el del panel del medio de la
Figura~\ref{fig-cuadricas}.

\textbf{b.}
\(Q = \zeta_1^2 + 2\zeta_1\zeta_2 + \zeta_2^2 + \zeta_3^2 = (\zeta_1+\zeta_2)^2 + \zeta_3^2\).
Solo \textbf{dos} cuadrados, los dos positivos, y la variable
\(\zeta_2\) ya no aparece sola: la signatura es \((2,0,1)\). Hay un
cero, así que es \textbf{parabólica}. Geométricamente \(Q=1\) es un
cilindro, como el panel derecho de la Figura~\ref{fig-cuadricas},
abierto en la dirección \(\zeta_1+\zeta_2=0\).

\textbf{c.} El término \(u_t\) es de orden uno y \textbf{no entra en la
matriz}. La parte principal es \(u_{tt}-u_{xx}-u_{yy}-u_{zz}\), con
signatura \((1,3,0)\): \textbf{hiperbólica}. Es la ecuación de onda con
un término de amortiguamiento ---el telegrafista---, y el
amortiguamiento no cambia el tipo, igual que en el
Ejercicio~\ref{exr-clasificar}. Físicamente tiene sentido: una onda
amortiguada sigue siendo una onda, sigue teniendo cono de luz, y sigue
propagándose a velocidad \(c\); lo único que hace el rozamiento es
apagarla.

\textbf{d.}
\(Q = \zeta_1^2 + 4\zeta_1\zeta_2 + 5\zeta_2^2 + 2\zeta_2\zeta_3 + 2\zeta_3^2\).
Empezamos por \(\zeta_1\):
\[\zeta_1^2 + 4\zeta_1\zeta_2 = (\zeta_1+2\zeta_2)^2 - 4\zeta_2^2 ,\] y
queda
\(5\zeta_2^2 - 4\zeta_2^2 + 2\zeta_2\zeta_3 + 2\zeta_3^2 = \zeta_2^2 + 2\zeta_2\zeta_3 + 2\zeta_3^2\).
Seguimos con \(\zeta_2\):
\[\zeta_2^2 + 2\zeta_2\zeta_3 = (\zeta_2+\zeta_3)^2 - \zeta_3^2 ,\] y
queda \(2\zeta_3^2 - \zeta_3^2 = \zeta_3^2\). En total
\[Q = (\zeta_1+2\zeta_2)^2 + (\zeta_2+\zeta_3)^2 + \zeta_3^2 ,\] tres
cuadrados positivos: signatura \((3,0,0)\), \textbf{elíptica}.

El método sirve siempre y es mucho más rápido que buscar autovalores: se
va eliminando una variable por vez, y al final se cuentan los signos.

\end{tcolorbox}

\begin{exercise}[]\protect\hypertarget{exr-cono-luz}{}\label{exr-cono-luz}

La ecuación de onda en el espacio, \(u_{tt} = c^2\nabla^2u\), tiene
direcciones nulas sobre el cono \(|\zeta| = \tau/c\).

\begin{enumerate}
\def\labelenumi{\alph{enumi}.}
\tightlist
\item
  Una fuente puntual emite una señal en el origen en \(t=0\). Usando el
  cono, decidí en qué región del espacio-tiempo la señal puede haber
  llegado, y dibujala para el caso de una sola dimensión espacial.
\item
  En un aula de \(10\) m, con \(c = 340\) m/s, ¿cuánto tarda un sonido
  emitido en el frente en llegar al fondo? ¿Y cuánto tardaría si la
  ecuación fuera parabólica?
\item
  ¿Qué le pasaría a la estructura causal si el término \(u_{tt}\)
  tuviera el mismo signo que el laplaciano, es decir si la ecuación
  fuera \(u_{tt} + c^2\nabla^2u = 0\)?
\end{enumerate}

\end{exercise}

\begin{tcolorbox}[enhanced jigsaw, arc=.35mm, bottomrule=.15mm, bottomtitle=1mm, breakable, colback=white, colbacktitle=quarto-callout-tip-color!10!white, colframe=quarto-callout-tip-color-frame, coltitle=black, left=2mm, leftrule=.75mm, opacityback=0, opacitybacktitle=0.6, rightrule=.15mm, title=\textcolor{quarto-callout-tip-color}{\faLightbulb}\hspace{0.5em}{Solución}, titlerule=0mm, toprule=.15mm, toptitle=1mm]

\textbf{a.} La señal solo puede haber llegado a los puntos que están
\textbf{dentro o sobre} el cono \(|x| \le c\,t\), con \(t>0\): el cono
de luz futuro. Fuera de él la solución sigue siendo cero. En una
dimensión espacial el cono es un par de rectas de pendiente \(\pm 1/c\)
saliendo del origen, y la región alcanzada es el triángulo que
encierran, que es exactamente el rango de influencia del
Ejercicio~\ref{exr-dominio-dependencia}.

\textbf{b.} \(t = 10/340 \approx 0{,}029\) s, unos \(29\) milisegundos.

Si la ecuación fuera parabólica, \textbf{cero}: llegaría
instantáneamente, aunque con una amplitud ridículamente chica. Ese es el
precio conceptual de la ecuación del calor ---predice velocidad de
propagación infinita, que es físicamente imposible--- y la razón de que
existan modelos hiperbólicos de conducción del calor para situaciones
donde eso importa, como pulsos de láser ultracortos.

\textbf{c.} Se convertiría en \(\nabla^2\) de cuatro variables, con
signatura \((4,0,0)\): \textbf{elíptica}. Y entonces no habría cono, no
habría direcciones nulas y no habría estructura causal en absoluto: el
valor del campo en cada punto del espacio-tiempo dependería de
\textbf{todo} el borde a la vez, pasado y futuro incluidos. Además, por
la sección~\ref{sec-hadamard}, el problema de valores iniciales estaría
mal planteado.

Dicho de otro modo: \textbf{el signo relativo entre \(u_{tt}\) y el
laplaciano es lo que hace que el tiempo sea el tiempo.} Cambiarlo no da
una física distinta, da una física imposible.

\end{tcolorbox}

\section{Transformadas integrales}\label{sec-transformadas}

Los tres métodos que faltan ---este, el de la
sección~\ref{sec-green-varias} y el de la
sección~\ref{sec-separacion}--- tienen algo en común que conviene decir
de entrada: \textbf{los tres necesitan que la ecuación sea lineal}. Se
acabó el territorio de las cuasi-lineales. A cambio, cada uno resuelve
problemas concretos de punta a punta.

\subsection{La idea: cambiar derivadas por
multiplicaciones}\label{sec-idea-transformada}

De los capítulos Capítulo~\ref{sec-laplace} y
Capítulo~\ref{sec-tfourier} nos llevamos una sola propiedad, y es toda
la sección:

\[
\mathcal{L}\{f'\}(p) = p\,F(p) - f(0), \qquad
\mathcal{F}\{f'\}(k) = i k\,\hat f(k) .
\]

\textbf{Derivar se convierte en multiplicar.} Y si una de las variables
de una EDP desaparece porque sus derivadas se volvieron factores, la
ecuación pierde una dimensión: una EDP en \((x,y,z,t)\) se transforma en
una EDP en \((x,y,z)\) solamente, con el parámetro de la transformada
mirando desde afuera.

\begin{tcolorbox}[enhanced jigsaw, arc=.35mm, bottomrule=.15mm, breakable, colback=white, colframe=quarto-callout-important-color-frame, left=2mm, leftrule=.75mm, opacityback=0, rightrule=.15mm, toprule=.15mm]
\begin{minipage}[t]{5.5mm}
\textcolor{quarto-callout-important-color}{\faExclamation}
\end{minipage}%
\begin{minipage}[t]{\textwidth - 5.5mm}

\vspace{-3mm}\textbf{Qué está pasando, en términos físicos}\vspace{3mm}

Transformar Fourier en el tiempo es \textbf{descomponer el forzado en
frecuencias puras}. Para cada frecuencia por separado uno pregunta: si
empujo este sistema con una fuerza que oscila a frecuencia \(\omega\),
¿cómo responde? Esa pregunta ya no tiene tiempo adentro ---la respuesta
oscila a la misma frecuencia, solo hay que averiguar su amplitud y su
fase en cada punto del espacio--- y por eso la ecuación que queda es más
simple.

Es exactamente lo que hace cualquiera que mide la \textbf{respuesta en
frecuencia} de un sistema: en vez de golpearlo y mirar qué pasa, lo
excita a una frecuencia por vez y anota la amplitud. Después, si el
sistema es lineal, la respuesta a cualquier forzado se arma sumando.

El precio es tener que volver: al final hay que
\textbf{antitransformar}, y ese suele ser el paso difícil.

\end{minipage}%
\end{tcolorbox}

Si \(u(x,0)=u_t(x,0)=0\), ¿en qué se convierte \(u_{tt}\) al aplicar la
transformada de Laplace en el tiempo?

\begin{itemize}
\tightlist
\item
  {\(p^2\,U\)}
\item
  {\(p\,U\)}
\item
  {\(U''\)}
\item
  {\(p^2 U - p\,u(x,0)\)}
\end{itemize}

La fórmula general es
\[\mathcal{L}\{u_{tt}\} = p^2U - p\,u(x,0) - u_t(x,0),\] y con las dos
condiciones iniciales nulas quedan solo \(p^2U\). La cuarta opción es la
fórmula a la que se le olvidó un término: es el error más común, y da
resultados que parecen razonables hasta que uno chequea la condición
inicial.

Lo importante es lo otro: \textbf{derivar dos veces se volvió
multiplicar por \(p^2\)}, y de paso las condiciones iniciales entraron
solas a la ecuación transformada. Esa es la gran ventaja de la
transformada de Laplace sobre la separación de variables de la
sección~\ref{sec-separacion}, donde las condiciones iniciales hay que
imponerlas a mano al final, calculando coeficientes de Fourier.

\subsection{El problema
general}\label{sec-problema-general-transformadas}

Consideremos, en un volumen acotado \(D\) con superficie cerrada \(S\),
la ecuación

\begin{equation}\protect\phantomsection\label{eq-general-transformadas}{
\nabla^2 u = a\,u + b\,\frac{\partial u}{\partial t} + c\,\frac{\partial^2 u}{\partial t^2}
- f(x,y,z,t) .
}\end{equation}

Vale la pena mirar qué es cada término, porque esta única ecuación
contiene a las tres modelo:

\begin{longtable}[]{@{}
  >{\raggedright\arraybackslash}p{(\linewidth - 4\tabcolsep) * \real{0.1600}}
  >{\raggedright\arraybackslash}p{(\linewidth - 4\tabcolsep) * \real{0.4400}}
  >{\raggedright\arraybackslash}p{(\linewidth - 4\tabcolsep) * \real{0.4000}}@{}}
\caption{Una ecuación, cuatro
físicas.}\label{tbl-terminos}\tabularnewline
\toprule\noalign{}
\begin{minipage}[b]{\linewidth}\raggedright
Término
\end{minipage} & \begin{minipage}[b]{\linewidth}\raggedright
Qué representa
\end{minipage} & \begin{minipage}[b]{\linewidth}\raggedright
Se queda solo en\ldots{}
\end{minipage} \\
\midrule\noalign{}
\endfirsthead
\toprule\noalign{}
\begin{minipage}[b]{\linewidth}\raggedright
Término
\end{minipage} & \begin{minipage}[b]{\linewidth}\raggedright
Qué representa
\end{minipage} & \begin{minipage}[b]{\linewidth}\raggedright
Se queda solo en\ldots{}
\end{minipage} \\
\midrule\noalign{}
\endhead
\bottomrule\noalign{}
\endlastfoot
\(c\,u_{tt}\) & inercia & la \textbf{ecuación de onda}, con
\(c = 1/v^2\) \\
\(b\,u_t\) & disipación, difusión & la \textbf{ecuación del calor}, con
\(b = 1/k\) \\
\(a\,u\) & un término restitutivo o de absorción & Helmholtz y
Klein--Gordon \\
\(f\) & la fuente o el forzado externo & \textbf{Poisson}, si los otros
tres son cero \\
\end{longtable}

Para determinar la solución hacen falta dos cosas más.

\textbf{Condiciones de contorno}, prescritas sobre \(S\) e
independientes de \(t\). Son las del Definición~\ref{def-condiciones}:
Dirichlet, Neumann, o la combinación de Robin

\begin{equation}\protect\phantomsection\label{eq-robin-t}{
\frac{\partial u}{\partial n} + \alpha\,u = 0 \qquad\text{sobre } S .
}\end{equation}

\textbf{Condiciones iniciales}, en \(t=0\), para la función y su
derivada temporal:

\begin{equation}\protect\phantomsection\label{eq-iniciales-t}{
u(x,y,z,0) = u_0(x,y,z), \qquad
\left.\frac{\partial u}{\partial t}\right|_{t=0} = v_0(x,y,z) ,
}\end{equation}

y se asume que para \(t<0\) tanto \(f\) como \(u\) son nulas: el sistema
estaba quieto y en \(t=0\) empezó a pasar algo.

\subsection{La transformación}\label{sec-transformar}

Definimos las transformadas en el tiempo

\begin{equation}\protect\phantomsection\label{eq-def-transformadas}{
\Phi(x,y,z,s) = \frac{1}{\sqrt{2\pi}}\int_0^{\infty} u(x,y,z,t)\,e^{-ist}\,dt, \qquad
F(x,y,z,s) = \frac{1}{\sqrt{2\pi}}\int_0^{\infty} f(x,y,z,t)\,e^{-ist}\,dt .
}\end{equation}

Como el laplaciano solo deriva respecto de las variables espaciales,
atraviesa la integral sin hacer nada:
\(\mathcal{T}\{\nabla^2u\} = \nabla^2\Phi\). Las derivadas temporales,
en cambio, se convierten en factores y arrastran las condiciones
iniciales. Aplicando la transformada a \ref{eq-general-transformadas}
queda

\begin{equation}\protect\phantomsection\label{eq-transformada-general}{
\nabla^2\Phi = a\,\Phi + b\left(is\,\Phi - \frac{u_0}{\sqrt{2\pi}}\right)
+ c\left(-s^2\,\Phi - \frac{is\,u_0 + v_0}{\sqrt{2\pi}}\right) - F ,
}\end{equation}

y agrupando todo lo que no depende de \(\Phi\) del lado derecho,

\begin{equation}\protect\phantomsection\label{eq-helmholtz-generado}{
\nabla^2\Phi + \big(cs^2 - ibs - a\big)\,\Phi = g(x,y,z,s) ,
}\end{equation}

donde \(g\) junta la fuente transformada y los datos iniciales. La
condición de contorno \ref{eq-robin-t}, al ser independiente de \(t\),
se transforma en sí misma:

\[
\frac{\partial\Phi}{\partial n} + \alpha\,\Phi = 0 \qquad\text{sobre } S .
\]

\begin{tcolorbox}[enhanced jigsaw, arc=.35mm, bottomrule=.15mm, breakable, colback=white, colframe=quarto-callout-important-color-frame, left=2mm, leftrule=.75mm, opacityback=0, rightrule=.15mm, toprule=.15mm]
\begin{minipage}[t]{5.5mm}
\textcolor{quarto-callout-important-color}{\faExclamation}
\end{minipage}%
\begin{minipage}[t]{\textwidth - 5.5mm}

\vspace{-3mm}\textbf{El tiempo desapareció}\vspace{3mm}

Comparemos. Antes teníamos una ecuación en cuatro variables con dos
condiciones iniciales y una condición de contorno. Ahora tenemos,
\textbf{para cada valor de \(s\) por separado}, una ecuación en tres
variables con una sola condición de contorno. El parámetro \(s\) quedó
como un espectador: no se deriva, no se integra, solo aparece en un
coeficiente.

Y esa ecuación tiene nombre.

\end{minipage}%
\end{tcolorbox}

\subsection{El problema de Helmholtz}\label{sec-helmholtz}

\begin{definition}[Problema de
Helmholtz]\protect\hypertarget{def-helmholtz}{}\label{def-helmholtz}

Se llama \textbf{problema de Helmholtz} al de hallar \(\Phi\) en un
dominio \(D\) tal que

\begin{equation}\protect\phantomsection\label{eq-helmholtz}{
\nabla^2\Phi + \lambda\,\Phi = g \quad\text{en } D, \qquad
\frac{\partial\Phi}{\partial n} + \alpha\,\Phi = 0 \quad\text{sobre } S .
}\end{equation}

La constante \(\lambda\) se llama \textbf{número de onda al cuadrado},
\(\lambda = k^2\), y \(g\) es la fuente.

\end{definition}

Toda la sección se reduce a esto: \textbf{transformar en el tiempo
convierte cualquier ecuación de la Tabla~\ref{tbl-terminos} en un
problema de Helmholtz.} Y Helmholtz es la ecuación de un solo modo, de
una sola frecuencia: la de una cavidad resonante, la de un modo guiado
en una fibra, la de un estado estacionario de la ecuación de
Schrödinger.

\begin{tcolorbox}[enhanced jigsaw, arc=.35mm, bottomrule=.15mm, breakable, colback=white, colframe=quarto-callout-tip-color-frame, left=2mm, leftrule=.75mm, opacityback=0, rightrule=.15mm, toprule=.15mm]
\begin{minipage}[t]{5.5mm}
\textcolor{quarto-callout-tip-color}{\faLightbulb}
\end{minipage}%
\begin{minipage}[t]{\textwidth - 5.5mm}

\vspace{-3mm}\textbf{Dónde ya vimos esto}\vspace{3mm}

Si sacamos las dependencias en \(y\) y \(z\), \ref{eq-helmholtz} se
vuelve
\[\Phi'' + \lambda\,\Phi = g, \qquad \text{con condiciones de contorno homogéneas},\]
que es exactamente el problema de Sturm--Liouville de la
Capítulo~\ref{sec-sturm} con su función de Green asociada de la
Capítulo~\ref{sec-green}. Todo lo que sabemos de allá ---que hay una
sucesión de autovalores \(\lambda_n\), que las autofunciones son
ortogonales y completas--- vale acá, con la diferencia de que ahora los
autovalores están indexados por varios enteros en vez de uno.

Los \(\lambda_n\) son las \textbf{frecuencias de resonancia} del
dominio, y las autofunciones son sus \textbf{modos normales}. Que un
tambor de una forma dada suene como suena, y no de otra manera, es la
lista de autovalores de su problema de Helmholtz.

\end{minipage}%
\end{tcolorbox}

\begin{exercise}[]\protect\hypertarget{exr-helmholtz-modos}{}\label{exr-helmholtz-modos}

~

\begin{enumerate}
\def\labelenumi{\alph{enumi}.}
\tightlist
\item
  Buscar soluciones de la ecuación de onda de la forma
  \(u(\mathbf{x},t) = v(\mathbf{x})\,e^{-i\omega t}\) y mostrar que
  \(v\) satisface un problema de Helmholtz. ¿Cuánto vale \(\lambda\)?
\item
  Resolverlo para una cuerda de longitud \(L\) con extremos fijos:
  hallar los \(\lambda_n\) y las frecuencias \(\omega_n\).
\item
  Hacer lo mismo para una membrana rectangular de lados \(a\) y \(b\),
  fija en todo el borde. Comparar las frecuencias con las de la cuerda.
  ¿Por qué una guitarra tiene una nota bien definida y un tambor no?
\end{enumerate}

\end{exercise}

\begin{tcolorbox}[enhanced jigsaw, arc=.35mm, bottomrule=.15mm, bottomtitle=1mm, breakable, colback=white, colbacktitle=quarto-callout-tip-color!10!white, colframe=quarto-callout-tip-color-frame, coltitle=black, left=2mm, leftrule=.75mm, opacityback=0, opacitybacktitle=0.6, rightrule=.15mm, title=\textcolor{quarto-callout-tip-color}{\faLightbulb}\hspace{0.5em}{Solución}, titlerule=0mm, toprule=.15mm, toptitle=1mm]

\textbf{a.} Reemplazando \(u = v\,e^{-i\omega t}\) en
\(u_{tt} = c^2\nabla^2u\), la exponencial sale de factor y
\(\partial_t^2\) aporta \(-\omega^2\): \[
-\omega^2 v = c^2\nabla^2 v \qquad\Longrightarrow\qquad
\nabla^2 v + \frac{\omega^2}{c^2}\,v = 0 ,
\] que es \ref{eq-helmholtz} con \(g=0\) y
\[\lambda = \frac{\omega^2}{c^2} = k^2 .\] Las condiciones de contorno
de \(u\) pasan tal cual a \(v\), porque no dependen del tiempo.

\textbf{b.} Para la cuerda, \(v'' + \lambda v = 0\) con \(v(0)=v(L)=0\).
Es el problema de Sturm--Liouville de la Capítulo~\ref{sec-sturm}: hay
solución no trivial solo para
\[\lambda_n = \left(\frac{n\pi}{L}\right)^2, \qquad v_n(x) = \operatorname{sen}\frac{n\pi x}{L},
\qquad n = 1,2,3,\dots\] y las frecuencias son
\[\omega_n = c\sqrt{\lambda_n} = \frac{n\pi c}{L} .\] \textbf{Son
múltiplos enteros de la fundamental.}

\textbf{c.} En la membrana, separando en \(x\) e \(y\), las
autofunciones son
\[v_{mn}(x,y) = \operatorname{sen}\frac{m\pi x}{a}\,\operatorname{sen}\frac{n\pi y}{b},
\qquad
\lambda_{mn} = \pi^2\left(\frac{m^2}{a^2} + \frac{n^2}{b^2}\right),\] y
por lo tanto
\[\omega_{mn} = \pi c\,\sqrt{\frac{m^2}{a^2} + \frac{n^2}{b^2}} .\]

\textbf{Y acá está la diferencia.} En la cuerda las frecuencias son
\(\omega_1, 2\omega_1, 3\omega_1,
\dots\): una serie armónica. El oído junta esa serie y la escucha como
\textbf{una sola nota} de altura \(\omega_1\). En la membrana, en
cambio, esa raíz cuadrada no da múltiplos enteros de nada ---con
\(a=b\), los primeros valores de \(\sqrt{m^2+n^2}\) son
\(\sqrt2, \sqrt5, \sqrt8, \sqrt{10},\dots\)---, así que el sonido no
tiene una serie armónica y el oído no le asigna una altura definida.

Por eso una guitarra toca melodías y un tambor marca el ritmo. La
diferencia no está en el material ni en cómo se los toca: está en la
lista de autovalores del problema de Helmholtz de su dominio. Y por eso
también los timbales, que sí tienen altura, se construyen con una
membrana acoplada a una caja de aire que corrige el espectro.

\end{tcolorbox}

\subsection{Un ejemplo completo: la cuerda cargada de
golpe}\label{sec-ejemplo-transformada}

\begin{example}[Una cuerda a la que se le aplica una fuerza de
repente]\protect\hypertarget{exm-cuerda-forzada}{}\label{exm-cuerda-forzada}

Una cuerda de longitud \(1\), fija en los dos extremos, está quieta
hasta \(t=0\). En ese instante se le aplica una fuerza distribuida
\(F_0\operatorname{sen}(\pi x)\) ---máxima en el medio, nula en los
extremos--- y se la deja así. ¿Qué hace?

El problema es \[
u_{tt} - c^2u_{xx} = F_0\operatorname{sen}(\pi x), \qquad
u(0,t)=u(1,t)=0, \qquad u(x,0)=u_t(x,0)=0 .
\]

\textbf{Paso 1: transformar.} Usamos la transformada de Laplace en el
tiempo, \(U(x,p) = \int_0^\infty u\,e^{-pt}dt\), que para un problema
que arranca en \(t=0\) es la más cómoda. Como las dos condiciones
iniciales son nulas, \(\mathcal{L}\{u_{tt}\} = p^2U\), y
\(\mathcal{L}\{F_0\operatorname{sen}\pi x\} = F_0\operatorname{sen}(\pi x)/p\)
porque la fuerza es constante en el tiempo desde \(t=0\). Queda \[
p^2 U - c^2 U_{xx} = \frac{F_0\operatorname{sen}(\pi x)}{p} ,
\qquad U(0,p)=U(1,p)=0 .
\] Una ecuación \textbf{ordinaria} en \(x\), con \(p\) de parámetro.

\textbf{Paso 2: resolverla.} Como el lado derecho es proporcional a
\(\operatorname{sen}(\pi x)\) y esa función ya cumple las condiciones de
contorno, probamos \(U = K(p)\operatorname{sen}(\pi x)\): \[
\big(p^2 + c^2\pi^2\big)\,K = \frac{F_0}{p}
\qquad\Longrightarrow\qquad
U(x,p) = \frac{F_0\operatorname{sen}(\pi x)}{p\,\big(p^2+\pi^2c^2\big)} .
\]

\textbf{Paso 3: antitransformar.} De las tablas de la
Capítulo~\ref{sec-laplace},
\[\mathcal{L}^{-1}\left\{\frac{1}{p(p^2+\omega^2)}\right\} = \frac{1-\cos\omega t}{\omega^2},\]
y con \(\omega = \pi c\) resulta \[
\boxed{\;u(x,t) = \frac{F_0}{\pi^2c^2}\,\operatorname{sen}(\pi x)\,\big(1 - \cos \pi c\,t\big).\;}
\]

\end{example}

\begin{tcolorbox}[enhanced jigsaw, arc=.35mm, bottomrule=.15mm, breakable, colback=white, colframe=quarto-callout-important-color-frame, left=2mm, leftrule=.75mm, opacityback=0, rightrule=.15mm, toprule=.15mm]
\begin{minipage}[t]{5.5mm}
\textcolor{quarto-callout-important-color}{\faExclamation}
\end{minipage}%
\begin{minipage}[t]{\textwidth - 5.5mm}

\vspace{-3mm}\textbf{Lo que dice la respuesta, y por qué a un ingeniero le importa}\vspace{3mm}

Separemos los dos factores.

\[
\underbrace{\frac{F_0}{\pi^2c^2}\operatorname{sen}(\pi x)}_{\text{la deformación estática}}
\;\times\;
\underbrace{\big(1-\cos\pi c\,t\big)}_{\text{oscila entre 0 y 2}}
\]

El primero es exactamente lo que valdría \(u\) si la fuerza se aplicara
\textbf{muy despacio}: la solución de
\(-c^2u_{xx} = F_0\operatorname{sen}\pi x\), o sea el equilibrio. El
segundo oscila entre \(0\) y \(2\) con la frecuencia fundamental
\(\omega = \pi c\) de la cuerda.

Conclusión: la cuerda no se queda en la posición de equilibrio.
\textbf{Oscila alrededor de ella y llega a deformarse el doble.} Aplicar
una carga de golpe produce el doble de deformación que aplicarla
lentamente, y ese factor \(2\) es un número que se usa todos los días en
ingeniería estructural con el nombre de \emph{factor de carga dinámica}.
Se paga por apurarse.

Y notá que el sistema no se detiene nunca, porque no pusimos ningún
término disipativo. Con un \(b\,u_t\) en la ecuación, el coseno se
apagaría y la cuerda terminaría en la posición estática, que es lo que
pasa de verdad.

\end{minipage}%
\end{tcolorbox}

\emph{(En la versión web hay acá una animación de la cuerda oscilando
entre la posición de reposo y el doble de la deformación estática.)}

Sobre una estructura elástica sin rozamiento se aplica una carga
\textbf{de golpe}. ¿Cuántas veces la deformación estática llega a valer
la deformación máxima? Escribí el número.

Dos. El factor temporal de la solución del
Ejemplo~\ref{exm-cuerda-forzada} es \(1-\cos\omega t\), que oscila entre
\(0\) y \(2\): la estructura pasa de largo la posición de equilibrio y
llega al doble antes de volver.

Se lo llama \textbf{factor de carga dinámica} y es un número que se usa
a diario: una grúa que suelta su carga de golpe, un ascensor que
arranca, un puente por el que pasa un camión. Diseñar para la carga
estática y no para el doble es un error caro.

Con rozamiento el coseno se apaga y la estructura termina en la posición
estática, pero el sobrepaso inicial ocurre igual. El factor \(2\) es el
peor caso, y por eso es el que se usa.

\begin{exercise}[]\protect\hypertarget{exr-cuerda-otra-fuerza}{}\label{exr-cuerda-otra-fuerza}

Sobre la misma cuerda del Ejemplo~\ref{exm-cuerda-forzada}:

\begin{enumerate}
\def\labelenumi{\alph{enumi}.}
\tightlist
\item
  ¿Qué pasa si la fuerza aplicada de golpe es
  \(F_0\operatorname{sen}(2\pi x)\) en lugar de
  \(F_0\operatorname{sen}(\pi x)\)? Comparar la deformación máxima y la
  frecuencia con las del ejemplo.
\item
  ¿Y si la fuerza es \(F_0\operatorname{sen}(\pi x)\cos(\Omega t)\), es
  decir si en vez de aplicarla de golpe se la hace oscilar a frecuencia
  \(\Omega\)? Resolver y estudiar qué pasa cuando \(\Omega\) se acerca a
  \(\pi c\).
\end{enumerate}

\end{exercise}

\begin{tcolorbox}[enhanced jigsaw, arc=.35mm, bottomrule=.15mm, bottomtitle=1mm, breakable, colback=white, colbacktitle=quarto-callout-tip-color!10!white, colframe=quarto-callout-tip-color-frame, coltitle=black, left=2mm, leftrule=.75mm, opacityback=0, opacitybacktitle=0.6, rightrule=.15mm, title=\textcolor{quarto-callout-tip-color}{\faLightbulb}\hspace{0.5em}{Solución}, titlerule=0mm, toprule=.15mm, toptitle=1mm]

\textbf{a.} Todo igual, cambiando \(\pi\) por \(2\pi\). Probando
\(U = K\operatorname{sen}(2\pi x)\) queda \((p^2+4\pi^2c^2)K = F_0/p\),
y antitransformando \[
u(x,t) = \frac{F_0}{4\pi^2c^2}\operatorname{sen}(2\pi x)\,\big(1-\cos 2\pi c\,t\big).
\] Dos diferencias: la deformación estática es \textbf{cuatro veces
menor} ---porque va con \(1/\lambda^2\)--- y la oscilación es
\textbf{dos veces más rápida}. El factor \(2\) del sobrepaso, en cambio,
no cambió: es el mismo para cualquier modo.

Es el segundo armónico de la cuerda: dos jorobas de signos opuestos con
un nodo en el medio.

\textbf{b.} Ahora \(\mathcal{L}\{\cos\Omega t\} = p/(p^2+\Omega^2)\),
así que con \(\omega = \pi c\) \[
\big(p^2+\omega^2\big)K = \frac{F_0\,p}{p^2+\Omega^2}
\qquad\Longrightarrow\qquad
U = \frac{F_0\,p\,\operatorname{sen}(\pi x)}{\big(p^2+\Omega^2\big)\big(p^2+\omega^2\big)} .
\] La antitransformada de \(p/[(p^2+\Omega^2)(p^2+\omega^2)]\) es
\((\cos\Omega t - \cos\omega t)/(\omega^2-\Omega^2)\), de modo que \[
u(x,t) = F_0\operatorname{sen}(\pi x)\,\frac{\cos\Omega t - \cos\omega t}{\omega^2-\Omega^2} .
\]

\textbf{Cuando \(\Omega \to \omega\)} el denominador se va a cero, pero
el numerador también, y el límite existe. Tomándolo con L'Hôpital
respecto de \(\Omega\): \[
u(x,t) \;\longrightarrow\; F_0\operatorname{sen}(\pi x)\,\frac{t\,\operatorname{sen}\omega t}{2\omega} .
\] La amplitud \textbf{crece linealmente con el tiempo}: eso es la
\textbf{resonancia}. La cuerda absorbe energía en cada ciclo porque la
fuerza está siempre en fase con su movimiento, y sin disipación no hay
nada que la detenga.

Y fijate cómo aparece en la fórmula: \(\omega = \pi c\) es el primer
autovalor del problema de Helmholtz asociado, o sea la primera
frecuencia de resonancia del dominio. \textbf{Forzar a un autovalor
rompe cosas}, y por eso el primer paso al diseñar cualquier estructura
es calcular sus modos y asegurarse de que ninguna excitación esperable
caiga cerca.

\end{tcolorbox}

\subsection{Qué transformada usar}\label{sec-que-transformada}

La elección no es libre: depende del dominio de la variable que uno
quiere sacarse de encima y de qué datos hay sobre ella.

\begin{longtable}[]{@{}
  >{\raggedright\arraybackslash}p{(\linewidth - 4\tabcolsep) * \real{0.2400}}
  >{\raggedright\arraybackslash}p{(\linewidth - 4\tabcolsep) * \real{0.3400}}
  >{\raggedright\arraybackslash}p{(\linewidth - 4\tabcolsep) * \real{0.4200}}@{}}
\caption{Qué transformada va con cada situación. La última fila es la
más importante de recordar.}\label{tbl-transformadas}\tabularnewline
\toprule\noalign{}
\begin{minipage}[b]{\linewidth}\raggedright
Variable en\ldots{}
\end{minipage} & \begin{minipage}[b]{\linewidth}\raggedright
Qué se conoce
\end{minipage} & \begin{minipage}[b]{\linewidth}\raggedright
Transformada
\end{minipage} \\
\midrule\noalign{}
\endfirsthead
\toprule\noalign{}
\begin{minipage}[b]{\linewidth}\raggedright
Variable en\ldots{}
\end{minipage} & \begin{minipage}[b]{\linewidth}\raggedright
Qué se conoce
\end{minipage} & \begin{minipage}[b]{\linewidth}\raggedright
Transformada
\end{minipage} \\
\midrule\noalign{}
\endhead
\bottomrule\noalign{}
\endlastfoot
\(t \in [0,\infty)\) & condiciones iniciales en \(t=0\) &
\textbf{Laplace} \\
\(x \in (-\infty,\infty)\) & que \(u\to 0\) en \(\pm\infty\) &
\textbf{Fourier} \\
\(x \in [0,\infty)\) & \(u(0)\) dado (Dirichlet) & Fourier \textbf{en
senos} \\
\(x \in [0,\infty)\) & \(u_x(0)\) dado (Neumann) & Fourier \textbf{en
cosenos} \\
\(x \in [0,L]\) acotado & contorno en los dos extremos & \emph{no
conviene}: usar separación de variables
(sección~\ref{sec-separacion}) \\
\end{longtable}

\begin{tcolorbox}[enhanced jigsaw, arc=.35mm, bottomrule=.15mm, breakable, colback=white, colframe=quarto-callout-warning-color-frame, left=2mm, leftrule=.75mm, opacityback=0, rightrule=.15mm, toprule=.15mm]
\begin{minipage}[t]{5.5mm}
\textcolor{quarto-callout-warning-color}{\faExclamationTriangle}
\end{minipage}%
\begin{minipage}[t]{\textwidth - 5.5mm}

\vspace{-3mm}\textbf{Los tres requisitos, y qué pasa si fallan}\vspace{3mm}

Para que el método funcione hacen falta tres cosas, y conviene
chequearlas antes de empezar:

\begin{enumerate}
\def\labelenumi{\arabic{enumi}.}
\tightlist
\item
  \textbf{La ecuación tiene que ser lineal.} Si no, no hay transformada
  que sirva: la transformada de un producto no es el producto de las
  transformadas.
\item
  \textbf{Los coeficientes no pueden depender de la variable que se
  transforma.} Si \(c = c(t)\), al transformar aparece una convolución y
  el problema no se simplifica.
\item
  \textbf{Los coeficientes de las condiciones de contorno tampoco.} El
  \emph{dato} del borde sí puede depender del tiempo ---se transforma
  junto con todo lo demás y listo---; lo que no puede es que cambie el
  \(\alpha\) de \ref{eq-robin-t}, y mucho menos que \textbf{se mueva la
  frontera}, porque entonces el dominio espacial dependería del tiempo
  que acabamos de eliminar.
\end{enumerate}

Cuando alguno falla, todavía quedan la función de Green y la separación
de variables, que son las dos secciones que siguen.

\end{minipage}%
\end{tcolorbox}

Querés resolver la ecuación del calor en una barra
\textbf{semi-infinita} \(x\ge 0\), con la temperatura del extremo
\(x=0\) fijada. ¿Qué transformada conviene aplicar en \(x\)?

\begin{itemize}
\tightlist
\item
  {Fourier en senos}
\item
  {Fourier en cosenos}
\item
  {Fourier completa}
\item
  {Laplace en \(x\)}
\end{itemize}

La transformada en senos es la que «se come» una condición de Dirichlet:
al integrar por partes dos veces, el término de borde que sobrevive
involucra \(u(0)\), que es justamente el dato que tenemos.

La de cosenos hace lo contrario ---necesita \(u_x(0)\), o sea una
condición de Neumann--- y por eso sería la elección si el extremo
estuviera aislado. La completa pide que la variable recorra toda la
recta, y acá va de \(0\) a \(\infty\). Y Laplace en \(x\) se puede
intentar, pero pediría conocer \(u_x(0)\) \textbf{además} de \(u(0)\),
es decir un dato de Cauchy, que no tenemos y que además, por la
sección~\ref{sec-hadamard}, sería mala idea pedir.

Está en la Tabla~\ref{tbl-transformadas}: la transformada se elige por
el dominio de la variable y por qué dato hay en el borde.

\begin{exercise}[]\protect\hypertarget{exr-transformada-requisitos}{}\label{exr-transformada-requisitos}

Para cada situación, decidir si el método de las transformadas
integrales en el tiempo se puede aplicar, y por qué.

\begin{enumerate}
\def\labelenumi{\alph{enumi}.}
\tightlist
\item
  Una cuerda con un extremo fijo y el otro oscilando: \(u(0,t)=0\),
  \(u(L,t)=A\operatorname{sen}\Omega t\).
\item
  Una cuerda cuya tensión se va aumentando mientras vibra, de modo que
  \(c = c(t)\).
\item
  Una barra que se alarga: uno de sus extremos está en \(x = L(t)\).
\item
  La ecuación de Burgers \(u_t + u\,u_x = 0\).
\item
  Difusión en un medio no homogéneo, \(u_t = \big(k(x)\,u_x\big)_x\).
\end{enumerate}

\end{exercise}

\begin{tcolorbox}[enhanced jigsaw, arc=.35mm, bottomrule=.15mm, bottomtitle=1mm, breakable, colback=white, colbacktitle=quarto-callout-tip-color!10!white, colframe=quarto-callout-tip-color-frame, coltitle=black, left=2mm, leftrule=.75mm, opacityback=0, opacitybacktitle=0.6, rightrule=.15mm, title=\textcolor{quarto-callout-tip-color}{\faLightbulb}\hspace{0.5em}{Solución}, titlerule=0mm, toprule=.15mm, toptitle=1mm]

\textbf{a. Sí.} Que el \textbf{dato} dependa del tiempo no molesta: se
transforma junto con todo lo demás y queda
\(U(L,p) = A\,\Omega/(p^2+\Omega^2)\), que sigue siendo una condición de
contorno perfecta para el problema transformado. Lo que no puede
depender del tiempo son los \textbf{coeficientes} de la condición, es
decir el \(\alpha\) de \ref{eq-robin-t}.

\textbf{b. No.} Acá falla el requisito de que los coeficientes no
dependan de la variable que se transforma. Al transformar, el producto
\(c(t)^2u_{xx}\) se convierte en una \textbf{convolución} en \(p\), y el
problema queda igual de difícil o peor.

\textbf{c.~No, y es el caso más feo.} Con la frontera en movimiento, el
dominio espacial cambia con el tiempo, así que ni siquiera se puede
decir «para cada \(p\), una ecuación en \(x\) en tal dominio»: el
dominio depende del tiempo que acabamos de eliminar. Los problemas de
frontera móvil ---el hielo que se derrite, el frente de una llama--- son
un capítulo aparte y difícil.

\textbf{d.~No.} No es lineal, y la transformada de un producto no es el
producto de las transformadas. Para Burgers hay que volver a las
características de la sección~\ref{sec-casi-lineales}.

\textbf{e. Sí.} \(k\) depende de \(x\), no de \(t\), así que se
transforma sin problema y queda
\[p\,U - u_0(x) = \big(k(x)\,U_x\big)_x ,\] una ecuación
\textbf{ordinaria} en \(x\), lineal, con coeficientes variables. Puede
ser difícil de resolver ---quizás no tenga solución en términos de
funciones elementales--- pero el método se aplica y el problema bajó una
dimensión, que es todo lo que le pedimos.

\end{tcolorbox}

\section{Función de Green en varias variables}\label{sec-green-varias}

El método de las transformadas nos dejó con un problema de Helmholtz,
que sigue siendo una EDP. Este método ataca ese problema ---y muchos
otros--- con una idea de una sola línea:

\begin{quote}
\textbf{Resolver una vez para una fuente puntual, y después sumar.}
\end{quote}

Es la superposición de la sección~\ref{sec-lineales} llevada al extremo.
Si el sistema es lineal, y si sabemos qué produce una fuente concentrada
en un punto, entonces sabemos qué produce cualquier fuente: la
descomponemos en puntos, sumamos las respuestas, y listo.

\begin{tcolorbox}[enhanced jigsaw, arc=.35mm, bottomrule=.15mm, breakable, colback=white, colframe=quarto-callout-important-color-frame, left=2mm, leftrule=.75mm, opacityback=0, rightrule=.15mm, toprule=.15mm]
\begin{minipage}[t]{5.5mm}
\textcolor{quarto-callout-important-color}{\faExclamation}
\end{minipage}%
\begin{minipage}[t]{\textwidth - 5.5mm}

\vspace{-3mm}\textbf{Esto ya lo sabés de electrostática}\vspace{3mm}

Un físico usa funciones de Green desde el primer curso de
electromagnetismo, aunque no las llame así.

El potencial de \textbf{una carga puntual} \(q\) en el origen es
\(\varphi = q/(4\pi\varepsilon_0 r)\). El potencial de \textbf{una
distribución cualquiera} de carga \(\rho(\mathbf{x}')\) es \[
\varphi(\mathbf{x}) = \frac{1}{4\pi\varepsilon_0}
\int \frac{\rho(\mathbf{x}')}{|\mathbf{x}-\mathbf{x}'|}\,d^3x' .
\] Nadie resuelve esa integral resolviendo una EDP: se la escribe
directamente, \textbf{sumando el aporte de cada pedacito de carga}. Y
funciona porque la ecuación de Poisson,
\(\nabla^2\varphi = -\rho/\varepsilon_0\), es lineal.

Pues bien: \(1/(4\pi r)\) \textbf{es} la función de Green de la ecuación
de Poisson en el espacio libre, y la integral de Coulomb \textbf{es} la
fórmula general de este método. Todo lo que sigue es darle nombre a eso
y extenderlo a otras ecuaciones y a dominios con contorno.

\end{minipage}%
\end{tcolorbox}

\subsection{Fuente adentro, contorno en cero}\label{sec-green-caso-1}

Empecemos por el caso más parecido al de electrostática. Sea \(L\) un
operador diferencial lineal ---por ejemplo
\(L[u] = \nabla^2u + \lambda u\), que es el que nos dejó la
sección~\ref{sec-transformadas}--- y consideremos el problema

\begin{equation}\protect\phantomsection\label{eq-problema-green-1}{
L[u(\mathbf{x})] = f(\mathbf{x}) \quad \text{en } D, \qquad
B[u(\mathbf{x})] = 0 \quad \text{en } S ,
}\end{equation}

donde \(B\) es la condición de contorno, \textbf{homogénea}: puede ser
\(B[u]=u\) (Dirichlet), \(B[u] = \partial u/\partial n\) (Neumann) o la
combinación de Robin \(B[u] = \alpha u + \nabla u\cdot\hat n\).

\begin{definition}[Función de
Green]\protect\hypertarget{def-green-varias}{}\label{def-green-varias}

La \textbf{función de Green} del problema \ref{eq-problema-green-1} es
la función \(G(\mathbf{x},\xi)\) que satisface, para cada \(\xi \in D\)
fijo,

\begin{equation}\protect\phantomsection\label{eq-def-green}{
L\big[G(\mathbf{x},\xi)\big] = \delta(\mathbf{x}-\xi) \quad\text{en } D, \qquad
B\big[G(\mathbf{x},\xi)\big] = 0 \quad\text{en } S .
}\end{equation}

Es decir: \textbf{\(G(\cdot,\xi)\) es la respuesta del sistema a una
fuente puntual unitaria colocada en \(\xi\)}, con el mismo contorno que
el problema original.

\end{definition}

\begin{theorem}[La solución, por
superposición]\protect\hypertarget{thm-green-varias}{}\label{thm-green-varias}

Si \(G\) satisface \ref{eq-def-green}, entonces

\begin{equation}\protect\phantomsection\label{eq-solucion-green}{
u(\mathbf{x}) = \int_D G(\mathbf{x},\xi)\,f(\xi)\,d\xi
}\end{equation}

resuelve el problema \ref{eq-problema-green-1}.

\end{theorem}

La verificación es corta y vale la pena verla, porque muestra
exactamente dónde se usa cada cosa. Como \(L\) actúa sobre
\(\mathbf{x}\) y la integral es en \(\xi\), se los puede intercambiar:

\[
L[u] = \int_D L\big[G(\mathbf{x},\xi)\big]\,f(\xi)\,d\xi
= \int_D \delta(\mathbf{x}-\xi)\,f(\xi)\,d\xi = f(\mathbf{x}) \ \checkmark
\]

usando la propiedad de la delta de la Capítulo~\ref{sec-distribuciones},
y del mismo modo

\[
B[u] = \int_D B\big[G(\mathbf{x},\xi)\big]\,f(\xi)\,d\xi = 0 \ \checkmark .
\]

\begin{tcolorbox}[enhanced jigsaw, arc=.35mm, bottomrule=.15mm, breakable, colback=white, colframe=quarto-callout-tip-color-frame, left=2mm, leftrule=.75mm, opacityback=0, rightrule=.15mm, toprule=.15mm]
\begin{minipage}[t]{5.5mm}
\textcolor{quarto-callout-tip-color}{\faLightbulb}
\end{minipage}%
\begin{minipage}[t]{\textwidth - 5.5mm}

\vspace{-3mm}\textbf{Lo mismo, en el idioma de sistemas}\vspace{3mm}

Si a \(f\) se la llama \emph{entrada} y a \(u\) \emph{salida}, entonces
\(G\) es la \textbf{respuesta al impulso} del sistema y
\ref{eq-solucion-green} es una \textbf{convolución}. Es literalmente el
mismo objeto que la respuesta al impulso de un circuito o de un filtro.

Y se conecta con la sección~\ref{sec-transformadas} de la manera
esperable: en el espacio libre, donde no hay contorno que rompa la
simetría, \(G\) depende solo de \(\mathbf{x}-\xi\), la integral es una
convolución de verdad, y transformando Fourier se vuelve un producto.
Los dos métodos son el mismo visto de dos lados.

\end{minipage}%
\end{tcolorbox}

\begin{example}[Las funciones de Green del espacio
libre]\protect\hypertarget{exm-green-libres}{}\label{exm-green-libres}

Sin contorno ---o con el contorno mandado al infinito--- las funciones
de Green se calculan de una vez y para siempre. En tres dimensiones, con
\(r = |\mathbf{x}-\xi|\):

{\def\LTcaptype{none} % do not increment counter
\begin{longtable}[]{@{}
  >{\raggedright\arraybackslash}p{(\linewidth - 4\tabcolsep) * \real{0.3333}}
  >{\raggedright\arraybackslash}p{(\linewidth - 4\tabcolsep) * \real{0.3333}}
  >{\raggedright\arraybackslash}p{(\linewidth - 4\tabcolsep) * \real{0.3333}}@{}}
\toprule\noalign{}
\begin{minipage}[b]{\linewidth}\raggedright
Operador
\end{minipage} & \begin{minipage}[b]{\linewidth}\raggedright
Ecuación
\end{minipage} & \begin{minipage}[b]{\linewidth}\raggedright
Función de Green
\end{minipage} \\
\midrule\noalign{}
\endhead
\bottomrule\noalign{}
\endlastfoot
Laplace & \(\nabla^2 G = \delta\) &
\(\displaystyle G = -\frac{1}{4\pi r}\) \\
Helmholtz & \(\nabla^2 G + k^2G = \delta\) &
\(\displaystyle G = -\frac{e^{ikr}}{4\pi r}\) \\
\end{longtable}
}

La primera es la ley de Coulomb: reemplazándola en
\ref{eq-solucion-green} con \(f = -\rho/\varepsilon_0\) sale exactamente
la integral del recuadro de más arriba.

La segunda es una \textbf{onda esférica saliente}: una fuente puntual
que oscila a frecuencia \(\omega = kc\) emite ondas que se alejan y cuya
amplitud decae como \(1/r\) ---de ahí que la intensidad, que va como la
amplitud al cuadrado, decaiga como \(1/r^2\)---. Si se toma la parte
real, \(-\cos(kr)/(4\pi r)\), se obtiene la solución estacionaria, que
es la superposición de una onda saliente y una entrante.

En cuanto hay contorno esto deja de alcanzar: hay que agregarle a \(G\)
una función sin singularidades que arregle la condición de borde. Para
dominios simples eso se hace con el \textbf{método de las imágenes}, que
es una carga ficticia colocada del otro lado del borde, y que se trata
en el capítulo 11.

\end{example}

¿Por qué se puede escribir el potencial de una distribución de carga
como una integral, sin resolver ninguna ecuación diferencial?

\begin{itemize}
\tightlist
\item
  {Porque la ecuación de Poisson es lineal y conocemos la respuesta a
  una carga puntual}
\item
  {Porque el potencial es un escalar y no un vector}
\item
  {Porque la carga eléctrica se conserva}
\item
  {Porque el campo electrostático es conservativo}
\end{itemize}

Por la superposición, que es la linealidad. La distribución \(\rho\) se
piensa como una suma de cargas puntuales, la respuesta a cada una se
conoce ---es \(1/(4\pi\varepsilon_0 r)\)--- y como la ecuación es
lineal, la respuesta total es la suma. La integral de Coulomb es
exactamente \ref{eq-solucion-green}.

Las otras tres afirmaciones son verdaderas y no tienen nada que ver. La
prueba está en que el mismo argumento funciona para la deformación de
una cuerda, que no es electrostática ni es un escalar conservado: solo
hace falta que la ecuación sea lineal.

Y al revés: para un medio no lineal ---un dieléctrico saturable, por
ejemplo--- la integral de Coulomb \textbf{deja de valer}, aunque la
carga se siga conservando.

\subsection{Fuente en el borde, ecuación
homogénea}\label{sec-green-caso-2}

Estamos acostumbrados a pensar en fuentes adentro del dominio. Pero
muchos problemas físicos son al revés: la ecuación no tiene fuente y lo
que impone las cosas es el \textbf{borde}. Una placa sin cargas cuyos
bordes están a potenciales dados; una barra sin calentadores internos
con los extremos a temperaturas fijas.

\begin{equation}\protect\phantomsection\label{eq-problema-green-2}{
L[u(\mathbf{x})] = 0 \quad\text{en } D, \qquad
B[u(\mathbf{x})] = h(\mathbf{x}) \quad\text{en } S .
}\end{equation}

El truco es exactamente el mismo, corriendo la delta de lugar:
\textbf{ahora la fuente puntual se pone sobre el borde}.

\begin{definition}[Función de Green
modificada]\protect\hypertarget{def-green-modificada}{}\label{def-green-modificada}

Se busca \(G(\mathbf{x},\xi)\) tal que

\begin{equation}\protect\phantomsection\label{eq-def-green-2}{
L\big[G(\mathbf{x},\xi)\big] = 0 \quad\text{en } D, \qquad
B\big[G(\mathbf{x},\xi)\big] = \delta(\mathbf{x}-\xi) \quad\text{en } S ,
}\end{equation}

con \(\xi \in S\). Entonces la solución de \ref{eq-problema-green-2} es

\begin{equation}\protect\phantomsection\label{eq-solucion-green-2}{
u(\mathbf{x}) = \int_S G(\mathbf{x},\xi)\,h(\xi)\,d\xi ,
}\end{equation}

donde la integral se hace \textbf{sobre la superficie}, no sobre el
volumen.

\end{definition}

La verificación es la misma cuenta de antes con los papeles cambiados:

\[
L[u] = \int_S L\big[G\big]\,h\,d\xi = 0 \ \checkmark, \qquad
B[u] = \int_S \delta(\mathbf{x}-\xi)\,h(\xi)\,d\xi = h(\mathbf{x}) \ \checkmark .
\]

\begin{tcolorbox}[enhanced jigsaw, arc=.35mm, bottomrule=.15mm, breakable, colback=white, colframe=quarto-callout-tip-color-frame, left=2mm, leftrule=.75mm, opacityback=0, rightrule=.15mm, toprule=.15mm]
\begin{minipage}[t]{5.5mm}
\textcolor{quarto-callout-tip-color}{\faLightbulb}
\end{minipage}%
\begin{minipage}[t]{\textwidth - 5.5mm}

\vspace{-3mm}\textbf{Y si hay las dos cosas}\vspace{3mm}

Un problema con fuente \textbf{y} con contorno inhomogéneo se parte en
dos por linealidad, tal como en la sección~\ref{sec-lineales}: \[
u = \underbrace{\int_D G_1(\mathbf{x},\xi)\,f(\xi)\,d\xi}_{\text{la fuente, con borde nulo}}
\;+\;
\underbrace{\int_S G_2(\mathbf{x},\xi)\,h(\xi)\,d\xi}_{\text{el borde, sin fuente}} .
\] Cada pedazo se resuelve con su función de Green y después se suman.
Es la razón práctica de que valga la pena separar los dos casos.

\end{minipage}%
\end{tcolorbox}

\begin{exercise}[]\protect\hypertarget{exr-green-fisico}{}\label{exr-green-fisico}

Para cada situación, decir qué representa físicamente la función de
Green, con qué experimento se la mediría, y qué papel juegan \(f\) y
\(h\):

\begin{enumerate}
\def\labelenumi{\alph{enumi}.}
\tightlist
\item
  Conducción del calor estacionaria en una placa, con los bordes a
  temperatura fija y una fuente de calor adentro.
\item
  El potencial electrostático dentro de una caja metálica conectada a
  tierra, con una distribución de carga adentro.
\item
  Una membrana de tambor tensada sobre un aro, con una carga distribuida
  encima.
\item
  La misma placa del caso \(a\), pero \textbf{sin} fuente adentro y con
  los bordes a temperaturas distintas.
\end{enumerate}

\end{exercise}

\begin{tcolorbox}[enhanced jigsaw, arc=.35mm, bottomrule=.15mm, bottomtitle=1mm, breakable, colback=white, colbacktitle=quarto-callout-tip-color!10!white, colframe=quarto-callout-tip-color-frame, coltitle=black, left=2mm, leftrule=.75mm, opacityback=0, opacitybacktitle=0.6, rightrule=.15mm, title=\textcolor{quarto-callout-tip-color}{\faLightbulb}\hspace{0.5em}{Solución}, titlerule=0mm, toprule=.15mm, toptitle=1mm]

\textbf{a.} \(G(\mathbf{x},\xi)\) es la temperatura en \(\mathbf{x}\)
producida por una fuente puntual de potencia unitaria en \(\xi\), con
los bordes mantenidos a cero. Se la mediría poniendo una resistencia
chiquita en \(\xi\) y midiendo con un termómetro en \(\mathbf{x}\). La
fuente \(f\) es la distribución de potencia calorífica; \(h=0\) porque
los bordes están a temperatura fija.

\textbf{b.} \(G\) es el potencial producido por una carga puntual
unitaria en \(\xi\) \textbf{dentro de la caja}, que no es
\(1/(4\pi\varepsilon_0 r)\) sino esa función corregida para que se anule
en las paredes. Esa corrección es lo que se calcula con el método de las
imágenes. Acá \(f\) es la densidad de carga.

\textbf{c.} \(G\) es la deformación de la membrana bajo una pesita
unitaria en \(\xi\) --- el análogo en dos dimensiones de la carpa del
Ejemplo~\ref{exm-green-cuerda}, aunque acá ya no es un cono perfecto
porque en dos dimensiones la solución fundamental es logarítmica y no
lineal.

\textbf{d.} Este es el segundo caso, el de la
sección~\ref{sec-green-caso-2}. Ahora la ecuación es homogénea ---no hay
fuente adentro--- y lo que manda es el borde: \(h\) es la temperatura
prescrita sobre el borde y \(G(\mathbf{x},\xi)\), con \(\xi\)
\textbf{sobre el borde}, dice cuánto influye en el punto interior
\(\mathbf{x}\) un pedacito caliente del contorno.

Es la misma placa, la misma ecuación y el mismo dominio, pero con una
función de Green distinta: la del caso 1 se integra sobre el volumen y
la del caso 2 sobre la superficie. Si la placa tuviera fuente \textbf{y}
bordes a distintas temperaturas, se resolverían los dos problemas por
separado y se sumarían.

\end{tcolorbox}

\subsection{El caso más simple, y el más útil de tener a
mano}\label{sec-green-cuerda}

Para ver la estructura sin que la geometría estorbe, hagamos el caso de
una variable, que además es un problema físico que conviene tener en la
cabeza.

\begin{example}[Una cuerda con una carga
puntual]\protect\hypertarget{exm-green-cuerda}{}\label{exm-green-cuerda}

Una cuerda de longitud \(L\), tensada con fuerza \(T\) y fija en los
extremos, soporta una carga distribuida \(f(x)\) por unidad de longitud.
Su deformación satisface

\[
T\,u'' = -f(x), \qquad u(0)=u(L)=0 .
\]

Pongamos primero \textbf{una sola carga puntual unitaria} en \(x=\xi\),
o sea \(f = \delta(x-\xi)\). Fuera del punto de aplicación no hay carga,
así que \(u''=0\) y la cuerda es \textbf{recta} a cada lado: queda un
triángulo con el vértice en \(\xi\). Imponiendo que se anule en los
extremos, que sea continua en \(\xi\) y que la pendiente pegue el salto
que le impone la delta, sale

\begin{equation}\protect\phantomsection\label{eq-green-cuerda}{
G(x,\xi) = \frac{1}{T}
\begin{cases}
\dfrac{x\,(L-\xi)}{L}, & x \le \xi, \\[8pt]
\dfrac{\xi\,(L-x)}{L}, & x \ge \xi .
\end{cases}
}\end{equation}

\textbf{Esa carpa es la función de Green.} No hay que interpretarla: se
la puede hacer con un hilo y una pesita.

Y entonces, para una carga distribuida cualquiera, \[
u(x) = \int_0^L G(x,\xi)\,f(\xi)\,d\xi ,
\] que se lee como: \textbf{la deformación es la suma de las carpas que
produciría cada pedacito de carga por separado.}

\end{example}

\emph{(En la versión web hay acá un gráfico interactivo: a la izquierda
la carpa de una sola carga puntual, cuya posición se mueve con un
slider, y a la derecha la superposición de \(N\) carpas aproximando la
solución bajo una carga distribuida, con \(N\) regulable. Al crecer
\(N\), la suma converge a la solución exacta.)}

En la cuerda del Ejemplo~\ref{exm-green-cuerda}, ¿qué es físicamente
\(G(x,\xi)\)?

\begin{itemize}
\tightlist
\item
  {La forma que toma la cuerda cuando se le cuelga una pesita unitaria
  en \(\xi\)}
\item
  {La fuerza que hay que aplicar en \(\xi\) para deformarla una unidad}
\item
  {La energía elástica almacenada en la cuerda}
\item
  {La solución general de la ecuación}
\end{itemize}

Es la deformación producida por una carga puntual unitaria en \(\xi\),
medida en el punto \(x\). Nada más y nada menos: se la puede armar con
un hilo, dos clavos y una arandela.

Y de ahí sale un hecho que sorprende la primera vez. La fórmula
\ref{eq-green-cuerda} es \textbf{simétrica}: \[G(x,\xi) = G(\xi,x).\]
Traducido: \emph{la deformación en \(x\) producida por una pesita
colgada en \(\xi\) es exactamente igual a la deformación en \(\xi\)
producida por la misma pesita colgada en \(x\).} Se llama
\textbf{reciprocidad de Maxwell--Betti}, vale para cualquier sistema
elástico lineal, y es de esas cosas que uno no creería si no las viera.

No es casualidad: es consecuencia de que el operador \(d^2/dx^2\) con
estas condiciones de contorno sea \textbf{autoadjunto}, que es lo que se
estudió en la Capítulo~\ref{sec-sturm}. Todas las funciones de Green de
operadores autoadjuntos son simétricas.

Una cuerda de longitud \(1\) y tensión \(1\), fija en los extremos,
soporta una carga puntual unitaria en \(\xi = 0{,}25\). ¿Cuánto vale la
deformación \textbf{en el punto de aplicación}?

De \ref{eq-green-cuerda} con \(x=\xi\), \(L=1\) y \(T=1\):
\[G(\xi,\xi) = \xi\,(1-\xi) = 0{,}25 \times 0{,}75 = 0{,}1875 .\]

Notá que la máxima deformación posible se obtiene con la carga
\textbf{en el medio}, donde \(G = 1/4\), y que se anula al acercarse a
cualquiera de los apoyos. Es lo que uno esperaría: una carga cerca del
apoyo la aguanta el apoyo.

\begin{exercise}[]\protect\hypertarget{exr-green-verificar}{}\label{exr-green-verificar}

Verificar que \ref{eq-green-cuerda} es efectivamente la función de Green
del problema, comprobando las cuatro cosas que la definen: que se anula
en los dos extremos, que es continua en \(x=\xi\), que satisface
\(u''=0\) a cada lado, y que la derivada pega en \(\xi\) el salto que le
impone la delta. Después comprobar la simetría \(G(x,\xi)=G(\xi,x)\).

\end{exercise}

\begin{tcolorbox}[enhanced jigsaw, arc=.35mm, bottomrule=.15mm, bottomtitle=1mm, breakable, colback=white, colbacktitle=quarto-callout-tip-color!10!white, colframe=quarto-callout-tip-color-frame, coltitle=black, left=2mm, leftrule=.75mm, opacityback=0, opacitybacktitle=0.6, rightrule=.15mm, title=\textcolor{quarto-callout-tip-color}{\faLightbulb}\hspace{0.5em}{Solución}, titlerule=0mm, toprule=.15mm, toptitle=1mm]

Tomemos \(T=1\) para aligerar.

\textbf{Se anula en los extremos.} En \(x=0\) estamos en la rama
\(x\le\xi\) y vale \(0\cdot(L-\xi)/L = 0\). En \(x=L\), rama
\(x\ge\xi\), vale \(\xi\,(L-L)/L = 0\). ✓

\textbf{Es continua en \(x=\xi\).} Las dos ramas evaluadas ahí dan
\(\xi(L-\xi)/L\). ✓

\textbf{Cumple \(u''=0\) a cada lado.} Las dos ramas son lineales en
\(x\). ✓ Que sea así no es un detalle: es la ecuación \(u''=0\) diciendo
que \textbf{donde no hay carga, la cuerda es recta}.

\textbf{El salto de la derivada.} A la izquierda la pendiente es
\((L-\xi)/L\); a la derecha es \(-\xi/L\). El salto vale \[
u'(\xi^+) - u'(\xi^-) = -\frac{\xi}{L} - \frac{L-\xi}{L} = -1 ,
\] que es justo lo que pide \(u'' = -\delta(x-\xi)\): integrando la
delta a través del punto se obtiene \(-1\). ✓

\textbf{La simetría.} Intercambiando \(x\) y \(\xi\), la rama
\(x\le\xi\) se convierte en la rama \(\xi \le x\) y la expresión
\(x(L-\xi)/L\) pasa a \(\xi(L-x)/L\), que es exactamente la otra rama.
Las dos ramas se intercambian entre sí, así que \(G(x,\xi)=G(\xi,x)\). ✓

\end{tcolorbox}

\begin{exercise}[]\protect\hypertarget{exr-green-uniforme}{}\label{exr-green-uniforme}

Usar la función de Green \ref{eq-green-cuerda} para hallar la
deformación de la cuerda bajo una carga \textbf{uniforme}
\(f(x) = f_0\), y comparar con lo que se obtiene integrando directamente
\(T u'' = -f_0\). ¿Cuánto vale la flecha máxima?

\end{exercise}

\begin{tcolorbox}[enhanced jigsaw, arc=.35mm, bottomrule=.15mm, bottomtitle=1mm, breakable, colback=white, colbacktitle=quarto-callout-tip-color!10!white, colframe=quarto-callout-tip-color-frame, coltitle=black, left=2mm, leftrule=.75mm, opacityback=0, opacitybacktitle=0.6, rightrule=.15mm, title=\textcolor{quarto-callout-tip-color}{\faLightbulb}\hspace{0.5em}{Solución}, titlerule=0mm, toprule=.15mm, toptitle=1mm]

Con \ref{eq-solucion-green}, partiendo la integral en las dos ramas: \[
u(x) = \frac{f_0}{T}\left[\int_0^{x}\frac{\xi(L-x)}{L}\,d\xi + \int_x^{L}\frac{x(L-\xi)}{L}\,d\xi\right]
= \frac{f_0}{T}\left[\frac{x^2(L-x)}{2L} + \frac{x(L-x)^2}{2L}\right] .
\] Sacando factor común \(x(L-x)/(2L)\), el corchete que queda es
\(x + (L-x) = L\), así que \[
\boxed{\;u(x) = \frac{f_0}{2T}\,x\,(L-x).\;}
\]

\textbf{Por el otro camino.} Integrando \(u'' = -f_0/T\) dos veces sale
\(u = -f_0x^2/(2T) + C_1x + C_2\), y las condiciones \(u(0)=u(L)=0\) dan
\(C_2=0\) y \(C_1 = f_0L/(2T)\), o sea lo mismo. ✓

\textbf{La flecha máxima} está en el medio y vale
\[u\left(\frac{L}{2}\right) = \frac{f_0L^2}{8T},\] la fórmula de la
flecha de un cable con carga uniforme, que aparece en cualquier manual
de estructuras. Notá que crece como \(L^2\): duplicar la luz cuadruplica
la flecha. Y que baja con \(T\): por eso se tensan los cables.

La forma es una \textbf{parábola}, no una catenaria. La catenaria
aparece cuando el peso está distribuido a lo largo del \textbf{cable} y
no a lo largo de la horizontal; para cables poco combados las dos son
casi iguales, y de hecho un puente colgante es una parábola porque lo
que pesa es el tablero.

\end{tcolorbox}

\begin{tcolorbox}[enhanced jigsaw, arc=.35mm, bottomrule=.15mm, breakable, colback=white, colframe=quarto-callout-warning-color-frame, left=2mm, leftrule=.75mm, opacityback=0, rightrule=.15mm, toprule=.15mm]
\begin{minipage}[t]{5.5mm}
\textcolor{quarto-callout-warning-color}{\faExclamationTriangle}
\end{minipage}%
\begin{minipage}[t]{\textwidth - 5.5mm}

\vspace{-3mm}\textbf{La recompensa y el precio}\vspace{3mm}

\textbf{La recompensa.} Una vez que se tiene \(G\), se tiene la solución
para \textbf{cualquier} fuente y cualquier dato de borde. No es la
solución de un problema: es la de una familia entera.

\textbf{El precio.} Hallar \(G\) suele ser tan difícil como el problema
original, y depende del dominio: la función de Green de una esfera no
sirve para un cubo. Se la consigue explícitamente en pocos casos ---el
espacio libre, la esfera, el semiespacio, el círculo--- y en el resto
hay que calcularla numéricamente o combinar el método con transformadas.

Por eso conviene pensar en la función de Green menos como una receta de
cálculo y más como una \textbf{manera de organizar el problema}: separa
lo que depende del dominio y del operador, que es \(G\), de lo que
depende del caso particular, que son \(f\) y \(h\).

\end{minipage}%
\end{tcolorbox}

\section{Separación de variables}\label{sec-separacion}

Este es el método más usado de los cuatro y el más limitado. No sirve
siempre: necesita que el dominio y las condiciones tengan la simetría
adecuada. Pero cuando sirve, da la solución completa y, sobre todo,
\textbf{la da escrita en el idioma en que el sistema físico habla de sí
mismo}.

\subsection{El hecho físico: los modos}\label{sec-modos}

Antes de cualquier cuenta, el fenómeno.

Si uno agarra una cuerda tensa y la excita a una frecuencia cualquiera,
no pasa gran cosa. Pero si la excita a ciertas frecuencias particulares,
la cuerda se pone a vibrar con mucha amplitud adoptando \textbf{una
forma fija que solo cambia de tamaño}: hay puntos que se quedan
perfectamente quietos ---los nodos--- y puntos que oscilan al máximo.
Cada una de esas formas es un \textbf{modo normal}, y las frecuencias en
que aparecen son las frecuencias propias del sistema.

Video recomendado: SMU Physics, \emph{Vibrating String -- Normal Modes},
\url{https://www.youtube.com/watch?v=BSIw5SgUirg}.

Lo que hace el método de separación de variables es \textbf{buscar
exactamente esos modos}, y después armar la solución general sumándolos.
La hipótesis de partida ---que la solución se factoriza como una forma
por una función del tiempo--- no es un truco algebraico: es la
definición matemática de «modo normal».

\[
u(x,t) = \underbrace{X(x)}_{\text{la forma}} \;\times\; \underbrace{T(t)}_{\text{cómo late}}
\]

\subsection{El esquema general}\label{sec-esquema-separacion}

Los pasos son siempre los mismos, en cualquier cantidad de variables.

\begin{enumerate}
\def\labelenumi{\arabic{enumi}.}
\tightlist
\item
  \textbf{Sustituir} \(u(x,y,z,\dots) = X(x)\,Y(y)\,Z(z)\cdots\) en la
  ecuación y ver si se puede separar, es decir si se puede llegar a que
  un lado dependa solo de una variable y el otro solo de las demás. Que
  se pueda depende de la ecuación \textbf{y del sistema de coordenadas}.
  Y no siempre el producto es lo adecuado: en la ecuación de
  Hamilton--Jacobi \(\varphi_t + H(\varphi_x,\varphi_y,x,y,t) = 0\), que
  es de primer orden y fuertemente no lineal, lo que conviene proponer
  es una \textbf{suma} de funciones de una variable.
\item
  \textbf{Traducir las condiciones de contorno} de \(u\) en condiciones
  para \(X\), \(Y\), \(Z\). Ojo: la continuidad y la acotación de \(u\)
  pueden agregar condiciones que no estaban escritas, y en coordenadas
  polares o esféricas eso pasa siempre.
\item
  \textbf{Detectar cuáles de los problemas resultantes son de
  Sturm--Liouville} y hallar sus autovalores y autofunciones. Este es el
  paso conceptual: los autovalores son las frecuencias propias y las
  autofunciones son los modos.
\item
  \textbf{Resolver las ecuaciones restantes} con las condiciones de
  contorno homogéneas.
\item
  \textbf{Reconstruir} las autofunciones completas,
  \(\phi_n = X_nY_nZ_n\cdots\), y escribir la solución como combinación
  lineal, \(u = \sum_n a_n\phi_n\). Esto se puede hacer porque la
  ecuación es lineal: es la superposición de la
  sección~\ref{sec-lineales}, otra vez.
\item
  \textbf{Hallar los coeficientes} \(a_n\) imponiendo las condiciones
  que quedaron ---típicamente las iniciales--- y usando la
  \textbf{ortogonalidad y la completitud} de las autofunciones del
  problema de Sturm--Liouville. Acá es donde entran las series de
  Fourier de la Capítulo~\ref{sec-fourier}.
\end{enumerate}

\begin{tcolorbox}[enhanced jigsaw, arc=.35mm, bottomrule=.15mm, breakable, colback=white, colframe=quarto-callout-tip-color-frame, left=2mm, leftrule=.75mm, opacityback=0, rightrule=.15mm, toprule=.15mm]
\begin{minipage}[t]{5.5mm}
\textcolor{quarto-callout-tip-color}{\faLightbulb}
\end{minipage}%
\begin{minipage}[t]{\textwidth - 5.5mm}

\vspace{-3mm}\textbf{Dónde está el trabajo}\vspace{3mm}

El paso 3 es el que decide todo, y es el que conecta este capítulo con
los primeros. Un problema de Sturm--Liouville garantiza que hay una
sucesión de autovalores, que las autofunciones son ortogonales y que son
\textbf{completas}: cualquier dato inicial razonable se puede escribir
como suma de ellas. Sin completitud, el paso 6 sería imposible y el
método no serviría para nada.

Por eso la primera parte del libro dedicó dos capítulos a
Sturm--Liouville y a las series de Fourier antes de tocar una sola EDP.

\end{minipage}%
\end{tcolorbox}

\subsection{La cuerda con extremos fijos}\label{sec-cuerda-separacion}

Hagamos el caso completo, que es el que hay que saber hacer con los ojos
cerrados.

\begin{equation}\protect\phantomsection\label{eq-cuerda-separacion}{
u_{tt} = a^2u_{xx}, \qquad
\begin{cases}
u(0,t) = 0, & u(x,0) = \varphi(x),\\
u(l,t) = 0, & u_t(x,0) = \psi(x).
\end{cases}
}\end{equation}

\textbf{Paso 1.} Proponemos \(u(x,t) = X(x)\,T(t)\). Derivando y
reemplazando, \[
X\,T'' = a^2\,T\,X'' \qquad\Longrightarrow\qquad \frac{X''}{X} = \frac{T''}{a^2T} .
\] El lado izquierdo depende solo de \(x\) y el derecho solo de \(t\).
Para que sean iguales \textbf{para todo \(x\) y todo \(t\)}, los dos
tienen que ser la misma constante, que llamamos \(-k^2\) por razones que
se van a ver enseguida: \[
\frac{X''}{X} = \frac{T''}{a^2T} = -k^2 .
\] La constante \(-k^2\) se llama \textbf{constante de separación}, y la
EDP se partió en dos ecuaciones \textbf{ordinarias}:
\begin{equation}\protect\phantomsection\label{eq-separadas}{
X'' + k^2X = 0, \qquad T'' + a^2k^2T = 0 .
}\end{equation}

\textbf{Paso 2.} De \(u(0,t)=0\) para todo \(t\) sale \(X(0)\,T(t)=0\),
y como \(T\) no es idénticamente nula, \(X(0)=0\). Igual con el otro
extremo: \(X(l)=0\).

\textbf{Paso 3.} Ahí está el problema de Sturm--Liouville:
\(X''+k^2X=0\) con \(X(0)=X(l)=0\). La solución general es
\(X = C\cos kx + D\operatorname{sen} kx\); la condición en \(x=0\) da
\(C=0\), y la condición en \(x=l\) pide \[
D\operatorname{sen}(kl) = 0 .
\] Con \(D=0\) tendríamos la solución nula, que no sirve. Así que
\(\operatorname{sen}(kl)=0\), y eso solo pasa para valores discretos: \[
k_n = \frac{n\pi}{l}, \qquad n = 1,2,3,\dots
\] \textbf{Acá aparece la cuantización}, y aparece sola: pedirle a una
forma que se anule en los dos extremos deja solo un conjunto numerable
de posibilidades. Es exactamente el mismo mecanismo por el que los
niveles de energía de un pozo cuántico son discretos.

(Y ahora se ve por qué elegimos \(-k^2\): con constante de separación
positiva las soluciones serían exponenciales reales, que no se pueden
anular en dos puntos sin ser idénticamente nulas.)

\textbf{Paso 4.} Para el tiempo, \(T'' + a^2k_n^2T = 0\) da \[
T_n = A_n\cos\!\left(\frac{n\pi a}{l}\,t\right) + B_n\operatorname{sen}\!\left(\frac{n\pi a}{l}\,t\right).
\]

\textbf{Paso 5.} Cada producto
\begin{equation}\protect\phantomsection\label{eq-modo-n}{
u_n(x,t) = \left[A_n\cos\frac{n\pi a\,t}{l} + B_n\operatorname{sen}\frac{n\pi a\,t}{l}\right]
\operatorname{sen}\frac{n\pi x}{l}
}\end{equation} es solución, y por linealidad también lo es su suma:
\begin{equation}\protect\phantomsection\label{eq-serie-cuerda}{
u(x,t) = \sum_{n=1}^{\infty}
\left[A_n\cos\frac{n\pi a\,t}{l} + B_n\operatorname{sen}\frac{n\pi a\,t}{l}\right]
\operatorname{sen}\frac{n\pi x}{l} .
}\end{equation}

\textbf{Paso 6.} Faltan los coeficientes, y para eso están las
condiciones iniciales. Evaluando \ref{eq-serie-cuerda} en \(t=0\), \[
u(x,0) = \sum_n A_n \operatorname{sen}\frac{n\pi x}{l} = \varphi(x),
\] que es la serie de Fourier en senos de \(\varphi\). Derivando
\ref{eq-serie-cuerda} respecto de \(t\) y evaluando en \(t=0\), \[
u_t(x,0) = \sum_n \frac{n\pi a}{l}\,B_n \operatorname{sen}\frac{n\pi x}{l} = \psi(x) .
\] Usando la ortogonalidad de los senos,
\begin{equation}\protect\phantomsection\label{eq-coeficientes-cuerda}{
A_n = \frac{2}{l}\int_0^l \varphi(z)\,\operatorname{sen}\frac{n\pi z}{l}\,dz ,
\qquad
B_n = \frac{2}{n\pi a}\int_0^l \psi(z)\,\operatorname{sen}\frac{n\pi z}{l}\,dz .
}\end{equation}

Y el problema está resuelto.

Al separar la cuerda llegamos a \(\;X''/X = T''/(a^2T)\). ¿Por qué esas
dos expresiones tienen que ser una \textbf{constante}?

\begin{itemize}
\tightlist
\item
  {Porque una depende solo de \(x\), la otra solo de \(t\), y son
  iguales para todo par \((x,t)\)}
\item
  {Porque la ecuación es lineal}
\item
  {Porque las condiciones de contorno son homogéneas}
\item
  {Porque elegimos \(-k^2\) por conveniencia}
\end{itemize}

Es el argumento clave del método y conviene tenerlo claro. Fijemos \(t\)
y movamos \(x\): el lado derecho no cambia, porque no depende de \(x\).
Pero es igual al izquierdo, así que \textbf{el izquierdo tampoco puede
cambiar}. Y al revés con la otra variable. Una función de \(x\) que es
igual a una función de \(t\) para todos los valores de las dos no tiene
más remedio que ser constante.

Las otras tres opciones son afirmaciones verdaderas que hacen falta en
\textbf{otros} pasos: la linealidad se usa en el paso 5 para sumar, y
las condiciones homogéneas en el paso 3 para que haya un problema de
Sturm--Liouville. Pero para este paso solo hace falta que las variables
estén separadas.

¿Qué pasaría si la constante de separación fuera \textbf{positiva}, o
sea \(X''/X = +k^2\) con \(k\) real y no nulo?

\begin{itemize}
\tightlist
\item
  {\(X\) sería combinación de exponenciales reales, y no puede anularse
  en los dos extremos salvo siendo idénticamente nula}
\item
  {\(X\) sería oscilante y habría infinitas soluciones}
\item
  {La ecuación se volvería parabólica}
\item
  {\(T\) crecería exponencialmente, lo cual es físicamente imposible}
\end{itemize}

Con \(X'' = k^2X\) la solución general es
\(X = C\cosh kx + D\operatorname{senh} kx\). La condición \(X(0)=0\) da
\(C=0\); la condición \(X(l)=0\) pide \(D\operatorname{senh}(kl)=0\), y
como \(\operatorname{senh}(kl)\neq 0\) para \(k\neq 0\), queda \(D=0\).
\textbf{Solo la solución nula.} Con \(k=0\) pasa lo mismo: \(X = Cx+D\)
tampoco puede anularse en dos puntos sin ser idénticamente nula.

O sea que el signo de la constante de separación \textbf{no es una
elección nuestra}: lo imponen las condiciones de contorno. Por eso la
escribimos \(-k^2\) desde el principio, sabiendo lo que iba a pasar.

La cuarta opción es la trampa: es cierto que \(T\) crecería
exponencialmente, y también es cierto que eso sería raro para una cuerda
con extremos fijos. Pero el argumento correcto es el otro, y es
puramente matemático: la parte espacial se muere antes de que uno llegue
a mirar el tiempo.

\subsection{Qué es cada cosa}\label{sec-interpretacion}

Cada término de \ref{eq-serie-cuerda} se puede reescribir juntando el
seno y el coseno en uno solo:

\[
u_n(x,t) = N_n\,\operatorname{sen}\!\left(\frac{n\pi a}{l}t + \delta_n\right)
\operatorname{sen}\frac{n\pi x}{l},
\]

que es una \textbf{onda estacionaria}: la forma espacial
\(\operatorname{sen}(n\pi x/l)\) no se mueve, solo late con el tiempo. Y
de ahí se lee todo.

\begin{itemize}
\tightlist
\item
  Las \textbf{frecuencias propias} son \[\omega_n = \frac{n\pi a}{l},\]
  múltiplos enteros de la fundamental \(\omega_1\). Esa serie armónica
  es la que hace que una cuerda tenga una nota definida (comparar con el
  Ejercicio~\ref{exr-helmholtz-modos}, donde la membrana no la tiene).
\item
  Los \textbf{nodos} del armónico \(n\) están donde
  \(\operatorname{sen}(n\pi x/l)=0\), o sea en
  \[x = 0,\ \frac{l}{n},\ \frac{2l}{n},\ \dots,\ l ,\] y son puntos que
  \textbf{no se mueven nunca}. Son los que se ven en el video.
\item
  Los \textbf{vientres} o antinodos, donde la amplitud es máxima, están
  en \[x = \frac{l}{2n},\ \frac{3l}{2n},\ \dots,\ \frac{(2n-1)l}{2n} .\]
\end{itemize}

\begin{tcolorbox}[enhanced jigsaw, arc=.35mm, bottomrule=.15mm, breakable, colback=white, colframe=quarto-callout-important-color-frame, left=2mm, leftrule=.75mm, opacityback=0, rightrule=.15mm, toprule=.15mm]
\begin{minipage}[t]{5.5mm}
\textcolor{quarto-callout-important-color}{\faExclamation}
\end{minipage}%
\begin{minipage}[t]{\textwidth - 5.5mm}

\vspace{-3mm}\textbf{El truco del guitarrista}\vspace{3mm}

Apoyar el dedo suavemente en la mitad de la cuerda y puntearla produce
un sonido una octava más agudo y con un timbre distinto: es el
\textbf{armónico natural}. La explicación está en la lista de nodos.

El dedo impone un nodo en \(x=l/2\). Los modos que tienen un nodo ahí
---los pares, \(n=2,4,6,\dots\)--- sobreviven; los que tendrían un
vientre ahí ---los impares--- se apagan. Lo que queda empieza en
\(\omega_2 = 2\omega_1\): la octava. Apoyando el dedo en \(l/3\) se
filtra distinto y sale la quinta, y así.

El guitarrista no está cambiando la cuerda: está \textbf{eligiendo qué
términos de \ref{eq-serie-cuerda} sobreviven}.

\end{minipage}%
\end{tcolorbox}

\emph{(En la versión web hay acá una animación de los cuatro primeros
modos vibrando, con sus nodos marcados.)}

\emph{(En la versión web hay acá un gráfico interactivo que arma una
cuerda punteada sumando modos, con sliders para la cantidad de modos, el
tiempo y el punto de punteo.)}

¿Cuántos nodos \textbf{interiores} ---sin contar los dos extremos---
tiene el cuarto modo de una cuerda con extremos fijos? Escribí el
número.

El modo \(n\) es \(\operatorname{sen}(n\pi x/l)\), que se anula en
\(x = 0, l/n, 2l/n,\dots,l\): son \(n+1\) nodos en total, de los cuales
dos son los extremos. Así que hay \(n-1\) interiores, y para \(n=4\) son
\textbf{tres}.

Se ven en la animación de los modos: en el cuarto panel hay tres puntos
negros entre medio que no se mueven nunca, y son los que un guitarrista
puede tocar suavemente para filtrar armónicos.

\begin{exercise}[]\protect\hypertarget{exr-cuerda-punteada}{}\label{exr-cuerda-punteada}

Una cuerda de longitud \(1\) se levanta una altura \(h\) en el punto
\(x=p\), quedando con forma de triángulo, y se la suelta desde el
reposo.

\begin{enumerate}
\def\labelenumi{\alph{enumi}.}
\tightlist
\item
  Calcular los coeficientes \(A_n\) de \ref{eq-coeficientes-cuerda}. (El
  resultado es
  \(A_n = \dfrac{2h\operatorname{sen}(n\pi p)}{n^2\pi^2\,p\,(1-p)}\).)
\item
  ¿Qué pasa si se puntea justo en el medio, \(p=1/2\)? ¿Qué armónicos
  desaparecen y por qué?
\item
  Comparar el punteo cerca del puente (\(p\) chico) con el punteo en el
  medio. ¿Cuál suena más brillante, y qué quiere decir eso en términos
  de la serie?
\end{enumerate}

\end{exercise}

\begin{tcolorbox}[enhanced jigsaw, arc=.35mm, bottomrule=.15mm, bottomtitle=1mm, breakable, colback=white, colbacktitle=quarto-callout-tip-color!10!white, colframe=quarto-callout-tip-color-frame, coltitle=black, left=2mm, leftrule=.75mm, opacityback=0, opacitybacktitle=0.6, rightrule=.15mm, title=\textcolor{quarto-callout-tip-color}{\faLightbulb}\hspace{0.5em}{Solución}, titlerule=0mm, toprule=.15mm, toptitle=1mm]

\textbf{a.} El dato es \(\varphi(x) = hx/p\) para \(x\le p\) y
\(h(1-x)/(1-p)\) para \(x\ge p\), y \(\psi \equiv 0\), así que todos los
\(B_n\) son cero y la cuerda solo tiene cosenos en el tiempo. Partiendo
la integral de \ref{eq-coeficientes-cuerda} en las dos ramas e
integrando por partes dos veces se llega al resultado del enunciado. Los
dos hechos que conviene retener:

\begin{itemize}
\tightlist
\item
  \(A_n\) decae como \(1/n^2\). Ese exponente es la firma de una función
  \textbf{continua con una esquina}: si el dato fuera suave decaería más
  rápido, y si tuviera un salto decaería como \(1/n\).
\item
  \(A_n\) lleva el factor \(\operatorname{sen}(n\pi p)\), que es el modo
  \(n\) evaluado en el punto de punteo.
\end{itemize}

\textbf{b.} Con \(p=1/2\), \(\operatorname{sen}(n\pi/2)\) vale \(0\)
para todo \(n\) \textbf{par}. Desaparecen todos los armónicos pares.

Y la razón es geométrica, no algebraica: los modos pares tienen un
\textbf{nodo} justo en \(x=1/2\). Levantar la cuerda desde un nodo no
puede excitar un modo que ahí no se mueve. En general,
\(\operatorname{sen}(n\pi p) = 0\) significa exactamente que \(p\) es un
nodo del modo \(n\): \textbf{puntear en un nodo apaga ese armónico}. Es
el mismo mecanismo que el armónico natural del guitarrista, al revés.

\textbf{c.} Cerca del puente. Con \(p\) chico,
\(\operatorname{sen}(n\pi p) \approx n\pi p\) para todos los \(n\) hasta
\(n \sim 1/p\), y entonces \[A_n \approx \frac{2h}{n\pi(1-p)} ,\] que
decae como \(1/n\) y no como \(1/n^2\): \textbf{hay mucha más energía en
los armónicos altos}. Eso es lo que el oído llama un sonido brillante o
metálico. Punteando en el medio, en cambio, se apagan todos los pares y
los impares decaen rápido: sonido redondo, casi puro.

Es exactamente lo que hace un guitarrista al mover la mano entre el
puente y el diapasón, y es la razón de que el mismo instrumento y la
misma nota puedan sonar de maneras tan distintas. La altura la fija
\(\omega_1\); el timbre lo fija la lista de los \(A_n\).

\end{tcolorbox}

\begin{exercise}[]\protect\hypertarget{exr-calor-barra}{}\label{exr-calor-barra}

Aplicar el esquema de seis pasos a la ecuación del calor:
\[u_t = k\,u_{xx}, \qquad u(0,t)=u(L,t)=0, \qquad u(x,0)=\varphi(x).\]

\begin{enumerate}
\def\labelenumi{\alph{enumi}.}
\tightlist
\item
  Hallar la solución en serie.
\item
  ¿Cuánto tarda en apagarse el modo \(n\) comparado con el fundamental?
  ¿Qué tiene que ver eso con lo que dijimos en la
  sección~\ref{sec-tres-mundos} sobre el suavizado instantáneo?
\item
  ¿Qué pasa si se intenta usar la misma fórmula para \(t<0\)?
\end{enumerate}

\end{exercise}

\begin{tcolorbox}[enhanced jigsaw, arc=.35mm, bottomrule=.15mm, bottomtitle=1mm, breakable, colback=white, colbacktitle=quarto-callout-tip-color!10!white, colframe=quarto-callout-tip-color-frame, coltitle=black, left=2mm, leftrule=.75mm, opacityback=0, opacitybacktitle=0.6, rightrule=.15mm, title=\textcolor{quarto-callout-tip-color}{\faLightbulb}\hspace{0.5em}{Solución}, titlerule=0mm, toprule=.15mm, toptitle=1mm]

\textbf{a.} Con \(u = X(x)T(t)\) queda \(T'/(kT) = X''/X = -\lambda^2\),
es decir \[X'' + \lambda^2X = 0, \qquad T' + k\lambda^2T = 0 .\] La
parte espacial es \textbf{el mismo problema de Sturm--Liouville de la
cuerda} ---misma ecuación, mismas condiciones---, así que
\(\lambda_n = n\pi/L\) y \(X_n = \operatorname{sen}(n\pi x/L)\). Lo que
cambia es el tiempo: ahora la ecuación es de primer orden y da
\(T_n = e^{-k\lambda_n^2 t}\), una \textbf{exponencial decreciente} en
vez de una oscilación. Entonces \[
u(x,t) = \sum_{n=1}^{\infty} A_n\,e^{-k n^2\pi^2 t/L^2}\,\operatorname{sen}\frac{n\pi x}{L},
\qquad
A_n = \frac{2}{L}\int_0^L \varphi(z)\operatorname{sen}\frac{n\pi z}{L}\,dz .
\] Los coeficientes son \textbf{los mismos} que los de la cuerda: la
diferencia entre calor y onda está toda en el factor temporal.

\textbf{b.} El modo \(n\) decae con tiempo característico
\[\tau_n = \frac{L^2}{k\,n^2\pi^2} = \frac{\tau_1}{n^2} :\] \textbf{el
modo \(n\) se apaga \(n^2\) veces más rápido que el fundamental.} El
décimo armónico dura cien veces menos.

Y ahí está el suavizado de la sección~\ref{sec-tres-mundos}, con una
fórmula. Los detalles finos de un dato inicial son justamente sus
armónicos altos, y son los primeros en morir. Apenas \(t>0\), la
exponencial \(e^{-kn^2\pi^2t/L^2}\) aplasta la cola de la serie tan
fuerte que la suma es infinitamente derivable, por más picuda que fuera
\(\varphi\). Después de un rato solo queda el fundamental, y la barra
adopta la forma \(\operatorname{sen}(\pi x/L)\) sin importar cómo
empezó.

\textbf{c.} Con \(t<0\) el factor se convierte en
\(e^{+k n^2\pi^2|t|/L^2}\), que \textbf{crece} con \(n\) y lo hace como
una exponencial de \(n^2\). La serie deja de converger salvo para datos
excepcionalmente suaves, y un error minúsculo en un armónico alto se
amplifica sin control.

Es el ejemplo de Hadamard de la sección~\ref{sec-hadamard} otra vez,
ahora en su versión más transparente: la ecuación del calor hacia atrás
está mal planteada, y la serie de Fourier lo muestra término a término.

\end{tcolorbox}

\subsection{Una que no suena a nada: la viga}\label{sec-viga}

El método no está atado a las ecuaciones de segundo orden. Veámoslo con
una viga elástica, que es de \textbf{cuarto} orden y da un resultado
sorprendente.

\begin{example}[Vibraciones de una viga
empotrada]\protect\hypertarget{exm-viga}{}\label{exm-viga}

Una viga fina y uniforme de longitud \(L\), empotrada en \(x=0\) y libre
en \(x=L\), tiene una deflexión \(w(x,t)\) que satisface

\begin{equation}\protect\phantomsection\label{eq-viga}{
\rho\,\frac{\partial^2 w}{\partial t^2} + \gamma\,\frac{\partial w}{\partial t}
+ EI\,\frac{\partial^4 w}{\partial x^4} = f(x,t) ,
}\end{equation}

donde \(\rho\) es la masa por unidad de longitud, \(\gamma\) el
coeficiente de fricción viscosa ---cuando la viga vibra en el aire, por
ejemplo---, \(EI\) el módulo de flexión y \(f\) la fuerza externa.

\textbf{Paso 1.} Sin fuerzas externas, proponemos \(w = T(t)\,\phi(x)\)
y separamos. Llamando \(\omega^4\) a la constante de separación,
\begin{equation}\protect\phantomsection\label{eq-viga-separada}{
\frac{d^4\phi}{dx^4} - \omega^4\phi = 0, \qquad
\rho\,\frac{d^2T}{dt^2} + \gamma\,\frac{dT}{dt} + EI\,\omega^4\,T = 0 .
}\end{equation}

\textbf{Paso 2.} Las condiciones de una viga empotrada en un extremo y
libre en el otro son cuatro: \[
w(0,t) = \frac{\partial w}{\partial x}(0,t) = 0, \qquad
\frac{\partial^2 w}{\partial x^2}(L,t) = \frac{\partial^3 w}{\partial x^3}(L,t) = 0 .
\] Las dos primeras dicen que en el empotramiento la viga no se mueve ni
se inclina; las dos últimas, que en el extremo libre no hay momento
flector ni fuerza de corte.

\textbf{Paso 3.} Aunque no estudiamos Sturm--Liouville para orden
cuatro, no es difícil ver que \(\mathcal{L} = d^4/dx^4\) con estas
condiciones es \textbf{autoadjunto}, y eso alcanza para garantizar que
las autofunciones formen un conjunto ortogonal completo de \(L^2[0,L]\).
La solución general de \(\phi'''' = \omega^4\phi\) es \[
\phi(x) = c_1\cosh\omega x + c_2\operatorname{senh}\omega x + c_3\cos\omega x + c_4\operatorname{sen}\omega x ,
\] y aplicando las cuatro condiciones se llega a las autofunciones
\begin{equation}\protect\phantomsection\label{eq-modos-viga}{
\phi_k(x) = \frac{\cosh\omega_kx - \cos\omega_kx}{\cosh\omega_kL + \cos\omega_kL}
- \frac{\operatorname{senh}\omega_kx - \operatorname{sen}\omega_kx}{\operatorname{senh}\omega_kL + \operatorname{sen}\omega_kL} ,
}\end{equation} con los autovalores dados por las raíces de
\begin{equation}\protect\phantomsection\label{eq-autovalores-viga}{
\cos\omega_kL = -\frac{1}{\cosh\omega_kL} .
}\end{equation}

\textbf{Paso 4.} Las condiciones iniciales y la ecuación para \(T(t)\)
cierran el problema. Como esa ecuación es ordinaria, lineal y con
coeficientes constantes, se la puede atacar con la transformada de
Laplace de la sección~\ref{sec-transformadas}.

\end{example}

\begin{tcolorbox}[enhanced jigsaw, arc=.35mm, bottomrule=.15mm, breakable, colback=white, colframe=quarto-callout-important-color-frame, left=2mm, leftrule=.75mm, opacityback=0, rightrule=.15mm, toprule=.15mm]
\begin{minipage}[t]{5.5mm}
\textcolor{quarto-callout-important-color}{\faExclamation}
\end{minipage}%
\begin{minipage}[t]{\textwidth - 5.5mm}

\vspace{-3mm}\textbf{Por qué un xilofón hay que afinarlo a mano}\vspace{3mm}

La \ref{eq-autovalores-viga} no se resuelve con lápiz y papel, pero sus
raíces son fáciles de obtener numéricamente. Las primeras son \[
\omega_1L \approx 1{,}875, \qquad \omega_2L \approx 4{,}694, \qquad \omega_3L \approx 7{,}855 .
\] Y como la frecuencia de vibración va con \(\omega^2\) ---de
\ref{eq-viga-separada} con \(\gamma=0\),
\(\Omega = \omega^2\sqrt{EI/\rho}\)---, las frecuencias están en la
proporción \[
1 \;:\; 6{,}27 \;:\; 17{,}55 \;:\; \dots
\]

\textbf{Nada que ver con \(1:2:3\).} Una viga no tiene serie armónica, y
por eso una barra metálica golpeada suena a «clonk» metálico y no a una
nota: sus sobretonos son inarmónicos y el oído no les encuentra un
fundamental común.

Los fabricantes de xilofones y marimbas resuelven esto \textbf{cavando
la barra por abajo}, con un arco que le saca material en el centro. Eso
cambia \(EI(x)\) a lo largo de la barra y corre los sobretonos hasta
ponerlos en \(1:4:10\), que sí es una relación que el oído acepta como
una nota con timbre. Es un problema de EDP resuelto con una gubia, y se
resolvió empíricamente siglos antes de que existiera esta teoría.

\end{minipage}%
\end{tcolorbox}

\begin{exercise}[]\protect\hypertarget{exr-viga-numeros}{}\label{exr-viga-numeros}

Sobre la viga del Ejemplo~\ref{exm-viga}:

\begin{enumerate}
\def\labelenumi{\alph{enumi}.}
\tightlist
\item
  Verificar gráfica o numéricamente que la primera raíz de
  \ref{eq-autovalores-viga} es \(\omega_1L \approx 1{,}875\).
\item
  Calcular la relación entre las tres primeras frecuencias sabiendo que
  \(\Omega \propto \omega^2\), y comparar con la serie armónica
  \(1:2:3\) de la cuerda.
\item
  Una regla de plástico sujeta al borde de una mesa y hecha vibrar: ¿por
  qué suena distinto según cuánto sobresalga, y cómo depende la
  frecuencia de la longitud libre?
\end{enumerate}

\end{exercise}

\begin{tcolorbox}[enhanced jigsaw, arc=.35mm, bottomrule=.15mm, bottomtitle=1mm, breakable, colback=white, colbacktitle=quarto-callout-tip-color!10!white, colframe=quarto-callout-tip-color-frame, coltitle=black, left=2mm, leftrule=.75mm, opacityback=0, opacitybacktitle=0.6, rightrule=.15mm, title=\textcolor{quarto-callout-tip-color}{\faLightbulb}\hspace{0.5em}{Solución}, titlerule=0mm, toprule=.15mm, toptitle=1mm]

\textbf{a.} Hay que buscar los cruces de \(\cos z\) con \(-1/\cosh z\).
Como \(\cosh z\) crece rápido, el lado derecho se pega a cero enseguida,
así que las raíces se acercan a los ceros de \(\cos z\), o sea a
\(z \approx (2m-1)\pi/2\). Para la primera, \(\pi/2 \approx 1{,}571\) es
una estimación grosera; refinando ---con bisección, o probando--- se
llega a \(z_1 = 1{,}8751\). Las siguientes, \(4{,}6941\) y \(7{,}8548\),
ya están muy cerca de \(3\pi/2 = 4{,}712\) y \(5\pi/2 = 7{,}854\).

\textbf{b.} Como \(\Omega_n \propto (\omega_nL)^2\), \[
\frac{\Omega_2}{\Omega_1} = \left(\frac{4{,}6941}{1{,}8751}\right)^2 = 6{,}27, \qquad
\frac{\Omega_3}{\Omega_1} = \left(\frac{7{,}8548}{1{,}8751}\right)^2 = 17{,}55 .
\] La serie es \(1 : 6{,}27 : 17{,}55\), muy lejos de \(1:2:3\). Los
sobretonos de una viga son \textbf{inarmónicos}, y por eso una barra
golpeada no tiene una altura definida como la tiene una cuerda.

\textbf{c.} La frecuencia sale de
\(\Omega = (\omega L)^2\sqrt{EI/\rho}\,/L^2\), y como \(\omega_1L\) es
un número fijo, \[\Omega_1 \propto \frac{1}{L^2} .\] Duplicar la parte
que sobresale baja la frecuencia \textbf{cuatro veces}, no dos. Es lo
que uno comprueba corriendo la regla sobre el borde de la mesa: el tono
cae muchísimo más rápido de lo que uno esperaría, y esa dependencia
cuadrática es la firma de una ecuación de cuarto orden. Una cuerda, que
es de segundo orden, tiene \(\Omega \propto 1/L\).

\end{tcolorbox}

\subsection{Cuándo se puede separar}\label{sec-cuando-separar}

Las condiciones son bastante exigentes, y conviene chequearlas antes de
empezar:

\begin{longtable}[]{@{}
  >{\raggedright\arraybackslash}p{(\linewidth - 2\tabcolsep) * \real{0.5800}}
  >{\raggedright\arraybackslash}p{(\linewidth - 2\tabcolsep) * \real{0.4200}}@{}}
\caption{Los cuatro requisitos.}\label{tbl-separacion}\tabularnewline
\toprule\noalign{}
\begin{minipage}[b]{\linewidth}\raggedright
Hace falta\ldots{}
\end{minipage} & \begin{minipage}[b]{\linewidth}\raggedright
Si no\ldots{}
\end{minipage} \\
\midrule\noalign{}
\endfirsthead
\toprule\noalign{}
\begin{minipage}[b]{\linewidth}\raggedright
Hace falta\ldots{}
\end{minipage} & \begin{minipage}[b]{\linewidth}\raggedright
Si no\ldots{}
\end{minipage} \\
\midrule\noalign{}
\endhead
\bottomrule\noalign{}
\endlastfoot
que la ecuación sea \textbf{lineal} & no vale la superposición del paso
5 \\
que el \textbf{dominio} sea una «caja» en algún sistema de coordenadas
separable & los bordes mezclarían las variables \\
que las condiciones de contorno en las direcciones separadas sean
\textbf{homogéneas} & el paso 3 no da un problema de Sturm--Liouville \\
que los \textbf{coeficientes} no mezclen variables & no se puede pasar
todo de un lado \\
\end{longtable}

Los sistemas de coordenadas en los que la ecuación de Helmholtz separa
son once, y son todos los que uno ya conoce: cartesianas, cilíndricas,
esféricas, parabólicas, elípticas y unas pocas más. No es casualidad que
los problemas resueltos de los libros sean siempre en cajas, cilindros y
esferas: \textbf{son los dominios donde el método funciona}. Un dominio
con forma de riñón no separa en nada, y ahí hay que ir a métodos
numéricos.

Y las condiciones de contorno inhomogéneas no son un impedimento
definitivo: se las traslada al paso 6 o se parte el problema en dos por
linealidad, tal como hicimos en la sección~\ref{sec-green-varias}. Lo
que no se puede es tenerlas en las direcciones que se separan.

\begin{exercise}[]\protect\hypertarget{exr-separacion-requisitos}{}\label{exr-separacion-requisitos}

Decidir cuáles de estos problemas se pueden atacar por separación de
variables, y por qué.

\begin{enumerate}
\def\labelenumi{\alph{enumi}.}
\tightlist
\item
  La ecuación de onda en un rectángulo, con todo el borde fijo.
\item
  La ecuación de Laplace en un disco, con el potencial dado sobre la
  circunferencia.
\item
  La ecuación del calor en un dominio con forma de riñón.
\item
  La ecuación de onda en \([0,L]\) con el extremo \(x=0\) fijo y el
  extremo \(x=L\) oscilando, \(u(L,t) = A\operatorname{sen}\Omega t\).
\item
  La ecuación de Burgers \(u_t + u\,u_x = 0\).
\end{enumerate}

\end{exercise}

\begin{tcolorbox}[enhanced jigsaw, arc=.35mm, bottomrule=.15mm, bottomtitle=1mm, breakable, colback=white, colbacktitle=quarto-callout-tip-color!10!white, colframe=quarto-callout-tip-color-frame, coltitle=black, left=2mm, leftrule=.75mm, opacityback=0, opacitybacktitle=0.6, rightrule=.15mm, title=\textcolor{quarto-callout-tip-color}{\faLightbulb}\hspace{0.5em}{Solución}, titlerule=0mm, toprule=.15mm, toptitle=1mm]

\textbf{a. Sí.} Lineal, dominio que es una caja en coordenadas
cartesianas, y condiciones homogéneas en las dos direcciones. Salen los
modos \(\operatorname{sen}(m\pi x/a)\operatorname{sen}(n\pi y/b)\) del
Ejercicio~\ref{exr-helmholtz-modos}.

\textbf{b. Sí}, en coordenadas polares. El disco es una «caja» en
\((r,\theta)\) ---es \(0\le r\le R\), \(0\le\theta<2\pi\)--- y ahí la
ecuación separa. Las condiciones que faltan aparecen en el paso 2 y no
estaban escritas: que \(u\) sea \textbf{periódica} en \(\theta\) y
\textbf{acotada} en \(r=0\). El dato sobre la circunferencia es
inhomogéneo, pero está en la dirección \(r\), que es la que no se separa
como problema de autovalores, así que no molesta: se usa en el paso 6.

\textbf{c.~No.} No hay ningún sistema de coordenadas en el que un riñón
sea una caja. Este problema va a métodos numéricos, y ahí es donde
entran los elementos finitos.

\textbf{d.~No directamente}, porque la condición en \(x=L\) es
inhomogénea y está justo en la dirección que hay que separar: el paso 3
no daría un problema de Sturm--Liouville. La salida estándar es partir
\(u = v + w\), con \(w\) una función elegida a mano que absorba el dato
del borde ---por ejemplo \(w = (x/L)A\operatorname{sen}\Omega t\)--- de
modo que \(v\) quede con condiciones homogéneas. La ecuación para \(v\)
tiene ahora una fuente, pero eso sí se puede manejar.

\textbf{e. No.} No es lineal, así que se cae el paso 5: la suma de dos
soluciones no es solución y no tiene sentido armar una serie. Para
Burgers hay que volver a la sección~\ref{sec-casi-lineales}.

\end{tcolorbox}

\section{Conclusiones}\label{sec-conclusiones}

Este capítulo abrió la segunda mitad del libro, y lo hizo cambiando de
objeto: de curvas a superficies, de constantes arbitrarias a funciones
arbitrarias, de condiciones puntuales a condiciones sobre curvas
enteras. Vale la pena recoger lo que quedó.

\textbf{La clasificación no es una taxonomía.} Aprendimos a calcular dos
cosas de una ecuación antes de intentar resolverla: cuán lineal es y de
qué tipo es. La primera decide \textbf{qué métodos están permitidos}
---la superposición, que es toda la primera parte del libro, vive solo
del lado lineal---. La segunda decide \textbf{cómo se comportan las
soluciones}: si la información viaja a velocidad finita conservando los
quiebres, si se suaviza al instante, o si no viaja en absoluto. Y las
dos se calculan mirando únicamente la parte principal.

\textbf{Los problemas se pueden plantear mal, y eso importa.} Existencia
y unicidad no alcanzan: si la solución no depende continuamente de los
datos, la que uno calcula no es la que ocurre. Cada tipo pide sus datos
y rechaza los otros, y el ejemplo de Hadamard mostró que darle datos de
Cauchy a una elíptica no es un error de cálculo sino una imposibilidad
de principio.

\textbf{Las soluciones clásicas no alcanzan.} La ecuación de Burgers
desarrolla discontinuidades sola, partiendo de un dato perfectamente
suave, y esas discontinuidades son ondas de choque reales. La salida fue
reconocer que la ecuación diferencial era un producto derivado de una
ley de conservación integral, y volver a ella.

\textbf{Y las características son el hilo conductor.} Aparecieron como
las curvas donde una ecuación de primer orden se vuelve ordinaria;
reaparecieron como las direcciones nulas de la forma cuadrática, que es
lo que define el tipo; volvieron como las coordenadas en las que la
ecuación adopta su forma canónica; y terminaron siendo el cono de luz de
la relatividad. Es el mismo objeto todas las veces.

\subsection{Cómo elegir el método}\label{sec-elegir-metodo}

Los cuatro métodos del capítulo no compiten: cada uno tiene su
territorio.

\begin{longtable}[]{@{}
  >{\raggedright\arraybackslash}p{(\linewidth - 4\tabcolsep) * \real{0.3400}}
  >{\raggedright\arraybackslash}p{(\linewidth - 4\tabcolsep) * \real{0.2600}}
  >{\raggedright\arraybackslash}p{(\linewidth - 4\tabcolsep) * \real{0.4000}}@{}}
\caption{El árbol de decisión del
capítulo.}\label{tbl-elegir}\tabularnewline
\toprule\noalign{}
\begin{minipage}[b]{\linewidth}\raggedright
Si el problema\ldots{}
\end{minipage} & \begin{minipage}[b]{\linewidth}\raggedright
conviene\ldots{}
\end{minipage} & \begin{minipage}[b]{\linewidth}\raggedright
porque
\end{minipage} \\
\midrule\noalign{}
\endfirsthead
\toprule\noalign{}
\begin{minipage}[b]{\linewidth}\raggedright
Si el problema\ldots{}
\end{minipage} & \begin{minipage}[b]{\linewidth}\raggedright
conviene\ldots{}
\end{minipage} & \begin{minipage}[b]{\linewidth}\raggedright
porque
\end{minipage} \\
\midrule\noalign{}
\endhead
\bottomrule\noalign{}
\endlastfoot
es de primer orden, aunque sea cuasi-lineal & \textbf{características}
(sección~\ref{sec-caracteristicas}) & es el único que no necesita
linealidad \\
es lineal, con una variable no acotada y datos iniciales &
\textbf{transformadas integrales} (sección~\ref{sec-transformadas}) &
esa variable se elimina y el problema baja una dimensión \\
es lineal y hay que resolverlo para muchas fuentes distintas &
\textbf{función de Green} (sección~\ref{sec-green-varias}) & se calcula
una vez y sirve para todas \\
es lineal, en un dominio con la simetría adecuada y con contornos
homogéneos & \textbf{separación de variables}
(sección~\ref{sec-separacion}) & da los modos, que es como el sistema se
describe a sí mismo \\
no es lineal y no es de primer orden & \emph{ninguno de los cuatro} &
hay que ir a métodos numéricos o a técnicas específicas \\
\end{longtable}

En la práctica se combinan. El camino más frecuente es:
\textbf{transformar} en el tiempo para quedarse con un problema de
Helmholtz, y después resolverlo con \textbf{Green} o con
\textbf{separación de variables} según cómo sea el dominio. Los
capítulos que siguen hacen exactamente eso, un tipo por capítulo:
hiperbólicas en el 9, parabólicas en el 10, elípticas en el 11.

\section{Resumen del capítulo}\label{sec-resumen}

Esto es una guía de temas para el final. No reemplaza al capítulo: lo
condensa. Están todas las definiciones y todos los enunciados,
\textbf{copiados palabra por palabra} de donde aparecen, y además las
fórmulas que conviene tener a mano y un procedimiento para cuando te dan
una ecuación y no sabés por dónde empezarla.

Los temas y el orden son los del capítulo 8 del libro de la cátedra. De
lo que este libro agrega por fuera de ese programa quedan solo dos
cosas, porque son demasiado útiles para dejarlas afuera: la lectura del
campo vectorial que guía a las características, y la lista de las
maneras en que el método puede fallar.

\begin{quote}
\textbf{Sobre la notación.} Para las ecuaciones de segundo orden en dos
variables escribimos \(A\,u_{xx} + 2B\,u_{xy} + C\,u_{yy}\) y
\(\Delta = B^2-AC\). González escribe
\(a\,u_{xx} + 2b\,u_{xy} + c\,u_{yy}\) y \(\Delta = b^2-ac\): es
exactamente lo mismo, solo cambian las letras.
\end{quote}

\subsection{Qué es una EDP, y qué es una solución}\label{sec-resumen-1}

\begin{tcolorbox}[enhanced jigsaw, arc=.35mm, bottomrule=.15mm, breakable, colback=white, colframe=quarto-callout-note-color-frame, left=2mm, leftrule=.75mm, opacityback=0, rightrule=.15mm, toprule=.15mm]
\begin{minipage}[t]{5.5mm}
\textcolor{quarto-callout-note-color}{\faInfo}
\end{minipage}%
\begin{minipage}[t]{\textwidth - 5.5mm}

\vspace{-3mm}\textbf{Definición~\ref{def-edp} --- Ecuación en derivadas parciales}\vspace{3mm}

Sea \(u\) una función incógnita de \(n\) variables \(x_1,\dots,x_n\).
Una \textbf{ecuación en derivadas parciales} es una relación de la forma

\[
\Phi\!\left(x_1,\dots,x_n,\;u,\;\frac{\partial u}{\partial x_1},\dots,\frac{\partial u}{\partial x_n},\;
\frac{\partial^2 u}{\partial x_1^2},\frac{\partial^2 u}{\partial x_1\partial x_2},\dots,\;
\frac{\partial^N u}{\partial x_n^N}\right) = 0 .
\]

Su \textbf{orden} es el mayor \(N\) tal que en la relación aparece
efectivamente una derivada de orden \(N\).

\end{minipage}%
\end{tcolorbox}

Notación:
\(u_{x_ix_j} = \dfrac{\partial^2u}{\partial x_i\partial x_j}\), y con
multiíndice \(\alpha=(\alpha_1,\dots,\alpha_n)\) y
\(|\alpha| = \alpha_1+\dots+\alpha_n\),

\[
D^{\alpha}u = \frac{\partial^{|\alpha|} u}{\partial x_1^{\alpha_1}\cdots\partial x_n^{\alpha_n}} .
\]

\begin{tcolorbox}[enhanced jigsaw, arc=.35mm, bottomrule=.15mm, breakable, colback=white, colframe=quarto-callout-note-color-frame, left=2mm, leftrule=.75mm, opacityback=0, rightrule=.15mm, toprule=.15mm]
\begin{minipage}[t]{5.5mm}
\textcolor{quarto-callout-note-color}{\faInfo}
\end{minipage}%
\begin{minipage}[t]{\textwidth - 5.5mm}

\vspace{-3mm}\textbf{Definición~\ref{def-solucion-clasica} --- Solución clásica}\vspace{3mm}

Una \textbf{solución clásica} ---o \textbf{fuerte}--- de una EDP de
orden \(N\) en un dominio \(D\) es una función \(u\) que es \(N\) veces
continuamente derivable en \(D\) y que satisface la ecuación en todos
los puntos de \(D\).

\end{minipage}%
\end{tcolorbox}

\textbf{Soluciones generalizadas o débiles.} No siempre alcanza con las
clásicas. Hay ecuaciones que solo admiten soluciones con
discontinuidades en \(u\) o en sus derivadas si se quiere que el
problema esté bien planteado, y esas soluciones corresponden a fenómenos
físicos reales: el ejemplo típico es la ecuación de conservación

\[
u_t + F(u)_x = 0 ,
\]

cuyas \textbf{ondas de choque} son superficies de discontinuidad. El
motivo de fondo es que las ecuaciones diferenciales suelen ser el
resultado local de un balance global escrito como ecuación integral, y
para una ecuación integral una discontinuidad no es ningún problema.

\subsection{Primera clasificación: el grado de
linealidad}\label{sec-resumen-2}

\begin{tcolorbox}[enhanced jigsaw, arc=.35mm, bottomrule=.15mm, breakable, colback=white, colframe=quarto-callout-note-color-frame, left=2mm, leftrule=.75mm, opacityback=0, rightrule=.15mm, toprule=.15mm]
\begin{minipage}[t]{5.5mm}
\textcolor{quarto-callout-note-color}{\faInfo}
\end{minipage}%
\begin{minipage}[t]{\textwidth - 5.5mm}

\vspace{-3mm}\textbf{Definición~\ref{def-parte-principal} --- Parte principal}\vspace{3mm}

Sea una EDP de orden \(N\) escrita como

\[
\sum_{|\alpha| = N} a_\alpha\, D^\alpha u \;+\; (\text{términos con derivadas de orden} < N) = 0 .
\]

Su \textbf{parte principal} es el operador formado solo por los términos
de orden exactamente \(N\):

\[
L_P[u] = \sum_{|\alpha| = N} a_\alpha\, D^\alpha u .
\]

\end{minipage}%
\end{tcolorbox}

\begin{tcolorbox}[enhanced jigsaw, arc=.35mm, bottomrule=.15mm, breakable, colback=white, colframe=quarto-callout-note-color-frame, left=2mm, leftrule=.75mm, opacityback=0, rightrule=.15mm, toprule=.15mm]
\begin{minipage}[t]{5.5mm}
\textcolor{quarto-callout-note-color}{\faInfo}
\end{minipage}%
\begin{minipage}[t]{\textwidth - 5.5mm}

\vspace{-3mm}\textbf{Definición~\ref{def-lineal} --- Ecuación lineal}\vspace{3mm}

Una EDP de orden \(N\) es \textbf{lineal} si se puede escribir como

\[
\sum_{|\alpha| \le N} a_\alpha(x)\, D^\alpha u = f(x),
\]

es decir, si es lineal en \(u\) y en todas sus derivadas, con
coeficientes que dependen únicamente de las variables independientes.
Escribiendo \(L[u] = \sum_{|\alpha|\le N} a_\alpha(x)D^\alpha u\), la
ecuación es \(L[u] = f\), y \(L\) es un \textbf{operador lineal}:
\(L[\lambda u + \mu v] = \lambda L[u] + \mu L[v]\).

Si \(f \equiv 0\) la ecuación se dice \textbf{homogénea}; si no,
\textbf{inhomogénea} o \textbf{con término independiente}.

\end{minipage}%
\end{tcolorbox}

\begin{tcolorbox}[enhanced jigsaw, arc=.35mm, bottomrule=.15mm, breakable, colback=white, colframe=quarto-callout-note-color-frame, left=2mm, leftrule=.75mm, opacityback=0, rightrule=.15mm, toprule=.15mm]
\begin{minipage}[t]{5.5mm}
\textcolor{quarto-callout-note-color}{\faInfo}
\end{minipage}%
\begin{minipage}[t]{\textwidth - 5.5mm}

\vspace{-3mm}\textbf{Definición~\ref{def-semilineal} --- Ecuación semilineal}\vspace{3mm}

Una EDP de orden \(N\) es \textbf{semilineal} si su parte principal es
lineal con coeficientes que dependen solo de \(x\), y toda la no
linealidad está confinada a derivadas de orden estrictamente menor que
\(N\):

\[
\sum_{|\alpha| = N} a_\alpha(x)\, D^\alpha u \;+\; a_0\big(x,\,u,\,Du,\,\dots,\,D^{N-1}u\big) = 0 .
\]

\end{minipage}%
\end{tcolorbox}

\begin{tcolorbox}[enhanced jigsaw, arc=.35mm, bottomrule=.15mm, breakable, colback=white, colframe=quarto-callout-note-color-frame, left=2mm, leftrule=.75mm, opacityback=0, rightrule=.15mm, toprule=.15mm]
\begin{minipage}[t]{5.5mm}
\textcolor{quarto-callout-note-color}{\faInfo}
\end{minipage}%
\begin{minipage}[t]{\textwidth - 5.5mm}

\vspace{-3mm}\textbf{Definición~\ref{def-cuasilineal} --- Ecuación cuasi-lineal}\vspace{3mm}

Una EDP de orden \(N\) es \textbf{cuasi-lineal} si es lineal en las
derivadas de orden \(N\) ---aparecen a la primera potencia y no
multiplicadas entre sí--- pero los coeficientes de esas derivadas pueden
depender de \(u\) y de sus derivadas de orden menor:

\[
\sum_{|\alpha| = N} a_\alpha\big(x,u,Du,\dots,D^{N-1}u\big)\, D^\alpha u
\;+\; a_0\big(x,u,Du,\dots,D^{N-1}u\big) = 0 .
\]

\end{minipage}%
\end{tcolorbox}

\begin{tcolorbox}[enhanced jigsaw, arc=.35mm, bottomrule=.15mm, breakable, colback=white, colframe=quarto-callout-note-color-frame, left=2mm, leftrule=.75mm, opacityback=0, rightrule=.15mm, toprule=.15mm]
\begin{minipage}[t]{5.5mm}
\textcolor{quarto-callout-note-color}{\faInfo}
\end{minipage}%
\begin{minipage}[t]{\textwidth - 5.5mm}

\vspace{-3mm}\textbf{Definición~\ref{def-no-lineal} --- Ecuación completamente no lineal}\vspace{3mm}

Una EDP es \textbf{completamente no lineal} si depende de las derivadas
de orden \(N\) de manera no lineal: alguna aparece elevada a una
potencia, dentro de una función, o multiplicada por otra derivada de
orden \(N\).

\end{minipage}%
\end{tcolorbox}

\begin{longtable}[]{@{}
  >{\raggedright\arraybackslash}p{(\linewidth - 4\tabcolsep) * \real{0.3333}}
  >{\raggedright\arraybackslash}p{(\linewidth - 4\tabcolsep) * \real{0.3333}}
  >{\raggedright\arraybackslash}p{(\linewidth - 4\tabcolsep) * \real{0.3333}}@{}}
\caption{La escala de linealidad. Las cuatro clases están encajadas:
toda lineal es semilineal, toda semilineal es
cuasi-lineal.}\tabularnewline
\toprule\noalign{}
\begin{minipage}[b]{\linewidth}\raggedright
Clase
\end{minipage} & \begin{minipage}[b]{\linewidth}\raggedright
Los coeficientes de la parte principal\ldots{}
\end{minipage} & \begin{minipage}[b]{\linewidth}\raggedright
La no linealidad está\ldots{}
\end{minipage} \\
\midrule\noalign{}
\endfirsthead
\toprule\noalign{}
\begin{minipage}[b]{\linewidth}\raggedright
Clase
\end{minipage} & \begin{minipage}[b]{\linewidth}\raggedright
Los coeficientes de la parte principal\ldots{}
\end{minipage} & \begin{minipage}[b]{\linewidth}\raggedright
La no linealidad está\ldots{}
\end{minipage} \\
\midrule\noalign{}
\endhead
\bottomrule\noalign{}
\endlastfoot
Lineal & dependen solo de \(x\) & no hay \\
Semilineal & dependen solo de \(x\) & en términos de orden \(< N\) \\
Cuasi-lineal & pueden depender de \(u\) y de \(Du,\dots,D^{N-1}u\) & en
los coeficientes y en el resto \\
Completamente no lineal & --- & en las derivadas de orden \(N\) \\
\end{longtable}

\begin{quote}
\textbf{Cuidado con el nombre.} González llama \textbf{«casi--lineales»}
a las que acá y en casi toda la bibliografía se llaman
\textbf{cuasi-lineales}.
\end{quote}

\subsection{Cuándo un problema está bien planteado}\label{sec-resumen-3}

\begin{tcolorbox}[enhanced jigsaw, arc=.35mm, bottomrule=.15mm, breakable, colback=white, colframe=quarto-callout-note-color-frame, left=2mm, leftrule=.75mm, opacityback=0, rightrule=.15mm, toprule=.15mm]
\begin{minipage}[t]{5.5mm}
\textcolor{quarto-callout-note-color}{\faInfo}
\end{minipage}%
\begin{minipage}[t]{\textwidth - 5.5mm}

\vspace{-3mm}\textbf{Definición~\ref{def-bien-planteado} --- Problema bien planteado}\vspace{3mm}

Un problema formado por una ecuación en derivadas parciales en un
dominio \(D\) junto con condiciones sobre la frontera de \(D\)
---espacial, temporal, o ambas--- está \textbf{bien planteado} en el
sentido de Hadamard si cumple las tres cosas:

\begin{enumerate}
\def\labelenumi{\arabic{enumi}.}
\tightlist
\item
  \textbf{Existencia.} Hay al menos una solución.
\item
  \textbf{Unicidad.} Hay a lo sumo una solución.
\item
  \textbf{Estabilidad.} La solución depende de manera continua de los
  datos: si los datos cambian poco, la solución cambia poco.
\end{enumerate}

Si falla cualquiera de las tres, el problema está \textbf{mal
planteado}.

\end{minipage}%
\end{tcolorbox}

La tercera es la que se olvida y la que decide si el problema sirve. Los
datos de un problema físico siempre tienen error; si el problema los
amplifica, la solución calculada no se parece a la que ocurre, y que
exista y sea única no ayuda en nada.

\textbf{El problema de Hadamard.} Tomando las condiciones de
Cauchy--Riemann

\[
\frac{\partial u}{\partial y} = \frac{\partial v}{\partial x}, \qquad
\frac{\partial v}{\partial y} = -\frac{\partial u}{\partial x} ,
\]

el par

\[
u_n = e^{-\sqrt{n}}\,e^{ny}\cos(nx), \qquad v_n = e^{-\sqrt{n}}\,e^{ny}\operatorname{sen}(nx)
\]

las satisface. Sobre \(y=0\) el dato es \(e^{-\sqrt n}\cos(nx)\), que
junto con todas sus derivadas tiende a cero cuando \(n\to\infty\); pero
para cualquier \(y>0\) fijo la amplitud \(e^{ny-\sqrt n}\) tiende a
infinito. El dato nulo tiene solución nula. \textbf{Hay solución, es
única, y el problema está mal planteado}, porque no depende
continuamente de los datos.

\subsubsection{Las condiciones de
contorno}\label{las-condiciones-de-contorno}

\begin{tcolorbox}[enhanced jigsaw, arc=.35mm, bottomrule=.15mm, breakable, colback=white, colframe=quarto-callout-note-color-frame, left=2mm, leftrule=.75mm, opacityback=0, rightrule=.15mm, toprule=.15mm]
\begin{minipage}[t]{5.5mm}
\textcolor{quarto-callout-note-color}{\faInfo}
\end{minipage}%
\begin{minipage}[t]{\textwidth - 5.5mm}

\vspace{-3mm}\textbf{Definición~\ref{def-condiciones} --- Dirichlet, Neumann y Robin}\vspace{3mm}

\begin{itemize}
\tightlist
\item
  \textbf{Condición de Dirichlet:} se prescribe el valor de la función
  sobre la frontera, \(u\big|_S = h\).
\item
  \textbf{Condición de Neumann:} se prescribe la derivada normal,
  \(\dfrac{\partial u}{\partial n}\Big|_S = h\). Un borde aislado es
  \(\partial u/\partial n = 0\).
\item
  \textbf{Condición de Robin} o \textbf{mixta:} una combinación de las
  dos, \(\dfrac{\partial u}{\partial n} + \alpha u\Big|_S = h\).
\end{itemize}

Una condición se llama \textbf{homogénea} si \(h \equiv 0\).

\end{minipage}%
\end{tcolorbox}

Cuando se prescriben \textbf{a la vez} el valor y la derivada normal
sobre una misma curva, se habla de \textbf{datos de Cauchy}. Es lo que
se le dio al sistema del problema de Hadamard, y es lo que lo arruina.

\subsection{Ecuaciones de primer orden: el método de las
características}\label{sec-resumen-4}

\subsubsection{La superficie solución}\label{la-superficie-soluciuxf3n}

\begin{tcolorbox}[enhanced jigsaw, arc=.35mm, bottomrule=.15mm, breakable, colback=white, colframe=quarto-callout-note-color-frame, left=2mm, leftrule=.75mm, opacityback=0, rightrule=.15mm, toprule=.15mm]
\begin{minipage}[t]{5.5mm}
\textcolor{quarto-callout-note-color}{\faInfo}
\end{minipage}%
\begin{minipage}[t]{\textwidth - 5.5mm}

\vspace{-3mm}\textbf{Definición~\ref{def-superficie-solucion} --- Superficie solución}\vspace{3mm}

Dada una EDP de primer orden en dos variables, una \textbf{superficie
solución} es el gráfico \[
S = \{(x,y,u(x,y)) : (x,y) \in D\}
\] de una solución \(u\) definida en un dominio \(D\) del plano.

Como \(S\) es la superficie de nivel cero de
\(\Phi(x,y,z) = z - u(x,y)\), un vector normal a \(S\) en cada punto es
\[
\vec n = \nabla\Phi = \big(-u_x,\, -u_y,\, 1\big),
\] o, cambiándole el signo, \(\big(u_x, u_y, -1\big)\).

\end{minipage}%
\end{tcolorbox}

Para la ecuación cuasi-lineal general en dos variables

\[
a(x,y,u)\,u_x + b(x,y,u)\,u_y = c(x,y,u) ,
\]

armando el campo \(\vec V = (a,b,c)\) resulta
\(\vec V\cdot(u_x,u_y,-1) = a\,u_x + b\,u_y - c\), así que \textbf{la
ecuación dice exactamente que \(\vec V\) es tangente a la superficie
solución} en cada uno de sus puntos.

\subsubsection{El campo que guía}\label{el-campo-que-guuxeda}

\begin{tcolorbox}[enhanced jigsaw, arc=.35mm, bottomrule=.15mm, breakable, colback=white, colframe=quarto-callout-note-color-frame, left=2mm, leftrule=.75mm, opacityback=0, rightrule=.15mm, toprule=.15mm]
\begin{minipage}[t]{5.5mm}
\textcolor{quarto-callout-note-color}{\faInfo}
\end{minipage}%
\begin{minipage}[t]{\textwidth - 5.5mm}

\vspace{-3mm}\textbf{Definición~\ref{def-campo-integral} --- Campo vectorial y curva integral}\vspace{3mm}

Un \textbf{campo vectorial} en un dominio \(D \subseteq \mathbb{R}^2\)
es una función que a cada punto le asigna un vector, \[
\vec v : D \to \mathbb{R}^2, \qquad \vec v(x,y) = \big(v_1(x,y),\, v_2(x,y)\big) .
\]

Una \textbf{curva integral} del campo es una curva
\(s \mapsto \big(x(s), y(s)\big)\) que en cada punto avanza en la
dirección que marca el campo, y con esa misma rapidez: \[
\frac{dx}{ds} = v_1\big(x(s), y(s)\big), \qquad
\frac{dy}{ds} = v_2\big(x(s), y(s)\big) .
\]

A la familia de todas las curvas integrales se la llama \textbf{flujo}
del campo, y a su dibujo, \textbf{retrato de fases}.

\end{minipage}%
\end{tcolorbox}

Las tres lecturas que hay que tener:

\begin{enumerate}
\def\labelenumi{\arabic{enumi}.}
\tightlist
\item
  \textbf{En el piso.} Si \(a\) y \(b\) dependen solo de \((x,y)\)
  ---caso lineal y semilineal---, el par \(\vec v = (a,b)\) es un campo
  vectorial sobre el dominio, y \textbf{la característica proyectada es
  una curva integral de \(\vec v\)}. Las rutas se pueden dibujar antes
  de resolver nada.
\item
  \textbf{La ecuación es una derivada direccional.} Como
  \(a\,u_x + b\,u_y = \vec v\cdot\nabla u\), la ecuación
  \(a\,u_x + b\,u_y = c\) dice que la derivada de \(u\) en la dirección
  del campo vale \(c\). Con \(c=0\): \(u\) es constante sobre las curvas
  integrales, que es la definición de primera integral.
\item
  \textbf{En el espacio.} El campo completo es \(\vec V = (a,b,c)\),
  cuya sombra sobre el piso es \(\vec v\). La tercera componente \(c\)
  es el \textbf{ritmo de ascenso del hilo}.
\end{enumerate}

\begin{longtable}[]{@{}ll@{}}
\caption{Qué hace la tercera componente. El piso es el mismo en las tres
filas.}\tabularnewline
\toprule\noalign{}
Si \(c\)\ldots{} & el hilo\ldots{} \\
\midrule\noalign{}
\endfirsthead
\toprule\noalign{}
Si \(c\)\ldots{} & el hilo\ldots{} \\
\midrule\noalign{}
\endhead
\bottomrule\noalign{}
\endlastfoot
es \(0\) & va a altura constante \\
depende solo de \((x,y)\) & sube a un ritmo prefijado en cada punto \\
depende de \(u\) & sube a un ritmo que depende de su propia altura \\
\end{longtable}

\textbf{En las cuasi-lineales no hay campo en el piso.} Si \(a\) o \(b\)
dependen de \(u\), la flecha del punto \((x,y)\) depende de a qué altura
estés: el campo existe pero vive en el espacio y no se puede bajar. Esa
es la razón de fondo de que el caso cuasi-lineal sea más difícil y de
que ahí aparezcan los choques.

\subsubsection{El sistema
característico}\label{el-sistema-caracteruxedstico}

Si un observador recorre el plano según \(\big(x(s),y(s)\big)\) y va
anotando \(\varphi(s) = u(x(s),y(s))\), la regla de la cadena da
\(\varphi' = x'\,u_x + y'\,u_y\). Eligiendo \(x' = a\) e \(y' = b\), la
ecuación dice que \(\varphi' = c\), y no queda ninguna derivada parcial:

\[
\boxed{\;\frac{dx}{ds} = a(x,y,u), \qquad
\frac{dy}{ds} = b(x,y,u), \qquad
\frac{du}{ds} = c(x,y,u).\;}
\]

\begin{tcolorbox}[enhanced jigsaw, arc=.35mm, bottomrule=.15mm, breakable, colback=white, colframe=quarto-callout-note-color-frame, left=2mm, leftrule=.75mm, opacityback=0, rightrule=.15mm, toprule=.15mm]
\begin{minipage}[t]{5.5mm}
\textcolor{quarto-callout-note-color}{\faInfo}
\end{minipage}%
\begin{minipage}[t]{\textwidth - 5.5mm}

\vspace{-3mm}\textbf{Definición~\ref{def-sistema-caracteristico} --- Sistema característico,
curva característica y característica proyectada}\vspace{3mm}

Al sistema anterior se lo llama \textbf{sistema característico} de la
ecuación. También se lo escribe, eliminando el parámetro, como

\[
\frac{dx}{a} = \frac{dy}{b} = \frac{du}{c} .
\]

Sus soluciones \(\big(x(s),y(s),u(s)\big)\) son curvas en el espacio
\((x,y,u)\) y se llaman \textbf{curvas características} de la ecuación:
son, por construcción, las curvas integrales del campo \(\vec V\).

La sombra de una curva característica sobre el plano \((x,y)\) ---o sea
el par \(\big(x(s),y(s)\big)\) sin la altura--- se llama
\textbf{característica proyectada}, o característica base. Es la
trayectoria del observador.

\end{minipage}%
\end{tcolorbox}

Como por cada punto del espacio pasa una sola curva característica,
\textbf{la superficie solución está formada por curvas características}:
se la construye largando una por cada punto de la curva inicial.

\subsubsection{Coeficientes constantes: la onda
simple}\label{coeficientes-constantes-la-onda-simple}

Para \(u_t + c\,u_x = 0\) el sistema característico es \(dx/dt = c\),
\(du/dt = 0\). Las características proyectadas son las rectas
\(x - ct = \text{const}\), sobre ellas \(u\) es constante, y con el dato
\(u(x,0)=f(x)\) resulta

\[
u(x,t) = f(x - ct) :
\]

el dato inicial transportado sin deformarse una distancia \(ct\) (onda
no dispersiva).

\subsubsection{Coeficientes variables}\label{coeficientes-variables}

Para \(\displaystyle\sum_{i=1}^{n} a_i(x)\,u_{x_i} = 0\) las
características se obtienen de

\[
\frac{dx_i}{ds} = a_i(x_1,\dots,x_n) ,
\]

y a lo largo de ellas \(du/ds = 0\): la solución no depende de \(s\).

\begin{tcolorbox}[enhanced jigsaw, arc=.35mm, bottomrule=.15mm, breakable, colback=white, colframe=quarto-callout-note-color-frame, left=2mm, leftrule=.75mm, opacityback=0, rightrule=.15mm, toprule=.15mm]
\begin{minipage}[t]{5.5mm}
\textcolor{quarto-callout-note-color}{\faInfo}
\end{minipage}%
\begin{minipage}[t]{\textwidth - 5.5mm}

\vspace{-3mm}\textbf{Definición~\ref{def-primera-integral} --- Primera integral}\vspace{3mm}

Una función \(\phi\) de clase \(C^1\), no constante, es una
\textbf{primera integral} del campo si es constante a lo largo de cada
curva integral del campo. Derivando \(\phi\) respecto de \(s\) y usando
el sistema característico, la condición se escribe

\[
\sum_i a_i\,\phi_{x_i} = 0 ,
\]

es decir: \textbf{\(\phi\) es primera integral si y solo si \(\phi\) es
solución de la ecuación.}

\end{minipage}%
\end{tcolorbox}

\begin{tcolorbox}[enhanced jigsaw, arc=.35mm, bottomrule=.15mm, breakable, colback=white, colframe=quarto-callout-note-color-frame, left=2mm, leftrule=.75mm, opacityback=0, rightrule=.15mm, toprule=.15mm]
\begin{minipage}[t]{5.5mm}
\textcolor{quarto-callout-note-color}{\faInfo}
\end{minipage}%
\begin{minipage}[t]{\textwidth - 5.5mm}

\vspace{-3mm}\textbf{Teorema~\ref{thm-general-homogenea} --- Solución general de la homogénea}\vspace{3mm}

Si \(\phi\) es una primera integral del campo \((a,b)\) en un dominio
\(D\), entonces

\[
u(x,y) = F\big(\phi(x,y)\big)
\]

es solución de \(a\,u_x + b\,u_y = 0\) en \(D\) para \textbf{cualquier}
función \(F\) de una variable de clase \(C^1\). Recíprocamente, toda
solución tiene esa forma.

\end{minipage}%
\end{tcolorbox}

Es el análogo exacto de «la solución general de una EDO de primer orden
tiene una constante arbitraria», con una diferencia importante: acá lo
arbitrario no es una constante sino una \textbf{función entera}. Por eso
una sola condición en un punto nunca alcanza, y hace falta un dato sobre
toda una curva.

\subsubsection{Casi--lineales, caso
homogéneo}\label{casilineales-caso-homoguxe9neo}

Para \(u_t + a(x,t,u)\,u_x = 0\) con \(u(x,0)=f(x)\), el sistema
característico es

\[
\frac{dx}{dt} = a(x,t,u), \qquad \frac{du}{dt} = 0 ,
\]

así que \(u\) sigue siendo constante sobre cada característica, pero
\textbf{la velocidad de avance depende del valor transportado}. Para la
ecuación de Burgers sin viscosidad, \(a(u)=u\):

\[
u_t + u\,u_x = 0 \qquad\Longrightarrow\qquad
x = x_0 + f(x_0)\,t, \quad u\big(x_0+f(x_0)t,\;t\big) = f(x_0) ,
\]

una solución de carácter implícito.

\subsubsection{Casi--lineales, caso
general}\label{casilineales-caso-general}

Para \(a_1(x,y,u)\,u_x + a_2(x,y,u)\,u_y = a(x,y,u)\), la dirección
\(\big(a_1,a_2,a\big)\) es tangente al gráfico de la solución, de modo
que \textbf{las características permanecen dentro de la superficie
solución o no la tocan}.

\begin{tcolorbox}[enhanced jigsaw, arc=.35mm, bottomrule=.15mm, breakable, colback=white, colframe=quarto-callout-note-color-frame, left=2mm, leftrule=.75mm, opacityback=0, rightrule=.15mm, toprule=.15mm]
\begin{minipage}[t]{5.5mm}
\textcolor{quarto-callout-note-color}{\faInfo}
\end{minipage}%
\begin{minipage}[t]{\textwidth - 5.5mm}

\vspace{-3mm}\textbf{Teorema~\ref{thm-transversalidad} --- Existencia y unicidad local}\vspace{3mm}

Sean \(a_1\), \(a_2\) y \(a\) de clase \(C^1\) en un dominio de
\(\mathbb{R}^3\) que contiene al punto
\(\big(x_0(\sigma_0),\,y_0(\sigma_0),\,u_0(\sigma_0)\big)\), y sea
\(\Gamma\) continuamente derivable. Si en ese punto

\[
\begin{vmatrix}
a_1 & x_0'(\sigma_0)\\[2pt]
a_2 & y_0'(\sigma_0)
\end{vmatrix}
= a_1\,y_0'(\sigma_0) - a_2\,x_0'(\sigma_0) \;\neq\; 0 ,
\]

entonces existe una única solución de clase \(C^1\) en un entorno de
\(\big(x_0(\sigma_0), y_0(\sigma_0)\big)\) que toma sobre \(\Gamma\) los
valores \(u_0\).

La condición se llama \textbf{condición de transversalidad}, y dice
exactamente que la proyección del campo \(\vec V=(a_1,a_2,a)\) sobre el
plano \((x,y)\) no es tangente a la curva inicial en ese punto.

\end{minipage}%
\end{tcolorbox}

\subsubsection{Cuándo el método
falla}\label{cuuxe1ndo-el-muxe9todo-falla}

Cada hipótesis tapa un agujero concreto. Conviene saber cuál tapa cuál,
porque los ejercicios de parcial suelen ser exactamente uno de estos
cinco casos.

\begin{longtable}[]{@{}
  >{\raggedright\arraybackslash}p{(\linewidth - 4\tabcolsep) * \real{0.3333}}
  >{\raggedright\arraybackslash}p{(\linewidth - 4\tabcolsep) * \real{0.3333}}
  >{\raggedright\arraybackslash}p{(\linewidth - 4\tabcolsep) * \real{0.3333}}@{}}
\caption{Las hipótesis del método y el agujero que tapa cada
una.}\tabularnewline
\toprule\noalign{}
\begin{minipage}[b]{\linewidth}\raggedright
Hipótesis
\end{minipage} & \begin{minipage}[b]{\linewidth}\raggedright
Qué garantiza
\end{minipage} & \begin{minipage}[b]{\linewidth}\raggedright
Qué pasa si falla
\end{minipage} \\
\midrule\noalign{}
\endfirsthead
\toprule\noalign{}
\begin{minipage}[b]{\linewidth}\raggedright
Hipótesis
\end{minipage} & \begin{minipage}[b]{\linewidth}\raggedright
Qué garantiza
\end{minipage} & \begin{minipage}[b]{\linewidth}\raggedright
Qué pasa si falla
\end{minipage} \\
\midrule\noalign{}
\endhead
\bottomrule\noalign{}
\endlastfoot
\(a\), \(b\), \(c\) de clase \(C^1\) & por cada punto pasa un solo hilo
& los hilos se bifurcan: \(u\) no queda determinada \\
\((a,b) \neq (0,0)\) & el observador efectivamente avanza & el punto es
singular del campo; las sombras se juntan ahí \\
transversalidad: \(a\,y_0' - b\,x_0' \neq 0\) & los hilos cruzan
\(\Gamma\) y barren una superficie & o no hay solución, o hay
infinitas \\
la conclusión es \textbf{local} & la superficie es el gráfico de una
función & las sombras se cruzan: choque, y hay que pasar a soluciones
generalizadas \\
\(\Gamma\) toca todas las características de la región & el dato alcanza
para todo el dominio & quedan zonas del piso sin tejer \\
\end{longtable}

Dos precisiones que se preguntan:

\begin{itemize}
\tightlist
\item
  \textbf{Los hilos no se cruzan; sus sombras sí pueden.} En el espacio
  la unicidad sigue valiendo y las curvas características nunca se
  tocan. Lo que se cruza son sus proyecciones, porque los dos hilos van
  a alturas distintas. Por eso en el caso lineal y semilineal el choque
  es imposible ---ahí las sombras son curvas integrales de un campo del
  piso--- y en el cuasi-lineal es lo típico.
\item
  \textbf{Cuando falla la transversalidad, cuál de las dos cosas pasa lo
  decide el dato.} Si el dato sobre \(\Gamma\) no es el que la ecuación
  misma produciría a lo largo de \(\Gamma\), no hay solución; si
  coincide, hay infinitas.
\end{itemize}

\textbf{Hamilton--Jacobi.} Hay ecuaciones de primer orden importantes
que \textbf{no} son cuasi-lineales. La más notable es

\[
u_t + H\big(u_x, u_y, x, y, t\big) = 0 ,
\]

que juega un papel central en la mecánica clásica y cuya generalización
cuántica es el puente entre las teorías corpuscular y ondulatoria de la
materia. También se resuelve por características, con un tratamiento
propio.

\subsection{Segundo orden semilineales en dos
variables}\label{sec-resumen-5}

La forma general es

\[
A(x,y)\,u_{xx} + 2B(x,y)\,u_{xy} + C(x,y)\,u_{yy} = F(x,y,u,u_x,u_y) .
\]

\begin{tcolorbox}[enhanced jigsaw, arc=.35mm, bottomrule=.15mm, breakable, colback=white, colframe=quarto-callout-note-color-frame, left=2mm, leftrule=.75mm, opacityback=0, rightrule=.15mm, toprule=.15mm]
\begin{minipage}[t]{5.5mm}
\textcolor{quarto-callout-note-color}{\faInfo}
\end{minipage}%
\begin{minipage}[t]{\textwidth - 5.5mm}

\vspace{-3mm}\textbf{Definición~\ref{def-discriminante} --- Discriminante}\vspace{3mm}

Se llama \textbf{discriminante} de la ecuación a
\(\;\Delta = B^2 - AC\).

\end{minipage}%
\end{tcolorbox}

\begin{tcolorbox}[enhanced jigsaw, arc=.35mm, bottomrule=.15mm, breakable, colback=white, colframe=quarto-callout-note-color-frame, left=2mm, leftrule=.75mm, opacityback=0, rightrule=.15mm, toprule=.15mm]
\begin{minipage}[t]{5.5mm}
\textcolor{quarto-callout-note-color}{\faInfo}
\end{minipage}%
\begin{minipage}[t]{\textwidth - 5.5mm}

\vspace{-3mm}\textbf{Teorema~\ref{thm-clasificacion} --- Clasificación}\vspace{3mm}

En un entorno de un punto donde \(A\), \(B\) y \(C\) son continuas, la
ecuación se clasifica según el signo de \(\Delta = B^2 - AC\):

{\def\LTcaptype{none} % do not increment counter
\begin{longtable}[]{@{}llll@{}}
\toprule\noalign{}
\(\Delta\) & Tipo & Características reales & Ecuación modelo \\
\midrule\noalign{}
\endhead
\bottomrule\noalign{}
\endlastfoot
\(>0\) & hiperbólica & dos familias & \(u_{tt} - c^2 u_{xx} = 0\) \\
\(=0\) & parabólica & una familia & \(u_t - k\,u_{xx} = 0\) \\
\(<0\) & elíptica & ninguna & \(u_{xx} + u_{yy} = 0\) \\
\end{longtable}
}

\end{minipage}%
\end{tcolorbox}

Si la ecuación es de un tipo en todo punto de un dominio, se dice que es
de ese tipo en el dominio. \textbf{La naturaleza de las soluciones está
ligada al tipo, y el tipo a la existencia y al número de características
reales.} Los tres ejemplos con nombre: las ondas acústicas en un gas
isótropo y homogéneo son hiperbólicas, la difusión del calor es
parabólica, y la ecuación de Laplace es elíptica.

\subsubsection{El cambio de variables}\label{el-cambio-de-variables}

Con \(\xi = \xi(x,y)\), \(\psi = \psi(x,y)\) invertible, la ecuación se
transforma en

\[
\alpha\,u_{\xi\xi} + \beta\,u_{\xi\psi} + \gamma\,u_{\psi\psi} = H(\xi,\psi,u,u_\xi,u_\psi),
\]

con

\[
\begin{aligned}
\alpha &= A\,\xi_x^2 + 2B\,\xi_x\xi_y + C\,\xi_y^2, \\
\beta  &= 2\big[A\,\xi_x\psi_x + B(\xi_x\psi_y + \xi_y\psi_x) + C\,\xi_y\psi_y\big], \\
\gamma &= A\,\psi_x^2 + 2B\,\psi_x\psi_y + C\,\psi_y^2 .
\end{aligned}
\]

Anular \(\alpha\) o \(\gamma\) equivale a pedir
\(A\phi_x^2 + 2B\phi_x\phi_y + C\phi_y^2 = 0\), cuyas raíces son

\[
\frac{\phi_x}{\phi_y} = \frac{-B \pm \sqrt{B^2-AC}}{A} ,
\]

y \textbf{el radicando es el discriminante}: hay dos direcciones reales
distintas si \(\Delta>0\), una doble si \(\Delta=0\), y ninguna real si
\(\Delta<0\).

\begin{tcolorbox}[enhanced jigsaw, arc=.35mm, bottomrule=.15mm, breakable, colback=white, colframe=quarto-callout-note-color-frame, left=2mm, leftrule=.75mm, opacityback=0, rightrule=.15mm, toprule=.15mm]
\begin{minipage}[t]{5.5mm}
\textcolor{quarto-callout-note-color}{\faInfo}
\end{minipage}%
\begin{minipage}[t]{\textwidth - 5.5mm}

\vspace{-3mm}\textbf{Definición~\ref{def-curvas-caracteristicas-2} --- Curvas características
de una ecuación de segundo orden}\vspace{3mm}

Las soluciones de la ecuación diferencial ordinaria

\[
A\left(\frac{dy}{dx}\right)^2 - 2B\,\frac{dy}{dx} + C = 0
\]

se llaman \textbf{curvas características} de la ecuación. Cuando
\(\Delta>0\) hay dos familias reales; si \(F_1(x,y) = \text{const}\) y
\(F_2(x,y) = \text{const}\) las describen, el cambio de variables
\(\xi = F_1\), \(\psi = F_2\) lleva la ecuación a la forma canónica
hiperbólica \(u_{\xi\psi} = G(\xi,\psi,u,u_\xi,u_\psi)\).

\end{minipage}%
\end{tcolorbox}

\subsubsection{Las tres formas
canónicas}\label{las-tres-formas-canuxf3nicas}

\begin{longtable}[]{@{}
  >{\raggedright\arraybackslash}p{(\linewidth - 6\tabcolsep) * \real{0.2500}}
  >{\raggedright\arraybackslash}p{(\linewidth - 6\tabcolsep) * \real{0.2500}}
  >{\raggedright\arraybackslash}p{(\linewidth - 6\tabcolsep) * \real{0.2500}}
  >{\raggedright\arraybackslash}p{(\linewidth - 6\tabcolsep) * \real{0.2500}}@{}}
\caption{La receta. En los tres casos el primer paso es el mismo:
resolver la EDO de las características.}\tabularnewline
\toprule\noalign{}
\begin{minipage}[b]{\linewidth}\raggedright
Tipo
\end{minipage} & \begin{minipage}[b]{\linewidth}\raggedright
Se anula
\end{minipage} & \begin{minipage}[b]{\linewidth}\raggedright
Cómo se elige el cambio
\end{minipage} & \begin{minipage}[b]{\linewidth}\raggedright
Forma canónica
\end{minipage} \\
\midrule\noalign{}
\endfirsthead
\toprule\noalign{}
\begin{minipage}[b]{\linewidth}\raggedright
Tipo
\end{minipage} & \begin{minipage}[b]{\linewidth}\raggedright
Se anula
\end{minipage} & \begin{minipage}[b]{\linewidth}\raggedright
Cómo se elige el cambio
\end{minipage} & \begin{minipage}[b]{\linewidth}\raggedright
Forma canónica
\end{minipage} \\
\midrule\noalign{}
\endhead
\bottomrule\noalign{}
\endlastfoot
Hiperbólica & \(\alpha = \gamma = 0\) & \(\xi\) y \(\psi\) constantes
sobre \textbf{cada} familia de características & \(u_{\xi\psi} = G\) \\
Parabólica & \(\beta = \gamma = 0\) & \(\psi\) sobre la \textbf{única}
familia; \(\xi\) libre e independiente & \(u_{\xi\xi} = G\) \\
Elíptica & \(\beta=0\), \(\alpha=\gamma\) & \(\sigma\) y \(\tau\): parte
real e imaginaria de las características complejas &
\(u_{\sigma\sigma}+u_{\tau\tau} = G\) \\
\end{longtable}

\textbf{Detalles de cada caso.}

\begin{itemize}
\tightlist
\item
  \emph{Hiperbólicas.} Las dos condiciones son
  \(\xi_x + \dfrac{B-\sqrt{B^2-AC}}{A}\,\xi_y = 0\) y
  \(\psi_x + \dfrac{B+\sqrt{B^2-AC}}{A}\,\psi_y = 0\), cada una una EDP
  de primer orden que se resuelve por características. Hay una segunda
  forma canónica usual, \(u_{\sigma\sigma} - u_{\tau\tau} = K\), que se
  obtiene con \(\sigma = \xi+\psi\), \(\tau = \xi-\psi\).
\item
  \emph{Parabólicas.} Con \(A\neq0\), la única dirección es
  \(\psi_x/\psi_y = -B/A\), y entonces \(\beta\) se anula solo: \(\xi\)
  se elige libremente, con la única condición de ser independiente de
  \(\psi\). Si \(A = 0\), la condición pasa a ser
  \(\psi_x + (B/C)\psi_y = 0\).
\item
  \emph{Elípticas.} Las direcciones son complejas conjugadas,
  \(\dfrac{\xi_x}{\xi_y} = \dfrac{-B \pm i\sqrt{AC-B^2}}{A}\), y se toma
  \(\sigma = \dfrac{\xi+\psi}{2}\), \(\tau = \dfrac{\xi-\psi}{2i}\) para
  volver a variables reales.
\end{itemize}

\textbf{Un caso que conviene saber de memoria.} Para
\(u_{tt} - c^2u_{xx} = s(x,t)\), las características son
\(x \pm ct = \text{const}\); con \(\xi = x+ct\) y \(\psi = x-ct\) la
ecuación queda \(u_{\xi\psi} = -\,s/(4c^2)\), y cuando \(s\equiv0\) se
obtiene la solución conocida

\[
u(x,t) = f(x+ct) + g(x-ct),
\]

con \(f\) y \(g\) arbitrarias de clase \(C^2\).

\subsection{Segundo orden semilineales en muchas
variables}\label{sec-resumen-6}

La parte principal es

\[
L_P[u] = \sum_{i=1}^{n}\sum_{j=1}^{n} a_{ij}(x)\,u_{x_ix_j} ,
\]

con coeficientes que se pueden suponer \textbf{simétricos},
\(a_{ij}=a_{ji}\), porque \(u_{x_ix_j} = u_{x_jx_i}\). Con el cambio
lineal

\[
y_i = \sum_{j=1}^{n} c_{ij}\,\big(x_j - x_j^0\big)
\]

la parte principal se transforma en
\(M[u] = \sum_{k,i} A_{ki}\,u_{y_iy_k}\) con
\(A_{ik} = \sum_{j,l} c_{ij}\,c_{kl}\,a_{jl}\), y esa matriz se puede
\textbf{diagonalizar} en el punto de modo que los términos de la
diagonal sean \(-1\), \(1\) o \(0\).

\begin{tcolorbox}[enhanced jigsaw, arc=.35mm, bottomrule=.15mm, breakable, colback=white, colframe=quarto-callout-note-color-frame, left=2mm, leftrule=.75mm, opacityback=0, rightrule=.15mm, toprule=.15mm]
\begin{minipage}[t]{5.5mm}
\textcolor{quarto-callout-note-color}{\faInfo}
\end{minipage}%
\begin{minipage}[t]{\textwidth - 5.5mm}

\vspace{-3mm}\textbf{Definición~\ref{def-signatura} --- Signatura}\vspace{3mm}

Diagonalizada la matriz de coeficientes en un punto, llamamos

\begin{itemize}
\tightlist
\item
  \(n_+\) = cantidad de unos positivos,
\item
  \(n_-\) = cantidad de unos negativos,
\item
  \(n_0\) = cantidad de ceros,
\end{itemize}

con \(n_+ + n_- + n_0 = n\). La terna \((n_+, n_-, n_0)\) se llama
\textbf{signatura} de la forma cuadrática en ese punto, y no depende de
cómo se haya diagonalizado.

\end{minipage}%
\end{tcolorbox}

\begin{tcolorbox}[enhanced jigsaw, arc=.35mm, bottomrule=.15mm, breakable, colback=white, colframe=quarto-callout-note-color-frame, left=2mm, leftrule=.75mm, opacityback=0, rightrule=.15mm, toprule=.15mm]
\begin{minipage}[t]{5.5mm}
\textcolor{quarto-callout-note-color}{\faInfo}
\end{minipage}%
\begin{minipage}[t]{\textwidth - 5.5mm}

\vspace{-3mm}\textbf{Definición~\ref{def-clasificacion-n} --- Clasificación en \(n\)
variables}\vspace{3mm}

Sea una ecuación semilineal de segundo orden en un entorno de
\(x^0 \in \mathbb{R}^n\). En \(x^0\) se dice

\begin{itemize}
\tightlist
\item
  \textbf{elíptica} si \(n_+ = n\) o \(n_- = n\) (todos los signos
  iguales, ningún cero);
\item
  \textbf{parabólica} si \(n_0 > 0\) (hay al menos un cero);
\item
  \textbf{hiperbólica} si \(n_0 = 0\) y \(n_+ = 1\) o \(n_- = 1\)
  (ningún cero y un solo signo distinto);
\item
  \textbf{ultrahiperbólica} si \(n_0 = 0\), \(n_+ > 1\) y \(n_- > 1\).
\end{itemize}

\end{minipage}%
\end{tcolorbox}

\subsubsection{La dificultad que aparece a partir de tres
variables}\label{la-dificultad-que-aparece-a-partir-de-tres-variables}

Todo lo anterior vale \textbf{en un punto}. Diagonalizar en todo un
\textbf{entorno} exige anular los \(n(n-1)/2\) elementos de fuera de la
diagonal, y solo hay \(n\) funciones libres ---una por coordenada
nueva--- para conseguirlo:

{\def\LTcaptype{none} % do not increment counter
\begin{longtable}[]{@{}ccc@{}}
\toprule\noalign{}
\(n\) & funciones disponibles & condiciones a imponer \\
\midrule\noalign{}
\endhead
\bottomrule\noalign{}
\endlastfoot
\(2\) & \(2\) & \(1\) \\
\(3\) & \(3\) & \(3\) \\
\(4\) & \(4\) & \(6\) \\
\(n\) & \(n\) & \(n(n-1)/2\) \\
\end{longtable}
}

Para \(n > 3\) el problema queda \textbf{sobredeterminado}. Para \(n=3\)
se puede diagonalizar, pero los elementos de la diagonal pueden salir
distintos en distintos puntos del entorno. En consecuencia \textbf{no se
pueden reducir a forma canónica en un entorno los problemas de dimensión
\(n\ge3\)} ---por ejemplo, más de una dimensión espacial más el
tiempo---, salvo que los coeficientes sean constantes. Para \(n=2\), en
cambio, el problema está bien determinado aunque los coeficientes sean
variables, y por eso ahí sí se consiguen las transformaciones de la
sección anterior.

\subsection{Transformadas integrales}\label{sec-resumen-7}

El problema de partida, en un volumen acotado \(D\) de superficie
cerrada \(S\):

\[
\nabla^2 u = a\,u + b\,\frac{\partial u}{\partial t} + c\,\frac{\partial^2 u}{\partial t^2} - f(x,y,z,t) ,
\]

con \textbf{condiciones de contorno} sobre \(S\), independientes de
\(t\), del tipo Dirichlet, Neumann o Robin,

\[
\frac{\partial u}{\partial n} + \alpha\,u = 0 ,
\]

y \textbf{condiciones iniciales} en \(D\),

\[
u(x,y,z,0) = u_0(x,y,z), \qquad
\left.\frac{\partial u}{\partial t}\right|_{t=0} = v_0(x,y,z) ,
\]

asumiendo \(f\) y \(u\) nulas para \(t<0\).

Definiendo las transformadas en el tiempo

\[
\Phi(x,y,z,s) = \frac{1}{\sqrt{2\pi}}\int_0^{\infty} u\,e^{-ist}\,dt, \qquad
F(x,y,z,s) = \frac{1}{\sqrt{2\pi}}\int_0^{\infty} f\,e^{-ist}\,dt ,
\]

---o, si se prefiere, la variante equivalente con transformada de
Laplace--- el laplaciano atraviesa la integral sin cambiar, las
derivadas temporales se vuelven factores y arrastran las condiciones
iniciales, y \textbf{el tiempo desaparece}: queda, para cada \(s\) por
separado, una ecuación en las variables espaciales solamente. La
condición de contorno, al ser independiente de \(t\), se transforma en
sí misma.

\begin{tcolorbox}[enhanced jigsaw, arc=.35mm, bottomrule=.15mm, breakable, colback=white, colframe=quarto-callout-note-color-frame, left=2mm, leftrule=.75mm, opacityback=0, rightrule=.15mm, toprule=.15mm]
\begin{minipage}[t]{5.5mm}
\textcolor{quarto-callout-note-color}{\faInfo}
\end{minipage}%
\begin{minipage}[t]{\textwidth - 5.5mm}

\vspace{-3mm}\textbf{Definición~\ref{def-helmholtz} --- Problema de Helmholtz}\vspace{3mm}

Se llama \textbf{problema de Helmholtz} al de hallar \(\Phi\) en un
dominio \(D\) tal que

\[
\nabla^2\Phi + \lambda\,\Phi = g \quad\text{en } D, \qquad
\frac{\partial\Phi}{\partial n} + \alpha\,\Phi = 0 \quad\text{sobre } S .
\]

\end{minipage}%
\end{tcolorbox}

\textbf{Lo que hay que retener:} transformar en el tiempo convierte el
problema original en un problema de Helmholtz. Ese problema es la
generalización a varias dimensiones del problema de autovalores de
Sturm--Liouville en una dimensión, y una vez resuelto queda
antitransformar para volver.

\subsection{Función de Green en varias variables}\label{sec-resumen-8}

\subsubsection{Caso 1: ecuación inhomogénea, contorno
homogéneo}\label{caso-1-ecuaciuxf3n-inhomoguxe9nea-contorno-homoguxe9neo}

\[
L[u(\mathbf{x})] = f(\mathbf{x}) \ \text{ en } D, \qquad
B[u(\mathbf{x})] = 0 \ \text{ en } S .
\]

\begin{tcolorbox}[enhanced jigsaw, arc=.35mm, bottomrule=.15mm, breakable, colback=white, colframe=quarto-callout-note-color-frame, left=2mm, leftrule=.75mm, opacityback=0, rightrule=.15mm, toprule=.15mm]
\begin{minipage}[t]{5.5mm}
\textcolor{quarto-callout-note-color}{\faInfo}
\end{minipage}%
\begin{minipage}[t]{\textwidth - 5.5mm}

\vspace{-3mm}\textbf{Definición~\ref{def-green-varias} --- Función de Green}\vspace{3mm}

La \textbf{función de Green} del problema es la función
\(G(\mathbf{x},\xi)\) que satisface, para cada \(\xi \in D\) fijo,

\[
L\big[G(\mathbf{x},\xi)\big] = \delta(\mathbf{x}-\xi) \ \text{ en } D, \qquad
B\big[G(\mathbf{x},\xi)\big] = 0 \ \text{ en } S .
\]

Es decir: \textbf{\(G(\cdot,\xi)\) es la respuesta del sistema a una
fuente puntual unitaria colocada en \(\xi\)}, con el mismo contorno que
el problema original.

\end{minipage}%
\end{tcolorbox}

\begin{tcolorbox}[enhanced jigsaw, arc=.35mm, bottomrule=.15mm, breakable, colback=white, colframe=quarto-callout-note-color-frame, left=2mm, leftrule=.75mm, opacityback=0, rightrule=.15mm, toprule=.15mm]
\begin{minipage}[t]{5.5mm}
\textcolor{quarto-callout-note-color}{\faInfo}
\end{minipage}%
\begin{minipage}[t]{\textwidth - 5.5mm}

\vspace{-3mm}\textbf{Teorema~\ref{thm-green-varias} --- La solución, por superposición}\vspace{3mm}

Si \(G\) satisface lo anterior, entonces

\[
u(\mathbf{x}) = \int_D G(\mathbf{x},\xi)\,f(\xi)\,d\xi
\]

resuelve el problema.

\end{minipage}%
\end{tcolorbox}

La verificación usa las propiedades de la delta:
\(L[u] = \int_D L[G]\,f\,d\xi = \int_D \delta(\mathbf{x}-\xi)f(\xi)\,d\xi = f(\mathbf{x})\),
y \(B[u] = \int_D B[G]\,f\,d\xi = 0\).

\subsubsection{Caso 2: ecuación homogénea, contorno
inhomogéneo}\label{caso-2-ecuaciuxf3n-homoguxe9nea-contorno-inhomoguxe9neo}

\[
L[u(\mathbf{x})] = 0 \ \text{ en } D, \qquad
B[u(\mathbf{x})] = h(\mathbf{x}) \ \text{ en } S .
\]

\begin{tcolorbox}[enhanced jigsaw, arc=.35mm, bottomrule=.15mm, breakable, colback=white, colframe=quarto-callout-note-color-frame, left=2mm, leftrule=.75mm, opacityback=0, rightrule=.15mm, toprule=.15mm]
\begin{minipage}[t]{5.5mm}
\textcolor{quarto-callout-note-color}{\faInfo}
\end{minipage}%
\begin{minipage}[t]{\textwidth - 5.5mm}

\vspace{-3mm}\textbf{Definición~\ref{def-green-modificada} --- Función de Green modificada}\vspace{3mm}

Se busca \(G(\mathbf{x},\xi)\) tal que

\[
L\big[G(\mathbf{x},\xi)\big] = 0 \ \text{ en } D, \qquad
B\big[G(\mathbf{x},\xi)\big] = \delta(\mathbf{x}-\xi) \ \text{ en } S ,
\]

con \(\xi \in S\). Entonces la solución es

\[
u(\mathbf{x}) = \int_S G(\mathbf{x},\xi)\,h(\xi)\,d\xi ,
\]

donde la integral se hace \textbf{sobre la superficie}, no sobre el
volumen.

\end{minipage}%
\end{tcolorbox}

\textbf{Lo que hay que retener:} el método es muy general y simple de
plantear; la dificultad está en hallar efectivamente \(G\), que depende
del operador \textbf{y} del dominio. Sin condiciones de contorno, la
ecuación de Helmholtz tiene la solución simple
\(\cos(\sqrt{\lambda}\,r)/(4\pi r)\), y hay que agregarle una función
con derivadas primera y segunda continuas para cumplir el contorno.
Cuando la ecuación de Helmholtz se reduce a la de Poisson, la integral
de Coulomb \textbf{es} la solución de Green.

\subsection{Separación de variables}\label{sec-resumen-9}

Es el método más usado y el más limitado: no se puede usar en general,
sino bajo condiciones de simetría en la ecuación y en las condiciones de
contorno e iniciales que permitan expresar el problema en un sistema de
coordenadas que admita la separación.

\subsubsection{El esquema, en seis
pasos}\label{el-esquema-en-seis-pasos}

\begin{enumerate}
\def\labelenumi{\arabic{enumi}.}
\tightlist
\item
  \textbf{Sustituir} \(u(x,y,z,\dots) = X(x)\,Y(y)\,Z(z)\cdots\) y ver
  si es posible construir con eso un conjunto de ecuaciones ordinarias
  separadas entre sí. La sustitución puede no ser un producto: en
  Hamilton--Jacobi conviene proponer una \textbf{suma}.
\item
  \textbf{Traducir las condiciones de contorno} de \(u\) en condiciones
  para \(X\), \(Y\), \(Z\). La continuidad y la acotación de \(u\)
  pueden agregar condiciones.
\item
  \textbf{Detectar para qué coordenadas los problemas resultantes son de
  Sturm--Liouville} y hallar sus autovalores y autofunciones.
\item
  \textbf{Resolver el resto de las ecuaciones} con condiciones de
  contorno homogéneas.
\item
  \textbf{Reconstruir} las autofunciones, \(\phi_n = X_nY_nZ_n\cdots\),
  y escribir la solución como combinación lineal,
  \(u = \sum_n a_n\phi_n\).
\item
  \textbf{Hallar los \(a_n\)} con las condiciones no homogéneas, usando
  la ortogonalidad y la completitud de las autofunciones del problema de
  Sturm--Liouville.
\end{enumerate}

El paso 3 es el que decide todo: sin la completitud de las
autofunciones, el paso 6 sería imposible.

\subsubsection{El caso testigo: la cuerda con extremos
fijos}\label{el-caso-testigo-la-cuerda-con-extremos-fijos}

\[
u_{tt} = a^2u_{xx}, \qquad
\begin{cases}
u(0,t) = 0, & u(x,0) = \varphi(x),\\
u(l,t) = 0, & u_t(x,0) = \psi(x).
\end{cases}
\]

\textbf{Paso 1.} Con \(u = X(x)T(t)\) queda
\(\dfrac{X''}{X} = \dfrac{T''}{a^2T} = \lambda \equiv -k^2\), donde
\(-k^2\) es la \textbf{constante de separación}, y la EDP se parte en
dos ecuaciones ordinarias:

\[
X'' + k^2X = 0, \qquad T'' + a^2k^2T = 0 .
\]

\textbf{Pasos 2 y 3.} De \(u(0,t)=0\) sale \(X(0)=0\) y de \(u(l,t)=0\)
sale \(X(l)=0\). Con \(X = C\cos kx + D\operatorname{sen} kx\), la
primera condición da \(C=0\) y la segunda pide
\(D\operatorname{sen}(kl)=0\), lo cual ocurre solo para

\[
k = \pm\frac{n\pi}{l} .
\]

\textbf{Paso 4.} \(T = A\cos akt + B\operatorname{sen} akt\).

\textbf{Paso 5.} Por linealidad, la solución es la suma de las
autofunciones:

\[
u(x,t) = \sum_{n=1}^{\infty}
\left[A_n\cos\frac{n\pi a\,t}{l} + B_n\operatorname{sen}\frac{n\pi a\,t}{l}\right]
\operatorname{sen}\frac{n\pi x}{l} .
\]

\textbf{Paso 6.} Evaluando en \(t=0\) la serie y su derivada temporal se
obtienen dos series de Fourier en senos, de donde

\[
A_n = \frac{2}{l}\int_0^l \varphi(z)\,\operatorname{sen}\frac{n\pi z}{l}\,dz ,
\qquad
B_n = \frac{2}{n\pi a}\int_0^l \psi(z)\,\operatorname{sen}\frac{n\pi z}{l}\,dz .
\]

\textbf{Interpretación física.} Cada término es una \textbf{onda
estacionaria}. Las frecuencias propias son

\[
\omega_n = \frac{n\pi a}{l} ,
\]

en los puntos \(x = 0,\; l/n,\; 2l/n,\;\dots,\;l\) la amplitud de la
\(n\)-ésima armónica se anula y están los \textbf{nodos}, y en
\(x = l/2n,\; 3l/2n,\;\dots,\;(2n-1)l/2n\) la amplitud es máxima y están
los \textbf{vientres} o antinodos.

\subsubsection{Un caso de cuarto orden: la viga
elástica}\label{un-caso-de-cuarto-orden-la-viga-eluxe1stica}

La deflexión \(w(x,t)\) de una viga fina uniforme de longitud \(L\)
satisface

\[
\rho\,\frac{\partial^2 w}{\partial t^2} + \gamma\,\frac{\partial w}{\partial t}
+ EI\,\frac{\partial^4 w}{\partial x^4} = f(x,t) ,
\]

con \(\rho\) la masa por unidad de longitud, \(\gamma\) el coeficiente
de fricción viscosa, \(EI\) el módulo de flexión y \(f\) la fuerza
externa. Sin fuerzas externas, con \(w = T(t)\phi(x)\):

\[
\frac{d^4\phi}{dx^4} - \omega^4\phi = 0, \qquad
\rho\,\frac{d^2T}{dt^2} + \gamma\,\frac{dT}{dt} + EI\,\omega^4\,T = 0 .
\]

Para una viga \textbf{empotrada en \(x=0\) y libre en \(x=L\)} las
condiciones de contorno son

\[
w(0,t) = \frac{\partial w}{\partial x}(0,t) = \frac{\partial^2 w}{\partial x^2}(L,t)
= \frac{\partial^3 w}{\partial x^3}(L,t) = 0 .
\]

El operador \(\mathcal{L} = d^4/dx^4\) con estas condiciones es
\textbf{autoadjunto}, así que sus autofunciones forman un conjunto
ortogonal completo de \(L^2[0,L]\). Partiendo de la solución general
\(\phi = c_1\cosh\omega x + c_2\operatorname{senh}\omega x + c_3\cos\omega x + c_4\operatorname{sen}\omega x\)
y aplicando las cuatro condiciones se obtienen las autofunciones

\[
\phi_k(x) = \frac{\cosh\omega_kx - \cos\omega_kx}{\cosh\omega_kL + \cos\omega_kL}
- \frac{\operatorname{senh}\omega_kx - \operatorname{sen}\omega_kx}{\operatorname{senh}\omega_kL + \operatorname{sen}\omega_kL} ,
\]

con los autovalores dados por las raíces de

\[
\cos\omega_kL = -\frac{1}{\cosh\omega_kL} .
\]

Las condiciones iniciales y la ecuación para \(T(t)\) cierran el
problema; su forma sugiere tratarla con transformada de Laplace.

\subsection{Formulario}\label{sec-resumen-10}

\textbf{Primer orden}

{\def\LTcaptype{none} % do not increment counter
\begin{longtable}[]{@{}
  >{\raggedright\arraybackslash}p{(\linewidth - 2\tabcolsep) * \real{0.5000}}
  >{\raggedright\arraybackslash}p{(\linewidth - 2\tabcolsep) * \real{0.5000}}@{}}
\toprule\noalign{}
\endhead
\bottomrule\noalign{}
\endlastfoot
Sistema característico &
\(\dfrac{dx}{a} = \dfrac{dy}{b} = \dfrac{du}{c}\) \\
El campo & \(\vec V = (a,b,c)\); su sombra en el piso,
\(\vec v = (a,b)\) \\
La ecuación, en una línea & \(\vec v\cdot\nabla u = c\) \\
Solución general de la homogénea & \(u = F(\phi)\), con \(\phi\) primera
integral \\
Onda simple & \(u_t + cu_x = 0 \Rightarrow u = f(x-ct)\) \\
Burgers &
\(u_t + uu_x = 0 \Rightarrow u\big(x_0+f(x_0)t,\,t\big) = f(x_0)\) \\
Normal a la superficie solución & \((u_x,\,u_y,\,-1)\) \\
Dirección tangente & \((a_1,\,a_2,\,a)\) \\
Transversalidad & \(a_1y_0' - a_2x_0' \neq 0\) \\
\end{longtable}
}

\textbf{Segundo orden, dos variables}

{\def\LTcaptype{none} % do not increment counter
\begin{longtable}[]{@{}
  >{\raggedright\arraybackslash}p{(\linewidth - 2\tabcolsep) * \real{0.5000}}
  >{\raggedright\arraybackslash}p{(\linewidth - 2\tabcolsep) * \real{0.5000}}@{}}
\toprule\noalign{}
\endhead
\bottomrule\noalign{}
\endlastfoot
Ecuación & \(A u_{xx} + 2B u_{xy} + C u_{yy} = F\) \\
Discriminante & \(\Delta = B^2 - AC\) \\
EDO de las características & \(A(y')^2 - 2B\,y' + C = 0\) \\
Coeficiente de \(u_{\xi\xi}\) &
\(\alpha = A\xi_x^2 + 2B\xi_x\xi_y + C\xi_y^2\) \\
Coeficiente de \(u_{\xi\psi}\) &
\(\beta = 2[A\xi_x\psi_x + B(\xi_x\psi_y+\xi_y\psi_x) + C\xi_y\psi_y]\) \\
Onda & \(x \pm ct = \text{const}\), \(\;u = f(x+ct)+g(x-ct)\) \\
\end{longtable}
}

\textbf{Muchas variables}

{\def\LTcaptype{none} % do not increment counter
\begin{longtable}[]{@{}
  >{\raggedright\arraybackslash}p{(\linewidth - 2\tabcolsep) * \real{0.5000}}
  >{\raggedright\arraybackslash}p{(\linewidth - 2\tabcolsep) * \real{0.5000}}@{}}
\toprule\noalign{}
\endhead
\bottomrule\noalign{}
\endlastfoot
Parte principal & \(L_P[u] = \sum_{i,j} a_{ij}u_{x_ix_j}\), con
\(a_{ij}=a_{ji}\) \\
Cambio lineal & \(A_{ik} = \sum_{j,l} c_{ij}c_{kl}a_{jl}\) \\
Clasificación & por la signatura \((n_+,n_-,n_0)\) \\
Obstrucción & \(n(n-1)/2\) condiciones contra \(n\) funciones \\
\end{longtable}
}

\textbf{Los tres métodos lineales}

{\def\LTcaptype{none} % do not increment counter
\begin{longtable}[]{@{}
  >{\raggedright\arraybackslash}p{(\linewidth - 2\tabcolsep) * \real{0.5000}}
  >{\raggedright\arraybackslash}p{(\linewidth - 2\tabcolsep) * \real{0.5000}}@{}}
\toprule\noalign{}
\endhead
\bottomrule\noalign{}
\endlastfoot
Transformada en el tiempo &
\(\Phi = \frac{1}{\sqrt{2\pi}}\int_0^\infty u\,e^{-ist}dt\) \\
Reduce a & \(\nabla^2\Phi + \lambda\Phi = g\),
\(\;\partial_n\Phi + \alpha\Phi = 0\) \\
Green, fuente adentro & \(L[G]=\delta\), \(B[G]=0\),
\(\;u = \int_D G f\,d\xi\) \\
Green, fuente en el borde & \(L[G]=0\), \(B[G]=\delta\),
\(\;u = \int_S G h\,d\xi\) \\
Separación, cuerda & \(\omega_n = n\pi a/l\),
\(\;X_n = \operatorname{sen}(n\pi x/l)\) \\
Coeficientes &
\(A_n = \frac{2}{l}\int_0^l\varphi\operatorname{sen}\frac{n\pi z}{l}dz\),
\(\;B_n = \frac{2}{n\pi a}\int_0^l\psi\operatorname{sen}\frac{n\pi z}{l}dz\) \\
\end{longtable}
}

\subsection{Me dan una ecuación: qué hago}\label{sec-resumen-11}

\begin{enumerate}
\def\labelenumi{\arabic{enumi}.}
\tightlist
\item
  \textbf{¿De qué orden es?} El mayor \(N\) que aparece efectivamente.
  Todo lo demás depende de esto.
\item
  \textbf{¿Cuál es la parte principal?} Solo los términos de orden
  exactamente \(N\). Es lo único que se mira para clasificar.
\item
  \textbf{¿Qué tan lineal es?} Mirar de qué dependen los coeficientes de
  la parte principal:

  \begin{itemize}
  \tightlist
  \item
    solo de \(x\), y todo lo demás también lineal → \textbf{lineal};
  \item
    solo de \(x\), con no linealidad en órdenes menores →
    \textbf{semilineal};
  \item
    de \(u\) o de sus derivadas de orden menor → \textbf{cuasi-lineal};
  \item
    las derivadas de orden \(N\) entran de forma no lineal →
    \textbf{completamente no lineal}.
  \end{itemize}
\item
  \textbf{¿De qué tipo es?}

  \begin{itemize}
  \tightlist
  \item
    \emph{Primer orden}: no hay tipos; hay características, y siempre
    las hay.
  \item
    \emph{Segundo orden, dos variables}: calcular \(\Delta = B^2-AC\)
    ---cuidado con el \(2B\)--- y ver el signo. Si \(A\), \(B\), \(C\)
    dependen del punto, el tipo puede cambiar de región en región.
  \item
    \emph{Segundo orden, más variables}: armar la matriz simétrica
    \(a_{ij}\), diagonalizar o completar cuadrados, contar
    \((n_+,n_-,n_0)\).
  \end{itemize}
\item
  \textbf{¿Qué datos me dan, y alcanzan?} Revisar existencia, unicidad y
  estabilidad. Si el dato es de Cauchy sobre una curva abierta y la
  ecuación es elíptica, sospechar del problema de Hadamard.
\item
  \textbf{¿Qué método corresponde?}
\end{enumerate}

{\def\LTcaptype{none} % do not increment counter
\begin{longtable}[]{@{}
  >{\raggedright\arraybackslash}p{(\linewidth - 2\tabcolsep) * \real{0.5000}}
  >{\raggedright\arraybackslash}p{(\linewidth - 2\tabcolsep) * \real{0.5000}}@{}}
\toprule\noalign{}
\begin{minipage}[b]{\linewidth}\raggedright
Si la ecuación\ldots{}
\end{minipage} & \begin{minipage}[b]{\linewidth}\raggedright
el método es\ldots{}
\end{minipage} \\
\midrule\noalign{}
\endhead
\bottomrule\noalign{}
\endlastfoot
es de primer orden, aunque sea casi--lineal &
\textbf{características} \\
es lineal y una variable se puede eliminar por transformada &
\textbf{transformadas integrales} (queda Helmholtz) \\
es lineal y hay que resolverla para muchas fuentes distintas &
\textbf{función de Green} \\
es lineal, el dominio tiene la simetría adecuada y los contornos son
homogéneos & \textbf{separación de variables} \\
\end{longtable}
}

En la práctica se combinan: lo más frecuente es transformar en el tiempo
para quedarse con un problema de Helmholtz, y resolver ése con Green o
con separación de variables.

\begin{enumerate}
\def\labelenumi{\arabic{enumi}.}
\setcounter{enumi}{6}
\tightlist
\item
  \textbf{Si es de primer orden y hay que resolverla:} escribir el
  sistema característico, parametrizar la curva inicial, verificar la
  transversalidad, integrar y despejar. Antes de dar la solución por
  buena, pasar la lista de las cinco fallas: ¿los coeficientes son
  \(C^1\)?, ¿el campo se anula en algún punto del dominio?, ¿el
  determinante de transversalidad se anula en algún \(\sigma\)?, ¿es
  cuasi-lineal y por lo tanto la respuesta vale solo hasta el quiebre?,
  ¿el dato alcanza para toda la región que piden?
\item
  \textbf{Si es de segundo orden en dos variables y hay que llevarla a
  forma canónica:} resolver \(A(y')^2 - 2By' + C = 0\), usar sus
  soluciones como coordenadas nuevas, y dividir.
\end{enumerate}

\section{Preguntas guía para el final}\label{sec-preguntas}

No son ejercicios de cálculo: son las preguntas que conviene poder
contestar hablando, sin papel. Cada una tiene desplegada la idea de lo
que se espera.

\begin{tcolorbox}[enhanced jigsaw, arc=.35mm, bottomrule=.15mm, bottomtitle=1mm, breakable, colback=white, colbacktitle=quarto-callout-tip-color!10!white, colframe=quarto-callout-tip-color-frame, coltitle=black, left=2mm, leftrule=.75mm, opacityback=0, opacitybacktitle=0.6, rightrule=.15mm, title=\textcolor{quarto-callout-tip-color}{\faLightbulb}\hspace{0.5em}{1. ¿Por qué una EDP necesita datos sobre curvas y no en puntos aislados?}, titlerule=0mm, toprule=.15mm, toptitle=1mm]

Porque su indeterminación es \textbf{funcional} y no de constantes. La
solución general de \(u_y = 0\) es \(u = f(x)\) con \(f\) arbitraria, y
una función arbitraria no se fija con dos números. La regla empírica es
que una EDP de orden \(n\) arrastra \(n\) funciones arbitrarias, así
como una EDO de orden \(n\) arrastra \(n\) constantes. Ver
sección~\ref{sec-porque}.

\end{tcolorbox}

\begin{tcolorbox}[enhanced jigsaw, arc=.35mm, bottomrule=.15mm, bottomtitle=1mm, breakable, colback=white, colbacktitle=quarto-callout-tip-color!10!white, colframe=quarto-callout-tip-color-frame, coltitle=black, left=2mm, leftrule=.75mm, opacityback=0, opacitybacktitle=0.6, rightrule=.15mm, title=\textcolor{quarto-callout-tip-color}{\faLightbulb}\hspace{0.5em}{2. ¿Qué es la parte principal y por qué se lleva toda la atención?}, titlerule=0mm, toprule=.15mm, toptitle=1mm]

Son los términos con derivadas de orden máximo. Deciden el tipo, las
características, la velocidad de propagación y si las discontinuidades
sobreviven. Los términos de orden menor cambian la solución pero no el
tipo: una onda con rozamiento sigue siendo hiperbólica. Ver
Definición~\ref{def-parte-principal}.

\end{tcolorbox}

\begin{tcolorbox}[enhanced jigsaw, arc=.35mm, bottomrule=.15mm, bottomtitle=1mm, breakable, colback=white, colbacktitle=quarto-callout-tip-color!10!white, colframe=quarto-callout-tip-color-frame, coltitle=black, left=2mm, leftrule=.75mm, opacityback=0, opacitybacktitle=0.6, rightrule=.15mm, title=\textcolor{quarto-callout-tip-color}{\faLightbulb}\hspace{0.5em}{3. Diferencia entre lineal, semilineal, cuasi-lineal y completamente no
lineal}, titlerule=0mm, toprule=.15mm, toptitle=1mm]

Es una escala de cuánto contamina la incógnita a su propia parte
principal. Lineal: coeficientes que dependen solo de \(x\), y vale la
superposición. Semilineal: parte principal lineal con coeficientes que
dependen solo de \(x\), no linealidad en órdenes menores, y el tipo
todavía se calcula sin resolver. Cuasi-lineal: los coeficientes de la
parte principal dependen de \(u\), y el tipo pasa a depender de la
solución. Completamente no lineal: las derivadas de orden máximo entran
de manera no lineal. Ver Figura~\ref{fig-linealidad} y
Tabla~\ref{tbl-linealidad}.

\end{tcolorbox}

\begin{tcolorbox}[enhanced jigsaw, arc=.35mm, bottomrule=.15mm, bottomtitle=1mm, breakable, colback=white, colbacktitle=quarto-callout-tip-color!10!white, colframe=quarto-callout-tip-color-frame, coltitle=black, left=2mm, leftrule=.75mm, opacityback=0, opacitybacktitle=0.6, rightrule=.15mm, title=\textcolor{quarto-callout-tip-color}{\faLightbulb}\hspace{0.5em}{4. ¿Qué significa que un problema esté bien planteado, y cuál de las
tres condiciones se olvida siempre?}, titlerule=0mm, toprule=.15mm, toptitle=1mm]

Existencia, unicidad y \textbf{estabilidad}. La que se olvida es la
tercera, y es la que hace que el problema sirva para algo: los datos
siempre tienen error, y si el problema los amplifica, tener solución
única no ayuda. Ver Definición~\ref{def-bien-planteado}.

\end{tcolorbox}

\begin{tcolorbox}[enhanced jigsaw, arc=.35mm, bottomrule=.15mm, bottomtitle=1mm, breakable, colback=white, colbacktitle=quarto-callout-tip-color!10!white, colframe=quarto-callout-tip-color-frame, coltitle=black, left=2mm, leftrule=.75mm, opacityback=0, opacitybacktitle=0.6, rightrule=.15mm, title=\textcolor{quarto-callout-tip-color}{\faLightbulb}\hspace{0.5em}{5. Contá el ejemplo de Hadamard y decí qué demuestra}, titlerule=0mm, toprule=.15mm, toptitle=1mm]

Un sistema de Cauchy--Riemann con dato \(e^{-\sqrt n}\cos(nx)\) sobre
\(y=0\): el dato y todas sus derivadas tienden a cero, y la solución, a
cualquier altura positiva, tiende a infinito. Como es la ecuación de
Laplace disfrazada, lo que demuestra no es que Laplace sea mala sino que
\textbf{darle datos de Cauchy sobre una curva abierta la arruina}: una
elíptica pide datos sobre una frontera cerrada. Ver
sección~\ref{sec-hadamard}.

\end{tcolorbox}

\begin{tcolorbox}[enhanced jigsaw, arc=.35mm, bottomrule=.15mm, bottomtitle=1mm, breakable, colback=white, colbacktitle=quarto-callout-tip-color!10!white, colframe=quarto-callout-tip-color-frame, coltitle=black, left=2mm, leftrule=.75mm, opacityback=0, opacitybacktitle=0.6, rightrule=.15mm, title=\textcolor{quarto-callout-tip-color}{\faLightbulb}\hspace{0.5em}{6. ¿Qué datos pide cada tipo?}, titlerule=0mm, toprule=.15mm, toptitle=1mm]

Hiperbólica: Cauchy en \(t=0\) ---posición y velocidad--- más contorno.
Parabólica: dato inicial más contorno, hacia adelante en el tiempo.
Elíptica: Dirichlet, Neumann o Robin sobre toda la frontera cerrada. Y
los tres errores clásicos son: Cauchy a una elíptica, ir hacia atrás con
una parabólica, y pedirle a una hiperbólica menos de lo que necesita.
Ver Tabla~\ref{tbl-datos}.

\end{tcolorbox}

\begin{tcolorbox}[enhanced jigsaw, arc=.35mm, bottomrule=.15mm, bottomtitle=1mm, breakable, colback=white, colbacktitle=quarto-callout-tip-color!10!white, colframe=quarto-callout-tip-color-frame, coltitle=black, left=2mm, leftrule=.75mm, opacityback=0, opacitybacktitle=0.6, rightrule=.15mm, title=\textcolor{quarto-callout-tip-color}{\faLightbulb}\hspace{0.5em}{7. ¿Qué es una solución débil y por qué hace falta?}, titlerule=0mm, toprule=.15mm, toptitle=1mm]

Una función que satisface la formulación \textbf{integral} de la ley de
conservación, sin necesidad de ser derivable ni continua. Hace falta
porque la ecuación de Burgers desarrolla discontinuidades sola a partir
de datos suaves, y esas discontinuidades son ondas de choque reales. La
definición se obtiene pasándole las derivadas a una función de prueba,
exactamente como en distribuciones. Ver
Definición~\ref{def-solucion-debil}.

\end{tcolorbox}

\begin{tcolorbox}[enhanced jigsaw, arc=.35mm, bottomrule=.15mm, bottomtitle=1mm, breakable, colback=white, colbacktitle=quarto-callout-tip-color!10!white, colframe=quarto-callout-tip-color-frame, coltitle=black, left=2mm, leftrule=.75mm, opacityback=0, opacitybacktitle=0.6, rightrule=.15mm, title=\textcolor{quarto-callout-tip-color}{\faLightbulb}\hspace{0.5em}{8. Explicá el método de las características con la imagen del observador}, titlerule=0mm, toprule=.15mm, toptitle=1mm]

Un observador camina por el plano y anota la altura de la superficie
solución sobre él. Si elige caminar en la dirección \((a,b)\) de los
coeficientes, la regla de la cadena convierte la EDP en tres EDOs. El
punto del techo que tiene encima dibuja un \textbf{hilo}, los hilos no
se cruzan, y barriendo la curva inicial \textbf{tejen la superficie
solución}. Ver sección~\ref{sec-observador} y Teorema~\ref{thm-tejido}.

\end{tcolorbox}

\begin{tcolorbox}[enhanced jigsaw, arc=.35mm, bottomrule=.15mm, bottomtitle=1mm, breakable, colback=white, colbacktitle=quarto-callout-tip-color!10!white, colframe=quarto-callout-tip-color-frame, coltitle=black, left=2mm, leftrule=.75mm, opacityback=0, opacitybacktitle=0.6, rightrule=.15mm, title=\textcolor{quarto-callout-tip-color}{\faLightbulb}\hspace{0.5em}{9. ¿Qué dice geométricamente la ecuación \(a\,u_x + b\,u_y = c\)?}, titlerule=0mm, toprule=.15mm, toptitle=1mm]

Que el campo vectorial \(\vec V = (a,b,c)\) es \textbf{tangente} a la
superficie solución en cada uno de sus puntos, porque
\(\vec V\cdot(u_x,u_y,-1) = a u_x + b u_y - c = 0\). Resolver la
ecuación es encontrar superficies tangentes al campo. Ver
\ref{eq-campo}.

\end{tcolorbox}

\begin{tcolorbox}[enhanced jigsaw, arc=.35mm, bottomrule=.15mm, bottomtitle=1mm, breakable, colback=white, colbacktitle=quarto-callout-tip-color!10!white, colframe=quarto-callout-tip-color-frame, coltitle=black, left=2mm, leftrule=.75mm, opacityback=0, opacitybacktitle=0.6, rightrule=.15mm, title=\textcolor{quarto-callout-tip-color}{\faLightbulb}\hspace{0.5em}{10. ¿Cuándo falla el método de las características y qué pasa entonces?}, titlerule=0mm, toprule=.15mm, toptitle=1mm]

Cuando la curva inicial es tangente a la dirección característica: el
hilo sale paralelo a ella en vez de despegarse, y no teje nada. Si el
dato no coincide con lo que la ecuación habría producido, \textbf{no hay
solución}; si coincide, hay \textbf{infinitas}. Es la condición de
transversalidad. Ver Teorema~\ref{thm-transversalidad} y
Teorema~\ref{thm-degenerado}.

\end{tcolorbox}

\begin{tcolorbox}[enhanced jigsaw, arc=.35mm, bottomrule=.15mm, bottomtitle=1mm, breakable, colback=white, colbacktitle=quarto-callout-tip-color!10!white, colframe=quarto-callout-tip-color-frame, coltitle=black, left=2mm, leftrule=.75mm, opacityback=0, opacitybacktitle=0.6, rightrule=.15mm, title=\textcolor{quarto-callout-tip-color}{\faLightbulb}\hspace{0.5em}{11. ¿Por qué las características de una cuasi-lineal se pueden cruzar, y
qué pasa ahí?}, titlerule=0mm, toprule=.15mm, toptitle=1mm]

Porque su dirección depende del valor de \(u\) que transportan, así que
dos características con valores distintos pueden converger. En el cruce
la solución tendría que valer dos cosas: la derivada explota y la onda
\textbf{rompe}, en el instante \(t^{*} = -1/\min f'\). Después de eso
hay que pasar a soluciones débiles. Ver Teorema~\ref{thm-quiebre}.

\end{tcolorbox}

\begin{tcolorbox}[enhanced jigsaw, arc=.35mm, bottomrule=.15mm, bottomtitle=1mm, breakable, colback=white, colbacktitle=quarto-callout-tip-color!10!white, colframe=quarto-callout-tip-color-frame, coltitle=black, left=2mm, leftrule=.75mm, opacityback=0, opacitybacktitle=0.6, rightrule=.15mm, title=\textcolor{quarto-callout-tip-color}{\faLightbulb}\hspace{0.5em}{12. ¿De dónde salen los nombres hiperbólica, parabólica y elíptica?}, titlerule=0mm, toprule=.15mm, toptitle=1mm]

De las cónicas. A la parte principal se le asocia la forma cuadrática
\(Q(\xi,\eta) = A\xi^2 + 2B\xi\eta + C\eta^2\), y el tipo de la curva
\(Q=1\) lo decide el mismo discriminante \(\Delta = B^2-AC\) que
clasifica la EDP. Las \textbf{asíntotas} de la cónica son las
\textbf{direcciones características}. No es una analogía: es la misma
clasificación. Ver sección~\ref{sec-conicas}.

\end{tcolorbox}

\begin{tcolorbox}[enhanced jigsaw, arc=.35mm, bottomrule=.15mm, bottomtitle=1mm, breakable, colback=white, colbacktitle=quarto-callout-tip-color!10!white, colframe=quarto-callout-tip-color-frame, coltitle=black, left=2mm, leftrule=.75mm, opacityback=0, opacitybacktitle=0.6, rightrule=.15mm, title=\textcolor{quarto-callout-tip-color}{\faLightbulb}\hspace{0.5em}{13. ¿Cuántas familias de características tiene cada tipo, y por qué?}, titlerule=0mm, toprule=.15mm, toptitle=1mm]

Las direcciones nulas son las raíces de \(Am^2+2Bm+C=0\), y el radicando
es \(\Delta\). Dos reales distintas si \(\Delta>0\) (hiperbólica), una
doble si \(\Delta=0\) (parabólica), ninguna real si \(\Delta<0\)
(elíptica). De ahí sale todo lo demás.

\end{tcolorbox}

\begin{tcolorbox}[enhanced jigsaw, arc=.35mm, bottomrule=.15mm, bottomtitle=1mm, breakable, colback=white, colbacktitle=quarto-callout-tip-color!10!white, colframe=quarto-callout-tip-color-frame, coltitle=black, left=2mm, leftrule=.75mm, opacityback=0, opacitybacktitle=0.6, rightrule=.15mm, title=\textcolor{quarto-callout-tip-color}{\faLightbulb}\hspace{0.5em}{14. ¿Qué es una forma canónica y cómo se la encuentra?}, titlerule=0mm, toprule=.15mm, toptitle=1mm]

El sistema de coordenadas en el que la parte principal queda lo más
simple posible. Se lo encuentra resolviendo la EDO de las
características \(A(y')^2 - 2By' + C = 0\) y usando sus soluciones como
coordenadas nuevas. Hiperbólica: \(u_{\xi\eta}=G\). Parabólica:
\(u_{\xi\xi}=G\). Elíptica: \(u_{\sigma\sigma}+u_{\tau\tau}=G\). Ver
Tabla~\ref{tbl-canonicas}.

\end{tcolorbox}

\begin{tcolorbox}[enhanced jigsaw, arc=.35mm, bottomrule=.15mm, bottomtitle=1mm, breakable, colback=white, colbacktitle=quarto-callout-tip-color!10!white, colframe=quarto-callout-tip-color-frame, coltitle=black, left=2mm, leftrule=.75mm, opacityback=0, opacitybacktitle=0.6, rightrule=.15mm, title=\textcolor{quarto-callout-tip-color}{\faLightbulb}\hspace{0.5em}{15. Deducí la solución de d'Alembert}, titlerule=0mm, toprule=.15mm, toptitle=1mm]

Con \(\xi = x+ct\), \(\eta = x-ct\), la ecuación de onda se vuelve
\(u_{\xi\eta}=0\), cuya solución general es \(u = f(\xi)+g(\eta)\), es
decir \(u = f(x+ct)+g(x-ct)\). Físicamente: cambiamos a las coordenadas
que viajan con la onda. Imponiendo las condiciones iniciales salen las
dos funciones y queda la fórmula con la semisuma de \(\varphi\) y la
integral de \(\psi\). Ver Ejercicio~\ref{exr-dalembert}.

\end{tcolorbox}

\begin{tcolorbox}[enhanced jigsaw, arc=.35mm, bottomrule=.15mm, bottomtitle=1mm, breakable, colback=white, colbacktitle=quarto-callout-tip-color!10!white, colframe=quarto-callout-tip-color-frame, coltitle=black, left=2mm, leftrule=.75mm, opacityback=0, opacitybacktitle=0.6, rightrule=.15mm, title=\textcolor{quarto-callout-tip-color}{\faLightbulb}\hspace{0.5em}{16. ¿Qué es el dominio de dependencia y por qué existe?}, titlerule=0mm, toprule=.15mm, toptitle=1mm]

El intervalo \([x_0-ct_0,\,x_0+ct_0]\) del que depende \(u(x_0,t_0)\), y
es la base del triángulo que forman las dos características que pasan
por el punto. Existe porque la ecuación es hiperbólica y la información
viaja a velocidad finita. La ecuación del calor no lo tiene: ahí el
dominio de dependencia es toda la recta. Ver
Ejercicio~\ref{exr-dominio-dependencia}.

\end{tcolorbox}

\begin{tcolorbox}[enhanced jigsaw, arc=.35mm, bottomrule=.15mm, bottomtitle=1mm, breakable, colback=white, colbacktitle=quarto-callout-tip-color!10!white, colframe=quarto-callout-tip-color-frame, coltitle=black, left=2mm, leftrule=.75mm, opacityback=0, opacitybacktitle=0.6, rightrule=.15mm, title=\textcolor{quarto-callout-tip-color}{\faLightbulb}\hspace{0.5em}{17. ¿Cómo se clasifica en muchas variables, y qué aparece de nuevo?}, titlerule=0mm, toprule=.15mm, toptitle=1mm]

Por la \textbf{signatura} \((n_+,n_-,n_0)\) de la matriz de
coeficientes, que por la ley de inercia de Sylvester no depende de cómo
se diagonalice. Aparece un tipo nuevo, el \textbf{ultrahiperbólico}, que
necesita al menos cuatro variables y prácticamente no ocurre en física
porque haría falta más de una dirección tipo tiempo. Ver
Definición~\ref{def-clasificacion-n}.

\end{tcolorbox}

\begin{tcolorbox}[enhanced jigsaw, arc=.35mm, bottomrule=.15mm, bottomtitle=1mm, breakable, colback=white, colbacktitle=quarto-callout-tip-color!10!white, colframe=quarto-callout-tip-color-frame, coltitle=black, left=2mm, leftrule=.75mm, opacityback=0, opacitybacktitle=0.6, rightrule=.15mm, title=\textcolor{quarto-callout-tip-color}{\faLightbulb}\hspace{0.5em}{18. ¿Qué relación hay entre la ecuación de onda y el cono de luz?}, titlerule=0mm, toprule=.15mm, toptitle=1mm]

La forma cuadrática de la ecuación de onda es \(\tau^2 - c^2|\zeta|^2\),
que es la métrica de Minkowski, y sus direcciones nulas forman el cono
de luz. Que la ecuación sea hiperbólica y que exista una velocidad
máxima de propagación son la misma afirmación. Ver
Figura~\ref{fig-cuadricas}.

\end{tcolorbox}

\begin{tcolorbox}[enhanced jigsaw, arc=.35mm, bottomrule=.15mm, bottomtitle=1mm, breakable, colback=white, colbacktitle=quarto-callout-tip-color!10!white, colframe=quarto-callout-tip-color-frame, coltitle=black, left=2mm, leftrule=.75mm, opacityback=0, opacitybacktitle=0.6, rightrule=.15mm, title=\textcolor{quarto-callout-tip-color}{\faLightbulb}\hspace{0.5em}{19. ¿Por qué a partir de tres variables no hay forma canónica en un
entorno?}, titlerule=0mm, toprule=.15mm, toptitle=1mm]

Porque diagonalizar en un entorno exige anular los \(n(n-1)/2\)
elementos de fuera de la diagonal y solo hay \(n\) funciones libres para
hacerlo. Para \(n\ge4\) el sistema está sobredeterminado. La excepción
es que los coeficientes sean constantes. Ver Tabla~\ref{tbl-conteo}.

\end{tcolorbox}

\begin{tcolorbox}[enhanced jigsaw, arc=.35mm, bottomrule=.15mm, bottomtitle=1mm, breakable, colback=white, colbacktitle=quarto-callout-tip-color!10!white, colframe=quarto-callout-tip-color-frame, coltitle=black, left=2mm, leftrule=.75mm, opacityback=0, opacitybacktitle=0.6, rightrule=.15mm, title=\textcolor{quarto-callout-tip-color}{\faLightbulb}\hspace{0.5em}{20. ¿Qué hace una transformada integral y qué requisitos tiene?}, titlerule=0mm, toprule=.15mm, toptitle=1mm]

Convierte derivadas en multiplicaciones, eliminando una variable: la EDP
baja una dimensión y queda un problema de Helmholtz. Requisitos:
ecuación lineal, coeficientes que no dependan de la variable
transformada, y frontera fija. Ver sección~\ref{sec-transformadas}.

\end{tcolorbox}

\begin{tcolorbox}[enhanced jigsaw, arc=.35mm, bottomrule=.15mm, bottomtitle=1mm, breakable, colback=white, colbacktitle=quarto-callout-tip-color!10!white, colframe=quarto-callout-tip-color-frame, coltitle=black, left=2mm, leftrule=.75mm, opacityback=0, opacitybacktitle=0.6, rightrule=.15mm, title=\textcolor{quarto-callout-tip-color}{\faLightbulb}\hspace{0.5em}{21. ¿Qué es la función de Green y en qué se apoya?}, titlerule=0mm, toprule=.15mm, toptitle=1mm]

La respuesta del sistema a una fuente puntual unitaria, con las mismas
condiciones de contorno. Se apoya en la superposición: cualquier fuente
es una suma de fuentes puntuales, así que la solución es la integral de
\(G\) contra la fuente. La integral de Coulomb es exactamente eso. Ver
Definición~\ref{def-green-varias}.

\end{tcolorbox}

\begin{tcolorbox}[enhanced jigsaw, arc=.35mm, bottomrule=.15mm, bottomtitle=1mm, breakable, colback=white, colbacktitle=quarto-callout-tip-color!10!white, colframe=quarto-callout-tip-color-frame, coltitle=black, left=2mm, leftrule=.75mm, opacityback=0, opacitybacktitle=0.6, rightrule=.15mm, title=\textcolor{quarto-callout-tip-color}{\faLightbulb}\hspace{0.5em}{22. ¿Por qué en separación de variables la constante de separación es
una constante?}, titlerule=0mm, toprule=.15mm, toptitle=1mm]

Porque un lado depende solo de una variable y el otro solo de las demás,
y son iguales para todos los valores. Fijando una y moviendo la otra,
ninguno de los dos puede cambiar. Ver
sección~\ref{sec-cuerda-separacion}.

\end{tcolorbox}

\begin{tcolorbox}[enhanced jigsaw, arc=.35mm, bottomrule=.15mm, bottomtitle=1mm, breakable, colback=white, colbacktitle=quarto-callout-tip-color!10!white, colframe=quarto-callout-tip-color-frame, coltitle=black, left=2mm, leftrule=.75mm, opacityback=0, opacitybacktitle=0.6, rightrule=.15mm, title=\textcolor{quarto-callout-tip-color}{\faLightbulb}\hspace{0.5em}{23. ¿Por qué las frecuencias son discretas?}, titlerule=0mm, toprule=.15mm, toptitle=1mm]

Porque pedirle a la parte espacial que se anule en los dos extremos deja
solo un conjunto numerable de constantes de separación posibles:
\(\operatorname{sen}(kl)=0\) obliga a \(k = n\pi/l\). La cuantización
sale de las condiciones de contorno, no de la ecuación, y es el mismo
mecanismo que discretiza los niveles de energía de un pozo cuántico.

\end{tcolorbox}

\begin{tcolorbox}[enhanced jigsaw, arc=.35mm, bottomrule=.15mm, bottomtitle=1mm, breakable, colback=white, colbacktitle=quarto-callout-tip-color!10!white, colframe=quarto-callout-tip-color-frame, coltitle=black, left=2mm, leftrule=.75mm, opacityback=0, opacitybacktitle=0.6, rightrule=.15mm, title=\textcolor{quarto-callout-tip-color}{\faLightbulb}\hspace{0.5em}{24. ¿Por qué una guitarra tiene una nota definida y un tambor no?}, titlerule=0mm, toprule=.15mm, toptitle=1mm]

Porque las frecuencias propias de la cuerda son \(\omega_n = n\pi a/l\),
múltiplos enteros de la fundamental, y el oído junta esa serie armónica
en una sola altura. Las de una membrana van como
\(\sqrt{m^2/a^2+n^2/b^2}\), que no son múltiplos de nada. Lo mismo pasa
con una viga, cuyas frecuencias están en la proporción
\(1:6{,}27:17{,}55\). La diferencia está en la lista de autovalores del
problema de Helmholtz del dominio. Ver
Ejercicio~\ref{exr-helmholtz-modos} y Ejemplo~\ref{exm-viga}.

\end{tcolorbox}

\begin{tcolorbox}[enhanced jigsaw, arc=.35mm, bottomrule=.15mm, bottomtitle=1mm, breakable, colback=white, colbacktitle=quarto-callout-tip-color!10!white, colframe=quarto-callout-tip-color-frame, coltitle=black, left=2mm, leftrule=.75mm, opacityback=0, opacitybacktitle=0.6, rightrule=.15mm, title=\textcolor{quarto-callout-tip-color}{\faLightbulb}\hspace{0.5em}{25. ¿Cuál de los cuatro métodos usarías, y por qué?}, titlerule=0mm, toprule=.15mm, toptitle=1mm]

Primero mirar la linealidad: si no es lineal y es de primer orden,
características, y si no es lineal ni de primer orden, no hay método
general. Si es lineal, mirar el dominio: una variable no acotada con
datos iniciales pide transformadas; muchas fuentes distintas sobre el
mismo dominio piden Green; un dominio simétrico con contornos homogéneos
pide separación. En la práctica se combinan: transformar en el tiempo y
después separar. Ver Tabla~\ref{tbl-elegir}.

\end{tcolorbox}

\chapter{Ecuaciones hiperbólicas}\label{sec-hiperbolicas}

\begin{tcolorbox}[enhanced jigsaw, arc=.35mm, bottomrule=.15mm, breakable, colback=white, colframe=quarto-callout-note-color-frame, left=2mm, leftrule=.75mm, opacityback=0, rightrule=.15mm, toprule=.15mm]
\begin{minipage}[t]{5.5mm}
\textcolor{quarto-callout-note-color}{\faInfo}
\end{minipage}%
\begin{minipage}[t]{\textwidth - 5.5mm}

\vspace{-3mm}\textbf{Qué hace falta para este capítulo}\vspace{3mm}

Este capítulo es la continuación directa del Capítulo~\ref{sec-edp}, y
lo da por leído: sobre todo la clasificación
(sección~\ref{sec-clasificacion}), las formas canónicas
(sección~\ref{sec-idea-canonica}) y la separación de variables en la
cuerda (sección~\ref{sec-cuerda-separacion}). Además se usan las series
de Fourier (Capítulo~\ref{sec-fourier}), la función de Green en una
variable (Capítulo~\ref{sec-green}), la delta de Dirac
(Capítulo~\ref{sec-distribuciones}) y la transformada de Fourier
(Capítulo~\ref{sec-tfourier}). Para la última sección hace falta cálculo
vectorial: divergencia, rotor y el teorema de la divergencia. Cada
sección avisa al empezar qué necesita.

\end{minipage}%
\end{tcolorbox}

\section{Todo lo que sabés del mundo te llegó por una
onda}\label{sec-hip-motivacion}

Contá los segundos entre el relámpago y el trueno, dividí por tres, y
tenés a cuántos kilómetros cayó el rayo. Es una regla que se aprende de
chico y que casi nadie se detiene a mirar. Pero para que funcione tienen
que pasar \textbf{tres cosas}, y las tres son el tema de este capítulo.

\textbf{Uno: el sonido tarda.} Si el aire respondiera de inmediato en
todos lados, el trueno llegaría junto con el relámpago y la regla no
serviría para nada. La perturbación viaja a una velocidad finita, unos
\(343\ \text{m/s}\), y por eso hay un \emph{antes} y un \emph{después}
de que llegue.

\textbf{Dos: el sonido llega entero.} Un trueno cercano es un estallido
seco, no un murmullo que crece despacio. La forma del pulso se conserva
durante el viaje: lo que sale del rayo es, con buena aproximación, lo
que llega al oído, solo que más débil. Si el aire se comportara como una
barra que conduce calor, el estallido se iría desparramando por el
camino hasta volverse un rumor sin bordes, y no habría ningún instante
que marcar.

\textbf{Tres: dos sonidos no se mezclan.} Durante la tormenta se sigue
oyendo la voz de quien está al lado, y se la oye igual que siempre. El
trueno y la voz se \emph{suman} en el aire y el oído los separa sin
esfuerzo. Eso es la superposición del sección~\ref{sec-lineales}: el
aire, a las amplitudes de la vida diaria, es un medio lineal.

Velocidad finita, forma que se conserva, superposición. Esas tres
propiedades juntas son la firma de la \textbf{ecuación de onda},

\begin{equation}\protect\phantomsection\label{eq-hip-onda-1d}{
u_{tt} = a^2\,u_{xx},
}\end{equation}

y de la familia a la que pertenece: las ecuaciones
\textbf{hiperbólicas}. En el sección~\ref{sec-tres-mundos} ya vimos, en
un solo gráfico, cómo se distingue de las otras dos: la de onda
transporta los quiebres intactos, la del calor los borra al instante, la
de Laplace ni siquiera tiene tiempo. Este capítulo se dedica entero a la
primera.

\subsection{Lo que viaja a velocidad finita}\label{sec-hip-catalogo}

La lista de cosas que obedecen \ref{eq-hip-onda-1d}, o su versión en más
dimensiones, es casi la lista de todo lo que sabemos del mundo sin
tocarlo.

\begin{itemize}
\tightlist
\item
  \textbf{La luz}, y con ella todo el espectro electromagnético. En 1865
  Maxwell escribió las ecuaciones del campo eléctrico y magnético, las
  combinó y le salió una ecuación de onda. La velocidad de esas ondas no
  la eligió él: salía de dos constantes medidas en el laboratorio con
  imanes y capacitores, y daba el valor medido para la velocidad de la
  luz. Así se descubrió qué es la luz. Lo rehacemos en el
  sección~\ref{sec-hip-maxwell}.
\item
  \textbf{El sonido}, en el aire, en el agua y en los sólidos.
\item
  \textbf{Las cuerdas y las membranas} de los instrumentos musicales.
\item
  \textbf{Los terremotos.} Un sismo emite dos ondas de volumen
  distintas: las P, de compresión, más rápidas, y las S, de corte, más
  lentas. Las S no atraviesan los líquidos, y la «zona de sombra» que
  eso deja del otro lado del planeta es la prueba de que el núcleo
  externo de la Tierra es líquido. Nadie bajó a mirarlo: se dedujo de
  dónde \emph{no} llegaban las ondas.
\item
  \textbf{La corriente en un cable} largo, que viaja como una onda de
  tensión (sección~\ref{sec-hip-telegrafo}).
\item
  \textbf{El espacio-tiempo.} Lejos de las fuentes, las ecuaciones de
  Einstein linealizadas dicen que las perturbaciones de la geometría se
  propagan como ondas a la velocidad de la luz. El 14 de septiembre de
  2015 los dos detectores de LIGO registraron una: la había emitido la
  fusión de dos agujeros negros a más de mil millones de años luz.
\end{itemize}

Video recomendado: LIGO Lab Caltech : MIT, \emph{The Sound of Two Black
Holes Colliding}, \url{https://www.youtube.com/watch?v=QyDcTbR-kEA}.

\begin{tcolorbox}[enhanced jigsaw, arc=.35mm, bottomrule=.15mm, breakable, colback=white, colframe=quarto-callout-warning-color-frame, left=2mm, leftrule=.75mm, opacityback=0, rightrule=.15mm, toprule=.15mm]
\begin{minipage}[t]{5.5mm}
\textcolor{quarto-callout-warning-color}{\faExclamationTriangle}
\end{minipage}%
\begin{minipage}[t]{\textwidth - 5.5mm}

\vspace{-3mm}\textbf{No todo lo que ondula es la ecuación de onda}\vspace{3mm}

Las olas del mar ondulan, pero no cumplen \ref{eq-hip-onda-1d}: las
largas viajan más rápido que las cortas, y por eso un paquete de olas se
desarma por el camino. A eso se lo llama \textbf{dispersión}, y lo
tratamos al final del capítulo (sección~\ref{sec-hip-telegrafo}). La
función de onda de la mecánica cuántica tampoco: la ecuación de
Schrödinger es de primer orden en el tiempo y tiene, a los ojos del
sección~\ref{sec-clasificacion}, la estructura de una parabólica.
«Hiperbólica» no es «que oscila»; es \textbf{que la información viaja a
velocidad finita}.

\end{minipage}%
\end{tcolorbox}

\subsection{El trueno y el sismógrafo}\label{sec-hip-trueno}

La regla del trueno tiene una versión profesional. Un sismógrafo no ve
el sismo: ve llegar primero las ondas P y, un rato después, las S. Como
las dos salieron juntas del foco, la demora entre ellas crece con la
distancia, exactamente como la demora entre el relámpago y el trueno.

\begin{figure}

\centering{

\pandocbounded{\includegraphics[keepaspectratio]{09-hiperbolicas_files/figure-pdf/fig-hip-trueno-output-1.pdf}}

}

\caption{\label{fig-hip-trueno}La misma idea en dos escalas. Cada recta
es el frente de una onda en el plano distancia--tiempo: la pendiente es
la inversa de la velocidad. \textbf{Izquierda:} el relámpago y el
trueno. La recta de la luz es prácticamente horizontal ---llega en
microsegundos---, así que la demora es, en la práctica, el tiempo que
tarda el sonido: unos tres segundos por kilómetro. \textbf{Derecha:} un
sismo registrado a 84 km. Las ondas P (\(6\) km/s) y las S (\(3{,}5\)
km/s) salieron juntas del foco, y la separación vertical entre las dos
rectas es la demora que ve el sismógrafo. Abajo, un sismograma sintético
en esa estación: primero llega la P, más débil, y diez segundos después
la S.}

\end{figure}%

Un sismógrafo registra la llegada de las ondas S \textbf{diez segundos}
después que la de las P. Tomando \(v_P = 6\ \text{km/s}\) y
\(v_S = 3{,}5\ \text{km/s}\), ¿a cuántos kilómetros está el foco?
Respondé en kilómetros, redondeado.

Si el foco está a distancia \(d\), las dos ondas llegan en \(d/v_P\) y
\(d/v_S\), así que \[
\Delta t = \frac{d}{v_S} - \frac{d}{v_P} = d\left(\frac{1}{v_S}-\frac{1}{v_P}\right)
\quad\Longrightarrow\quad
d = \frac{\Delta t}{1/v_S - 1/v_P} = \frac{10}{0{,}2857 - 0{,}1667} \approx 84\ \text{km}.
\] Una sola estación da la distancia pero no la dirección: el foco puede
estar en cualquier punto de una circunferencia de radio \(84\) km. Con
tres estaciones, las tres circunferencias se cortan en un punto. Es la
triangulación con la que se localiza un sismo, y es la regla del trueno
con dos ondas lentas en vez de una lenta y una instantánea.

Las tres propiedades del principio reaparecen en el cálculo: la
velocidad finita es lo que hace que haya un instante de llegada; que la
forma se conserve es lo que permite reconocer \emph{cuándo} llegó; y la
superposición es lo que deja ver las dos llegadas separadas en la misma
traza.

\subsection{La pregunta del capítulo}\label{sec-hip-pregunta}

Todo lo que sigue responde una sola pregunta:

\begin{quote}
\emph{Se sabe cómo empieza algo ---su forma y su velocidad en el
instante inicial---, cómo está sujeto en los bordes y qué fuerzas actúan
sobre él. ¿Cómo sigue?}
\end{quote}

Y para responderla hay tres herramientas, que ya conocemos del
Capítulo~\ref{sec-edp}:

\begin{enumerate}
\def\labelenumi{\arabic{enumi}.}
\tightlist
\item
  \textbf{Las características}, que siguen a la onda viajando con ella.
  Son el método natural cuando no hay bordes, o cuando los bordes están
  lejos. Dan la fórmula de d'Alembert.
\item
  \textbf{Los modos}, que descomponen el movimiento en vibraciones que
  no viajan sino que \emph{laten} en su lugar. Son el método natural
  cuando hay bordes. Dan la separación de variables.
\item
  \textbf{La función de Green}, que responde a un golpe puntual y
  después suma. Es el método natural cuando hay fuentes. Da los
  potenciales retardados del electromagnetismo.
\end{enumerate}

Las tres dan la misma respuesta cuando se pueden usar las tres, y buena
parte del capítulo es ver \textbf{por qué dan la misma respuesta}: en el
sección~\ref{sec-hip-reflexiones} vamos a ver que una cuerda con
extremos fijos se puede resolver con ondas que rebotan o con modos que
laten, y que las dos imágenes son la misma cuenta mirada de dos lados.

\section{Qué traemos del capítulo 8 y qué falta}\label{sec-hip-puente}

La ecuación de onda no es nueva. En el Capítulo~\ref{sec-edp} apareció
una y otra vez, siempre como \emph{ejemplo} de otra cosa: de una
hiperbólica, de una forma canónica, de un problema separable. Acá pasa a
ser el tema. Antes de arrancar conviene ordenar qué ya tenemos, qué le
falta a lo que tenemos, y qué vamos a agregar.

\subsection{Lo que ya vimos}\label{sec-hip-ya-vimos}

Se recuerda, no se repite. Si algo de esta lista no te suena, volvé a la
sección indicada antes de seguir.

\begin{longtable}[]{@{}
  >{\raggedright\arraybackslash}p{(\linewidth - 4\tabcolsep) * \real{0.3000}}
  >{\raggedright\arraybackslash}p{(\linewidth - 4\tabcolsep) * \real{0.2000}}
  >{\raggedright\arraybackslash}p{(\linewidth - 4\tabcolsep) * \real{0.5000}}@{}}
\caption{Lo que el Capítulo~\ref{sec-edp} ya dejó hecho sobre la
ecuación de onda.}\label{tbl-hip-ya-vimos}\tabularnewline
\toprule\noalign{}
\begin{minipage}[b]{\linewidth}\raggedright
Qué
\end{minipage} & \begin{minipage}[b]{\linewidth}\raggedright
Dónde
\end{minipage} & \begin{minipage}[b]{\linewidth}\raggedright
La idea en una línea
\end{minipage} \\
\midrule\noalign{}
\endfirsthead
\toprule\noalign{}
\begin{minipage}[b]{\linewidth}\raggedright
Qué
\end{minipage} & \begin{minipage}[b]{\linewidth}\raggedright
Dónde
\end{minipage} & \begin{minipage}[b]{\linewidth}\raggedright
La idea en una línea
\end{minipage} \\
\midrule\noalign{}
\endhead
\bottomrule\noalign{}
\endlastfoot
La ecuación de onda es hiperbólica & sección~\ref{sec-clasificacion} &
Su discriminante es positivo: tiene dos familias de características
reales. \\
Las características son \(x\pm at=\text{const}\) &
sección~\ref{sec-canonica-hiperbolica} & Son las rectas por las que
viaja la información, a velocidad \(a\) para cada lado. \\
La forma canónica es \(u_{\xi\eta}=0\) & sección~\ref{sec-idea-canonica}
& En las coordenadas que viajan con la onda, la ecuación no tiene nada
que decir. \\
Solución general \(u=f(x+at)+g(x-at)\) & sección~\ref{sec-idea-canonica}
& Un perfil que va para la izquierda más uno que va para la derecha. \\
La fórmula de d'Alembert & Ejercicio~\ref{exr-dalembert} & Resuelve el
problema en toda la recta a partir de \(\varphi\) y \(\psi\). \\
Dominio de dependencia y rango de influencia &
Ejercicio~\ref{exr-dominio-dependencia} & \(u(x_0,t_0)\) solo depende
del dato en \([x_0-at_0,\,x_0+at_0]\). \\
El cono de luz & sección~\ref{sec-cuadricas} & En más dimensiones, las
características forman un cono. \\
Qué datos pide una hiperbólica & sección~\ref{sec-condiciones} &
Posición y velocidad iniciales, y una condición en cada borde. \\
Separación de variables en la cuerda fija &
sección~\ref{sec-cuerda-separacion} & La solución es una suma de modos
\(\operatorname{sen}(n\pi x/l)\) que laten a frecuencias
\(n\pi a/l\). \\
La cuerda cargada de golpe & sección~\ref{sec-ejemplo-transformada} &
Con transformadas: una fuerza súbita produce un sobrepaso de factor
2. \\
Green del espacio libre, incluida Helmholtz &
sección~\ref{sec-green-caso-1} & La respuesta a una fuente puntual que
oscila es una onda esférica \(e^{ikr}/r\). \\
\end{longtable}

\subsection{Lo que falta, y por qué importa}\label{sec-hip-falta}

Mirada entera, la lista de arriba tiene tres agujeros. Todos tienen la
misma forma: \textbf{resolvimos el problema en un mundo más limpio que
el real}.

\textbf{No hay bordes.} d'Alembert resuelve la cuerda infinita. La
separación de variables resuelve la cuerda con extremos fijos. Pero no
vimos qué pasa \emph{en} un borde, qué le pasa a una onda cuando llega a
él, ni cómo se pasa de una imagen a la otra. Una cuerda de guitarra real
es las dos cosas a la vez: una onda que viaja y rebota, y un conjunto de
modos que suenan.

\textbf{No hay fuentes.} Salvo la cuerda cargada de golpe, todas las
ondas del Capítulo~\ref{sec-edp} se movían solas, a partir de lo que
tenían en el instante inicial. Pero las ondas interesantes las fabrica
algo: una antena, un parlante, una mano que sigue empujando. En el
Ejemplo~\ref{exm-onda-fuente} quedó prometida la fórmula de d'Alembert
\emph{con fuente}; hay que escribirla.

\textbf{No hay garantías.} En el sección~\ref{sec-bien-planteado}
dijimos que un problema tiene que tener una única solución que dependa
con continuidad de los datos. Para la ecuación de onda todavía no lo
probamos: la fórmula de d'Alembert da \emph{una} solución, pero ¿es la
única? ¿Y si la cuerda no es homogénea, y no hay fórmula?

\subsection{Lo que agregamos y lo que
expandimos}\label{sec-hip-agregamos}

\begin{longtable}[]{@{}
  >{\raggedright\arraybackslash}p{(\linewidth - 4\tabcolsep) * \real{0.1800}}
  >{\raggedright\arraybackslash}p{(\linewidth - 4\tabcolsep) * \real{0.4600}}
  >{\raggedright\arraybackslash}p{(\linewidth - 4\tabcolsep) * \real{0.3600}}@{}}
\caption{Mapa del capítulo.}\label{tbl-hip-mapa}\tabularnewline
\toprule\noalign{}
\begin{minipage}[b]{\linewidth}\raggedright
Sección
\end{minipage} & \begin{minipage}[b]{\linewidth}\raggedright
Qué hacemos
\end{minipage} & \begin{minipage}[b]{\linewidth}\raggedright
Por qué
\end{minipage} \\
\midrule\noalign{}
\endfirsthead
\toprule\noalign{}
\begin{minipage}[b]{\linewidth}\raggedright
Sección
\end{minipage} & \begin{minipage}[b]{\linewidth}\raggedright
Qué hacemos
\end{minipage} & \begin{minipage}[b]{\linewidth}\raggedright
Por qué
\end{minipage} \\
\midrule\noalign{}
\endhead
\bottomrule\noalign{}
\endlastfoot
sección~\ref{sec-hip-deduccion} & \textbf{Deducimos} la ecuación de la
cuerda con Newton, y vemos que la barra, el aire, el cable y la membrana
dan la misma. & Para saber qué es \(a\) y cuándo deja de valer el
modelo. \\
sección~\ref{sec-hip-energia} & \textbf{Demostramos} la unicidad con la
energía. & El libro de la cátedra enuncia el teorema y remite a otro
texto. La demostración es corta y es física pura. \\
sección~\ref{sec-hip-dalembert} & \textbf{Rehacemos} d'Alembert con
cuidado: la lectura geométrica y la estabilidad. & Era una promesa del
Capítulo~\ref{sec-edp}, y es la base de todo lo demás. \\
sección~\ref{sec-hip-reflexiones} & \textbf{Agregamos} los bordes:
reflexión en extremos fijos y libres, y la cuerda finita como rebotes. &
Es el puente entre ondas que viajan y modos que laten. \\
sección~\ref{sec-hip-fuentes} & \textbf{Agregamos} las fuentes:
d'Alembert con fuente, cuerda forzada, resonancia, bordes que se mueven.
& Es lo que convierte la teoría en ingeniería. \\
sección~\ref{sec-hip-membranas} & \textbf{Expandimos} a dos dimensiones:
membranas rectangulares y circulares. & Aparecen las líneas nodales, la
degeneración y un adelanto de las funciones de Bessel. \\
sección~\ref{sec-hip-espacio} & \textbf{Expandimos} a tres dimensiones:
ondas esféricas y el principio de Huygens. & En tres dimensiones un
chasquido suena seco; en dos, deja estela. \\
sección~\ref{sec-hip-maxwell} & \textbf{Expandimos} al
electromagnetismo: de Maxwell a los potenciales retardados y la
difracción. & Es la aplicación más importante de todo el capítulo. \\
sección~\ref{sec-hip-telegrafo} & \textbf{Agregamos} pérdidas y
dispersión. & Para saber qué se rompe cuando el medio no es ideal. \\
\end{longtable}

\begin{tcolorbox}[enhanced jigsaw, arc=.35mm, bottomrule=.15mm, breakable, colback=white, colframe=quarto-callout-tip-color-frame, left=2mm, leftrule=.75mm, opacityback=0, rightrule=.15mm, toprule=.15mm]
\begin{minipage}[t]{5.5mm}
\textcolor{quarto-callout-tip-color}{\faLightbulb}
\end{minipage}%
\begin{minipage}[t]{\textwidth - 5.5mm}

\vspace{-3mm}\textbf{Cómo se lee junto al libro de la cátedra}\vspace{3mm}

El orden de este capítulo es pedagógico, no el de González, pero
\textbf{todo lo del libro de la cátedra está acá}. La correspondencia es
esta:

{\def\LTcaptype{none} % do not increment counter
\begin{longtable}[]{@{}
  >{\raggedright\arraybackslash}p{(\linewidth - 2\tabcolsep) * \real{0.5000}}
  >{\raggedright\arraybackslash}p{(\linewidth - 2\tabcolsep) * \real{0.5000}}@{}}
\toprule\noalign{}
\begin{minipage}[b]{\linewidth}\raggedright
González
\end{minipage} & \begin{minipage}[b]{\linewidth}\raggedright
Acá
\end{minipage} \\
\midrule\noalign{}
\endhead
\bottomrule\noalign{}
\endlastfoot
9.1 Introducción, Teorema 9.1.1 & sección~\ref{sec-hip-deduccion} y
sección~\ref{sec-hip-energia} \\
9.2.1 Fórmula de d'Alembert & sección~\ref{sec-hip-dalembert} \\
9.2.2 Estabilidad de las soluciones &
sección~\ref{sec-hip-estabilidad} \\
9.2.3 Propagación libre de ondas en el espacio &
sección~\ref{sec-hip-espacio} \\
9.3.1 Membrana rectangular & sección~\ref{sec-hip-membranas} \\
9.3.2 Vibraciones forzadas & sección~\ref{sec-hip-fuentes} \\
9.4.1--9.4.2 Maxwell, potenciales retardados &
sección~\ref{sec-hip-maxwell} \\
9.4.3--9.4.4 Fórmulas de Green, Kirchhoff &
sección~\ref{sec-hip-kirchhoff} \\
9.6 Problemas 9.1 a 9.12 & sección~\ref{sec-hip-problemas}, todos
resueltos \\
\end{longtable}
}

\end{minipage}%
\end{tcolorbox}

\begin{tcolorbox}[enhanced jigsaw, arc=.35mm, bottomrule=.15mm, breakable, colback=white, colframe=quarto-callout-warning-color-frame, left=2mm, leftrule=.75mm, opacityback=0, rightrule=.15mm, toprule=.15mm]
\begin{minipage}[t]{5.5mm}
\textcolor{quarto-callout-warning-color}{\faExclamationTriangle}
\end{minipage}%
\begin{minipage}[t]{\textwidth - 5.5mm}

\vspace{-3mm}\textbf{Una convención: la velocidad se llama \(a\)}\vspace{3mm}

En el Capítulo~\ref{sec-edp} la velocidad de la onda aparece a veces
como \(c\) y a veces como \(a\). En este capítulo es siempre \(a\), como
en el libro de la cátedra, y \(c\) queda reservada para la velocidad de
la luz, que aparece en el sección~\ref{sec-hip-maxwell}.

\end{minipage}%
\end{tcolorbox}

\section{De dónde sale la ecuación}\label{sec-hip-deduccion}

En el 1 vimos nacer tres ecuaciones de primer orden en el tiempo a
partir de una ley de conservación. La de onda es de segundo orden en el
tiempo, y eso ya dice de dónde viene: no de un balance de
\emph{cantidad} sino de un balance de \emph{fuerzas}. Es la segunda ley
de Newton, escrita para un medio continuo.

\subsection{La cuerda}\label{sec-hip-cuerda}

Una cuerda tensa, sujeta en los extremos, en reposo sobre el eje \(x\).
Se la aparta un poco y se la suelta. Llamamos \(u(x,t)\) al
desplazamiento vertical del punto que en reposo estaba en \(x\).

Para escribir la ecuación hay que decir qué se supone. Son cuatro cosas,
y conviene tenerlas a la vista porque cada una, cuando falla, da una
ecuación distinta:

\begin{enumerate}
\def\labelenumi{\arabic{enumi}.}
\tightlist
\item
  \textbf{Los puntos se mueven solo en vertical.} No hay desplazamiento
  a lo largo de la cuerda.
\item
  \textbf{La cuerda es perfectamente flexible.} No resiste la flexión:
  la única fuerza interna es la tensión, y va tangente a la cuerda.
\item
  \textbf{Las pendientes son chicas}, \(|u_x| \ll 1\). Todo lo que sea
  de orden \(u_x^2\) se desprecia.
\item
  \textbf{La densidad lineal \(\rho(x)\) y la tensión \(T\)} son las de
  la cuerda en reposo.
\end{enumerate}

\begin{figure}

\centering{

\pandocbounded{\includegraphics[keepaspectratio]{09-hiperbolicas_files/figure-pdf/fig-hip-elemento-output-1.pdf}}

}

\caption{\label{fig-hip-elemento}Un pedacito de cuerda entre \(x\) y
\(x+\Delta x\), muy exagerado. Sobre cada extremo tira el resto de la
cuerda con una fuerza de módulo \(T\) tangente a la curva. Las
componentes horizontales, \(T\cos\theta\), son casi iguales y se
cancelan. Las verticales, \(T\operatorname{sen}\theta\), no: si la
cuerda está curvada, un extremo tira más hacia abajo que lo que el otro
tira hacia arriba, y la diferencia es lo que acelera el pedacito. Por
eso la fuerza neta es proporcional a la \textbf{curvatura}, y no a la
pendiente.}

\end{figure}%

Sobre el pedacito \([x, x+\Delta x]\) actúan las tensiones de los dos
lados, cada una tangente a la cuerda. Si \(\theta(x)\) es el ángulo de
la tangente, las componentes horizontales son \(T\cos\theta \approx T\)
en los dos extremos, con signos opuestos, y se cancelan: por eso el
pedacito no se mueve en horizontal, que es la hipótesis 1. Las
verticales son \(T\operatorname{sen}\theta\), y como la pendiente es
chica,

\[
\operatorname{sen}\theta \approx \tan\theta = u_x .
\]

La fuerza vertical neta es entonces
\(T\,u_x(x+\Delta x,t) - T\,u_x(x,t)\). Si además actúa una fuerza
externa de densidad \(F(x,t)\) ---fuerza por unidad de longitud---, el
pedacito de masa \(\rho\,\Delta x\) cumple

\[
\rho\,\Delta x\;u_{tt} \;=\; T\big[u_x(x+\Delta x,t) - u_x(x,t)\big] + F\,\Delta x .
\]

Dividiendo por \(\Delta x\) y haciendo \(\Delta x\to 0\),

\begin{equation}\protect\phantomsection\label{eq-hip-cuerda}{
\rho\,u_{tt} = T\,u_{xx} + F .
}\end{equation}

Con \(\rho\) constante y llamando

\begin{equation}\protect\phantomsection\label{eq-hip-velocidad}{
a = \sqrt{\frac{T}{\rho}}, \qquad f = \frac{F}{\rho},
}\end{equation}

queda \(u_{tt} = a^2u_{xx} + f\), que es exactamente la ecuación (9.1)
del libro de la cátedra.

\begin{tcolorbox}[enhanced jigsaw, arc=.35mm, bottomrule=.15mm, breakable, colback=white, colframe=quarto-callout-important-color-frame, left=2mm, leftrule=.75mm, opacityback=0, rightrule=.15mm, toprule=.15mm]
\begin{minipage}[t]{5.5mm}
\textcolor{quarto-callout-important-color}{\faExclamation}
\end{minipage}%
\begin{minipage}[t]{\textwidth - 5.5mm}

\vspace{-3mm}\textbf{Qué dice la ecuación, sin símbolos}\vspace{3mm}

\textbf{La aceleración de cada punto es proporcional a la curvatura de
la cuerda en ese punto.} Donde la cuerda es cóncava hacia abajo ---la
cresta de una joroba--- la tensión la empuja hacia abajo; donde es
cóncava hacia arriba, hacia arriba. Un tramo recto, por inclinado que
esté, no siente fuerza neta: tira lo mismo de los dos lados.

Y \(a\) tiene unidades de velocidad:
\(\sqrt{\text{N}/(\text{kg/m})} = \text{m/s}\). Más tensión, ondas más
rápidas; más masa, ondas más lentas.

\end{minipage}%
\end{tcolorbox}

\emph{(En la versión web hay acá un gráfico interactivo: una joroba
sobre la cuerda con flechas que muestran la fuerza neta \(T\,u_{xx}\) en
cada punto. Se puede cambiar el ancho de la joroba, la tensión, y
sumarle una rampa recta. Al achicar el ancho la fuerza crece como
\(1/\text{ancho}^2\); la rampa no cambia ninguna flecha.)}

La cuerda más aguda de una guitarra es de acero, con densidad lineal
\(\rho = 0{,}40\ \text{g/m}\), y está tensada con \(T = 70\ \text{N}\).
¿A qué velocidad viajan las ondas por ella? Respondé en m/s.

Pasando todo a unidades del SI,
\(\rho = 4{,}0\times 10^{-4}\ \text{kg/m}\), y
\[a = \sqrt{\frac{T}{\rho}} = \sqrt{\frac{70}{4\times10^{-4}}} = \sqrt{175\,000} \approx 418\ \text{m/s}.\]
Es más rápido que el sonido en el aire, y no hay contradicción: son dos
medios distintos. La cuerda empuja el aire, y lo que llega al oído es
una onda \textbf{en el aire}, que viaja a \(343\ \text{m/s}\) con la
frecuencia que le impuso la cuerda.

Si se duplica la tensión de una cuerda sin cambiar nada más, las ondas
viajan\ldots{}

\begin{itemize}
\tightlist
\item
  {el doble de rápido}
\item
  {\(\sqrt 2 \approx 1{,}41\) veces más rápido}
\item
  {a la misma velocidad, porque la masa no cambió}
\item
  {la mitad de rápido}
\end{itemize}

Por \ref{eq-hip-velocidad}, \(a\propto\sqrt T\). Y como en el
sección~\ref{sec-cuerda-separacion} vimos que las frecuencias de la
cuerda son \(\omega_n = n\pi a/l\), la nota sube en el mismo factor
\(\sqrt2\). En música eso es media octava, seis semitonos. Por eso se
afina una guitarra girando las clavijas: se cambia \(T\) para cambiar
\(a\).

\begin{tcolorbox}[enhanced jigsaw, arc=.35mm, bottomrule=.15mm, breakable, colback=white, colframe=quarto-callout-note-color-frame, left=2mm, leftrule=.75mm, opacityback=0, rightrule=.15mm, toprule=.15mm]
\begin{minipage}[t]{5.5mm}
\textcolor{quarto-callout-note-color}{\faInfo}
\end{minipage}%
\begin{minipage}[t]{\textwidth - 5.5mm}

\vspace{-3mm}\textbf{Cuándo deja de valer}\vspace{3mm}

Cada hipótesis tiene su precio. Si se abandona la \textbf{3}
---amplitudes grandes--- el \(\operatorname{sen}\theta\approx u_x\) deja
de valer y la ecuación se vuelve no lineal: la tensión misma cambia al
estirarse la cuerda, y los modos empiezan a intercambiar energía. Si se
abandona la \textbf{2} ---la cuerda tiene rigidez--- aparece un término
de cuarto orden, \(-EI\,u_{xxxx}\), y la cuerda se va pareciendo a la
viga del sección~\ref{sec-viga}: los armónicos altos quedan un poco más
agudos de lo que deberían. Los afinadores de piano lo conocen bien,
porque las cuerdas graves del piano son gruesas y rígidas.

\end{minipage}%
\end{tcolorbox}

\subsection{Si la cuerda no es homogénea}\label{sec-hip-inhomogenea}

Si la densidad cambia de un punto a otro ---una cuerda que se va
afinando--- la cuenta es la misma, con \(\rho(x)\) en lugar de \(\rho\).
Y si además la cuerda es un medio elástico cuya rigidez cambia, como una
barra hecha de materiales distintos, la tensión se reemplaza por una
rigidez \(k(x)\) que queda \textbf{adentro} de la derivada, porque lo
que se equilibra son los flujos de fuerza a los dos lados del pedacito:

\begin{equation}\protect\phantomsection\label{eq-hip-inhomogenea}{
\rho(x)\,u_{tt} = \frac{\partial}{\partial x}\Big(k(x)\,\frac{\partial u}{\partial x}\Big) + F(x,t),
\qquad \rho(x)>0,\quad k(x)>0 .
}\end{equation}

Esta es la ecuación del Teorema 9.1.1 de González, y es la que vamos a
usar para demostrar la unicidad en el sección~\ref{sec-hip-energia}.
Fijate que el operador de la derecha es el de Sturm--Liouville del
Capítulo~\ref{sec-sturm}: no es casualidad, y es la razón por la que la
separación de variables funciona también en medios inhomogéneos.

\begin{definition}[Ecuación de
onda]\protect\hypertarget{def-hip-onda}{}\label{def-hip-onda}

Se llama \textbf{ecuación de onda} en \(n\) dimensiones espaciales a

\begin{equation}\protect\phantomsection\label{eq-hip-onda-n}{
u_{tt} = a^2\,\nabla^2 u + f(\mathbf{x},t), \qquad \mathbf{x}\in\mathbb{R}^n,\ t>0,
}\end{equation}

donde \(a>0\) es la \textbf{velocidad de propagación} y \(f\) la
\textbf{fuente} o \textbf{forzante}. Es \textbf{homogénea} si
\(f\equiv0\) y \textbf{no homogénea} en caso contrario. Con \(n=1\) es
la ecuación de la cuerda; con \(n=2\), la de la membrana; con \(n=3\),
la del sonido y la de la luz.

La versión con coeficientes variables, \ref{eq-hip-inhomogenea}, se
llama ecuación de onda \textbf{en un medio inhomogéneo}.

\end{definition}

\subsection{La misma ecuación en otros medios}\label{sec-hip-analogias}

Que la cuerda dé \ref{eq-hip-onda-1d} no sorprende demasiado. Lo que
sorprende es que \textbf{medios que no se parecen en nada} den
exactamente la misma ecuación. En todos los casos la estructura es la
misma: hay una magnitud que acumula inercia y otra que acumula una
fuerza restitutiva, y la velocidad es la raíz del cociente entre la
rigidez y la inercia.

\begin{longtable}[]{@{}
  >{\raggedright\arraybackslash}p{(\linewidth - 10\tabcolsep) * \real{0.1600}}
  >{\raggedright\arraybackslash}p{(\linewidth - 10\tabcolsep) * \real{0.2000}}
  >{\raggedright\arraybackslash}p{(\linewidth - 10\tabcolsep) * \real{0.1400}}
  >{\raggedright\arraybackslash}p{(\linewidth - 10\tabcolsep) * \real{0.1700}}
  >{\raggedright\arraybackslash}p{(\linewidth - 10\tabcolsep) * \real{0.1300}}
  >{\raggedright\arraybackslash}p{(\linewidth - 10\tabcolsep) * \real{0.2000}}@{}}
\caption{Las analogías físicas del temario de la cátedra. En cada fila,
\(a = \sqrt{\text{rigidez}/\text{inercia}}\). Los valores son órdenes de
magnitud.}\label{tbl-hip-analogias}\tabularnewline
\toprule\noalign{}
\begin{minipage}[b]{\linewidth}\raggedright
Medio
\end{minipage} & \begin{minipage}[b]{\linewidth}\raggedright
Incógnita \(u\)
\end{minipage} & \begin{minipage}[b]{\linewidth}\raggedright
Rigidez
\end{minipage} & \begin{minipage}[b]{\linewidth}\raggedright
Inercia
\end{minipage} & \begin{minipage}[b]{\linewidth}\raggedright
\(a\)
\end{minipage} & \begin{minipage}[b]{\linewidth}\raggedright
Valor típico
\end{minipage} \\
\midrule\noalign{}
\endfirsthead
\toprule\noalign{}
\begin{minipage}[b]{\linewidth}\raggedright
Medio
\end{minipage} & \begin{minipage}[b]{\linewidth}\raggedright
Incógnita \(u\)
\end{minipage} & \begin{minipage}[b]{\linewidth}\raggedright
Rigidez
\end{minipage} & \begin{minipage}[b]{\linewidth}\raggedright
Inercia
\end{minipage} & \begin{minipage}[b]{\linewidth}\raggedright
\(a\)
\end{minipage} & \begin{minipage}[b]{\linewidth}\raggedright
Valor típico
\end{minipage} \\
\midrule\noalign{}
\endhead
\bottomrule\noalign{}
\endlastfoot
Cuerda & desplazamiento transversal & tensión \(T\) & densidad lineal
\(\rho\) & \(\sqrt{T/\rho}\) & \(100\)--\(500\) m/s \\
Barra, vibración longitudinal & desplazamiento a lo largo de la barra &
módulo de Young \(E\) & densidad \(\rho\) & \(\sqrt{E/\rho}\) & acero:
\(5000\) m/s \\
Gas (sonido) & presión o densidad, respecto del reposo & \(\gamma p_0\)
& densidad \(\rho_0\) & \(\sqrt{\gamma p_0/\rho_0}\) & aire: \(343\)
m/s \\
Línea de transmisión & tensión eléctrica \(V\) & \(1/C\) (capacidad por
metro) & \(L\) (inductancia por metro) & \(1/\sqrt{LC}\) & coaxial:
\(2\times10^8\) m/s \\
Membrana & desplazamiento transversal & tensión por unidad de longitud
\(\tau\) & densidad superficial \(\sigma\) & \(\sqrt{\tau/\sigma}\) &
parche de tambor: \(100\) m/s \\
Vacío & cada componente de \(\mathbf E\) y \(\mathbf B\) &
\(1/\varepsilon_0\) & \(\mu_0\) & \(1/\sqrt{\varepsilon_0\mu_0}\) &
\(3\times10^8\) m/s \\
\end{longtable}

Una cuenta de control con la tercera fila: el aire tiene
\(\gamma = 1{,}4\), \(p_0 = 101\,325\ \text{Pa}\) y
\(\rho_0 = 1{,}2\ \text{kg/m}^3\), así que
\(a = \sqrt{1{,}4\cdot 101\,325/1{,}2} \approx 344\ \text{m/s}\). Es el
número del trueno. El \(\gamma\) está porque las compresiones del sonido
son tan rápidas que el aire no llega a intercambiar calor, y se comprime
adiabáticamente. Newton, en 1687, puso \(p_0\) sin el \(\gamma\)
---compresión a temperatura constante--- y le dio un 15 \% menos. El
error lo corrigió Laplace más de un siglo después.

\begin{exercise}[]\protect\hypertarget{exr-hip-barra}{}\label{exr-hip-barra}

Deducir la ecuación de las \textbf{oscilaciones longitudinales} de una
barra elástica de sección \(S\), densidad \(\rho\) y módulo de Young
\(E\), siendo \(u(x,t)\) el desplazamiento, a lo largo de la barra, del
punto que en reposo estaba en \(x\). Usar la ley de Hooke: la tensión
normal ---fuerza por unidad de área--- es \(E\) por el estiramiento
relativo.

Después, decir qué condición de contorno corresponde a un
\textbf{extremo libre} y cuál a un \textbf{extremo empotrado}.
\emph{(Esto hace falta para los problemas 9.7 y 9.8 de la cátedra.)}

\end{exercise}

\begin{tcolorbox}[enhanced jigsaw, arc=.35mm, bottomrule=.15mm, bottomtitle=1mm, breakable, colback=white, colbacktitle=quarto-callout-tip-color!10!white, colframe=quarto-callout-tip-color-frame, coltitle=black, left=2mm, leftrule=.75mm, opacityback=0, opacitybacktitle=0.6, rightrule=.15mm, title=\textcolor{quarto-callout-tip-color}{\faLightbulb}\hspace{0.5em}{Solución}, titlerule=0mm, toprule=.15mm, toptitle=1mm]

\textbf{El estiramiento.} El pedacito que en reposo ocupaba
\([x, x+\Delta x]\) ahora ocupa
\([x+u(x,t),\ x+\Delta x + u(x+\Delta x,t)]\). Su largo pasó de
\(\Delta x\) a \(\Delta x + u(x+\Delta x,t) - u(x,t)\), así que el
estiramiento relativo es
\[\frac{u(x+\Delta x,t)-u(x,t)}{\Delta x} \;\longrightarrow\; u_x .\]

\textbf{La fuerza.} Por Hooke, la sección en \(x\) está sometida a una
tensión \(E\,u_x\) y la fuerza que el material de la derecha hace sobre
el de la izquierda es \(N(x,t) = E\,S\,u_x(x,t)\).

\textbf{Newton.} El pedacito tiene masa \(\rho S\Delta x\), y sobre él
actúan \(N(x+\Delta x)\) hacia la derecha y \(N(x)\) hacia la izquierda:
\[\rho S\,\Delta x\;u_{tt} = E S\big[u_x(x+\Delta x,t) - u_x(x,t)\big]
\quad\Longrightarrow\quad u_{tt} = \frac{E}{\rho}\,u_{xx}, \qquad a = \sqrt{E/\rho}.\]
Es la misma ecuación que la cuerda, aunque el movimiento es a lo largo
de la barra y no transversal. Lo que cumple el papel de la tensión es la
rigidez del material.

\textbf{Los bordes.} Un extremo \textbf{empotrado} no se mueve:
\(u = 0\), condición de Dirichlet. En un extremo \textbf{libre} no hay
nada que haga fuerza sobre la sección final, así que la fuerza interna
tiene que anularse ahí: \(E S\,u_x = 0\), o sea \(u_x = 0\), condición
de Neumann. Físicamente: el extremo libre se mueve todo lo que quiere,
pero \textbf{no está estirado ni comprimido}.

\textbf{Los números.} Para el acero, \(E \approx 200\ \text{GPa}\) y
\(\rho\approx 7850\ \text{kg/m}^3\), así que
\(a\approx 5000\ \text{m/s}\): unas quince veces la velocidad del sonido
en el aire. Por eso se oye el tren por el riel mucho antes de oírlo por
el aire.

\end{tcolorbox}

\subsection{El problema: qué datos hacen falta}\label{sec-hip-planteo}

La ecuación sola no determina el movimiento: una misma cuerda puede
hacer infinitas cosas distintas. Lo que falta es lo que el
sección~\ref{sec-condiciones} ya anticipó para una hiperbólica:
\textbf{dos condiciones iniciales y una condición en cada borde}.

\begin{definition}[Problema mixto de la
cuerda]\protect\hypertarget{def-hip-problema-mixto}{}\label{def-hip-problema-mixto}

Hallar \(u(x,t)\), definida y de clase \(C^2\) en \(0\le x\le l\),
\(t\ge 0\), tal que

\begin{equation}\protect\phantomsection\label{eq-hip-mixto}{
\begin{aligned}
&u_{tt} = a^2u_{xx} + f(x,t), && 0<x<l,\ t>0, \\
&u(0,t) = \mu_1(t),\quad u(l,t) = \mu_2(t), && t\ge 0 \quad\text{(condiciones de contorno)},\\
&u(x,0) = \varphi(x),\quad u_t(x,0) = \psi(x), && 0\le x\le l \quad\text{(condiciones iniciales)}.
\end{aligned}
}\end{equation}

Se llama \textbf{mixto} porque combina condiciones iniciales y de
contorno.

\end{definition}

Hacen falta \textbf{dos} condiciones iniciales porque la ecuación es de
segundo orden en el tiempo: es lo mismo que en la mecánica de una
partícula, donde para predecir la trayectoria hay que dar la posición y
la velocidad. Acá hay que darlas \emph{en cada punto} de la cuerda: la
forma inicial \(\varphi\) y la distribución de velocidades \(\psi\).

En los bordes, en cambio, alcanza con \textbf{una} condición en cada
extremo, y la física decide cuál:

\begin{longtable}[]{@{}
  >{\raggedright\arraybackslash}p{(\linewidth - 6\tabcolsep) * \real{0.3000}}
  >{\raggedright\arraybackslash}p{(\linewidth - 6\tabcolsep) * \real{0.1700}}
  >{\raggedright\arraybackslash}p{(\linewidth - 6\tabcolsep) * \real{0.2000}}
  >{\raggedright\arraybackslash}p{(\linewidth - 6\tabcolsep) * \real{0.3300}}@{}}
\caption{Condiciones de contorno de la cuerda, leídas físicamente. Los
signos corresponden al extremo derecho, \(x=l\); en el izquierdo
cambian, porque la tensión que actúa sobre el extremo tira hacia el otro
lado.}\label{tbl-hip-bordes}\tabularnewline
\toprule\noalign{}
\begin{minipage}[b]{\linewidth}\raggedright
El extremo\ldots{}
\end{minipage} & \begin{minipage}[b]{\linewidth}\raggedright
Condición
\end{minipage} & \begin{minipage}[b]{\linewidth}\raggedright
Nombre
\end{minipage} & \begin{minipage}[b]{\linewidth}\raggedright
Por qué
\end{minipage} \\
\midrule\noalign{}
\endfirsthead
\toprule\noalign{}
\begin{minipage}[b]{\linewidth}\raggedright
El extremo\ldots{}
\end{minipage} & \begin{minipage}[b]{\linewidth}\raggedright
Condición
\end{minipage} & \begin{minipage}[b]{\linewidth}\raggedright
Nombre
\end{minipage} & \begin{minipage}[b]{\linewidth}\raggedright
Por qué
\end{minipage} \\
\midrule\noalign{}
\endhead
\bottomrule\noalign{}
\endlastfoot
está atado a un punto fijo & \(u = 0\) & Dirichlet & no se mueve \\
se mueve de una forma dada & \(u = \mu(t)\) & Dirichlet no homogénea &
alguien lo mueve a mano \\
es un anillo sin masa que desliza por una varilla lisa & \(u_x = 0\) &
Neumann & la varilla no hace fuerza vertical, así que la tensión tiene
que llegar horizontal \\
se tira con una fuerza vertical dada \(\nu(t)\) & \(T u_x = \nu(t)\) &
Neumann no homogénea & la componente vertical de la tensión equilibra la
fuerza \\
está atado a un resorte de constante \(\kappa\) & \(T u_x = -\kappa u\)
& Robin & la tensión equilibra la fuerza del resorte \\
\end{longtable}

La tercera fila es la del problema 9.6 de la cátedra, y vale la pena
entenderla: si el anillo no tiene masa, cualquier fuerza vertical neta
sobre él le daría aceleración infinita. Así que la fuerza vertical neta
es cero, y como la varilla lisa solo empuja en horizontal, la componente
vertical de la tensión, \(T u_x\), tiene que ser cero. \textbf{Un
extremo libre llega siempre horizontal.}

\subsection{Tres casos límite}\label{sec-hip-casos-limite}

El problema mixto \ref{eq-hip-mixto} es el caso general. Pero hay
situaciones en que parte de los datos todavía no importan, y el problema
se simplifica. La escala que decide es el \textbf{tiempo que tarda la
onda en recorrer la cuerda}, \(l/a\).

\begin{enumerate}
\def\labelenumi{\arabic{enumi}.}
\tightlist
\item
  \textbf{Tiempos cortos, lejos de los bordes: el problema de Cauchy.}
  Si \(t \ll l/a\), a un punto del interior todavía no le llegó ninguna
  noticia de los extremos, y la cuerda se comporta como si fuera
  infinita. Solo quedan las condiciones iniciales:
  \[u_{tt} = a^2u_{xx} + f,\quad -\infty<x<\infty;\qquad u(x,0)=\varphi(x),\quad u_t(x,0)=\psi(x).\]
  Lo resuelve la fórmula de d'Alembert: sección~\ref{sec-hip-dalembert}.
\item
  \textbf{Cerca de un borde, lejos del otro: la cuerda semi-infinita.}
  Se estudia lo que pasa cerca de \(x=0\) mientras la influencia de
  \(x=l\) no llega:
  \[u(0,t) = \mu(t),\quad t\ge0;\qquad u(x,0)=\varphi(x),\quad u_t(x,0)=\psi(x),\quad 0\le x<\infty .\]
  Lo resolvemos con reflexiones en el sección~\ref{sec-hip-reflexiones}.
\item
  \textbf{Tiempos largos: el régimen permanente.} Si \(t\gg l/a\), se
  supone que el sistema \emph{perdió memoria} de las condiciones
  iniciales y solo importan los bordes y las fuentes. Es la respuesta
  estacionaria a un forzante periódico, y la tratamos en el
  sección~\ref{sec-hip-fuentes}.
\end{enumerate}

\begin{tcolorbox}[enhanced jigsaw, arc=.35mm, bottomrule=.15mm, breakable, colback=white, colframe=quarto-callout-warning-color-frame, left=2mm, leftrule=.75mm, opacityback=0, rightrule=.15mm, toprule=.15mm]
\begin{minipage}[t]{5.5mm}
\textcolor{quarto-callout-warning-color}{\faExclamationTriangle}
\end{minipage}%
\begin{minipage}[t]{\textwidth - 5.5mm}

\vspace{-3mm}\textbf{Ojo con el caso 3}\vspace{3mm}

Una cuerda \textbf{ideal} no pierde nunca la memoria: en
\ref{eq-hip-cuerda} no hay nada que disipe energía, y en el
sección~\ref{sec-hip-energia} vamos a demostrar que la energía se
conserva exactamente. Si se la pulsa y se la deja, suena para siempre.
El caso 3 tiene sentido solo cuando hay \textbf{algo de rozamiento},
aunque sea muy poco: alcanza para que, después de mucho tiempo, las
oscilaciones libres se hayan apagado y quede solo la respuesta a lo que
la sigue forzando. Lo precisamos en el sección~\ref{sec-hip-telegrafo}.

\end{minipage}%
\end{tcolorbox}

Una cuerda de \(l = 1\ \text{m}\) con \(a = 400\ \text{m/s}\) se golpea
en el medio. ¿Cuánto tarda el punto medio en «enterarse» de que la
cuerda tiene extremos?

\begin{itemize}
\tightlist
\item
  {\(2{,}5\) ms: lo que tarda la perturbación en ir hasta un extremo y
  volver}
\item
  {\(1{,}25\) ms: lo que tarda la perturbación en llegar a un extremo}
\item
  {Nunca: los extremos están fijos y no hacen nada}
\item
  {Nada: desde el primer instante siente los extremos}
\end{itemize}

El punto medio se entera de que hay bordes cuando le llega una onda
\textbf{reflejada} en uno de ellos. La perturbación tiene que ir del
centro al extremo (\(0{,}5\) m) y volver (\(0{,}5\) m): un metro a
\(400\) m/s, o sea \(2{,}5\) ms. En \(1{,}25\) ms la onda llega al
extremo, pero el rebote todavía tiene que volver. Y la última opción es
la de la ecuación del calor, que no tiene velocidad finita.

La moraleja es la escala: \textbf{milisegundos}. En una cuerda de
guitarra el caso 1 se termina casi de inmediato, y el sonido que se oye
es el de los modos.

\section{Energía y unicidad}\label{sec-hip-energia}

Ya tenemos el problema \ref{eq-hip-mixto}. Antes de resolverlo conviene
saber si \textbf{tiene sentido resolverlo}: si los datos determinan una
sola solución, o si podría haber dos movimientos distintos compatibles
con la misma forma inicial, la misma velocidad inicial y los mismos
bordes. Es el pedido de unicidad de Hadamard
(Definición~\ref{def-bien-planteado}).

González enuncia el resultado como Teorema 9.1.1 y remite la
demostración a otro libro. Acá la hacemos, porque es corta y porque la
herramienta que usa, la \textbf{energía}, es de las que más rinden en
todo el curso: sirve cuando no hay ninguna fórmula para la solución.

\subsection{La energía de una cuerda}\label{sec-hip-energia-def}

Una cuerda que vibra tiene dos formas de guardar energía. Cada pedacito
se mueve, y eso es energía \textbf{cinética}:
\(\tfrac12(\rho\,dx)\,u_t^2\). Y la cuerda está más estirada que en
reposo, y eso es energía \textbf{potencial elástica}: el pedacito que en
reposo medía \(dx\) ahora mide \[
\sqrt{dx^2 + du^2} = \sqrt{1+u_x^2}\;dx \approx \Big(1 + \tfrac12u_x^2\Big)dx ,
\] así que se estiró \(\tfrac12u_x^2\,dx\), y para estirarlo contra la
tensión \(T\) hubo que hacer un trabajo \(\tfrac12T\,u_x^2\,dx\). En un
medio con rigidez variable, \(T\) se reemplaza por \(k(x)\).

\begin{definition}[Energía de la
cuerda]\protect\hypertarget{def-hip-energia}{}\label{def-hip-energia}

Para una solución de \ref{eq-hip-inhomogenea} en \([0,l]\) se llama
\textbf{energía} a

\begin{equation}\protect\phantomsection\label{eq-hip-energia}{
E(t) = \int_0^l \Big[\tfrac12\,\rho(x)\,u_t^2 + \tfrac12\,k(x)\,u_x^2\Big]\,dx .
}\end{equation}

El primer término es la energía \textbf{cinética} y el segundo la
\textbf{potencial}. Los integrandos se llaman \textbf{densidades de
energía}.

\end{definition}

Las dos densidades son no negativas, porque \(\rho>0\) y \(k>0\). Esto
se va a usar enseguida.

\subsection{Adónde va la energía}\label{sec-hip-balance}

\begin{proposition}[Balance de
energía]\protect\hypertarget{prp-hip-balance}{}\label{prp-hip-balance}

Si \(u\) es una solución de clase \(C^2\) de \ref{eq-hip-inhomogenea} en
\([0,l]\), entonces

\begin{equation}\protect\phantomsection\label{eq-hip-balance}{
\frac{dE}{dt} \;=\; \int_0^l F\,u_t\,dx \;+\; \Big[\,k\,u_x\,u_t\,\Big]_{x=0}^{x=l} .
}\end{equation}

\end{proposition}

\begin{proof}
Derivando bajo el signo integral, cosa que se puede hacer porque \(u\)
es \(C^2\), \[
\frac{dE}{dt} = \int_0^l \big[\rho\,u_t\,u_{tt} + k\,u_x\,u_{xt}\big]\,dx .
\] En el primer término reemplazamos \(\rho\,u_{tt}\) usando la
ecuación, \(\rho u_{tt} = (k u_x)_x + F\): \[
\frac{dE}{dt} = \int_0^l \big[u_t\,(k u_x)_x + k\,u_x\,u_{xt}\big]\,dx + \int_0^l F\,u_t\,dx .
\] El corchete es la derivada de un producto:
\(\partial_x\big(k\,u_x\,u_t\big) = (k u_x)_x\,u_t + k\,u_x\,u_{tx}\).
Así que la primera integral es \(\big[k\,u_x\,u_t\big]_0^l\), por el
teorema fundamental del cálculo.
\end{proof}

Cada término de \ref{eq-hip-balance} tiene una lectura física, y
conviene decirla porque es la que se recuerda:

\begin{itemize}
\tightlist
\item
  \(\int F\,u_t\,dx\) es la \textbf{potencia de la fuerza externa}:
  fuerza por velocidad, sumada sobre toda la cuerda.
\item
  \(k\,u_x\,u_t\) en un extremo es la \textbf{potencia que entra o sale
  por el borde}: la fuerza que la cuerda hace sobre lo que la sujeta
  (\(\pm k u_x\), la componente vertical de la tensión) por la velocidad
  del extremo.
\end{itemize}

Si no hay fuerza externa y \textbf{cada extremo está fijo o libre}, los
dos términos se anulan: en un extremo fijo, \(u=0\) para todo \(t\) y
entonces \(u_t=0\); en un extremo libre, \(u_x=0\). Queda

\[
\frac{dE}{dt} = 0 .
\]

\begin{tcolorbox}[enhanced jigsaw, arc=.35mm, bottomrule=.15mm, breakable, colback=white, colframe=quarto-callout-important-color-frame, left=2mm, leftrule=.75mm, opacityback=0, rightrule=.15mm, toprule=.15mm]
\begin{minipage}[t]{5.5mm}
\textcolor{quarto-callout-important-color}{\faExclamation}
\end{minipage}%
\begin{minipage}[t]{\textwidth - 5.5mm}

\vspace{-3mm}\textbf{La energía de una cuerda ideal se conserva}\vspace{3mm}

Sin fuerzas externas y con extremos fijos o libres, \textbf{la energía
total no cambia nunca}. Lo que cambia es cómo se reparte: entre cinética
y potencial, y entre un lugar de la cuerda y otro. Es la razón por la
que una cuerda ideal no deja de sonar, y la razón del recuadro «Ojo con
el caso 3» del sección~\ref{sec-hip-casos-limite}.

\end{minipage}%
\end{tcolorbox}

\emph{(En la versión web hay acá un gráfico interactivo: una cuerda
pulsada, con un slider de tiempo. Debajo de la forma se ven la densidad
de energía cinética y la potencial a lo largo de la cuerda, que son
escalones que viajan con las esquinas, y al costado tres barras: la
energía cinética, la potencial y la total. Las dos primeras suben y
bajan; la total no se mueve.)}

\subsection{El teorema de unicidad}\label{sec-hip-unicidad}

\begin{theorem}[Unicidad del problema mixto (Teorema 9.1.1 de
González)]\protect\hypertarget{thm-hip-unicidad}{}\label{thm-hip-unicidad}

Supongamos que

\begin{enumerate}
\def\labelenumi{\arabic{enumi}.}
\tightlist
\item
  los coeficientes \(\rho(x)\) y \(k(x)\) son continuos en \([0,l]\),
  con \(\rho(x)>0\) y \(k(x)>0\);
\item
  se busca \(u\) de clase \(C^2\) en \([0,l]\times[0,\infty)\).
\end{enumerate}

Entonces existe \textbf{a lo sumo una} función \(u(x,t)\) que satisface

\[
\rho(x)\,u_{tt} = \big(k(x)\,u_x\big)_x + F(x,t), \qquad 0<x<l,\ t>0,
\]

con las condiciones iniciales \(u(x,0) = \varphi(x)\),
\(u_t(x,0) = \psi(x)\) y las de contorno \(u(0,t) = \mu_1(t)\),
\(u(l,t) = \mu_2(t)\).

\end{theorem}

\begin{proof}
Supongamos que hay dos, \(u_1\) y \(u_2\), y llamemos \(w = u_1 - u_2\).
Como la ecuación es lineal, \(w\) cumple la ecuación \textbf{sin fuente}
---las \(F\) se restan y se cancelan---, con datos \textbf{nulos}: \[
\rho\,w_{tt} = (k\,w_x)_x, \qquad w(x,0) = 0,\quad w_t(x,0) = 0, \qquad w(0,t) = w(l,t) = 0 .
\] Sea \(E(t)\) la energía de \(w\). Como los extremos están fijos y no
hay fuente, por Proposición~\ref{prp-hip-balance} \(E\) es constante. Y
en \(t=0\) vale cero: \(w_t(x,0) = 0\) por la segunda condición inicial,
y \(w_x(x,0) = 0\) porque \(w(x,0)\equiv0\) se puede derivar en \(x\).
Entonces \[
E(t) = \int_0^l \Big[\tfrac12\rho\,w_t^2 + \tfrac12 k\,w_x^2\Big]dx = 0 \qquad\text{para todo } t\ge0 .
\] El integrando es continuo y no negativo, y su integral es cero: tiene
que ser idénticamente cero. Como \(\rho>0\) y \(k>0\), eso obliga a
\(w_t\equiv0\) y \(w_x\equiv0\). Una función con las dos derivadas
parciales nulas es constante, y como \(w(x,0)=0\), la constante es cero.
Luego \(u_1 = u_2\).
\end{proof}

La demostración es corta, pero cada hipótesis está usada en un lugar
preciso. Vale la pena verlo con la tabla del
sección~\ref{sec-patologias} en mente:

\begin{longtable}[]{@{}
  >{\raggedright\arraybackslash}p{(\linewidth - 4\tabcolsep) * \real{0.2200}}
  >{\raggedright\arraybackslash}p{(\linewidth - 4\tabcolsep) * \real{0.3800}}
  >{\raggedright\arraybackslash}p{(\linewidth - 4\tabcolsep) * \real{0.4000}}@{}}
\caption{Qué hace cada hipótesis del
Teorema~\ref{thm-hip-unicidad}.}\label{tbl-hip-unicidad}\tabularnewline
\toprule\noalign{}
\begin{minipage}[b]{\linewidth}\raggedright
Hipótesis
\end{minipage} & \begin{minipage}[b]{\linewidth}\raggedright
Dónde se usa
\end{minipage} & \begin{minipage}[b]{\linewidth}\raggedright
Qué pasa si falta
\end{minipage} \\
\midrule\noalign{}
\endfirsthead
\toprule\noalign{}
\begin{minipage}[b]{\linewidth}\raggedright
Hipótesis
\end{minipage} & \begin{minipage}[b]{\linewidth}\raggedright
Dónde se usa
\end{minipage} & \begin{minipage}[b]{\linewidth}\raggedright
Qué pasa si falta
\end{minipage} \\
\midrule\noalign{}
\endhead
\bottomrule\noalign{}
\endlastfoot
La ecuación es \textbf{lineal} & la diferencia de dos soluciones es
solución de la homogénea & en una no lineal, \(u_1-u_2\) no cumple
ninguna ecuación útil \\
\(u\) es \(C^2\) & para derivar bajo la integral e integrar por partes &
con quiebres, como las cuerdas pulsadas, hay que pasar a soluciones
débiles (sección~\ref{sec-generalizadas}) \\
\(\rho > 0\) & para concluir \(w_t = 0\) de \(E = 0\) & un pedazo sin
masa podría moverse sin guardar energía cinética \\
\(k > 0\) & para concluir \(w_x = 0\) de \(E = 0\) & un pedazo sin
rigidez podría deformarse sin guardar energía potencial \\
extremos fijos (o libres) & para que el término de borde de
\ref{eq-hip-balance} sea cero & con un borde que entrega energía a su
antojo, la cuenta no cierra \\
\end{longtable}

\begin{tcolorbox}[enhanced jigsaw, arc=.35mm, bottomrule=.15mm, breakable, colback=white, colframe=quarto-callout-tip-color-frame, left=2mm, leftrule=.75mm, opacityback=0, rightrule=.15mm, toprule=.15mm]
\begin{minipage}[t]{5.5mm}
\textcolor{quarto-callout-tip-color}{\faLightbulb}
\end{minipage}%
\begin{minipage}[t]{\textwidth - 5.5mm}

\vspace{-3mm}\textbf{Por qué este es un método y no un truco}\vspace{3mm}

No resolvimos nada. No hay fórmula para la solución de
\ref{eq-hip-inhomogenea} con \(\rho\) y \(k\) cualesquiera, y sin
embargo sabemos que si existe, es una sola. El argumento funciona igual
en dos y tres dimensiones, cambiando \(u_x^2\) por \(|\nabla u|^2\) y el
teorema fundamental del cálculo por el teorema de la divergencia. Y en
el capítulo 10 lo vamos a repetir para la ecuación del calor, donde la
«energía» \(\int u^2\) ya no se conserva sino que \textbf{decrece}: esa
diferencia es la irreversibilidad del calor.

\end{minipage}%
\end{tcolorbox}

El extremo \(x=l\) de una cuerda está unido a un \textbf{amortiguador}:
un pistón en aceite que hace una fuerza vertical proporcional a la
velocidad, en contra del movimiento. La condición de contorno es
\(T\,u_x(l,t) = -\beta\,u_t(l,t)\) con \(\beta>0\), y el otro extremo
está fijo. ¿Qué pasa con la energía?

\begin{itemize}
\tightlist
\item
  {Se conserva, como siempre que no hay fuerza externa}
\item
  {Decrece: \(dE/dt = -\beta\,u_t(l,t)^2 \le 0\)}
\item
  {Crece, porque el amortiguador empuja la cuerda}
\item
  {Depende del signo de \(u_x(l,t)\)}
\end{itemize}

En \ref{eq-hip-balance}, con \(F = 0\) y el extremo izquierdo fijo,
queda solo el término del borde derecho:
\(T\,u_x(l,t)\,u_t(l,t) = -\beta\,u_t(l,t)^2\). Es un cuadrado con signo
menos: la energía nunca crece y baja cada vez que el extremo se mueve.
Es lo que hace el amortiguador: convierte energía de la cuerda en calor
del aceite.

Hay un valor especial, \(\beta = \sqrt{T\rho}\), en el que el
amortiguador \textbf{se traga la onda entera} sin reflejar nada: la
cuerda se comporta como si siguiera infinitamente más allá de \(x=l\).
Lo vemos en el sección~\ref{sec-hip-reflexiones}.

\begin{exercise}[]\protect\hypertarget{exr-hip-robin}{}\label{exr-hip-robin}

El extremo \(x=l\) de una cuerda homogénea está atado a un
\textbf{resorte} de constante \(\kappa>0\), de modo que
\(T\,u_x(l,t) = -\kappa\,u(l,t)\), y el extremo \(x=0\) está fijo. No
hay fuerza externa.

\begin{enumerate}
\def\labelenumi{\alph{enumi})}
\item
  Mostrar que la energía de la cuerda \textbf{no} se conserva, pero que
  \(\;\mathcal E(t) = E(t) + \tfrac12\,\kappa\,u(l,t)^2\;\) sí.
  Interpretar el término agregado.
\item
  Deducir que el problema con esta condición de contorno también tiene a
  lo sumo una solución.
\end{enumerate}

\end{exercise}

\begin{tcolorbox}[enhanced jigsaw, arc=.35mm, bottomrule=.15mm, bottomtitle=1mm, breakable, colback=white, colbacktitle=quarto-callout-tip-color!10!white, colframe=quarto-callout-tip-color-frame, coltitle=black, left=2mm, leftrule=.75mm, opacityback=0, opacitybacktitle=0.6, rightrule=.15mm, title=\textcolor{quarto-callout-tip-color}{\faLightbulb}\hspace{0.5em}{Solución}, titlerule=0mm, toprule=.15mm, toptitle=1mm]

\textbf{a)} Por \ref{eq-hip-balance}, con \(F = 0\) y \(u_t(0,t) = 0\),
\[
\frac{dE}{dt} = T\,u_x(l,t)\,u_t(l,t) = -\kappa\,u(l,t)\,u_t(l,t) = -\frac{d}{dt}\Big[\tfrac12\kappa\,u(l,t)^2\Big].
\] Pasando de miembro,
\(\dfrac{d}{dt}\Big[E + \tfrac12\kappa\,u(l,t)^2\Big] = 0\).

El término agregado es la \textbf{energía potencial del resorte}. La
energía de la cuerda sola no se conserva porque la cuerda y el resorte
se la pasan de uno al otro; la del sistema completo sí. Es lo que pasa
siempre con el término de borde de \ref{eq-hip-balance}: mide lo que se
intercambia con algo que no estamos modelando como parte de la cuerda.

\textbf{b)} Si \(u_1\) y \(u_2\) son dos soluciones, \(w = u_1 - u_2\)
cumple la ecuación homogénea con datos iniciales nulos, \(w(0,t) = 0\) y
\(T\,w_x(l,t) = -\kappa\,w(l,t)\). Su \(\mathcal E\) se conserva y vale
cero en \(t=0\), así que \(\mathcal E(t) = 0\) para todo \(t\). Como
\(\kappa > 0\), los tres sumandos son no negativos, y cada uno es cero.
En particular la energía de la cuerda es cero, y se concluye igual que
en el Teorema~\ref{thm-hip-unicidad}.

\textbf{El signo importa.} Si fuera \(\kappa<0\) ---un «resorte» que
empuja hacia afuera---, \(\mathcal E\) podría ser cero con la cuerda
moviéndose, porque el término del resorte sería negativo, y la
demostración no funcionaría.

\end{tcolorbox}

\begin{exercise}[]\protect\hypertarget{exr-hip-equiparticion}{}\label{exr-hip-equiparticion}

Mostrar que en una onda viajera \(u(x,t) = g(x-at)\) sobre una cuerda
homogénea la densidad de energía cinética es igual a la potencial
\textbf{en cada punto y en cada instante}. ¿Qué dice esto sobre la
energía de las dos mitades en que se parte una cuerda pulsada?

\end{exercise}

\begin{tcolorbox}[enhanced jigsaw, arc=.35mm, bottomrule=.15mm, bottomtitle=1mm, breakable, colback=white, colbacktitle=quarto-callout-tip-color!10!white, colframe=quarto-callout-tip-color-frame, coltitle=black, left=2mm, leftrule=.75mm, opacityback=0, opacitybacktitle=0.6, rightrule=.15mm, title=\textcolor{quarto-callout-tip-color}{\faLightbulb}\hspace{0.5em}{Solución}, titlerule=0mm, toprule=.15mm, toptitle=1mm]

Con \(u = g(x-at)\) se tiene \(u_t = -a\,g'(x-at)\) y
\(u_x = g'(x-at)\). Entonces \[
\tfrac12\rho\,u_t^2 = \tfrac12\rho\,a^2\,g'^2 = \tfrac12 T\,g'^2 = \tfrac12 T\,u_x^2 ,
\] usando \(\rho a^2 = T\), que es \ref{eq-hip-velocidad}. Las dos
densidades coinciden punto a punto. Lo mismo pasa con \(g(x+at)\).

Una cuerda pulsada arranca con toda la energía en forma potencial,
porque está quieta. Por d'Alembert se parte en dos mitades,
\(\tfrac12\varphi(x+at)\) y \(\tfrac12\varphi(x-at)\), que cuando se
separan son dos ondas viajeras. Cada una tiene energía potencial
\(\tfrac14\) de la inicial ---la pendiente es la mitad, y la energía va
con el cuadrado--- y la misma cantidad de cinética. En total,
\(4\times\tfrac14 = 1\): \textbf{la energía inicial, que era toda
potencial, termina mitad cinética y mitad potencial, repartida en partes
iguales entre las dos ondas.} Es un caso de \emph{equipartición}.

Ojo: esto vale para un pulso localizado en una cuerda infinita, una vez
que las dos mitades se separaron. En la cuerda finita del gráfico
interactivo de más arriba el triángulo ocupa toda la cuerda, las mitades
se pisan y rebotan todo el tiempo, y la proporción entre cinética y
potencial va cambiando: la equipartición no se ve ahí.

\end{tcolorbox}

\section{La fórmula de d'Alembert, a fondo}\label{sec-hip-dalembert}

En el Ejercicio~\ref{exr-dalembert} dedujimos la fórmula a partir de la
forma canónica. La cuenta es exactamente la de la sección 9.2.1 de
González y no la repetimos. Lo que hacemos acá es lo que quedó
pendiente: enunciarla como teorema, con hipótesis, leerla con cuidado, y
ver qué más se puede sacar de ella.

\begin{tcolorbox}[enhanced jigsaw, arc=.35mm, bottomrule=.15mm, breakable, colback=white, colframe=quarto-callout-warning-color-frame, left=2mm, leftrule=.75mm, opacityback=0, rightrule=.15mm, toprule=.15mm]
\begin{minipage}[t]{5.5mm}
\textcolor{quarto-callout-warning-color}{\faExclamationTriangle}
\end{minipage}%
\begin{minipage}[t]{\textwidth - 5.5mm}

\vspace{-3mm}\textbf{Ojo con el libro de la cátedra (sección 9.2.1)}\vspace{3mm}

Dos erratas de la deducción que conviene tener marcadas si se estudia de
ahí:

\begin{itemize}
\tightlist
\item
  Al despejar las variables originales dice
  \(t = \dfrac{\xi+\eta}{2a}\). Tiene que ser
  \(t = \dfrac{\xi-\eta}{2a}\), porque
  \(\xi - \eta = (x+at)-(x-at) = 2at\).
\item
  En las fórmulas de \(f_1\), \(f_2\) y en la (9.10) aparece
  \(\dfrac{1}{2\alpha}\). Es \(\dfrac{1}{2a}\): la \(\alpha\) es la
  variable de integración y no puede estar afuera de la integral. El
  mismo error se repite en la sección 9.2.2.
\end{itemize}

\end{minipage}%
\end{tcolorbox}

\begin{theorem}[Fórmula de
d'Alembert]\protect\hypertarget{thm-hip-dalembert}{}\label{thm-hip-dalembert}

Sean \(\varphi\in C^2(\mathbb R)\) y \(\psi\in C^1(\mathbb R)\). El
problema de Cauchy

\[
u_{tt} = a^2u_{xx},\quad x\in\mathbb R,\ t>0;\qquad u(x,0) = \varphi(x),\quad u_t(x,0) = \psi(x),
\]

tiene una \textbf{única} solución de clase \(C^2\), que es

\begin{equation}\protect\phantomsection\label{eq-hip-dalembert}{
u(x,t) \;=\; \underbrace{\frac{\varphi(x+at)+\varphi(x-at)}{2}}_{u_1:\ \text{la forma se parte}}
\;+\; \underbrace{\frac{1}{2a}\int_{x-at}^{x+at}\psi(s)\,ds}_{u_2:\ \text{la velocidad se acumula}} .
}\end{equation}

\end{theorem}

\begin{proof}
\textbf{Unicidad.} Toda solución \(C^2\) cumple \(u_{\xi\eta}=0\) en las
coordenadas \(\xi=x+at\), \(\eta = x-at\)
(sección~\ref{sec-idea-canonica}), así que es de la forma
\(f(x+at)+g(x-at)\). Las condiciones iniciales determinan \(f\) y \(g\)
salvo una constante que se cancela (Ejercicio~\ref{exr-dalembert}). No
hay lugar para una segunda solución.

\textbf{Existencia.} Hay que verificar que \ref{eq-hip-dalembert} cumple
la ecuación y las condiciones: es el Ejercicio~\ref{exr-hip-verificar}.
\end{proof}

Con las hipótesis de regularidad pasa algo que vale la pena señalar:
\(u\) tiene \textbf{exactamente} la regularidad de los datos, ni más ni
menos. Si \(\varphi\) tiene una esquina, \(u\) tiene esquinas; si
\(\psi\) tiene un salto, \(u\) tiene esquinas. Eso es lo contrario de la
ecuación del calor, que suaviza todo al instante
(sección~\ref{sec-tres-mundos}), y tiene un enunciado preciso:

\begin{tcolorbox}[enhanced jigsaw, arc=.35mm, bottomrule=.15mm, breakable, colback=white, colframe=quarto-callout-important-color-frame, left=2mm, leftrule=.75mm, opacityback=0, rightrule=.15mm, toprule=.15mm]
\begin{minipage}[t]{5.5mm}
\textcolor{quarto-callout-important-color}{\faExclamation}
\end{minipage}%
\begin{minipage}[t]{\textwidth - 5.5mm}

\vspace{-3mm}\textbf{Las singularidades viajan por las características}\vspace{3mm}

Si el dato tiene una singularidad ---una esquina, un salto--- en un
punto \(x_0\), la solución tiene singularidades \textbf{solo} sobre las
dos características que salen de \(x_0\), \(x = x_0 \pm at\), y en
ningún otro lado. Se lee directamente de \ref{eq-hip-dalembert}:
\(\varphi(x\pm at)\) es suave salvo cuando \(x\pm at = x_0\).

Esto es lo que hace que un trueno sea un estallido y no un rumor, y lo
que permite medir la distancia a un sismo: el frente, que es la
singularidad, llega entero.

\end{minipage}%
\end{tcolorbox}

\subsection{Para jugar: la fórmula por
dentro}\label{sec-hip-lienzo-dalembert}

El gráfico de abajo pone la fórmula y su dominio de dependencia uno
arriba del otro. Arriba, la cuerda en el instante \(t\). Abajo, el plano
\((x,t)\): el punto elegido, las dos características que bajan de él, y
el pedazo del eje \(t=0\) del que depende su valor.

\emph{(En la versión web hay acá un gráfico interactivo en dos paneles.
Arriba, la cuerda en el instante \(t\), con la forma inicial ---un
triángulo--- partida en dos mitades punteadas y la contribución de una
velocidad inicial en escalón. Abajo, el plano \((x,t)\) con un punto
\((x_0,t)\), las dos características que bajan de él y su dominio de
dependencia marcado sobre el eje; el valor de \(u(x_0,t)\) se descompone
en sus dos términos. Se pueden mover \(t\), \(x_0\) y las alturas de la
forma y de la velocidad iniciales.)}

Una cuerda infinita con \(a = 1\) está en reposo y sin deformar, y en
\(t=0\) recibe un golpe que le da velocidad \(\psi = 1\) en \([-1, 1]\)
y cero en el resto. ¿Cuánto vale \(u(3;\ 2{,}5)\)?

Como \(\varphi = 0\), solo queda la integral:
\[u(3;\,2{,}5) = \frac12\int_{3-2{,}5}^{3+2{,}5}\psi(s)\,ds = \frac12\int_{0{,}5}^{5{,}5}\psi(s)\,ds .\]
El dominio de dependencia es \([0{,}5;\ 5{,}5]\) y \(\psi\) vale \(1\)
solo en \([-1,1]\). Se pisan en \([0{,}5;\ 1]\), de largo \(0{,}5\), así
que \(u = \tfrac12\cdot 0{,}5 = 0{,}25\).

Fijate lo que dice esto: el punto \(x=3\) ya se movió, pero \textbf{no a
la altura máxima}. Esa altura es \(\tfrac12\int_{-1}^1\psi = 1\) y la
alcanzan los puntos cuyo dominio de dependencia contiene \textbf{entero}
al intervalo golpeado. Es la meseta que se ve en el gráfico de arriba.

¿Qué datos iniciales producen una solución que se queda \textbf{quieta
en \(x=0\) para siempre}, es decir \(u(0,t) = 0\) para todo \(t\)?

\begin{itemize}
\tightlist
\item
  {Los que tienen \(\varphi\) y \(\psi\) impares}
\item
  {Los que tienen \(\varphi\) y \(\psi\) pares}
\item
  {Los que tienen \(\varphi(0) = 0\) y \(\psi(0) = 0\)}
\item
  {Ninguno: una cuerda infinita no tiene puntos fijos}
\end{itemize}

Poniendo \(x = 0\) en \ref{eq-hip-dalembert},
\[u(0,t) = \frac{\varphi(at)+\varphi(-at)}{2} + \frac{1}{2a}\int_{-at}^{at}\psi(s)\,ds .\]
Si \(\varphi\) es impar, el primer término es cero; si \(\psi\) es
impar, la integral sobre un intervalo simétrico es cero. Así que datos
impares dan \(u(0,t) = 0\) siempre.

La tercera opción no alcanza: que el origen arranque quieto no impide
que después le llegue una onda desde otro lado. Y la última opción es
falsa: con datos impares el origen es un punto fijo aunque la cuerda sea
infinita.

Guardá esta observación: es \textbf{la} idea del
sección~\ref{sec-hip-reflexiones}. Si queremos una cuerda con un extremo
fijo en \(x=0\), alcanza con inventar del otro lado una cuerda
imaginaria con datos impares.

\subsection{La regla del paralelogramo}\label{sec-hip-paralelogramo}

La fórmula de d'Alembert es una consecuencia de que la solución general
sea \(f(x+at)+g(x-at)\). Hay otra consecuencia, geométrica, que no
necesita conocer ni \(f\) ni \(g\), y que va a servir para resolver
problemas con bordes.

\begin{proposition}[Regla del
paralelogramo]\protect\hypertarget{prp-hip-paralelogramo}{}\label{prp-hip-paralelogramo}

Sea \(ABCD\) un paralelogramo en el plano \((x,t)\) con los lados sobre
características: \(AB\) y \(DC\) paralelos a \(x-at=\text{const}\), y
\(AD\) y \(BC\) paralelos a \(x+at=\text{const}\). Si \(u\) es solución
de \(u_{tt}=a^2u_{xx}\) en una región que lo contiene,

\begin{equation}\protect\phantomsection\label{eq-hip-paralelogramo}{
u(A) + u(C) = u(B) + u(D) .
}\end{equation}

\end{proposition}

\begin{proof}
Escribamos \(u = f(\xi) + g(\eta)\) con \(\xi = x+at\), \(\eta=x-at\).
Sobre un lado paralelo a \(\eta=\text{const}\) el valor de \(\eta\) no
cambia, y lo mismo con \(\xi\). Si \(A\) tiene coordenadas
características \((\xi_1,\eta_1)\) y \(C\) tiene \((\xi_2,\eta_2)\),
entonces los otros dos vértices son \((\xi_2,\eta_1)\) y
\((\xi_1,\eta_2)\). Sumando, \[
u(A) + u(C) = f(\xi_1)+g(\eta_1)+f(\xi_2)+g(\eta_2) = u(B) + u(D) .
\]
\end{proof}

\begin{figure}

\centering{

\pandocbounded{\includegraphics[keepaspectratio]{09-hiperbolicas_files/figure-pdf/fig-hip-paralelogramo-output-1.pdf}}

}

\caption{\label{fig-hip-paralelogramo}La regla del paralelogramo.
\textbf{Izquierda:} en el plano \((x,t)\) el paralelogramo tiene los
lados sobre las dos familias de características, que acá son rectas de
pendiente \(\pm1/a\). \textbf{Derecha:} el mismo paralelogramo en las
coordenadas características \((\xi,\eta)\) es un rectángulo con lados
paralelos a los ejes. Como \(u = f(\xi)+g(\eta)\), cada vértice lleva un
valor de \(f\) y uno de \(g\), y los vértices opuestos suman lo mismo.}

\end{figure}%

Parece una curiosidad, pero es una herramienta de cálculo: si se conoce
\(u\) en tres vértices, se conoce en el cuarto sin resolver nada. En el
sección~\ref{sec-hip-reflexiones} la vamos a usar para saber qué pasa
cerca de un borde, donde la fórmula de d'Alembert no se puede aplicar
tal cual.

\begin{exercise}[]\protect\hypertarget{exr-hip-goursat}{}\label{exr-hip-goursat}

\textbf{El problema de Goursat.} Hallar la solución de
\(u_{tt} = a^2u_{xx}\) en el cono \(t>|x|/a\) sabiendo sus valores
\textbf{sobre las dos características} que salen del origen:
\[u(x,x/a) = \alpha(x)\ \ (x\ge0), \qquad u(x,-x/a) = \beta(x)\ \ (x\le0), \qquad \alpha(0)=\beta(0).\]
Usar la regla del paralelogramo con un vértice en el origen. ¿Por qué
hace falta la condición \(\alpha(0)=\beta(0)\)?

\end{exercise}

\begin{tcolorbox}[enhanced jigsaw, arc=.35mm, bottomrule=.15mm, bottomtitle=1mm, breakable, colback=white, colbacktitle=quarto-callout-tip-color!10!white, colframe=quarto-callout-tip-color-frame, coltitle=black, left=2mm, leftrule=.75mm, opacityback=0, opacitybacktitle=0.6, rightrule=.15mm, title=\textcolor{quarto-callout-tip-color}{\faLightbulb}\hspace{0.5em}{Solución}, titlerule=0mm, toprule=.15mm, toptitle=1mm]

Sea \(C=(x,t)\) un punto del cono. Armamos el paralelogramo con un
vértice \(A\) en el origen y los otros dos, \(B\) y \(D\), sobre las
características dadas. Desde \(C\) bajamos por la recta
\(x-at=\text{const}\) hasta cortar la característica \(x = -at\), y por
la recta \(x+at = \text{const}\) hasta cortar la característica
\(x=at\).

\begin{itemize}
\tightlist
\item
  La recta \(x' + at' = x + at\) corta a \(x' = at'\) en
  \(x' = \tfrac{x+at}{2}\). Ahí \(u = \alpha\big(\tfrac{x+at}{2}\big)\).
\item
  La recta \(x' - at' = x - at\) corta a \(x' = -at'\) en
  \(x' = \tfrac{x-at}{2}\). Ahí \(u = \beta\big(\tfrac{x-at}{2}\big)\).
\end{itemize}

Por \ref{eq-hip-paralelogramo}, con \(u(A) = u(0,0) = \alpha(0)\), \[
\boxed{\;u(x,t) = \alpha\!\Big(\frac{x+at}{2}\Big) + \beta\!\Big(\frac{x-at}{2}\Big) - \alpha(0).\;}
\]

\textbf{Verificación.} Sobre \(x=at\):
\(\alpha(x) + \beta(0) - \alpha(0) = \alpha(x)\). Sobre \(x=-at\):
\(\alpha(0)+\beta(x)-\alpha(0) = \beta(x)\). Y es de la forma
\(f(x+at)+g(x-at)\), así que cumple la ecuación.

\textbf{La condición de compatibilidad.} Las dos características se
cruzan en el origen, y ahí \(u\) tiene un solo valor. Si
\(\alpha(0)\neq\beta(0)\) los datos se contradicen y no hay solución.

\textbf{Lo que se aprende.} En la ecuación de onda se puede dar el dato
\emph{sobre} las características, y alcanza con \textbf{uno} en cada
una. En el sección~\ref{sec-patologia-transversalidad} vimos que en
primer orden un dato sobre una característica era un desastre; en
segundo orden, dar el valor sobre las dos familias es un problema bien
planteado. No es contradictorio: lo que no se puede es dar el valor
\textbf{y la derivada} sobre una característica, que es lo que sería un
problema de Cauchy ahí.

\end{tcolorbox}

\subsection{Estabilidad}\label{sec-hip-estabilidad}

El Teorema~\ref{thm-hip-dalembert} da existencia y unicidad. Falta el
tercer pedido de Hadamard: que la solución dependa con continuidad de
los datos. Es la sección 9.2.2 de González.

\begin{proposition}[Dependencia continua de los
datos]\protect\hypertarget{prp-hip-estabilidad}{}\label{prp-hip-estabilidad}

Sean \(u_1\) y \(u_2\) las soluciones del problema de Cauchy con datos
\((\varphi_1,\psi_1)\) y \((\varphi_2,\psi_2)\), y sea \(t_0>0\). Si
para todo \(x\)

\[
|\varphi_1(x)-\varphi_2(x)| < \delta, \qquad |\psi_1(x)-\psi_2(x)| < \delta,
\]

entonces para todo \(x\) y todo \(0\le t\le t_0\)

\begin{equation}\protect\phantomsection\label{eq-hip-estabilidad}{
|u_1(x,t) - u_2(x,t)| \;\le\; \delta\,(1+t_0) .
}\end{equation}

En particular, dado \(\varepsilon>0\), alcanza con tomar
\(\delta = \varepsilon/(1+t_0)\) para que las soluciones difieran en
menos de \(\varepsilon\) hasta el tiempo \(t_0\).

\end{proposition}

\begin{proof}
Restando las dos fórmulas de d'Alembert y usando la desigualdad
triangular, \[
|u_1-u_2| \le \frac{|\varphi_1-\varphi_2|(x+at)}{2} + \frac{|\varphi_1-\varphi_2|(x-at)}{2}
+ \frac{1}{2a}\int_{x-at}^{x+at}|\psi_1-\psi_2|\,ds
< \frac\delta2 + \frac\delta2 + \frac{1}{2a}\,\delta\,(2at) = \delta(1+t) \le \delta(1+t_0).
\]
\end{proof}

La cota \textbf{crece con el tiempo}, y no es un defecto de la
demostración. Si los datos coinciden salvo que
\(\psi_2 = \psi_1 + \delta\) en todos lados, la diferencia \(u_2 - u_1\)
es la solución con \(\varphi = 0\) y \(\psi = \delta\), que es
exactamente \(\delta\,t\): una cuerda entera que se va desplazando a
velocidad \(\delta\). Un error chico en la velocidad produce un error en
la posición que crece como el tiempo, igual que en la mecánica de una
partícula. Lo que importa es que \textbf{crece linealmente y no
exponencialmente}, y ahí está toda la diferencia con el ejemplo de
Hadamard.

\begin{figure}

\centering{

\pandocbounded{\includegraphics[keepaspectratio]{09-hiperbolicas_files/figure-pdf/fig-hip-estabilidad-output-1.pdf}}

}

\caption{\label{fig-hip-estabilidad}Dos formas de depender de los datos.
\textbf{Izquierda:} para la ecuación de onda, una perturbación de tamaño
\(\delta = 10^{-3}\) en los datos. Si se perturba solo la forma, el
error nunca pasa de \(\delta\) (las ondas se llevan la perturbación,
pero no la agrandan). Si se perturba la velocidad, crece como
\(\delta\,t\). La cota \ref{eq-hip-estabilidad}, en punteado, acota los
dos casos. \textbf{Derecha:} la misma perturbación en el ejemplo de
Hadamard del sección~\ref{sec-hadamard}, con la distancia a la
superficie en el lugar del tiempo, para una ondulación de longitud de
onda \(2\pi/3\). En escala logarítmica, la recta de la izquierda es la
curva de abajo, casi plana; la de Hadamard es una recta que cruza más de
doce órdenes de magnitud.}

\end{figure}%

\begin{tcolorbox}[enhanced jigsaw, arc=.35mm, bottomrule=.15mm, breakable, colback=white, colframe=quarto-callout-note-color-frame, left=2mm, leftrule=.75mm, opacityback=0, rightrule=.15mm, toprule=.15mm]
\begin{minipage}[t]{5.5mm}
\textcolor{quarto-callout-note-color}{\faInfo}
\end{minipage}%
\begin{minipage}[t]{\textwidth - 5.5mm}

\vspace{-3mm}\textbf{Lo que la estimación no dice}\vspace{3mm}

\ref{eq-hip-estabilidad} pide que los datos estén cerca \textbf{en todos
lados} y promete que las soluciones están cerca \textbf{en todos lados},
en la norma del supremo. Es la forma más simple de estabilidad, y
alcanza para la ecuación de onda en una dimensión. Pero no es la única
manera de medir cercanía: la energía del sección~\ref{sec-hip-energia}
es otra, y en tres dimensiones es la única que funciona bien, porque la
solución puede ser más «puntuda» que los datos ---lo vemos en el
sección~\ref{sec-hip-espacio}---.

\end{minipage}%
\end{tcolorbox}

\begin{exercise}[]\protect\hypertarget{exr-hip-verificar}{}\label{exr-hip-verificar}

Verificar que \ref{eq-hip-dalembert} es solución: que cumple la ecuación
y las dos condiciones iniciales. Mostrar además que si
\(\varphi\in C^2\) y \(\psi\in C^1\), entonces \(u\in C^2\), y que con
menos regularidad eso puede fallar.

\end{exercise}

\begin{tcolorbox}[enhanced jigsaw, arc=.35mm, bottomrule=.15mm, bottomtitle=1mm, breakable, colback=white, colbacktitle=quarto-callout-tip-color!10!white, colframe=quarto-callout-tip-color-frame, coltitle=black, left=2mm, leftrule=.75mm, opacityback=0, opacitybacktitle=0.6, rightrule=.15mm, title=\textcolor{quarto-callout-tip-color}{\faLightbulb}\hspace{0.5em}{Solución}, titlerule=0mm, toprule=.15mm, toptitle=1mm]

\textbf{El término de \(\varphi\).}
\(u_1 = \tfrac12[\varphi(x+at)+\varphi(x-at)]\) es de la forma
\(f(x+at)+g(x-at)\), y cualquier función de esa forma cumple la
ecuación: \(\partial_t^2 f(x+at) = a^2f''\) y
\(\partial_x^2 f(x+at) = f''\).

\textbf{El término de \(\psi\).} Sea \(\Psi\) una primitiva de \(\psi\).
Entonces \(u_2 = \tfrac{1}{2a}\big[\Psi(x+at) - \Psi(x-at)\big]\), que
también es de la forma \(f(x+at)+g(x-at)\). Cumple la ecuación.

\textbf{Las condiciones iniciales.} En \(t=0\):
\(u_1(x,0) = \varphi(x)\) y \(u_2(x,0) = 0\), porque la integral es
sobre un intervalo de largo cero. Derivando en \(t\) con la regla de
Leibniz, \[
u_t = \frac{a\varphi'(x+at) - a\varphi'(x-at)}{2} + \frac{1}{2a}\big[a\psi(x+at) + a\psi(x-at)\big],
\] y en \(t=0\) queda \(0 + \psi(x)\). Las dos condiciones se cumplen.

\textbf{La regularidad.} \(u_1\) tiene dos derivadas continuas si
\(\varphi\) las tiene. \(u_2\) es \(\Psi(x\pm at)\), y \(\Psi\) tiene
una derivada más que \(\psi\): si \(\psi\in C^1\), \(\Psi\in C^2\).

Si \(\varphi\) tiene una esquina ---la cuerda pulsada---, \(u_1\) tiene
esquinas y \(u_{xx}\) no existe sobre las características que salen de
la esquina. La fórmula sigue teniendo sentido y describe perfectamente
la física, pero \(u\) no es solución clásica: es una \textbf{solución
débil}, en el sentido del sección~\ref{sec-generalizadas}. Por eso el
Teorema~\ref{thm-hip-dalembert} tiene hipótesis de regularidad.

\end{tcolorbox}

\begin{exercise}[]\protect\hypertarget{exr-hip-estacionaria}{}\label{exr-hip-estacionaria}

Resolver el problema de Cauchy con
\(\varphi(x) = \operatorname{sen}(kx)\) y \(\psi(x) = 0\). Mostrar que
la solución es una \textbf{onda estacionaria} ---una forma fija que
cambia de tamaño--- y escribirla como suma de dos ondas viajeras. ¿Dónde
están los nodos?

\end{exercise}

\begin{tcolorbox}[enhanced jigsaw, arc=.35mm, bottomrule=.15mm, bottomtitle=1mm, breakable, colback=white, colbacktitle=quarto-callout-tip-color!10!white, colframe=quarto-callout-tip-color-frame, coltitle=black, left=2mm, leftrule=.75mm, opacityback=0, opacitybacktitle=0.6, rightrule=.15mm, title=\textcolor{quarto-callout-tip-color}{\faLightbulb}\hspace{0.5em}{Solución}, titlerule=0mm, toprule=.15mm, toptitle=1mm]

Por d'Alembert, y usando
\(\operatorname{sen}(A)+\operatorname{sen}(B) = 2\operatorname{sen}\tfrac{A+B}2\cos\tfrac{A-B}2\),
\[
u = \frac{\operatorname{sen}\big(k(x+at)\big) + \operatorname{sen}\big(k(x-at)\big)}{2}
= \operatorname{sen}(kx)\,\cos(kat) .
\] Es una forma fija, \(\operatorname{sen}(kx)\), multiplicada por un
factor que oscila en el tiempo con frecuencia \(\omega = ka\). Los nodos
son los ceros de \(\operatorname{sen}(kx)\), \(x = n\pi/k\): esos puntos
no se mueven nunca.

Leída de derecha a izquierda, la cuenta dice lo contrario: \textbf{una
onda estacionaria es la suma de dos ondas iguales que viajan en sentidos
opuestos.} Es el vínculo entre las dos herramientas de la
sección~\ref{sec-hip-pregunta}. Si ponemos \(k = n\pi/l\), los nodos
caen en \(x = 0\) y \(x = l\), y la onda estacionaria es el modo \(n\)
de la cuerda fija del sección~\ref{sec-cuerda-separacion}. La cuerda
finita que «late» en sus modos es una cuerda infinita donde dos ondas
viajeras se cruzan todo el tiempo. Lo desarrollamos en el
sección~\ref{sec-hip-reflexiones}.

\end{tcolorbox}

\section{Bordes: las ondas rebotan}\label{sec-hip-reflexiones}

La fórmula de d'Alembert resuelve la cuerda infinita. Una cuerda de
verdad termina en algún lado, y cuando una onda llega al final tiene que
pasar algo. Esta sección responde qué, y de paso une las dos imágenes de
la sección~\ref{sec-hip-pregunta}: la onda que viaja y el modo que late.

Video recomendado: \emph{Waves on springs --- reflection and
transmission} (NCPQ), \url{https://www.youtube.com/watch?v=ocLGb28_8UQ}.

\subsection{La cuerda semi-infinita con el extremo
fijo}\label{sec-hip-semi-fijo}

Es el caso 2 del sección~\ref{sec-hip-casos-limite} con \(\mu\equiv0\):

\begin{equation}\protect\phantomsection\label{eq-hip-semi}{
u_{tt} = a^2u_{xx}\quad (x>0,\ t>0), \qquad u(0,t) = 0, \qquad u(x,0) = \varphi(x),\quad u_t(x,0) = \psi(x)\quad (x\ge 0).
}\end{equation}

La fórmula de d'Alembert no se puede aplicar tal cual, por dos razones.
Pide datos en toda la recta, y acá solo los hay para \(x\ge0\). Y aunque
los tuviéramos, no sabe nada del borde.

La salida es la observación del segundo ejercicio del
sección~\ref{sec-hip-dalembert}: \textbf{si los datos son impares, la
solución de d'Alembert se queda quieta en \(x=0\)}. Entonces inventamos
una cuerda del otro lado, en \(x<0\), con los datos reflejados y
cambiados de signo. La cuerda completa es infinita, d'Alembert la
resuelve, y la solución vale cero en el origen para siempre: hace de
extremo fijo sin que nadie lo haya fijado.

\begin{definition}[Extensión impar y
par]\protect\hypertarget{def-hip-extension}{}\label{def-hip-extension}

Dada \(\varphi\) definida en \([0,\infty)\), su \textbf{extensión impar}
y su \textbf{extensión par} son las funciones definidas en toda la recta
por

\[
\varphi_{\text{i}}(x) = \begin{cases}\varphi(x), & x\ge0,\\ -\varphi(-x), & x<0,\end{cases}
\qquad
\varphi_{\text{p}}(x) = \begin{cases}\varphi(x), & x\ge0,\\ \varphi(-x), & x<0.\end{cases}
\]

\end{definition}

\begin{proposition}[Cuerda semi-infinita con extremo
fijo]\protect\hypertarget{prp-hip-semi-fijo}{}\label{prp-hip-semi-fijo}

La solución de \ref{eq-hip-semi} es la de d'Alembert con los datos
extendidos en forma impar, restringida a \(x\ge0\). Escrita solo con los
datos originales,

\begin{equation}\protect\phantomsection\label{eq-hip-semi-fijo}{
u(x,t) =
\begin{cases}
\dfrac{\varphi(x+at)+\varphi(x-at)}{2} + \dfrac{1}{2a}\displaystyle\int_{x-at}^{x+at}\psi(s)\,ds, & x\ge at,\\[3mm]
\dfrac{\varphi(x+at)-\varphi(at-x)}{2} + \dfrac{1}{2a}\displaystyle\int_{at-x}^{x+at}\psi(s)\,ds, & 0\le x< at.
\end{cases}
}\end{equation}

\end{proposition}

\begin{proof}
La extensión impar cumple la ecuación en toda la recta y, por el segundo
ejercicio del sección~\ref{sec-hip-dalembert}, vale cero en \(x=0\). En
\(x\ge0\) tiene los datos que pedimos. Por la unicidad, es \textbf{la}
solución.

Para escribirla sin la extensión: si \(x\ge at\), los dos puntos
\(x\pm at\) están en \(x\ge0\) y no hay nada que cambiar. Si \(x<at\),
el punto \(x-at\) es negativo, y ahí
\(\varphi_{\text{i}}(x-at) = -\varphi(at-x)\). En la integral, el tramo
\([x-at,\,at-x]\) es simétrico y la extensión impar de \(\psi\) integra
cero en él; queda solo el tramo \([at-x,\,x+at]\).
\end{proof}

La línea \(x = at\) que separa los dos casos es la característica que
sale del borde en \(t=0\). Por debajo, el punto no se enteró de que hay
un borde: su dominio de dependencia no llega al origen. Por encima, sí:
parte de lo que ve es un \textbf{eco}, lo que salió hacia la izquierda y
rebotó.

\begin{tcolorbox}[enhanced jigsaw, arc=.35mm, bottomrule=.15mm, breakable, colback=white, colframe=quarto-callout-important-color-frame, left=2mm, leftrule=.75mm, opacityback=0, rightrule=.15mm, toprule=.15mm]
\begin{minipage}[t]{5.5mm}
\textcolor{quarto-callout-important-color}{\faExclamation}
\end{minipage}%
\begin{minipage}[t]{\textwidth - 5.5mm}

\vspace{-3mm}\textbf{Qué le pasa a un pulso en un extremo fijo}\vspace{3mm}

Un pulso que viaja hacia el extremo fijo tiene, del otro lado del
espejo, un gemelo \textbf{invertido} que viaja hacia él. Cuando se
cruzan en \(x=0\) se cancelan exactamente ---por eso el extremo no se
mueve---, y después el gemelo sale a la cuerda real mientras el pulso
original sigue hacia la parte imaginaria. Lo que se ve es que \textbf{el
pulso vuelve dado vuelta}.

\end{minipage}%
\end{tcolorbox}

\subsection{El extremo libre}\label{sec-hip-semi-libre}

Si en vez de atado el extremo es un anillo que desliza, la condición es
\(u_x(0,t) = 0\) (Tabla~\ref{tbl-hip-bordes}). La receta es la misma
cambiando la imagen: ahora el gemelo \textbf{no} se invierte.

\begin{proposition}[Cuerda semi-infinita con extremo
libre]\protect\hypertarget{prp-hip-semi-libre}{}\label{prp-hip-semi-libre}

La solución de \ref{eq-hip-semi} con \(u_x(0,t)=0\) en lugar de
\(u(0,t)=0\) es la de d'Alembert con los datos extendidos en forma
\textbf{par}. Para \(0\le x<at\),

\begin{equation}\protect\phantomsection\label{eq-hip-semi-libre}{
u(x,t) = \frac{\varphi(x+at)+\varphi(at-x)}{2} + \frac{1}{2a}\left[\int_{0}^{x+at}\psi(s)\,ds + \int_{0}^{at-x}\psi(s)\,ds\right],
}\end{equation}

y para \(x\ge at\) es la fórmula de d'Alembert sin cambios.

\end{proposition}

Que funcione se ve derivando: si los datos son pares, la solución es par
en \(x\) ---por la unicidad, porque \(u(-x,t)\) también es solución con
los mismos datos---, y la derivada de una función par se anula en el
origen. El Ejercicio~\ref{exr-hip-libre} pide hacer la cuenta.

\subsection{Cualquier extremo: el coeficiente de
reflexión}\label{sec-hip-coef-reflexion}

Entre el extremo fijo y el libre hay toda una familia: el amortiguador
del ejercicio del sección~\ref{sec-hip-unicidad}. Pongámoslo en el
extremo izquierdo de una cuerda que ocupa \(x>0\). Ahí la cuerda tira
del extremo hacia arriba con la componente vertical de la tensión,
\(T\,u_x\), y el amortiguador frena con \(-\beta\,u_t\); como el extremo
no tiene masa, las dos fuerzas se equilibran: \[
T\,u_x(0,t) = \beta\,u_t(0,t).
\] (En el extremo derecho el signo de la tensión se invierte, y queda el
\(T u_x = -\beta u_t\) del ejercicio.) Para ver qué refleja conviene
dejar las extensiones y hacer la cuenta directa.

Una onda \(f(x+at)\) llega al extremo \(x=0\) viajando hacia la
izquierda. Buscamos la onda reflejada, que viaja hacia la derecha:
\(u = f(x+at) + g(x-at)\). La condición de borde en \(x=0\) da \[
T\,\big[f'(at)+g'(-at)\big] = \beta\,a\,\big[f'(at) - g'(-at)\big]
\quad\Longrightarrow\quad
g'(-at) = \frac{\beta a - T}{\beta a + T}\;f'(at).
\] Integrando en \(t\) ---y fijando la constante con que \(g\) es cero
antes de que llegue nada---, la onda reflejada es \(g(s) = R\,f(-s)\):
\textbf{la imagen especular de la que llegó, multiplicada por}

\begin{equation}\protect\phantomsection\label{eq-hip-reflexion}{
R = \frac{Z-\beta}{Z+\beta}, \qquad Z = \frac{T}{a} = \sqrt{T\rho}.
}\end{equation}

\begin{definition}[Impedancia y coeficiente de
reflexión]\protect\hypertarget{def-hip-impedancia}{}\label{def-hip-impedancia}

La cantidad \(Z = \sqrt{T\rho}\) se llama \textbf{impedancia} de la
cuerda. El número \(R\) de \ref{eq-hip-reflexion} es el
\textbf{coeficiente de reflexión}: la amplitud de la onda reflejada
dividida por la de la incidente.

\end{definition}

Los casos que ya conocíamos salen solos, y aparece uno nuevo:

\begin{longtable}[]{@{}
  >{\raggedright\arraybackslash}p{(\linewidth - 8\tabcolsep) * \real{0.2200}}
  >{\raggedright\arraybackslash}p{(\linewidth - 8\tabcolsep) * \real{0.1400}}
  >{\raggedright\arraybackslash}p{(\linewidth - 8\tabcolsep) * \real{0.1200}}
  >{\raggedright\arraybackslash}p{(\linewidth - 8\tabcolsep) * \real{0.2200}}
  >{\raggedright\arraybackslash}p{(\linewidth - 8\tabcolsep) * \real{0.3000}}@{}}
\caption{El coeficiente de reflexión para cada tipo de
extremo.}\label{tbl-hip-reflexion}\tabularnewline
\toprule\noalign{}
\begin{minipage}[b]{\linewidth}\raggedright
Extremo
\end{minipage} & \begin{minipage}[b]{\linewidth}\raggedright
\(\beta\)
\end{minipage} & \begin{minipage}[b]{\linewidth}\raggedright
\(R\)
\end{minipage} & \begin{minipage}[b]{\linewidth}\raggedright
Qué vuelve
\end{minipage} & \begin{minipage}[b]{\linewidth}\raggedright
Energía reflejada, \(R^2\)
\end{minipage} \\
\midrule\noalign{}
\endfirsthead
\toprule\noalign{}
\begin{minipage}[b]{\linewidth}\raggedright
Extremo
\end{minipage} & \begin{minipage}[b]{\linewidth}\raggedright
\(\beta\)
\end{minipage} & \begin{minipage}[b]{\linewidth}\raggedright
\(R\)
\end{minipage} & \begin{minipage}[b]{\linewidth}\raggedright
Qué vuelve
\end{minipage} & \begin{minipage}[b]{\linewidth}\raggedright
Energía reflejada, \(R^2\)
\end{minipage} \\
\midrule\noalign{}
\endhead
\bottomrule\noalign{}
\endlastfoot
libre & \(0\) & \(+1\) & el pulso derecho & toda \\
fijo & \(\infty\) & \(-1\) & el pulso invertido & toda \\
\textbf{adaptado} & \(Z\) & \(0\) & \textbf{nada} & ninguna \\
amortiguador cualquiera & \(0<\beta<\infty\) & entre \(-1\) y \(1\) &
una copia achicada & \(R^2\); el resto, \(1-R^2\), se va en calor \\
\end{longtable}

La tercera fila es de ingeniería pura. Un amortiguador con \(\beta = Z\)
hace, sobre el extremo, \textbf{exactamente la misma fuerza que haría el
resto de una cuerda infinita}: la onda llega y no se entera de que la
cuerda terminó. Es lo que se hace en los cables: una \textbf{resistencia
de terminación} igual a la impedancia del cable, para que las señales no
reboten en la punta y se superpongan con las siguientes.

\emph{(En la versión web hay acá un gráfico interactivo: un pulso que
viaja hacia el extremo de una cuerda y, del otro lado del extremo, sobre
una cuerda imaginaria, su gemelo especular multiplicado por el
coeficiente de reflexión \(R\). Un slider lleva el extremo de libre
(\(R=1\)) a fijo (\(R=-1\)), pasando por el amortiguador adaptado
(\(R=0\)), y otro mueve el tiempo.)}

Un pulso hacia arriba llega al extremo \textbf{fijo} de una cuerda. ¿Qué
vuelve?

\begin{itemize}
\tightlist
\item
  {Un pulso hacia abajo, del mismo tamaño}
\item
  {Un pulso hacia arriba, del mismo tamaño}
\item
  {Nada: el extremo fijo absorbe el pulso}
\item
  {Un pulso hacia arriba de la mitad del tamaño}
\end{itemize}

En el extremo fijo \(R = -1\): la onda reflejada es la imagen especular
de la incidente con el signo cambiado. La razón física es la tercera ley
de Newton: cuando el pulso llega, tira del soporte hacia arriba; el
soporte, que no se mueve, tira de la cuerda hacia abajo, y ese tirón
genera el pulso invertido. Con un extremo libre no hay soporte que tire,
el anillo sube de más ---el doble de la amplitud del pulso, en el
instante de la llegada--- y el pulso vuelve derecho.

La tercera opción es la del amortiguador adaptado. La cuarta no
corresponde a un extremo fijo, pero sí a otro extremo: un pulso derecho
de la mitad del tamaño vuelve de un amortiguador con \(\beta = Z/3\),
para el que \(R = 1/2\).

\subsection{La cuerda finita: rebotes infinitos}\label{sec-hip-finita}

Con dos extremos fijos, en \(x=0\) y en \(x=l\), el truco se aplica dos
veces. La extensión impar respecto de \(0\) arregla el extremo
izquierdo; pero entonces hace falta que también sea impar respecto de
\(l\), y eso obliga a reflejar la imagen, y la imagen de la imagen, sin
fin. Lo que queda es una función definida en toda la recta que es impar
y \textbf{periódica de período \(2l\)}.

\begin{theorem}[La cuerda fija por
reflexiones]\protect\hypertarget{thm-hip-finita}{}\label{thm-hip-finita}

Sean \(\tilde\varphi\) y \(\tilde\psi\) las extensiones impares y
\(2l\)-periódicas de \(\varphi\) y \(\psi\). La solución de la cuerda
con extremos fijos,
\[u_{tt} = a^2u_{xx},\quad u(0,t) = u(l,t) = 0,\quad u(x,0) = \varphi(x),\quad u_t(x,0) = \psi(x),\]
es la fórmula de d'Alembert aplicada a \(\tilde\varphi\) y
\(\tilde\psi\):

\begin{equation}\protect\phantomsection\label{eq-hip-finita}{
u(x,t) = \frac{\tilde\varphi(x+at)+\tilde\varphi(x-at)}{2} + \frac{1}{2a}\int_{x-at}^{x+at}\tilde\psi(s)\,ds .
}\end{equation}

En particular, el movimiento es \textbf{periódico en el tiempo} con
período \(2l/a\).

\end{theorem}

\begin{proof}
Una función impar respecto de \(0\) y \(2l\)-periódica es también impar
respecto de \(l\):
\(\tilde\varphi(l+y) = \tilde\varphi(y-l) = -\tilde\varphi(l-y)\),
usando primero la periodicidad y después la imparidad. Con datos impares
respecto de \(0\) y de \(l\), la solución de d'Alembert se anula en los
dos puntos. Es la solución por unicidad
(Teorema~\ref{thm-hip-unicidad}).

Para la periodicidad: si \(t\) aumenta en \(2l/a\), los argumentos
\(x\pm at\) aumentan o disminuyen en \(2l\), y \(\tilde\varphi\) no
cambia. La integral de \(\tilde\psi\) sobre un período es cero porque
\(\tilde\psi\) es impar, así que el segundo término tampoco cambia.
\end{proof}

Ahora juntemos esto con el sección~\ref{sec-cuerda-separacion}. Allá la
misma cuerda se resolvió con modos: \[
u(x,t) = \sum_{n\ge1}\Big(a_n\cos\omega_nt + b_n\operatorname{sen}\omega_nt\Big)\operatorname{sen}\frac{n\pi x}{l},
\qquad \omega_n = \frac{n\pi a}{l}.
\] ¿Son la misma solución? Tienen que serlo, por la unicidad. Pero se
puede ver directamente. La extensión impar y \(2l\)-periódica de
\(\varphi\) tiene una serie de Fourier \textbf{solo de senos}, y es
exactamente la serie de senos de \(\varphi\) en \([0,l]\)
(Capítulo~\ref{sec-fourier}): \[
\tilde\varphi(x) = \sum_{n\ge1} a_n \operatorname{sen}\frac{n\pi x}{l}.
\] Poniéndola en el primer término de \ref{eq-hip-finita} y usando
\(\operatorname{sen}A + \operatorname{sen}B = 2\operatorname{sen}\tfrac{A+B}{2}\cos\tfrac{A-B}{2}\),
\[
\frac{\tilde\varphi(x+at)+\tilde\varphi(x-at)}{2} = \sum_{n\ge1} a_n\operatorname{sen}\frac{n\pi x}{l}\cos\frac{n\pi at}{l},
\] que es la parte en coseno de la serie de modos. Con \(\psi\) pasa lo
mismo, y sale la parte en seno.

\begin{tcolorbox}[enhanced jigsaw, arc=.35mm, bottomrule=.15mm, breakable, colback=white, colframe=quarto-callout-important-color-frame, left=2mm, leftrule=.75mm, opacityback=0, rightrule=.15mm, toprule=.15mm]
\begin{minipage}[t]{5.5mm}
\textcolor{quarto-callout-important-color}{\faExclamation}
\end{minipage}%
\begin{minipage}[t]{\textwidth - 5.5mm}

\vspace{-3mm}\textbf{Dos imágenes, una cuerda}\vspace{3mm}

\textbf{d'Alembert} dice: la cuerda es una onda que viaja, rebota en un
extremo dada vuelta, viaja, rebota en el otro, y a los \(2l/a\) segundos
está donde empezó.

\textbf{Los modos} dicen: la cuerda es una suma de formas fijas
\(\operatorname{sen}(n\pi x/l)\) que laten, cada una a su frecuencia
\(\omega_n = n\pi a/l\), todas múltiplos de la fundamental.

Son la misma cuenta. Cada modo es la suma de dos ondas que viajan en
sentidos opuestos (Ejercicio~\ref{exr-hip-estacionaria}), y la suma de
todos los modos rearma los pulsos que rebotan. Las frecuencias son
múltiplos enteros de \(\pi a/l\) \textbf{porque} el viaje de ida y
vuelta dura \(2l/a\): todo lo que tiene período \(2l/a\) se escribe con
esas frecuencias y solo con ellas. Por eso una cuerda tiene una nota
definida, con armónicos afinados (sección~\ref{sec-interpretacion}).

Qué imagen conviene depende de la pregunta. Para saber \textbf{cuándo}
llega algo a un punto, d'Alembert. Para saber \textbf{cómo suena}, los
modos.

\end{minipage}%
\end{tcolorbox}

\emph{(En la versión web hay acá un gráfico interactivo en dos paneles
sobre la misma cuerda pulsada. Arriba, la construcción de d'Alembert: la
extensión impar y periódica partida en dos mitades que viajan, y la
cuerda real como la franja entre \(0\) y \(1\). Abajo, la suma de los
primeros \(N\) modos encima de la solución exacta. Se pueden mover el
tiempo, la cantidad de modos y el punto donde se pulsa.)}

La cuerda aguda de la guitarra del ejercicio del
sección~\ref{sec-hip-cuerda}, con \(a \approx 418\ \text{m/s}\), tiene
un largo vibrante de \(l = 0{,}648\ \text{m}\). ¿Cuál es su frecuencia
fundamental, en Hz?

El período es el tiempo de ida y vuelta, \(2l/a\), así que la frecuencia
fundamental es
\[f_1 = \frac{a}{2l} = \frac{418}{1{,}296} \approx 322{,}5\ \text{Hz}.\]
Es la misma fórmula que \(\omega_1 = \pi a/l\) dividida por \(2\pi\). La
nota de esa cuerda es un mi de \(329{,}6\) Hz: con \(70\) N de tensión
quedó un poco baja. Para afinarla hay que subir la tensión en un factor
\((329{,}6/322{,}5)^2 \approx 1{,}045\), o sea a unos \(73\) N.

\subsection{Un extremo que se mueve}\label{sec-hip-extremo-movil}

Queda el caso 2 del sección~\ref{sec-hip-casos-limite} con toda su
generalidad: el extremo no está fijo sino que alguien lo mueve,
\(u(0,t) = \mu(t)\). Es una persona sacudiendo la punta de una soga
larga.

\begin{proposition}[Cuerda semi-infinita con el extremo
movido]\protect\hypertarget{prp-hip-extremo-movil}{}\label{prp-hip-extremo-movil}

Con \(\varphi = \psi = 0\) ---la cuerda en reposo--- y \(\mu(0) = 0\),
la solución de
\[u_{tt} = a^2u_{xx}\ (x>0),\qquad u(0,t) = \mu(t),\qquad u(x,0) = u_t(x,0) = 0\]
es

\begin{equation}\protect\phantomsection\label{eq-hip-extremo-movil}{
u(x,t) = \begin{cases} \mu\!\left(t - \dfrac{x}{a}\right), & x < at,\\[2mm] 0, & x\ge at.\end{cases}
}\end{equation}

Con datos iniciales no nulos, se suma esta función a la solución de
\ref{eq-hip-semi-fijo}.

\end{proposition}

\begin{proof}
La función es de la forma \(g(x-at)\), con \(g(s) = \mu(-s/a)\) para
\(s<0\) y \(g = 0\) para \(s\ge0\): cumple la ecuación. En \(x=0\) vale
\(\mu(t)\). En \(t=0\), todo \(x>0\) cumple \(x\ge at\) y \(u=0\); y lo
mismo para \(u_t\). La suma con la solución de datos \((\varphi,\psi)\)
y extremo fijo cumple todo, por linealidad.
\end{proof}

La lectura es la más directa posible: \textbf{lo que hace el extremo en
el instante \(t\) llega al punto \(x\) con una demora \(x/a\)}. El
extremo es una fuente que emite hacia la derecha, y la cuerda reproduce
su historia, retrasada.

\begin{exercise}[]\protect\hypertarget{exr-hip-libre}{}\label{exr-hip-libre}

Deducir \ref{eq-hip-semi-libre} extendiendo los datos en forma par, y
verificar directamente que cumple \(u_x(0,t) = 0\).

\end{exercise}

\begin{tcolorbox}[enhanced jigsaw, arc=.35mm, bottomrule=.15mm, bottomtitle=1mm, breakable, colback=white, colbacktitle=quarto-callout-tip-color!10!white, colframe=quarto-callout-tip-color-frame, coltitle=black, left=2mm, leftrule=.75mm, opacityback=0, opacitybacktitle=0.6, rightrule=.15mm, title=\textcolor{quarto-callout-tip-color}{\faLightbulb}\hspace{0.5em}{Solución}, titlerule=0mm, toprule=.15mm, toptitle=1mm]

Con la extensión par, para \(0\le x<at\) el punto \(x-at\) es negativo y
\(\varphi_{\text{p}}(x-at) = \varphi(at-x)\). En la integral, partimos
en \([x-at,\,0]\) y \([0,\,x+at]\), y en el primer tramo cambiamos
\(s\to-s\): \[
\int_{x-at}^{0}\psi_{\text{p}}(s)\,ds = \int_{0}^{at-x}\psi(s)\,ds .
\] Eso da \ref{eq-hip-semi-libre}.

\textbf{Verificación.} Derivando en \(x\), con la regla de Leibniz, \[
u_x = \frac{\varphi'(x+at) - \varphi'(at-x)}{2} + \frac{1}{2a}\big[\psi(x+at) - \psi(at-x)\big],
\] y en \(x=0\) los dos corchetes se anulan: \(u_x(0,t) = 0\).

\textbf{Una condición que no hay que olvidar.} Para que la solución sea
\(C^1\) sobre la característica \(x = at\), donde se pegan las dos
fórmulas, hace falta que los datos sean compatibles con el borde:
\(\varphi'(0) = 0\). Si no, la extensión par tiene una esquina en el
origen y la solución tiene una esquina que viaja por \(x = at\). Es la
misma idea del Ejercicio~\ref{exr-hip-verificar}.

\end{tcolorbox}

\begin{exercise}[]\protect\hypertarget{exr-hip-soga}{}\label{exr-hip-soga}

Una soga muy larga está en reposo. En \(t=0\) alguien empieza a sacudir
la punta, \(u(0,t) = A\operatorname{sen}(\Omega t)\). Escribir
\(u(x,t)\) y describir qué se ve. ¿La solución es de clase \(C^2\)?

\end{exercise}

\begin{tcolorbox}[enhanced jigsaw, arc=.35mm, bottomrule=.15mm, bottomtitle=1mm, breakable, colback=white, colbacktitle=quarto-callout-tip-color!10!white, colframe=quarto-callout-tip-color-frame, coltitle=black, left=2mm, leftrule=.75mm, opacityback=0, opacitybacktitle=0.6, rightrule=.15mm, title=\textcolor{quarto-callout-tip-color}{\faLightbulb}\hspace{0.5em}{Solución}, titlerule=0mm, toprule=.15mm, toptitle=1mm]

Por \ref{eq-hip-extremo-movil}, \[
u(x,t) = \begin{cases} A\operatorname{sen}\big(\Omega t - kx\big), & x<at,\\ 0, & x\ge at,\end{cases}
\qquad k = \frac{\Omega}{a}.
\] Detrás del frente \(x = at\) hay una onda sinusoidal que viaja hacia
la derecha, de longitud de onda \(\lambda = 2\pi/k = 2\pi a/\Omega\);
delante del frente, la soga todavía está quieta.

\textbf{La regularidad.} En el frente, \(u\) es continua (vale \(0\) de
los dos lados) pero \(u_x\) salta: de un lado es \(-Ak\cos(0) = -Ak\),
del otro \(0\). La solución tiene una \textbf{esquina} que viaja por la
característica \(x = at\). No es \(C^2\), ni siquiera \(C^1\).

La razón es una incompatibilidad entre los datos: en el origen, en
\(t=0\), la condición de borde dice que el extremo arranca con velocidad
\(\mu'(0) = A\Omega\), y la condición inicial dice que arranca con
velocidad \(\psi(0) = 0\). Las dos cosas no pueden ser ciertas, y el
conflicto viaja por la característica. Si la mano arrancara despacio
---con \(\mu'(0) = 0\)--- la esquina desaparecería. Es la misma lección
de siempre: \textbf{las singularidades viajan por las características},
y a veces las fabrican los datos.

\end{tcolorbox}

\begin{exercise}[]\protect\hypertarget{exr-hip-paralelogramo-borde}{}\label{exr-hip-paralelogramo-borde}

Deducir la segunda línea de \ref{eq-hip-semi-fijo} \textbf{sin usar la
extensión impar}, aplicando la regla del paralelogramo
(Proposición~\ref{prp-hip-paralelogramo}) a un paralelogramo con un
vértice en el punto \((x,t)\), con \(x<at\), y otro sobre el borde
\(x=0\).

\end{exercise}

\begin{tcolorbox}[enhanced jigsaw, arc=.35mm, bottomrule=.15mm, bottomtitle=1mm, breakable, colback=white, colbacktitle=quarto-callout-tip-color!10!white, colframe=quarto-callout-tip-color-frame, coltitle=black, left=2mm, leftrule=.75mm, opacityback=0, opacitybacktitle=0.6, rightrule=.15mm, title=\textcolor{quarto-callout-tip-color}{\faLightbulb}\hspace{0.5em}{Solución}, titlerule=0mm, toprule=.15mm, toptitle=1mm]

Sea \(C = (x,t)\) con \(0<x<at\). Armamos el paralelogramo bajando desde
\(C\) por las dos familias de características:

\begin{itemize}
\tightlist
\item
  Por la recta \(x' - at' = x - at\), que baja hacia la izquierda, se
  llega al borde \(x'=0\) en \(D = \big(0,\ t - x/a\big)\). Como
  \(x<at\), ese punto está en \(t'>0\).
\item
  Desde \(D\), por la otra familia, \(x' + at' = at - x\), se baja hasta
  el eje: \(A = (at - x,\ 0)\).
\item
  Desde \(C\), por \(x' + at' = x + at\), y desde \(A\), por
  \(x' - at' = at - x\), las dos rectas se cortan en
  \(B = \big(at,\ x/a\big)\).
\end{itemize}

Los cuatro vértices \(A\), \(B\), \(C\), \(D\) tienen lados sobre
características, y la regla dice \[
u(C) = u(B) + u(D) - u(A) = u\big(at,\ x/a\big) + 0 - \varphi(at-x),
\] usando \(u(D) = 0\) por el borde fijo. El punto \(B\) tiene
\(x' = at \ge a\,t' = x\), así que está en la región donde vale
d'Alembert sin cambios: \[
u\big(at,\,x/a\big) = \frac{\varphi(at + x) + \varphi(at - x)}{2} + \frac{1}{2a}\int_{at-x}^{at+x}\psi .
\] Sumando \(-\varphi(at-x)\) queda \[
u(x,t) = \frac{\varphi(x+at) - \varphi(at-x)}{2} + \frac{1}{2a}\int_{at-x}^{x+at}\psi(s)\,ds ,
\] que es \ref{eq-hip-semi-fijo}. Sin inventar ninguna cuerda
imaginaria: solo con la geometría de las características.

\end{tcolorbox}

\section{Fuentes: ondas que alguien fabrica}\label{sec-hip-fuentes}

Hasta acá las ondas se movían solas, con lo que tenían en \(t=0\). Ahora
hay algo que las sigue empujando: una fuerza \(F(x,t)\) sobre la cuerda,
o un parlante en el aire, o una corriente en una antena. La ecuación es
la no homogénea,

\[
u_{tt} = a^2u_{xx} + f(x,t), \qquad f = F/\rho .
\]

Como la ecuación es lineal, la respuesta es \textbf{lo que harían los
datos solos más lo que hace la fuente sola}, y para lo segundo hay una
sola idea, con tres nombres según dónde se la mire: principio de
Duhamel, función de Green, respuesta al impulso. La idea es la del
Capítulo~\ref{sec-green}: una fuerza que dura un rato es una sucesión de
golpes cortos, y la respuesta a cada golpe ya la conocemos.

\subsection{d'Alembert con fuente}\label{sec-hip-duhamel}

En el Ejemplo~\ref{exm-onda-fuente} quedó prometida esta fórmula. Acá
está.

\begin{theorem}[Fórmula de d'Alembert con
fuente]\protect\hypertarget{thm-hip-dalembert-fuente}{}\label{thm-hip-dalembert-fuente}

Sean \(\varphi\in C^2(\mathbb R)\), \(\psi\in C^1(\mathbb R)\) y \(f\)
continua con \(f_x\) continua. La única solución \(C^2\) de
\[u_{tt} = a^2u_{xx} + f(x,t),\quad x\in\mathbb R,\ t>0;\qquad u(x,0)=\varphi(x),\quad u_t(x,0)=\psi(x)\]
es

\begin{equation}\protect\phantomsection\label{eq-hip-dalembert-fuente}{
u(x,t) = \frac{\varphi(x+at)+\varphi(x-at)}{2} + \frac{1}{2a}\int_{x-at}^{x+at}\psi(s)\,ds
+ \frac{1}{2a}\iint_{\Delta(x,t)} f(s,\tau)\,ds\,d\tau ,
}\end{equation}

donde \(\Delta(x,t)\) es el \textbf{triángulo característico} del punto:
el que tiene vértice en \((x,t)\) y base \([x-at,\,x+at]\) sobre el eje
\(\tau=0\). Escrita como integral iterada, \[
\iint_{\Delta(x,t)} f\,ds\,d\tau = \int_0^t\int_{x-a(t-\tau)}^{x+a(t-\tau)} f(s,\tau)\,ds\,d\tau .
\]

\end{theorem}

\begin{proof}
Hay una demostración que da la fórmula entera de una sola vez, con el
teorema de Green en el plano. Sea \(u\) una solución \(C^2\) e
integremos la ecuación sobre el triángulo: \[
\iint_{\Delta}\big(u_{tt} - a^2u_{xx}\big)\,dx\,dt = \iint_\Delta f\,dx\,dt .
\] Por el teorema de Green, con \(P = -u_t\) y \(Q = -a^2u_x\), el
primer miembro es la circulación
\(\oint_{\partial\Delta}\big(-u_t\,dx - a^2u_x\,dt\big)\) en sentido
antihorario. Recorramos los tres lados:

\begin{itemize}
\tightlist
\item
  \textbf{La base}, de \((x-at,0)\) a \((x+at,0)\): \(dt = 0\), y queda
  \(-\int u_t\,dx = -\int_{x-at}^{x+at}\psi\).
\item
  \textbf{El lado derecho}, de \((x+at,0)\) al vértice: es una
  característica, \(dx = -a\,dt\). El integrando se vuelve
  \(a\,(u_x\,dx + u_t\,dt) = a\,du\), y su integral es
  \(a\big[u(x,t) - \varphi(x+at)\big]\).
\item
  \textbf{El lado izquierdo}, del vértice a \((x-at,0)\): ahora
  \(dx = a\,dt\), el integrando es \(-a\,du\), y la integral es
  \(-a\big[\varphi(x-at) - u(x,t)\big]\).
\end{itemize}

Sumando,
\(2a\,u(x,t) - a\varphi(x+at) - a\varphi(x-at) - \int\psi = \iint_\Delta f\),
que es \ref{eq-hip-dalembert-fuente} despejando \(u(x,t)\).

Esto prueba que \textbf{si hay} solución \(C^2\), es esa: la unicidad.
Que la fórmula efectivamente resuelve el problema es el
Ejercicio~\ref{exr-hip-duhamel}.
\end{proof}

La cuenta de los lados del triángulo es la misma idea que el método de
las características del sección~\ref{sec-caracteristicas}: sobre una
característica, \(u_x\,dx + u_t\,dt\) es la derivada de \(u\) a lo largo
del camino. El triángulo característico ya había aparecido como dominio
de dependencia; ahora se ve que es también el \textbf{dominio de
dependencia de la fuente}. El valor de \(u\) en \((x,t)\) depende de lo
que la fuerza hizo \textbf{dentro del triángulo}, y de nada más: una
fuerza aplicada lejos, o hace poco, todavía no tuvo tiempo de llegar.

\begin{tcolorbox}[enhanced jigsaw, arc=.35mm, bottomrule=.15mm, breakable, colback=white, colframe=quarto-callout-note-color-frame, left=2mm, leftrule=.75mm, opacityback=0, rightrule=.15mm, toprule=.15mm]
\begin{minipage}[t]{5.5mm}
\textcolor{quarto-callout-note-color}{\faInfo}
\end{minipage}%
\begin{minipage}[t]{\textwidth - 5.5mm}

\vspace{-3mm}\textbf{El principio de Duhamel, en palabras}\vspace{3mm}

La integral doble se lee así. Entre \(\tau\) y \(\tau+d\tau\) la fuerza
le da a la cuerda un impulso, que cambia su velocidad en
\(f(s,\tau)\,d\tau\). Eso es una velocidad inicial, y por d'Alembert, al
tiempo \(t\) ---o sea \(t-\tau\) después--- produce el desplazamiento
\[\frac{1}{2a}\int_{x-a(t-\tau)}^{x+a(t-\tau)} f(s,\tau)\,d\tau\,ds .\]
Sumando sobre todos los \(\tau\) entre \(0\) y \(t\) se obtiene la
integral doble. \textbf{Una fuerza continua es una lluvia de golpes, y
la respuesta es la suma de las respuestas a cada golpe.}

\end{minipage}%
\end{tcolorbox}

\subsection{La función de Green de la cuerda
infinita}\label{sec-hip-green-1d}

El caso extremo de la lluvia de golpes es \textbf{un solo golpe},
concentrado en un punto y en un instante:
\(f(x,t) = \delta(x-s)\,\delta(t-\tau)\). Por
\ref{eq-hip-dalembert-fuente}, la integral de la delta doble sobre el
triángulo vale \(1\) si el punto \((s,\tau)\) cae adentro y \(0\) si no.

\begin{definition}[Función de Green de la ecuación de onda en una
dimensión]\protect\hypertarget{def-hip-green-1d}{}\label{def-hip-green-1d}

La solución de \(u_{tt} - a^2u_{xx} = \delta(x-s)\,\delta(t-\tau)\) con
la cuerda en reposo antes de \(\tau\) es

\begin{equation}\protect\phantomsection\label{eq-hip-green-1d}{
G(x,t;s,\tau) = \frac{1}{2a}\,H\!\big(a(t-\tau) - |x-s|\big),
}\end{equation}

donde \(H\) es la función de Heaviside. Con ella, la contribución de la
fuente en \ref{eq-hip-dalembert-fuente} es \[
\int_0^\infty\!\!\int_{-\infty}^{\infty} G(x,t;s,\tau)\,f(s,\tau)\,ds\,d\tau .
\]

\end{definition}

La respuesta a un golpe puntual es una \textbf{meseta} de altura
\(1/(2a)\) que se ensancha a velocidad \(a\) para los dos lados. Es la
cuerda golpeada de la GIF del Ejercicio~\ref{exr-dalembert}, en el
límite de un martillo infinitamente angosto. Y la forma es la del
\textbf{cono de luz}: \(G\) vale una constante adentro del cono que sale
de \((s,\tau)\) y cero afuera. En el sección~\ref{sec-hip-maxwell} vamos
a ver que en tres dimensiones la función de Green ya no llena el cono:
vive \textbf{solo sobre su superficie}, y esa diferencia es la que hace
que el sonido se oiga como se oye.

Una cuerda muy larga, con \(\rho = 0{,}01\ \text{kg/m}\) y
\(a = 100\ \text{m/s}\), recibe un golpe seco de impulso
\(I = 0{,}01\ \text{N·s}\) en un punto. ¿Cuánto sube la meseta? Respondé
en milímetros.

Un impulso \(I\) concentrado en un punto es una fuerza
\(F = I\,\delta(x-s)\,\delta(t)\), así que
\(f = F/\rho = (I/\rho)\,\delta(x-s)\,\delta(t)\). Por
\ref{eq-hip-green-1d}, la respuesta es \(I/\rho\) veces la función de
Green: una meseta de altura \[
\frac{I}{2a\rho} = \frac{0{,}01}{2\cdot100\cdot0{,}01} = 0{,}005\ \text{m} = 5\ \text{mm}.
\] Lo mismo sale de pensarlo como velocidad inicial: el golpe le da a un
tramo de cuerda un momento \(I\), es decir
\(\psi = (I/\rho)\,\delta(x-s)\), y la integral de d'Alembert da
\(\tfrac{1}{2a}\cdot I/\rho\). Es el problema 9.5 de la cátedra en una
cuerda infinita.

\subsection{La cuerda finita forzada}\label{sec-hip-forzada}

Esta es la sección 9.3.2 de González. La cuerda tiene extremos fijos,
está sometida a una fuerza \(F(x,t)\), y arranca con datos \(\varphi\) y
\(\psi\):

\begin{equation}\protect\phantomsection\label{eq-hip-forzada}{
u_{tt} = a^2u_{xx} + f(x,t),\qquad u(0,t) = u(l,t) = 0,\qquad u(x,0) = \varphi(x),\quad u_t(x,0) = \psi(x).
}\end{equation}

Por la linealidad partimos el problema en dos, \(u = v + w\):

\begin{itemize}
\tightlist
\item
  \(w\) son las \textbf{oscilaciones libres}: la ecuación sin fuente,
  con los datos \(\varphi\) y \(\psi\). Ya sabemos resolverla, por modos
  (sección~\ref{sec-cuerda-separacion}) o por reflexiones
  (Teorema~\ref{thm-hip-finita}).
\item
  \(v\) son las \textbf{oscilaciones forzadas}: la ecuación con la
  fuente, pero con la cuerda arrancando en reposo,
  \(v(x,0) = v_t(x,0) = 0\).
\end{itemize}

\begin{tcolorbox}[enhanced jigsaw, arc=.35mm, bottomrule=.15mm, breakable, colback=white, colframe=quarto-callout-warning-color-frame, left=2mm, leftrule=.75mm, opacityback=0, rightrule=.15mm, toprule=.15mm]
\begin{minipage}[t]{5.5mm}
\textcolor{quarto-callout-warning-color}{\faExclamationTriangle}
\end{minipage}%
\begin{minipage}[t]{\textwidth - 5.5mm}

\vspace{-3mm}\textbf{Ojo con el libro de la cátedra (sección 9.3.2)}\vspace{3mm}

En (9.17) y (9.18) las condiciones están escritas para \(u\), pero son
las de \(v\) y \(w\) respectivamente. Y en (9.24)--(9.26) la variable de
integración está mal puesta: dice \(f_n(t)\) y \(f(z,t)\) adentro de
\(\int_0^t\cdots d\tau\), donde tiene que decir \(f_n(\tau)\) y
\(f(z,\tau)\). Con la \(t\) la integral no tiene sentido: la variable
que se integra no aparecería en el integrando. La versión correcta es
\ref{eq-hip-forzada-sol}.

\end{minipage}%
\end{tcolorbox}

Para \(v\) la idea es la de siempre en este capítulo: escribir todo en
la base de modos de la cuerda, \(\operatorname{sen}(n\pi x/l)\), que
cumplen las condiciones de contorno.

\[
v(x,t) = \sum_{n\ge1} T_n(t)\operatorname{sen}\frac{n\pi x}{l}, \qquad
f(x,t) = \sum_{n\ge1} f_n(t)\operatorname{sen}\frac{n\pi x}{l}, \quad
f_n(t) = \frac2l\int_0^l f(z,t)\operatorname{sen}\frac{n\pi z}{l}\,dz .
\]

Reemplazando en la ecuación y comparando coeficiente a coeficiente, cada
modo cumple \textbf{la ecuación de un oscilador forzado}:

\begin{equation}\protect\phantomsection\label{eq-hip-oscilador}{
T_n'' + \omega_n^2\,T_n = f_n(t), \qquad T_n(0) = T_n'(0) = 0, \qquad \omega_n = \frac{n\pi a}{l}.
}\end{equation}

Es la ecuación que se resolvió con la función de Green en el
Capítulo~\ref{sec-green}:

\begin{theorem}[Oscilaciones forzadas de la cuerda
fija]\protect\hypertarget{thm-hip-forzada}{}\label{thm-hip-forzada}

La solución de \ref{eq-hip-forzada} con \(\varphi = \psi = 0\) es

\begin{equation}\protect\phantomsection\label{eq-hip-forzada-sol}{
v(x,t) = \sum_{n\ge1}\left[\frac{1}{\omega_n}\int_0^t f_n(\tau)\operatorname{sen}\omega_n(t-\tau)\,d\tau\right]\operatorname{sen}\frac{n\pi x}{l}
= \sum_{n\ge1}\left[\frac{2}{l\,\omega_n}\int_0^t\!\!\int_0^l f(z,\tau)\operatorname{sen}\frac{n\pi z}{l}\operatorname{sen}\omega_n(t-\tau)\,dz\,d\tau\right]\operatorname{sen}\frac{n\pi x}{l}.
}\end{equation}

La solución general de \ref{eq-hip-forzada} es \(u = v + w\).

\end{theorem}

La lectura es la que ya teníamos: \textbf{la cuerda es un banco de
osciladores independientes}, uno por modo, cada uno con su frecuencia
\(\omega_n\). La fuerza se reparte entre ellos según cuánto se parezca a
cada modo ---eso es \(f_n\)---, y cada oscilador responde por su cuenta.

\begin{exercise}[]\protect\hypertarget{exr-hip-comparar-transformada}{}\label{exr-hip-comparar-transformada}

González propone comparar \ref{eq-hip-forzada-sol} con lo obtenido por
transformadas. Hacerlo con el ejemplo del
sección~\ref{sec-ejemplo-transformada}: una fuerza constante \(F_0(x)\)
que se aplica de golpe en \(t=0\) a la cuerda en reposo. Mostrar que
cada modo llega, en algún instante, al \textbf{doble} de su deflexión
estática.

\end{exercise}

\begin{tcolorbox}[enhanced jigsaw, arc=.35mm, bottomrule=.15mm, bottomtitle=1mm, breakable, colback=white, colbacktitle=quarto-callout-tip-color!10!white, colframe=quarto-callout-tip-color-frame, coltitle=black, left=2mm, leftrule=.75mm, opacityback=0, opacitybacktitle=0.6, rightrule=.15mm, title=\textcolor{quarto-callout-tip-color}{\faLightbulb}\hspace{0.5em}{Solución}, titlerule=0mm, toprule=.15mm, toptitle=1mm]

Si la fuerza no depende del tiempo para \(t>0\), los coeficientes
\(f_n\) son constantes y \ref{eq-hip-oscilador} se integra directamente:
\[
T_n(t) = \frac{f_n}{\omega_n}\int_0^t\operatorname{sen}\omega_n(t-\tau)\,d\tau
= \frac{f_n}{\omega_n^2}\big(1 - \cos\omega_nt\big).
\] La \textbf{deflexión estática} del modo ---la que tendría la cuerda
si la fuerza se aplicara despacio y se esperara al equilibrio--- es la
solución constante de \ref{eq-hip-oscilador}, \(f_n/\omega_n^2\). La
respuesta a la carga súbita oscila entre \(0\) y \(2f_n/\omega_n^2\):
\textbf{cada modo llega al doble de su valor estático} cuando
\(\cos\omega_nt = -1\).

Es exactamente el factor 2 del sección~\ref{sec-ejemplo-transformada},
que allá salió de invertir una transformada. Acá sale modo por modo, y
se entiende mejor por qué: cada modo es un oscilador, y un oscilador al
que se le aplica una fuerza constante de golpe oscila alrededor de la
nueva posición de equilibrio, con amplitud igual a la distancia a la que
estaba. Arranca a una distancia \(f_n/\omega_n^2\) por debajo, así que
llega a la misma distancia por arriba. Es lo que pasa cuando uno se sube
de golpe a una balanza de resorte: la aguja pasa del peso verdadero
antes de asentarse.

En el sección~\ref{sec-ejemplo-transformada} la fuerza era
\(F_0\operatorname{sen}(\pi x)\), que excita un solo modo, y el factor 2
vale para la cuerda entera. Con una fuerza cualquiera hay una sutileza:
cada modo llega a su doble en su propio instante, cuando
\(\cos\omega_nt = -1\), y esos instantes no tienen por qué coincidir. Si
la carga es \textbf{simétrica} respecto del medio, los modos pares no se
excitan, y en \(t = l/a\) todos los impares tienen \(\cos(n\pi) = -1\) a
la vez: la cuerda entera llega al doble de su deflexión estática. Si no
es simétrica, el máximo de la cuerda completa hay que buscarlo.

\end{tcolorbox}

\subsection{Resonancia}\label{sec-hip-resonancia}

El caso más importante de fuerza es la \textbf{periódica}: un parlante,
un motor desbalanceado, el viento que desprende vórtices a ritmo
regular. Tomemos una fuerza que en cada punto oscila con la misma
frecuencia \(\Omega\), \(f(x,t) = g(x)\operatorname{sen}\Omega t\). Sus
coeficientes son \(f_n(t) = g_n\operatorname{sen}\Omega t\), y
\ref{eq-hip-oscilador} es un oscilador forzado sinusoidalmente. Si
\(\Omega\neq\omega_n\),

\begin{equation}\protect\phantomsection\label{eq-hip-resonancia}{
T_n(t) = \frac{g_n}{\omega_n^2-\Omega^2}\left(\operatorname{sen}\Omega t - \frac{\Omega}{\omega_n}\operatorname{sen}\omega_nt\right).
}\end{equation}

El primer término oscila con la frecuencia de la fuerza: es el
\textbf{régimen permanente}. El segundo oscila con la frecuencia propia
del modo: es el \textbf{transitorio}, que con algo de rozamiento se
apagaría (sección~\ref{sec-hip-casos-limite}).

El denominador lo dice todo. Cuando la frecuencia de la fuerza se acerca
a una frecuencia propia, la amplitud de ese modo se dispara. Y en el
límite:

\begin{definition}[Resonancia]\protect\hypertarget{def-hip-resonancia}{}\label{def-hip-resonancia}

Se dice que hay \textbf{resonancia} en el modo \(n\) cuando la
frecuencia del forzante coincide con su frecuencia propia,
\(\Omega = \omega_n\), y el coeficiente \(g_n\) no es cero. En ese caso

\begin{equation}\protect\phantomsection\label{eq-hip-resonancia-exacta}{
T_n(t) = \frac{g_n}{2\omega_n^2}\big(\operatorname{sen}\omega_nt - \omega_nt\cos\omega_nt\big),
}\end{equation}

cuya amplitud \textbf{crece linealmente con el tiempo} sin cota.

\end{definition}

\ref{eq-hip-resonancia-exacta} sale como el límite de
\ref{eq-hip-resonancia} cuando \(\Omega\to\omega_n\), con la regla de
L'Hôpital, o resolviendo directamente el oscilador. Dos observaciones
cambian cómo se lee:

\begin{itemize}
\tightlist
\item
  \textbf{Hace falta \(g_n\neq0\).} Si la fuerza no se parece en nada al
  modo, el modo no se entera. Una fuerza puntual aplicada justo en un
  nodo del modo \(n\) tiene \(g_n = 0\) y no lo excita, por más que la
  frecuencia sea la justa.
\item
  \textbf{El crecimiento lineal es una idealización.} Sin rozamiento, la
  energía que entra en cada ciclo no tiene por dónde salir. Con
  rozamiento la amplitud se estabiliza en un valor grande pero finito, y
  la resonancia se ve como un pico angosto en la curva de respuesta.
\end{itemize}

Video recomendado: CK-12, \emph{Standing Waves Part I: Demonstration},
\url{https://www.youtube.com/watch?v=-gr7KmTOrx0}.

\emph{(En la versión web hay acá un gráfico interactivo: una cuerda con
extremos fijos empujada en un punto por una fuerza que oscila a
frecuencia \(\Omega\). Arriba, la cuerda en el instante \(t\); abajo, la
curva de respuesta en escala logarítmica, con picos en las frecuencias
propias \(\omega_n = n\pi\). Se pueden mover \(\Omega\), el punto de
aplicación y el tiempo. Cerca de una frecuencia propia la amplitud crece
sin cota; si la fuerza se aplica en un nodo del modo, el pico
correspondiente desaparece.)}

Una cuerda con extremos fijos se empuja \textbf{en su punto medio} con
una fuerza que oscila. ¿Qué modos pueden entrar en resonancia?

\begin{itemize}
\tightlist
\item
  {Solo los impares: \(n = 1, 3, 5, \dots\)}
\item
  {Solo los pares: \(n = 2, 4, 6, \dots\)}
\item
  {Todos}
\item
  {Solo el fundamental}
\end{itemize}

Una fuerza puntual en \(x_f\) tiene coeficientes proporcionales a
\(\operatorname{sen}(n\pi x_f/l)\). En \(x_f = l/2\) eso es
\(\operatorname{sen}(n\pi/2)\), que vale cero para todo \(n\) par: los
modos pares tienen un nodo en el medio. Los impares tienen ahí un
vientre y se excitan todos.

Es el mismo mecanismo que el \textbf{truco del guitarrista} del
sección~\ref{sec-interpretacion}, visto del otro lado: allá, apoyar el
dedo en un nodo \textbf{elegía} qué modos sobrevivían; acá, empujar en
un nodo elige qué modos \textbf{no} se pueden excitar.

\begin{tcolorbox}[enhanced jigsaw, arc=.35mm, bottomrule=.15mm, breakable, colback=white, colframe=quarto-callout-note-color-frame, left=2mm, leftrule=.75mm, opacityback=0, rightrule=.15mm, toprule=.15mm]
\begin{minipage}[t]{5.5mm}
\textcolor{quarto-callout-note-color}{\faInfo}
\end{minipage}%
\begin{minipage}[t]{\textwidth - 5.5mm}

\vspace{-3mm}\textbf{Una advertencia sobre el puente de Tacoma}\vspace{3mm}

El derrumbe del puente de Tacoma Narrows, en 1940, aparece en muchos
libros como ejemplo de resonancia. No lo es, o no en el sentido de esta
sección. El viento no empujaba al puente a una frecuencia fija que
coincidiera con la del puente: el propio movimiento del tablero
modificaba el flujo del aire de manera que la fuerza lo realimentaba.
Eso se llama \textbf{flameo}, o autoexcitación aeroelástica, y es un
fenómeno distinto, que no se describe con una fuerza \(f(x,t)\) dada de
antemano. La aclaración se volvió conocida por un artículo de K. Y.
Billah y R. H. Scanlan de 1991 en el \emph{American Journal of Physics}.

\end{minipage}%
\end{tcolorbox}

\subsection{Bordes que se mueven}\label{sec-hip-bordes-moviles}

Falta un caso del problema \ref{eq-hip-mixto}: condiciones de contorno
no homogéneas, \(u(0,t) = \mu_1(t)\) y \(u(l,t) = \mu_2(t)\). La
separación de variables no las acepta directamente, porque necesita
condiciones homogéneas para que los modos sean los senos. La solución es
\textbf{pasar el problema del borde al interior}.

\begin{proposition}[Reducción a bordes
fijos]\protect\hypertarget{prp-hip-bordes-moviles}{}\label{prp-hip-bordes-moviles}

Sea \[
U(x,t) = \mu_1(t)\Big(1 - \frac{x}{l}\Big) + \mu_2(t)\,\frac{x}{l},
\] la recta que une los dos extremos en cada instante. Entonces
\(w = u - U\) cumple

\[
\begin{aligned}
&w_{tt} = a^2w_{xx} + f(x,t) - U_{tt}(x,t), && w(0,t) = w(l,t) = 0,\\
&w(x,0) = \varphi(x) - U(x,0), && w_t(x,0) = \psi(x) - U_t(x,0).
\end{aligned}
\]

\end{proposition}

\begin{proof}
\(U\) es lineal en \(x\), así que \(U_{xx} = 0\), y \(u = w+U\) en la
ecuación da \(w_{tt} + U_{tt} = a^2w_{xx} + f\). En los bordes,
\(U(0,t) = \mu_1(t)\) y \(U(l,t) = \mu_2(t)\), que se cancelan con los
de \(u\). Las condiciones iniciales salen de restar.
\end{proof}

El precio es que aparece una fuente, \(-U_{tt}\): \textbf{mover los
bordes equivale a empujar la cuerda desde adentro}. Visto desde un
sistema de referencia que acompaña a la recta \(U\), la cuerda siente
una fuerza de inercia, como un pasajero en un colectivo que frena.

La elección de \(U\) no es única: cualquier función que cumpla las
condiciones de borde sirve. La recta es la más simple, pero a veces
conviene otra que además cumpla la ecuación, y entonces \(w\) no tiene
fuente. Es lo que se hace en el ejercicio que sigue.

\begin{exercise}[]\protect\hypertarget{exr-hip-extremo-oscilante}{}\label{exr-hip-extremo-oscilante}

El extremo \(x=0\) de una cuerda está fijo y el extremo \(x=l\) se
sacude, \(u(l,t) = A\operatorname{sen}\Omega t\).

\begin{enumerate}
\def\labelenumi{\alph{enumi})}
\item
  Buscar una solución de la forma
  \(u_p = X(x)\operatorname{sen}\Omega t\) que cumpla la ecuación y las
  dos condiciones de contorno. ¿Para qué valores de \(\Omega\) no
  existe?
\item
  Explicar cómo se completa la solución si la cuerda arranca en reposo.
\end{enumerate}

\end{exercise}

\begin{tcolorbox}[enhanced jigsaw, arc=.35mm, bottomrule=.15mm, bottomtitle=1mm, breakable, colback=white, colbacktitle=quarto-callout-tip-color!10!white, colframe=quarto-callout-tip-color-frame, coltitle=black, left=2mm, leftrule=.75mm, opacityback=0, opacitybacktitle=0.6, rightrule=.15mm, title=\textcolor{quarto-callout-tip-color}{\faLightbulb}\hspace{0.5em}{Solución}, titlerule=0mm, toprule=.15mm, toptitle=1mm]

\textbf{a)} Reemplazando \(u_p = X(x)\operatorname{sen}\Omega t\) en
\(u_{tt} = a^2u_{xx}\) queda \(-\Omega^2X = a^2X''\), o sea
\(X'' + k^2X = 0\) con \(k = \Omega/a\). Con \(X(0) = 0\),
\(X = C\operatorname{sen}kx\); con \(X(l) = A\),
\(C = A/\operatorname{sen}kl\). Entonces \[
u_p(x,t) = A\,\frac{\operatorname{sen}(\Omega x/a)}{\operatorname{sen}(\Omega l/a)}\,\operatorname{sen}\Omega t .
\] No existe cuando \(\operatorname{sen}(\Omega l/a) = 0\), es decir
\(\Omega = n\pi a/l = \omega_n\): \textbf{exactamente en las frecuencias
propias}. Es la resonancia otra vez, ahora producida por un borde que se
mueve. Cerca de esas frecuencias, la amplitud en el interior es enorme
comparada con la del extremo: es lo que se ve en el video de más arriba,
donde el vibrador mueve la punta muy poco y la cuerda se sacude mucho.

\textbf{b)} \(u_p\) cumple la ecuación y los bordes, pero en \(t=0\)
tiene \(u_p(x,0) = 0\) y
\(\partial_tu_p(x,0) = A\Omega\operatorname{sen}(\Omega x/a)/\operatorname{sen}(\Omega l/a)\),
que no es la velocidad nula pedida. Así que \(u = u_p + w\), donde \(w\)
cumple la ecuación homogénea, bordes fijos, \(w(x,0) = 0\) y
\(w_t(x,0) = -\partial_tu_p(x,0)\). Ese \(w\) son oscilaciones libres, y
se resuelve por modos. Es el transitorio: con rozamiento se apaga, y
queda \(u_p\), el régimen permanente del caso 3 del
sección~\ref{sec-hip-casos-limite}.

Fijate que este \(u_p\) hace el papel de \(U\) en la
Proposición~\ref{prp-hip-bordes-moviles}, pero mejor elegido: como
cumple la ecuación, \(w\) quedó \textbf{sin fuente}.

\end{tcolorbox}

\begin{exercise}[]\protect\hypertarget{exr-hip-duhamel}{}\label{exr-hip-duhamel}

Verificar que el término de la fuente en \ref{eq-hip-dalembert-fuente},
\[v(x,t) = \frac{1}{2a}\int_0^t\int_{x-a(t-\tau)}^{x+a(t-\tau)} f(s,\tau)\,ds\,d\tau,\]
cumple \(v_{tt} = a^2v_{xx} + f\) y \(v(x,0) = v_t(x,0) = 0\).

\end{exercise}

\begin{tcolorbox}[enhanced jigsaw, arc=.35mm, bottomrule=.15mm, bottomtitle=1mm, breakable, colback=white, colbacktitle=quarto-callout-tip-color!10!white, colframe=quarto-callout-tip-color-frame, coltitle=black, left=2mm, leftrule=.75mm, opacityback=0, opacitybacktitle=0.6, rightrule=.15mm, title=\textcolor{quarto-callout-tip-color}{\faLightbulb}\hspace{0.5em}{Solución}, titlerule=0mm, toprule=.15mm, toptitle=1mm]

Llamemos
\(I(x,t,\tau) = \int_{x-a(t-\tau)}^{x+a(t-\tau)} f(s,\tau)\,ds\), de
modo que \(v = \frac{1}{2a}\int_0^t I\,d\tau\). Notar que
\(I(x,t,t) = 0\), porque el intervalo tiene largo cero.

\textbf{Primera derivada en \(t\).} Por la regla de Leibniz, \[
v_t = \frac{1}{2a}\Big[I(x,t,t) + \int_0^t I_t\,d\tau\Big] = \frac{1}{2a}\int_0^t I_t\,d\tau,
\qquad I_t = a\big[f(x+a(t-\tau),\tau) + f(x-a(t-\tau),\tau)\big].
\] Así que
\(v_t = \tfrac12\int_0^t\big[f(x+a(t-\tau),\tau) + f(x-a(t-\tau),\tau)\big]d\tau\).

\textbf{Segunda derivada en \(t\).} Otra vez Leibniz: el término de
borde, con \(\tau = t\), es \(\tfrac12[f(x,t)+f(x,t)] = f(x,t)\), y
queda \[
v_{tt} = f(x,t) + \frac a2\int_0^t\big[f_x(x+a(t-\tau),\tau) - f_x(x-a(t-\tau),\tau)\big]d\tau .
\]

\textbf{Derivadas en \(x\).}
\(I_x = f(x+a(t-\tau),\tau) - f(x-a(t-\tau),\tau)\), y derivando otra
vez, \[
v_{xx} = \frac{1}{2a}\int_0^t\big[f_x(x+a(t-\tau),\tau) - f_x(x-a(t-\tau),\tau)\big]d\tau .
\] Comparando, \(v_{tt} = f + a^2v_{xx}\).

\textbf{Condiciones iniciales.} En \(t = 0\) la integral en \(\tau\) es
sobre un intervalo de largo cero: \(v = 0\) y \(v_t = 0\).

Para derivar bajo la integral hace falta \(f_x\) continua: es la
hipótesis del Teorema~\ref{thm-hip-dalembert-fuente}.

\end{tcolorbox}

\section{Membranas}\label{sec-hip-membranas}

Una cuerda es una línea, y sus modos son senos. Un parche de tambor es
una superficie, y lo que cambia al pasar de una dimensión a dos no es
solo la cuenta. Aparecen tres fenómenos que la cuerda no tiene: las
frecuencias dejan de ser múltiplos de la fundamental, los nodos se
vuelven \textbf{líneas}, y aparecen modos distintos con la \textbf{misma
frecuencia}. Esta es la sección 9.3.1 de González, ampliada, con un
adelanto de la membrana circular.

\subsection{La ecuación de la membrana}\label{sec-hip-membrana-ec}

Una membrana es una lámina elástica muy delgada, sin rigidez propia,
estirada sobre un marco con una \textbf{tensión superficial} \(\tau\):
la fuerza por unidad de longitud que hace cada lado de un corte sobre el
otro. Llamamos \(\sigma\) a su densidad superficial de masa y
\(u(x,y,t)\) al desplazamiento vertical. El argumento es el de la cuerda
repetido en dos direcciones: sobre un rectangulito \(dx\times dy\), los
bordes paralelos al eje \(y\) reciben una fuerza vertical neta
\(\tau\,u_{xx}\,dx\,dy\), y los paralelos al eje \(x\) otra
\(\tau\,u_{yy}\,dx\,dy\). Con Newton,

\begin{equation}\protect\phantomsection\label{eq-hip-membrana}{
\sigma\,u_{tt} = \tau\,(u_{xx}+u_{yy}) + F
\qquad\Longrightarrow\qquad
u_{tt} = a^2\,\nabla^2u + f, \qquad a = \sqrt{\tau/\sigma}.
}\end{equation}

Las hipótesis son las mismas que las de la cuerda: desplazamientos
chicos, pendientes chicas, tensión uniforme. La lectura también:
\textbf{la aceleración de cada punto es proporcional a la curvatura
media de la superficie ahí}. Un borde fijo es \(u=0\) sobre el contorno.

\subsection{La membrana rectangular}\label{sec-hip-rectangular}

Sea la membrana \(0\le x\le l_1\), \(0\le y\le l_2\), con el borde fijo
y sin fuerzas externas. Seguimos los seis pasos del
sección~\ref{sec-esquema-separacion}.

\textbf{Separar el tiempo.} Con \(u = U(x,y)\,T(t)\), \[
\frac{T''}{a^2T} = \frac{\nabla^2U}{U} = -k^2 ,
\] porque el primer miembro depende solo de \(t\) y el segundo solo de
\((x,y)\). Quedan \(T'' + a^2k^2T = 0\) y la \textbf{ecuación de
Helmholtz}

\begin{equation}\protect\phantomsection\label{eq-hip-helmholtz-rect}{
\nabla^2U + k^2U = 0 \quad\text{en el rectángulo}, \qquad U = 0 \quad\text{en el borde}.
}\end{equation}

Es un problema de autovalores en dos variables: hay que encontrar los
\(k^2\) para los que tiene solución no nula. Es la ecuación de Helmholtz
del sección~\ref{sec-helmholtz}, ahora con borde.

\textbf{Separar el espacio.} Con \(U = X(x)\,Y(y)\), \[
\frac{X''}{X} + \frac{Y''}{Y} = -k^2 \quad\Longrightarrow\quad
X'' + \lambda^2X = 0,\ \ X(0)=X(l_1)=0; \qquad Y'' + \mu^2Y = 0,\ \ Y(0)=Y(l_2)=0,
\] con \(\lambda^2+\mu^2 = k^2\). Son dos problemas de la cuerda
(sección~\ref{sec-cuerda-separacion}), y ya sabemos sus soluciones:
\(X_m = \operatorname{sen}(m\pi x/l_1)\) y
\(Y_n = \operatorname{sen}(n\pi y/l_2)\).

\begin{tcolorbox}[enhanced jigsaw, arc=.35mm, bottomrule=.15mm, breakable, colback=white, colframe=quarto-callout-warning-color-frame, left=2mm, leftrule=.75mm, opacityback=0, rightrule=.15mm, toprule=.15mm]
\begin{minipage}[t]{5.5mm}
\textcolor{quarto-callout-warning-color}{\faExclamationTriangle}
\end{minipage}%
\begin{minipage}[t]{\textwidth - 5.5mm}

\vspace{-3mm}\textbf{Ojo con el libro de la cátedra (sección 9.3.1)}\vspace{3mm}

\begin{itemize}
\tightlist
\item
  La ecuación de Helmholtz aparece como \(U_{xx}+U_{yy}+k^2 = 0\): falta
  la \(U\) en el último término. Es \(U_{xx} + U_{yy} + k^2U = 0\).
\item
  En (9.13) y (9.14) dice \(\cos\varpi_{\sigma\tau}\) y
  \(\operatorname{sen}\varpi_{\sigma\tau}\) sin la \(t\): es
  \(\cos(\varpi_{\sigma\tau}t)\) y
  \(\operatorname{sen}(\varpi_{\sigma\tau}t)\).
\item
  La igualdad
  \(\omega_{\sigma\tau} = \tfrac{a\pi}{r}\sqrt{\sigma^2+\tau^2}\equiv a\sqrt{\sigma^2+\tau^2}\),
  y con ella \(\omega_{11} = a\sqrt2\), solo vale si el lado es
  \(r = \pi\). Con lado cualquiera, \(\omega_{11} = a\pi\sqrt2/r\).
\item
  Los índices se llaman \(\sigma\) y \(\tau\), y los lados \(l\) y
  \(m\). Acá usamos \(m,n\) para los índices y \(l_1,l_2\) para los
  lados, como en los problemas 9.9 y 9.10, para que \(m\) no sea a la
  vez un lado y un índice.
\end{itemize}

\end{minipage}%
\end{tcolorbox}

\begin{theorem}[Vibraciones libres de la membrana
rectangular]\protect\hypertarget{thm-hip-rectangular}{}\label{thm-hip-rectangular}

La solución de \(u_{tt} = a^2\nabla^2u\) en \([0,l_1]\times[0,l_2]\),
con \(u=0\) en el borde, \(u(x,y,0) = \varphi(x,y)\) y
\(u_t(x,y,0) = \psi(x,y)\), es

\begin{equation}\protect\phantomsection\label{eq-hip-rectangular}{
u(x,y,t) = \sum_{m,n\ge1}\big(A_{mn}\cos\omega_{mn}t + B_{mn}\operatorname{sen}\omega_{mn}t\big)
\operatorname{sen}\frac{m\pi x}{l_1}\,\operatorname{sen}\frac{n\pi y}{l_2},
}\end{equation}

con las \textbf{frecuencias propias}

\begin{equation}\protect\phantomsection\label{eq-hip-frecuencias-rect}{
\omega_{mn} = a\pi\sqrt{\frac{m^2}{l_1^2} + \frac{n^2}{l_2^2}}
}\end{equation}

y los coeficientes de la \textbf{serie doble de Fourier} de los datos:

\[
A_{mn} = \frac{4}{l_1l_2}\int_0^{l_1}\!\!\int_0^{l_2}\varphi(x,y)\operatorname{sen}\frac{m\pi x}{l_1}\operatorname{sen}\frac{n\pi y}{l_2}\,dy\,dx,
\qquad
\omega_{mn}B_{mn} = \frac{4}{l_1l_2}\int_0^{l_1}\!\!\int_0^{l_2}\psi(x,y)\operatorname{sen}\frac{m\pi x}{l_1}\operatorname{sen}\frac{n\pi y}{l_2}\,dy\,dx .
\]

\end{theorem}

\begin{proof}
Cada producto
\(\operatorname{sen}(m\pi x/l_1)\operatorname{sen}(n\pi y/l_2)\) cumple
\ref{eq-hip-helmholtz-rect} con \(k^2 = \pi^2(m^2/l_1^2 + n^2/l_2^2)\),
y el factor temporal cumple \(T'' + a^2k^2T=0\), de modo que
\(\omega_{mn} = ak\). La suma cumple la ecuación por linealidad.
Evaluando en \(t=0\),
\(\varphi = \sum A_{mn}\operatorname{sen}\operatorname{sen}\) y
\(\psi = \sum\omega_{mn}B_{mn}\operatorname{sen}\operatorname{sen}\);
multiplicando por
\(\operatorname{sen}(m'\pi x/l_1)\operatorname{sen}(n'\pi y/l_2)\) e
integrando, la ortogonalidad de los senos en cada variable deja solo el
término \(m=m'\), \(n=n'\), y cada integral de un seno al cuadrado
aporta un factor \(l_i/2\). De ahí el \(4/(l_1l_2)\).
\end{proof}

Es la cuerda dos veces, pero el resultado no se parece al de la cuerda.
Hay tres cosas nuevas.

\subsection{Uno: la membrana no tiene nota}\label{sec-hip-sin-nota}

En la cuerda, \(\omega_n = n\,\omega_1\): todas las frecuencias son
múltiplos enteros de la fundamental, y por eso la cuerda suena con una
altura definida (sección~\ref{sec-interpretacion}). En la membrana,
\ref{eq-hip-frecuencias-rect} tiene una \textbf{raíz cuadrada}, y los
cocientes dejan de ser enteros.

\begin{longtable}[]{@{}
  >{\raggedright\arraybackslash}p{(\linewidth - 2\tabcolsep) * \real{0.2600}}
  >{\raggedright\arraybackslash}p{(\linewidth - 2\tabcolsep) * \real{0.7400}}@{}}
\caption{Por qué una guitarra tiene nota y un tambor no. En la membrana
cuadrada los cocientes son \(\sqrt{(m^2+n^2)/2}\); en la circular salen
de los ceros de las funciones de Bessel
(sección~\ref{sec-hip-circular}).}\label{tbl-hip-cocientes}\tabularnewline
\toprule\noalign{}
\begin{minipage}[b]{\linewidth}\raggedright
Sistema
\end{minipage} & \begin{minipage}[b]{\linewidth}\raggedright
Cocientes \(\omega/\omega_{\text{fundamental}}\) de los primeros modos
\end{minipage} \\
\midrule\noalign{}
\endfirsthead
\toprule\noalign{}
\begin{minipage}[b]{\linewidth}\raggedright
Sistema
\end{minipage} & \begin{minipage}[b]{\linewidth}\raggedright
Cocientes \(\omega/\omega_{\text{fundamental}}\) de los primeros modos
\end{minipage} \\
\midrule\noalign{}
\endhead
\bottomrule\noalign{}
\endlastfoot
Cuerda & \(1,\ 2,\ 3,\ 4,\ 5,\ 6\) \\
Membrana cuadrada & \(1,\ 1{,}581,\ 2,\ 2{,}236,\ 2{,}550,\ 2{,}915\) \\
Membrana circular &
\(1,\ 1{,}593,\ 2{,}136,\ 2{,}295,\ 2{,}653,\ 2{,}917\) \\
\end{longtable}

Un sonido con componentes a frecuencias \(1, 1{,}58, 2, 2{,}24\dots\) no
es armónico: el oído no lo agrupa en una sola nota. Es la respuesta
completa a la pregunta 24 del Capítulo~\ref{sec-edp}. Los timbales de
orquesta, que sí tienen altura definida, lo logran con trucos físicos
---el aire encerrado en la caldera y la carga del aire sobre el parche
corren las frecuencias hacia cocientes casi armónicos--- que están fuera
de este modelo.

En una membrana \textbf{cuadrada}, ¿cuánto vale el cociente
\(\omega_{12}/\omega_{11}\)? Dá tres decimales.

Con \(l_1 = l_2 = l\), \ref{eq-hip-frecuencias-rect} da
\(\omega_{mn} = \tfrac{a\pi}{l}\sqrt{m^2+n^2}\), así que
\[\frac{\omega_{12}}{\omega_{11}} = \frac{\sqrt{1+4}}{\sqrt{1+1}} = \sqrt{\tfrac52} \approx 1{,}581 .\]
No es entero ni un cociente simple: el segundo modo del cuadrado no es
un armónico del primero. Fijate además que \(\omega_{21}\) da lo mismo.
Eso es lo que sigue.

\subsection{Dos: los nodos son líneas}\label{sec-hip-lineas-nodales}

\begin{definition}[Línea
nodal]\protect\hypertarget{def-hip-linea-nodal}{}\label{def-hip-linea-nodal}

En un modo \(u = U(x,y)\,T(t)\), se llama \textbf{línea nodal} a cada
curva del conjunto \(\{(x,y): U(x,y) = 0\}\) en el interior de la
membrana. Sobre ella la membrana no se mueve nunca.

\end{definition}

En el rectángulo, el modo \((m,n)\) tiene \(m-1\) líneas nodales
verticales, en \(x = j\,l_1/m\), y \(n-1\) horizontales, en
\(y = j\,l_2/n\): una grilla. Las líneas nodales parten la membrana en
\(mn\) celdas, y celdas vecinas se mueven en contrafase: cuando una
sube, las de al lado bajan.

\subsection{Tres: modos distintos con la misma
frecuencia}\label{sec-hip-degeneracion}

En el cuadrado, \ref{eq-hip-frecuencias-rect} depende de \(m^2+n^2\),
así que los modos \((1,2)\) y \((2,1)\) tienen \textbf{la misma
frecuencia} aunque sean formas distintas. Y como la ecuación es lineal,
cualquier combinación \[
U = \alpha\,\operatorname{sen}\frac{\pi x}{l}\operatorname{sen}\frac{2\pi y}{l} + \beta\,\operatorname{sen}\frac{2\pi x}{l}\operatorname{sen}\frac{\pi y}{l}
\] también es un modo: cambia de tamaño en el tiempo con la frecuencia
común \(\omega_{12}\) y conserva la forma.

\begin{definition}[Autovalor
degenerado]\protect\hypertarget{def-hip-degeneracion}{}\label{def-hip-degeneracion}

Un autovalor \(k^2\) de \ref{eq-hip-helmholtz-rect} es
\textbf{degenerado} si le corresponden dos o más autofunciones
linealmente independientes. La cantidad máxima de independientes es su
\textbf{multiplicidad}, y cualquier combinación lineal de ellas es
también autofunción.

\end{definition}

Esto es lo que dice González con «por combinación lineal podemos lograr
infinitas vibraciones de la misma frecuencia». Lo que no dice es lo más
llamativo: \textbf{las líneas nodales de la combinación no son las de
ninguno de los dos modos}. Con \(\alpha = \beta\) la línea nodal es una
diagonal; con \(\alpha = -\beta\), la otra; con valores intermedios,
curvas. Es el Ejercicio~\ref{exr-hip-diagonal}, y se ve en el gráfico
que sigue.

La degeneración viene de una \textbf{simetría}: el cuadrado no cambia si
se intercambian \(x\) e \(y\), y entonces tampoco sus frecuencias. En un
rectángulo de lados distintos esa simetría se rompe, las dos frecuencias
se separan, y la combinación deja de ser un modo: las dos partes laten a
ritmos distintos y la forma cambia todo el tiempo. En la mecánica
cuántica, donde la ecuación de Schrödinger estacionaria es también una
ecuación de Helmholtz, la degeneración de los niveles de energía tiene
exactamente este origen.

\emph{(En la versión web hay acá un gráfico interactivo en tres
dimensiones y girable: la membrana rectangular en un modo \((m,n)\) o en
una mezcla de los modos \((m,n)\) y \((n,m)\), con sus líneas nodales
dibujadas en el piso. En el cuadrado la mezcla de \((1,2)\) y \((2,1)\)
tiene una diagonal como línea nodal, que no se mueve; si se estira un
lado, las dos frecuencias se separan y las líneas nodales empiezan a
moverse.)}

\begin{tcolorbox}[enhanced jigsaw, arc=.35mm, bottomrule=.15mm, breakable, colback=white, colframe=quarto-callout-note-color-frame, left=2mm, leftrule=.75mm, opacityback=0, rightrule=.15mm, toprule=.15mm]
\begin{minipage}[t]{5.5mm}
\textcolor{quarto-callout-note-color}{\faInfo}
\end{minipage}%
\begin{minipage}[t]{\textwidth - 5.5mm}

\vspace{-3mm}\textbf{Las figuras de Chladni son de placas, no de membranas}\vspace{3mm}

González llama figuras de Chladni a las líneas nodales de la membrana.
El nombre viene de Ernst Chladni, que a fines del siglo XVIII echaba
arena sobre \textbf{placas metálicas} y las hacía vibrar con un arco de
violín: la arena se junta sobre las líneas nodales y dibuja la figura.
Pero una placa no es una membrana: resiste la flexión, como la viga del
sección~\ref{sec-viga}, y su ecuación es de cuarto orden,
\(u_{tt} + \kappa^2\nabla^4u = 0\). Además la placa de Chladni suele
estar sujeta solo en el centro, con los bordes libres. Las figuras se
parecen a las de la membrana en espíritu ---líneas donde no hay
movimiento, degeneraciones que producen curvas--- pero no son las mismas
funciones.

\end{minipage}%
\end{tcolorbox}

Video recomendado: \emph{Chladni Figures --- random couscous snaps into
beautiful patterns}, \url{https://www.youtube.com/watch?v=CR_XL192wXw}.

\begin{exercise}[]\protect\hypertarget{exr-hip-diagonal}{}\label{exr-hip-diagonal}

En la membrana cuadrada \([0,1]\times[0,1]\), hallar las líneas nodales
de
\(U_\pm = \operatorname{sen}\pi x\operatorname{sen}2\pi y \pm \operatorname{sen}2\pi x\operatorname{sen}\pi y\).

\end{exercise}

\begin{tcolorbox}[enhanced jigsaw, arc=.35mm, bottomrule=.15mm, bottomtitle=1mm, breakable, colback=white, colbacktitle=quarto-callout-tip-color!10!white, colframe=quarto-callout-tip-color-frame, coltitle=black, left=2mm, leftrule=.75mm, opacityback=0, opacitybacktitle=0.6, rightrule=.15mm, title=\textcolor{quarto-callout-tip-color}{\faLightbulb}\hspace{0.5em}{Solución}, titlerule=0mm, toprule=.15mm, toptitle=1mm]

Usando
\(\operatorname{sen}2\theta = 2\operatorname{sen}\theta\cos\theta\), \[
U_\pm = 2\operatorname{sen}\pi x\operatorname{sen}\pi y\,\big(\cos\pi y \pm \cos\pi x\big).
\] En el interior del cuadrado, \(\operatorname{sen}\pi x\) y
\(\operatorname{sen}\pi y\) no se anulan, así que las líneas nodales son
los ceros del paréntesis.

\begin{itemize}
\tightlist
\item
  \textbf{Con el signo menos}, \(\cos\pi y = \cos\pi x\). Como \(\pi x\)
  y \(\pi y\) están en \((0,\pi)\), donde el coseno es inyectivo, eso es
  \(x = y\): \textbf{la diagonal}.
\item
  \textbf{Con el signo más},
  \(\cos\pi y = -\cos\pi x = \cos(\pi - \pi x)\), o sea \(y = 1 - x\):
  \textbf{la otra diagonal}.
\end{itemize}

Ninguno de los dos modos originales tiene líneas nodales diagonales: el
\((1,2)\) tiene la recta \(y = \tfrac12\) y el \((2,1)\) la recta
\(x=\tfrac12\). La diagonal aparece solo en la combinación, y solo
porque el cuadrado es simétrico respecto de ella.

\end{tcolorbox}

\subsection{Un adelanto: la membrana circular}\label{sec-hip-circular}

Los tambores son redondos. En un disco de radio \(R_0\) con el borde
fijo, lo natural es usar coordenadas polares \((r,\theta)\), y los
problemas 9.11 y 9.12 de la cátedra piden el caso en que todo tiene
\textbf{simetría radial}: \(u = u(r,t)\), sin dependencia en \(\theta\).
Entonces el laplaciano en polares se reduce a su parte radial y la
ecuación es

\begin{equation}\protect\phantomsection\label{eq-hip-circular}{
u_{tt} = a^2\Big(u_{rr} + \frac1r\,u_r\Big), \qquad 0\le r<R_0, \qquad u(R_0,t) = 0 .
}\end{equation}

Separando, \(u = R(r)\,T(t)\), sale otra vez \(T'' + a^2k^2T = 0\) y,
para la parte radial,

\begin{equation}\protect\phantomsection\label{eq-hip-bessel-radial}{
R'' + \frac1r\,R' + k^2R = 0 .
}\end{equation}

El término \(R'/r\) es lo único nuevo, y alcanza para que las soluciones
\textbf{no} sean senos ni cosenos. Con el cambio \(s = kr\) la ecuación
queda \(R'' + R'/s + R = 0\), que es la \textbf{ecuación de Bessel de
orden cero}. El capítulo 12 está dedicado a estas funciones; acá va lo
justo para resolver los problemas.

\begin{definition}[Función de Bessel de orden
cero]\protect\hypertarget{def-hip-bessel-j0}{}\label{def-hip-bessel-j0}

Se llama \textbf{función de Bessel de primera especie de orden cero} a

\begin{equation}\protect\phantomsection\label{eq-hip-j0}{
J_0(s) = \sum_{j=0}^{\infty}\frac{(-1)^j}{(j!)^2}\Big(\frac s2\Big)^{2j}
= 1 - \frac{s^2}{4} + \frac{s^4}{64} - \frac{s^6}{2304} + \cdots
}\end{equation}

La serie converge para todo \(s\). \(J_0\) es la única solución de
\(R'' + R'/s + R = 0\), salvo múltiplos, que es \textbf{acotada en
\(s=0\)}; la otra solución independiente, \(Y_0\), tiende a \(-\infty\)
como \(\tfrac2\pi\ln s\). Se usa también \(J_1 = -J_0'\).

\end{definition}

Lo que hay que saber de \(J_0\), sin demostración (capítulo 12):

\begin{enumerate}
\def\labelenumi{\arabic{enumi}.}
\tightlist
\item
  \textbf{Se parece a un coseno que se amortigua.} Para \(s\) grande,
  \(J_0(s)\approx\sqrt{2/(\pi s)}\,\cos(s - \pi/4)\): oscila con
  amplitud que decae como \(1/\sqrt s\). Es el «coseno del círculo»: la
  amplitud decrece porque la misma energía se reparte en anillos cada
  vez más largos.
\item
  \textbf{Tiene infinitos ceros positivos},
  \(j_{0,1} < j_{0,2} < \cdots\), que no son equiespaciados al principio
  y se acercan a una separación \(\pi\) después:
  \(j_{0,1}\approx 2{,}4048\), \(j_{0,2}\approx 5{,}5201\),
  \(j_{0,3}\approx 8{,}6537\), \(j_{0,4}\approx 11{,}7915\).
\item
  \textbf{Ortogonalidad con peso \(r\).} La ecuación
  \ref{eq-hip-bessel-radial} multiplicada por \(r\) es
  \((rR')' + k^2\,rR = 0\), un problema de Sturm--Liouville con peso
  \(r\) (Capítulo~\ref{sec-sturm}). Por eso \[
  \int_0^{R_0} J_0\!\Big(\frac{j_{0,p}\,r}{R_0}\Big)J_0\!\Big(\frac{j_{0,q}\,r}{R_0}\Big)\,r\,dr =
  \begin{cases}0, & p\neq q,\\[1mm] \tfrac{R_0^2}{2}\,J_1(j_{0,p})^2, & p = q.\end{cases}
  \]
\end{enumerate}

\begin{figure}

\centering{

\pandocbounded{\includegraphics[keepaspectratio]{09-hiperbolicas_files/figure-pdf/fig-hip-bessel-output-1.pdf}}

}

\caption{\label{fig-hip-bessel}La función de Bessel \(J_0\) y los modos
radiales del tambor. \textbf{Izquierda:} \(J_0(s)\) con sus primeros
ceros marcados, junto a la aproximación
\(\sqrt{2/(\pi s)}\cos(s-\pi/4)\) en punteado, que ya es muy buena a
partir de \(s\approx 3\). \textbf{Derecha:} los tres primeros modos
radiales de la membrana circular, \(J_0(j_{0,p}\,r/R_0)\). El borde
tiene que caer en un cero de \(J_0\), así que el modo \(p\) usa el
pedazo de la curva que va hasta el cero número \(p\), y tiene \(p-1\)
círculos nodales.}

\end{figure}%

La condición de borde, \(R(R_0) = 0\), obliga a que \(kR_0\) sea un cero
de \(J_0\). Así se discretizan las frecuencias, igual que en la cuerda,
con la diferencia de que los ceros de \(J_0\) no son múltiplos de
\(\pi\):

\begin{theorem}[Vibraciones radiales de la membrana
circular]\protect\hypertarget{thm-hip-circular}{}\label{thm-hip-circular}

La solución de \ref{eq-hip-circular}, acotada en \(r=0\), con
\(u(r,0) = \varphi(r)\) y \(u_t(r,0) = \psi(r)\), es

\begin{equation}\protect\phantomsection\label{eq-hip-circular-sol}{
u(r,t) = \sum_{p\ge1}\big(A_p\cos\omega_pt + B_p\operatorname{sen}\omega_pt\big)\,J_0\!\Big(\frac{j_{0,p}\,r}{R_0}\Big),
\qquad \omega_p = \frac{a\,j_{0,p}}{R_0},
}\end{equation}

con

\[
A_p = \frac{2}{R_0^2J_1(j_{0,p})^2}\int_0^{R_0}\varphi(r)\,J_0\!\Big(\frac{j_{0,p}\,r}{R_0}\Big)\,r\,dr,
\qquad
B_p = \frac{2}{a\,j_{0,p}\,R_0\,J_1(j_{0,p})^2}\int_0^{R_0}\psi(r)\,J_0\!\Big(\frac{j_{0,p}\,r}{R_0}\Big)\,r\,dr .
\]

\end{theorem}

\begin{proof}
Es el esquema del sección~\ref{sec-esquema-separacion}. Los modos son
\(J_0(j_{0,p}r/R_0)\,T_p(t)\): la parte radial cumple
\ref{eq-hip-bessel-radial} con \(k = j_{0,p}/R_0\), es acotada en el
centro y se anula en el borde. La acotación en \(r = 0\) es la condición
que el paso 2 del esquema advertía que aparece sin estar escrita:
descarta \(Y_0\). Los coeficientes salen de multiplicar los datos por
\(J_0(j_{0,q}r/R_0)\,r\), integrar y usar la ortogonalidad con peso
\(r\); el \(B_p\) lleva además el factor \(\omega_p\) que aparece al
derivar en \(t\). La completitud de estas funciones ---que cualquier
dato razonable se puede desarrollar en ellas--- es un teorema del
capítulo 12.
\end{proof}

Los modos radiales son \textbf{anillos}: el modo \(p\) tiene \(p-1\)
círculos nodales, en los radios \(r = R_0\,j_{0,q}/j_{0,p}\) con
\(q<p\). Y los cocientes de frecuencias son
\(j_{0,2}/j_{0,1}\approx 2{,}295\), \(j_{0,3}/j_{0,1}\approx 3{,}598\):
nada de enteros. Un tambor golpeado en el centro excita sobre todo estos
modos, y suena sin altura definida.

\begin{tcolorbox}[enhanced jigsaw, arc=.35mm, bottomrule=.15mm, breakable, colback=white, colframe=quarto-callout-note-color-frame, left=2mm, leftrule=.75mm, opacityback=0, rightrule=.15mm, toprule=.15mm]
\begin{minipage}[t]{5.5mm}
\textcolor{quarto-callout-note-color}{\faInfo}
\end{minipage}%
\begin{minipage}[t]{\textwidth - 5.5mm}

\vspace{-3mm}\textbf{Los modos que no son radiales}\vspace{3mm}

Sin simetría radial la separación es \(u = R(r)\,\Theta(\theta)\,T(t)\),
la parte angular da \(\cos(n\theta)\) y \(\operatorname{sen}(n\theta)\),
y la radial es la ecuación de Bessel de orden \(n\), con soluciones
\(J_n\). Aparecen entonces \textbf{diámetros nodales} además de los
círculos, y cada modo con \(n\ge1\) es degenerado ---\(\cos n\theta\) y
\(\operatorname{sen}n\theta\) tienen la misma frecuencia---, por la
simetría de rotación del disco. Los cocientes de la
Tabla~\ref{tbl-hip-cocientes} para la membrana circular incluyen esos
modos. Todo esto es del capítulo 12.

\end{minipage}%
\end{tcolorbox}

¿Cuántas veces más aguda es la frecuencia del segundo modo radial de un
tambor que la del primero? Usá los ceros de \(J_0\) de más arriba y dá
tres decimales.

Por \ref{eq-hip-circular-sol}, \(\omega_p = a\,j_{0,p}/R_0\), así que
\[\frac{\omega_2}{\omega_1} = \frac{j_{0,2}}{j_{0,1}} = \frac{5{,}5201}{2{,}4048}\approx 2{,}295 .\]
En una cuerda el segundo modo es exactamente una octava arriba, un
cociente \(2\). En el tambor es un poco más que una octava, y los
siguientes se alejan cada vez más de cualquier serie armónica.

\begin{exercise}[]\protect\hypertarget{exr-hip-serie-j0}{}\label{exr-hip-serie-j0}

Buscar una solución de \(R'' + \tfrac1s R' + R = 0\) en forma de serie
de potencias, \(R = \sum_{j\ge0}c_js^j\), con \(c_0 = 1\). Mostrar que
los coeficientes impares son nulos y obtener \ref{eq-hip-j0}.

\end{exercise}

\begin{tcolorbox}[enhanced jigsaw, arc=.35mm, bottomrule=.15mm, bottomtitle=1mm, breakable, colback=white, colbacktitle=quarto-callout-tip-color!10!white, colframe=quarto-callout-tip-color-frame, coltitle=black, left=2mm, leftrule=.75mm, opacityback=0, opacitybacktitle=0.6, rightrule=.15mm, title=\textcolor{quarto-callout-tip-color}{\faLightbulb}\hspace{0.5em}{Solución}, titlerule=0mm, toprule=.15mm, toptitle=1mm]

Multiplicando la ecuación por \(s\) para no tener potencias negativas,
\(sR'' + R' + sR = 0\). Con la serie, \[
sR'' + R' = \sum_{j\ge1}\big[j(j-1) + j\big]c_j\,s^{j-1} = \sum_{j\ge1}j^2c_j\,s^{j-1},
\qquad sR = \sum_{j\ge0}c_j\,s^{j+1}.
\] Igualando a cero el coeficiente de cada potencia \(s^{j-1}\):

\begin{itemize}
\tightlist
\item
  Para \(j = 1\): \(1^2\,c_1 = 0\), así que \(c_1 = 0\).
\item
  Para \(j\ge2\): \(j^2c_j + c_{j-2} = 0\), es decir
  \(c_j = -c_{j-2}/j^2\).
\end{itemize}

Como \(c_1 = 0\), la recursión anula todos los impares. Para los pares,
con \(j = 2i\), \[
c_{2i} = -\frac{c_{2i-2}}{4i^2} = \frac{(-1)^i}{4^i\,(i!)^2}\,c_0 = \frac{(-1)^i}{(i!)^2}\Big(\frac12\Big)^{2i},
\] que con \(c_0 = 1\) es \ref{eq-hip-j0}. La serie converge para todo
\(s\) por el criterio del cociente:
\(|c_{2i}s^{2i}/c_{2i-2}s^{2i-2}| = s^2/(4i^2)\to0\).

La otra solución no se puede obtener así: el método de Frobenius
(capítulo 12) muestra que lleva un \(\ln s\), y por eso no es una serie
de potencias. Es la \(Y_0\), la que descartamos por no ser acotada en el
centro del tambor.

\end{tcolorbox}

\section{Ondas en el espacio}\label{sec-hip-espacio}

El sonido y la luz no viajan por una cuerda: se expanden en tres
dimensiones. La ecuación es

\begin{equation}\protect\phantomsection\label{eq-hip-onda-3d}{
\nabla^2u = \frac{1}{a^2}\,u_{tt} - f(\mathbf x,t), \qquad \mathbf x\in\mathbb R^3,
}\end{equation}

que es \ref{eq-hip-onda-n} con \(n=3\) escrita como la escribe González
en (9.11). Se pueden plantear problemas de Cauchy en todo el espacio o
problemas de contorno en una región acotada ---una sala, una cavidad
resonante, una guía de ondas---. Esta sección hace el primer tipo, y
descubre que el número de dimensiones cambia la física de una manera que
nadie se esperaría.

\subsection{Ondas esféricas}\label{sec-hip-esfericas}

El caso más simple es el de simetría central: \(u = u(r,t)\), con
\(r = |\mathbf x|\). Es el sonido de un globo que revienta, o de una
fuente puntual que emite igual en todas las direcciones. En esféricas,
sin dependencia angular, el laplaciano tiene una forma que parece hecha
a propósito:

\[
\nabla^2u = u_{rr} + \frac2r\,u_r = \frac1r\,\frac{\partial^2}{\partial r^2}\big(r\,u\big).
\]

(La segunda igualdad se verifica derivando:
\((ru)_{rr} = r\,u_{rr} + 2u_r\).) Entonces, sin fuente, la ecuación es
\(\tfrac1r(ru)_{rr} = \tfrac1{a^2}u_{tt}\), y multiplicando por \(r\):

\begin{proposition}[Reducción a una
dimensión]\protect\hypertarget{prp-hip-esfericas}{}\label{prp-hip-esfericas}

Si \(u(r,t)\) es solución radial de la ecuación de onda en
\(\mathbb R^3\), entonces \(v = r\,u\) cumple la ecuación de la cuerda,
\[
v_{rr} = \frac1{a^2}\,v_{tt}, \qquad r>0,
\] y si \(u\) es acotada en el origen, \(v(0,t) = 0\). En consecuencia

\begin{equation}\protect\phantomsection\label{eq-hip-esfericas}{
u(r,t) = \frac{F(t - r/a)}{r} + \frac{G(t + r/a)}{r},
}\end{equation}

con \(G = -F\) si la solución es regular en el origen.

\end{proposition}

\begin{proof}
La primera parte es la cuenta de arriba. Si \(u\) es acotada cerca de
\(r=0\), \(v = ru\to0\). La solución general de la cuerda es
\(v = F(t-r/a) + G(t+r/a)\), escrita como función de \(t\mp r/a\) en
lugar de \(r\mp at\), que es lo mismo salvo un cambio de nombre. La
condición \(v(0,t) = 0\) da \(F(t) + G(t) = 0\) para todo \(t\).
\end{proof}

\begin{definition}[Ondas
esféricas]\protect\hypertarget{def-hip-onda-esferica}{}\label{def-hip-onda-esferica}

Las soluciones \(\dfrac{F(t-r/a)}{r}\) y \(\dfrac{G(t+r/a)}{r}\) se
llaman \textbf{ondas esféricas} \textbf{divergente} (o saliente) y
\textbf{convergente} (o entrante), respectivamente.

\end{definition}

Hay dos cosas para leer en \ref{eq-hip-esfericas}.

\textbf{La amplitud cae como \(1/r\).} El perfil \(F\) viaja sin
deformarse, como en la cuerda, pero achicándose. La razón es la
conservación de la energía: la energía de la onda está repartida sobre
una esfera de área \(4\pi r^2\), así que la intensidad ---energía por
unidad de área--- cae como \(1/r^2\), y como la energía va con el
cuadrado de la amplitud, la amplitud cae como \(1/r\). Es la onda
esférica \(e^{ikr}/r\) de la tabla de Green del
sección~\ref{sec-green-caso-1}, ahora con cualquier perfil.

\textbf{El origen es un extremo fijo.} La condición \(v(0,t) = 0\) es la
de la cuerda semi-infinita del sección~\ref{sec-hip-semi-fijo}. Una onda
convergente que llega al centro \textbf{rebota como en un extremo fijo}:
sale como onda divergente, con \(F = -G\). Es un eco en el centro de la
esfera.

\begin{tcolorbox}[enhanced jigsaw, arc=.35mm, bottomrule=.15mm, breakable, colback=white, colframe=quarto-callout-warning-color-frame, left=2mm, leftrule=.75mm, opacityback=0, rightrule=.15mm, toprule=.15mm]
\begin{minipage}[t]{5.5mm}
\textcolor{quarto-callout-warning-color}{\faExclamationTriangle}
\end{minipage}%
\begin{minipage}[t]{\textwidth - 5.5mm}

\vspace{-3mm}\textbf{Ojo con el libro de la cátedra (sección 9.2.3)}\vspace{3mm}

Al volver a las variables originales, González escribe
\(u = \tfrac1r f_1(t-r/a) + \tfrac1r f_1(t+r/a)\), con \(f_1\) dos
veces: el segundo es \(f_2\). Y en el último renglón, con
\(f_2 = -f_1 = f\), el resultado correcto es
\[u(r,t) = \frac{f(t+r/a) - f(t-r/a)}{r},\] con un \textbf{signo menos}
que en el libro no aparece. Se puede comprobar que es el signo bueno:
con el signo menos, el cociente tiende a \(2f'(t)/a\) cuando \(r\to0\) y
la solución es regular; con un signo más, divergiría como \(2f(t)/r\) en
el origen, que es justo lo que la condición \(v(0,t)=0\) tenía que
impedir.

\end{minipage}%
\end{tcolorbox}

\subsection{El problema de Cauchy con simetría
esférica}\label{sec-hip-cauchy-esferico}

Con datos radiales, \(u(r,0) = \varphi(r)\) y \(u_t(r,0) = \psi(r)\), la
función \(v = ru\) es una cuerda semi-infinita con extremo fijo y datos
\(r\varphi(r)\), \(r\psi(r)\). La Proposición~\ref{prp-hip-semi-fijo} la
resuelve entera. Para \(r\ge at\),

\begin{equation}\protect\phantomsection\label{eq-hip-cauchy-esferico}{
u(r,t) = \frac{(r+at)\,\varphi(r+at) + (r-at)\,\varphi(r-at)}{2r} + \frac{1}{2ar}\int_{r-at}^{r+at}s\,\psi(s)\,ds ,
}\end{equation}

y para \(r<at\) vale lo mismo reemplazando \((r-at)\varphi(r-at)\) por
\(-(at-r)\varphi(at-r)\) y el límite inferior de la integral por
\(at-r\). Ese reemplazo es exactamente el eco en el centro.

Y acá aparece algo que en una dimensión no podía pasar. En el centro,
tomando el límite \(r\to0\),

\begin{equation}\protect\phantomsection\label{eq-hip-centro}{
u(0,t) = \frac{d}{ds}\Big[s\,\varphi(s)\Big]_{s=at} + t\,\psi(at) = \varphi(at) + at\,\varphi'(at) + t\,\psi(at) .
}\end{equation}

\textbf{En el centro aparece la derivada del dato.} Si \(\varphi\) tiene
un escalón ---un globo con aire a más presión que afuera, que
revienta---, \(\varphi'\) es una delta, y la solución en el centro tiene
un pico infinito en el instante en que el frente, que viajaba hacia
adentro, llega al centro. Es el \textbf{enfoque} de una onda
convergente: toda la energía que estaba repartida en una esfera se
concentra en un punto. En una dimensión la solución nunca es peor que el
dato; en tres, sí. Es lo que anunciaba el recuadro «Lo que la estimación
no dice» del sección~\ref{sec-hip-estabilidad}.

\begin{figure}

\centering{

\pandocbounded{\includegraphics[keepaspectratio]{09-hiperbolicas_files/figure-pdf/fig-hip-globo-output-1.pdf}}

}

\caption{\label{fig-hip-globo}Un globo que revienta, con \(a=1\). El
dato inicial es un exceso de presión \(\varphi(r)\) que vale \(1\)
adentro de una esfera de radio \(b=1\) y cae a cero en un borde suave;
la velocidad inicial es nula. \textbf{Izquierda:} el perfil \(u(r,t)\)
en cuatro instantes. La parte que viaja hacia afuera es una \textbf{onda
N} ---una compresión seguida de una rarefacción--- cuya amplitud cae
como \(1/r\). \textbf{Derecha:} la presión en el centro,
\ref{eq-hip-centro}. Se queda en \(1\) mientras la noticia del borde no
llega, y en \(t = b/a\) la parte convergente del frente se enfoca en el
centro y produce un pico negativo, tanto más alto cuanto más abrupto es
el borde del globo.}

\end{figure}%

Una fuente puntual de sonido emite en todas direcciones. A \(10\) m se
miden \(60\) dB. ¿Cuánto se mide a \(20\) m, si nada absorbe el sonido?

\begin{itemize}
\tightlist
\item
  {\(54\) dB: la intensidad cae a la cuarta parte}
\item
  {\(57\) dB: la intensidad cae a la mitad}
\item
  {\(30\) dB: la intensidad cae a la mitad}
\item
  {\(60\) dB: el sonido no se pierde, solo viaja}
\end{itemize}

La amplitud cae como \(1/r\) y la intensidad como \(1/r^2\): al duplicar
la distancia la intensidad se divide por \(4\). Los decibeles miden
\(10\log_{10}\) del cociente de intensidades, y
\(10\log_{10}4 \approx 6\): se pierden \(6\) dB. La segunda opción
confunde amplitud con intensidad ---sería lo que se pierde si la
intensidad cayera como \(1/r\)---. La tercera confunde decibeles con una
escala lineal. La cuarta tiene razón en que la energía no se pierde,
pero se reparte en una esfera cuatro veces más grande.

\subsection{El principio de Huygens}\label{sec-hip-huygens}

Para datos sin simetría hay una fórmula general, que se debe a Kirchhoff
y a Poisson. Llamemos \(M_\rho[g](\mathbf x)\) al \textbf{promedio} de
una función \(g\) sobre la esfera de centro \(\mathbf x\) y radio
\(\rho\).

\begin{theorem}[Fórmula de Kirchhoff para el problema de Cauchy en
\(\mathbb R^3\)]\protect\hypertarget{thm-hip-kirchhoff-cauchy}{}\label{thm-hip-kirchhoff-cauchy}

Si \(\varphi\in C^3\) y \(\psi\in C^2\), la solución de
\(u_{tt} = a^2\nabla^2u\) en \(\mathbb R^3\) con
\(u(\mathbf x,0) = \varphi\) y \(u_t(\mathbf x,0) = \psi\) es

\begin{equation}\protect\phantomsection\label{eq-hip-kirchhoff-cauchy}{
u(\mathbf x,t) = \frac{\partial}{\partial t}\Big(t\,M_{at}[\varphi](\mathbf x)\Big) + t\,M_{at}[\psi](\mathbf x).
}\end{equation}

\end{theorem}

La demostración usa la idea de las ondas esféricas: el promedio sobre
esferas de cualquier solución cumple, como función del radio, la
ecuación radial de más arriba (lema de Darboux), y se aplica lo que ya
sabemos. Se puede ver completa en el capítulo 9 de Strauss o en la
sección 2.4 de Evans. Lo que importa es \textbf{qué dice}.

Para calcular \(u\) en el punto \(\mathbf x\) en el instante \(t\) hay
que promediar los datos \textbf{sobre la superficie de la esfera} de
radio \(at\) centrada en \(\mathbf x\): no sobre la bola entera, solo
sobre la cáscara. Así que si los datos son nulos fuera de una región
chica alrededor de un punto \(\mathbf x_0\), en un punto \(\mathbf x\) a
distancia \(d\) la solución es cero salvo cuando la esfera pasa por la
región, es decir cuando \(at\approx d\). \textbf{La perturbación llega,
pasa y se va, y detrás no queda nada.} Eso se llama \textbf{principio de
Huygens} en sentido fuerte.

En dos dimensiones no es así. La fórmula de Poisson, que se obtiene de
la de Kirchhoff tomando datos que no dependen de \(z\) ---el
\emph{método de descenso}---, promedia los datos sobre \textbf{todo el
disco} de radio \(at\), con un peso que se hace infinito en el borde. El
frente llega nítido, pero cuando pasa, la región sigue adentro del
disco, y la solución no vuelve a cero: queda una \textbf{estela} que
decae despacio. En una dimensión la integral de d'Alembert también es
sobre todo el intervalo, y la estela no decae nada: es la meseta.

La manera más limpia de verlo es comparar las funciones de Green, la
respuesta a un golpe puntual en el origen en \(t=0\):

\begin{longtable}[]{@{}
  >{\centering\arraybackslash}p{(\linewidth - 6\tabcolsep) * \real{0.1000}}
  >{\raggedright\arraybackslash}p{(\linewidth - 6\tabcolsep) * \real{0.4000}}
  >{\raggedright\arraybackslash}p{(\linewidth - 6\tabcolsep) * \real{0.2400}}
  >{\raggedright\arraybackslash}p{(\linewidth - 6\tabcolsep) * \real{0.2600}}@{}}
\caption{Las funciones de Green de la ecuación de onda en el espacio
libre. La de una dimensión es
\ref{eq-hip-green-1d}.}\label{tbl-hip-green-dimensiones}\tabularnewline
\toprule\noalign{}
\begin{minipage}[b]{\linewidth}\centering
Dimensión
\end{minipage} & \begin{minipage}[b]{\linewidth}\raggedright
\(G(\mathbf x, t)\) para
\(u_{tt} - a^2\nabla^2u = \delta(\mathbf x)\,\delta(t)\)
\end{minipage} & \begin{minipage}[b]{\linewidth}\raggedright
Dónde vive
\end{minipage} & \begin{minipage}[b]{\linewidth}\raggedright
Qué se oye
\end{minipage} \\
\midrule\noalign{}
\endfirsthead
\toprule\noalign{}
\begin{minipage}[b]{\linewidth}\centering
Dimensión
\end{minipage} & \begin{minipage}[b]{\linewidth}\raggedright
\(G(\mathbf x, t)\) para
\(u_{tt} - a^2\nabla^2u = \delta(\mathbf x)\,\delta(t)\)
\end{minipage} & \begin{minipage}[b]{\linewidth}\raggedright
Dónde vive
\end{minipage} & \begin{minipage}[b]{\linewidth}\raggedright
Qué se oye
\end{minipage} \\
\midrule\noalign{}
\endhead
\bottomrule\noalign{}
\endlastfoot
1 & \(\dfrac{1}{2a}\,H(at - |x|)\) & todo el interior del cono & un
escalón que no se apaga \\
2 & \(\dfrac{1}{2\pi a}\,\dfrac{H(at-r)}{\sqrt{a^2t^2-r^2}}\) & todo el
interior del cono & un pico y una cola larga \\
3 & \(\dfrac{1}{4\pi a^2}\,\dfrac{\delta(t - r/a)}{r}\) & \textbf{solo
la superficie} del cono & un chasquido limpio \\
\end{longtable}

\begin{tcolorbox}[enhanced jigsaw, arc=.35mm, bottomrule=.15mm, breakable, colback=white, colframe=quarto-callout-important-color-frame, left=2mm, leftrule=.75mm, opacityback=0, rightrule=.15mm, toprule=.15mm]
\begin{minipage}[t]{5.5mm}
\textcolor{quarto-callout-important-color}{\faExclamation}
\end{minipage}%
\begin{minipage}[t]{\textwidth - 5.5mm}

\vspace{-3mm}\textbf{Por qué vivimos en un mundo donde se puede hablar}\vspace{3mm}

Si el espacio tuviera dos dimensiones, cada palabra dejaría una cola que
se mezclaría con la siguiente, y hablar sería como gritar en una
catedral que no deja de resonar. En tres dimensiones cada sonido llega,
pasa, y deja el aire limpio para el siguiente. El principio de Huygens
vale en todas las dimensiones \textbf{impares} a partir de tres, y falla
en todas las pares.

\end{minipage}%
\end{tcolorbox}

\emph{(En la versión web hay acá un gráfico interactivo con seis
paneles: la respuesta a un golpe corto en una, dos y tres dimensiones,
como foto en un instante y como señal registrada a una distancia fija.
En una dimensión queda una meseta; en dos, un pico con una cola larga;
en tres, solo una copia del golpe.)}

\begin{exercise}[]\protect\hypertarget{exr-hip-globo}{}\label{exr-hip-globo}

Deducir \ref{eq-hip-centro} a partir de \ref{eq-hip-cauchy-esferico} y
la regla para \(r<at\). Para el globo de la Figura~\ref{fig-hip-globo}
con borde abrupto ---\(\varphi = 1\) para \(r<b\) y \(0\) afuera,
\(\psi = 0\)---, decir qué registra un micrófono puesto \textbf{en el
centro}.

\end{exercise}

\begin{tcolorbox}[enhanced jigsaw, arc=.35mm, bottomrule=.15mm, bottomtitle=1mm, breakable, colback=white, colbacktitle=quarto-callout-tip-color!10!white, colframe=quarto-callout-tip-color-frame, coltitle=black, left=2mm, leftrule=.75mm, opacityback=0, opacitybacktitle=0.6, rightrule=.15mm, title=\textcolor{quarto-callout-tip-color}{\faLightbulb}\hspace{0.5em}{Solución}, titlerule=0mm, toprule=.15mm, toptitle=1mm]

Para \(r<at\), con \(\psi\) y el reemplazo del eco, \[
u(r,t) = \frac{(r+at)\,\varphi(r+at) - (at-r)\,\varphi(at-r)}{2r} + \frac{1}{2ar}\int_{at-r}^{r+at}s\,\psi(s)\,ds .
\] Llamando \(V(s) = s\varphi(s)\), el primer término es
\(\dfrac{V(at+r) - V(at-r)}{2r}\), un cociente incremental centrado, que
tiende a \(V'(at) = \varphi(at) + at\,\varphi'(at)\). La integral es
sobre un intervalo de largo \(2r\) alrededor de \(at\), y dividida por
\(2ar\) tiende a \(\tfrac1a\,at\,\psi(at) = t\,\psi(at)\). Eso es
\ref{eq-hip-centro}.

\textbf{El globo.} Con \(\psi = 0\) y el escalón, \(\varphi(at) = 1\)
para \(t<b/a\) y \(0\) después, y \(\varphi'(s) = -\delta(s-b)\).
Entonces \[
u(0,t) = H\big(b - at\big) - at\,\delta(at - b) = H\big(b-at\big) - \frac{b}{a}\,\delta\!\Big(t - \frac ba\Big).
\] El micrófono del centro registra la sobrepresión del globo sin
cambios mientras el frente de rarefacción viaja hacia adentro, y en el
instante \(t = b/a\), cuando llega, registra una \textbf{delta
negativa}: toda la rarefacción que estaba repartida sobre la esfera de
radio \(b\) se concentra en un punto. Después, presión normal. Con un
borde suave, como en la figura, la delta es un pico alto y angosto.

\end{tcolorbox}

\section{Ondas electromagnéticas}\label{sec-hip-maxwell}

Esta es la sección 9.4 de González, y la más importante del capítulo por
sus consecuencias. En la intro dijimos que Maxwell combinó sus
ecuaciones y le salió una ecuación de onda con la velocidad de la luz.
Acá lo hacemos, y después usamos todo lo del capítulo ---Green, Huygens,
el cono de luz--- para responder la pregunta central del
electromagnetismo: \textbf{qué campo produce una carga que se mueve}.

Para esta sección hace falta cálculo vectorial: divergencia, rotor, la
identidad
\(\nabla\times(\nabla\times\mathbf F) = \nabla(\nabla\cdot\mathbf F) - \nabla^2\mathbf F\),
y el teorema de la divergencia.

\subsection{Las ecuaciones de Maxwell}\label{sec-hip-maxwell-ec}

Trabajamos en el sistema internacional (SI). El libro de la cátedra usa
el sistema gaussiano (cgs); la traducción está en la
Tabla~\ref{tbl-hip-si-cgs} al final de la sección. En el vacío, con una
densidad de carga \(\rho\) y una densidad de corriente \(\mathbf J\),
los campos eléctrico \(\mathbf E\) y magnético \(\mathbf B\) cumplen

\begin{equation}\protect\phantomsection\label{eq-hip-maxwell}{
\begin{aligned}
\nabla\cdot\mathbf E &= \frac{\rho}{\varepsilon_0} && \text{(Gauss)}, &
\nabla\times\mathbf E &= -\frac{\partial\mathbf B}{\partial t} && \text{(Faraday)},\\
\nabla\cdot\mathbf B &= 0 && \text{(no hay monopolos)}, &
\nabla\times\mathbf B &= \mu_0\mathbf J + \mu_0\varepsilon_0\frac{\partial\mathbf E}{\partial t} && \text{(Ampère–Maxwell)}.
\end{aligned}
}\end{equation}

El último término, \(\mu_0\varepsilon_0\,\partial_t\mathbf E\), es la
\textbf{corriente de desplazamiento}, y es el aporte de Maxwell: sin él
las ecuaciones son inconsistentes con la conservación de la carga, y con
él aparecen las ondas.

\begin{proposition}[Los campos en el vacío son
ondas]\protect\hypertarget{prp-hip-ondas-em}{}\label{prp-hip-ondas-em}

Donde no hay cargas ni corrientes, cada componente cartesiana de
\(\mathbf E\) y de \(\mathbf B\) cumple la ecuación de onda

\begin{equation}\protect\phantomsection\label{eq-hip-ondas-em}{
\nabla^2\mathbf E = \frac{1}{c^2}\,\frac{\partial^2\mathbf E}{\partial t^2},\qquad
\nabla^2\mathbf B = \frac{1}{c^2}\,\frac{\partial^2\mathbf B}{\partial t^2},\qquad
c = \frac{1}{\sqrt{\mu_0\varepsilon_0}} .
}\end{equation}

\end{proposition}

\begin{proof}
Tomando rotor en la ley de Faraday y usando la identidad del rotor del
rotor, \[
\nabla(\nabla\cdot\mathbf E) - \nabla^2\mathbf E = -\frac{\partial}{\partial t}(\nabla\times\mathbf B).
\] Sin cargas, \(\nabla\cdot\mathbf E = 0\); sin corrientes,
\(\nabla\times\mathbf B = \mu_0\varepsilon_0\,\partial_t\mathbf E\).
Queda
\(-\nabla^2\mathbf E = -\mu_0\varepsilon_0\,\partial_t^2\mathbf E\).
Para \(\mathbf B\) es la misma cuenta empezando por la ley de
Ampère--Maxwell.
\end{proof}

Con \(\mu_0 = 4\pi\times10^{-7}\ \text{N/A}^2\) y
\(\varepsilon_0 = 8{,}854\times10^{-12}\ \text{F/m}\) ---dos constantes
que se miden con imanes y capacitores, sin nada que ver con la luz---,
calculá \(c = 1/\sqrt{\mu_0\varepsilon_0}\). Respondé en unidades de
\(10^8\ \text{m/s}\), con tres decimales.

\(\mu_0\varepsilon_0 = 4\pi\times10^{-7}\cdot8{,}854\times10^{-12} \approx 1{,}1127\times10^{-17}\ \text{s}^2/\text{m}^2\),
y
\(c = 1/\sqrt{1{,}1127\times10^{-17}} \approx 2{,}998\times10^8\ \text{m/s}\).

Es el valor medido para la velocidad de la luz. Cuando Maxwell hizo esta
cuenta, con los valores de su época, la coincidencia fue la prueba de
que la luz es una onda electromagnética. Una constante de la óptica
salía de dos constantes de la electricidad.

\begin{exercise}[]\protect\hypertarget{exr-hip-onda-plana}{}\label{exr-hip-onda-plana}

Mostrar que \(\mathbf E = E_0\cos(kx - \omega t)\,\hat{\mathbf y}\)
cumple \ref{eq-hip-ondas-em} si y solo si \(\omega = ck\). Hallar el
\(\mathbf B\) que la acompaña, usando la ley de Faraday, y comprobar que
es perpendicular a \(\mathbf E\) y a la dirección de propagación, y que
\(|\mathbf B| = |\mathbf E|/c\).

\end{exercise}

\begin{tcolorbox}[enhanced jigsaw, arc=.35mm, bottomrule=.15mm, bottomtitle=1mm, breakable, colback=white, colbacktitle=quarto-callout-tip-color!10!white, colframe=quarto-callout-tip-color-frame, coltitle=black, left=2mm, leftrule=.75mm, opacityback=0, opacitybacktitle=0.6, rightrule=.15mm, title=\textcolor{quarto-callout-tip-color}{\faLightbulb}\hspace{0.5em}{Solución}, titlerule=0mm, toprule=.15mm, toptitle=1mm]

La única componente es \(E_y = E_0\cos(kx-\omega t)\), con
\(\nabla^2E_y = -k^2E_y\) y \(\partial_t^2E_y = -\omega^2E_y\). La
ecuación de onda pide \(-k^2 = -\omega^2/c^2\), o sea \(\omega = ck\)
(con \(\omega,k>0\)). Es una onda que viaja hacia \(+x\) a velocidad
\(\omega/k = c\). Además \(\nabla\cdot\mathbf E = \partial_yE_y = 0\):
cumple Gauss sin cargas.

Por Faraday, \(\partial_t\mathbf B = -\nabla\times\mathbf E\). El rotor
de \(E_y(x,t)\,\hat{\mathbf y}\) es
\(\partial_xE_y\,\hat{\mathbf z} = -kE_0\operatorname{sen}(kx-\omega t)\,\hat{\mathbf z}\),
así que \[
\partial_t B_z = kE_0\operatorname{sen}(kx-\omega t)\quad\Longrightarrow\quad
B_z = \frac{k}{\omega}E_0\cos(kx-\omega t) = \frac{E_0}{c}\cos(kx-\omega t),
\] descartando la constante de integración, que sería un campo estático
agregado. Entonces
\(\mathbf B = (E_0/c)\cos(kx-\omega t)\,\hat{\mathbf z}\): perpendicular
a \(\mathbf E\) (que apunta en \(\hat{\mathbf y}\)), perpendicular a la
dirección de propagación (\(\hat{\mathbf x}\)), en fase con
\(\mathbf E\), y con \(|\mathbf B| = |\mathbf E|/c\). La onda
electromagnética es \textbf{transversal}, como la de la cuerda y a
diferencia del sonido.

\end{tcolorbox}

\subsection{Los potenciales}\label{sec-hip-potenciales}

Con cargas y corrientes el cálculo directo con los campos se complica:
son seis componentes acopladas. Los \textbf{potenciales} reducen la
cantidad de incógnitas a cambio de subir el orden de las ecuaciones.

Como \(\nabla\cdot\mathbf B = 0\), existe un campo vectorial
\(\mathbf A\) con \(\mathbf B = \nabla\times\mathbf A\). Poniéndolo en
Faraday, \(\nabla\times(\mathbf E + \partial_t\mathbf A) = 0\), así que
existe un escalar \(\Phi\) con
\(\mathbf E + \partial_t\mathbf A = -\nabla\Phi\).

\begin{definition}[Potenciales
electromagnéticos]\protect\hypertarget{def-hip-potenciales}{}\label{def-hip-potenciales}

El \textbf{potencial vector} \(\mathbf A\) y el \textbf{potencial
escalar} \(\Phi\) son campos tales que

\begin{equation}\protect\phantomsection\label{eq-hip-potenciales}{
\mathbf B = \nabla\times\mathbf A, \qquad \mathbf E = -\nabla\Phi - \frac{\partial\mathbf A}{\partial t}.
}\end{equation}

\end{definition}

Con esta definición, Faraday y la ausencia de monopolos se cumplen
solas. Las otras dos ecuaciones de Maxwell se convierten en

\begin{equation}\protect\phantomsection\label{eq-hip-potenciales-ec}{
\begin{aligned}
\nabla^2\Phi + \frac{\partial}{\partial t}(\nabla\cdot\mathbf A) &= -\frac{\rho}{\varepsilon_0},\\
\nabla^2\mathbf A - \frac{1}{c^2}\frac{\partial^2\mathbf A}{\partial t^2}
- \nabla\Big(\nabla\cdot\mathbf A + \frac{1}{c^2}\frac{\partial\Phi}{\partial t}\Big) &= -\mu_0\mathbf J .
\end{aligned}
}\end{equation}

Están acopladas y son feas. Pero los potenciales no están unívocamente
determinados, y esa libertad es la que las arregla.

\begin{definition}[Transformación de gauge (de
medida)]\protect\hypertarget{def-hip-gauge}{}\label{def-hip-gauge}

Para cualquier función escalar \(\Lambda(\mathbf x,t)\), los potenciales
\[
\mathbf A' = \mathbf A + \nabla\Lambda, \qquad \Phi' = \Phi - \frac{\partial\Lambda}{\partial t}
\] producen los mismos \(\mathbf E\) y \(\mathbf B\). Esta
indeterminación se llama \textbf{libertad de gauge} o \textbf{de
medida}, y se fija imponiendo una condición adicional. La
\textbf{condición de Lorenz} es

\begin{equation}\protect\phantomsection\label{eq-hip-lorenz}{
\nabla\cdot\mathbf A + \frac{1}{c^2}\frac{\partial\Phi}{\partial t} = 0 .
}\end{equation}

\end{definition}

Que \(\mathbf E\) y \(\mathbf B\) no cambian se verifica en una línea:
\(\nabla\times\nabla\Lambda = 0\), y en \(\mathbf E\) los dos términos
nuevos, \(+\nabla\partial_t\Lambda\) y \(-\partial_t\nabla\Lambda\), se
cancelan. (La condición suele llamarse «de Lorentz», como en González;
se debe a Ludvig Lorenz, sin t, y no a Hendrik Lorentz, el de las
transformaciones.)

Con \ref{eq-hip-lorenz} el paréntesis de la segunda ecuación se anula y
la primera se reescribe usando
\(\nabla\cdot\mathbf A = -\partial_t\Phi/c^2\). Queda lo que buscábamos:
\textbf{cuatro ecuaciones de onda desacopladas, una para cada
potencial},

\begin{equation}\protect\phantomsection\label{eq-hip-potenciales-onda}{
\nabla^2\Phi - \frac{1}{c^2}\frac{\partial^2\Phi}{\partial t^2} = -\frac{\rho}{\varepsilon_0},
\qquad
\nabla^2\mathbf A - \frac{1}{c^2}\frac{\partial^2\mathbf A}{\partial t^2} = -\mu_0\mathbf J .
}\end{equation}

Son \ref{eq-hip-onda-3d} con \(a = c\) y las fuentes puestas por las
cargas y las corrientes. Todo lo que sabemos de la ecuación de onda se
aplica ahora al electromagnetismo.

\begin{tcolorbox}[enhanced jigsaw, arc=.35mm, bottomrule=.15mm, breakable, colback=white, colframe=quarto-callout-warning-color-frame, left=2mm, leftrule=.75mm, opacityback=0, rightrule=.15mm, toprule=.15mm]
\begin{minipage}[t]{5.5mm}
\textcolor{quarto-callout-warning-color}{\faExclamationTriangle}
\end{minipage}%
\begin{minipage}[t]{\textwidth - 5.5mm}

\vspace{-3mm}\textbf{Ojo con el libro de la cátedra (secciones 9.4.1 y 9.4.2)}\vspace{3mm}

\begin{itemize}
\tightlist
\item
  Presenta la ley de Faraday (9.30) como «la ecuación de Faraday con la
  corriente de desplazamiento». La corriente de desplazamiento está en
  la ley de Ampère--Maxwell (9.28), no en la de Faraday.
\item
  En (9.33) el signo del segundo término está cambiado: en cgs es
  \(\nabla^2\Phi + \tfrac1c\,\partial_t(\nabla\cdot\mathbf A) = -4\pi\rho\).
  Sale de poner
  \(\mathbf E = -\nabla\Phi - \tfrac1c\partial_t\mathbf A\) en
  \(\nabla\cdot\mathbf E = 4\pi\rho\).
\item
  En (9.36) falta un signo menos: es
  \(\nabla^2\mathbf A - \tfrac1{c^2}\partial_t^2\mathbf A = -\tfrac{4\pi}{c}\,\mathbf j\),
  consistente con (9.34), que está bien.
\item
  Dice «transformando Fourier a (9.36)» pero transforma (9.37), la del
  potencial escalar. Y en (9.41) aparece \(r\) donde va \(R\).
\end{itemize}

\end{minipage}%
\end{tcolorbox}

\begin{exercise}[]\protect\hypertarget{exr-hip-lorenz-existe}{}\label{exr-hip-lorenz-existe}

Mostrar que siempre se puede elegir el gauge de Lorenz: dados
\(\mathbf A\) y \(\Phi\) cualesquiera, hallar qué ecuación tiene que
cumplir \(\Lambda\) para que \(\mathbf A'\) y \(\Phi'\) cumplan
\ref{eq-hip-lorenz}. ¿Queda alguna libertad después?

\end{exercise}

\begin{tcolorbox}[enhanced jigsaw, arc=.35mm, bottomrule=.15mm, bottomtitle=1mm, breakable, colback=white, colbacktitle=quarto-callout-tip-color!10!white, colframe=quarto-callout-tip-color-frame, coltitle=black, left=2mm, leftrule=.75mm, opacityback=0, opacitybacktitle=0.6, rightrule=.15mm, title=\textcolor{quarto-callout-tip-color}{\faLightbulb}\hspace{0.5em}{Solución}, titlerule=0mm, toprule=.15mm, toptitle=1mm]

Hay que pedir \[
\nabla\cdot\mathbf A' + \frac{1}{c^2}\partial_t\Phi' = \nabla\cdot\mathbf A + \nabla^2\Lambda + \frac{1}{c^2}\partial_t\Phi - \frac{1}{c^2}\partial_t^2\Lambda = 0,
\] es decir \[
\nabla^2\Lambda - \frac{1}{c^2}\,\partial_t^2\Lambda = -\Big(\nabla\cdot\mathbf A + \frac{1}{c^2}\partial_t\Phi\Big).
\] Es una \textbf{ecuación de onda no homogénea} para \(\Lambda\), con
fuente conocida. Siempre tiene solución: la da la función de Green, que
es lo próximo que vamos a construir. Así que el gauge de Lorenz siempre
se puede alcanzar.

La libertad no se agota: si \(\Lambda_0\) cumple la ecuación de onda
\textbf{homogénea}, se le puede sumar a \(\Lambda\) sin romper la
condición. Esa libertad residual es la que en electrodinámica permite,
por ejemplo, anular \(\Phi\) en regiones sin cargas.

\end{tcolorbox}

\subsection{Los potenciales retardados}\label{sec-hip-retardados}

Ya tenemos cuatro ecuaciones de onda con fuente, sin bordes. Las
resolvemos con la función de Green, en el espíritu del
sección~\ref{sec-green-varias}: transformar Fourier en el tiempo,
resolver Helmholtz, volver. Es la cuenta de la sección 9.4.2 de
González, completada.

\textbf{Transformar.} Con la convención de González,
\(\hat\Phi(\mathbf x,\omega) = \int\Phi(\mathbf x,t)\,e^{i\omega t}\,dt\),
la derivada \(\partial_t^2\) se vuelve una multiplicación por
\(-\omega^2\), y la ecuación de \(\Phi\) queda

\[
\big(\nabla^2 + k^2\big)\hat\Phi = -\frac{\hat\rho}{\varepsilon_0}, \qquad k = \frac{\omega}{c}.
\]

Es la ecuación de Helmholtz del sección~\ref{sec-helmholtz}, ahora en
tres dimensiones y con una fuente.

\textbf{Green de Helmholtz.} Buscamos \(G\) con
\((\nabla^2+k^2)G = -\delta(\mathbf x-\mathbf x')\). Sin bordes no hay
dirección privilegiada, así que \(G\) depende solo de
\(R = |\mathbf x-\mathbf x'|\), y con el laplaciano radial del
sección~\ref{sec-hip-esfericas}, para \(R>0\), \[
\frac1R\frac{d^2}{dR^2}(RG) + k^2G = 0 \quad\Longrightarrow\quad RG = Ae^{ikR} + Be^{-ikR}.
\] Cerca de \(R = 0\) las dos se comportan como \((A+B)/R\), y para que
la delta tenga peso \(1\) hace falta que \(G\approx1/(4\pi R)\), la de
Coulomb. Así que \(A + B = 1/(4\pi)\):

\begin{equation}\protect\phantomsection\label{eq-hip-green-helmholtz}{
G = \frac{A\,e^{ikR} + B\,e^{-ikR}}{R}, \qquad A+B = \frac{1}{4\pi}.
}\end{equation}

\textbf{Volver al tiempo.} Cada término se antitransforma solo. Como
\(\tfrac{1}{2\pi}\int e^{i\omega R/c}\,e^{-i\omega\tau}\,d\omega = \delta(\tau - R/c)\),
el término con \(e^{ikR}\) da una delta en \(\tau = R/c\) y el término
con \(e^{-ikR}\) una en \(\tau = -R/c\):

\[
G(R,\tau) = 4\pi A\,\frac{\delta(\tau - R/c)}{4\pi R} + 4\pi B\,\frac{\delta(\tau + R/c)}{4\pi R},
\]

donde \(\tau\) es el tiempo transcurrido desde la fuente. González se
detiene acá, con \(A\) y \(B\) libres. Pero la física los fija.

\begin{tcolorbox}[enhanced jigsaw, arc=.35mm, bottomrule=.15mm, breakable, colback=white, colframe=quarto-callout-important-color-frame, left=2mm, leftrule=.75mm, opacityback=0, rightrule=.15mm, toprule=.15mm]
\begin{minipage}[t]{5.5mm}
\textcolor{quarto-callout-important-color}{\faExclamation}
\end{minipage}%
\begin{minipage}[t]{\textwidth - 5.5mm}

\vspace{-3mm}\textbf{La causalidad elige}\vspace{3mm}

El término con \(\delta(\tau + R/c)\) dice que el campo en \(\mathbf x\)
responde a la fuente \textbf{antes} de que la fuente haga algo, un
tiempo \(R/c\) antes. Matemáticamente es tan buena solución como la otra
---la ecuación de onda es simétrica ante \(t\to-t\)---, pero no describe
nada que se haya observado. La condición física es que \textbf{el efecto
no preceda a la causa}, y eso pone \(B = 0\), \(A = 1/(4\pi)\). La
solución que queda se llama \textbf{retardada}; la descartada,
\textbf{avanzada}.

\end{minipage}%
\end{tcolorbox}

\begin{theorem}[Potenciales
retardados]\protect\hypertarget{thm-hip-retardados}{}\label{thm-hip-retardados}

En el gauge de Lorenz y sin bordes, la solución causal de
\ref{eq-hip-potenciales-onda} es

\begin{equation}\protect\phantomsection\label{eq-hip-retardados}{
\Phi(\mathbf x,t) = \frac{1}{4\pi\varepsilon_0}\int\frac{\rho\big(\mathbf x',\,t_r\big)}{|\mathbf x-\mathbf x'|}\,dV',
\qquad
\mathbf A(\mathbf x,t) = \frac{\mu_0}{4\pi}\int\frac{\mathbf J\big(\mathbf x',\,t_r\big)}{|\mathbf x-\mathbf x'|}\,dV',
}\end{equation}

donde

\[
t_r = t - \frac{|\mathbf x - \mathbf x'|}{c}
\]

es el \textbf{tiempo retardado}: el instante en que tuvo que salir de
\(\mathbf x'\) una señal a velocidad \(c\) para llegar a \(\mathbf x\)
en el instante \(t\).

\end{theorem}

\begin{proof}
La función de Green en el tiempo es
\(G = \delta(t - t' - R/c)/(4\pi R)\), que cumple
\(\nabla^2G - \tfrac1{c^2}\partial_t^2G = -\delta(\mathbf x-\mathbf x')\,\delta(t-t')\).
Por superposición, \[
\Phi(\mathbf x,t) = \int\!\!\int G\,\frac{\rho(\mathbf x',t')}{\varepsilon_0}\,dt'\,dV'
= \frac{1}{4\pi\varepsilon_0}\int\frac{1}{R}\left[\int\delta\Big(t-t'-\frac Rc\Big)\rho(\mathbf x',t')\,dt'\right]dV',
\] y la integral en \(t'\) evalúa \(\rho\) en \(t' = t - R/c\). Lo mismo
para cada componente de \(\mathbf A\).
\end{proof}

Hay que mirar la fórmula con cuidado, porque se parece mucho a algo que
ya sabemos. Si la carga no depende del tiempo,
\(\rho(\mathbf x', t_r) = \rho(\mathbf x')\), y \ref{eq-hip-retardados}
es \textbf{la ley de Coulomb}. La única diferencia, y es toda la
diferencia, es que cada pedazo de carga contribuye con el valor que
tenía \textbf{cuando salió la señal}, no con el que tiene ahora. Es la
Coulomb de siempre escrita sobre la superficie del cono de luz del
pasado: la función de Green de la Tabla~\ref{tbl-hip-green-dimensiones}
en tres dimensiones, que vive solo sobre el cono.

\textbf{Mirar lejos es mirar al pasado.} El Sol que vemos es el de hace
ocho minutos; la estrella más cercana, la de hace cuatro años; la
radiación de fondo de microondas, el universo de hace casi catorce mil
millones de años. Todos son potenciales retardados.

\emph{(En la versión web hay acá un gráfico interactivo: una carga que
oscila sobre una recta y un observador quieto. En el plano \((x,t)\) se
ven la línea de universo de la carga, el cono de luz hacia el pasado del
observador y el punto donde se cortan, que es el instante retardado.
Debajo, la posición verdadera de la carga y la posición que ve el
observador, que es una copia atrasada y deformada por el efecto
Doppler.)}

La distancia media de la Tierra al Sol es \(1{,}496\times10^{11}\) m.
¿Cuántos segundos de atraso tiene la imagen del Sol que vemos?

\(t - t_r = R/c = 1{,}496\times10^{11}/2{,}998\times10^8 \approx 499\)
s, es decir 8 minutos y 19 segundos. Si el Sol desapareciera, lo
seguiríamos viendo ---y la Tierra seguiría girando a su alrededor,
porque la gravedad también se propaga a la velocidad de la luz---
durante ese tiempo.

\subsection{Fórmulas de Green y fórmula de
Kirchhoff}\label{sec-hip-kirchhoff}

La función de Green del espacio libre resuelve el problema sin bordes.
Con bordes ---una pantalla con un agujero, una antena dentro de una
caja--- hace falta una herramienta que relacione lo que pasa adentro de
una región con lo que pasa en su frontera. Son las \textbf{identidades
de Green}, que González presenta en 9.4.3.

\begin{theorem}[Identidades de
Green]\protect\hypertarget{thm-hip-identidades-green}{}\label{thm-hip-identidades-green}

Sea \(V\) una región acotada con frontera \(S\) regular y normal
exterior \(\mathbf n\), y sean \(\phi,\psi\) funciones de clase \(C^2\)
en \(V\) y \(C^1\) hasta el borde. Entonces

\begin{equation}\protect\phantomsection\label{eq-hip-green-1}{
\int_V\big(\phi\,\nabla^2\psi + \nabla\phi\cdot\nabla\psi\big)\,dV = \oint_S\phi\,\frac{\partial\psi}{\partial n}\,dS
\qquad\text{(primera identidad)},
}\end{equation}

\begin{equation}\protect\phantomsection\label{eq-hip-green-2}{
\int_V\big(\phi\,\nabla^2\psi - \psi\,\nabla^2\phi\big)\,dV = \oint_S\Big(\phi\,\frac{\partial\psi}{\partial n} - \psi\,\frac{\partial\phi}{\partial n}\Big)\,dS
\qquad\text{(segunda identidad)}.
}\end{equation}

\end{theorem}

\begin{proof}
La primera es el teorema de la divergencia aplicado al campo
\(\phi\,\nabla\psi\), cuya divergencia es
\(\nabla\phi\cdot\nabla\psi + \phi\,\nabla^2\psi\) y cuyo flujo es
\(\phi\,\partial_n\psi\). La segunda sale de escribir la primera con
\(\phi\) y \(\psi\) intercambiadas y restar: el término
\(\nabla\phi\cdot\nabla\psi\) es simétrico y se cancela.
\end{proof}

La primera identidad es el teorema fundamental del cálculo en varias
variables, y es la que usamos, en una dimensión, para el balance de
energía del sección~\ref{sec-hip-balance}. La segunda es la que dice
algo nuevo: el miembro izquierdo es simétrico en la forma en que
aparecen \(\phi\) y \(\psi\), y si las dos cumplen ecuaciones conocidas,
el volumen se puede cambiar por la frontera.

Eso es lo que hace la \textbf{fórmula de Kirchhoff}. Tomamos una
solución \(\psi\) de la ecuación de Helmholtz sin fuentes adentro de
\(V\), y en la segunda identidad ponemos \(\phi = G\), la función de
Green del espacio libre con el polo en un punto \(\mathbf x\) interior.

\begin{theorem}[Fórmula integral de
Kirchhoff]\protect\hypertarget{thm-hip-kirchhoff}{}\label{thm-hip-kirchhoff}

Si \((\nabla^2+k^2)\psi = 0\) en \(V\) y
\(G(\mathbf x,\mathbf x') = \dfrac{e^{ik|\mathbf x-\mathbf x'|}}{4\pi|\mathbf x-\mathbf x'|}\),
entonces para todo \(\mathbf x\) en el interior de \(V\)

\begin{equation}\protect\phantomsection\label{eq-hip-kirchhoff}{
\psi(\mathbf x) = \oint_S\left[G(\mathbf x,\mathbf x')\,\frac{\partial\psi}{\partial n'}(\mathbf x') - \psi(\mathbf x')\,\frac{\partial G}{\partial n'}(\mathbf x,\mathbf x')\right]dS',
}\end{equation}

con \(\mathbf n'\) la normal \textbf{exterior} a \(V\).

\end{theorem}

\begin{proof}
En la segunda identidad, con \(\phi = G\) y las variables de integración
\(\mathbf x'\), se usa \(\nabla'^2\psi = -k^2\psi\) y
\(\nabla'^2G = -k^2G - \delta(\mathbf x-\mathbf x')\). El integrando del
volumen es
\(G(-k^2\psi) - \psi(-k^2G - \delta) = \psi\,\delta(\mathbf x-\mathbf x')\),
cuya integral es \(\psi(\mathbf x)\).

Para hacer esto sin deltas, se saca de \(V\) una bolita de radio
\(\varepsilon\) alrededor de \(\mathbf x\), donde \(G\) es suave y el
volumen da cero, y se calcula el flujo por la esferita: con
\(G\approx1/(4\pi\varepsilon)\) y
\(\partial G/\partial n'\approx+1/(4\pi\varepsilon^2)\) (la normal
exterior a la región que queda apunta hacia \(\mathbf x\)), el término
\(-\psi\,\partial_{n'}G\) integrado sobre un área \(4\pi\varepsilon^2\)
tiende a \(-\psi(\mathbf x)\) y el término con \(G\) tiende a cero.
Pasando al otro miembro sale \ref{eq-hip-kirchhoff}.
\end{proof}

Lo que dice la fórmula es el principio de Huygens escrito con precisión:
\textbf{el campo adentro está determinado por el campo y su derivada
normal sobre el borde}, y cada punto del borde actúa como una fuente
secundaria que emite ondas esféricas \(G\). Es la formalización de la
construcción de Huygens de 1678, con las «ondículas» que salen de cada
punto del frente.

\begin{tcolorbox}[enhanced jigsaw, arc=.35mm, bottomrule=.15mm, breakable, colback=white, colframe=quarto-callout-warning-color-frame, left=2mm, leftrule=.75mm, opacityback=0, rightrule=.15mm, toprule=.15mm]
\begin{minipage}[t]{5.5mm}
\textcolor{quarto-callout-warning-color}{\faExclamationTriangle}
\end{minipage}%
\begin{minipage}[t]{\textwidth - 5.5mm}

\vspace{-3mm}\textbf{Ojo con el libro de la cátedra (sección 9.4.4) y las convenciones}\vspace{3mm}

En (9.49), González escribe
\(\psi(\mathbf x) = \oint\big[\psi\,\mathbf n'\cdot\nabla'G - G\,\mathbf n'\cdot\nabla'\psi\big]dS'\),
que tiene el signo opuesto a \ref{eq-hip-kirchhoff} y no tiene ningún
factor \(4\pi\). Con \textbf{su} \(G\), la de
\((\nabla^2+k^2)\bar G = -4\pi\delta\), que es \(e^{ikR}/R\) sin el
\(4\pi\) abajo, la cuenta de la demostración de arriba da un factor
\(\tfrac{1}{4\pi}\) delante de la integral. El signo de (9.49) es el
correcto si la normal apunta \textbf{hacia adentro} de \(V\), que es la
convención de varios libros de óptica. Es decir: el signo puede ser una
convención, pero el \(\tfrac{1}{4\pi}\) falta.

\end{minipage}%
\end{tcolorbox}

\begin{longtable}[]{@{}
  >{\raggedright\arraybackslash}p{(\linewidth - 6\tabcolsep) * \real{0.2200}}
  >{\raggedright\arraybackslash}p{(\linewidth - 6\tabcolsep) * \real{0.2400}}
  >{\raggedright\arraybackslash}p{(\linewidth - 6\tabcolsep) * \real{0.2600}}
  >{\raggedright\arraybackslash}p{(\linewidth - 6\tabcolsep) * \real{0.2800}}@{}}
\caption{Las tres convenciones que conviven. Todas describen la misma
física; cambian el signo de la delta y dónde se pone el
\(4\pi\).}\label{tbl-hip-convenciones-green}\tabularnewline
\toprule\noalign{}
\begin{minipage}[b]{\linewidth}\raggedright
\end{minipage} & \begin{minipage}[b]{\linewidth}\raggedright
Capítulo 8
\end{minipage} & \begin{minipage}[b]{\linewidth}\raggedright
Este capítulo (SI)
\end{minipage} & \begin{minipage}[b]{\linewidth}\raggedright
González (cgs)
\end{minipage} \\
\midrule\noalign{}
\endfirsthead
\toprule\noalign{}
\begin{minipage}[b]{\linewidth}\raggedright
\end{minipage} & \begin{minipage}[b]{\linewidth}\raggedright
Capítulo 8
\end{minipage} & \begin{minipage}[b]{\linewidth}\raggedright
Este capítulo (SI)
\end{minipage} & \begin{minipage}[b]{\linewidth}\raggedright
González (cgs)
\end{minipage} \\
\midrule\noalign{}
\endhead
\bottomrule\noalign{}
\endlastfoot
Ecuación de Green de Helmholtz & \((\nabla^2+k^2)G = \delta\) &
\((\nabla^2+k^2)G = -\delta\) &
\((\nabla^2+k^2)\bar G = -4\pi\delta\) \\
\(G\) saliente & \(-\dfrac{e^{ikR}}{4\pi R}\) &
\(\dfrac{e^{ikR}}{4\pi R}\) & \(\dfrac{e^{ikR}}{R}\) \\
Green retardada & --- & \(\dfrac{\delta(t-t'-R/c)}{4\pi R}\) &
\(\dfrac{\delta(\tau-R/c)}{R}\) \\
Kirchhoff (normal exterior) &
\(\psi = \oint(\psi\,\partial_nG - G\,\partial_n\psi)\) &
\(\psi = \oint(G\,\partial_n\psi - \psi\,\partial_nG)\) &
\(\psi = \tfrac{1}{4\pi}\oint(\bar G\,\partial_n\psi - \psi\,\partial_n\bar G)\) \\
\end{longtable}

\subsection{Difracción}\label{sec-hip-difraccion}

La fórmula de Kirchhoff tiene una virtud y un defecto, y González los
señala a los dos.

\textbf{La virtud}: si se conocen \(\psi\) y \(\partial_n\psi\) sobre el
borde, se conoce \(\psi\) en todas partes. \textbf{El defecto}: conocer
las dos cosas sobre el borde es, en general, haber resuelto el problema.
Para una pantalla opaca con un agujero, Kirchhoff propuso una
\textbf{aproximación}: sobre el agujero, \(\psi\) y \(\partial_n\psi\)
son los de la onda incidente, como si la pantalla no estuviera; sobre la
pantalla, las dos son cero.

La aproximación funciona notablemente bien cuando el agujero es grande
comparado con la longitud de onda. Pero es \textbf{inconsistente}: una
solución de Helmholtz con \(\psi = 0\) y \(\partial_n\psi=0\) sobre un
pedazo de superficie es idénticamente nula ---es un problema de Cauchy
para una elíptica, el ejemplo de Hadamard otra vez---, así que la
fórmula no puede reproducir los valores que se le pusieron en el borde.
La manera correcta es usar una función de Green que ya cumpla la
condición de la pantalla, construida con el método de las imágenes
(capítulo 11), que da las fórmulas de Rayleigh--Sommerfeld.

\emph{(En la versión web hay acá un gráfico interactivo: la intensidad
detrás de una rendija calculada sumando 400 ondículas de Huygens,
comparada con la fórmula de Fraunhofer. Se pueden mover el ancho de la
rendija y la distancia a la pantalla, en longitudes de onda.)}

\begin{tcolorbox}[enhanced jigsaw, arc=.35mm, bottomrule=.15mm, breakable, colback=white, colframe=quarto-callout-note-color-frame, left=2mm, leftrule=.75mm, opacityback=0, rightrule=.15mm, toprule=.15mm]
\begin{minipage}[t]{5.5mm}
\textcolor{quarto-callout-note-color}{\faInfo}
\end{minipage}%
\begin{minipage}[t]{\textwidth - 5.5mm}

\vspace{-3mm}\textbf{La tabla de traducción SI ↔ cgs}\vspace{3mm}

\begin{longtable}[]{@{}
  >{\raggedright\arraybackslash}p{(\linewidth - 4\tabcolsep) * \real{0.2800}}
  >{\raggedright\arraybackslash}p{(\linewidth - 4\tabcolsep) * \real{0.3600}}
  >{\raggedright\arraybackslash}p{(\linewidth - 4\tabcolsep) * \real{0.3600}}@{}}
\caption{Con \(\square = \nabla^2 - \tfrac{1}{c^2}\partial_t^2\). En el
vacío de González \(\mathbf D = \mathbf E\) y
\(\mathbf H = \mathbf B\).}\label{tbl-hip-si-cgs}\tabularnewline
\toprule\noalign{}
\begin{minipage}[b]{\linewidth}\raggedright
\end{minipage} & \begin{minipage}[b]{\linewidth}\raggedright
SI (este libro)
\end{minipage} & \begin{minipage}[b]{\linewidth}\raggedright
cgs gaussiano (González)
\end{minipage} \\
\midrule\noalign{}
\endfirsthead
\toprule\noalign{}
\begin{minipage}[b]{\linewidth}\raggedright
\end{minipage} & \begin{minipage}[b]{\linewidth}\raggedright
SI (este libro)
\end{minipage} & \begin{minipage}[b]{\linewidth}\raggedright
cgs gaussiano (González)
\end{minipage} \\
\midrule\noalign{}
\endhead
\bottomrule\noalign{}
\endlastfoot
Gauss & \(\nabla\cdot\mathbf E = \rho/\varepsilon_0\) &
\(\nabla\cdot\mathbf E = 4\pi\rho\) \\
Faraday & \(\nabla\times\mathbf E = -\partial_t\mathbf B\) &
\(\nabla\times\mathbf E = -\tfrac1c\partial_t\mathbf B\) \\
Ampère--Maxwell &
\(\nabla\times\mathbf B = \mu_0\mathbf J + \tfrac1{c^2}\partial_t\mathbf E\)
&
\(\nabla\times\mathbf B = \tfrac{4\pi}{c}\mathbf J + \tfrac1c\partial_t\mathbf E\) \\
Campo en función de potenciales &
\(\mathbf E = -\nabla\Phi - \partial_t\mathbf A\) &
\(\mathbf E = -\nabla\Phi - \tfrac1c\partial_t\mathbf A\) \\
Gauge de Lorenz &
\(\nabla\cdot\mathbf A + \tfrac1{c^2}\partial_t\Phi = 0\) &
\(\nabla\cdot\mathbf A + \tfrac1c\partial_t\Phi = 0\) \\
Ecuación de \(\Phi\) & \(\square\Phi = -\rho/\varepsilon_0\) &
\(\square\Phi = -4\pi\rho\) \\
Ecuación de \(\mathbf A\) & \(\square\mathbf A = -\mu_0\mathbf J\) &
\(\square\mathbf A = -\tfrac{4\pi}{c}\mathbf J\) \\
Potencial retardado &
\(\Phi = \dfrac{1}{4\pi\varepsilon_0}\displaystyle\int\frac{\rho(t_r)}{R}\,dV'\)
& \(\Phi = \displaystyle\int\frac{\rho(t_r)}{R}\,dV'\) \\
\end{longtable}

\end{minipage}%
\end{tcolorbox}

\section{Más allá: pérdidas y dispersión}\label{sec-hip-telegrafo}

Todo lo anterior supone un medio ideal: nada disipa energía y todas las
frecuencias viajan a la misma velocidad. Ningún medio real es así. Esta
última sección, que no está en el libro de la cátedra, mira qué se rompe
cuando se abandonan esas dos idealizaciones. Es la que justifica el caso
3 del sección~\ref{sec-hip-casos-limite} y la advertencia sobre las olas
de la intro.

\subsection{La ecuación del telégrafo}\label{sec-hip-ec-telegrafo}

Un cable largo de dos conductores ---un par telefónico, un coaxial---
tiene, por metro, una inductancia \(L\), una capacidad \(C\) entre los
conductores, una resistencia \(R\) de los conductores y una conductancia
de fuga \(G\) del aislante. Si \(V(x,t)\) es la tensión entre los
conductores e \(I(x,t)\) la corriente, aplicar las leyes de Kirchhoff a
un tramo \(dx\) da

\[
V_x = -L\,I_t - R\,I, \qquad I_x = -C\,V_t - G\,V .
\]

Derivando la primera en \(x\) y eliminando \(I\) con la segunda,

\begin{equation}\protect\phantomsection\label{eq-hip-telegrafo}{
V_{xx} = LC\,V_{tt} + (RC + GL)\,V_t + RG\,V .
}\end{equation}

\begin{definition}[Ecuación del
telégrafo]\protect\hypertarget{def-hip-telegrafo}{}\label{def-hip-telegrafo}

Se llama \textbf{ecuación del telégrafo} a \[
u_{tt} + (\alpha+\beta)\,u_t + \alpha\beta\,u = a^2\,u_{xx},
\] que es \ref{eq-hip-telegrafo} dividida por \(LC\), con
\(a = 1/\sqrt{LC}\), \(\alpha = R/L\) y \(\beta = G/C\).

\end{definition}

Su parte principal es la de la ecuación de onda, así que \textbf{es
hiperbólica}: la clasificación del sección~\ref{sec-clasificacion} solo
mira los términos de mayor orden. Las características son las mismas, la
información viaja a la misma velocidad \(a\), y los frentes llegan
nítidos. Lo que cambia es qué pasa \textbf{detrás} del frente.

Con \(R = G = 0\) es la ecuación de onda: la línea de transmisión de la
Tabla~\ref{tbl-hip-analogias}. Con \(R\) y \(G\) distintos de cero hay
dos efectos, y conviene separarlos.

\subsection{Pérdidas: la memoria se apaga}\label{sec-hip-perdidas}

El término \((\alpha+\beta)u_t\) es un \textbf{rozamiento}: una fuerza
proporcional a la velocidad y opuesta a ella. En una cuerda que vibra en
el aire, \(u_{tt} + 2\gamma u_t = a^2u_{xx}\), y con el mismo cálculo
del sección~\ref{sec-hip-balance}

\[
\frac{dE}{dt} = -2\gamma\int_0^l\rho\,u_t^2\,dx \;\le\; 0 .
\]

La energía \textbf{decrece} mientras la cuerda se mueva. Separando
variables en la cuerda fija, cada modo cumple ahora
\(T_n'' + 2\gamma T_n' + \omega_n^2T_n = 0\), que es un oscilador
amortiguado:

\[
T_n(t) = e^{-\gamma t}\big(a_n\cos\tilde\omega_nt + b_n\operatorname{sen}\tilde\omega_nt\big),
\qquad \tilde\omega_n = \sqrt{\omega_n^2-\gamma^2} .
\]

Todos los modos se apagan con el factor \(e^{-\gamma t}\). \textbf{Esa
es la justificación del caso 3 del sección~\ref{sec-hip-casos-limite}}:
después de un tiempo largo comparado con \(1/\gamma\), las oscilaciones
libres ---las que dependen de los datos iniciales--- desaparecieron, y
queda solo la respuesta a lo que siga forzando la cuerda. Y la
resonancia del sección~\ref{sec-hip-resonancia} deja de ser infinita: el
régimen permanente con forzante \(g_n\operatorname{sen}\Omega t\) tiene
amplitud

\[
\frac{|g_n|}{\sqrt{(\omega_n^2-\Omega^2)^2 + 4\gamma^2\Omega^2}},
\]

finita aun cuando \(\Omega = \omega_n\), con un pico tanto más alto y
angosto cuanto menor es \(\gamma\).

\subsection{Dispersión: la forma se desarma}\label{sec-hip-dispersion}

El segundo efecto es más sutil. Probemos una onda de una sola
frecuencia, \(u = e^{i(kx-\omega t)}\), en la ecuación del telégrafo.
Reemplazando, \[
-\omega^2 - i(\alpha+\beta)\,\omega + \alpha\beta = -a^2k^2 .
\]

\begin{definition}[Relación de
dispersión]\protect\hypertarget{def-hip-dispersion}{}\label{def-hip-dispersion}

La relación entre la frecuencia \(\omega\) y el número de onda \(k\) que
tiene que cumplir una onda \(e^{i(kx-\omega t)}\) para ser solución de
una ecuación lineal con coeficientes constantes se llama
\textbf{relación de dispersión}. La \textbf{velocidad de fase} es
\(\omega/k\). Una ecuación es \textbf{no dispersiva} si la velocidad de
fase es la misma para todos los \(k\), y \textbf{dispersiva} si no.

\end{definition}

Para la ecuación de onda, \(\omega^2 = a^2k^2\): la velocidad de fase es
\(a\) para todas las longitudes de onda, y por eso un pulso, que es una
suma de muchas, viaja sin deformarse. Esa es la segunda de las tres
propiedades de la intro, «la forma se conserva», leída en el lenguaje de
Fourier.

Para el telégrafo, resolviendo la cuadrática,

\[
\omega = -i\,\frac{\alpha+\beta}{2} \pm \sqrt{a^2k^2 - \Big(\frac{\alpha-\beta}{2}\Big)^2} .
\]

Para \(ak > |\alpha-\beta|/2\) la raíz es real: la parte imaginaria de
\(\omega\) es la atenuación, igual para todos esos \(k\), y la parte
real, dividida por \(k\), \textbf{sí depende de \(k\)}. (Para \(k\) más
chicos la raíz es imaginaria y esas componentes no se propagan: solo
decaen.) La dependencia de \(k\) dice que las ondas largas (con \(k\)
chico) viajan más despacio que las cortas. Un pulso, que es una suma de
ondas de todas las longitudes, se desarma: sus componentes se van
separando y detrás del frente queda una cola. En un cable telegráfico
largo eso significa que un punto y una raya enviados seguidos llegan
pegados uno con otro.

\begin{tcolorbox}[enhanced jigsaw, arc=.35mm, bottomrule=.15mm, breakable, colback=white, colframe=quarto-callout-important-color-frame, left=2mm, leftrule=.75mm, opacityback=0, rightrule=.15mm, toprule=.15mm]
\begin{minipage}[t]{5.5mm}
\textcolor{quarto-callout-important-color}{\faExclamation}
\end{minipage}%
\begin{minipage}[t]{\textwidth - 5.5mm}

\vspace{-3mm}\textbf{La condición de Heaviside}\vspace{3mm}

Hay un caso en que la raíz no depende de \(k\) de manera molesta: si
\(\alpha = \beta\), o sea

\[
\frac{R}{L} = \frac{G}{C},
\]

la raíz es exactamente \(ak\), y la relación de dispersión es
\(\omega = ak - i\alpha\). Todas las ondas viajan a la velocidad \(a\) y
se atenúan igual, y la solución es

\[
u(x,t) = e^{-\alpha t}\,f(x - at) :
\]

\textbf{el pulso llega más débil, pero con la misma forma}. Oliver
Heaviside lo descubrió en la década de 1880 y propuso algo que parecía
absurdo: para mejorar las líneas largas, que tenían \(G\) muy chico y
distorsionaban mucho, había que \textbf{aumentar la inductancia}
agregándoles bobinas. Pasó más de una década hasta que se hizo, y
funcionó.

\end{minipage}%
\end{tcolorbox}

\emph{(En la versión web hay acá un gráfico interactivo: un pulso que
viaja por una línea de transmisión con pérdidas, calculado con una serie
de Fourier exacta. Con \(R/L\) distinto de \(G/C\) el pulso se atenúa,
se ensancha y deja una cola; con \(R/L = G/C\) ---la condición de
Heaviside--- llega atenuado pero con la misma forma.)}

En la ecuación del telégrafo con \(\alpha\neq\beta\), un pulso que viaja
deja una cola detrás. ¿El frente del pulso llega \textbf{antes} que en
el cable sin pérdidas?

\begin{itemize}
\tightlist
\item
  {No: el frente viaja exactamente a velocidad \(a\), como sin pérdidas}
\item
  {Sí: las ondas cortas viajan más rápido que \(a\) y se adelantan}
\item
  {No: el frente llega después, porque las pérdidas lo frenan}
\item
  {Depende de si \(\alpha\) es mayor o menor que \(\beta\)}
\end{itemize}

El frente es una singularidad, y \textbf{las singularidades viajan por
las características} (sección~\ref{sec-hip-dalembert}). Las
características las decide la parte principal, que es la misma que la de
la ecuación de onda: \(x\pm at = \text{const}\). Así que el frente viaja
a velocidad \(a\), ni más ni menos.

La segunda opción tiene un grano de verdad mal leído: en la relación de
dispersión la velocidad de fase de las ondas cortas \textbf{tiende} a
\(a\) desde abajo cuando \(k\to\infty\), nunca la supera. Lo que la
dispersión cambia es lo que viene \textbf{detrás} del frente, no el
frente.

\begin{tcolorbox}[enhanced jigsaw, arc=.35mm, bottomrule=.15mm, breakable, colback=white, colframe=quarto-callout-note-color-frame, left=2mm, leftrule=.75mm, opacityback=0, rightrule=.15mm, toprule=.15mm]
\begin{minipage}[t]{5.5mm}
\textcolor{quarto-callout-note-color}{\faInfo}
\end{minipage}%
\begin{minipage}[t]{\textwidth - 5.5mm}

\vspace{-3mm}\textbf{Otras ecuaciones dispersivas}\vspace{3mm}

La dispersión no necesita pérdidas. La \textbf{ecuación de
Klein--Gordon}, \(u_{tt} = a^2u_{xx} - m^2u\), conserva la energía y es
dispersiva: su relación \(\omega^2 = a^2k^2 + m^2\) hace que las ondas
largas tengan velocidad de fase mayor que \(a\) y velocidad de grupo
---la del paquete, \(d\omega/dk\)--- menor que \(a\). La viga del
sección~\ref{sec-viga} es dispersiva, con \(\omega\propto k^2\). Y las
olas en aguas profundas cumplen \(\omega^2 = gk\): las olas largas son
más rápidas, y por eso el mar de fondo que llega a una playa después de
una tormenta lejana trae primero las olas más largas.

\end{minipage}%
\end{tcolorbox}

\section{Conclusiones}\label{sec-hip-conclusiones}

González cierra el capítulo con tres ideas: el problema hiperbólico
admite en general el método de las características; la presencia de
contornos puede sugerir sistemas de coordenadas convenientes en los que
el problema se separa; y las inhomogeneidades se tratan, en principio,
con la función de Green. Las tres están en este capítulo, y ahora se
puede decir algo más: \textbf{no son tres métodos distintos sino tres
maneras de mirar lo mismo}.

\begin{itemize}
\tightlist
\item
  \textbf{Las características} siguen a la onda. Dan d'Alembert, el
  dominio de dependencia, las reflexiones y la regla del paralelogramo.
  Son la herramienta cuando la pregunta es \textbf{cuándo llega algo}.
\item
  \textbf{Los modos} se quedan quietos y laten. Dan las frecuencias
  propias, las líneas nodales, la resonancia. Son la herramienta cuando
  la pregunta es \textbf{cómo suena} o \textbf{a qué frecuencia se
  rompe}.
\item
  \textbf{La función de Green} responde a un golpe y suma. Da Duhamel,
  la meseta de la cuerda, la cáscara del cono de luz en tres
  dimensiones, los potenciales retardados. Es la herramienta cuando hay
  \textbf{fuentes}.
\end{itemize}

Y las tres están unidas por dos hechos. Uno es la \textbf{linealidad},
que permite sumar: sumar ondas viajeras, sumar modos, sumar respuestas a
golpes. El otro es la \textbf{velocidad finita}, que es lo que hace
hiperbólica a la ecuación y lo que aparece en cada una de las tres
imágenes: como pendiente de las características, como el período
\(2l/a\) que fija las frecuencias de la cuerda, como el tiempo
retardado.

\subsection{Cómo elegir el método}\label{sec-hip-elegir}

\begin{longtable}[]{@{}
  >{\raggedright\arraybackslash}p{(\linewidth - 6\tabcolsep) * \real{0.3000}}
  >{\raggedright\arraybackslash}p{(\linewidth - 6\tabcolsep) * \real{0.2400}}
  >{\raggedright\arraybackslash}p{(\linewidth - 6\tabcolsep) * \real{0.2800}}
  >{\raggedright\arraybackslash}p{(\linewidth - 6\tabcolsep) * \real{0.1800}}@{}}
\caption{Guía para decidir por dónde
empezar.}\label{tbl-hip-elegir}\tabularnewline
\toprule\noalign{}
\begin{minipage}[b]{\linewidth}\raggedright
Si el problema tiene\ldots{}
\end{minipage} & \begin{minipage}[b]{\linewidth}\raggedright
Conviene\ldots{}
\end{minipage} & \begin{minipage}[b]{\linewidth}\raggedright
Porque\ldots{}
\end{minipage} & \begin{minipage}[b]{\linewidth}\raggedright
Ver
\end{minipage} \\
\midrule\noalign{}
\endfirsthead
\toprule\noalign{}
\begin{minipage}[b]{\linewidth}\raggedright
Si el problema tiene\ldots{}
\end{minipage} & \begin{minipage}[b]{\linewidth}\raggedright
Conviene\ldots{}
\end{minipage} & \begin{minipage}[b]{\linewidth}\raggedright
Porque\ldots{}
\end{minipage} & \begin{minipage}[b]{\linewidth}\raggedright
Ver
\end{minipage} \\
\midrule\noalign{}
\endhead
\bottomrule\noalign{}
\endlastfoot
una dimensión, sin bordes o con bordes lejos & d'Alembert & da la
solución cerrada y muestra el dominio de dependencia &
sección~\ref{sec-hip-dalembert} \\
un borde, en una dimensión & extender los datos (impar si es fijo, par
si es libre) & convierte el borde en una cuerda infinita &
sección~\ref{sec-hip-reflexiones} \\
dos bordes, y la pregunta es cuándo llega algo & reflexiones & se ven
los ecos & Teorema~\ref{thm-hip-finita} \\
dos bordes, y la pregunta es qué frecuencias hay & modos & la respuesta
sale directamente & sección~\ref{sec-cuerda-separacion} \\
una fuerza aplicada & modos + Duhamel en cada uno & cada modo es un
oscilador forzado & sección~\ref{sec-hip-forzada} \\
bordes que se mueven & restar una función que cumpla los bordes & pasa
el borde al interior como fuente &
sección~\ref{sec-hip-bordes-moviles} \\
dos dimensiones y un borde rectangular o circular & separación en
cartesianas o polares & el borde es una curva coordenada &
sección~\ref{sec-hip-membranas} \\
tres dimensiones, simetría esférica & \(v = ru\) & reduce a una cuerda
semi-infinita & sección~\ref{sec-hip-esfericas} \\
tres dimensiones, fuentes, sin bordes & potenciales retardados & Green
del espacio libre & sección~\ref{sec-hip-retardados} \\
una pregunta de existencia o unicidad, sin fórmula & energía & no
necesita resolver nada & sección~\ref{sec-hip-energia} \\
\end{longtable}

\section{Problemas del libro de la cátedra}\label{sec-hip-problemas}

Son los doce problemas de la sección 9.6 de González, con los enunciados
transcriptos y las soluciones completas. Salvo que se diga otra cosa, la
cuerda tiene largo \(l\), densidad \(\rho\), tensión \(T\) y velocidad
\(a = \sqrt{T/\rho}\), y se escribe \(\omega_n = n\pi a/l\).

\begin{exercise}[]\protect\hypertarget{exr-hip-p1}{}\label{exr-hip-p1}

\textbf{Problema 9.1.} Hallar las vibraciones de una cuerda con extremos
\(x = 0\) y \(x = l\) rígidamente fijos, perturbada por la desviación
inicial representada en la figura ---un triángulo de altura \(h\) con el
vértice en \(x_0\)---. Considerar nulas las velocidades iniciales.

\end{exercise}

\begin{tcolorbox}[enhanced jigsaw, arc=.35mm, bottomrule=.15mm, bottomtitle=1mm, breakable, colback=white, colbacktitle=quarto-callout-tip-color!10!white, colframe=quarto-callout-tip-color-frame, coltitle=black, left=2mm, leftrule=.75mm, opacityback=0, opacitybacktitle=0.6, rightrule=.15mm, title=\textcolor{quarto-callout-tip-color}{\faLightbulb}\hspace{0.5em}{Solución}, titlerule=0mm, toprule=.15mm, toptitle=1mm]

La desviación inicial es \[
\varphi(x) = \begin{cases} \dfrac{h\,x}{x_0}, & 0\le x\le x_0,\\[2mm] \dfrac{h\,(l-x)}{l-x_0}, & x_0\le x\le l,\end{cases}
\qquad \psi = 0 .
\] Por sección~\ref{sec-cuerda-separacion},
\(u = \sum_n b_n\cos(\omega_nt)\operatorname{sen}(n\pi x/l)\) con
\(b_n = \tfrac2l\int_0^l\varphi(x)\operatorname{sen}(n\pi x/l)\,dx\).
Integrando por partes en cada tramo ---los términos de borde se cancelan
porque \(\varphi\) es continua y nula en los extremos, y queda solo el
salto de \(\varphi'\) en \(x_0\)---, \[
b_n = \frac{2h\,l^2}{n^2\pi^2\,x_0\,(l-x_0)}\,\operatorname{sen}\frac{n\pi x_0}{l} .
\] Entonces \[
u(x,t) = \frac{2h\,l^2}{\pi^2x_0(l-x_0)}\sum_{n\ge1}\frac{1}{n^2}\operatorname{sen}\frac{n\pi x_0}{l}\,\cos\frac{n\pi at}{l}\,\operatorname{sen}\frac{n\pi x}{l}.
\] Los coeficientes decaen como \(1/n^2\), por la esquina. Y los modos
con un nodo en \(x_0\) ---los \(n\) con \(nx_0/l\) entero--- no
aparecen: es el truco del guitarrista del
sección~\ref{sec-interpretacion}. La misma solución en la imagen de
d'Alembert es la del gráfico «La misma cuerda, vista de dos maneras» del
sección~\ref{sec-hip-finita}.

\end{tcolorbox}

\begin{exercise}[]\protect\hypertarget{exr-hip-p2}{}\label{exr-hip-p2}

\textbf{Problema 9.2.} La cuerda \(0\le x\le l\) de extremos rígidamente
fijos se encontraba, hasta el momento \(t = 0\), en estado de equilibrio
bajo la acción de una fuerza transversal constante \(F_0\), aplicada al
punto \(x_0\) de la cuerda perpendicularmente a la posición no
perturbada de la misma. En \(t = 0\) la acción de la fuerza \(F_0\) cesa
instantáneamente. Hallar las oscilaciones de la cuerda para \(t>0\).

\end{exercise}

\begin{tcolorbox}[enhanced jigsaw, arc=.35mm, bottomrule=.15mm, bottomtitle=1mm, breakable, colback=white, colbacktitle=quarto-callout-tip-color!10!white, colframe=quarto-callout-tip-color-frame, coltitle=black, left=2mm, leftrule=.75mm, opacityback=0, opacitybacktitle=0.6, rightrule=.15mm, title=\textcolor{quarto-callout-tip-color}{\faLightbulb}\hspace{0.5em}{Solución}, titlerule=0mm, toprule=.15mm, toptitle=1mm]

\textbf{La forma inicial.} En equilibrio, lejos de \(x_0\) no hay
fuerza, así que la cuerda es recta a cada lado: un triángulo con el
vértice en \(x_0\) desplazado una cantidad \(-h\) (hacia abajo, como en
la figura del libro). En el vértice, las componentes verticales de la
tensión equilibran a \(F_0\): \[
T\,\frac{h}{x_0} + T\,\frac{h}{l-x_0} = F_0 \quad\Longrightarrow\quad h = \frac{F_0\,x_0\,(l-x_0)}{T\,l}.
\] (Con pendientes chicas,
\(\operatorname{sen}\alpha\approx\tan\alpha = h/x_0\).)

\textbf{Las oscilaciones.} Es el problema 9.1 con altura \(-h\).
Reemplazando, \[
u(x,t) = -\frac{2F_0\,l}{\pi^2T}\sum_{n\ge1}\frac{1}{n^2}\operatorname{sen}\frac{n\pi x_0}{l}\,\cos\frac{n\pi at}{l}\,\operatorname{sen}\frac{n\pi x}{l}.
\] Fijate que la amplitud depende de \(F_0/T\) y no de \(\rho\): la
forma inicial la decide el equilibrio de fuerzas, y la densidad solo
entra en la frecuencia, a través de \(a\).

\end{tcolorbox}

\begin{exercise}[]\protect\hypertarget{exr-hip-p3}{}\label{exr-hip-p3}

\textbf{Problema 9.3.} Los extremos de una cuerda están rígidamente
fijos y la desviación inicial tiene la forma de una parábola cuadrada,
simétrica con respecto a la perpendicular en el punto medio de la
cuerda. Hallar las oscilaciones de la cuerda si las velocidades
iniciales son idénticamente nulas.

\end{exercise}

\begin{tcolorbox}[enhanced jigsaw, arc=.35mm, bottomrule=.15mm, bottomtitle=1mm, breakable, colback=white, colbacktitle=quarto-callout-tip-color!10!white, colframe=quarto-callout-tip-color-frame, coltitle=black, left=2mm, leftrule=.75mm, opacityback=0, opacitybacktitle=0.6, rightrule=.15mm, title=\textcolor{quarto-callout-tip-color}{\faLightbulb}\hspace{0.5em}{Solución}, titlerule=0mm, toprule=.15mm, toptitle=1mm]

Llamando \(h\) a la desviación máxima, en el centro, la parábola que se
anula en los extremos es \(\varphi(x) = \dfrac{4h}{l^2}\,x(l-x)\). Con
dos integraciones por partes, \[
\int_0^l x(l-x)\operatorname{sen}\frac{n\pi x}{l}\,dx = \frac{2l^3}{n^3\pi^3}\big(1-(-1)^n\big),
\] así que \[
b_n = \frac2l\cdot\frac{4h}{l^2}\cdot\frac{2l^3}{n^3\pi^3}\big(1-(-1)^n\big) =
\begin{cases}\dfrac{32h}{n^3\pi^3}, & n \text{ impar},\\[2mm] 0, & n\text{ par},\end{cases}
\] y \[
u(x,t) = \frac{32h}{\pi^3}\sum_{n\ \text{impar}}\frac{1}{n^3}\cos\frac{n\pi at}{l}\,\operatorname{sen}\frac{n\pi x}{l}.
\] Los pares faltan por la simetría: la parábola es simétrica respecto
del centro y los modos pares son antisimétricos. Y los coeficientes
decaen como \(1/n^3\), más rápido que en el 9.1, porque la parábola no
tiene esquinas: el primer coeficiente es \(32/\pi^3\approx1{,}03\) veces
\(h\), y el segundo es \(27\) veces más chico. La cuerda suena casi como
un seno puro.

\end{tcolorbox}

\begin{exercise}[]\protect\hypertarget{exr-hip-p4}{}\label{exr-hip-p4}

\textbf{Problema 9.4.} Una cuerda de extremos fijos se perturba mediante
el golpe de un martillo plano y rígido, el cual le comunica la
distribución inicial de velocidades \[
u_t(x,0) = \psi(x) = \begin{cases} 0, & 0\le x\le x_0-\delta,\\ v_0, & x_0-\delta\le x\le x_0+\delta,\\ 0, & x_0+\delta\le x\le l.\end{cases}
\] Hallar las oscilaciones \(u(x,t)\) de la cuerda, si la desviación
inicial es nula.

\end{exercise}

\begin{tcolorbox}[enhanced jigsaw, arc=.35mm, bottomrule=.15mm, bottomtitle=1mm, breakable, colback=white, colbacktitle=quarto-callout-tip-color!10!white, colframe=quarto-callout-tip-color-frame, coltitle=black, left=2mm, leftrule=.75mm, opacityback=0, opacitybacktitle=0.6, rightrule=.15mm, title=\textcolor{quarto-callout-tip-color}{\faLightbulb}\hspace{0.5em}{Solución}, titlerule=0mm, toprule=.15mm, toptitle=1mm]

Con \(\varphi = 0\),
\(u = \sum B_n\operatorname{sen}(\omega_nt)\operatorname{sen}(n\pi x/l)\)
y
\(\omega_nB_n = \tfrac2l\int_0^l\psi\operatorname{sen}(n\pi x/l)\,dx\).
La integral es \[
\frac{2v_0}{l}\int_{x_0-\delta}^{x_0+\delta}\operatorname{sen}\frac{n\pi x}{l}\,dx
= \frac{2v_0}{n\pi}\Big[\cos\frac{n\pi(x_0-\delta)}{l} - \cos\frac{n\pi(x_0+\delta)}{l}\Big]
= \frac{4v_0}{n\pi}\operatorname{sen}\frac{n\pi x_0}{l}\operatorname{sen}\frac{n\pi\delta}{l}.
\] Dividiendo por \(\omega_n = n\pi a/l\), \[
u(x,t) = \frac{4v_0\,l}{\pi^2a}\sum_{n\ge1}\frac{1}{n^2}\operatorname{sen}\frac{n\pi x_0}{l}\operatorname{sen}\frac{n\pi\delta}{l}\,\operatorname{sen}\frac{n\pi at}{l}\,\operatorname{sen}\frac{n\pi x}{l}.
\] El factor \(\operatorname{sen}(n\pi\delta/l)\) lleva la información
del ancho del martillo: se anula cuando \(n\delta/l\) es entero, es
decir cuando el ancho del martillo, \(2\delta\), es un número entero de
longitudes de onda del modo, \(2l/n\). Esos modos no se excitan, porque
el martillo empuja por igual sus partes que suben y sus partes que
bajan. Y para \(n\) grande el factor ya no ayuda: un martillo ancho no
excita bien los armónicos altos.

\end{tcolorbox}

\begin{exercise}[]\protect\hypertarget{exr-hip-p5}{}\label{exr-hip-p5}

\textbf{Problema 9.5.} Una cuerda con extremos fijos se perturba
mediante un golpe de martillo, el cual le comunica un impulso \(I\) en
el punto \(x_0\). Hallar las oscilaciones de la cuerda si la desviación
inicial es nula.

\end{exercise}

\begin{tcolorbox}[enhanced jigsaw, arc=.35mm, bottomrule=.15mm, bottomtitle=1mm, breakable, colback=white, colbacktitle=quarto-callout-tip-color!10!white, colframe=quarto-callout-tip-color-frame, coltitle=black, left=2mm, leftrule=.75mm, opacityback=0, opacitybacktitle=0.6, rightrule=.15mm, title=\textcolor{quarto-callout-tip-color}{\faLightbulb}\hspace{0.5em}{Solución}, titlerule=0mm, toprule=.15mm, toptitle=1mm]

Un impulso \(I\) concentrado en \(x_0\) cambia el momento de la cuerda
en \(I\) en ese punto, así que \(\rho\,\psi(x) = I\,\delta(x-x_0)\).
Entonces \[
\omega_nB_n = \frac2l\cdot\frac{I}{\rho}\operatorname{sen}\frac{n\pi x_0}{l}
\quad\Longrightarrow\quad
B_n = \frac{2I}{\rho\,n\pi a}\operatorname{sen}\frac{n\pi x_0}{l},
\] y \[
u(x,t) = \frac{2I}{\pi\rho a}\sum_{n\ge1}\frac1n\operatorname{sen}\frac{n\pi x_0}{l}\,\operatorname{sen}\frac{n\pi at}{l}\,\operatorname{sen}\frac{n\pi x}{l}.
\]

\textbf{Control.} Es el límite del 9.4 cuando el martillo se hace
angosto con impulso fijo, \(I = \rho\cdot2\delta\cdot v_0\): con
\(\operatorname{sen}(n\pi\delta/l)\approx n\pi\delta/l\), el coeficiente
del 9.4 tiende a
\(\frac{4v_0\delta}{n\pi a}\operatorname{sen}\frac{n\pi x_0}{l} = \frac{2I}{\rho n\pi a}\operatorname{sen}\frac{n\pi x_0}{l}\).

Los coeficientes decaen como \(1/n\), más despacio que en la cuerda
pulsada: el golpe puntual tiene mucha energía en los armónicos altos.
Mientras la onda no llega a los extremos, la solución es la meseta del
ejercicio numérico del sección~\ref{sec-hip-green-1d}.

\end{tcolorbox}

\begin{exercise}[]\protect\hypertarget{exr-hip-p6}{}\label{exr-hip-p6}

\textbf{Problema 9.6.} Una cuerda homogénea de longitud \(l\) está fija
en el extremo \(x = 0\). El otro extremo está sujeto a un anillo, de
masa despreciable, que puede resbalar por una varilla lisa. El anillo
está apenas desviado una distancia \(h\) de su posición de equilibrio y
en el instante \(t = 0\) se deja libre. Hallar la desviación \(u(x,t)\)
de la cuerda para cualquier instante posterior.

\end{exercise}

\begin{tcolorbox}[enhanced jigsaw, arc=.35mm, bottomrule=.15mm, bottomtitle=1mm, breakable, colback=white, colbacktitle=quarto-callout-tip-color!10!white, colframe=quarto-callout-tip-color-frame, coltitle=black, left=2mm, leftrule=.75mm, opacityback=0, opacitybacktitle=0.6, rightrule=.15mm, title=\textcolor{quarto-callout-tip-color}{\faLightbulb}\hspace{0.5em}{Solución}, titlerule=0mm, toprule=.15mm, toptitle=1mm]

\textbf{Las condiciones.} Extremo fijo en \(x = 0\): \(u(0,t) = 0\).
Anillo sin masa en una varilla lisa: \(u_x(l,t) = 0\)
(Tabla~\ref{tbl-hip-bordes}). Antes de soltarlo, alguien sostiene el
anillo en \(h\), y la cuerda en equilibrio sin cargas es recta:
\(\varphi(x) = hx/l\), \(\psi = 0\).

\textbf{Los modos.} Con \(X'' + \lambda^2X = 0\), \(X(0) = 0\) da
\(X = \operatorname{sen}\lambda x\) y \(X'(l) = 0\) da
\(\cos\lambda l = 0\): \[
\lambda_n = \frac{(2n+1)\pi}{2l},\qquad X_n = \operatorname{sen}\frac{(2n+1)\pi x}{2l},\qquad \omega_n = \frac{(2n+1)\pi a}{2l},\qquad n = 0,1,2,\dots
\] Son los modos de un tubo cerrado en un extremo y abierto en el otro:
solo armónicos \textbf{impares} de la fundamental \(\pi a/(2l)\), que es
una octava más grave que la de la cuerda con los dos extremos fijos.

\textbf{Los coeficientes.} Con
\(\int_0^lx\operatorname{sen}\lambda_nx\,dx = -\frac{l\cos\lambda_nl}{\lambda_n} + \frac{\operatorname{sen}\lambda_nl}{\lambda_n^2} = \frac{(-1)^n}{\lambda_n^2}\),
\[
b_n = \frac2l\cdot\frac hl\cdot\frac{(-1)^n}{\lambda_n^2} = \frac{8h\,(-1)^n}{(2n+1)^2\pi^2},
\] \[
u(x,t) = \frac{8h}{\pi^2}\sum_{n\ge0}\frac{(-1)^n}{(2n+1)^2}\operatorname{sen}\frac{(2n+1)\pi x}{2l}\,\cos\frac{(2n+1)\pi at}{2l}.
\] \textbf{Control.} En \(x = l\), \(t = 0\):
\(\sum(-1)^n(-1)^n/(2n+1)^2 = \pi^2/8\), y queda \(u = h\).

\textbf{La imagen de d'Alembert.} El extremo libre refleja sin invertir
y el fijo invirtiendo, así que la extensión que resuelve el problema es
impar respecto de \(0\) y par respecto de \(l\), con período \(4l\). Por
eso la fundamental tiene período \(4l/a\) en vez de \(2l/a\).

\end{tcolorbox}

\begin{exercise}[]\protect\hypertarget{exr-hip-p7}{}\label{exr-hip-p7}

\textbf{Problema 9.7.} Hallar las oscilaciones longitudinales de una
barra elástica con los extremos libres, si las velocidades iniciales y
las desviaciones iniciales en la dirección longitudinal son arbitrarias.

\end{exercise}

\begin{tcolorbox}[enhanced jigsaw, arc=.35mm, bottomrule=.15mm, bottomtitle=1mm, breakable, colback=white, colbacktitle=quarto-callout-tip-color!10!white, colframe=quarto-callout-tip-color-frame, coltitle=black, left=2mm, leftrule=.75mm, opacityback=0, opacitybacktitle=0.6, rightrule=.15mm, title=\textcolor{quarto-callout-tip-color}{\faLightbulb}\hspace{0.5em}{Solución}, titlerule=0mm, toprule=.15mm, toptitle=1mm]

Por el Ejercicio~\ref{exr-hip-barra}, \(u_{tt} = a^2u_{xx}\) con
\(a = \sqrt{E/\rho}\) y, en extremos libres,
\(u_x(0,t) = u_x(l,t) = 0\). Datos \(\varphi\) y \(\psi\) arbitrarios.

\textbf{Los modos.} \(X'' + \lambda^2X = 0\) con \(X'(0) = X'(l) = 0\)
da \(X_n = \cos(n\pi x/l)\) con \(n = 0, 1, 2,\dots\) \textbf{El
\(n = 0\) está}, y es nuevo: \(X_0 = 1\), con \(\lambda_0 = 0\). Su
parte temporal cumple \(T_0'' = 0\), así que \(T_0 = A_0 + B_0t\).

\textbf{La solución.} \[
u(x,t) = A_0 + B_0\,t + \sum_{n\ge1}\big(A_n\cos\omega_nt + B_n\operatorname{sen}\omega_nt\big)\cos\frac{n\pi x}{l},
\] con \[
A_0 = \frac1l\int_0^l\varphi,\qquad B_0 = \frac1l\int_0^l\psi,\qquad
A_n = \frac2l\int_0^l\varphi\cos\frac{n\pi x}{l},\qquad \omega_nB_n = \frac2l\int_0^l\psi\cos\frac{n\pi x}{l}.
\] (El \(1/l\) en lugar de \(2/l\) del modo cero es porque
\(\int_0^l1^2 = l\) y no \(l/2\).)

\textbf{La física del modo cero.} \(A_0 + B_0t\) es un
\textbf{movimiento rígido}: la barra entera se traslada con la velocidad
media inicial \(B_0\), que es la velocidad del centro de masa. Como los
extremos están libres, nada la frena. Las oscilaciones son lo que pasa
encima de esa traslación.

\end{tcolorbox}

\begin{exercise}[]\protect\hypertarget{exr-hip-p8}{}\label{exr-hip-p8}

\textbf{Problema 9.8.} Hallar las oscilaciones de una barra elástica de
extremos libres, la cual recibió en el instante \(t = 0\) un impulso
longitudinal \(I\) en uno de los extremos.

\end{exercise}

\begin{tcolorbox}[enhanced jigsaw, arc=.35mm, bottomrule=.15mm, bottomtitle=1mm, breakable, colback=white, colbacktitle=quarto-callout-tip-color!10!white, colframe=quarto-callout-tip-color-frame, coltitle=black, left=2mm, leftrule=.75mm, opacityback=0, opacitybacktitle=0.6, rightrule=.15mm, title=\textcolor{quarto-callout-tip-color}{\faLightbulb}\hspace{0.5em}{Solución}, titlerule=0mm, toprule=.15mm, toptitle=1mm]

Sea \(S\) la sección y \(M = \rho Sl\) la masa. El impulso en el extremo
\(x = 0\) equivale a una velocidad inicial concentrada ahí:
\(\rho S\,\psi(x) = I\,\delta(x)\), tomando toda la delta adentro de la
barra, y \(\varphi = 0\). Con las fórmulas del 9.7, \[
B_0 = \frac1l\cdot\frac{I}{\rho S} = \frac IM,\qquad
\omega_nB_n = \frac2l\cdot\frac{I}{\rho S}\cos 0 = \frac{2I}{M},\qquad A_n = 0,
\] y \[
u(x,t) = \frac IM\,t + \frac{2I\,l}{\pi a M}\sum_{n\ge1}\frac1n\,\operatorname{sen}\frac{n\pi at}{l}\,\cos\frac{n\pi x}{l}.
\] El primer término es el centro de masa, que por conservación del
momento se mueve con velocidad \(I/M\). El segundo es un pulso de
compresión que recorre la barra de punta a punta, a velocidad \(a\),
rebotando en los extremos libres.

\end{tcolorbox}

\begin{exercise}[]\protect\hypertarget{exr-hip-p9}{}\label{exr-hip-p9}

\textbf{Problema 9.9.} Hallar las oscilaciones transversales de una
membrana rectangular \(0\le x\le l_1\), \(0\le y\le l_2\) con el borde
fijo, causadas por la desviación inicial
\(u(x,y,0) = Axy(l_1-x)(l_2-y)\), si se puede despreciar la reacción del
medio ambiente.

\end{exercise}

\begin{tcolorbox}[enhanced jigsaw, arc=.35mm, bottomrule=.15mm, bottomtitle=1mm, breakable, colback=white, colbacktitle=quarto-callout-tip-color!10!white, colframe=quarto-callout-tip-color-frame, coltitle=black, left=2mm, leftrule=.75mm, opacityback=0, opacitybacktitle=0.6, rightrule=.15mm, title=\textcolor{quarto-callout-tip-color}{\faLightbulb}\hspace{0.5em}{Solución}, titlerule=0mm, toprule=.15mm, toptitle=1mm]

Por el Teorema~\ref{thm-hip-rectangular}, con \(\psi = 0\) solo hay
cosenos, y la integral doble se factoriza: \[
A_{mn} = \frac{4A}{l_1l_2}\left[\int_0^{l_1}x(l_1-x)\operatorname{sen}\frac{m\pi x}{l_1}dx\right]\left[\int_0^{l_2}y(l_2-y)\operatorname{sen}\frac{n\pi y}{l_2}dy\right].
\] Cada corchete es la integral del problema 9.3:
\(\frac{2l_i^3}{k^3\pi^3}(1-(-1)^k)\), que vale
\(\frac{4l_i^3}{k^3\pi^3}\) si \(k\) es impar y \(0\) si es par.
Entonces \[
A_{mn} = \frac{64A\,l_1^2l_2^2}{\pi^6m^3n^3}\quad (m,n\text{ impares}),\qquad A_{mn} = 0\ \text{ si no},
\] \[
u(x,y,t) = \frac{64A\,l_1^2l_2^2}{\pi^6}\sum_{m,n\ \text{impares}}\frac{1}{m^3n^3}\cos(\omega_{mn}t)\operatorname{sen}\frac{m\pi x}{l_1}\operatorname{sen}\frac{n\pi y}{l_2},
\] con \(\omega_{mn}\) de \ref{eq-hip-frecuencias-rect}. El término
\((1,1)\) domina: el siguiente, \((1,3)\) o \((3,1)\), es \(27\) veces
más chico.

\end{tcolorbox}

\begin{exercise}[]\protect\hypertarget{exr-hip-p10}{}\label{exr-hip-p10}

\textbf{Problema 9.10.} Hallar las oscilaciones transversales de una
membrana rectangular \(0\le x\le l_1\), \(0\le y\le l_2\) con el borde
fijo, causadas por la distribución inicial de velocidades
\(u_t(x,y,0) = Axy(l_1-x)(l_2-y)\), si se puede despreciar la reacción
del medio ambiente.

\end{exercise}

\begin{tcolorbox}[enhanced jigsaw, arc=.35mm, bottomrule=.15mm, bottomtitle=1mm, breakable, colback=white, colbacktitle=quarto-callout-tip-color!10!white, colframe=quarto-callout-tip-color-frame, coltitle=black, left=2mm, leftrule=.75mm, opacityback=0, opacitybacktitle=0.6, rightrule=.15mm, title=\textcolor{quarto-callout-tip-color}{\faLightbulb}\hspace{0.5em}{Solución}, titlerule=0mm, toprule=.15mm, toptitle=1mm]

Es el 9.9 con los papeles cambiados: ahora \(\varphi = 0\) y la integral
da \(\omega_{mn}B_{mn}\). Así que \[
u(x,y,t) = \frac{64A\,l_1^2l_2^2}{\pi^6}\sum_{m,n\ \text{impares}}\frac{1}{m^3n^3\,\omega_{mn}}\operatorname{sen}(\omega_{mn}t)\operatorname{sen}\frac{m\pi x}{l_1}\operatorname{sen}\frac{n\pi y}{l_2}.
\] El factor \(1/\omega_{mn}\) adicional hace que los modos altos pesen
todavía menos que en el 9.9: una velocidad inicial se integra una vez
para dar desplazamiento, y la integración suaviza.

\end{tcolorbox}

\begin{exercise}[]\protect\hypertarget{exr-hip-p11}{}\label{exr-hip-p11}

\textbf{Problema 9.11.} Hallar las oscilaciones transversales de una
membrana circular con borde fijo, causadas por la distribución inicial
radialmente simétrica de las desviaciones y de las velocidades, y
considerando que la reacción del medio ambiente es despreciable.

\end{exercise}

\begin{tcolorbox}[enhanced jigsaw, arc=.35mm, bottomrule=.15mm, bottomtitle=1mm, breakable, colback=white, colbacktitle=quarto-callout-tip-color!10!white, colframe=quarto-callout-tip-color-frame, coltitle=black, left=2mm, leftrule=.75mm, opacityback=0, opacitybacktitle=0.6, rightrule=.15mm, title=\textcolor{quarto-callout-tip-color}{\faLightbulb}\hspace{0.5em}{Solución}, titlerule=0mm, toprule=.15mm, toptitle=1mm]

Es exactamente el Teorema~\ref{thm-hip-circular}. Con radio \(R_0\),
\(u(r,0) = \varphi(r)\) y \(u_t(r,0) = \psi(r)\), \[
u(r,t) = \sum_{p\ge1}\Big(A_p\cos\frac{a\,j_{0,p}\,t}{R_0} + B_p\operatorname{sen}\frac{a\,j_{0,p}\,t}{R_0}\Big)J_0\!\Big(\frac{j_{0,p}\,r}{R_0}\Big),
\] \[
A_p = \frac{2}{R_0^2J_1(j_{0,p})^2}\int_0^{R_0}\varphi(r)J_0\!\Big(\frac{j_{0,p}r}{R_0}\Big)r\,dr,
\qquad
B_p = \frac{2}{a\,j_{0,p}R_0J_1(j_{0,p})^2}\int_0^{R_0}\psi(r)J_0\!\Big(\frac{j_{0,p}r}{R_0}\Big)r\,dr .
\] La simetría radial de los datos garantiza que no aparecen modos con
dependencia angular: los coeficientes de los \(J_n(\cdot)\cos n\theta\)
con \(n\ge1\) son integrales de \(\cos n\theta\) en una vuelta, que dan
cero.

\end{tcolorbox}

\begin{exercise}[]\protect\hypertarget{exr-hip-p12}{}\label{exr-hip-p12}

\textbf{Problema 9.12.} Resolver el ejercicio anterior suponiendo que la
desviación inicial tiene la forma de un paraboloide de revolución y que
las velocidades iniciales son nulas.

\end{exercise}

\begin{tcolorbox}[enhanced jigsaw, arc=.35mm, bottomrule=.15mm, bottomtitle=1mm, breakable, colback=white, colbacktitle=quarto-callout-tip-color!10!white, colframe=quarto-callout-tip-color-frame, coltitle=black, left=2mm, leftrule=.75mm, opacityback=0, opacitybacktitle=0.6, rightrule=.15mm, title=\textcolor{quarto-callout-tip-color}{\faLightbulb}\hspace{0.5em}{Solución}, titlerule=0mm, toprule=.15mm, toptitle=1mm]

Con altura \(h\) en el centro y cero en el borde,
\(\varphi(r) = h\,(1 - r^2/R_0^2)\), y \(B_p = 0\). Cambiando a
\(s = j_{0,p}r/R_0\) y escribiendo \(j = j_{0,p}\), \[
A_p = \frac{2h}{j^2J_1(j)^2}\int_0^{j}\Big(s - \frac{s^3}{j^2}\Big)J_0(s)\,ds .
\] Hacen falta dos integrales de Bessel, que se demuestran en el
capítulo 12 derivando los segundos miembros: \[
\int s\,J_0(s)\,ds = s\,J_1(s), \qquad \int s^3J_0(s)\,ds = s^3J_1(s) - 2s^2J_2(s),
\] y la relación \(J_2(s) = \tfrac2sJ_1(s) - J_0(s)\). Con ellas, \[
\int_0^{j}\Big(s - \frac{s^3}{j^2}\Big)J_0\,ds = jJ_1(j) - \frac{j^3J_1(j) - 2j^2J_2(j)}{j^2} = 2J_2(j) = \frac{4J_1(j)}{j},
\] donde en el último paso se usó \(J_0(j) = 0\). Entonces \[
A_p = \frac{8h}{j_{0,p}^3\,J_1(j_{0,p})},\qquad
u(r,t) = 8h\sum_{p\ge1}\frac{J_0\big(j_{0,p}r/R_0\big)}{j_{0,p}^3\,J_1(j_{0,p})}\cos\frac{a\,j_{0,p}\,t}{R_0}.
\] \textbf{Control numérico.} Los primeros coeficientes son
\(A_1\approx1{,}108\,h\), \(A_2\approx-0{,}140\,h\),
\(A_3\approx0{,}045\,h\), y coinciden con la integración numérica
directa. La suma en \(r=0\), \(t=0\) da \(h\), porque \(J_0(0) = 1\) y
la serie reproduce el paraboloide. El primer modo solo ya es el
\(111\,\%\) de la altura: el segundo corrige en sentido contrario.

\end{tcolorbox}

\section{Resumen del capítulo}\label{sec-hip-resumen}

Esto es una guía de temas para el final. No reemplaza al capítulo: lo
condensa. Están todas las definiciones y todos los enunciados,
\textbf{copiados palabra por palabra} de donde aparecen, y además las
fórmulas que conviene tener a mano y un procedimiento para cuando te dan
un problema de ondas y no sabés por dónde empezarlo.

Los temas y el orden son los del capítulo 9 del libro de la cátedra. Al
final hay una sección con lo que este libro agrega por fuera de ese
programa, pero solo lo que hace falta para resolver los problemas 9.1 a
9.12 o para entender por qué los métodos funcionan.

\begin{quote}
\textbf{Sobre la notación.} Llamamos \(a\) a la velocidad de
propagación, como González, y reservamos \(c\) para la velocidad de la
luz. El electromagnetismo está en unidades SI; la traducción a las
fórmulas cgs del libro de la cátedra está en la
Tabla~\ref{tbl-hip-si-cgs}. González llama \(\sigma,\tau\) a los índices
de la membrana y \(l,m\) a sus lados; acá son \(m,n\) y \(l_1,l_2\).
\end{quote}

\subsection{El problema tipo y la unicidad}\label{sec-hip-resumen-1}

La ecuación sale de aplicar Newton a un pedacito de cuerda: la
aceleración es proporcional a la curvatura, con \(a = \sqrt{T/\rho}\).
La misma ecuación, con otro \(a\), describe la barra, el sonido, la
línea de transmisión, la membrana y el campo electromagnético
(Tabla~\ref{tbl-hip-analogias}).

\begin{tcolorbox}[enhanced jigsaw, arc=.35mm, bottomrule=.15mm, breakable, colback=white, colframe=quarto-callout-note-color-frame, left=2mm, leftrule=.75mm, opacityback=0, rightrule=.15mm, toprule=.15mm]
\begin{minipage}[t]{5.5mm}
\textcolor{quarto-callout-note-color}{\faInfo}
\end{minipage}%
\begin{minipage}[t]{\textwidth - 5.5mm}

\vspace{-3mm}\textbf{Definición~\ref{def-hip-onda} --- Ecuación de onda}\vspace{3mm}

Se llama \textbf{ecuación de onda} en \(n\) dimensiones espaciales a

\[
u_{tt} = a^2\,\nabla^2 u + f(\mathbf{x},t), \qquad \mathbf{x}\in\mathbb{R}^n,\ t>0,
\]

donde \(a>0\) es la \textbf{velocidad de propagación} y \(f\) la
\textbf{fuente} o \textbf{forzante}. Es \textbf{homogénea} si
\(f\equiv0\) y \textbf{no homogénea} en caso contrario. Con \(n=1\) es
la ecuación de la cuerda; con \(n=2\), la de la membrana; con \(n=3\),
la del sonido y la de la luz.

La versión con coeficientes variables, \ref{eq-hip-inhomogenea}, se
llama ecuación de onda \textbf{en un medio inhomogéneo}.

\end{minipage}%
\end{tcolorbox}

\begin{tcolorbox}[enhanced jigsaw, arc=.35mm, bottomrule=.15mm, breakable, colback=white, colframe=quarto-callout-note-color-frame, left=2mm, leftrule=.75mm, opacityback=0, rightrule=.15mm, toprule=.15mm]
\begin{minipage}[t]{5.5mm}
\textcolor{quarto-callout-note-color}{\faInfo}
\end{minipage}%
\begin{minipage}[t]{\textwidth - 5.5mm}

\vspace{-3mm}\textbf{Definición~\ref{def-hip-problema-mixto} --- Problema mixto de la cuerda}\vspace{3mm}

Hallar \(u(x,t)\), definida y de clase \(C^2\) en \(0\le x\le l\),
\(t\ge 0\), tal que

\[
\begin{aligned}
&u_{tt} = a^2u_{xx} + f(x,t), && 0<x<l,\ t>0, \\
&u(0,t) = \mu_1(t),\quad u(l,t) = \mu_2(t), && t\ge 0 \quad\text{(condiciones de contorno)},\\
&u(x,0) = \varphi(x),\quad u_t(x,0) = \psi(x), && 0\le x\le l \quad\text{(condiciones iniciales)}.
\end{aligned}
\]

Se llama \textbf{mixto} porque combina condiciones iniciales y de
contorno.

\end{minipage}%
\end{tcolorbox}

\textbf{Los tres casos límite} (González, 9.1). La escala que decide es
\(l/a\), el tiempo que tarda la onda en recorrer la cuerda:

\begin{enumerate}
\def\labelenumi{\arabic{enumi}.}
\tightlist
\item
  Tiempos cortos, lejos de los bordes: \textbf{problema de Cauchy}, solo
  condiciones iniciales en toda la recta.
\item
  Cerca de un borde, lejos del otro: \textbf{cuerda semi-infinita}, con
  \(u(0,t) = \mu(t)\).
\item
  Tiempos largos: \textbf{régimen permanente}, solo importan bordes y
  fuentes. Requiere algo de rozamiento: una cuerda ideal no pierde la
  memoria de los datos iniciales.
\end{enumerate}

\begin{tcolorbox}[enhanced jigsaw, arc=.35mm, bottomrule=.15mm, breakable, colback=white, colframe=quarto-callout-note-color-frame, left=2mm, leftrule=.75mm, opacityback=0, rightrule=.15mm, toprule=.15mm]
\begin{minipage}[t]{5.5mm}
\textcolor{quarto-callout-note-color}{\faInfo}
\end{minipage}%
\begin{minipage}[t]{\textwidth - 5.5mm}

\vspace{-3mm}\textbf{Teorema~\ref{thm-hip-unicidad} --- Unicidad del problema mixto (Teorema
9.1.1 de González)}\vspace{3mm}

Supongamos que

\begin{enumerate}
\def\labelenumi{\arabic{enumi}.}
\tightlist
\item
  los coeficientes \(\rho(x)\) y \(k(x)\) son continuos en \([0,l]\),
  con \(\rho(x)>0\) y \(k(x)>0\);
\item
  se busca \(u\) de clase \(C^2\) en \([0,l]\times[0,\infty)\).
\end{enumerate}

Entonces existe \textbf{a lo sumo una} función \(u(x,t)\) que satisface

\[
\rho(x)\,u_{tt} = \big(k(x)\,u_x\big)_x + F(x,t), \qquad 0<x<l,\ t>0,
\]

con las condiciones iniciales \(u(x,0) = \varphi(x)\),
\(u_t(x,0) = \psi(x)\) y las de contorno \(u(0,t) = \mu_1(t)\),
\(u(l,t) = \mu_2(t)\).

\end{minipage}%
\end{tcolorbox}

La demostración usa la energía:

\begin{tcolorbox}[enhanced jigsaw, arc=.35mm, bottomrule=.15mm, breakable, colback=white, colframe=quarto-callout-note-color-frame, left=2mm, leftrule=.75mm, opacityback=0, rightrule=.15mm, toprule=.15mm]
\begin{minipage}[t]{5.5mm}
\textcolor{quarto-callout-note-color}{\faInfo}
\end{minipage}%
\begin{minipage}[t]{\textwidth - 5.5mm}

\vspace{-3mm}\textbf{Definición~\ref{def-hip-energia} --- Energía de la cuerda}\vspace{3mm}

Para una solución de \ref{eq-hip-inhomogenea} en \([0,l]\) se llama
\textbf{energía} a

\[
E(t) = \int_0^l \Big[\tfrac12\,\rho(x)\,u_t^2 + \tfrac12\,k(x)\,u_x^2\Big]\,dx .
\]

El primer término es la energía \textbf{cinética} y el segundo la
\textbf{potencial}. Los integrandos se llaman \textbf{densidades de
energía}.

\end{minipage}%
\end{tcolorbox}

\begin{tcolorbox}[enhanced jigsaw, arc=.35mm, bottomrule=.15mm, breakable, colback=white, colframe=quarto-callout-note-color-frame, left=2mm, leftrule=.75mm, opacityback=0, rightrule=.15mm, toprule=.15mm]
\begin{minipage}[t]{5.5mm}
\textcolor{quarto-callout-note-color}{\faInfo}
\end{minipage}%
\begin{minipage}[t]{\textwidth - 5.5mm}

\vspace{-3mm}\textbf{Proposición~\ref{prp-hip-balance} --- Balance de energía}\vspace{3mm}

Si \(u\) es una solución de clase \(C^2\) de \ref{eq-hip-inhomogenea} en
\([0,l]\), entonces

\[
\frac{dE}{dt} \;=\; \int_0^l F\,u_t\,dx \;+\; \Big[\,k\,u_x\,u_t\,\Big]_{x=0}^{x=l} .
\]

\end{minipage}%
\end{tcolorbox}

\subsection{La fórmula de d'Alembert}\label{sec-hip-resumen-2}

Con \(\xi = x+at\), \(\eta = x-at\) la ecuación se vuelve
\(u_{\xi\eta} = 0\), de solución general \(u = f(x+at) + g(x-at)\)
(capítulo 8). Imponiendo los datos:

\begin{tcolorbox}[enhanced jigsaw, arc=.35mm, bottomrule=.15mm, breakable, colback=white, colframe=quarto-callout-note-color-frame, left=2mm, leftrule=.75mm, opacityback=0, rightrule=.15mm, toprule=.15mm]
\begin{minipage}[t]{5.5mm}
\textcolor{quarto-callout-note-color}{\faInfo}
\end{minipage}%
\begin{minipage}[t]{\textwidth - 5.5mm}

\vspace{-3mm}\textbf{Teorema~\ref{thm-hip-dalembert} --- Fórmula de d'Alembert}\vspace{3mm}

Sean \(\varphi\in C^2(\mathbb R)\) y \(\psi\in C^1(\mathbb R)\). El
problema de Cauchy

\[
u_{tt} = a^2u_{xx},\quad x\in\mathbb R,\ t>0;\qquad u(x,0) = \varphi(x),\quad u_t(x,0) = \psi(x),
\]

tiene una \textbf{única} solución de clase \(C^2\), que es

\[
u(x,t) \;=\; \underbrace{\frac{\varphi(x+at)+\varphi(x-at)}{2}}_{u_1:\ \text{la forma se parte}}
\;+\; \underbrace{\frac{1}{2a}\int_{x-at}^{x+at}\psi(s)\,ds}_{u_2:\ \text{la velocidad se acumula}} .
\]

\end{minipage}%
\end{tcolorbox}

\begin{itemize}
\tightlist
\item
  El primer término parte la forma inicial en dos mitades que viajan en
  sentidos opuestos; el segundo acumula la velocidad inicial sobre el
  \textbf{dominio de dependencia} \([x-at,\,x+at]\).
\item
  \textbf{Las singularidades viajan por las características}: si el dato
  tiene una esquina o un salto en \(x_0\), la solución los tiene solo
  sobre \(x = x_0\pm at\).
\end{itemize}

\begin{tcolorbox}[enhanced jigsaw, arc=.35mm, bottomrule=.15mm, breakable, colback=white, colframe=quarto-callout-note-color-frame, left=2mm, leftrule=.75mm, opacityback=0, rightrule=.15mm, toprule=.15mm]
\begin{minipage}[t]{5.5mm}
\textcolor{quarto-callout-note-color}{\faInfo}
\end{minipage}%
\begin{minipage}[t]{\textwidth - 5.5mm}

\vspace{-3mm}\textbf{Proposición~\ref{prp-hip-paralelogramo} --- Regla del paralelogramo}\vspace{3mm}

Sea \(ABCD\) un paralelogramo en el plano \((x,t)\) con los lados sobre
características: \(AB\) y \(DC\) paralelos a \(x-at=\text{const}\), y
\(AD\) y \(BC\) paralelos a \(x+at=\text{const}\). Si \(u\) es solución
de \(u_{tt}=a^2u_{xx}\) en una región que lo contiene,

\[
u(A) + u(C) = u(B) + u(D) .
\]

\end{minipage}%
\end{tcolorbox}

\subsection{Estabilidad de las soluciones}\label{sec-hip-resumen-3}

\begin{tcolorbox}[enhanced jigsaw, arc=.35mm, bottomrule=.15mm, breakable, colback=white, colframe=quarto-callout-note-color-frame, left=2mm, leftrule=.75mm, opacityback=0, rightrule=.15mm, toprule=.15mm]
\begin{minipage}[t]{5.5mm}
\textcolor{quarto-callout-note-color}{\faInfo}
\end{minipage}%
\begin{minipage}[t]{\textwidth - 5.5mm}

\vspace{-3mm}\textbf{Proposición~\ref{prp-hip-estabilidad} --- Dependencia continua de los
datos}\vspace{3mm}

Sean \(u_1\) y \(u_2\) las soluciones del problema de Cauchy con datos
\((\varphi_1,\psi_1)\) y \((\varphi_2,\psi_2)\), y sea \(t_0>0\). Si
para todo \(x\)

\[
|\varphi_1(x)-\varphi_2(x)| < \delta, \qquad |\psi_1(x)-\psi_2(x)| < \delta,
\]

entonces para todo \(x\) y todo \(0\le t\le t_0\)

\[
|u_1(x,t) - u_2(x,t)| \;\le\; \delta\,(1+t_0) .
\]

En particular, dado \(\varepsilon>0\), alcanza con tomar
\(\delta = \varepsilon/(1+t_0)\) para que las soluciones difieran en
menos de \(\varepsilon\) hasta el tiempo \(t_0\).

\end{minipage}%
\end{tcolorbox}

La cota crece linealmente con el tiempo, no exponencialmente como en el
ejemplo de Hadamard: un error uniforme \(\delta\) en la velocidad
inicial desplaza toda la cuerda en \(\delta t\).

\subsection{Propagación libre de ondas en el
espacio}\label{sec-hip-resumen-4}

\begin{tcolorbox}[enhanced jigsaw, arc=.35mm, bottomrule=.15mm, breakable, colback=white, colframe=quarto-callout-note-color-frame, left=2mm, leftrule=.75mm, opacityback=0, rightrule=.15mm, toprule=.15mm]
\begin{minipage}[t]{5.5mm}
\textcolor{quarto-callout-note-color}{\faInfo}
\end{minipage}%
\begin{minipage}[t]{\textwidth - 5.5mm}

\vspace{-3mm}\textbf{Proposición~\ref{prp-hip-esfericas} --- Reducción a una dimensión}\vspace{3mm}

Si \(u(r,t)\) es solución radial de la ecuación de onda en
\(\mathbb R^3\), entonces \(v = r\,u\) cumple la ecuación de la cuerda,
\[
v_{rr} = \frac1{a^2}\,v_{tt}, \qquad r>0,
\] y si \(u\) es acotada en el origen, \(v(0,t) = 0\). En consecuencia

\[
u(r,t) = \frac{F(t - r/a)}{r} + \frac{G(t + r/a)}{r},
\]

con \(G = -F\) si la solución es regular en el origen.

\end{minipage}%
\end{tcolorbox}

\begin{tcolorbox}[enhanced jigsaw, arc=.35mm, bottomrule=.15mm, breakable, colback=white, colframe=quarto-callout-note-color-frame, left=2mm, leftrule=.75mm, opacityback=0, rightrule=.15mm, toprule=.15mm]
\begin{minipage}[t]{5.5mm}
\textcolor{quarto-callout-note-color}{\faInfo}
\end{minipage}%
\begin{minipage}[t]{\textwidth - 5.5mm}

\vspace{-3mm}\textbf{Definición~\ref{def-hip-onda-esferica} --- Ondas esféricas}\vspace{3mm}

Las soluciones \(\dfrac{F(t-r/a)}{r}\) y \(\dfrac{G(t+r/a)}{r}\) se
llaman \textbf{ondas esféricas} \textbf{divergente} (o saliente) y
\textbf{convergente} (o entrante), respectivamente.

\end{minipage}%
\end{tcolorbox}

\begin{itemize}
\tightlist
\item
  La amplitud cae como \(1/r\) porque la energía se reparte en esferas
  de área \(4\pi r^2\).
\item
  El origen hace de \textbf{extremo fijo}: una onda convergente rebota
  en el centro y sale divergente.
\item
  Con datos radiales, \(v = ru\) es una cuerda semi-infinita con datos
  \(r\varphi\), \(r\psi\). En el centro,
  \(u(0,t) = \varphi(at) + at\,\varphi'(at) + t\,\psi(at)\): aparece la
  \textbf{derivada} del dato (enfoque).
\end{itemize}

\begin{tcolorbox}[enhanced jigsaw, arc=.35mm, bottomrule=.15mm, breakable, colback=white, colframe=quarto-callout-note-color-frame, left=2mm, leftrule=.75mm, opacityback=0, rightrule=.15mm, toprule=.15mm]
\begin{minipage}[t]{5.5mm}
\textcolor{quarto-callout-note-color}{\faInfo}
\end{minipage}%
\begin{minipage}[t]{\textwidth - 5.5mm}

\vspace{-3mm}\textbf{Teorema~\ref{thm-hip-kirchhoff-cauchy} --- Fórmula de Kirchhoff para el
problema de Cauchy en \(\mathbb R^3\)}\vspace{3mm}

Si \(\varphi\in C^3\) y \(\psi\in C^2\), la solución de
\(u_{tt} = a^2\nabla^2u\) en \(\mathbb R^3\) con
\(u(\mathbf x,0) = \varphi\) y \(u_t(\mathbf x,0) = \psi\) es

\[
u(\mathbf x,t) = \frac{\partial}{\partial t}\Big(t\,M_{at}[\varphi](\mathbf x)\Big) + t\,M_{at}[\psi](\mathbf x).
\]

\end{minipage}%
\end{tcolorbox}

\textbf{Principio de Huygens.} En tres dimensiones la solución en un
punto depende de los datos solo sobre una esfera, y una perturbación
localizada pasa sin dejar nada detrás. En dos dimensiones depende de
todo el disco y queda una estela; en una, una meseta
(Tabla~\ref{tbl-hip-green-dimensiones}).

\subsection{Membrana rectangular}\label{sec-hip-resumen-5}

Separando \(u = U(x,y)\,T(t)\) sale la ecuación de Helmholtz
\(\nabla^2U + k^2U = 0\) con \(U = 0\) en el borde, y separando otra
vez, \(U = X(x)\,Y(y)\), dos problemas de la cuerda.

\begin{tcolorbox}[enhanced jigsaw, arc=.35mm, bottomrule=.15mm, breakable, colback=white, colframe=quarto-callout-note-color-frame, left=2mm, leftrule=.75mm, opacityback=0, rightrule=.15mm, toprule=.15mm]
\begin{minipage}[t]{5.5mm}
\textcolor{quarto-callout-note-color}{\faInfo}
\end{minipage}%
\begin{minipage}[t]{\textwidth - 5.5mm}

\vspace{-3mm}\textbf{Teorema~\ref{thm-hip-rectangular} --- Vibraciones libres de la membrana
rectangular}\vspace{3mm}

La solución de \(u_{tt} = a^2\nabla^2u\) en \([0,l_1]\times[0,l_2]\),
con \(u=0\) en el borde, \(u(x,y,0) = \varphi(x,y)\) y
\(u_t(x,y,0) = \psi(x,y)\), es

\[
u(x,y,t) = \sum_{m,n\ge1}\big(A_{mn}\cos\omega_{mn}t + B_{mn}\operatorname{sen}\omega_{mn}t\big)
\operatorname{sen}\frac{m\pi x}{l_1}\,\operatorname{sen}\frac{n\pi y}{l_2},
\]

con las \textbf{frecuencias propias}

\[
\omega_{mn} = a\pi\sqrt{\frac{m^2}{l_1^2} + \frac{n^2}{l_2^2}}
\]

y los coeficientes de la \textbf{serie doble de Fourier} de los datos:

\[
A_{mn} = \frac{4}{l_1l_2}\int_0^{l_1}\!\!\int_0^{l_2}\varphi(x,y)\operatorname{sen}\frac{m\pi x}{l_1}\operatorname{sen}\frac{n\pi y}{l_2}\,dy\,dx,
\qquad
\omega_{mn}B_{mn} = \frac{4}{l_1l_2}\int_0^{l_1}\!\!\int_0^{l_2}\psi(x,y)\operatorname{sen}\frac{m\pi x}{l_1}\operatorname{sen}\frac{n\pi y}{l_2}\,dy\,dx .
\]

\end{minipage}%
\end{tcolorbox}

\begin{tcolorbox}[enhanced jigsaw, arc=.35mm, bottomrule=.15mm, breakable, colback=white, colframe=quarto-callout-note-color-frame, left=2mm, leftrule=.75mm, opacityback=0, rightrule=.15mm, toprule=.15mm]
\begin{minipage}[t]{5.5mm}
\textcolor{quarto-callout-note-color}{\faInfo}
\end{minipage}%
\begin{minipage}[t]{\textwidth - 5.5mm}

\vspace{-3mm}\textbf{Definición~\ref{def-hip-linea-nodal} --- Línea nodal}\vspace{3mm}

En un modo \(u = U(x,y)\,T(t)\), se llama \textbf{línea nodal} a cada
curva del conjunto \(\{(x,y): U(x,y) = 0\}\) en el interior de la
membrana. Sobre ella la membrana no se mueve nunca.

\end{minipage}%
\end{tcolorbox}

\begin{tcolorbox}[enhanced jigsaw, arc=.35mm, bottomrule=.15mm, breakable, colback=white, colframe=quarto-callout-note-color-frame, left=2mm, leftrule=.75mm, opacityback=0, rightrule=.15mm, toprule=.15mm]
\begin{minipage}[t]{5.5mm}
\textcolor{quarto-callout-note-color}{\faInfo}
\end{minipage}%
\begin{minipage}[t]{\textwidth - 5.5mm}

\vspace{-3mm}\textbf{Definición~\ref{def-hip-degeneracion} --- Autovalor degenerado}\vspace{3mm}

Un autovalor \(k^2\) de \ref{eq-hip-helmholtz-rect} es
\textbf{degenerado} si le corresponden dos o más autofunciones
linealmente independientes. La cantidad máxima de independientes es su
\textbf{multiplicidad}, y cualquier combinación lineal de ellas es
también autofunción.

\end{minipage}%
\end{tcolorbox}

\begin{itemize}
\tightlist
\item
  Las frecuencias \textbf{no} son múltiplos de la fundamental: la
  membrana no tiene nota definida.
\item
  En el cuadrado, \(\omega_{mn} = \omega_{nm}\); la mezcla de \((1,2)\)
  y \((2,1)\) tiene líneas nodales diagonales, que no tiene ninguno de
  los dos modos.
\item
  Las figuras de Chladni del libro de la cátedra son, en rigor, de
  placas (cuarto orden).
\end{itemize}

\subsection{Vibraciones forzadas de la cuerda
fija}\label{sec-hip-resumen-6}

Por linealidad, \(u = v + w\): \(w\) son las oscilaciones libres con los
datos iniciales; \(v\), las forzadas desde el reposo. Desarrollando
\(v\) y la fuerza en los modos \(\operatorname{sen}(n\pi x/l)\), cada
coeficiente cumple la ecuación de un oscilador forzado, \[
T_n'' + \omega_n^2T_n = f_n(t),\qquad T_n(0) = T_n'(0) = 0,\qquad f_n(t) = \frac2l\int_0^lf(z,t)\operatorname{sen}\frac{n\pi z}{l}\,dz .
\]

\begin{tcolorbox}[enhanced jigsaw, arc=.35mm, bottomrule=.15mm, breakable, colback=white, colframe=quarto-callout-note-color-frame, left=2mm, leftrule=.75mm, opacityback=0, rightrule=.15mm, toprule=.15mm]
\begin{minipage}[t]{5.5mm}
\textcolor{quarto-callout-note-color}{\faInfo}
\end{minipage}%
\begin{minipage}[t]{\textwidth - 5.5mm}

\vspace{-3mm}\textbf{Teorema~\ref{thm-hip-forzada} --- Oscilaciones forzadas de la cuerda
fija}\vspace{3mm}

La solución de \ref{eq-hip-forzada} con \(\varphi = \psi = 0\) es

\[
v(x,t) = \sum_{n\ge1}\left[\frac{1}{\omega_n}\int_0^t f_n(\tau)\operatorname{sen}\omega_n(t-\tau)\,d\tau\right]\operatorname{sen}\frac{n\pi x}{l}
= \sum_{n\ge1}\left[\frac{2}{l\,\omega_n}\int_0^t\!\!\int_0^l f(z,\tau)\operatorname{sen}\frac{n\pi z}{l}\operatorname{sen}\omega_n(t-\tau)\,dz\,d\tau\right]\operatorname{sen}\frac{n\pi x}{l}.
\]

La solución general de \ref{eq-hip-forzada} es \(u = v + w\).

\end{minipage}%
\end{tcolorbox}

\subsection{Ecuaciones de Maxwell y
potenciales}\label{sec-hip-resumen-7}

\begin{tcolorbox}[enhanced jigsaw, arc=.35mm, bottomrule=.15mm, breakable, colback=white, colframe=quarto-callout-note-color-frame, left=2mm, leftrule=.75mm, opacityback=0, rightrule=.15mm, toprule=.15mm]
\begin{minipage}[t]{5.5mm}
\textcolor{quarto-callout-note-color}{\faInfo}
\end{minipage}%
\begin{minipage}[t]{\textwidth - 5.5mm}

\vspace{-3mm}\textbf{Proposición~\ref{prp-hip-ondas-em} --- Los campos en el vacío son ondas}\vspace{3mm}

Donde no hay cargas ni corrientes, cada componente cartesiana de
\(\mathbf E\) y de \(\mathbf B\) cumple la ecuación de onda

\[
\nabla^2\mathbf E = \frac{1}{c^2}\,\frac{\partial^2\mathbf E}{\partial t^2},\qquad
\nabla^2\mathbf B = \frac{1}{c^2}\,\frac{\partial^2\mathbf B}{\partial t^2},\qquad
c = \frac{1}{\sqrt{\mu_0\varepsilon_0}} .
\]

\end{minipage}%
\end{tcolorbox}

\begin{tcolorbox}[enhanced jigsaw, arc=.35mm, bottomrule=.15mm, breakable, colback=white, colframe=quarto-callout-note-color-frame, left=2mm, leftrule=.75mm, opacityback=0, rightrule=.15mm, toprule=.15mm]
\begin{minipage}[t]{5.5mm}
\textcolor{quarto-callout-note-color}{\faInfo}
\end{minipage}%
\begin{minipage}[t]{\textwidth - 5.5mm}

\vspace{-3mm}\textbf{Definición~\ref{def-hip-potenciales} --- Potenciales electromagnéticos}\vspace{3mm}

El \textbf{potencial vector} \(\mathbf A\) y el \textbf{potencial
escalar} \(\Phi\) son campos tales que

\[
\mathbf B = \nabla\times\mathbf A, \qquad \mathbf E = -\nabla\Phi - \frac{\partial\mathbf A}{\partial t}.
\]

\end{minipage}%
\end{tcolorbox}

\begin{tcolorbox}[enhanced jigsaw, arc=.35mm, bottomrule=.15mm, breakable, colback=white, colframe=quarto-callout-note-color-frame, left=2mm, leftrule=.75mm, opacityback=0, rightrule=.15mm, toprule=.15mm]
\begin{minipage}[t]{5.5mm}
\textcolor{quarto-callout-note-color}{\faInfo}
\end{minipage}%
\begin{minipage}[t]{\textwidth - 5.5mm}

\vspace{-3mm}\textbf{Definición~\ref{def-hip-gauge} --- Transformación de gauge (de medida)}\vspace{3mm}

Para cualquier función escalar \(\Lambda(\mathbf x,t)\), los potenciales
\[
\mathbf A' = \mathbf A + \nabla\Lambda, \qquad \Phi' = \Phi - \frac{\partial\Lambda}{\partial t}
\] producen los mismos \(\mathbf E\) y \(\mathbf B\). Esta
indeterminación se llama \textbf{libertad de gauge} o \textbf{de
medida}, y se fija imponiendo una condición adicional. La
\textbf{condición de Lorenz} es

\[
\nabla\cdot\mathbf A + \frac{1}{c^2}\frac{\partial\Phi}{\partial t} = 0 .
\]

\end{minipage}%
\end{tcolorbox}

En el gauge de Lorenz, los cuatro potenciales cumplen ecuaciones de onda
desacopladas: \(\square\Phi = -\rho/\varepsilon_0\),
\(\square\mathbf A = -\mu_0\mathbf J\), con
\(\square = \nabla^2 - \tfrac1{c^2}\partial_t^2\).

\subsection{Potenciales retardados}\label{sec-hip-resumen-8}

Transformando Fourier en el tiempo queda la ecuación de Helmholtz con
fuente. Su función de Green sin bordes es
\(\dfrac{Ae^{ikR}+Be^{-ikR}}{R}\) con \(A+B = \tfrac{1}{4\pi}\); al
volver al tiempo, el término con \(e^{-ikR}\) responde \textbf{antes} de
la causa, y la causalidad lo descarta.

\begin{tcolorbox}[enhanced jigsaw, arc=.35mm, bottomrule=.15mm, breakable, colback=white, colframe=quarto-callout-note-color-frame, left=2mm, leftrule=.75mm, opacityback=0, rightrule=.15mm, toprule=.15mm]
\begin{minipage}[t]{5.5mm}
\textcolor{quarto-callout-note-color}{\faInfo}
\end{minipage}%
\begin{minipage}[t]{\textwidth - 5.5mm}

\vspace{-3mm}\textbf{Teorema~\ref{thm-hip-retardados} --- Potenciales retardados}\vspace{3mm}

En el gauge de Lorenz y sin bordes, la solución causal de
\ref{eq-hip-potenciales-onda} es

\[
\Phi(\mathbf x,t) = \frac{1}{4\pi\varepsilon_0}\int\frac{\rho\big(\mathbf x',\,t_r\big)}{|\mathbf x-\mathbf x'|}\,dV',
\qquad
\mathbf A(\mathbf x,t) = \frac{\mu_0}{4\pi}\int\frac{\mathbf J\big(\mathbf x',\,t_r\big)}{|\mathbf x-\mathbf x'|}\,dV',
\]

donde

\[
t_r = t - \frac{|\mathbf x - \mathbf x'|}{c}
\]

es el \textbf{tiempo retardado}: el instante en que tuvo que salir de
\(\mathbf x'\) una señal a velocidad \(c\) para llegar a \(\mathbf x\)
en el instante \(t\).

\end{minipage}%
\end{tcolorbox}

\subsection{Fórmulas de Green y fórmula de
Kirchhoff}\label{sec-hip-resumen-9}

\begin{tcolorbox}[enhanced jigsaw, arc=.35mm, bottomrule=.15mm, breakable, colback=white, colframe=quarto-callout-note-color-frame, left=2mm, leftrule=.75mm, opacityback=0, rightrule=.15mm, toprule=.15mm]
\begin{minipage}[t]{5.5mm}
\textcolor{quarto-callout-note-color}{\faInfo}
\end{minipage}%
\begin{minipage}[t]{\textwidth - 5.5mm}

\vspace{-3mm}\textbf{Teorema~\ref{thm-hip-identidades-green} --- Identidades de Green}\vspace{3mm}

Sea \(V\) una región acotada con frontera \(S\) regular y normal
exterior \(\mathbf n\), y sean \(\phi,\psi\) funciones de clase \(C^2\)
en \(V\) y \(C^1\) hasta el borde. Entonces

\[
\int_V\big(\phi\,\nabla^2\psi + \nabla\phi\cdot\nabla\psi\big)\,dV = \oint_S\phi\,\frac{\partial\psi}{\partial n}\,dS
\qquad\text{(primera identidad)},
\]

\[
\int_V\big(\phi\,\nabla^2\psi - \psi\,\nabla^2\phi\big)\,dV = \oint_S\Big(\phi\,\frac{\partial\psi}{\partial n} - \psi\,\frac{\partial\phi}{\partial n}\Big)\,dS
\qquad\text{(segunda identidad)}.
\]

\end{minipage}%
\end{tcolorbox}

\begin{tcolorbox}[enhanced jigsaw, arc=.35mm, bottomrule=.15mm, breakable, colback=white, colframe=quarto-callout-note-color-frame, left=2mm, leftrule=.75mm, opacityback=0, rightrule=.15mm, toprule=.15mm]
\begin{minipage}[t]{5.5mm}
\textcolor{quarto-callout-note-color}{\faInfo}
\end{minipage}%
\begin{minipage}[t]{\textwidth - 5.5mm}

\vspace{-3mm}\textbf{Teorema~\ref{thm-hip-kirchhoff} --- Fórmula integral de Kirchhoff}\vspace{3mm}

Si \((\nabla^2+k^2)\psi = 0\) en \(V\) y
\(G(\mathbf x,\mathbf x') = \dfrac{e^{ik|\mathbf x-\mathbf x'|}}{4\pi|\mathbf x-\mathbf x'|}\),
entonces para todo \(\mathbf x\) en el interior de \(V\)

\[
\psi(\mathbf x) = \oint_S\left[G(\mathbf x,\mathbf x')\,\frac{\partial\psi}{\partial n'}(\mathbf x') - \psi(\mathbf x')\,\frac{\partial G}{\partial n'}(\mathbf x,\mathbf x')\right]dS',
\]

con \(\mathbf n'\) la normal \textbf{exterior} a \(V\).

\end{minipage}%
\end{tcolorbox}

\begin{itemize}
\tightlist
\item
  Es el principio de Huygens formalizado: cada punto del borde emite una
  onda esférica secundaria.
\item
  \textbf{Aproximación de Kirchhoff} para difracción: en el agujero, la
  onda incidente; en la pantalla, cero. Funciona bien si el agujero es
  grande frente a \(\lambda\), pero es inconsistente, porque una
  solución con valor y derivada normal nulos sobre un pedazo de
  superficie es idénticamente nula.
\item
  Las convenciones de signo y de \(4\pi\) cambian entre libros:
  Tabla~\ref{tbl-hip-convenciones-green}.
\end{itemize}

\subsection{Lo que agrega este libro}\label{sec-hip-resumen-10}

\subsubsection{Bordes y reflexiones}\label{bordes-y-reflexiones}

Hacen falta para el problema 9.6 (extremo libre) y para entender por qué
la cuerda finita tiene modos.

\begin{tcolorbox}[enhanced jigsaw, arc=.35mm, bottomrule=.15mm, breakable, colback=white, colframe=quarto-callout-note-color-frame, left=2mm, leftrule=.75mm, opacityback=0, rightrule=.15mm, toprule=.15mm]
\begin{minipage}[t]{5.5mm}
\textcolor{quarto-callout-note-color}{\faInfo}
\end{minipage}%
\begin{minipage}[t]{\textwidth - 5.5mm}

\vspace{-3mm}\textbf{Definición~\ref{def-hip-extension} --- Extensión impar y par}\vspace{3mm}

Dada \(\varphi\) definida en \([0,\infty)\), su \textbf{extensión impar}
y su \textbf{extensión par} son las funciones definidas en toda la recta
por

\[
\varphi_{\text{i}}(x) = \begin{cases}\varphi(x), & x\ge0,\\ -\varphi(-x), & x<0,\end{cases}
\qquad
\varphi_{\text{p}}(x) = \begin{cases}\varphi(x), & x\ge0,\\ \varphi(-x), & x<0.\end{cases}
\]

\end{minipage}%
\end{tcolorbox}

\begin{tcolorbox}[enhanced jigsaw, arc=.35mm, bottomrule=.15mm, breakable, colback=white, colframe=quarto-callout-note-color-frame, left=2mm, leftrule=.75mm, opacityback=0, rightrule=.15mm, toprule=.15mm]
\begin{minipage}[t]{5.5mm}
\textcolor{quarto-callout-note-color}{\faInfo}
\end{minipage}%
\begin{minipage}[t]{\textwidth - 5.5mm}

\vspace{-3mm}\textbf{Proposición~\ref{prp-hip-semi-fijo} --- Cuerda semi-infinita con extremo
fijo}\vspace{3mm}

La solución de \ref{eq-hip-semi} es la de d'Alembert con los datos
extendidos en forma impar, restringida a \(x\ge0\). Escrita solo con los
datos originales,

\[
u(x,t) =
\begin{cases}
\dfrac{\varphi(x+at)+\varphi(x-at)}{2} + \dfrac{1}{2a}\displaystyle\int_{x-at}^{x+at}\psi(s)\,ds, & x\ge at,\\[3mm]
\dfrac{\varphi(x+at)-\varphi(at-x)}{2} + \dfrac{1}{2a}\displaystyle\int_{at-x}^{x+at}\psi(s)\,ds, & 0\le x< at.
\end{cases}
\]

\end{minipage}%
\end{tcolorbox}

\begin{tcolorbox}[enhanced jigsaw, arc=.35mm, bottomrule=.15mm, breakable, colback=white, colframe=quarto-callout-note-color-frame, left=2mm, leftrule=.75mm, opacityback=0, rightrule=.15mm, toprule=.15mm]
\begin{minipage}[t]{5.5mm}
\textcolor{quarto-callout-note-color}{\faInfo}
\end{minipage}%
\begin{minipage}[t]{\textwidth - 5.5mm}

\vspace{-3mm}\textbf{Proposición~\ref{prp-hip-semi-libre} --- Cuerda semi-infinita con
extremo libre}\vspace{3mm}

La solución de \ref{eq-hip-semi} con \(u_x(0,t)=0\) en lugar de
\(u(0,t)=0\) es la de d'Alembert con los datos extendidos en forma
\textbf{par}. Para \(0\le x<at\),

\[
u(x,t) = \frac{\varphi(x+at)+\varphi(at-x)}{2} + \frac{1}{2a}\left[\int_{0}^{x+at}\psi(s)\,ds + \int_{0}^{at-x}\psi(s)\,ds\right],
\]

y para \(x\ge at\) es la fórmula de d'Alembert sin cambios.

\end{minipage}%
\end{tcolorbox}

\begin{tcolorbox}[enhanced jigsaw, arc=.35mm, bottomrule=.15mm, breakable, colback=white, colframe=quarto-callout-note-color-frame, left=2mm, leftrule=.75mm, opacityback=0, rightrule=.15mm, toprule=.15mm]
\begin{minipage}[t]{5.5mm}
\textcolor{quarto-callout-note-color}{\faInfo}
\end{minipage}%
\begin{minipage}[t]{\textwidth - 5.5mm}

\vspace{-3mm}\textbf{Definición~\ref{def-hip-impedancia} --- Impedancia y coeficiente de
reflexión}\vspace{3mm}

La cantidad \(Z = \sqrt{T\rho}\) se llama \textbf{impedancia} de la
cuerda. El número \(R\) de \ref{eq-hip-reflexion} es el
\textbf{coeficiente de reflexión}: la amplitud de la onda reflejada
dividida por la de la incidente.

\end{minipage}%
\end{tcolorbox}

\begin{tcolorbox}[enhanced jigsaw, arc=.35mm, bottomrule=.15mm, breakable, colback=white, colframe=quarto-callout-note-color-frame, left=2mm, leftrule=.75mm, opacityback=0, rightrule=.15mm, toprule=.15mm]
\begin{minipage}[t]{5.5mm}
\textcolor{quarto-callout-note-color}{\faInfo}
\end{minipage}%
\begin{minipage}[t]{\textwidth - 5.5mm}

\vspace{-3mm}\textbf{Teorema~\ref{thm-hip-finita} --- La cuerda fija por reflexiones}\vspace{3mm}

Sean \(\tilde\varphi\) y \(\tilde\psi\) las extensiones impares y
\(2l\)-periódicas de \(\varphi\) y \(\psi\). La solución de la cuerda
con extremos fijos,
\[u_{tt} = a^2u_{xx},\quad u(0,t) = u(l,t) = 0,\quad u(x,0) = \varphi(x),\quad u_t(x,0) = \psi(x),\]
es la fórmula de d'Alembert aplicada a \(\tilde\varphi\) y
\(\tilde\psi\):

\[
u(x,t) = \frac{\tilde\varphi(x+at)+\tilde\varphi(x-at)}{2} + \frac{1}{2a}\int_{x-at}^{x+at}\tilde\psi(s)\,ds .
\]

En particular, el movimiento es \textbf{periódico en el tiempo} con
período \(2l/a\).

\end{minipage}%
\end{tcolorbox}

\begin{tcolorbox}[enhanced jigsaw, arc=.35mm, bottomrule=.15mm, breakable, colback=white, colframe=quarto-callout-note-color-frame, left=2mm, leftrule=.75mm, opacityback=0, rightrule=.15mm, toprule=.15mm]
\begin{minipage}[t]{5.5mm}
\textcolor{quarto-callout-note-color}{\faInfo}
\end{minipage}%
\begin{minipage}[t]{\textwidth - 5.5mm}

\vspace{-3mm}\textbf{Proposición~\ref{prp-hip-extremo-movil} --- Cuerda semi-infinita con el
extremo movido}\vspace{3mm}

Con \(\varphi = \psi = 0\) ---la cuerda en reposo--- y \(\mu(0) = 0\),
la solución de
\[u_{tt} = a^2u_{xx}\ (x>0),\qquad u(0,t) = \mu(t),\qquad u(x,0) = u_t(x,0) = 0\]
es

\[
u(x,t) = \begin{cases} \mu\!\left(t - \dfrac{x}{a}\right), & x < at,\\[2mm] 0, & x\ge at.\end{cases}
\]

Con datos iniciales no nulos, se suma esta función a la solución de
\ref{eq-hip-semi-fijo}.

\end{minipage}%
\end{tcolorbox}

\subsubsection{Fuentes, Green y
resonancia}\label{fuentes-green-y-resonancia}

\begin{tcolorbox}[enhanced jigsaw, arc=.35mm, bottomrule=.15mm, breakable, colback=white, colframe=quarto-callout-note-color-frame, left=2mm, leftrule=.75mm, opacityback=0, rightrule=.15mm, toprule=.15mm]
\begin{minipage}[t]{5.5mm}
\textcolor{quarto-callout-note-color}{\faInfo}
\end{minipage}%
\begin{minipage}[t]{\textwidth - 5.5mm}

\vspace{-3mm}\textbf{Teorema~\ref{thm-hip-dalembert-fuente} --- Fórmula de d'Alembert con
fuente}\vspace{3mm}

Sean \(\varphi\in C^2(\mathbb R)\), \(\psi\in C^1(\mathbb R)\) y \(f\)
continua con \(f_x\) continua. La única solución \(C^2\) de
\[u_{tt} = a^2u_{xx} + f(x,t),\quad x\in\mathbb R,\ t>0;\qquad u(x,0)=\varphi(x),\quad u_t(x,0)=\psi(x)\]
es

\[
u(x,t) = \frac{\varphi(x+at)+\varphi(x-at)}{2} + \frac{1}{2a}\int_{x-at}^{x+at}\psi(s)\,ds
+ \frac{1}{2a}\iint_{\Delta(x,t)} f(s,\tau)\,ds\,d\tau ,
\]

donde \(\Delta(x,t)\) es el \textbf{triángulo característico} del punto:
el que tiene vértice en \((x,t)\) y base \([x-at,\,x+at]\) sobre el eje
\(\tau=0\). Escrita como integral iterada, \[
\iint_{\Delta(x,t)} f\,ds\,d\tau = \int_0^t\int_{x-a(t-\tau)}^{x+a(t-\tau)} f(s,\tau)\,ds\,d\tau .
\]

\end{minipage}%
\end{tcolorbox}

\begin{tcolorbox}[enhanced jigsaw, arc=.35mm, bottomrule=.15mm, breakable, colback=white, colframe=quarto-callout-note-color-frame, left=2mm, leftrule=.75mm, opacityback=0, rightrule=.15mm, toprule=.15mm]
\begin{minipage}[t]{5.5mm}
\textcolor{quarto-callout-note-color}{\faInfo}
\end{minipage}%
\begin{minipage}[t]{\textwidth - 5.5mm}

\vspace{-3mm}\textbf{Definición~\ref{def-hip-green-1d} --- Función de Green de la ecuación de
onda en una dimensión}\vspace{3mm}

La solución de \(u_{tt} - a^2u_{xx} = \delta(x-s)\,\delta(t-\tau)\) con
la cuerda en reposo antes de \(\tau\) es

\[
G(x,t;s,\tau) = \frac{1}{2a}\,H\!\big(a(t-\tau) - |x-s|\big),
\]

donde \(H\) es la función de Heaviside. Con ella, la contribución de la
fuente en \ref{eq-hip-dalembert-fuente} es \[
\int_0^\infty\!\!\int_{-\infty}^{\infty} G(x,t;s,\tau)\,f(s,\tau)\,ds\,d\tau .
\]

\end{minipage}%
\end{tcolorbox}

\begin{tcolorbox}[enhanced jigsaw, arc=.35mm, bottomrule=.15mm, breakable, colback=white, colframe=quarto-callout-note-color-frame, left=2mm, leftrule=.75mm, opacityback=0, rightrule=.15mm, toprule=.15mm]
\begin{minipage}[t]{5.5mm}
\textcolor{quarto-callout-note-color}{\faInfo}
\end{minipage}%
\begin{minipage}[t]{\textwidth - 5.5mm}

\vspace{-3mm}\textbf{Definición~\ref{def-hip-resonancia} --- Resonancia}\vspace{3mm}

Se dice que hay \textbf{resonancia} en el modo \(n\) cuando la
frecuencia del forzante coincide con su frecuencia propia,
\(\Omega = \omega_n\), y el coeficiente \(g_n\) no es cero. En ese caso

\[
T_n(t) = \frac{g_n}{2\omega_n^2}\big(\operatorname{sen}\omega_nt - \omega_nt\cos\omega_nt\big),
\]

cuya amplitud \textbf{crece linealmente con el tiempo} sin cota.

\end{minipage}%
\end{tcolorbox}

\begin{tcolorbox}[enhanced jigsaw, arc=.35mm, bottomrule=.15mm, breakable, colback=white, colframe=quarto-callout-note-color-frame, left=2mm, leftrule=.75mm, opacityback=0, rightrule=.15mm, toprule=.15mm]
\begin{minipage}[t]{5.5mm}
\textcolor{quarto-callout-note-color}{\faInfo}
\end{minipage}%
\begin{minipage}[t]{\textwidth - 5.5mm}

\vspace{-3mm}\textbf{Proposición~\ref{prp-hip-bordes-moviles} --- Reducción a bordes fijos}\vspace{3mm}

Sea \[
U(x,t) = \mu_1(t)\Big(1 - \frac{x}{l}\Big) + \mu_2(t)\,\frac{x}{l},
\] la recta que une los dos extremos en cada instante. Entonces
\(w = u - U\) cumple

\[
\begin{aligned}
&w_{tt} = a^2w_{xx} + f(x,t) - U_{tt}(x,t), && w(0,t) = w(l,t) = 0,\\
&w(x,0) = \varphi(x) - U(x,0), && w_t(x,0) = \psi(x) - U_t(x,0).
\end{aligned}
\]

\end{minipage}%
\end{tcolorbox}

\subsubsection{La membrana circular}\label{la-membrana-circular}

Hace falta para los problemas 9.11 y 9.12. La teoría completa de Bessel
es del capítulo 12.

\begin{tcolorbox}[enhanced jigsaw, arc=.35mm, bottomrule=.15mm, breakable, colback=white, colframe=quarto-callout-note-color-frame, left=2mm, leftrule=.75mm, opacityback=0, rightrule=.15mm, toprule=.15mm]
\begin{minipage}[t]{5.5mm}
\textcolor{quarto-callout-note-color}{\faInfo}
\end{minipage}%
\begin{minipage}[t]{\textwidth - 5.5mm}

\vspace{-3mm}\textbf{Definición~\ref{def-hip-bessel-j0} --- Función de Bessel de orden cero}\vspace{3mm}

Se llama \textbf{función de Bessel de primera especie de orden cero} a

\[
J_0(s) = \sum_{j=0}^{\infty}\frac{(-1)^j}{(j!)^2}\Big(\frac s2\Big)^{2j}
= 1 - \frac{s^2}{4} + \frac{s^4}{64} - \frac{s^6}{2304} + \cdots
\]

La serie converge para todo \(s\). \(J_0\) es la única solución de
\(R'' + R'/s + R = 0\), salvo múltiplos, que es \textbf{acotada en
\(s=0\)}; la otra solución independiente, \(Y_0\), tiende a \(-\infty\)
como \(\tfrac2\pi\ln s\). Se usa también \(J_1 = -J_0'\).

\end{minipage}%
\end{tcolorbox}

\begin{tcolorbox}[enhanced jigsaw, arc=.35mm, bottomrule=.15mm, breakable, colback=white, colframe=quarto-callout-note-color-frame, left=2mm, leftrule=.75mm, opacityback=0, rightrule=.15mm, toprule=.15mm]
\begin{minipage}[t]{5.5mm}
\textcolor{quarto-callout-note-color}{\faInfo}
\end{minipage}%
\begin{minipage}[t]{\textwidth - 5.5mm}

\vspace{-3mm}\textbf{Teorema~\ref{thm-hip-circular} --- Vibraciones radiales de la membrana
circular}\vspace{3mm}

La solución de \ref{eq-hip-circular}, acotada en \(r=0\), con
\(u(r,0) = \varphi(r)\) y \(u_t(r,0) = \psi(r)\), es

\[
u(r,t) = \sum_{p\ge1}\big(A_p\cos\omega_pt + B_p\operatorname{sen}\omega_pt\big)\,J_0\!\Big(\frac{j_{0,p}\,r}{R_0}\Big),
\qquad \omega_p = \frac{a\,j_{0,p}}{R_0},
\]

con

\[
A_p = \frac{2}{R_0^2J_1(j_{0,p})^2}\int_0^{R_0}\varphi(r)\,J_0\!\Big(\frac{j_{0,p}\,r}{R_0}\Big)\,r\,dr,
\qquad
B_p = \frac{2}{a\,j_{0,p}\,R_0\,J_1(j_{0,p})^2}\int_0^{R_0}\psi(r)\,J_0\!\Big(\frac{j_{0,p}\,r}{R_0}\Big)\,r\,dr .
\]

\end{minipage}%
\end{tcolorbox}

\subsubsection{Pérdidas y dispersión}\label{puxe9rdidas-y-dispersiuxf3n}

\begin{tcolorbox}[enhanced jigsaw, arc=.35mm, bottomrule=.15mm, breakable, colback=white, colframe=quarto-callout-note-color-frame, left=2mm, leftrule=.75mm, opacityback=0, rightrule=.15mm, toprule=.15mm]
\begin{minipage}[t]{5.5mm}
\textcolor{quarto-callout-note-color}{\faInfo}
\end{minipage}%
\begin{minipage}[t]{\textwidth - 5.5mm}

\vspace{-3mm}\textbf{Definición~\ref{def-hip-telegrafo} --- Ecuación del telégrafo}\vspace{3mm}

Se llama \textbf{ecuación del telégrafo} a \[
u_{tt} + (\alpha+\beta)\,u_t + \alpha\beta\,u = a^2\,u_{xx},
\] que es \ref{eq-hip-telegrafo} dividida por \(LC\), con
\(a = 1/\sqrt{LC}\), \(\alpha = R/L\) y \(\beta = G/C\).

\end{minipage}%
\end{tcolorbox}

\begin{tcolorbox}[enhanced jigsaw, arc=.35mm, bottomrule=.15mm, breakable, colback=white, colframe=quarto-callout-note-color-frame, left=2mm, leftrule=.75mm, opacityback=0, rightrule=.15mm, toprule=.15mm]
\begin{minipage}[t]{5.5mm}
\textcolor{quarto-callout-note-color}{\faInfo}
\end{minipage}%
\begin{minipage}[t]{\textwidth - 5.5mm}

\vspace{-3mm}\textbf{Definición~\ref{def-hip-dispersion} --- Relación de dispersión}\vspace{3mm}

La relación entre la frecuencia \(\omega\) y el número de onda \(k\) que
tiene que cumplir una onda \(e^{i(kx-\omega t)}\) para ser solución de
una ecuación lineal con coeficientes constantes se llama
\textbf{relación de dispersión}. La \textbf{velocidad de fase} es
\(\omega/k\). Una ecuación es \textbf{no dispersiva} si la velocidad de
fase es la misma para todos los \(k\), y \textbf{dispersiva} si no.

\end{minipage}%
\end{tcolorbox}

Con rozamiento, cada modo decae como \(e^{-\gamma t}\) ---eso justifica
el régimen permanente--- y la resonancia queda finita. Con \(R/L = G/C\)
(\textbf{condición de Heaviside}) el telégrafo no dispersa:
\(u = e^{-\alpha t}f(x-at)\).

\subsection{Formulario}\label{sec-hip-resumen-11}

\textbf{La cuerda}

{\def\LTcaptype{none} % do not increment counter
\begin{longtable}[]{@{}
  >{\raggedright\arraybackslash}p{(\linewidth - 2\tabcolsep) * \real{0.5000}}
  >{\raggedright\arraybackslash}p{(\linewidth - 2\tabcolsep) * \real{0.5000}}@{}}
\toprule\noalign{}
\endhead
\bottomrule\noalign{}
\endlastfoot
Ecuación & \(\rho\,u_{tt} = T\,u_{xx} + F\), \(\;a = \sqrt{T/\rho}\) \\
Energía &
\(E = \int_0^l\big(\tfrac12\rho u_t^2 + \tfrac12 T u_x^2\big)dx\) \\
Balance & \(\dot E = \int F u_t\,dx + \big[T u_x u_t\big]_0^l\) \\
Extremo fijo / libre / resorte & \(u = 0\) / \(u_x = 0\) /
\(T u_x = \kappa u\) en \(x=0\), \(\;T u_x = -\kappa u\) en \(x=l\) \\
Modos, extremos fijos & \(X_n = \operatorname{sen}\frac{n\pi x}{l}\),
\(\;\omega_n = \frac{n\pi a}{l}\) \\
Modos, fijo--libre & \(X_n = \operatorname{sen}\frac{(2n+1)\pi x}{2l}\),
\(\;\omega_n = \frac{(2n+1)\pi a}{2l}\) \\
Modos, libre--libre & \(X_n = \cos\frac{n\pi x}{l}\), con \(n = 0\)
(traslación rígida) \\
Coeficientes &
\(A_n = \frac2l\int_0^l\varphi\operatorname{sen}\frac{n\pi x}{l}dx\),
\(\;\omega_nB_n = \frac2l\int_0^l\psi\operatorname{sen}\frac{n\pi x}{l}dx\) \\
Período & \(2l/a\) \\
\end{longtable}
}

\textbf{d'Alembert y compañía}

{\def\LTcaptype{none} % do not increment counter
\begin{longtable}[]{@{}
  >{\raggedright\arraybackslash}p{(\linewidth - 2\tabcolsep) * \real{0.5000}}
  >{\raggedright\arraybackslash}p{(\linewidth - 2\tabcolsep) * \real{0.5000}}@{}}
\toprule\noalign{}
\endhead
\bottomrule\noalign{}
\endlastfoot
Cauchy &
\(u = \frac{\varphi(x+at)+\varphi(x-at)}{2} + \frac1{2a}\int_{x-at}^{x+at}\psi\) \\
Con fuente &
\(+\ \frac{1}{2a}\int_0^t\int_{x-a(t-\tau)}^{x+a(t-\tau)}f(s,\tau)\,ds\,d\tau\) \\
Estabilidad & \(\lvert u_1-u_2\rvert \le \delta(1+t_0)\) \\
Paralelogramo & \(u(A)+u(C) = u(B)+u(D)\) \\
Extremo fijo, \(x<at\) &
\(\frac{\varphi(x+at)-\varphi(at-x)}{2} + \frac1{2a}\int_{at-x}^{x+at}\psi\) \\
Extremo movido & \(u = \mu(t - x/a)\) para \(x<at\) \\
Reflexión & \(R = \frac{Z-\beta}{Z+\beta}\), \(\;Z = \sqrt{T\rho}\) \\
Oscilador forzado &
\(T_n = \frac1{\omega_n}\int_0^tf_n(\tau)\operatorname{sen}\omega_n(t-\tau)\,d\tau\) \\
Resonancia &
\(T_n = \frac{g_n}{2\omega_n^2}\big(\operatorname{sen}\omega_nt - \omega_nt\cos\omega_nt\big)\) \\
\end{longtable}
}

\textbf{Dos y tres dimensiones}

{\def\LTcaptype{none} % do not increment counter
\begin{longtable}[]{@{}
  >{\raggedright\arraybackslash}p{(\linewidth - 2\tabcolsep) * \real{0.5000}}
  >{\raggedright\arraybackslash}p{(\linewidth - 2\tabcolsep) * \real{0.5000}}@{}}
\toprule\noalign{}
\endhead
\bottomrule\noalign{}
\endlastfoot
Membrana rectangular &
\(\omega_{mn} = a\pi\sqrt{m^2/l_1^2 + n^2/l_2^2}\),
\(\;A_{mn} = \frac{4}{l_1l_2}\iint\varphi\operatorname{sen}\operatorname{sen}\) \\
Membrana circular radial & \(\omega_p = a\,j_{0,p}/R_0\),
\(\;j_{0,1}\approx2{,}405\), \(j_{0,2}\approx5{,}520\),
\(j_{0,3}\approx8{,}654\) \\
Coeficientes de Bessel &
\(A_p = \frac{2}{R_0^2J_1(j_{0,p})^2}\int_0^{R_0}\varphi\,J_0\big(\tfrac{j_{0,p}r}{R_0}\big)r\,dr\) \\
Laplaciano radial en 3D & \(\nabla^2u = \frac1r(ru)_{rr}\) \\
Ondas esféricas & \(u = \frac{F(t-r/a)}{r} + \frac{G(t+r/a)}{r}\) \\
Green 1D / 2D / 3D & \(\frac{H(at-\lvert x\rvert)}{2a}\) /
\(\frac{H(at-r)}{2\pi a\sqrt{a^2t^2-r^2}}\) /
\(\frac{\delta(t-r/a)}{4\pi a^2 r}\) \\
Kirchhoff (Cauchy 3D) &
\(u = \partial_t\big(tM_{at}[\varphi]\big) + tM_{at}[\psi]\) \\
\end{longtable}
}

\textbf{Electromagnetismo (SI)}

{\def\LTcaptype{none} % do not increment counter
\begin{longtable}[]{@{}
  >{\raggedright\arraybackslash}p{(\linewidth - 2\tabcolsep) * \real{0.5000}}
  >{\raggedright\arraybackslash}p{(\linewidth - 2\tabcolsep) * \real{0.5000}}@{}}
\toprule\noalign{}
\endhead
\bottomrule\noalign{}
\endlastfoot
Velocidad & \(c = 1/\sqrt{\mu_0\varepsilon_0}\) \\
Potenciales & \(\mathbf B = \nabla\times\mathbf A\),
\(\;\mathbf E = -\nabla\Phi - \partial_t\mathbf A\) \\
Gauge de Lorenz &
\(\nabla\cdot\mathbf A + \frac1{c^2}\partial_t\Phi = 0\) \\
Retardados &
\(\Phi = \frac{1}{4\pi\varepsilon_0}\int\frac{\rho(\mathbf x',\,t-R/c)}{R}dV'\) \\
Green de Helmholtz & \((\nabla^2+k^2)G = -\delta\),
\(\;G = \frac{e^{ikR}}{4\pi R}\) \\
Kirchhoff &
\(\psi(\mathbf x) = \oint\big(G\,\partial_{n'}\psi - \psi\,\partial_{n'}G\big)dS'\) \\
\end{longtable}
}

\subsection{Me dan un problema de ondas: qué
hago}\label{sec-hip-resumen-12}

\begin{enumerate}
\def\labelenumi{\arabic{enumi}.}
\tightlist
\item
  \textbf{Escribir la ecuación y los datos.} Identificar \(a\), la
  fuerza por unidad de masa \(f\), las dos condiciones iniciales y una
  condición por cada borde. Traducir cada borde físico con la
  Tabla~\ref{tbl-hip-bordes}: fijo, libre (anillo), resorte,
  amortiguador, movido.
\item
  \textbf{Mirar la geometría.}

  \begin{itemize}
  \tightlist
  \item
    Sin bordes, en una dimensión → \textbf{d'Alembert}, con fuente si la
    hay.
  \item
    Un borde → \textbf{extensión} impar (fijo) o par (libre), y
    d'Alembert.
  \item
    Dos bordes → \textbf{modos}; o reflexiones si la pregunta es
    \emph{cuándo} llega algo.
  \item
    Rectángulo → senos dobles. Círculo con simetría radial → \(J_0\).
  \item
    Tres dimensiones con simetría esférica → \(v = ru\). Sin simetría →
    Kirchhoff.
  \end{itemize}
\item
  \textbf{Si los bordes no son homogéneos}, restar una función que los
  cumpla (la recta entre los extremos, o una solución estacionaria) y
  seguir con bordes fijos y una fuente nueva.
\item
  \textbf{Si hay fuerza}, partir \(u = v + w\) y resolver un oscilador
  forzado por modo. Revisar si alguna frecuencia del forzante coincide
  con una \(\omega_n\) \textbf{y} la fuerza tiene componente sobre ese
  modo: eso es resonancia.
\item
  \textbf{Calcular los coeficientes} con la ortogonalidad: senos (peso
  \(1\)), cosenos (con el modo cero aparte) o Bessel (peso \(r\)).
  Acordarse del \(\omega_n\) que acompaña a los coeficientes de
  \(\psi\).
\item
  \textbf{Controlar.} En \(t=0\) la serie tiene que reproducir los
  datos; los modos que tienen un nodo donde se aplica la perturbación no
  tienen que aparecer; y la velocidad de decaimiento de los coeficientes
  tiene que coincidir con la regularidad del dato: \(1/n\) si hay un
  golpe puntual, \(1/n^2\) si hay una esquina, \(1/n^3\) si además la
  pendiente es continua.
\item
  \textbf{Si la pregunta es de existencia o de unicidad}, no hace falta
  resolver: \textbf{energía}.
\end{enumerate}

\section{Preguntas guía para el final}\label{sec-hip-preguntas}

No son ejercicios de cálculo: son las preguntas que conviene poder
contestar hablando, sin papel. Cada una tiene desplegada la idea de lo
que se espera.

\begin{tcolorbox}[enhanced jigsaw, arc=.35mm, bottomrule=.15mm, bottomtitle=1mm, breakable, colback=white, colbacktitle=quarto-callout-tip-color!10!white, colframe=quarto-callout-tip-color-frame, coltitle=black, left=2mm, leftrule=.75mm, opacityback=0, opacitybacktitle=0.6, rightrule=.15mm, title=\textcolor{quarto-callout-tip-color}{\faLightbulb}\hspace{0.5em}{1. ¿Qué tres propiedades distinguen a la ecuación de onda, y cómo se ven
en la regla del trueno?}, titlerule=0mm, toprule=.15mm, toptitle=1mm]

Velocidad finita (hay un instante de llegada), forma que se conserva (el
frente llega nítido, porque la ecuación no es dispersiva) y
superposición (dos sonidos se suman sin mezclarse, porque es lineal).
Ver sección~\ref{sec-hip-motivacion}.

\end{tcolorbox}

\begin{tcolorbox}[enhanced jigsaw, arc=.35mm, bottomrule=.15mm, bottomtitle=1mm, breakable, colback=white, colbacktitle=quarto-callout-tip-color!10!white, colframe=quarto-callout-tip-color-frame, coltitle=black, left=2mm, leftrule=.75mm, opacityback=0, opacitybacktitle=0.6, rightrule=.15mm, title=\textcolor{quarto-callout-tip-color}{\faLightbulb}\hspace{0.5em}{2. Deducí la ecuación de la cuerda. ¿Qué se supone y qué pasa si falla
cada hipótesis?}, titlerule=0mm, toprule=.15mm, toptitle=1mm]

Newton sobre un pedacito: la fuerza vertical neta es la diferencia de
\(T\operatorname{sen}\theta\approx Tu_x\) en los dos extremos, así que
la aceleración es proporcional a la curvatura. Se supone movimiento solo
vertical, cuerda perfectamente flexible, pendientes chicas y tensión
uniforme. Sin pendientes chicas la ecuación es no lineal; con rigidez
aparece un término de cuarto orden, como en la viga. Ver
sección~\ref{sec-hip-cuerda}.

\end{tcolorbox}

\begin{tcolorbox}[enhanced jigsaw, arc=.35mm, bottomrule=.15mm, bottomtitle=1mm, breakable, colback=white, colbacktitle=quarto-callout-tip-color!10!white, colframe=quarto-callout-tip-color-frame, coltitle=black, left=2mm, leftrule=.75mm, opacityback=0, opacitybacktitle=0.6, rightrule=.15mm, title=\textcolor{quarto-callout-tip-color}{\faLightbulb}\hspace{0.5em}{3. ¿Qué es la velocidad \(a\) y de qué depende en cada medio?}, titlerule=0mm, toprule=.15mm, toptitle=1mm]

Es \(\sqrt{\text{rigidez}/\text{inercia}}\): \(\sqrt{T/\rho}\) en la
cuerda, \(\sqrt{E/\rho}\) en la barra, \(\sqrt{\gamma p_0/\rho_0}\) en
el aire, \(1/\sqrt{LC}\) en un cable, \(1/\sqrt{\varepsilon_0\mu_0}\)
para la luz. Ver Tabla~\ref{tbl-hip-analogias}.

\end{tcolorbox}

\begin{tcolorbox}[enhanced jigsaw, arc=.35mm, bottomrule=.15mm, bottomtitle=1mm, breakable, colback=white, colbacktitle=quarto-callout-tip-color!10!white, colframe=quarto-callout-tip-color-frame, coltitle=black, left=2mm, leftrule=.75mm, opacityback=0, opacitybacktitle=0.6, rightrule=.15mm, title=\textcolor{quarto-callout-tip-color}{\faLightbulb}\hspace{0.5em}{4. ¿Qué datos necesita el problema de la cuerda, y por qué dos
condiciones iniciales?}, titlerule=0mm, toprule=.15mm, toptitle=1mm]

Forma y velocidad iniciales en cada punto, y una condición en cada
extremo. Dos iniciales porque la ecuación es de segundo orden en el
tiempo, como en la mecánica de una partícula. Las de borde las decide la
física del extremo: fijo, libre (anillo en una varilla), con resorte.
Ver Definición~\ref{def-hip-problema-mixto} y
Tabla~\ref{tbl-hip-bordes}.

\end{tcolorbox}

\begin{tcolorbox}[enhanced jigsaw, arc=.35mm, bottomrule=.15mm, bottomtitle=1mm, breakable, colback=white, colbacktitle=quarto-callout-tip-color!10!white, colframe=quarto-callout-tip-color-frame, coltitle=black, left=2mm, leftrule=.75mm, opacityback=0, opacitybacktitle=0.6, rightrule=.15mm, title=\textcolor{quarto-callout-tip-color}{\faLightbulb}\hspace{0.5em}{5. ¿Cuáles son los tres casos límite y por qué el tercero necesita
rozamiento?}, titlerule=0mm, toprule=.15mm, toptitle=1mm]

Problema de Cauchy (tiempos cortos, lejos de los bordes), cuerda
semi-infinita (cerca de un borde) y régimen permanente (tiempos largos).
El tercero supone que el sistema olvidó los datos iniciales, pero la
energía de una cuerda ideal se conserva y las oscilaciones libres no se
apagan nunca: hace falta disipación. Ver
sección~\ref{sec-hip-casos-limite} y sección~\ref{sec-hip-perdidas}.

\end{tcolorbox}

\begin{tcolorbox}[enhanced jigsaw, arc=.35mm, bottomrule=.15mm, bottomtitle=1mm, breakable, colback=white, colbacktitle=quarto-callout-tip-color!10!white, colframe=quarto-callout-tip-color-frame, coltitle=black, left=2mm, leftrule=.75mm, opacityback=0, opacitybacktitle=0.6, rightrule=.15mm, title=\textcolor{quarto-callout-tip-color}{\faLightbulb}\hspace{0.5em}{6. ¿Qué es la energía de una cuerda y cuándo se conserva?}, titlerule=0mm, toprule=.15mm, toptitle=1mm]

\(E = \int(\tfrac12\rho u_t^2 + \tfrac12ku_x^2)\,dx\), cinética más
potencial elástica. Su derivada es la potencia de la fuerza externa más
la potencia que entra por los bordes, \([ku_xu_t]_0^l\). Sin fuerza y
con extremos fijos o libres, se conserva. Ver
Proposición~\ref{prp-hip-balance}.

\end{tcolorbox}

\begin{tcolorbox}[enhanced jigsaw, arc=.35mm, bottomrule=.15mm, bottomtitle=1mm, breakable, colback=white, colbacktitle=quarto-callout-tip-color!10!white, colframe=quarto-callout-tip-color-frame, coltitle=black, left=2mm, leftrule=.75mm, opacityback=0, opacitybacktitle=0.6, rightrule=.15mm, title=\textcolor{quarto-callout-tip-color}{\faLightbulb}\hspace{0.5em}{7. Demostrá el teorema de unicidad 9.1.1. ¿Dónde se usa cada hipótesis?}, titlerule=0mm, toprule=.15mm, toptitle=1mm]

La diferencia de dos soluciones cumple el problema homogéneo con datos
nulos (linealidad); su energía se conserva (regularidad \(C^2\),
extremos fijos) y vale cero al principio; el integrando es no negativo y
se anula, y como \(\rho,k>0\) eso obliga a \(w_t = w_x = 0\). Ver
Teorema~\ref{thm-hip-unicidad} y Tabla~\ref{tbl-hip-unicidad}.

\end{tcolorbox}

\begin{tcolorbox}[enhanced jigsaw, arc=.35mm, bottomrule=.15mm, bottomtitle=1mm, breakable, colback=white, colbacktitle=quarto-callout-tip-color!10!white, colframe=quarto-callout-tip-color-frame, coltitle=black, left=2mm, leftrule=.75mm, opacityback=0, opacitybacktitle=0.6, rightrule=.15mm, title=\textcolor{quarto-callout-tip-color}{\faLightbulb}\hspace{0.5em}{8. Leé la fórmula de d'Alembert término por término}, titlerule=0mm, toprule=.15mm, toptitle=1mm]

El primer término parte la forma inicial en dos mitades que viajan en
sentidos opuestos; el segundo acumula la velocidad inicial sobre el
dominio de dependencia \([x-at,x+at]\). Cuerda pulsada contra cuerda
golpeada. Ver Teorema~\ref{thm-hip-dalembert}.

\end{tcolorbox}

\begin{tcolorbox}[enhanced jigsaw, arc=.35mm, bottomrule=.15mm, bottomtitle=1mm, breakable, colback=white, colbacktitle=quarto-callout-tip-color!10!white, colframe=quarto-callout-tip-color-frame, coltitle=black, left=2mm, leftrule=.75mm, opacityback=0, opacitybacktitle=0.6, rightrule=.15mm, title=\textcolor{quarto-callout-tip-color}{\faLightbulb}\hspace{0.5em}{9. ¿Por dónde viajan las singularidades, y por qué importa?}, titlerule=0mm, toprule=.15mm, toptitle=1mm]

Solo por las características que salen del punto singular del dato. Es
lo que hace que un frente llegue entero ---el trueno, la onda P--- y es
lo contrario de la ecuación del calor, que suaviza al instante. Ver
sección~\ref{sec-hip-dalembert}.

\end{tcolorbox}

\begin{tcolorbox}[enhanced jigsaw, arc=.35mm, bottomrule=.15mm, bottomtitle=1mm, breakable, colback=white, colbacktitle=quarto-callout-tip-color!10!white, colframe=quarto-callout-tip-color-frame, coltitle=black, left=2mm, leftrule=.75mm, opacityback=0, opacitybacktitle=0.6, rightrule=.15mm, title=\textcolor{quarto-callout-tip-color}{\faLightbulb}\hspace{0.5em}{10. Enunciá y demostrá la estabilidad. ¿Por qué la cota crece con el
tiempo?}, titlerule=0mm, toprule=.15mm, toptitle=1mm]

Si los datos difieren menos que \(\delta\), las soluciones difieren a lo
sumo \(\delta(1+t_0)\) hasta \(t_0\). Crece porque un error uniforme en
la velocidad desplaza la cuerda entera, y ese error crece como \(t\). Lo
importante es que es lineal y no exponencial, como en el ejemplo de
Hadamard. Ver Proposición~\ref{prp-hip-estabilidad}.

\end{tcolorbox}

\begin{tcolorbox}[enhanced jigsaw, arc=.35mm, bottomrule=.15mm, bottomtitle=1mm, breakable, colback=white, colbacktitle=quarto-callout-tip-color!10!white, colframe=quarto-callout-tip-color-frame, coltitle=black, left=2mm, leftrule=.75mm, opacityback=0, opacitybacktitle=0.6, rightrule=.15mm, title=\textcolor{quarto-callout-tip-color}{\faLightbulb}\hspace{0.5em}{11. ¿Qué dice la regla del paralelogramo y para qué sirve?}, titlerule=0mm, toprule=.15mm, toptitle=1mm]

En un paralelogramo con lados sobre características, los vértices
opuestos suman lo mismo. Sale de que \(u = f(\xi)+g(\eta)\). Sirve para
calcular \(u\) cerca de un borde o con datos sobre características
(Goursat). Ver Proposición~\ref{prp-hip-paralelogramo}.

\end{tcolorbox}

\begin{tcolorbox}[enhanced jigsaw, arc=.35mm, bottomrule=.15mm, bottomtitle=1mm, breakable, colback=white, colbacktitle=quarto-callout-tip-color!10!white, colframe=quarto-callout-tip-color-frame, coltitle=black, left=2mm, leftrule=.75mm, opacityback=0, opacitybacktitle=0.6, rightrule=.15mm, title=\textcolor{quarto-callout-tip-color}{\faLightbulb}\hspace{0.5em}{12. ¿Cómo se resuelve la cuerda semi-infinita, y qué le pasa a un pulso
en el extremo?}, titlerule=0mm, toprule=.15mm, toptitle=1mm]

Se extienden los datos del otro lado: en forma impar si el extremo es
fijo, par si es libre, y se aplica d'Alembert. En el extremo fijo el
pulso vuelve invertido; en el libre, derecho. Ver
sección~\ref{sec-hip-reflexiones}.

\end{tcolorbox}

\begin{tcolorbox}[enhanced jigsaw, arc=.35mm, bottomrule=.15mm, bottomtitle=1mm, breakable, colback=white, colbacktitle=quarto-callout-tip-color!10!white, colframe=quarto-callout-tip-color-frame, coltitle=black, left=2mm, leftrule=.75mm, opacityback=0, opacitybacktitle=0.6, rightrule=.15mm, title=\textcolor{quarto-callout-tip-color}{\faLightbulb}\hspace{0.5em}{13. ¿Qué es la impedancia y qué es una terminación adaptada?}, titlerule=0mm, toprule=.15mm, toptitle=1mm]

\(Z = \sqrt{T\rho}\). Un extremo con un amortiguador \(\beta\) refleja
con coeficiente \(R = (Z-\beta)/(Z+\beta)\). Con \(\beta = Z\) no
refleja nada: hace la misma fuerza que el resto de una cuerda infinita.
Es la resistencia de terminación de los cables. Ver
Definición~\ref{def-hip-impedancia}.

\end{tcolorbox}

\begin{tcolorbox}[enhanced jigsaw, arc=.35mm, bottomrule=.15mm, bottomtitle=1mm, breakable, colback=white, colbacktitle=quarto-callout-tip-color!10!white, colframe=quarto-callout-tip-color-frame, coltitle=black, left=2mm, leftrule=.75mm, opacityback=0, opacitybacktitle=0.6, rightrule=.15mm, title=\textcolor{quarto-callout-tip-color}{\faLightbulb}\hspace{0.5em}{14. ¿Por qué la cuerda finita se puede ver como rebotes o como modos, y
qué tiene que ver con el período?}, titlerule=0mm, toprule=.15mm, toptitle=1mm]

La solución es d'Alembert con la extensión impar y \(2l\)-periódica, que
rebota con período \(2l/a\). Al desarrollar esa extensión en serie de
senos, cada término es un modo. Como todo lo que tiene período \(2l/a\)
se escribe con frecuencias múltiplos de \(\pi a/l\), la cuerda tiene
armónicos afinados. Ver Teorema~\ref{thm-hip-finita}.

\end{tcolorbox}

\begin{tcolorbox}[enhanced jigsaw, arc=.35mm, bottomrule=.15mm, bottomtitle=1mm, breakable, colback=white, colbacktitle=quarto-callout-tip-color!10!white, colframe=quarto-callout-tip-color-frame, coltitle=black, left=2mm, leftrule=.75mm, opacityback=0, opacitybacktitle=0.6, rightrule=.15mm, title=\textcolor{quarto-callout-tip-color}{\faLightbulb}\hspace{0.5em}{15. ¿Qué dice el principio de Duhamel?}, titlerule=0mm, toprule=.15mm, toptitle=1mm]

Una fuerza continua es una sucesión de golpes; cada golpe es una
velocidad inicial, y la respuesta es la suma de las respuestas. En la
recta da d'Alembert con fuente: la integral de la fuerza sobre el
triángulo característico. Ver Teorema~\ref{thm-hip-dalembert-fuente}.

\end{tcolorbox}

\begin{tcolorbox}[enhanced jigsaw, arc=.35mm, bottomrule=.15mm, bottomtitle=1mm, breakable, colback=white, colbacktitle=quarto-callout-tip-color!10!white, colframe=quarto-callout-tip-color-frame, coltitle=black, left=2mm, leftrule=.75mm, opacityback=0, opacitybacktitle=0.6, rightrule=.15mm, title=\textcolor{quarto-callout-tip-color}{\faLightbulb}\hspace{0.5em}{16. ¿Cómo se resuelve la cuerda forzada con extremos fijos?}, titlerule=0mm, toprule=.15mm, toptitle=1mm]

\(u = v + w\): \(w\) las oscilaciones libres con los datos, \(v\) las
forzadas desde el reposo. Desarrollando \(v\) y la fuerza en modos, cada
coeficiente cumple la ecuación de un oscilador forzado, que se resuelve
con la función de Green del oscilador. Ver
Teorema~\ref{thm-hip-forzada}.

\end{tcolorbox}

\begin{tcolorbox}[enhanced jigsaw, arc=.35mm, bottomrule=.15mm, bottomtitle=1mm, breakable, colback=white, colbacktitle=quarto-callout-tip-color!10!white, colframe=quarto-callout-tip-color-frame, coltitle=black, left=2mm, leftrule=.75mm, opacityback=0, opacitybacktitle=0.6, rightrule=.15mm, title=\textcolor{quarto-callout-tip-color}{\faLightbulb}\hspace{0.5em}{17. ¿Qué es la resonancia, cuándo no aparece y qué la limita?}, titlerule=0mm, toprule=.15mm, toptitle=1mm]

Cuando la frecuencia del forzante coincide con una frecuencia propia y
la fuerza tiene componente sobre ese modo, la amplitud crece
linealmente. No aparece si la fuerza se aplica en un nodo del modo. Con
rozamiento la amplitud queda finita. Ver
Definición~\ref{def-hip-resonancia}.

\end{tcolorbox}

\begin{tcolorbox}[enhanced jigsaw, arc=.35mm, bottomrule=.15mm, bottomtitle=1mm, breakable, colback=white, colbacktitle=quarto-callout-tip-color!10!white, colframe=quarto-callout-tip-color-frame, coltitle=black, left=2mm, leftrule=.75mm, opacityback=0, opacitybacktitle=0.6, rightrule=.15mm, title=\textcolor{quarto-callout-tip-color}{\faLightbulb}\hspace{0.5em}{18. ¿Qué se hace si los extremos se mueven?}, titlerule=0mm, toprule=.15mm, toptitle=1mm]

Se resta una función que cumpla las condiciones de borde, por ejemplo la
recta entre los extremos. El resto tiene bordes fijos y una fuente
nueva: mover los bordes equivale a empujar la cuerda desde adentro. Ver
Proposición~\ref{prp-hip-bordes-moviles}.

\end{tcolorbox}

\begin{tcolorbox}[enhanced jigsaw, arc=.35mm, bottomrule=.15mm, bottomtitle=1mm, breakable, colback=white, colbacktitle=quarto-callout-tip-color!10!white, colframe=quarto-callout-tip-color-frame, coltitle=black, left=2mm, leftrule=.75mm, opacityback=0, opacitybacktitle=0.6, rightrule=.15mm, title=\textcolor{quarto-callout-tip-color}{\faLightbulb}\hspace{0.5em}{19. Membrana rectangular: frecuencias, líneas nodales, degeneración}, titlerule=0mm, toprule=.15mm, toptitle=1mm]

\(\omega_{mn} = a\pi\sqrt{m^2/l_1^2+n^2/l_2^2}\). Los nodos son líneas,
que parten la membrana en celdas que se mueven en contrafase. En el
cuadrado, \((m,n)\) y \((n,m)\) tienen la misma frecuencia y sus mezclas
son modos con líneas nodales nuevas, por la simetría. Ver
Teorema~\ref{thm-hip-rectangular} y
Definición~\ref{def-hip-degeneracion}.

\end{tcolorbox}

\begin{tcolorbox}[enhanced jigsaw, arc=.35mm, bottomrule=.15mm, bottomtitle=1mm, breakable, colback=white, colbacktitle=quarto-callout-tip-color!10!white, colframe=quarto-callout-tip-color-frame, coltitle=black, left=2mm, leftrule=.75mm, opacityback=0, opacitybacktitle=0.6, rightrule=.15mm, title=\textcolor{quarto-callout-tip-color}{\faLightbulb}\hspace{0.5em}{20. ¿Por qué un tambor no tiene una nota definida?}, titlerule=0mm, toprule=.15mm, toptitle=1mm]

Porque sus frecuencias no son múltiplos enteros de la fundamental:
tienen raíces cuadradas (rectángulo) o ceros de funciones de Bessel
(círculo). Ver Tabla~\ref{tbl-hip-cocientes}.

\end{tcolorbox}

\begin{tcolorbox}[enhanced jigsaw, arc=.35mm, bottomrule=.15mm, bottomtitle=1mm, breakable, colback=white, colbacktitle=quarto-callout-tip-color!10!white, colframe=quarto-callout-tip-color-frame, coltitle=black, left=2mm, leftrule=.75mm, opacityback=0, opacitybacktitle=0.6, rightrule=.15mm, title=\textcolor{quarto-callout-tip-color}{\faLightbulb}\hspace{0.5em}{21. ¿Por qué aparecen las funciones de Bessel en el tambor, y por qué
\(J_0\) y no \(Y_0\)?}, titlerule=0mm, toprule=.15mm, toptitle=1mm]

El laplaciano en polares tiene el término \(u_r/r\), y la parte radial
es la ecuación de Bessel. De sus dos soluciones, \(Y_0\) diverge en el
centro y el parche no puede tener un desplazamiento infinito; queda
\(J_0\). Las frecuencias salen de sus ceros. Ver
sección~\ref{sec-hip-circular}.

\end{tcolorbox}

\begin{tcolorbox}[enhanced jigsaw, arc=.35mm, bottomrule=.15mm, bottomtitle=1mm, breakable, colback=white, colbacktitle=quarto-callout-tip-color!10!white, colframe=quarto-callout-tip-color-frame, coltitle=black, left=2mm, leftrule=.75mm, opacityback=0, opacitybacktitle=0.6, rightrule=.15mm, title=\textcolor{quarto-callout-tip-color}{\faLightbulb}\hspace{0.5em}{22. Ondas esféricas: ¿por qué \(1/r\), y qué pasa en el centro?}, titlerule=0mm, toprule=.15mm, toptitle=1mm]

Con \(v = ru\) la ecuación radial es la de la cuerda. El \(1/r\) es la
conservación de la energía repartida en esferas cada vez más grandes. El
centro hace de extremo fijo: una onda convergente rebota ahí, y si el
dato tiene un salto, la solución en el centro tiene un pico (enfoque).
Ver sección~\ref{sec-hip-esfericas}.

\end{tcolorbox}

\begin{tcolorbox}[enhanced jigsaw, arc=.35mm, bottomrule=.15mm, bottomtitle=1mm, breakable, colback=white, colbacktitle=quarto-callout-tip-color!10!white, colframe=quarto-callout-tip-color-frame, coltitle=black, left=2mm, leftrule=.75mm, opacityback=0, opacitybacktitle=0.6, rightrule=.15mm, title=\textcolor{quarto-callout-tip-color}{\faLightbulb}\hspace{0.5em}{23. ¿Qué es el principio de Huygens y por qué falla en dos dimensiones?}, titlerule=0mm, toprule=.15mm, toptitle=1mm]

En tres dimensiones la solución en un punto depende de los datos solo
sobre una esfera, no sobre la bola: la perturbación pasa y no deja nada.
En dos dimensiones depende de todo el disco, y queda una estela. Las
funciones de Green: cáscara del cono en 3D, interior del cono en 1D y
2D. Ver Tabla~\ref{tbl-hip-green-dimensiones}.

\end{tcolorbox}

\begin{tcolorbox}[enhanced jigsaw, arc=.35mm, bottomrule=.15mm, bottomtitle=1mm, breakable, colback=white, colbacktitle=quarto-callout-tip-color!10!white, colframe=quarto-callout-tip-color-frame, coltitle=black, left=2mm, leftrule=.75mm, opacityback=0, opacitybacktitle=0.6, rightrule=.15mm, title=\textcolor{quarto-callout-tip-color}{\faLightbulb}\hspace{0.5em}{24. ¿Cómo salen las ondas electromagnéticas de las ecuaciones de
Maxwell?}, titlerule=0mm, toprule=.15mm, toptitle=1mm]

Rotor de Faraday, identidad del rotor del rotor, y Ampère--Maxwell sin
corrientes: cada componente cumple la ecuación de onda con
\(c = 1/\sqrt{\mu_0\varepsilon_0}\), que coincide con la velocidad de la
luz. Sin la corriente de desplazamiento no saldría. Ver
Proposición~\ref{prp-hip-ondas-em}.

\end{tcolorbox}

\begin{tcolorbox}[enhanced jigsaw, arc=.35mm, bottomrule=.15mm, bottomtitle=1mm, breakable, colback=white, colbacktitle=quarto-callout-tip-color!10!white, colframe=quarto-callout-tip-color-frame, coltitle=black, left=2mm, leftrule=.75mm, opacityback=0, opacitybacktitle=0.6, rightrule=.15mm, title=\textcolor{quarto-callout-tip-color}{\faLightbulb}\hspace{0.5em}{25. ¿Para qué sirven los potenciales y el gauge de Lorenz?}, titlerule=0mm, toprule=.15mm, toptitle=1mm]

Reducen las incógnitas; la libertad de gauge permite imponer
\(\nabla\cdot\mathbf A + \partial_t\Phi/c^2 = 0\), y con eso las
ecuaciones se desacoplan en cuatro ecuaciones de onda con fuente. Ver
Definición~\ref{def-hip-gauge}.

\end{tcolorbox}

\begin{tcolorbox}[enhanced jigsaw, arc=.35mm, bottomrule=.15mm, bottomtitle=1mm, breakable, colback=white, colbacktitle=quarto-callout-tip-color!10!white, colframe=quarto-callout-tip-color-frame, coltitle=black, left=2mm, leftrule=.75mm, opacityback=0, opacitybacktitle=0.6, rightrule=.15mm, title=\textcolor{quarto-callout-tip-color}{\faLightbulb}\hspace{0.5em}{26. ¿Qué son los potenciales retardados y qué decide que sean
retardados?}, titlerule=0mm, toprule=.15mm, toptitle=1mm]

Son Coulomb evaluada en el tiempo retardado \(t - R/c\). La cuenta con
Fourier y Helmholtz deja dos soluciones, retardada y avanzada; la
causalidad descarta la avanzada. Ver Teorema~\ref{thm-hip-retardados}.

\end{tcolorbox}

\begin{tcolorbox}[enhanced jigsaw, arc=.35mm, bottomrule=.15mm, bottomtitle=1mm, breakable, colback=white, colbacktitle=quarto-callout-tip-color!10!white, colframe=quarto-callout-tip-color-frame, coltitle=black, left=2mm, leftrule=.75mm, opacityback=0, opacitybacktitle=0.6, rightrule=.15mm, title=\textcolor{quarto-callout-tip-color}{\faLightbulb}\hspace{0.5em}{27. ¿Qué dice la fórmula de Kirchhoff y en qué falla la aproximación de
Kirchhoff?}, titlerule=0mm, toprule=.15mm, toptitle=1mm]

Con la segunda identidad de Green y la \(G\) de Helmholtz, el campo
adentro de una región queda escrito con su valor y su derivada normal en
el borde: el principio de Huygens formalizado. La aproximación para
difracción pone la onda incidente en el agujero y cero en la pantalla,
pero una solución con valor y derivada nulos en un pedazo de superficie
es nula: es inconsistente. Ver Teorema~\ref{thm-hip-kirchhoff}.

\end{tcolorbox}

\begin{tcolorbox}[enhanced jigsaw, arc=.35mm, bottomrule=.15mm, bottomtitle=1mm, breakable, colback=white, colbacktitle=quarto-callout-tip-color!10!white, colframe=quarto-callout-tip-color-frame, coltitle=black, left=2mm, leftrule=.75mm, opacityback=0, opacitybacktitle=0.6, rightrule=.15mm, title=\textcolor{quarto-callout-tip-color}{\faLightbulb}\hspace{0.5em}{28. ¿Qué cambia con pérdidas? ¿Qué es la dispersión y qué es la
condición de Heaviside?}, titlerule=0mm, toprule=.15mm, toptitle=1mm]

Las pérdidas apagan las oscilaciones libres y hacen finita la
resonancia; el tipo sigue siendo hiperbólico y el frente viaja a
velocidad \(a\). La dispersión es que la velocidad de fase dependa de la
longitud de onda, y desarma los pulsos. Con \(R/L = G/C\) el telégrafo
no dispersa: el pulso llega atenuado pero entero. Ver
sección~\ref{sec-hip-telegrafo}.

\end{tcolorbox}

\subsection*{Bibliografía del
capítulo}\label{bibliografuxeda-del-capuxedtulo}
\addcontentsline{toc}{subsection}{Bibliografía del capítulo}

\begin{itemize}
\tightlist
\item
  A. G. González, \emph{Métodos Matemáticos de la Física}, UNCPBA,
  capítulo 9. Es el libro de la cátedra.
\item
  W. A. Strauss, \emph{Partial Differential Equations: An Introduction},
  2.ª ed., Wiley, 2008. Capítulos 2 y 3 (d'Alembert, energía,
  reflexiones, Duhamel) y 9 (ondas en el espacio, Kirchhoff, Huygens).
  La mejor referencia para seguir este capítulo.
\item
  L. C. Evans, \emph{Partial Differential Equations}, 2.ª ed., AMS,
  2010. Sección 2.4, con las demostraciones completas de Kirchhoff y
  Poisson y el método de descenso.
\item
  E. Butkov, \emph{Mathematical Physics}, Addison-Wesley, 1968.
  Capítulos sobre la ecuación de onda y funciones de Bessel, con muchos
  problemas de membranas.
\item
  J. D. Jackson, \emph{Classical Electrodynamics}, 3.ª ed., Wiley, 1999.
  Capítulo 6 para potenciales, gauge y potenciales retardados; capítulo
  10 para difracción y la aproximación de Kirchhoff.
\item
  K. Y. Billah y R. H. Scanlan, «Resonance, Tacoma Narrows bridge
  failure, and undergraduate physics textbooks», \emph{American Journal
  of Physics} 59 (1991) 118--124.
\end{itemize}

\chapter{Ecuaciones parabólicas}\label{sec-parabolicas}

\begin{tcolorbox}[enhanced jigsaw, arc=.35mm, bottomrule=.15mm, breakable, colback=white, colframe=quarto-callout-note-color-frame, left=2mm, leftrule=.75mm, opacityback=0, rightrule=.15mm, toprule=.15mm]
\begin{minipage}[t]{5.5mm}
\textcolor{quarto-callout-note-color}{\faInfo}
\end{minipage}%
\begin{minipage}[t]{\textwidth - 5.5mm}

\vspace{-3mm}\textbf{Capítulo en construcción}\vspace{3mm}

Abajo está el temario. Después hay un \textbf{adelanto}, con un ejemplo
resuelto.

\end{minipage}%
\end{tcolorbox}

\section*{Temario}\label{temario-7}
\addcontentsline{toc}{section}{Temario}

\markright{Temario}

Propagación del calor. Ecuación de difusión. Condiciones de contorno.
Separación de variables. Función influencia o de fuente puntual.
Funciones de Green para las condiciones de contorno e iniciales.

\section*{Adelanto: por qué el calor borra los
detalles}\label{adelanto-por-quuxe9-el-calor-borra-los-detalles}
\addcontentsline{toc}{section}{Adelanto: por qué el calor borra los
detalles}

\markright{Adelanto: por qué el calor borra los detalles}

Una barra de longitud \(1\), aislada en los costados, con los dos
extremos mantenidos a temperatura cero. Si \(u(x,t)\) es la temperatura,
el problema es

\begin{equation}\protect\phantomsection\label{eq-calor}{
u_t = k\,u_{xx}, \qquad u(0,t) = u(1,t) = 0, \qquad u(x,0) = \varphi(x).
}\end{equation}

Separando variables, \(u = X(x)T(t)\), las condiciones de contorno dejan
como únicos modos posibles \(X_n(x) = \sin(n\pi x)\), y para cada uno la
parte temporal decae exponencialmente:

\begin{equation}\protect\phantomsection\label{eq-calor-sol}{
u(x,t) \;=\; \sum_{n=1}^{\infty} b_n\, e^{-(n\pi)^2 k t}\, \sin(n\pi x),
\qquad b_n = 2\int_0^1 \varphi(x)\,\sin(n\pi x)\,dx .
}\end{equation}

Los \(b_n\) son los coeficientes de la serie de Fourier en senos de la
condición inicial: es el Capítulo~\ref{sec-fourier} entrando por la
puerta de atrás. Lo nuevo es el factor \(e^{-(n\pi)^2kt}\), y ahí está
toda la física del problema.

Fijate en el exponente: \textbf{crece con el cuadrado del número de
modo}. El modo \(n=10\) se apaga \(100\) veces más rápido que el
\(n=1\). Y los detalles finos de la condición inicial ---una esquina, un
escalón, un pico--- son justamente los que necesitan modos altos para
representarse. El resultado es que la difusión \textbf{borra los
detalles antes que las formas grandes}, y lo hace de inmediato: apenas
\(t>0\), la serie de la \ref{eq-calor-sol} converge tan rápido que la
solución es infinitamente derivable, por más fea que fuera \(\varphi\).

\begin{figure}

\centering{

\pandocbounded{\includegraphics[keepaspectratio]{10-parabolicas_files/figure-pdf/fig-calor-output-1.pdf}}

}

\caption{\label{fig-calor}La \ref{eq-calor-sol} con \(k=1\) y una
condición inicial discontinua: un escalón que vale \(1\) en el interior
y \(0\) en los extremos. En \(t=0\) se ve el fenómeno de Gibbs del
Capítulo~\ref{sec-fourier}, porque la serie está truncada. Un instante
después no queda ni rastro: los modos altos, que son los que producían
las oscilaciones, ya se apagaron.}

\end{figure}%

Esta es la diferencia de fondo con las ecuaciones hiperbólicas del
Capítulo~\ref{sec-hiperbolicas}, donde una discontinuidad en el dato
inicial viaja intacta para siempre. También es la razón por la que el
problema del calor \textbf{hacia atrás en el tiempo} está mal planteado:
si los modos altos se apagan hacia adelante, hacia atrás explotan, y dos
condiciones iniciales muy distintas producen temperaturas de hoy
indistinguibles.

Video recomendado: Grant Sanderson (3Blue1Brown), \emph{Solving the heat
equation}, \url{https://www.youtube.com/watch?v=ToIXSwZ1pJU}.

Tomemos la condición inicial más simple posible,
\(\varphi(x) = \sin(\pi x)\), con \(k = 1\). La serie de la
\ref{eq-calor-sol} tiene un solo término y la solución es
\(u(x,t) = e^{-\pi^2 t}\sin(\pi x)\).

\textbf{¿En qué tiempo la temperatura del centro de la barra cae a la
mitad de su valor inicial?} Escribí el resultado con tres decimales.

Hay que resolver \(e^{-\pi^2 t} = \tfrac12\), o sea
\(t = \ln 2 / \pi^2 = 0{,}0702\).

Lo interesante es cómo cambia esto con el modo. Para
\(\varphi(x) = \sin(n\pi x)\) el mismo cálculo da
\(t_n = \ln 2/(n\pi)^2\): el modo \(n\) se apaga a la mitad \(n^2\)
veces más rápido. Para \(n=10\) ya estamos en \(t \approx 0{,}0007\).

\chapter{Ecuaciones elípticas}\label{sec-elipticas}

\begin{tcolorbox}[enhanced jigsaw, arc=.35mm, bottomrule=.15mm, breakable, colback=white, colframe=quarto-callout-note-color-frame, left=2mm, leftrule=.75mm, opacityback=0, rightrule=.15mm, toprule=.15mm]
\begin{minipage}[t]{5.5mm}
\textcolor{quarto-callout-note-color}{\faInfo}
\end{minipage}%
\begin{minipage}[t]{\textwidth - 5.5mm}

\vspace{-3mm}\textbf{Capítulo en construcción}\vspace{3mm}

Abajo está el temario. El contenido se agrega en orden de necesidad.

\end{minipage}%
\end{tcolorbox}

\section*{Temario}\label{temario-8}
\addcontentsline{toc}{section}{Temario}

\markright{Temario}

Ecuación de Laplace en coordenadas curvilíneas, soluciones particulares.
Funciones armónicas. Separación de variables. Integral de Poisson.
Función de la fuente. Teoría del potencial. Desarrollo en multipolos.
Teorema de Helmholtz. Transformación conforme. Ecuación biarmónica.

\chapter{Funciones especiales}\label{sec-especiales}

\begin{tcolorbox}[enhanced jigsaw, arc=.35mm, bottomrule=.15mm, breakable, colback=white, colframe=quarto-callout-note-color-frame, left=2mm, leftrule=.75mm, opacityback=0, rightrule=.15mm, toprule=.15mm]
\begin{minipage}[t]{5.5mm}
\textcolor{quarto-callout-note-color}{\faInfo}
\end{minipage}%
\begin{minipage}[t]{\textwidth - 5.5mm}

\vspace{-3mm}\textbf{Capítulo en construcción}\vspace{3mm}

Abajo está el temario. El contenido se agrega en orden de necesidad.

\end{minipage}%
\end{tcolorbox}

\section*{Temario}\label{temario-9}
\addcontentsline{toc}{section}{Temario}

\markright{Temario}

Funciones cilíndricas de Bessel, Hankel y Neumann. Funciones esféricas:
polinomios de Legendre y armónicos esféricos. Polinomios de Hermite y de
Laguerre. Funciones hipergeométrica e hipergeométrica confluente.
Funciones Gamma, error y relacionadas.

\chapter{Desarrollos asintóticos}\label{sec-asintoticos}

\begin{tcolorbox}[enhanced jigsaw, arc=.35mm, bottomrule=.15mm, breakable, colback=white, colframe=quarto-callout-note-color-frame, left=2mm, leftrule=.75mm, opacityback=0, rightrule=.15mm, toprule=.15mm]
\begin{minipage}[t]{5.5mm}
\textcolor{quarto-callout-note-color}{\faInfo}
\end{minipage}%
\begin{minipage}[t]{\textwidth - 5.5mm}

\vspace{-3mm}\textbf{Capítulo en construcción}\vspace{3mm}

Abajo está el temario. El contenido se agrega en orden de necesidad.

\end{minipage}%
\end{tcolorbox}

\section*{Temario}\label{temario-10}
\addcontentsline{toc}{section}{Temario}

\markright{Temario}

Expansión asintótica de funciones integrales. Método de la fase
estacionaria. Orden dominante en ecuaciones diferenciales. Problemas
singulares irregulares. Empalme de soluciones asintóticas.

\chapter{Ecuaciones integrales}\label{sec-integrales}

\begin{tcolorbox}[enhanced jigsaw, arc=.35mm, bottomrule=.15mm, breakable, colback=white, colframe=quarto-callout-note-color-frame, left=2mm, leftrule=.75mm, opacityback=0, rightrule=.15mm, toprule=.15mm]
\begin{minipage}[t]{5.5mm}
\textcolor{quarto-callout-note-color}{\faInfo}
\end{minipage}%
\begin{minipage}[t]{\textwidth - 5.5mm}

\vspace{-3mm}\textbf{Capítulo en construcción}\vspace{3mm}

Abajo está el temario. El contenido se agrega en orden de necesidad.

\end{minipage}%
\end{tcolorbox}

\section*{Temario}\label{temario-11}
\addcontentsline{toc}{section}{Temario}

\markright{Temario}

Ecuaciones del tipo de Fredholm y de Volterra. Ecuaciones integrales
resueltas por transformadas.

\bookmarksetup{startatroot}

\chapter*{Bibliografía}\label{bibliografuxeda}
\addcontentsline{toc}{chapter}{Bibliografía}

\markboth{Bibliografía}{Bibliografía}

Los textos que uso, y para qué sirve cada uno.

\textbf{González, A. G.} \emph{Métodos Matemáticos de la Física.}
UNCPBA. El programa de la cátedra y el orden de los temas siguen a este
texto.

\textbf{Boyce, W. E. y DiPrima, R. C.} \emph{Elementary Differential
Equations and Boundary Value Problems.} Wiley. El más accesible para
todo lo de ecuaciones ordinarias y problemas de contorno.

\textbf{Zachmanoglou, E. C. y Thoe, D. W.} \emph{Introduction to Partial
Differential Equations with Applications.} Dover. Clasificación,
características y formas canónicas. Es el que mejor explica el capítulo
8.

\textbf{Strauss, W. A.} \emph{Partial Differential Equations: An
Introduction.} Wiley. El mejor primer libro de EDP que conozco. Corto,
honesto sobre lo que no demuestra, y con capítulos separados para onda,
difusión y Laplace.

\textbf{Evans, L. C.} \emph{Partial Differential Equations.} American
Mathematical Society. La referencia de posgrado. El capítulo 3, sobre
ecuaciones no lineales de primer orden y bandas características, es el
mejor tratamiento del tema que existe.

\textbf{Dallos, D.} \emph{Notas de clase de Ecuaciones Diferenciales
II.} 2024. Notas de un curso de Licenciatura en Matemática. Cubre curvas
y superficies integrales de campos vectoriales antes de las EDP de
primer orden, que es el orden que hace que el método de las
características se entienda en vez de memorizarse. En español.

\textbf{Butkov, E.} \emph{Mathematical Physics.} Addison-Wesley. Cubre
casi todo el programa desde el lado del físico. Bueno para ver de dónde
sale cada herramienta.

\textbf{Stakgold, I.} \emph{Boundary Value Problems of Mathematical
Physics.} SIAM. La referencia seria para funciones de Green y
distribuciones.

\textbf{Jackson, J. D.} \emph{Classical Electrodynamics.} Wiley. La
referencia de electromagnetismo. Los capítulos 6 ---potenciales, gauge,
potenciales retardados--- y 10 ---difracción y la aproximación de
Kirchhoff--- son el complemento natural del capítulo 9.

\textbf{Billah, K. Y. y Scanlan, R. H.} «Resonance, Tacoma Narrows
bridge failure, and undergraduate physics textbooks.» \emph{American
Journal of Physics} 59 (1991). El artículo que explica por qué el puente
de Tacoma no se cayó por resonancia.

\textbf{Tijonov, A. N. y Samarsky, A. A.} \emph{Ecuaciones de la Física
Matemática.} Mir. Clásico, con muchísimos problemas físicos resueltos en
detalle.

\textbf{Bender, C. M. y Orszag, S. A.} \emph{Advanced Mathematical
Methods for Scientists and Engineers.} Springer. Para desarrollos
asintóticos. No tiene rival en ese tema.

\textbf{Lighthill, M. J.} \emph{An Introduction to Fourier Analysis and
Generalised Functions.} Cambridge University Press. Corto y luminoso. La
mejor entrada a distribuciones.

\textbf{Courant, R. y Hilbert, D.} \emph{Methods of Mathematical
Physics.} Wiley. La obra de referencia. No para aprender de cero, sí
para consultar.

\textbf{Spiegel, M., Lipschutz, S. y Liu, J.} \emph{Mathematical
Handbook of Formulas and Tables.} McGraw-Hill. Tablas de transformadas y
de funciones especiales.


\backmatter


\end{document}
